\documentclass[11pt]{article}
\usepackage[a4paper,margin=28mm]{geometry}
\usepackage{amsmath,amssymb,amsthm}
\usepackage[hidelinks]{hyperref}
\hypersetup{
 pdftitle={Exact Zero Detection and Compact Arithmetic Realizations for the Riemann Xi Function},
 pdfauthor={Leonard van Hemert},
 pdfsubject={Weil forms, zero detection, and complete arithmetic realizations}
}
\setcounter{tocdepth}{2}
\newtheorem{theorem}{Theorem}
\newtheorem{lemma}[theorem]{Lemma}
\newtheorem{corollary}[theorem]{Corollary}
\newtheorem{proposition}[theorem]{Proposition}
\theoremstyle{remark}\newtheorem{remark}[theorem]{Remark}
\newcommand{\R}{\mathbb R}
\newcommand{\C}{\mathbb C}
\newcommand{\Z}{\mathbb Z}
\newcommand{\dd}{\,\mathrm d}
\newcommand{\sinhc}{\operatorname{sinhc}}
\newcommand{\dist}{\mathrm{dist}}
\title{Exact Zero Detection and Compact Arithmetic Realizations\\
for the Riemann Xi Function}
\author{Leonard van Hemert}
\date{11 October 2026}
\begin{document}
\maketitle
\begin{abstract}
We prove exact zero-detection and arithmetic realization theorems for
the Riemann xi function, without assuming the Riemann Hypothesis or
simplicity of its zeros. The exhausted negative index of the complete
semilocal Weil form is \(2K_+\), where \(K_+\) counts distinct off-line
mirror pairs at positive height and may be infinite. For the
logistic-smoothed zero kernel, the same index is attained on every
nonempty bounded open interval, after any fixed finite collection
of exponential moment and jet conditions. We give an explicit
norm cost for increasing moment depth.

On the complete closed arithmetic range, we construct a compact,
injective operator with dense range and spectrum
\[
 \{0,1\}\ \cup\
 \left\{\frac{p^2-1/4}{p^2+z^2}:\Xi(z)=0\right\},
 \qquad p\in\mathbb N,\quad p\ge1.
\]
Every nonzero eigenvalue is semisimple with multiplicity two; the
original zero orders remain in the arithmetic weights. The
fractional continuation retains all prime powers, the full gamma
term and both poles. Its full factorial Dirichlet--Hankel block is
positive, while its pool-null part has positive index \(2K_+\) at every finite
smoothing parameter despite uniform trace-norm decay.

We prove all-radius inverse-completion bounds and individual
growth lower bounds for adaptive finite M\"obius sources. For a fixed
nonzero compact input, every nonempty fixed Fourier band retains
every zero pole and its original order. A zero of displacement
\(\delta>0\) and order \(m\) forces cumulative squared source energy
of limsup size at least \(Y^{2\delta}(\log Y)^{2m-2}\), with an
explicit positive Laurent-coefficient constant. A single compact
input canceling both poles has numerator norm at least
\(c/(1+|\Im\rho|)\) in a fixed band.

Finally, an integer cap \(M\asymp X^3\), chosen independently of
the input, identifies the complete corrected gamma-head energy
with the logarithmic M\"obius covariance up to absolute \(o(1)\)
for inputs of norm \(\exp(O((\log X)^{1/3}))\). An elementary
squarefree estimate gives an explicit primitive-sector remainder.
These results quantify zero detection and arithmetic transfer.
The global bound excluding off-line modes remains open; the
Riemann Hypothesis is unresolved.
\end{abstract}

\clearpage
\tableofcontents
\clearpage

\section{Introduction and main results}\label{sec:introduction}
For $\Re s>1$ let $\zeta(s)=\sum_{n\ge1}n^{-s}$ and thereafter use its
meromorphic continuation. Set
\[
 \xi(s)=\tfrac12s(s-1)\pi^{-s/2}\Gamma(s/2)\zeta(s),\qquad
 \Xi(z)=\xi(\tfrac12+iz).
\]
We use the unconditional functional equation, conjugation symmetry, and
the critical strip for the nontrivial zeros; see
\cite[\S\S25.2,25.4,25.10]{DLMF}. The all-height sampling bounds use
$N(T)=O(T\log(2+T))$, with multiplicity. It follows, for example,
by summing the unconditional unit-interval estimate of
\cite[Lemma 9.4 and Remark 9.7]{Kedlaya}.
The critical-frame obstruction separately uses the classical
unconditional positive critical-line proportion in
\cite[Theorem 1]{Conrey}. No simplicity assumption occurs.
All sums indexed by $\rho$ run over distinct nontrivial zeros, with an
explicit multiplicity $m_\rho$.

Our Fourier convention is
$\widehat h(t)=\int h(x)e^{-itx}\dd x$.
The packet variable is $y$, and its spectral kernel uses $e^{2ity}$.
No complex-entire extension of that real-line kernel is asserted.
The general analytic results below require no numerical certificate.
The separately identified first- and third-head theorems use outward-rounded
Arb/Acb computations, with exact contour or disk geometry and reproduced
certificates. They concern approximating heads, not zeta zeros.
The paper develops three complementary ways to measure an individual
hypothetical zero outside the critical line: the inertia of the
complete Weil form, the spectrum of a compact arithmetic operator,
and the growth of a completed finite M\"obius source. The principal
results are as follows.

\begin{enumerate}
\item \textbf{Exact inertia, including infinite exceptions.}
Theorem~\ref{thm:negative-index} proves
\(\sup_{L\ge1}n_-(A_{e^L})=2K_+\). Here \(K_+\) counts distinct
conjugate mirror pairs of Xi zeros at positive height. Every
individual pair has a detector, with both pole values canceled and
the full arithmetic expression retained. For the genuinely
logistic-smoothed kernel, Theorem~\ref{thm:local-smoothed-index}
attains this index on every bounded open interval after any fixed
finite collection of real exponential moment and jet constraints.
Proposition~\ref{prop:local-moment-price} gives the cost of
increasing the moment depth on that interval.

\item \textbf{A complete compact arithmetic spectrum.}
Theorem~\ref{thm:compact-complete-arithmetic-spectrum} constructs
\(R_{p,t}\) on the actual closed arithmetic range \(\mathcal C_t\).
For every fixed integer \(p\ge1\) and admissible heat \(t\),
\[
 \sigma(R_{p,t})=\{0,1\}\cup
  \{(p^2-1/4)/(p^2+z^2):\Xi(z)=0\}.
\]
All nonzero eigenvalues are semisimple, with two-dimensional
eigenspaces. This operator is not asserted selfadjoint; realness
of its entire spectrum is a separate question.

\item \textbf{Fractional arithmetic with a positive full block.}
Theorem~\ref{thm:fractional-complete-arithmetic} identifies the
whole fractional family, proves the full Dirichlet--Hankel
inequality, and pays its complex-parameter norm. The pool-null
interaction decays in trace norm uniformly over normalized heat,
but its positive index remains \(2K_+\) at every finite parameter.
Thus small absolute norm does not remove an off-line direction.

\item \textbf{Adaptive completion and individual source growth.}
Theorem~\ref{thm:adaptive-completion-radius} gives the exact
all-radius completion cost. If a selected Xi zero has displacement
\(\delta>0\), adaptive compact inputs of radius \(o(\log X)\),
subpolynomial input norm and subpolynomially nonvanishing selected
value force a source limsup exponent at least \(\delta\).
The original output heat is arbitrary; the preparation price has
its own explicit support--heat condition.

\item \textbf{Fixed-band detection with a multiplicity-sensitive energy cost.}
Theorem~\ref{thm:fixed-band-cumulative} uses one compact input
and an exact operator-valued Mellin identity. Every hypothetical
off-line zero remains visible in every suitable fixed Fourier band,
even when its height lies outside that band. Its order determines
the logarithmic power in the cumulative energy lower bound.
Corollary~\ref{cor:source-exception-measure} gives a quantitative
lower bound on the exceptional set in the source parameter \(X\);
Proposition~\ref{prop:dyadic-source-exceptions} states the precise
stronger almost-everywhere bound that would imply RH.

\item \textbf{A finite gamma realization of the signed covariance.}
Theorem~\ref{thm:joint-integer-gamma-completion} constructs one
integer cap for the unweighted and logarithmic sources, uniformly
over all inputs in the weighted space \(E\). Theorem~
\ref{thm:log-primitive-remainder} proves
\(O_q(C_h(0)X^{-1/2}H^{3/2}\log^2(2H))\) for the low
primitive-sector remainder. The signed main term and its
complement retain their original coefficients.
\end{enumerate}

The finite-exception negative-index phenomenon is recorded in
Bombieri's work on Weil's quadratic functional
\cite{BombieriWeil}. Our proofs state the full-domain and
infinite-index quantifiers explicitly. The classical semilocal
normalization is related to Connes--Consani
\cite{ConnesConsani}; Suzuki's screw-function approach
\cite{Suzuki2026} provides further operator-theoretic context.
The packet, Xi-defect and fixed-parameter operator foundations
developed in the author's earlier papers
\cite{PacketSource,OscillationSource,XiSource,OperatorSource,KernelSource}
are reconstructed where needed. The earlier scalar packet criterion,
its exact exponential strip rate and unconditional oscillation
are prior results; the squared-energy, finite-source and
arithmetic realization statements have their own proofs here.
The quantitative scope here includes explicit constrained local
indices, completion costs, the full compact arithmetic range, and
fixed-band cumulative source bounds. These statements are distinct
from identifying a cofinal family of positive approximants with Xi.
In particular, Suzuki's unconditional real-zero construction and
its conditional limiting identification
\cite[Theorem 1.5 and Corollary 1.6]{Suzuki2026} do not supply that
identification for the arithmetic operators constructed here.

The proofs are organized by mechanism. The first part establishes
packet growth, adaptive moments and arithmetic extraction bounds.
The second establishes the complete form domain, interpolation and
exact indices. The third constructs the arithmetic operators,
their compact and fractional realizations, and finite-source
transfers. The fourth treats theta approximation and certified
zero control. Conditional implications and open questions are
labeled as such. The final section states the remaining global
estimate once, separately from the proved results.

\part{Packets, adaptive moments, and arithmetic obstructions}

\section{A packet whose off-line coefficients never vanish}
Let $a_n=2^{-n-1}$, $n\ge1$, and let $g$ be the infinite convolution of
the uniform probability densities on $[-a_n,a_n]$. Equivalently it is
the density of the absolutely convergent sum of independent bounded
random variables with those distributions. Put $c_g=\|g\|_2$,
$h=g/c_g$, $K=h*h$, and
\[
 L(w)=\int_{-1}^{1}K(r)e^{wr}\dd r.
\]

\begin{lemma}\label{lem:packet}
The function $g$ is nonzero, real, even and smooth, supported in
$[-1/2,1/2]$. The entire function $L$ has zero set
\[
 Z(L)=4\pi i\Z\setminus\{0\}.
\]
For each $A>0$, uniformly for $|\sigma|\le1/2$,
$L(\sigma+it)=O_A((1+|t|)^{-A})$. Also
$L(it)=|\widehat h(t)|^2\ge0$ for real $t$.
\end{lemma}
\begin{proof}
The sum of support radii is $1/2$. Its distribution has a density by
convolving its first uniform density with the distribution of the
remaining sum. Its characteristic function is the locally uniformly
convergent product
\[
 \widehat g(z)=\prod_{n\ge1}\frac{\sin(a_nz)}{a_nz}.
\]
On the real axis each factor has modulus at most one. Bounding any
fixed number $q$ of factors by $1/(a_n|t|)$ gives decay $O_q(|t|^{-q})$.
Fourier inversion with arbitrarily many powers of $t$ proves smoothness.
The density integrates to one, so $0<c_g<\infty$.
Evenness gives
\[
 L(w)=c_g^{-2}\prod_{n\ge1}
      \left(\frac{\sinh(a_nw)}{a_nw}\right)^2.
\]
On each compact set the deviations of the tail factors from one are
bounded by a constant times $a_n^2$, whose sum is finite. The tail has
no additional zeros. The union of the factor zero sets is exactly the
stated zero set. In particular $L(w)\ne0$ if $\Re w\ne0$.
Since $K\in C_c^\infty(\R)$, integration by parts proves the vertical
decay, uniformly in the bounded real strip. The convolution theorem
and the fact that $h$ is real and even give $L(it)=|\widehat h(t)|^2$.
\end{proof}

Write $w_\rho=\rho-1/2$ and $c_\rho=m_\rho L(w_\rho)$.
Lemma~\ref{lem:packet} and zero counting imply
\[
 A:=\sum_\rho^\dist|c_\rho|<\infty.
\]
The same argument with any fixed power of $|w_\rho|$ also converges.
Thus
\[
 Q(y)=\sum_\rho^\dist c_\rho e^{2w_\rho y}\qquad(y\in\R)
\]
and all fixed real derivatives converge absolutely and locally uniformly.
Conjugation makes $Q$ real; the functional equation and evenness of $L$
make it even.

\begin{lemma}[Individual pole detection]\label{lem:strip}
For each fixed $\kappa\ge0$,
\[
 Q(y)=O(e^{\kappa y})\quad(y\longrightarrow+\infty)
 \quad\Longleftrightarrow\quad
 \Re\rho\le\tfrac12+\tfrac\kappa2\quad\hbox{for every }\rho.
\]
The bound on the left is on the full sufficiently large real tail.
\end{lemma}
\begin{proof}
Fix $Y_0\ge0$. On $\Re s>1/2$, absolute convergence permits integration
term by term:
\[
 F_{Y_0}(s):=\int_{Y_0}^\infty Q(y)e^{-2sy}\dd y
 =\tfrac12\sum_\rho^\dist
   \frac{c_\rho e^{-2(s-w_\rho)Y_0}}{s-w_\rho}.
\]
The right side is normally meromorphic on $\C$. On a compact set away
from its candidate poles, the exponential factors are uniformly bounded;
all sufficiently high ordinates have $|s-w_\rho|^{-1}\le1$. Summability
is a Weierstrass majorant. Each off-line zero gives an isolated genuine
simple pole of residue $c_\rho/2\ne0$. The remaining terms are holomorphic
near it. Multiplicity scales the residue and does not create a higher
order pole.

The assumed tail bound makes the integral holomorphic on $\Re s>\kappa/2$,
by locally uniform dominated convergence including derivatives in $s$.
It agrees there with the meromorphic expression by the identity theorem,
starting on the overlap with $\Re s>1/2$. A zero with
$\Re w_\rho>\kappa/2$ would give a nonremovable pole in that holomorphic
domain, a contradiction. Conversely, the stated zero bound gives
$|Q(y)|\le A e^{\kappa y}$ for $y\ge0$. No rightmost zero was used.
\end{proof}

\section{The cost of positive spectral approximation}
Let $\mathcal P_{\rm fin}$ consist of all finite sums
\[
 P(y)=\sum_{j=1}^n b_j e^{2it_jy},\qquad b_j\ge0,\quad t_j\in\R,
\]
including the zero sum. Define
\[
 B=\sup_\rho^\dist(\Re\rho-1/2)\in[0,1/2],\quad
 M(Y)=\max_{|y|\le Y}|Q(y)|,\quad
 d(Y)=\inf_{P\in\mathcal P_{\rm fin}}\|Q-P\|_{[-Y,Y]}.
\]

\begin{theorem}\label{thm:distance}
For every $Y\ge0$,
\[
 \tfrac12(M(Y)-Q(0))_+\le d(Y)\le M(Y).
\]
Moreover,
\[
 \limsup_{Y\to\infty}\frac{\log(1+d(Y))}{Y}=2B.
\]
For every fixed $\kappa\ge0$,
\[
 d(Y)=O(e^{\kappa Y})
 \quad\Longleftrightarrow\quad
 \Re\rho\le\tfrac12+\tfrac\kappa2\quad\hbox{for every }\rho.
\]
If $B=0$, then $d(Y)=0$ for every finite $Y$.
\end{theorem}
\begin{proof}
Each $P$ is positive definite. Its two-point Gram matrix implies
$P(0)\ge0$ and $|P(y)|\le P(0)$. If $e=\|Q-P\|_{[-Y,Y]}$, evaluation
at zero gives $P(0)\le Q(0)+e$. Therefore $M(Y)\le Q(0)+2e$.
Taking infima proves the lower bound; the zero sum proves the upper bound.
This uses the signed value $Q(0)$ and remains valid if it is negative.

Summability gives $M(Y)\le A e^{2BY}$. For $B>0$, a logarithmic limsup
for $M$ smaller than $2B$ would give $Q=O(e^{\kappa y})$ for some
$0\le\kappa<2B$. Lemma~\ref{lem:strip} would contradict the definition
of $B$. Hence that limsup is $2B$. In this case $M$, being nondecreasing,
is unbounded, and the two bounds transfer the exponent to $d$.

If $B=0$, zero symmetry gives RH. Then $c_\rho=m_\rho
|\widehat h(\Im\rho)|^2\ge0$. Finite symmetric truncations of the zero
series lie in $\mathcal P_{\rm fin}$ and approximate $Q$ uniformly on
all of $\R$, with error at most the omitted positive mass. Thus $d(Y)=0$.
Finally the two bounds transfer $d=O(e^{\kappa Y})$ to the same tail
bound for $Q$, and conversely; apply Lemma~\ref{lem:strip}.
\end{proof}

\begin{corollary}
If for every $\varepsilon>0$ there exist $C_\varepsilon,Y_\varepsilon$
such that $d(Y)\le C_\varepsilon e^{\varepsilon Y}$ for every
$Y\ge Y_\varepsilon$, then RH holds. None of those upper bounds has
been established here. A zero-free half-plane with boundary $7/8$, if
established, would give only the rate $3/4$; boundary $11/12$ would give
only $5/6$. Neither rate is zero.
\end{corollary}

\begin{remark}
An actual finite semilocal construction belongs to a subclass of
$\mathcal P_{\rm fin}$. The lower bound applies to that construction.
The abstract infimum being zero does not prove its realizability by
those particular operators. The kernel positivity here concerns $Q$,
not positivity of $\Xi$ or pointwise nonnegativity of $Q$.
\end{remark}

\section{Cumulative energy detects every individual off-line pole}
Fix $Y_0>1/2$ and, for $Y\ge Y_0$, put
\[
 G(Y)=\int_{Y_0}^{Y}|Q(y)|^2\dd y,\qquad
 V_2(Y)=\frac{G(Y)}{1+Y}.
\]

\begin{theorem}[Squared-energy bridge]\label{thm:l2}
With the same $B$ as in Theorem~\ref{thm:distance},
\[
 \limsup_{Y\to\infty}\frac{\log(1+G(Y))}{Y}
 =\limsup_{Y\to\infty}\frac{\log(1+V_2(Y))}{Y}=4B.
\]
For each fixed $\kappa\ge0$,
\[
 V_2(Y)=O(e^{\kappa Y})
 \quad\Longleftrightarrow\quad
 \Re\rho\le\tfrac12+\tfrac\kappa4\quad\hbox{for every }\rho.
\]
In particular bounded $V_2$, or subexponential $G$ with the full
quantifier \emph{for every $\epsilon>0$}, is equivalent to RH.
Bounded $G$ itself is not asserted to be equivalent to RH.
\end{theorem}
\begin{proof}
Absolute summability gives $|Q(y)|\le A e^{2By}$.
Consequently $G(Y)=O(e^{4BY})$ if $B>0$, and $G(Y)=O(1+Y)$ if $B=0$.
For $a\ge0$ write
\[
 I(a)=\int_{Y_0}^{\infty}e^{-ay}|Q(y)|^2\dd y.
\]
If it is finite, Cauchy--Schwarz implies, for $\Re s>a/4$,
\[
 \left|\int_{Y_0}^{\infty}Q(y)e^{-2sy}\dd y\right|
 \le I(a)^{1/2}
 \left(\int_{Y_0}^{\infty}e^{-(4\Re s-a)y}\dd y\right)^{1/2}.
\]
The same bound with any fixed power of $y$ proves locally uniform
convergence of every derivative in $s$. The transform is therefore
holomorphic there. By Lemma~\ref{lem:strip}, an off-line zero with
$\Re w_\rho>a/4$ would put a genuine pole in that domain.
Thus $I(a)<\infty$ implies $a\ge4B$, by testing each pole and
taking the supremum, without choosing a rightmost zero.
Conversely the absolute bound gives $I(a)<\infty$ for every $a>4B$.
Hence its nonnegative abscissa of convergence is exactly $4B$.

If $r=\limsup\log(1+G(Y))/Y<a$, then integration by parts, using
$G(Y)=O(e^{(a-\eta)Y})$ for some $\eta>0$, makes $I(a)$ finite.
If $r>a$, a sequence with $G(Y_n)\ge e^{(a+\eta)Y_n}$ gives
$I(a)\ge e^{-aY_n}G(Y_n)\to\infty$.
This proves $r=4B$, including $B=0$. Division by $1+Y$ changes
the logarithm by at most $\log(1+Y)$, so its rate is the same.

If $V_2=O(e^{\kappa Y})$, then $G=O((1+Y)e^{\kappa Y})$,
and $I(a)$ is finite for every $a>\kappa$. The pole argument gives
$4B\le\kappa$. In the other direction $4B\le\kappa$ and the
absolute bound give the claimed estimate; when $\kappa=0$,
$G=O(1+Y)$ supplies the required endpoint. Symmetry turns
$B=0$ into RH, with all multiplicities allowed.
\end{proof}

\section{The exact centered prime-power target}
Let
\[
 P_h^2=L(1/2)=L(-1/2),\qquad
 S_h(t)=\sum_{n\ge2}\frac{\Lambda(n)}{\sqrt n}K(t-\log n),
 \qquad E_h(t)=S_h(t)-e^{t/2}P_h^2.
\]
The sum is finite for each fixed $t$, because $K$ is supported on $[-1,1]$.

\begin{lemma}[Complete prime identity]\label{lem:prime}
For every $t>1$,
\[
 Q(t/2)=-E_h(t)-T_h(t/2),\qquad
 T_h(y)=\sum_{j\ge1}L(2j+1/2)e^{-(4j+1)y}.
\]
For every fixed $Y_0>1/2$, this last sum converges uniformly on
$y\ge Y_0$ and $T_h(y)=O_{h,Y_0}(e^{-5y})$.
The same assertion holds for every real, even, nonnegative
$K_0\in C_c^\infty([-1,1])$, with
$L_0(w)=\int K_0(r)e^{wr}\dd r$ and all of $Q,S,E,T$
defined using $K_0,L_0$ consistently; $Q$ is its absolutely
convergent zero series with the same multiplicities.
\end{lemma}
\begin{proof}
Work with any such $K_0$, written $K$ below. Its smooth
compact support gives arbitrary-power vertical decay on
each fixed strip, so zero counting makes the zero series
absolutely convergent. The proof uses no dyadic product
property or zero-avoidance property of $L$.
The identity agrees with \cite[\texttt{04\_explicit\_formula.tex},
equation \texttt{J-explicit}, after reflection]{PacketSource};
here is a contour derivation including its convergence conditions.
For $c>1$, consider
\[
 \frac1{2\pi i}\int_{\Re s=c}
 -\frac{\zeta'}{\zeta}(s)L(s-1/2)e^{t(s-1/2)}\dd s.
\]
The Dirichlet series is absolutely convergent. Fourier inversion of
$r\mapsto K(r)e^{(c-1/2)r}$ makes the integral exactly $S_h(t)$.
The interchange is justified by that series and the rapid vertical
decay of $L$.

Shift first to $\Re s=-2N-1$, with $N$ fixed. The residues of
$-\zeta'/\zeta$ are $+1$ at $s=1$, $-m_\rho$ at each nontrivial
zero, and $-1$ at $s=-2j$, $1\le j\le N$.
There is no pole at $s=0$. Heights tending to infinity can be
chosen at distance $\gg1/\log T$ from every zero ordinate in a
unit neighborhood: the unit-interval zero count bounds the total
length excluded by such intervals by less than one.
This is the zero-avoiding-height lemma in
\cite[\texttt{04\_explicit\_formula.tex},
Lemma \texttt{zero-avoiding-heights}]{PacketSource}.
On this sequence $\zeta'/\zeta=O_N(\log^2 T)$ uniformly on
the fixed horizontal segment. To verify this last bound,
pair the unconditional centered Hadamard terms at
$w_\rho$ and $-w_\rho$. Ordinates below $T/2$ and above $2T$
contribute $O_N(\log T)$ by zero counting and the paired
$O(T/|\gamma|^2)$ tail. In the middle range, unit-interval
counting and summing $1/k$ over distance bands give
$O_N(\log^2 T)$; the nearest band uses the chosen gap.
The gamma logarithmic derivative adds only $O_N(\log T)$.
The rapid vertical decay of $L$ now sends the horizontal
integrals to zero. The zero residues converge absolutely.

On the remaining left line the functional equation and the
absolutely convergent series at $1-s$ give
$\zeta'/\zeta(s)=O(\log(2+N+|\Im s|))$.
The odd negative real part stays away from the trivial zeros.
For every fixed integer $q$, integration by parts in the $L$-integral
gives
\[
 |L(\sigma+i\tau)|\le C_q e^{|\sigma|}
 (1+|\sigma|)^q(1+|\tau|)^{-q}.
\]
Thus the left integral is bounded by a polynomial in $N$
times $e^{-(t-1)(2N+3/2)}$ and tends to zero as $N\to\infty$.
This yields
$S_h(t)=e^{t/2}L(1/2)-Q(t/2)-T_h(t/2)$.
The left-line estimates are uniform when $t$ ranges in any
compact subset of $(1,\infty)$.
Finally $0\le L(2j+1/2)\le L(0)e^{2j+1/2}$, so on
$y\ge Y_0>1/2$ the tail after factoring out $e^{-5y}$
is dominated by a convergent geometric series.
\end{proof}

Define
\[
 \mathcal A(Y)=\int_{2Y_0}^{2Y}|E_h(t)|^2\dd t,\quad
 W_Y(u,v)=\int_{2Y_0}^{2Y}K(t-u)K(t-v)\dd t.
\]
For $t>1$, substitution $x=e^{t-r}$ gives exactly
\[
 \int_1^\infty x^{-1/2}K(t-\log x)\dd x=e^{t/2}P_h^2.
\]
Consequently finite expansion and Tonelli give
\begin{align*}
 \mathcal A(Y)={}&
 \sum_{m,n\ge2}\frac{\Lambda(m)\Lambda(n)}{\sqrt{mn}}
       W_Y(\log m,\log n)\\
 &-2\sum_{m\ge2}\frac{\Lambda(m)}{\sqrt m}
       \int_1^\infty\frac{W_Y(\log m,\log x)}{\sqrt x}\dd x\\
 &+\int_1^\infty\int_1^\infty
       \frac{W_Y(\log x,\log x')}{\sqrt{xx'}}\dd x\dd x'.
\end{align*}
All displayed sums and integrals have finite support for fixed $Y$.
In particular $m,n\le e^{2Y+1}$ and $W_Y(u,v)=0$ when $|u-v|\ge2$.
This is a centered correlation of every rational prime power.
None of its three terms may be replaced by an independently
subexponential positive bound: the last term alone is exactly
$(P_h^2)^2(e^{2Y}-e^{2Y_0})$.

For clarity, the elementary unconditional bound is only
\[
 \psi(X):=\sum_{n\le X}\Lambda(n)=O(X),\quad
 E_h(t)=O_h(e^{t/2}),\quad\mathcal A(Y)=O_h(e^{2Y}).
\]
The first bound follows without RH: every prime in $(n,2n]$
divides $\binom{2n}{n}$, so its logarithmic sum is at most
$2n\log2$. Dyadic summation gives the corresponding bound for
$\sum_{p\le X}\log p$; summing its values at $X^{1/m}$ over
$m\le\log_2 X$ gives $\psi(X)=O(X)$.
The support of $K$ and the lower bound $n\ge e^{t-1}$ in its
prime sum prove the bound for $E_h$.
Lemma~\ref{lem:prime} and these estimates show
\[
 \mathcal A(Y)=2G(Y)+O_{h,Y_0}(1):
\]
the cross term is integrable since
$E_h(t)T_h(t/2)=O(e^{-2t})$ and $T_h(t/2)^2=O(e^{-5t})$.
Therefore Theorem~\ref{thm:l2} gives the rigorous target
\[
 \text{RH}\quad\Longleftrightarrow\quad
 \forall\epsilon>0\ \exists C_\epsilon,Y_\epsilon\
 \forall Y\ge Y_\epsilon:\quad
 \mathcal A(Y)\le C_\epsilon e^{\epsilon Y}.
\]
No bound of this strength has been derived from the OpenAI
row-energy or Claude trace estimates. It is a remaining proof
obligation, not an established arithmetic lemma.

\begin{proposition}[Conditional joint strip budget]\label{prop:budget}
If a zero-free half-plane $\Re s>b$, $b\ge1/2$, is proved,
then for this same frozen packet
\[
 Q(y)=O(e^{(2b-1)y}),\quad
 d(Y)=O(e^{(2b-1)Y}),\quad
 V_2(Y)=O(e^{(4b-2)Y}).
\]
Thus the claimed boundaries give the respective rate pairs
$(3/4,3/2)$ and $(5/6,5/3)$ for $(Q,V_2)$.
Neither pair is subexponential. Adding a positive asymptotic
critical-line proportion supplies no estimate for the
individual nonzero residues of Lemma~\ref{lem:strip}.
\end{proposition}
\begin{proof}
Apply Lemma~\ref{lem:strip} and Theorems~\ref{thm:distance}
and \ref{thm:l2}. The rates are upper bounds, not statements
that the zero supremum is attained or equals the boundary.
A proportion statement counts zeros and does not bound
these separate weighted residues.
\end{proof}

\section{A quantitative height-tail bridge for the frozen packet}
The qualitative Vinogradov--Korobov theorem is unconditional:
for some $c_0,T_0>0$, every nontrivial zero with $|\gamma|\ge T_0$
satisfies
\[
 1-\beta\ge
 \frac{c_0}{(\log|\gamma|)^{2/3}
                     (\log\log|\gamma|)^{1/3}}.
\]
One precise primary reference is \cite[p.~1, (1.1) and Theorem 1]{FordZF},
in its corrected version. Theorem 1 uses the unconditional
zeta-growth estimate (1.2) stated there, supplied by
\cite[Theorem 1]{FordGrowth}; we use only existence of
positive constants. The latter theorem holds uniformly for
$t\ge3$ and $1/2\le\sigma\le1$. No all-height numerical zero verification or
the disputed original Cheng--Graham estimate (1.6) is used.

\begin{theorem}[Dyadic height-tail budget]\label{thm:height-tail}
For some $b,C>0$ and every $U\ge2$,
\[
 \sum_{|\Im\rho|>e^U}m_\rho|L(\rho-1/2)|\le Ce^{-bU^2}.
\]
Whenever $\beta\le1-\eta(U)$ for all zeros with
$|\gamma|\le e^U$, where $\eta(U)>0$, the same frozen packet satisfies
for all $y\ge0$
\[
 |Q(y)|\le Ce^y
       \bigl(e^{-2y\eta(U)}+e^{-bU^2}\bigr).
\]
In particular, unconditionally for all sufficiently large $y,Y$,
with possibly different positive constants,
\[
 |Q(y)|\le C\exp\{y-cy^{3/4}(\log y)^{-1/4}\},
 \qquad
 \mathcal A(Y)+G(Y)\le
 C\exp\{2Y-cY^{3/4}(\log Y)^{-1/4}\}.
\]
\end{theorem}
\begin{proof}
The product in Lemma~\ref{lem:packet} gives
\[
 L(w)=c_g^{-2}\prod_{n\ge1}
                 \left(\frac{\sinh(a_nw)}{a_nw}\right)^2.
\]
For $a>0$, $w=\sigma+it$, $|\sigma|\le1/2$, and $T=|t|>0$,
the integral representation and $|\sinh(aw)|\le\cosh(a\sigma)$
give
\[
 \left|\frac{\sinh(aw)}{aw}\right|
 \le e^{a|\sigma|}\min\{1,(aT)^{-1}\}.
\]
The exponential factors over the whole squared product multiply
to at most $e^{|\sigma|}$, because $\sum a_n=1/2$.
For $T\ge4$ choose $N=\lfloor\log_2T\rfloor-1$. Then
$a_nT\ge1$ for $1\le n\le N$, and
\[
 \log|L(\sigma+it)|
 \le\log(c_g^{-2})+\tfrac12
       +(\log2)N(N+3)-2N\log T
 =-\frac{(\log T)^2}{\log2}+O(\log T).
\]
This bounds every factor beyond $N$ as well; there is no discarded
infinite-product tail. Local uniform product convergence and
continuity for bounded $T$ therefore give
\[
 |L(\sigma+it)|\le C_0
       e^{-b_0(\log(3+|t|))^2},\qquad |\sigma|\le1/2,
\]
for some $C_0,b_0>0$.

On $e^j<|\gamma|\le e^{j+1}$, unconditional zero counting with
multiplicity bounds the total coefficient modulus by
$C e^{j+1}(j+1)e^{-b_0j^2}$. Summing the Gaussian tail from
$j=\lfloor U\rfloor$, and absorbing finitely many small $j$,
proves the first assertion with a smaller $b>0$.
The total coefficient modulus is finite by Lemma~\ref{lem:packet}.
For the zeros below $e^U$ use
$e^{2(\beta-1/2)y}\le e^{(1-2\eta(U))y}$; above that height
use only $\beta<1$ and the new coefficient-tail bound.
This proves the displayed budget without choosing a rightmost zero.

The Vinogradov--Korobov theorem gives, for sufficiently large $U$,
a budget
$\eta(U)=c_1U^{-2/3}(\log U)^{-1/3}$ for \emph{all} zeros below
$e^U$. Indeed the height function is decreasing there.
The finitely many zeros below its starting height have a strictly
positive minimum distance from $\beta=1$, by discreteness of the
zeros of entire $\xi$ and the unconditional zero-free line at $1$.
Shrinking $c_1$ includes that finite set; no critical-line
verification of it is assumed.
Take $U=y^{3/8}(\log y)^{-1/8}$. Then
$U^2=y^{3/4}(\log y)^{-1/4}$ and
$yU^{-2/3}(\log U)^{-1/3}$ is comparable with $U^2$.
Both terms of the budget yield the stated bound for $Q$.

Put $r(Y)=Y^{3/4}(\log Y)^{-1/4}$. Since $r'(Y)\to0$,
$(2Y-2cr(Y))'\ge1$ for all sufficiently large $Y$.
Integrating $|Q(y)|^2$ from a fixed large height to $Y$
is thus at most a constant times $e^{2Y-2cr(Y)}$.
The initial finite interval contributes a constant.
The exact identity $\mathcal A(Y)=2G(Y)+O(1)$ proves the energy
assertion after changing constants.
\end{proof}

The saving in this theorem is $o(Y)$ in the energy exponent.
It proves an upper rate at most $2$, not that the actual rate
is positive or equals $2$. Theorem~\ref{thm:l2} still identifies
that actual rate as $4B$. The bound does not exclude a fixed
$0<B<1/2$, and does not close the subexponential target.

\begin{proposition}[A power-tail smoothing variant]\label{prop:power-packet}
Fix $p>1$, put $a_{n,p}=(2\zeta(p))^{-1}n^{-p}$, and let
$g_p=*_{n\ge1}b_{a_{n,p}}$, $h_p=g_p/\|g_p\|_2$,
$K_p=h_p*h_p$. Both $g_p,h_p$ are smooth and supported in
$[-1/2,1/2]$; $K_p$ is smooth and supported in $[-1,1]$. The transform
$L_p(w)=\int K_p(r)e^{wr}\dd r$ is entire and has no zero
off the imaginary axis. Define
\[
 Q_p(y)=\sum_\rho^{\rm dist}m_\rho L_p(\rho-1/2)
                                      e^{2(\rho-1/2)y}.
\]
Fix $Y_0>1/2$ and define explicitly
\begin{align*}
 S_p(t)&=\sum_{n\ge2}\frac{\Lambda(n)}{\sqrt n}K_p(t-\log n),
 &E_p(t)&=S_p(t)-e^{t/2}L_p(1/2),\\
 G_p(Y)&=\int_{Y_0}^Y|Q_p(y)|^2\dd y,
 &\mathcal A_p(Y)&=\int_{2Y_0}^{2Y}|E_p(t)|^2\dd t.
\end{align*}
For constants depending on this fixed $p$, uniformly for
$|\sigma|\le1/2$,
\[
 |L_p(\sigma+it)|\le C_p e^{-b_p|t|^{1/p}}.
\]
Consequently, writing
$\mathcal V(y)=(\log y)^{2/3}(\log\log y)^{1/3}$,
unconditionally for all sufficiently large $y,Y$,
\[
 |Q_p(y)|\le C_p\exp\{y-c_p y/\mathcal V(y)\},\qquad
 \mathcal A_p(Y)+G_p(Y)
   \le C_p\exp\{2Y-c_pY/\mathcal V(Y)\}.
\]
The energy constants may also depend on the fixed $Y_0$.
Here the prime energy and $G_p$ use this same new packet;
the original frozen dyadic packet remains unchanged.
\end{proposition}
\begin{proof}
Since $\sum a_{n,p}=1/2$, the sums of independent uniform
variables on $[-a_{n,p},a_{n,p}]$ converge absolutely to
a variable supported on $[-1/2,1/2]$. Convolving the first
uniform density with the law of the remaining sum gives a
bounded probability density $g_p$ and $0<\|g_p\|_2<\infty$.
Its bilateral transform is the locally uniformly convergent product
\[
 L_p(w)=\|g_p\|_2^{-2}
      \prod_{n\ge1}\left(\frac{\sinh(a_{n,p}w)}{a_{n,p}w}\right)^2.
\]
Indeed each factor is $1+O_R(a_{n,p}^2)$ for $|w|\le R$.
For all sufficiently large indices these errors have modulus
at most $1/2$. Their analytic logarithms converge absolutely
and uniformly on that disk, since $\sum a_{n,p}^2<\infty$.
The resulting tail is an exponential of a holomorphic function
and has no zero. Thus every product zero comes from a finite
factor, at $w=\pi i k/a_{n,p}$, $k\in\Z\setminus\{0\}$.

Use the factor estimate from Theorem~\ref{thm:height-tail},
valid for every $a>0$. Write $d_p=(2\zeta(p))^{-1}$ and
$N=\lfloor(d_p|t|/2)^{1/p}\rfloor$. For $n\le N$,
$a_{n,p}|t|\ge2$. The whole squared product is bounded by
$\|g_p\|_2^{-2}e^{|\sigma|}4^{-N}$; every factor beyond $N$
is included in $e^{2|\sigma|\sum a_{n,p}}=e^{|\sigma|}$.
This proves the uniform strip bound after changing constants.
On the real Fourier axis the unsquared sinc product is bounded
by $2^{-N}$. Multiplication by every fixed power of $|t|$
remains integrable, so Fourier inversion gives
$g_p,K_p\in C_c^\infty$.

Zero counting and partial summation give
\[
 \sum_{|\gamma|>T}m_\rho|L_p(\rho-1/2)|
        \le C_p e^{-b'_pT^{1/p}},
\]
for all sufficiently large $T$, and give total coefficient
summability, including every fixed derivative weight.
For $y$ large choose $T=y^p$. The qualitative
Vinogradov--Korobov theorem, with the finite low-height
set handled as in Theorem~\ref{thm:height-tail}, gives
$\beta\le1-c'_p/\mathcal V(y)$ for every $|\gamma|\le T$.
The low-height part of $Q_p$ is at most
$C_pe^{y-2c'_py/\mathcal V(y)}$, while the high-height part
is at most $C_pe^{y-b'_py}$. Since $\mathcal V(y)\to\infty$,
their sum gives the asserted packet bound.

Lemma~\ref{lem:prime} applies to this new $K_p$ with its
stated general hypotheses. In particular the same contour
limits, prime powers and trivial-zero correction give
$\mathcal A_p(Y)=2G_p(Y)+O_{p,Y_0}(1)$.
Finally $(Y/\mathcal V(Y))'\to0$, so integrating the
squared bound proves the energy assertion just as before.
\end{proof}

Every off-line zero has a nonzero $L_p$-coefficient, so the
individual-pole growth test also applies to $Q_p$. The saving
is still $o(Y)$: changing the smoothing does not turn this
unconditional estimate into the required subexponential one.

\begin{proposition}[Degeneration of a moving power filter]
\label{prop:moving-filter}
For the normalized power-packet family above, as $p\downarrow1$,
\[
 \|g_p\|_2^2\asymp(p-1)^{-1},\qquad
 c_R(p-1)\le |L_p(w)|\le C_R(p-1)
 \quad(|w|\le R,\ 1<p<1+\delta_R)
\]
for every fixed $R<\infty$, with positive constants
$c_R,C_R,\delta_R$. Thus nonvanishing of each individual
$L_p(w)$ off the imaginary axis does not give a uniform
coefficient lower bound when $p$ tends to $1$.

In particular take any fixed $0<\delta<1/2$, $\gamma>0$,
and the four-point model $\mathcal W=\{\delta+i\gamma,
\delta-i\gamma,-\delta+i\gamma,-\delta-i\gamma\}$.
Although $L_p(w)\ne0$ for every $w\in\mathcal W$ and $p>1$,
the moving model
\[
 Q_{\rm mod}(y)=\sum_{w\in\mathcal W}
             L_{1+e^{-2y}}(w)e^{2wy}
\]
satisfies $|Q_{\rm mod}(y)|\le C e^{-(2-2\delta)y}$
for large $y$. It has off-axis model frequencies.
This is a counterexample to a moving-filter growth inference,
not a model of the zeta zero set.
\end{proposition}
\begin{proof}
Let $X_p$ be the sum of the independent uniform variables
with radii $a_{n,p}$. Its support lies in $[-1/2,1/2]$ and
\[
 v_p=\operatorname{Var}(X_p)
 =\frac{\zeta(2p)}{12\zeta(p)^2}
 \asymp(p-1)^2.
\]
The pole of $\zeta$ at $1$ gives
$a_{1,p}\sim(p-1)/2$.
Convolution with the first uniform density gives
$\|g_p\|_\infty\le(2a_{1,p})^{-1}$, so
$\|g_p\|_2^2\le(2a_{1,p})^{-1}$.
Chebyshev gives probability at least $3/4$ in
$[-2\sqrt{v_p},2\sqrt{v_p}]$; Cauchy--Schwarz then gives
$\|g_p\|_2^2\ge9/(64\sqrt{v_p})$.
These two bounds prove the norm assertion.
For $|w|\le R$, put $M_p(w)=\mathbb E e^{wX_p}$.
Then
\[
 |M_p(w)-1|\le R e^{R/2}\mathbb E|X_p|
                      \le R e^{R/2}\sqrt{v_p}=O_R(p-1).
\]
Since $L_p(w)=\|g_p\|_2^{-2}M_p(w)^2$, the compact
two-sided bound follows. In the finite four-point model,
each coefficient is $O(e^{-2y})$ and each exponential
has magnitude at most $e^{2\delta y}$, giving the claimed
bound. No zeta zero is prescribed or inferred.
\end{proof}

\section{Uniform endpoint partition error in the Xi defect}
The exact Xi Fourier kernel in \cite[eqs. (\texttt{phi}),
(\texttt{decay})]{XiSource} is
\[
 \Phi(u)=\sum_{n\ge1}(4\pi^2n^4e^{9|u|/2}-6\pi n^2e^{5|u|/2})
             e^{-\pi n^2e^{2|u|}}.
\]
It is positive, even and smooth, and is bounded by
$C e^{9|u|/2}\exp(-\frac\pi2e^{2|u|})$.
For completeness, these properties and the Fourier representation can be
reconstructed directly. Put $\psi(t)=\sum_{n\ge1}e^{-\pi n^2t}$ and
$\theta(t)=1+2\psi(t)$. Periodizing a Gaussian and computing its Fourier
coefficients gives $\theta(t)=t^{-1/2}\theta(1/t)$: both Fourier series
converge absolutely, and the coefficient integral is
$t^{-1/2}e^{-\pi k^2/t}$.
Termwise Mellin integration on $\Re s>1$, followed by splitting at one
and using that identity, gives
\[
 \xi(s)=\tfrac12+\tfrac12s(s-1)
 \int_1^\infty\psi(t)
      \{t^{s/2}+t^{(1-s)/2}\}\frac{\dd t}{t}.
\]
The last integral converges locally uniformly for all $s\in\C$, so the
identity continues meromorphically, with both sides in fact entire.
For $f(u)=e^{u/2}\psi(e^{2u})$ and $s=1/2+iz$ it becomes
\[
 \Xi(z)=\tfrac12-2(z^2+1/4)\int_0^\infty f(u)\cos(zu)\dd u.
\]
The theta identity gives $f'(0)=-1/4$ and
$f(-u)=f(u)+\sinh(u/2)$. Twice integrating by parts shows that
$\Phi=f''-f/4$ has the stated expression and satisfies
$\Xi(z)=\int_\R\Phi(u)e^{izu}\dd u$.
The same identity makes $\Phi$ smoothly even. Each summand is positive
for $u\ge0$ since $2\pi n^2e^{2u}>3$. Finally one may use
$C=4\pi^2\sum_{n\ge1}n^4e^{-\pi(n^2-1/2)}<\infty$ in the displayed
decay bound. All integrations by parts are justified by that decay and
the corresponding derivative bounds of the explicit Gaussian series.
The arguments of this section apply to any positive even kernel having
all the exponential moments used below; for the displayed kernel all
integrability claims follow from this explicit bound.
Put $p=u+v$, $q=u-v$, and, for $0\le y\le1/2$,
\[
 W_y(u,v)=p^2\sinhc(yp)\Phi(u)\Phi(v),\quad
 C_y(x)=\iint W_y\cos(xq)\dd u\dd v,
\]
\[
 A_{2j}(y)=\iint W_yq^{2j}\dd u\dd v,\qquad
 G_y(R)=\iint W_y\cosh(Rq)\dd u\dd v.
\]
Here $\sinhc(0)=1$. The formula
$\sinhc(t)=\int_0^1\cosh(vt)\dd v$ makes $W_y$ nonnegative and
nondecreasing in $y$. It also gives
\[
 |\partial_yW_y|\le |p|^3e^{|p|/2}\Phi(u)\Phi(v).
\]
The double exponential bound justifies Tonelli, differentiation under
these integrals, and all forthcoming compact uniform limits.

Define finite constants, for each fixed $R$, by
\begin{align*}
 a&=A_0(0)>0,\qquad H_R=G_{1/2}(R),\\
 D_R&=\sup_{0\le y\le1/2}\iint |\partial_yW_y|\cosh(Rq)\dd u\dd v,\\
 D_0&:=\sup_{0\le y\le1/2}|A_0'(y)|<\infty,\\
 C_R&=D_R/a+H_RD_0/a^2.
\end{align*}
The definition of $D_0$ agrees with $D_R$ at $R=0$, since
$\partial_yW_y\ge0$. Its finiteness follows from the same integrable
derivative majorant.
For $I=[\alpha,\beta]\subset[0,1/2]$, put
\[
 \ell_j(I)=\frac{A_{2j}(\alpha)}{A_0(\beta)},\qquad
 u_j(I)=\frac{A_{2j}(\beta)}{A_0(\alpha)},\qquad
 \mu_j(y)=\frac{A_{2j}(y)}{A_0(y)}.
\]
Then $\ell_j\le\mu_j(y)\le u_j$ for $y\in I$. Set
\[
 S_m(y,x)=\sum_{j=0}^{2m+1}(-1)^j\mu_j(y)\frac{x^{2j}}{(2j)!}
\]
and let $P_{m,I}(x)$ have constant term $1$, lower coefficients $\ell_j$
for positive even $j\ge1$ and upper coefficients $u_j$ for negative odd $j$.

\begin{theorem}[Order-independent mesh control]\label{thm:mesh}
For every $m\ge0$, $y\in I$, and $|x|\le R$,
\[
 0\le S_m(y,x)-P_{m,I}(x)\le C_R(\beta-\alpha).
\]
For any finite partition $\mathcal I$ of a subinterval of $[0,1/2]$,
with mesh $\delta$, there are $T_m(R)\to0$ such that
\[
 \sup_{I\in\mathcal I,y\in I,|x|\le R}
 \left|P_{m,I}(x)-\frac{C_y(x)}{A_0(y)}\right|
 \le T_m(R)+C_R\delta.
\]
Thus every diagonal $m_n\to\infty$, $\delta_n\to0$ converges uniformly
on that fixed compact rectangle, without a rate coupling.
\end{theorem}
\begin{proof}
Coefficient signs give the first inequality. Majorize the finite error
by the nonnegative infinite sum; Tonelli gives exactly
\[
 \sum_{j\ge0}(u_j-\ell_j)\frac{R^{2j}}{(2j)!}
 =\frac{G_\beta(R)}{A_0(\alpha)}-\frac{G_\alpha(R)}{A_0(\beta)}.
\]
Including $j=0$ is harmless: the actual polynomial constant is exactly
one. Split the right side as
\[
 \frac{G_\beta-G_\alpha}{A_0(\alpha)}+
 G_\alpha\left(\frac1{A_0(\alpha)}-\frac1{A_0(\beta)}\right).
\]
The derivative bounds and $A_0\ge a$ prove the claimed $C_R|I|$ bound,
simultaneously for all orders.

For the Taylor remainder one may take
\[
 T_m(R)=\frac1a\iint W_{1/2}
       \sum_{j>2m+1}\frac{R^{2j}|q|^{2j}}{(2j)!}\dd u\dd v.
\]
It tends to zero by dominated convergence with majorant
$W_{1/2}\cosh(Rq)/a$. This proves the second estimate and the diagonal
assertion. The estimates are compact in $R$, not uniform as $R\to\infty$.
\end{proof}

The polynomials are genuine lower bounds. For real $t$ the Taylor
polynomial $T_{4m+2}(t)$ of cosine is at most $\cos t$.
For $m=0$ this follows by integrating $1-\cos t\ge0$ twice.
For $m\ge1$, the fourth derivative of the difference is
$\cos t-T_{4m-2}(t)$ and its four initial values vanish; induction and
evenness prove the inequality for all real $t$. Integrating it against
$W_y/A_0(y)$ proves $P_{m,I}\le S_m\le C_y/A_0(y)$.

\begin{corollary}[Positive-margin exhaustion]
Fix $0<\varepsilon<1/2$ and $R<\infty$. If $C_y(x)/A_0(y)$ has positive
minimum $\eta$ on $[\varepsilon,1/2]\times[-R,R]$, there are
$m_0,\delta_0>0$ such that every $m\ge m_0$ and every finite partition
of mesh at most $\delta_0$ has $P_{m,I}(x)\ge\eta/3$ for every cell and
every $|x|\le R$. Conversely such a positive lower certificate implies
positivity on that compact rectangle.
\end{corollary}
\begin{proof}
Choose $T_{m_0}(R)<\eta/3$ and $C_R\delta_0<\eta/3$; the tail bound is
nonincreasing in $m$. Apply Theorem~\ref{thm:mesh}. The converse uses
the lower-bound inequality.
\end{proof}

\section{The global positivity step is still missing}
Multiplying the absolutely convergent Xi Fourier integrals, using evenness
and symmetrizing, gives
\[
 |\Xi(x+iy)|^2=\iint\Phi(u)\Phi(v)\cosh(yp)\cos(xq)\dd u\dd v.
\]
Differentiation therefore gives
$\partial_y|\Xi(x+iy)|^2=yC_y(x)$.
If $\Xi(x+iy)=0$ with $y>0$, its squared modulus has zero first
derivative, so $C_y(x)=0$. Thus strict defect positivity in the entire
open upper strip excludes all off-line zeros by zero symmetry.

Conversely, under RH the order-one Hadamard factorization
\cite[pp.~1--2, equation (1), and p.~5, Exercise 2]{Elkies}, paired by
evenness, gives
$\Xi(z)=\Xi(0)\prod_{\gamma>0}(1-z^2/\gamma^2)^{m_\gamma}$.
The cited notes use $\pi^{-s/2}\Gamma(s/2)\zeta(s)$ for their meromorphic
$\xi$; multiplication by $s(s-1)/2$ converts their product to ours.
The zero-counting bound makes the product and its paired logarithmic
derivative locally convergent away from its zeros. For $y>0$ it gives
\[
 \Im\frac{\Xi'(x+iy)}{\Xi(x+iy)}
 =-\sum_{\gamma>0}m_\gamma
 \left(\frac{y}{(x-\gamma)^2+y^2}+
       \frac{y}{(x+\gamma)^2+y^2}\right)<0.
\]
Here the paired order-one Hadamard product has no residual linear
exponential factor, by evenness. It follows that $C_y(x)>0$ for $y>0$.
Consequently the still-unproved positive-margin statement for every
$\varepsilon>0$ and every finite $R$ is RH-equivalent. At $y=0$,
\[
 C_0(x)=2(\Xi'(x)^2-\Xi(x)\Xi''(x)).
\]
At a multiple real zero this vanishes. Strict boundary positivity would
silently add simplicity, so it is not required here.

\section{A compact-positive control with exact off-real zeros}
This section concerns a model Fourier kernel, not the zeta function.
Fix $0<\delta<1/2$, $R,Y\ge0$, and put
\[
 a=\frac1{4(1+R+Y)},\quad c=\cosh(a\delta),\quad
 F_a(z)=\frac{c+\cos(az)}{c+1},\quad X_a(z)=e^{-z^2/2}F_a(z).
\]
Writing $\varphi(u)=(2\pi)^{-1/2}e^{-u^2/2}$, the Fourier density is
\[
 \Phi_a(u)=\frac{c}{c+1}\varphi(u)+
 \frac{\varphi(u-a)+\varphi(u+a)}{2(c+1)}>0.
\]
It is smooth and even with all exponential moments. Its transform has
exact simple zeros
\[
 z=\frac{(2k+1)\pi}{a}\pm i\delta\qquad(k\in\Z).
\]
To solve $\cos(az)=-\cosh(a\delta)$, the imaginary part forces the real
part of $az$ to be an odd multiple of $\pi$, and then the imaginary part
has magnitude $a\delta$. The derivative is nonzero at these points.
All these zeros lie outside $|x|\le R$.

Yet its defect is strictly positive throughout $|x|\le R$, $0\le y\le Y$.
Here are explicit bounds, including $y=0$. On that rectangle
$|az|\le1/4$ and $a\le1/4$. The elementary series bound
$e^{1/4}<4/3$ gives
\[
 |F_a-1|\le1/48,\qquad |F_a'|\le1/24,\qquad |F_a''|\le1/24.
\]
Set $H(x,y)=|F_a(x+iy)|^2$. Since $F_a$ is real entire,
$\partial_yH(x,0)=0$, and
\[
 |\partial_y^2H|\le2(|F_a'|^2+|F_a''||F_a|)
 \le2\left((1/24)^2+(1/24)(49/48)\right).
\]
The defect extending continuously to zero is
\[
 C_a(x,y)=e^{-x^2+y^2}
 \left(2H(x,y)+\frac{\partial_yH(x,y)}{y}\right).
\]
By the fundamental theorem of calculus the divided derivative has the
same bound as $\partial_y^2H$, also at zero. Hence
\[
 C_a(x,y)\ge e^{-x^2+y^2}
 \left[2(47/48)^2-2((1/24)^2+(49/48)/24)\right]
 >e^{-x^2+y^2}>0.
\]
This is an exact analytic countermodel to a generic inference from
compact defect positivity to global real zeros. It does not have the
zeta Euler product or the exact Riemann Fourier kernel and therefore
is not a counterexample to RH.

\section{An order-at-most-one control from the exact Riemann density}
The order-two growth of the preceding control is not essential to its
obstruction. Here the model is built from the exact density $\Phi$, but
its density and arithmetic function are changed by scaling and convolution.
Put
\[
 M_0=\int\Phi(u)\dd u=\Xi(0)>0,\quad
 \dd\mu(u)=\Phi(u)\dd u/M_0,\quad
 v=\int u^2\dd\mu(u)>0.
\]
Fix $0<\delta<1/2$ and, for $a>0$, set $b=a^2$, $c=\cosh(b\delta)$ and
\[
 H_a(z)=\Xi(az)/M_0,\quad
 J_a(z)=\frac{c+\cos(bz)}{c+1},\quad X_a(z)=H_a(z)J_a(z).
\]
Its exact positive even probability density is
\[
 \phi_a(t)=\frac{\Phi(t/a)}{aM_0},\qquad
 p_a(t)=\frac{c}{c+1}\phi_a(t)
       +\frac{\phi_a(t-b)+\phi_a(t+b)}{2(c+1)}.
\]
It is smooth, strictly positive and doubly exponentially decreasing for
each fixed $a$. The cosine factor proves that every point
\[
 z=(2k+1)\pi/a^2\pm i\delta\qquad(k\in\Z)
\]
is a zero of $X_a$. The product can have other zeros or larger
multiplicities; no simplicity assertion is made.

For completeness, the density bound proves entire order at most one.
For $T=|\Im z|$ put $A=9/2+aT$ and $c_0=\pi/2$. Then
\[
 |H_a(z)|\le (2C/M_0)\int_0^\infty e^{Au-c_0e^{2u}}\dd u.
\]
With $r=e^{2u}$ the last integral is
$\tfrac12\int_1^\infty r^{A/2-1}e^{-c_0r}\dd r$.
Bounding $r^{A/2}e^{-c_0r/2}$ by its maximum
$(A/(c_0e))^{A/2}$ leaves the finite constant
$\tfrac12\int_1^\infty r^{-1}e^{-c_0r/2}\dd r$.
Thus $\log\max_{|z|\le r}|H_a(z)|=O_a(r\log(2+r))$.
The cosine factor has exponential type $b$, so $X_a$ has order at most one.

We next quantify its compact-positive defect. Fix $R,Y\ge0$, put
$\Lambda=1+R+Y$ and
\[
 E_4=\int |u|^4e^{\Lambda|u|}\dd\mu(u)<\infty.
\]
For $0<a\le1$, $|x|\le R$, $0\le y\le Y$, and $z=x+iy$, Taylor's
series for cosine and sine, dominated by $e^{|w|}$, give
\begin{align*}
 |H_a-1+va^2z^2/2|&\le a^4A_0,&
 |H_a'+va^2z|&\le a^4A_1,&
 |H_a''+va^2|&\le a^4A_2,\\
 A_0&=\Lambda^4E_4/24,& A_1&=\Lambda^3E_4/6,& A_2&=\Lambda^2E_4/2.
\end{align*}
The same dominated integrals justify these derivatives. Likewise,
$|J_a-1|\le a^4B_0$, $|J_a'|\le a^4B_1$, $|J_a''|\le a^4B_2$, where
\[
 B_0=\Lambda^2e^\Lambda/4,\quad
 B_1=\Lambda e^\Lambda/2,\quad B_2=e^\Lambda/2.
\]
Write $d_0=v\Lambda^2/2$, $d_1=v\Lambda$, $B_H=1+d_0+A_0$ and define
\begin{align*}
 K_0&=A_0(1+B_0)+(1+d_0)B_0,\\
 K_1&=A_1(1+B_0)+d_1B_0+B_HB_1,\\
 K_2&=A_2(1+B_0)+vB_0+2(d_1+A_1)B_1+B_HB_2.
\end{align*}
Product differentiation now proves
\[
 |X_a-1+va^2z^2/2|\le a^4K_0,\quad
 |X_a'+va^2z|\le a^4K_1,\quad
 |X_a''+va^2|\le a^4K_2.
\]
For $S_a(x,y)=|X_a(x+iy)|^2$ use
$S_{a,yy}=2|X_a'|^2-2\Re(X_a''\overline{X_a})$.
The displayed bounds give
\[
 |S_{a,yy}/a^2-2v|\le a^2K_S,
\]
where the finite, explicitly defined constant
\[
 K_S=2(d_1+K_1)^2+
      2\{v(d_0+K_0)+K_2(1+d_0+K_0)\}
\]
suffices. Since $X_a$ is real entire, $S_{a,y}(x,0)=0$. Its continuous
defect is therefore
\[
 C_a(x,y)=\int_0^1S_{a,yy}(x,ty)\dd t\quad(0\le y\le Y),
\]
equal to $S_{a,y}/y$ when $y>0$, and $S_{a,yy}(x,0)$ at zero. Consequently
\[
 \sup_{|x|\le R,\,0\le y\le Y}|C_a(x,y)/a^2-2v|\le a^2K_S.
\]
For $0<a\le\min\{1,\sqrt{v/(2K_S)}\}$ this is strictly positive
throughout the prescribed rectangle. If necessary decreasing $a$ further
ensures $\pi/a^2>R$. The case $K_S=0$ requires no further restriction.
This supplies exact off-real model zeros with positive compact defect,
order at most one and positive rapidly decreasing density. It still does
not preserve the zeta arithmetic function and is not an RH counterexample.

\section{Exact universal witnesses retain the sextic cost}
Use the unrestricted prime-ideal denominator family of
\cite[p.~4, \S2.2]{OpenAI1112}. Write
$\mathcal O=\mathbb Z[\omega]$, $E=\mathbb Q(\omega)$.
Outside primes dividing $6u$, let $\chi_{\mathfrak p}(u)$ be
the sixth root of unity congruent to
$u^{(N_E\mathfrak p-1)/6}$ modulo $\mathfrak p$.
It agrees with the source's primary-generator convention.
All prime ideals outside the fixed finite exclusions are
allowed, including those above inert rational primes.
When this statement is coupled to $A_u(D)$, those exclusions
include the source's fixed set $S$.

\begin{theorem}[Universal sextic projectors]\label{thm:projector}
For nonzero $u\in\mathcal O$, the identity
$\chi_{\mathfrak p}(u)=1$ for every prime ideal outside
a fixed finite set containing the divisors of $6u$
is equivalent to $u=a^6$ with nonzero $a\in\mathcal O$.

Let $U$ be a finite set of nonzero elements and suppose
\[
 \sum_{u\in U}w_u\chi_{\mathfrak p}(u)=1
\]
for every prime ideal outside one fixed finite set
containing the divisors of $6\prod_{u\in U}u$.
For the represented classes $c\in E^\times/E^{\times6}$,
\[
 \sum_{[u]=[1]}w_u=1,\qquad
 \sum_{[u]=c}w_u=0\quad(c\ne[1]).
\]
If $N_Eu\le H$ throughout $U$, then
\[
 \sum_{u\in U}|w_u|^2\ge M_H^{-1},\qquad
 M_H=\frac{\pi}{3\sqrt3}H^{1/6}+O(H^{1/12}),
\]
where $M_H$ counts nonzero sixth-power values of norm at most $H$.
The bound is attained on the full cutoff set by uniform weights
on its sixth powers and zero weights elsewhere.
\end{theorem}
\begin{proof}
Choose $\alpha_u^6=u$ and let $F=E(\alpha_u:u\in U)$.
Since $E$ contains $\mu_6$, this is a finite Galois extension.
Its group injects into $\mu_6^U$ by the characters
$\psi_u(\sigma)=\sigma(\alpha_u)/\alpha_u$, so it is abelian.
The characters $\psi_u,\psi_v$ agree exactly when
$u/v\in E^{\times6}$: equality makes $\alpha_u/\alpha_v$
fixed by the group, and the converse follows by choosing the
roots consistently. At unramified primes not dividing
$6\prod u$, Frobenius gives
$\psi_u(\operatorname{Frob}_{\mathfrak p})
=\chi_{\mathfrak p}(u)$. Reduction proves this because
$\alpha_u^{N_E\mathfrak p-1}
=u^{(N_E\mathfrak p-1)/6}$ and reduction is injective on
$\mu_6$. See also \cite[VIII, \S\S5.2--5.3,
pp.~244--245]{Milne}.

Chebotarev \cite[V, Theorem 3.23, p.~164]{Milne}
realizes every group element outside any fixed finite
exception set. Thus the weighted primewise hypothesis gives
$\sum_c W_c\psi_c(\sigma)=1$ for every $\sigma$ in the group,
where $W_c=\sum_{[u]=c}w_u$.
Distinct finite-group characters are orthogonal: their
quotient has zero group sum, since multiplication by one
of its values different from one leaves the sum unchanged.
Averaging against each conjugate character therefore gives
the asserted $W_c$. In particular the unweighted
single-element hypothesis forces its sixth root into $E$.
A sixth root of an integral $u$ is integral over $\mathcal O$;
integral closedness puts it in $\mathcal O$.
This proves the first statement and identifies class $[1]$.

Cauchy--Schwarz on the at most $M_H$ sixth-power rows
now proves the norm bound. To count them, the Eisenstein
lattice has covolume $\sqrt3/2$, hence
$\#\{a:N_Ea\le X\}=(2\pi/\sqrt3)X+O(\sqrt X+1)$.
This elementary lattice estimate follows by covering with
fundamental cells; only a boundary strip of fixed width
changes the disk area. Take $X=H^{1/6}$ and divide by six:
for nonzero elements $a^6=b^6$ exactly when $b/a\in\mu_6$.
Each value has six preimages. Uniform weights attain equality
outside the finite set of their prime divisors.
The fiber statement does not mean that every unit is a
sixth power: $\varepsilon^6=1$ for $\varepsilon\in\mu_6$.
\end{proof}

The least exact universal projector norm is therefore of
order $H^{-1/12}$. This statement alone does not bound projectors
required only on finite denominator sets. The next result gives
a separate finite-denominator bound, including approximate extraction.

\section{Smooth completion bounds finite sextic extractors}
Let $\mathcal C_D$ consist of the squarefree ideals prime to a fixed
finite set $S$ containing the primes above $6$, with norms in $[D,2D]$,
each represented once by a primary generator. Write
$N_D=|\mathcal C_D|$, $\mathcal U_H=\{u\in\mathcal O:0<N_Eu\le H\}$,
and
\[
 M_{\mathfrak n,u}=\chi_{\mathfrak n}(u),
       \qquad D\ge2,\quad H\ge1.
\]
These are raw multiplicative sextic symbols, with zero extension;
no Hecke-character identification is a premise of this section.

\begin{theorem}[Finite Gram and extraction bounds]\label{thm:finite-gram}
Choose once a nonnegative radial $C_c^\infty$ majorant equal to one
on the unit disk. A constant $C$, depending only on this fixed
majorant and independent of $D,H,S$, satisfies
\[
 \|M^*a\|_2^2\le C(H+DN_D)\|a\|_2^2
                  \quad(a\in\C^{\mathcal C_D}).
\]
If $N_D>0$, $Mc=\mathbf1+r$ and
$|\sum_{\mathfrak n}r_{\mathfrak n}|\le\eta N_D$
for $0\le\eta<1$, then
\[
 \|c\|_2^2\ge
   \frac{(1-\eta)^2N_D}{C(H+DN_D)}.
\]
For fixed $S$, $N_D\asymp_SD$ as $D\to\infty$, so this is
$\gg_S(1-\eta)^2D/(H+D^2)$.
In particular for each fixed $0<\vartheta<1$,
$H=D^{1+\vartheta}$, and fixed $\eta<1$, no such extractor has
$\|c\|_2^2\le H^{-1+o(1)}$.
\end{theorem}
\begin{proof}
Use the additive pairing and area measure
\[
 e(z)=\exp(4\pi i\Im z/\sqrt3),\qquad
 dz_*=(2/\sqrt3)\dd x\dd y.
\]
If $z=a+b\omega$ and $w=c+d\omega$, then
$(2/\sqrt3)\Im(zw)=ad+bc-bd$.
This integral pairing has determinant $-1$. Thus $\mathcal O$
is self-dual, has covolume one in $dz_*$, and the annihilator
of $q\mathcal O$ is $q^{-1}\mathcal O$.

For squarefree $q\ne1$ prime to $6$, let $\chi$ be a zero-extended
multiplicative character nontrivial at each prime dividing $q$.
The complete Gauss sums
\[
 \tau_\chi(\alpha)=\sum_{b\bmod q}\chi(b)e(b\alpha/q)
\]
satisfy $\tau_\chi(0)=0$ and $|\tau_\chi(\alpha)|\le\sqrt{N_Eq}$.
Here is a direct check. The Chinese remainder theorem splits
the sum into finite-field Gauss sums. A trivial local additive
character gives zero. For a nontrivial additive character $\psi$
of a field of size $Q_p$ and a nontrivial $\chi_p$,
\[
 \left|\sum_x\chi_p(x)\psi(x)\right|^2
 =\sum_{v\ne0}\chi_p(v)\sum_{y\ne0}\psi((v-1)y)
 =(Q_p-1)-\sum_{v\ne0,1}\chi_p(v)=Q_p.
\]
The self-dual pairing makes the local additive character
nontrivial exactly when the local frequency is a unit.
Multiplication over the primes proves the claim.

Fix one radial $F\in C_c^\infty(\C)$ with $F\ge0$,
$F=1$ on $|z|^2\le1$ and $F=0$ on $|z|^2\ge2$.
Let $\widehat F(w)=\int F(z)e(-zw)\,dz_*$.
For every $R>0$, periodization on $q\mathcal O$ gives
\[
 \sum_{v\in\mathcal O}\chi(v)F(v/\sqrt R)
 =\frac{R}{Q}\sum_{\alpha\in\mathcal O}
       \tau_\chi(\alpha)\widehat F(\sqrt R\,\alpha/q),
       \qquad Q=N_Eq.
\]
To see its normalization, the Fourier coefficient at $\alpha/q$
of the periodization is
$R\widehat F(\sqrt R\,\alpha/q)/Q$; integrate over one
fundamental domain. Schwartz decay makes the Fourier series
absolutely convergent. Summing its values against $\chi(b)$
gives precisely the displayed complete Gauss sum.

For $A>2$, integration by parts bounds $\widehat F(w)$ by
$C_A(1+|w|)^{-A}$. With $v_0=\sqrt{R/Q}$ the missing zero
frequency therefore gives
\[
 \left|\sum_v\chi(v)F(v/\sqrt R)\right|
 \le C_A\frac{R}{\sqrt Q}
     \sum_{\alpha\ne0}(1+v_0|\alpha|)^{-A}
 \le C_F\sqrt Q.
\]
Indeed lattice annuli bound the last sum by $O(v_0^{-2})$
if $v_0\le1$, and by $O(v_0^{-A})$ if $v_0\ge1$.
The bound is uniform for every positive scale $R$.

Consider the positive semidefinite majorizing Gram matrix
\[
 G^F_{\mathfrak n,\mathfrak m}
 =\sum_{u\in\mathcal O}\chi_{\mathfrak n}(u)
       \overline{\chi_{\mathfrak m}(u)}F(u/\sqrt H).
\]
Its diagonal is at most $C_FH$ by lattice counting.
The $u=0$ term is zero since $\mathfrak n$ is not the unit ideal.
For distinct $\mathfrak n,\mathfrak m$, put
$\mathfrak d=(\mathfrak n,\mathfrak m)$,
$\mathfrak n=\mathfrak d\mathfrak a$,
$\mathfrak m=\mathfrak d\mathfrak b$.
Squarefreeness makes $\mathfrak d,\mathfrak a,\mathfrak b$
pairwise coprime. The character
$\chi=\chi_{\mathfrak a}\overline{\chi_{\mathfrak b}}$
is nontrivial at each prime of
$\mathfrak q=\mathfrak a\mathfrak b\ne1$:
each local symbol has order six.
On the whole lattice, including nonunits,
\[
 \chi_{\mathfrak n}(u)\overline{\chi_{\mathfrak m}(u)}
       =\mathbf1_{(u,\mathfrak d)=1}\chi(u).
\]
Inclusion-exclusion over $\mathfrak k\mid\mathfrak d$ and
$u=kv$, using one generator of each ideal, gives
\[
 |G^F_{\mathfrak n,\mathfrak m}|
 \le C_F\tau(\mathfrak d)\sqrt{N_E\mathfrak q}
 =C_F\frac{\tau(\mathfrak d)}{N_E\mathfrak d}
             \sqrt{N_E\mathfrak n\,N_E\mathfrak m}
 \le2C_FD.
\]
Radiality of $F$ changes the inner scale to $H/N_E\mathfrak k$,
where the preceding uniform bound applies. There is no added
unit multiplicity. The last inequality uses
$\tau(\mathfrak d)/N_E\mathfrak d
=\prod_{\mathfrak p\mid\mathfrak d}2/N_E\mathfrak p\le1$.
Schur's bound now gives
$\|G^F\|\le C(H+DN_D)$.
Since $F\ge1$ on the original row set, this proves the energy bound.

For the extractor, Cauchy--Schwarz gives
\[
 (1-\eta)N_D\le|\langle\mathbf1,Mc\rangle|
 \le\|M^*\mathbf1\|_2\|c\|_2.
\]
Applying the energy bound with $a=\mathbf1$ proves the norm bound.
Both $\|r\|_\infty\le\eta$ and
$\|r\|_2\le\eta\sqrt{N_D}$ imply its stated residual hypothesis.

For completeness, squarefree column counting needs no prime ideal
theorem. The lattice count divided by six gives the number of
nonzero integral ideals of norm at most $X$ as
$\kappa X+O(\sqrt X+1)$, $\kappa=\pi/(3\sqrt3)$.
Fixed finite inclusion-exclusion at $S$ gives
$\kappa_SX+O_S(\sqrt X+1)$ with $\kappa_S>0$.
Using $\mathbf1_{\mathrm{squarefree}}(\mathfrak n)
=\sum_{\mathfrak d^2\mid\mathfrak n}\mu_E(\mathfrak d)$,
the count of squarefree $S$-coprime ideals is
\[
 \kappa_SX
  \sum_{(\mathfrak d,S)=1}\frac{\mu_E(\mathfrak d)}
                                  {(N_E\mathfrak d)^2}
       +O_S(\sqrt X\log(2+X)).
\]
The error follows from
$\sum_{N_E\mathfrak d\le\sqrt X}(N_E\mathfrak d)^{-1}
=O(\log(2+X))$ and an $O(X^{-1/2})$ tail for the absolutely
convergent squared-norm series, both by partial summation of
the elementary ideal count. The coefficient is positive by
its convergent Euler product
$\prod_{\mathfrak p\notin S}(1-(N_E\mathfrak p)^{-2})$.
Taking the difference at $2D$ and $D$ proves $N_D\asymp_SD$.
For fixed $0<\vartheta<1$ the resulting lower bound is
$\gg_SD^{-1}$, whereas
$H^{-1+o(1)}=D^{-1-\vartheta+o(1)}$, which is smaller.
\end{proof}

This new finite bound does not prohibit the scale $D^{-1+o(1)}$.
It also does not rule out $\vartheta=\vartheta(D)\to0$, for which
$D^{-1}=H^{-1+o(1)}$, or a functional designed to work on only
one particular M\"obius/Hecke vector. It does not use or prove
the source's Proposition 3.1. The next theorem strengthens this
finite obstruction by a separate raw-symbol recursion.

\section{A raw prefix sieve strengthens the finite obstruction}
We retain the literal Eisenstein residue symbols, without identifying
them with a globally unit-invariant Hecke family.
We use the following existing algebraic facts in precisely the form
given in \cite[\S2.1, PDF p.~3]{GaoZhao}: each ideal prime to $6$ has a
unique $E$-primary generator, products of such generators are
$E$-primary, and for coprime $E$-primary $m,n$
\[
 \chi_n(m)=\chi_m(n)
 (-1)^{((N_Em-1)/2)((N_En-1)/2)}.                 \tag{R}
\]
These are algebraic reciprocity facts, not a large-sieve input.
Continue to fix a finite prime-ideal set $S$ containing the primes
above $6$. Let $\mathcal E_X$ contain the $E$-primary, squarefree,
$S$-coprime generators with norm at most $X$, including $1$.
Set $\chi_1(a)=1$. Write $J=\{1,2,4,5\}$.
All divisor sums below run over integral ideal divisors, represented
once by their unique $E$-primary generators. This adds no factor
from the six units. Choose the representatives of the earlier
finite matrix to be $E$-primary as well; its residue symbols depend
on the denominator ideal, so that choice does not change the matrix.

\begin{theorem}[Raw prefix sieve]\label{thm:raw-sieve}
For every $\epsilon>0$, $M,N\ge2$, $j\in J$, and arbitrary complex
coefficients $a_n$ supported on $\mathcal E_N$,
\[
 \sum_{0<N_Eu\le M}
 \left|\sum_{n\in\mathcal E_N}a_n\chi_n(u)^j\right|^2
 \ll_{S,\epsilon}(MN)^\epsilon
 \{M+M^{1/3}N+(MN)^{2/3}\}\sum_n|a_n|^2.          \tag{RS}
\]
If the outer sum is also restricted to $\mathcal E_M$, its bound is
\[
 \ll_{S,\epsilon}(MN)^\epsilon
 \{M+N+(MN)^{2/3}\}\sum_n|a_n|^2.                \tag{SF}
\]
Neither the general Hecke conductor-factorization claim nor any
published large-sieve norm inequality is a premise of this proof.
\end{theorem}
\begin{proof}
Let $B_1^j(M,N)$ and $B_2^j(M,N)$ denote the squared operator norms of
the matrices in (SF) and (RS), respectively, before multiplication
by the displayed bounds. Define these norms for $M,N\ge1$.
Both row and column cutoffs are prefixes.
If $1\le M<2$, there are at most six unit rows, and
$B_i(M,N)\ll N$ by Cauchy--Schwarz.
If $1\le N<2$, only the unit column occurs and $B_i(M,N)\ll M$.
These elementary bounds cover all small cutoffs used below.
Put $B_i=\max_{j\in J}B_i^j$.
Literal inclusion makes both norms monotone in both cutoffs, and
$B_1\le B_2$. Transpose and (R), after splitting both variables into
the two norm classes modulo $4$, give
\[
 B_1(M,N)\ll B_1(N,M).                            \tag{D}
\]
The adjoint conjugates the symbol and sends $j$ to $6-j$, still in
$J$. The finitely many block constants are absolute; no assertion
of dyadic monotonicity is used.

The Gram proof of Theorem~\ref{thm:finite-gram} works also with
prefix columns and every $j\in J$. Each local power has order $6$
or $3$ and is nontrivial. Thus, for columns of norm at most $N$,
the off-diagonal entry is at most $C_FN$, the diagonal is at most
$C_FM$, and there are $O(N)$ columns. The unit column contributes
an $O(M)$ diagonal; its cross entry with a nonunit column follows
from the same nonprincipal completion. Consequently
\[
 B_2(M,N)\ll M+N^2.
\]
Using (D) gives the simultaneous starting hypothesis
\[
 E_2:\quad B_1(M,N)\ll_\epsilon(MN)^\epsilon
 \{M^2+N+(MN)^{2/3}\}.                            \tag{E2}
\]
For example if $M<N$, transpose the preceding Gram bound; if
$M\ge N$, its $M+N^2$ is $O(M^2)$.

We prove the completion needed for the iteration. Keep the fixed radial
majorant $F$ from the previous theorem. For $V\ge1$, fix one norm
class modulo $4$, and take coefficients $\lambda_n$ supported on
$V<N_En\le2V$ in that class. Let $B_3^j(U,V)$ be the supremum of
\[
 \left|\sum_{(n,m)=1}\lambda_n\overline{\lambda_m}
     \sum_{u\in\mathcal O}F(u/\sqrt U)
       \chi_n(u)^j\overline{\chi_m(u)^j}\right|
       \bigg/\sum_n|\lambda_n|^2,                \tag{C0}
\]
and maximize it over the two classes and $J$.
Its columns are dyadic throughout. Its rows are the full lattice,
with no primary or residue-class filter.
For every $a>0$ and $U\ge1/2,V\ge1$ we claim
\[
 B_3(U,V)\ll_{F,a}(2UV)^a
 \left\{U+\frac UV
  \max_{\substack{1\le K\le C(2UV)^aV^2/U}}B_2(K,2V)\right\}.
                                                               \tag{C}
\]
An empty maximum is zero.

Here are the normalizations and the proof. For coprime $n,m$ the
character $\psi=\chi_n^j\overline{\chi_m^j}$ is primitive and
nonprincipal modulo $q=nm$. The complete Gauss calculation in the
preceding theorem also gives
$\tau_\psi(\alpha)=\overline{\psi(\alpha)}\tau_\psi(1)$,
including value zero at nonunit $\alpha$, and
$|\tau_\psi(1)|=\sqrt{N_Eq}$.
The map $(x\bmod n,y\bmod m)\mapsto b=xm+yn\bmod nm$
is a bijection when $(n,m)=1$. It gives
$e(b/(nm))=e(x/n)e(y/m)$ and hence exactly
\[
 \tau_\psi(1)=\chi_n(m)^j\overline{\chi_m(n)^j}
             \tau_{\chi_n^j}(1)\tau_{\overline{\chi_m^j}}(1).
\]
Also
$\tau_{\overline{\chi_m^j}}(1)
=\chi_m(-1)^j\overline{\tau_{\chi_m^j}(1)}$.
At each prime, the local definition gives
$\chi_\pi(-1)=(-1)^{(N_E\pi-1)/6}
=(-1)^{(N_E\pi-1)/2}$.
Multiplying over squarefree $m$ gives
$\chi_m(-1)=(-1)^{(N_Em-1)/2}$.
If $t=(N_En-1)/2\bmod2=(N_Em-1)/2\bmod2$, the reciprocity factor
for the $j$th power is $(-1)^{jt^2}=(-1)^{jt}$.
It cancels the factor $\chi_m(-1)^j=(-1)^{jt}$.
Writing $\gamma_n=\tau_{\chi_n^j}(1)/\sqrt{N_En}$
and setting $\gamma_1=1$, we obtain
\[
 \tau_\psi(1)=\sqrt{N_En\,N_Em}\,
                       \gamma_n\overline{\gamma_m},
 \qquad |\gamma_n|=1.
\]
The same self-dual periodization as before therefore yields
\[
 \sum_u F(u/\sqrt U)\psi(u)
 =\frac{U}{\sqrt{N_En\,N_Em}}\gamma_n\overline{\gamma_m}
   \sum_{\alpha\ne0}\overline{\chi_n(\alpha)^j}\chi_m(\alpha)^j
       \Phi\!\left(\frac{U N_E\alpha}{N_En\,N_Em}\right),          \tag{P}
\]
where radiality defines $\Phi(x)=\widehat F(\sqrt x)$.
This is smooth at $0$, bounded there, and rapidly decreasing at
infinity. Each derivative has the corresponding rapid decrease.

Put $Z=(2UV)^\delta$, for a small fixed $\delta>0$, and first
discard spatial frequencies $N_E\alpha>4ZV^2/U$, the same cutoff
for every pair. The sum of absolute pair coefficients is
$O(V)\sum|\lambda|^2$. Lattice annuli and Schwartz decrease give
the normalized discarded part at most
\[
 C_{F,A}V^2Z^{1-A}.                              \tag{T}
\]
For completeness, when $4ZV^2/U\ge1$, an annulus of norm $K$
has $O(K)$ points and contributes at most
$CUK(UK/(4V^2))^{-A}$; sum its geometric tail.
When $4ZV^2/U<1$, sum from the first nonzero lattice norm;
$U>4ZV^2$ gives the same bound.
For any requested $r>0$, choosing
$A>1+(r+2)/\delta$ makes (T) $O((2UV)^{-r})$,
since $U\ge1/2$ and $V\le2UV$. This error is $O(U)$.

For a fixed $0<\sigma<1$, define
\[
 \widetilde\Phi(s)=\int_0^\infty\Phi(x)x^{s-1}\dd x,
 \qquad \Re s=\sigma.
\]
After $x=e^v$, all derivatives of $e^{\sigma v}\Phi(e^v)$
are integrable. Repeated integration by parts gives
$|\widetilde\Phi(\sigma+it)|\ll_{F,\sigma,A}(1+|t|)^{-A}$.
Fourier inversion of this integrable function is Mellin inversion.
We integrate the whole line; after the spatial cutoff only finitely
many $\alpha,n,m$ remain, and
$\int|\widetilde\Phi(\sigma+it)|\dd t<\infty$.

To remove coprimality after (P), for $\alpha\ne0$ define on
\emph{every} pair the
separable expression
\[
 H(n,m;\alpha,s)=\lambda_n\overline{\lambda_m}
  \gamma_n\overline{\gamma_m}(N_En\,N_Em)^{s-1/2}
  (N_E\alpha)^{-s}
  \overline{\chi_n(\alpha)^j}\chi_m(\alpha)^j.
\]
The integrand is $U^{1-s}H$ on coprime pairs. On other pairs this
is a chosen algebraic extension; no product Gauss identity is
asserted there. The identity
\[
 \sum_{(n,m)=1}H=\sum_d\mu_E(d)\sum_{d\mid n,m}H                 \tag{M}
\]
holds for any such extension. It introduces no imprimitive
Poisson formula.

On a frequency block $K/2<N_E\alpha\le K$, fix $d$ in (M).
Write $n=dn',m=dm'$. Zero extension gives
\[
 \overline{\chi_n(\alpha)^j}\chi_m(\alpha)^j
 =\mathbf1_{(\alpha,d)=1}
             \overline{\chi_{n'}(\alpha)^j}\chi_{m'}(\alpha)^j.
\]
The two free vectors indexed by $n'=n/d$ have entries
\[
 A_{n'}=\lambda_n\gamma_n(N_En)^{\sigma+it-1/2},
 \qquad
 B_{n'}=\overline{\lambda_n}\,\overline{\gamma_n}
                          (N_En)^{\sigma+it-1/2}.
\]
Their norm product is at most
$C V^{2\sigma-1}\sum_{n:\,d\mid n}|\lambda_n|^2$.
Cauchy--Schwarz in the frequency variable, dropping its
coprimality mask only after this nonnegative bound, controls
the row product by
\[
 \bigl(B_2^{6-j}(K,2V)B_2^j(K,2V)\bigr)^{1/2}
                \|A\|_2\|B\|_2
 \le B_2(K,2V)\|A\|_2\|B\|_2.
\]
The residual unit column stays inside this norm. It is not
removed as an $O(U)$ term after dualization.
Every squarefree divisor count is $O_b((N_En)^b)$ for every $b>0$:
the primes with $N_E\pi<2^{1/b}$ are finite, and at each remaining
prime $2\le(N_E\pi)^b$.
Summing over $d$ thus costs only $O_b(V^b)$.
The resulting frequency-block prefactor is exactly
\[
 U^{1-\sigma}V^{2\sigma-1}K^{-\sigma}
 =\frac UV\left(\frac{V^2}{UK}\right)^\sigma.
\]
Its additional factor is at most $(2V^2)^\sigma$ because
$U\ge1/2$ and $K\ge1$. There are $O_\delta(\log(2UV))$ frequency
blocks. Choose $\delta,\sigma,b$ sufficiently small relative
to $a$, absorb this logarithm in $(2UV)^a$, and integrate the
uniform-in-$t$ norm bound on the full Mellin line. Together
with (T) this proves (C). There are no shifted row weights.
If $U<1/2$, the support of $F$ and the vanishing at $u=0$
make the nonprincipal sum zero; the compact range
$1/2\le U<1$ was already included.

We now give the original gcd reduction separately.
Split the column prefix of $B_2(M,N)$ first into the unit
column, dyadic blocks $N_0<N_En\le2N_0$, and norm classes.
Cauchy--Schwarz costs $O(\log(2N))$.
In each nonunit block replace the sharp nonnegative row energy
by its majorant $F(u/\sqrt M)$. For a pair write
$n=dn',m=dm'$ with $d=(n,m)$, so $(n',m')=1$.
Squarefreeness and multiplicativity give the original row mask
$\mathbf1_{(u,d)=1}$. Remove it by
$\sum_{k\mid d}\mu_E(k)\mathbf1_{k\mid u}$ and set $u=kv$.
The coefficient phase $\chi_{n'}(k)^j$ is retained in the free
coefficient vector. Radiality leaves
\[
 U=M/N_Ek,\qquad V=N_0/N_Ed,\qquad N_Ek\le N_Ed.                \tag{G}
\]
Residual columns stay in one norm class and in the same dyadic
block $(V,2V]$. The original diagonal $n=m=d$ is bounded by
$CM\sum|\lambda_n|^2$ \emph{before} this expansion.
Equivalently it is the only principal residual pair.
If $V<1$, the residual block is empty or contains only the unit
generator, so it was already part of that diagonal.
For the other pairs apply (C).
The total coefficient-norm multiplicity from $d$ and $k\mid d$
is at most $\sum_n|\lambda_n|^2\tau(n)^2$, hence $O_b(N^b)$.
This proves the stated gcd reduction with no missing
principal frequency.

Suppose for some fixed $\alpha>4/3$, simultaneously for $J$
and every $\epsilon>0$, we have
\[
 E_\alpha:\quad B_1(X,N)\ll_\epsilon(XN)^\epsilon
                   \{X^\alpha+N+(XN)^{2/3}\}.                \tag{E}
\]
An elementary all-row transfer follows. Factor each nonzero row
uniquely as
$u=\varepsilon s\,u_1u_2^2u_3^3u_4^4u_5^5u_6^6$,
where $s$ has only prime factors in $S$,
$u_1,\ldots,u_5$ are squarefree and $S$-coprime,
and $u_6$ is arbitrary and $S$-coprime.
Choose their $E$-primary generators; the remaining factor
$\varepsilon$ is one of six units.
There are $O_S((1+\log M)^{|S|})$ possibilities for $s$,
and $O((1+\log M)^6)$ dyadic norm choices.
In a block with lengths $X_i\ge1$ let
$\ell=1$ if $X_1\ge X_2$ and $\ell=2$ otherwise, and put
$R=\prod_{i\ne\ell}X_i$. From
$X_1X_2^2\cdots X_6^6\ll M$ one has
\[
 R\ll\min\{M^{1/3},(M/X_\ell)^{2/3}\},
 \qquad RX_\ell\ll M.                            \tag{F}
\]
Indeed for $\ell=1$, $R^3\le
M(X_2/X_1)X_4^{-1}X_5^{-2}X_6^{-3}$, and
$R^{3/2}\le X_2^2X_3^3\cdots X_6^6\le M/X_1$.
For $\ell=2$ the corresponding factors are
$(X_1^2/X_2^2)X_4^{-1}X_5^{-2}X_6^{-3}$ and
$X_1^{3/2}\le X_1X_2$.
The last inequality in (F) is just comparison of exponents.
Fix all other factors and absorb their symbols into the
coefficients, decreasing their moduli when a gcd gives zero.
The selected symbol has exponent $j\ell\bmod6$, always in $J$.
Apply (E) to that squarefree row and count the other factors,
using (F). This gives
\[
 B_2(M,N)\ll_\epsilon(MN)^\epsilon
                 \{M^\alpha+M^{1/3}N+(MN)^{2/3}\}.             \tag{A}
\]
No assertion about only squarefree all-rows has been substituted here.

Insert (A) into (C). For $K\le V$,
$(KV)^{2/3}\le K^{1/3}V$; for $K\ge V$ it is at most
$K^\alpha$. The cap in (C) therefore gives, with arbitrarily
small total losses,
\[
 B_3(U,V)\ll_\epsilon(UV)^\epsilon
           \{U+V^{2\alpha-1}U^{1-\alpha}+(UV)^{2/3}\}.
\]
Retain the scales (G) until this substitution.
The dangerous middle term becomes
\[
 N_0^{2\alpha-1}M^{1-\alpha}
       (N_Ek)^{\alpha-1}(N_Ed)^{1-2\alpha}
 \le N_0^{2\alpha-1}M^{1-\alpha},
\]
because $k\mid d$ and $\alpha>1$.
The other terms are at most $M$ and $(MN_0)^{2/3}$.
Only now use $N_0\le N$ and sum the gcd multiplicities:
\[
 B_2(M,N)\ll_\epsilon(MN)^\epsilon
       \{M+N^{2\alpha-1}M^{1-\alpha}+(MN)^{2/3}\}.             \tag{B}
\]
All losses can be made smaller than a prescribed $\epsilon$.
For example choose each auxiliary loss smaller than
$\epsilon/[100(1+\alpha)]$: every dual cutoff is bounded
by a fixed polynomial in $MN$, and divisor and dyadic factors
have arbitrarily small power loss. Constants depend on the
fixed $S,\alpha,\epsilon,F$.

Set $\tau=2-3/(3\alpha-1)$,
$\beta=2-2/(3\alpha-1)$, and $T=N^\tau$.
For $M\ge T$, the middle term in (B) is at most $(MN)^{2/3}$.
For $M<T$ use literal prefix monotonicity of $B_1$ at $T$:
both $N^{2\alpha-1}T^{1-\alpha}$ and $(TN)^{2/3}$ equal
$N^\beta$, and $T\le N^\beta$.
After decreasing the auxiliary loss to account for $TN\le N^3$,
this proves
$B_1(M,N)\ll_\epsilon(MN)^\epsilon
\{M+N^\beta+(MN)^{2/3}\}$.
Transpose by (D) to obtain $E_\beta$.
Starting with (E2), the finite recurrence
$\alpha\mapsto2-2/(3\alpha-1)$ decreases from $2$
to $4/3$ from above. For every prescribed $d>0$ a
\emph{finite} number of steps gives $E_\alpha$ with
$4/3<\alpha<4/3+d$. No uniform constant at an infinite
iteration limit is claimed.

For $M\ge N$, use that exponent in (B) and note
\[
 \frac{N^{2\alpha-1}M^{1-\alpha}}{(MN)^{2/3}}
 \le N^{\alpha-4/3}.
\]
Choose $d$ and all auxiliary losses smaller than the requested
final power loss. This bounds $B_2(M,N)$ by
$(MN)^\epsilon\{M+(MN)^{2/3}\}$ in that regime.
For squarefree rows transpose the other regime.
This proves (SF) for all $M,N$.
Finally repeat the powerfree transfer with (SF) itself.
Now its first term is $RX_\ell\le CM$, rather than
$RX_\ell^\alpha$; the other two terms are still bounded
by (F). This proves (RS) in \emph{both} regimes.
\end{proof}

\begin{corollary}[A stronger finite extractor cost]\label{cor:raw-extractor}
For the same matrix $M$ and residual condition as in
Theorem~\ref{thm:finite-gram}, with $N_D>0$ and $H\ge D\ge2$,
every $\epsilon>0$
gives
\[
 \|c\|_2^2\gg_{S,\epsilon}
 (1-\eta)^2(HD)^{-\epsilon}
                  \frac{D}{H+(HD)^{2/3}}.                     \tag{EX}
\]
In particular for any fixed $0\le\vartheta<1$,
$H=D^{1+\vartheta}$, and fixed $\eta<1$, its lower bound is
$D^{-(1+2\vartheta)/3-\epsilon}$ for every $\epsilon>0$,
up to a positive constant. Thus a dense universal extractor
of squared norm $D^{-1+o(1)}$ is impossible in that regime.
The exclusion also holds uniformly if
$0\le\vartheta(D)\le1-\kappa$ for a fixed $\kappa>0$.
\end{corollary}
\begin{proof}
Apply (RS) with $N=2D$ and $M=H$ to the adjoint Gram energy.
Conjugation is covered by $J$. Since $H\ge D$,
$H^{1/3}D\le(HD)^{2/3}$.
Cauchy--Schwarz on $\langle\mathbf1,Mc\rangle$ and
$N_D\asymp_SD$ gives (EX).
The asymptotic count gives $N_D\ge c_SD$ for sufficiently large $D$.
On the remaining bounded range, $N_D>0$ gives $N_D\ge1$,
so decreasing $c_S$ extends this lower bound to every nonempty
column family with $D\ge2$.
For $H=D^{1+\vartheta}$ with $\vartheta<1$,
$(HD)^{2/3}$ dominates $H$ and yields the exponent stated,
by taking the loss in (EX) to be $\epsilon_0=\epsilon/3$,
since $HD=D^{2+\vartheta}\le D^3$.
For the uniform variant take $\epsilon_0=\kappa/10$.
The final exponent is then at most
$1-2\kappa/3+3\kappa/10=1-11\kappa/30$,
a fixed gap from $1$.
\end{proof}

This result concerns a universal residual condition on the full
finite denominator family. It neither disproves the estimate for
one specified M\"obius/Hecke vector nor transfers that estimate
to the centered rational $\Lambda$-correlation target.
The analytic norm bound is reconstructed here from raw symbols;
we make no claim of priority for the known squarefree-sieve
shape. No zeta zero is excluded by (RS) or (EX).

\part{Semilocal Weil forms and exact zero detection}

\section{A positivity obstruction for the actual semilocal operator}
Use the concrete form of \cite[\S3, eqs.~(3.2)--(3.5)]{OperatorSource},
reconstructed from \cite[Proposition 2.1 and Corollary 2.4]{ConnesConsani}.
In centered coordinates $\ell=\log u$ put
$\mathcal H_\lambda=L^2([-\log\lambda,\log\lambda],\dd\ell)$;
all functions are zero extended. The inner product is
antilinear in its first argument. Let
$C_{f,g}(r)=\int\overline{f(\ell-r)}g(\ell)\dd\ell$.
The form is
\begin{align*}
 q_\lambda(f,g)={}&
 \int_{\R}\overline{\widehat f(t)}\,\widehat g(t)
       \frac{a(t)}{2\pi}\dd t\\
 &+\overline{\widehat f(i/2)}\,\widehat g(-i/2)
       +\overline{\widehat f(-i/2)}\,\widehat g(i/2)\\
 &-\sum_{2\le n\le\lambda^2}\frac{\Lambda(n)}{\sqrt n}
       \{C_{f,g}(\log n)+C_{f,g}(-\log n)\},
\end{align*}
where $a(t)=\Re(\Gamma'/\Gamma)(1/4+it/2)-\log\pi$.
Its domain consists of functions in $\mathcal H_\lambda$
whose extensions by zero outside the displayed interval satisfy
$\int|\widehat f(t)|^2(1+\log(1+t^2))\dd t<\infty$.
The multiplier is bounded below and grows logarithmically,
by the digamma asymptotic \cite[eq.~(5.11.2)]{DLMFGamma}.
After a scalar shift its nonnegative weighted Fourier norm
is complete, by Plancherel and preservation of the support
condition under $L^2$ limits. The compact-support pole
evaluations and finite sum of translations are bounded.
Their addition preserves closedness and lower
semiboundedness. Denote the associated selfadjoint operator
by $A_\lambda$. Interior $C_c^\infty$ functions are in its
form domain. No hypothesis on its ground state is used.

\begin{theorem}[Failure of positivity preservation]\label{thm:semigroup}
For every $\lambda\ge2$ there are nonzero, nonnegative,
disjointly supported smooth form-domain elements $f,g$
with $q_\lambda(f,g)>0$. For those functions there is $t_0>0$
such that $\langle f,e^{-tA_\lambda}g\rangle<0$ whenever
$0<t<t_0$. Thus the semigroup is not positivity-preserving.
\end{theorem}
\begin{proof}
The digamma identity \cite[eq.~(5.9.16), $\Re z>0$]{DLMFGamma},
with integration variable $v=2r$, gives
\[
 \Re\frac{\Gamma'}{\Gamma}(1/4+it/2)
 =-\gamma_E+\int_0^\infty
 \frac{2e^{-2r}-2e^{-r/2}\cos(tr)}{1-e^{-2r}}\dd r.
\]
For disjoint compact supports, scalar terms vanish by
$\langle f,g\rangle=0$ and correlations vanish near $r=0$.
Plancherel sends the cosine multiplier to
$\{C_{f,g}(r)+C_{f,g}(-r)\}/2$.
Truncate first, then pass to the limit using the rapid
Fourier decay and separated supports. The resulting
off-diagonal physical kernel is
\[
 -k(|\ell-\ell'|),\qquad
 k(r)=\frac{e^{-r/2}}{1-e^{-2r}}
     =\frac{e^{r/2}}{e^r-e^{-r}}.
\]
The pole term gives $2\cosh((\ell-\ell')/2)$,
with the difference of coordinates.

Put $a_*=\tfrac12\log(3/2)$ and choose $\eta>0$ with
\[
 2\eta<\min\{\log(3/2)-\log(7/5),
                         \log2-\log(3/2)\}.
\]
Take nonnegative nonzero smooth bumps supported inside
$[-a_*-\eta,-a_*+\eta]$ and
$[a_*-\eta,a_*+\eta]$. They lie in the Hilbert interval
for every $\lambda\ge2$, and all cross distances satisfy
$\log(7/5)<|\ell-\ell'|<\log2$.
Every prime-power translation vanishes, including the
first at $\log2$. For a distance $r$ in this interval,
writing $x=e^r$ gives
\[
 2\cosh(r/2)-k(r)
 =\frac{x^3-x-1}{\sqrt x(x^2-1)}>0.
\]
At $x=7/5$ the numerator is $43/125$, and its derivative
$3x^2-1$ is positive for $x\ge7/5$.
Integrating against the two bumps proves $q_\lambda(f,g)>0$.

Suppose the semigroup preserves positivity for every $t>0$.
Shift $T_\lambda=A_\lambda+c\ge0$; its semigroup also
preserves positivity. Disjoint supports then give
$t^{-1}\langle f,(1-e^{-tT_\lambda})g\rangle\le0$.
The spectral multiplier $(1-e^{-tx})/(tx)$ is between
zero and one for $x\ge0$ and tends to one. Dominated
convergence on $T_\lambda^{1/2}f,T_\lambda^{1/2}g$ gives
\[
 \lim_{t\downarrow0}t^{-1}\langle
 f,(1-e^{-tT_\lambda})g\rangle
 =q_\lambda(f,g)+c\langle f,g\rangle=q_\lambda(f,g)>0,
\]
a contradiction.
In fact the same positive form limit makes
$\langle f,e^{-tT_\lambda}g\rangle<0$ for every sufficiently
small positive $t$, without assuming positivity preservation.
Multiplication by $e^{ct}>0$ gives the stated stronger conclusion.
\end{proof}

This excludes positivity improvement of that semigroup
as a justification for a simple positive even ground vector.
It proves neither an odd ground state nor degeneracy;
another simplicity or parity argument may still exist.
It is not an RH counterexample. Positivity of a quadratic
form and preservation of positivity by its semigroup are
different properties.

\section{Removing the pole term gives an irreducible jump form}
Keep the same coordinates and measure; put $L=\log\lambda$,
$c_\lambda(\ell)=\cosh(\ell/2)$ and
$s_\lambda(\ell)=\sinh(\ell/2)$. The pole operator is exactly
\[
 P_\lambda=2c_\lambda\langle c_\lambda,\cdot\rangle
              -2s_\lambda\langle s_\lambda,\cdot\rangle,
 \qquad B_\lambda=A_\lambda-P_\lambda.
\]
Its kernel relative to $\dd\ell'$ is
$2\cosh((\ell-\ell')/2)$. Under the unitary change of
variables $\ell=2x$ to Lebesgue measure $\dd x$, this kernel
would instead be $4\cosh(x-x')$; no such rescaling is used below.

\begin{theorem}[Pole-free ground vector]\label{thm:pole-free}
For every fixed $\lambda>1$, $B_\lambda$ has compact resolvent
and a simple lowest eigenvalue. Its normalized ground vector
can be chosen real, strictly positive almost everywhere on
$(-L,L)$, and even. No ground-state assumption on $A_\lambda$
is required.
\end{theorem}
\begin{proof}
The digamma integral used above gives exactly
\[
 a(t)-a(0)=2\int_0^\infty k(r)(1-\cos tr)\dd r,
 \qquad k(r)=\frac{e^{r/2}}{e^r-e^{-r}}>0.
\]
Write $(\tau_r f)(\ell)=f(\ell-r)$ for translations of the
zero extension on $\R$ and put
\[
 b_\lambda^{\rm shift}
   =2\sum_{2\le n\le\lambda^2}\frac{\Lambda(n)}{\sqrt n}-a(0).
\]
Tonelli and Plancherel, followed by the identity
$\|f-\tau_r f\|_2^2=2\|f\|_2^2
-2\Re\langle\tau_r f,f\rangle$, prove on the form domain
\begin{align*}
 \mathcal E_\lambda(f)
 &:=q_{B_\lambda}(f)+b_\lambda^{\rm shift}\|f\|_2^2\\
 &=\int_0^\infty k(r)\|f-\tau_r f\|_2^2\dd r
   +\sum_{2\le n\le\lambda^2}
      \frac{\Lambda(n)}{\sqrt n}
                      \|f-\tau_{\log n}f\|_2^2\ge0.
\end{align*}
All prime powers, including an endpoint $n=\lambda^2$,
are retained. Closedness follows from the weighted Fourier
norm described above and boundedness of the finite translations.
This proves a lower bound after the displayed shift, not
nonnegativity of the unshifted $B_\lambda$.

The form-domain inclusion in $L^2([-L,L])$ is compact. Indeed, on a
bounded form-norm set, Fourier mass outside $[-R,R]$ is
$O(1/\log R)$ as $R\to\infty$. Fourier projection to this
band, applied to functions supported on $[-L,L]$, has a
square-integrable kernel on $[-L,L]\times\R$ and hence is
compact. Restrict their output to $[-L,L]$. These compact maps
converge in norm on the form-norm
unit ball to its inclusion in $L^2$. Compactness and the
spectral theorem give a lowest eigenvalue of finite multiplicity.

The modulus contraction
$||z|-|w||\le|z-w|$ implies
$\mathcal E_\lambda(|f|)\le\mathcal E_\lambda(f)$ and
preservation of the form domain. For nonzero real
$f,g\ge0$ in that domain with disjoint supports, polarization gives
\[
 \mathcal E_{\lambda,{\rm arch}}(f,g)
 =-\int_0^\infty k(r)
       \bigl(\langle\tau_r f,g\rangle
                  +\langle f,\tau_r g\rangle\bigr)\dd r<0.
\]
The integral is finite by Cauchy--Schwarz in the jump norm.
Strictness follows because the unweighted integral of the
nonnegative parenthesis over $r>0$ is
$\|f\|_1\|g\|_1>0$, and $k(r)>0$ everywhere.
The finite prime terms are nonpositive cross terms and cannot
cancel this strict inequality.

The operator is real. Modulus contraction and minimization of
the Rayleigh quotient therefore give a real nonnegative
ground vector $\phi$. Every real ground vector $v$ has a fixed
sign: otherwise its nonzero positive and negative parts are
in the form domain and
\[
 \mathcal E_\lambda(v)-\mathcal E_\lambda(|v|)
       =-4\mathcal E_\lambda(v_+,v_-)>0,
\]
contrary to ground minimality. If the ground space had dimension
at least two, choose a real ground vector orthogonal to $\phi$
and change its sign to make it nonnegative. The two vectors
then have disjoint supports, but their cross form is zero by
the eigenvalue equation, contradicting the strict inequality.
The ground eigenvalue is simple.

For completeness, strict positivity follows directly from the
form. Set $T=B_\lambda+b_\lambda^{\rm shift}I\ge0$.
For $a>0$ and $F\ge0$, the functional
\[
 J(v)=\mathcal E_\lambda(v)+a\|v\|_2^2
                              -2\Re\langle F,v\rangle
\]
has unique minimizer $u=(T+a)^{-1}F$.
Modulus contraction gives $J(|u|)\le J(u)$; uniqueness gives
$u\ge0$. Every such resolvent preserves positivity, and so
does $e^{-tT}$ by the strong spectral Euler-product limit
$e^{-tT}=\lim_{n\to\infty}(I+tT/n)^{-n}$.
Write $T\phi=\mu_0\phi$, $\mu_0=\min\sigma(T)$.
If $E_0=\{\phi=0\}$ had positive measure, then for every
$g\ge0$ supported in $E_0$,
\[
 \langle\phi,e^{-tT}g\rangle
     =e^{-t\mu_0}\langle\phi,g\rangle=0.
\]
Positivity forces $e^{-tT}g$ to remain supported in $E_0$.
Linearity and selfadjointness make $L^2(E_0)$ a reducing
subspace for every $e^{-tT}$. Its multiplication projection
then commutes with $(T+a)^{-1}$ by the Laplace formula
and with $T^{1/2}$ by the spectral calculus. It preserves
the form domain and makes the cross form between the two
parts zero. Split a smooth flat bump, positive on $(-L,L)$,
by this projection. Both parts are nonzero, nonnegative
form-domain functions with disjoint supports, contradicting
the strictly negative cross form. Hence $\phi>0$ almost everywhere.
Reflection $\ell\mapsto-\ell$ preserves the form. Simplicity,
strict positivity and reflection squared equal to one give
$\phi(-\ell)=\phi(\ell)$.
\end{proof}

\begin{proposition}[Exact pole restoration]\label{prop:pole-restoration}
Let $R_B(z)=(B_\lambda-z)^{-1}$ for $z\in\rho(B_\lambda)$,
$U(a,b)=\sqrt2(ac_\lambda+bs_\lambda)$ and
$J=\operatorname{diag}(1,-1)$. Whenever the displayed
matrix is invertible,
\[
 (A_\lambda-z)^{-1}
 =R_B(z)-R_B(z)U\bigl(J+U^*R_B(z)U\bigr)^{-1}U^*R_B(z).
\]
The matrix is diagonal by parity; its two scalar factors are
\[
 D_e(z)=1+2\langle c_\lambda,R_B(z)c_\lambda\rangle,
 \qquad D_o(z)=1-2\langle s_\lambda,R_B(z)s_\lambda\rangle.
\]
Writing $b_\lambda=\min\sigma(B_\lambda)$, $D_e(z)>0$ for
real $z<b_\lambda$. Exactly one eigenvalue of $A_\lambda$
lies below $b_\lambda$ if and only if
\[
 \lim_{z\uparrow b_\lambda}
         2\langle s_\lambda,R_B(z)s_\lambda\rangle>1;
\]
that eigenvalue is odd and simple. Finally,
\[
 \|P_\lambda\|=2(L+\sinh L)\asymp\lambda
                       \quad(\lambda\to\infty).
\]
\end{proposition}
\begin{proof}
Since $P_\lambda=UJU^*$ and $J^{-1}=J$, direct
multiplication proves the finite-rank resolvent formula;
its matrix is invertible exactly when
$z\in\rho(A_\lambda)\cap\rho(B_\lambda)$.
Parity makes the cross matrix elements vanish.
For real $z<b_\lambda$ the two diagonal resolvent elements
are positive and strictly increasing, with derivative
$\|(B_\lambda-z)^{-1}h\|_2^2$ for $h=c_\lambda,s_\lambda$.
They tend to zero as $z\to-\infty$. Thus the even factor
cannot vanish, and the odd factor has exactly one zero
below $b_\lambda$ precisely under the stated strict inequality.
The rank-one eigenvalue equation gives an odd eigenvector
proportional to $R_B(z)s_\lambda$ and multiplicity one.
Since the ground vector of $B_\lambda$ is simple and even,
the limiting odd resolvent element at $b_\lambda$ is finite.
Finally $c_\lambda\perp s_\lambda$ and direct integration gives
$\|c_\lambda\|_2^2=L+\sinh L$,
$\|s_\lambda\|_2^2=\sinh L-L$.
\end{proof}

These are fixed-$\lambda$ operator results. The even pole term
raises the even sector and the odd term lowers the odd sector.
The theorem therefore supplies no even-ground-state conclusion
for $A_\lambda$. The perturbation grows with $\lambda$;
the resolvent formula supplies neither a uniform cofinal limit
nor identification with the arithmetical packet transform.

\section{The full negative ground value controls packet energy}
Write
\[
 \alpha_\lambda=\inf\sigma(A_\lambda),\qquad
 \nu_\lambda=\max\{0,-\alpha_\lambda\},\qquad \lambda>1.
\]
These quantities refer to the full operator, including its pole term.
They require no ground-vector parity, simplicity, or Fourier-divisor
assumption. Bounded perturbation of the preceding compact form
embedding also gives compact resolvent for $A_\lambda$.
Set
\[
 D_h=\frac1{2\pi}\int_{\R}a(t)L(it)\dd t
          +2L(1/2)-\frac{2\log2}{\sqrt2}K(\log2).
\]
The integral is finite since $\widehat h$ is rapidly decreasing.

\begin{theorem}[Negative ground value and packet compression]
\label{thm:negative-ground}
The function $\lambda\mapsto\nu_\lambda$ is nondecreasing.
For $y>1/2$ and $\log\lambda\ge y+1/2$, the compression of
the full form to the orthonormal packets
$h_y^R(\ell)=h(\ell-y)$ and $h_y^L(\ell)=h(\ell+y)$ is
\[
 \begin{pmatrix}D_h&Q(y)\\Q(y)&D_h\end{pmatrix}.
\]
In particular
\[
 (|Q(y)|-D_h)_+\le\nu_{\exp(y+1/2)}.
 \tag{NG}
\]
For $Y_0>1/2$ define
\[
 \mathcal N_h(Y)=\int_{Y_0}^{Y}(|Q(y)|-D_h)_+^2\dd y .
\]
Then, with all multiplicities included in $B$,
\[
 \limsup_{Y\to\infty}\frac{\log(1+\mathcal N_h(Y))}{Y}=4B,
 \qquad
 2B\le
 \limsup_{\lambda\to\infty}
       \frac{\log(1+\nu_\lambda)}{\log\lambda}\le1.
 \tag{NR}
\]
For every fixed $\eta\ge0$, an entire-tail bound
$\nu_\lambda=O(\lambda^\eta)$ implies
$\Re\rho\le1/2+\eta/2$ for every nontrivial zero.
Consequently
\[
 \forall\epsilon>0\ \exists C_\epsilon,\lambda_\epsilon\
 \forall\lambda\ge\lambda_\epsilon:\quad
 \nu_\lambda\le C_\epsilon\lambda^\epsilon
 \quad\Longrightarrow\quad\mathrm{RH}.
 \tag{NP}
\]
The hypothesis in (NP) is not asserted to hold.
\end{theorem}
\begin{proof}
If $1<\lambda_1<\lambda_2$, zero extension embeds the smaller
form domain into the larger one. Every new prime translation has
$\log n>2\log\lambda_1$ and hence zero correlation on the smaller
interval. The multiplier and pole terms are unchanged.
The two forms therefore agree there. The variational principle
gives $\alpha_{\lambda_2}\le\alpha_{\lambda_1}$ and the stated
monotonicity.

Translations preserve the multiplier and pole-difference kernels.
The self-correlation of $h$ is $K$ and has support $[-1,1]$;
once $\lambda\ge\sqrt e$, only $n=2$ can contribute to the
prime part of a translated diagonal. Both diagonals are therefore
exactly $D_h$. For $y>1/2$ the two supports are disjoint and the
packets are orthonormal. Their prime cross-term is
$-S_h(2y)$: every contributing $n$ has
$\log n<2y+1\le2\log\lambda$; equality contributes zero.
The pole cross-term is $L(1/2)(e^y+e^{-y})$.
The separated-support archimedean kernel gives, by its positive
geometric expansion and Tonelli,
\[
 -\int_0^\infty k(r)K(2y-r)\dd r
 =-\sum_{j\ge0}L(2j+1/2)e^{-(4j+1)y}
 =-L(1/2)e^{-y}-T_h(y).
\]
Thus the complete cross-term is
$e^yL(1/2)-S_h(2y)-T_h(y)=Q(y)$ by
Lemma~\ref{lem:prime}. It is real. The smaller compression
eigenvalue is $D_h-|Q(y)|$, and the full variational infimum
cannot exceed it. This proves (NG).

Put $n(y)=(|Q(y)|-D_h)_+$. Regardless of the sign of $D_h$,
$n(y)\le |Q(y)|+|D_h|$ and $|Q(y)|\le n(y)+|D_h|$.
Hence
\[
 G(Y)\le2\mathcal N_h(Y)+2D_h^2(Y-Y_0),\qquad
 \mathcal N_h(Y)\le2G(Y)+2D_h^2(Y-Y_0).
\]
Theorem~\ref{thm:l2} proves the first equality in (NR).
Monotonicity and (NG) also give
$\mathcal N_h(Y)\le(Y-Y_0)\nu_{\exp(Y+1/2)}^2$,
and therefore
\[
 1+\mathcal N_h(Y)\le(1+Y)(1+\nu_{\exp(Y+1/2)})^2.
\]
Taking logarithms proves its lower growth inequality, including
parameters at which $\nu_\lambda=0$.

For the upper inequality, the multiplier is at least $a(0)$,
each prime correlation has magnitude at most $\|f\|_2^2$,
and the negative pole part is $-2|s_\lambda\rangle
\langle s_\lambda|$. Consequently
\[
 q_\lambda(f,f)\ge
 \left[a(0)-2\sum_{n\le\lambda^2}\frac{\Lambda(n)}{\sqrt n}
                -2(\sinh(\log\lambda)-\log\lambda)\right]\|f\|_2^2.
\]
Partial summation of the already proved Chebyshev bound
$\psi(X)=O(X)$ gives
$\sum_{n\le\lambda^2}\Lambda(n)/\sqrt n=O(\lambda)$.
Thus $\nu_\lambda=O(\lambda)$.
Finally (NG) with the exact parameter $\lambda=e^{y+1/2}$
turns an entire-tail $O(\lambda^\eta)$ estimate into
$Q(y)=O(e^{\eta y})$. Lemma~\ref{lem:strip} proves the
claimed zero bound. Apply it for every $\epsilon>0$ in (NP);
functional-equation symmetry then gives RH.
\end{proof}

\begin{corollary}[A controlled parameter exhaustion]
\label{cor:ground-exhaustion}
Let $\lambda_j$ increase to infinity and suppose
$\log\lambda_{j+1}\le C\log\lambda_j$ eventually for some fixed
$C<\infty$. If, for every $\epsilon>0$,
$\nu_{\lambda_j}\le C_\epsilon\lambda_j^\epsilon$
eventually, then RH holds.
\end{corollary}
\begin{proof}
For $\lambda_j\le\lambda<\lambda_{j+1}$, monotonicity and the
estimate with exponent $\epsilon/C$ give
$\nu_\lambda\le C_{\epsilon/C}\lambda_{j+1}^{\epsilon/C}
\le C_{\epsilon/C}\lambda^\epsilon$.
This supplies the whole tail in (NP), with explicit control
of the gaps. Arbitrary cofinality alone is not used.
\end{proof}

The unconditional bound here is $O(\lambda)$, not the
subpolynomial hypothesis. The theorem replaces a divisor-identification
obligation by a precise lower-bound obligation on the full form;
it does not establish that lower bound.

\section{The natural positive multiplier measures have a different limit}
For $\lambda\ge\sqrt e$, set
\[
 A_{\rm pr}(\lambda)=\sum_{n\le\lambda^2}\frac{\Lambda(n)}{\sqrt n},
 \quad
 b_\lambda(t)=a(t)-a(0)+
     2\sum_{n\le\lambda^2}\frac{\Lambda(n)}{\sqrt n}
                          (1-\cos(t\log n)),
 \quad
 d\mu_\lambda(t)=\frac{b_\lambda(t)L(it)}{2\pi}\dd t .
\]
Here $b_\lambda\ge0$ is the multiplier of the shifted
pole-free form $\mathcal E_\lambda$, not of unshifted $B_\lambda$.
The measures $\mu_\lambda$ are finite and positive.

\begin{theorem}[Mass and the normalized multiplier limit]
\label{thm:multiplier-limit}
There is a fixed real $C_h$ such that
\[
 \mu_\lambda(\R)=2A_{\rm pr}(\lambda)+C_h,\qquad
 A_{\rm pr}(\lambda)\longrightarrow\infty .
\]
The measures $v_\lambda=\mu_\lambda/(2A_{\rm pr}(\lambda))$
have uniformly vanishing tails with every fixed polynomial weight:
for each $m\ge0$,
\[
 \lim_{R\to\infty}\sup_{\lambda\ge\sqrt e}
   \int_{|t|>R}(1+|t|)^m\dd v_\lambda(t)=0.
\]
They converge weakly to
$L(it)\dd t/(2\pi)$, a probability measure.
Their finite-parameter masses are
$1+C_h/(2A_{\rm pr}(\lambda))$ and need not equal one.
Their Fourier
kernels converge locally uniformly to $K(2y)$.
For $1/2<y\le\log\lambda-1/2$, exactly
\[
 Q(y)=\int_\R e^{2ity}\dd\mu_\lambda(t)
                         +L(1/2)(e^y+e^{-y}).
 \tag{MNL}
\]
In particular the normalized positive limit is the packet
autocorrelation, not the arithmetical interaction.
\end{theorem}
\begin{proof}
Fourier inversion gives $\int L(it)\dd t/(2\pi)=K(0)=1$
and $\int\cos(t\log n)L(it)\dd t/(2\pi)=K(\log n)$.
Only $n=2$ contributes to the latter sum once its cutoff
is at least $e$. Hence the exact constant is
\[
 C_h=\frac1{2\pi}\int(a(t)-a(0))L(it)\dd t
                       -\frac{2\log2}{\sqrt2}K(\log2).
\]
For completeness, $\sum_p1/p=\infty$ follows because the
Euler product over $p\le X$ is at least the harmonic sum
over integers at most $X$, whereas its logarithm is
$\sum_{p\le X}1/p+O(1)$.
As $(\log p)/\sqrt p\ge1/p$ for primes $p\ge3$, this implies
$A_{\rm pr}(\lambda)\to\infty$.

Since $a(t)-a(0)\ge0$ and $0\le1-\cos\le2$,
\[
 0\le\frac{b_\lambda(t)}{2A_{\rm pr}(\lambda)}
 \le2+\frac{a(t)-a(0)}{2A_{\rm pr}(\sqrt e)}.
\]
The right side grows only logarithmically; multiplication
by $L(it)=|\widehat h(t)|^2$ is integrable with every
polynomial weight. This proves the uniform tail assertion.

For a bounded continuous $F$, let
\[
 \Theta_F(r)=\frac1{2\pi}\int_\R F(t)L(it)\cos(tr)\dd t.
\]
It tends to zero as $r\to\infty$ by the Riemann--Lebesgue
lemma: approximate the integrable function $F(t)L(it)$ in
$L^1$ by a smooth compactly supported function, integrate
the latter by parts, and bound the approximation error by
its $L^1$ norm. Expansion of $b_\lambda$ gives
\[
 \int F\dd v_\lambda
 =\frac1{2\pi}\int F(t)L(it)\dd t+O_F(A_{\rm pr}(\lambda)^{-1})
 -\frac1{A_{\rm pr}(\lambda)}
     \sum_{n\le\lambda^2}\frac{\Lambda(n)}{\sqrt n}
                                  \Theta_F(\log n).
\]
The last term tends to zero: split at any fixed $n_0$,
let $\lambda\to\infty$ for the finite head, and then use
$\sup_{n>n_0}|\Theta_F(\log n)|\to0$ for the tail.
This proves the weak convergence.

More explicitly, for every fixed $R$ and sufficiently large
$\lambda$, Fourier inversion yields, uniformly for $|y|\le R$,
\[
 \int e^{2ity}\dd\mu_\lambda(t)
 =2A_{\rm pr}(\lambda)K(2y)
   +\frac1{2\pi}\int(a(t)-a(0))L(it)e^{2ity}\dd t
   -\sum_{n\ge2}\frac{\Lambda(n)}{\sqrt n}
                  [K(2y-\log n)+K(2y+\log n)] .
\]
The last sum involves only $n\le e^{2R+1}$, so its norm
on that compact is bounded independently of $\lambda$.
Division by $2A_{\rm pr}(\lambda)$ proves local uniform
kernel convergence to $K(2y)$.
Finally, when $y>1/2$, the shift contribution
$(2A_{\rm pr}(\lambda)-a(0))K(2y)$ vanishes.
The multiplier integral is then exactly
$q_{B_\lambda}(h_y^R,h_y^L)$.
Restoring the complete pole cross-term and applying the
preceding compression calculation proves (MNL).
\end{proof}

Thus uniform weighted tails can be proved for this explicit normalized
family, but the proved limiting kernel vanishes for $y>1/2$.
Recovering $Q$ requires both the growing normalization and the pole
correction in (MNL). Neither operation identifies the normalized
positive limit with $Q$.

\section{Continuous prime compensation is not small in operator norm}
Write $L=\log\lambda$ and $X=e^{2L}$. With zero extension from
$[-L,L]$, define the bounded selfadjoint operators
\[
 \Pi_L=\sum_{2\le n\le X}\frac{\Lambda(n)}{\sqrt n}
                          (\tau_{\log n}+\tau_{-\log n}),\qquad
 C_L f=\int_1^X x^{-1/2}
                       (\tau_{\log x}f+\tau_{-\log x}f)\dd x,
 \quad \mathfrak D_L=\Pi_L-C_L .
\]
Here the integral is an $L^2$ integral for each $f$;
strong continuity of translations and the integrable weight
make it well defined. Translations are compressed back to
the interval. The difference $\mathfrak D_L$ denotes an
operator, not the earlier scalar $D_h$.

\begin{theorem}[Sharp operator-norm obstruction to prime compensation]
\label{thm:prime-norm}
There is $c>0$ such that, for every sufficiently large $L$,
\[
 \|\mathfrak D_L\|\ge
 2\sum_{n\le e^{2L}}\frac{\Lambda(n)}{\sqrt n}
           \left(1-\frac{\log n}{2L}\right)
 \ge c\,\frac{e^L}{L}.
 \tag{ON}
\]
In fact there are fixed $c_1,C_1>0$ and $L_1$ such that
\[
 c_1e^L\le\|\mathfrak D_L\|\le C_1e^L
 \qquad(L\ge L_1).                                      \tag{SN}
\]
Nevertheless the exact full-form decomposition is
\[
 A_{e^L}=A_{\rm arch,L}+R_L-\mathfrak D_L,
 \qquad R_L(\ell,\ell')=e^{-|\ell-\ell'|/2},
 \qquad 0\le R_L,\quad \|R_L\|\le4 ,
 \tag{FD}
\]
in the sense of closed forms on the stated log-Fourier domain.
Thus replacing the centered prime operator by a subpolynomial
operator-norm remainder is impossible. This does not give a
negative eigenvalue of the full form: its archimedean multiplier
is retained in (FD).
\end{theorem}
\begin{proof}
For fixed $L$ and real $T$ use the unit vector
$f_T(\ell)=(2L)^{-1/2}e^{iT\ell}\mathbf1_{[-L,L]}(\ell)$.
Its Fourier transform is bounded near $T$ and
$O_{L,T}((1+|t|)^{-1})$ at infinity, so $f_T$ belongs
to the form domain. Direct overlap gives, for $0\le r\le2L$,
\[
 \langle f_T,(\tau_r+\tau_{-r})f_T\rangle
       =2(1-r/(2L))\cos(Tr).
\]
For the continuum term substitute $r=\log x$.
Its resulting cosine integral has the integrable weight
$2e^{r/2}(1-r/(2L))\mathbf1_{[0,2L]}(r)$ and tends
to zero as $T\to\infty$.

There are only finitely many primes at most $X$. Simultaneous
recurrence supplies $T_k\to\infty$ with
$e^{iT_k\log p}\to1$ for all those primes.
Here is an elementary construction: partition the torus for
the vector $(\log p/(2\pi))_{p\le X}$ into $Q^d$ cubes,
where $d$ is the number of primes. Two of its first $Q^d+1$
integer multiples occupy one cube, giving a positive integer
$q_Q\le Q^d$ with coordinate errors at most $1/Q$.
If these integers have an unbounded subsequence, use it.
Otherwise some fixed $q$ has zero coordinate errors; the
multiples $kq$ give exact recurrence.
Every retained prime power has a bounded exponent for this
fixed $L$, so $e^{iT_k\log n}\to1$ for all retained $n$.
It follows that
\[
 \langle f_{T_k},\mathfrak D_Lf_{T_k}\rangle
 \longrightarrow
 2\sum_{n\le X}\frac{\Lambda(n)}{\sqrt n}
                      (1-\log n/\log X),
\]
proving the first inequality of (ON).

The required lower Chebyshev estimate is elementary.
For each prime $p$, every summand in
$v_p\binom{2m}{m}=\sum_{j\ge1}
(\lfloor2m/p^j\rfloor-2\lfloor m/p^j\rfloor)$
is zero or one. Thus
$\log\binom{2m}{m}\le\psi(2m)$.
The largest coefficient of $(1+z)^{2m}$ is the central one,
so $\binom{2m}{m}\ge4^m/(2m+1)$.
Choosing $m=\lfloor x/2\rfloor$ proves
$\psi(x)\ge c_0x$ for all sufficiently large $x$, with
some fixed $c_0>0$.
Set $w_X(x)=x^{-1/2}(1-\log x/\log X)$.
It is nonnegative on $[1,X]$, vanishes at $X$, and
$-w_X'(x)\ge x^{-3/2}/\log X$ there.
Stieltjes integration by parts gives
\[
 \sum_{n\le X}\Lambda(n)w_X(n)
 =-\int_1^X\psi(x)w_X'(x)\dd x
 \ge\frac{c_0}{\log X}\int_{X/4}^{X/2}x^{-1/2}\dd x
 \gg\frac{\sqrt X}{\log X}.
\]
This proves the second inequality, including every prime power.

To obtain the sharper lower bound, keep the constants
$c_0Y\le\psi(Y)\le C_0Y$ valid for all sufficiently large $Y$.
Fix $\delta>0$ with $e^\delta>2C_0/c_0$ and put
\[
 g_\delta(u)=\delta^{-1/2}\mathbf1_{[0,\delta]}(u),\qquad
 G_\delta(s)=(g_\delta*g_\delta)(s).
\]
The function $G_\delta$ is triangular, with values
$s/\delta$ on $[0,\delta]$ and $(2\delta-s)/\delta$
on $[\delta,2\delta]$. In particular it is at least $1/2$
on $[\delta/2,3\delta/2]$. For $L>\delta$ the functions
$r_L(x)=g_\delta(L-x)$ and $\ell_L(x)=g_\delta(L+x)$
are disjoint unit vectors. Set $\phi_L=(r_L+\ell_L)/\sqrt2$.
Every diagonal translation overlap is nonnegative, while
the two cross terms together give
\[
 \langle\phi_L,\Pi_L\phi_L\rangle
 \ge\sum_{n\le X}\frac{\Lambda(n)}{\sqrt n}
                           G_\delta(2L-\log n).
\]
With $U=Xe^{-\delta/2}$ and $V=Xe^{-3\delta/2}$,
both large when $L$ is large, the fixed choice of $\delta$ gives
\[
 \psi(U)-\psi(V)\ge c_0U-C_0V\ge\tfrac12c_0U.
\]
Therefore the displayed Rayleigh quotient is at least
$c_0\sqrt U/4=c_0e^{-\delta/4}e^L/4$.
Modulate $\phi_L$ by $e^{iT_kx}$ along the same simultaneous
prime-phase recurrence already constructed above.
For each fixed $L$, its $\Pi_L$ quotient tends to the
unmodulated quotient. The continuous kernel
$e^{|x-x'|/2}$ makes $C_L$ Hilbert--Schmidt on the fixed
interval. The modulated vectors converge weakly to zero:
their inner products with any $L^2$ vector are Fourier
transforms of $L^1$ products. Compactness of $C_L$ hence
makes its quotient tend to zero. This proves the lower
bound in (SN). These step-function vectors also belong
to the form domain, since their Fourier tails are $O(1/|t|)$.
Finally
\[
 \|\Pi_L\|\le2\sum_{n\le X}\frac{\Lambda(n)}{\sqrt n}
       =O(e^L),\qquad
 \|C_L\|\le2\int_1^Xx^{-1/2}\dd x=4(e^L-1),
\]
which proves its upper bound. The endpoint width $\delta$
is fixed before $L$ tends to infinity; no PNT remainder
or phase-recurrence height bound is used.

Finally the kernel of $C_L$ is $e^{|\ell-\ell'|/2}$
away from the diagonal, while the pole kernel is
$e^{|\ell-\ell'|/2}+e^{-|\ell-\ell'|/2}$.
Their difference is exactly $R_L$.
The Fourier transform of $e^{-|r|/2}$ is
$(t^2+1/4)^{-1}$, proving positivity and the norm bound
$4$ by zero extension and Plancherel.
The identity $A=A_{\rm arch}+P-\Pi_L$ now proves (FD).
Although the resonant trial vectors make $\mathfrak D_L$
large, their frequencies tend to infinity, where the
archimedean multiplier grows logarithmically. No upper
bound on that cost was used, so (ON) cannot be read as a
negative-eigenvalue estimate for $A$.
\end{proof}

The next lower-bound problem is therefore a frequency-sensitive
form estimate retaining $A_{\rm arch,L}+R_L$ in (FD), rather
than an unweighted norm estimate for $\mathfrak D_L$.

\section{Frequency blocks retain the archimedean cost}
The multiplier of $\mathfrak D_L$ on zero extensions is
\[
 d_L(t)=2\sum_{n\le e^{2L}}\frac{\Lambda(n)}{\sqrt n}
                          \cos(t\log n)
             -2\int_1^{e^{2L}}x^{-1/2}\cos(t\log x)\dd x.
\]
For $h>0$ set $I_j=[jh,(j+1)h]$, $j\in\Z$, and define
\[
 \begin{split}
 g_L(t)&=d_L(t)-a(t),\qquad
 m_j=\frac1h\int_{I_j}g_L(t)\dd t,\\
 H_j(t)&=\int_{jh}^{t}(g_L(u)-m_j)\dd u,\quad t\in I_j,\\
 M_{L,h}&=\sup_j(m_j)_+,\qquad
 V_{L,h}=\sup_j\sup_{t\in I_j}|H_j(t)|.
 \end{split}
\]
The constant $h$ is the frequency-block width, not the earlier
packet function; only this section uses that scalar notation.

\begin{theorem}[Block means and centered primitives]
\label{thm:frequency-blocks}
If $M_{L,h},V_{L,h}<\infty$, then every element of the full
form domain on $[-L,L]$ satisfies
\[
 q_{A_{e^L}}(f,f)\ge
          -(M_{L,h}+2LV_{L,h})\|f\|_2^2.                 \tag{FB}
\]
For $L\ge1$, the unconditional elementary bounds give
\[
 M_{L,1/L}=O(e^L),\qquad V_{L,1/L}=O(e^L/L).
\]
Consequently the following are sufficient, uniformly for
all sufficiently large $L$, to imply RH:
\[
 \forall\epsilon>0:\quad
 M_{L,1/L}=O_\epsilon(e^{\epsilon L}),\qquad
 V_{L,1/L}=O_\epsilon(e^{\epsilon L}/L).                 \tag{BF}
\]
The two estimates in (BF) are not proved here and are not
asserted to be necessary.
\end{theorem}
\begin{proof}
Write $F=\widehat f$. Compact support gives
$F'=-i\widehat{xf}$, and Plancherel gives
\[
 \|F\|_2=\sqrt{2\pi}\|f\|_2,
 \qquad \|F'\|_2=\sqrt{2\pi}\|xf\|_2
                        \le\sqrt{2\pi}L\|f\|_2.
\]
Thus $(|F|^2)'$ is integrable and has $L^1$ norm at most
$4\pi L\|f\|_2^2$. Both endpoint values of $H_j$ are zero.
Integration by parts on each block therefore gives
\[
 \int_{I_j}(g_L-m_j)|F|^2\dd t
                  =-\int_{I_j}H_j(|F|^2)'\dd t.
\]
The infinite summation is legitimate. For fixed $L,h$,
$d_L$ is bounded and
$|m_j|\le C_{L,h}(1+\log(2+|t|))$ for $t\in I_j$;
the defining form-domain weight makes all mean terms
absolutely integrable. The centered terms are absolutely
summable by $V_{L,h}\|(|F|^2)'\|_1<\infty$.
Summing and dividing by $2\pi$ gives
\[
 \frac1{2\pi}\int g_L(t)|F(t)|^2\dd t
          \le(M_{L,h}+2LV_{L,h})\|f\|_2^2.
\]
Since (FD) subtracts this integral and adds $q_{R_L}\ge0$,
this proves (FB) on the entire domain, rather than just
on smooth test functions.

For the elementary estimates, partial summation of
$\psi(X)=O(X)$ gives
$\|d_L\|_\infty=O(e^L)$ and
$\|d_L'\|_\infty=O(Le^L)$ for $L\ge1$.
The digamma series shows $\sup_t|a'(t)|<\infty$:
with $x_n=n+1/4$ and $u=t/2$, its differentiated summands
are $x_nu/(x_n^2+u^2)^2$, bounded in magnitude by
$C/x_n^2$. This summable bound also justifies differentiation.
As $a$ is bounded below, $M_{L,1/L}=O(e^L)$.
For a Lipschitz function with constant $K$, its difference
from a mean on an interval of length $h$ is at most $Kh$;
the centered primitive is thus at most $Kh^2$.
Applying this with $K=O(Le^L)$ proves
$V_{L,1/L}=O(e^L/L)$. Under (BF), (FB) instead gives
the full subpower bound (NP), with all its parameter
quantifiers; Theorem~\ref{thm:negative-ground} then implies RH.
\end{proof}

This is a sufficient arithmetic criterion using one-sided
block means and absolute centered primitives. It avoids a
pointwise bound in its proof, but a large negative dip can
enlarge $V_{L,h}$ without harming the desired lower bound.
Thus it is not claimed to be weaker than every sufficient
one-sided pointwise criterion.

\begin{proposition}[Raw block averaging does not remove resonance]
\label{prop:raw-block-resonance}
For some fixed $c>0$ and all sufficiently large $L$,
\[
 \sup_{j\in\Z}L\int_{j/L}^{(j+1)/L}d_L(t)\dd t\ge ce^L.
\]
In particular the archimedean centering in (BF) cannot be
discarded merely by averaging at the scale $1/L$.
\end{proposition}
\begin{proof}
For a block of width $1/L$ and center $S$, its exact mean is
\[
 2\sum_{n\le e^{2L}}\frac{\Lambda(n)}{\sqrt n}
     \cos(S\log n)\operatorname{sinc}\!\left(\frac{\log n}{2L}\right)
 -2\int_1^{e^{2L}}x^{-1/2}\cos(S\log x)
     \operatorname{sinc}\!\left(\frac{\log x}{2L}\right)\dd x,
\]
where $\operatorname{sinc}(u)=\sin(u)/u$ with value one at zero.
Use the prime-phase recurrence from Theorem~\ref{thm:prime-norm}
and replace each $T_k$ by the nearest center $S_k$ of the fixed
grid. The error is at most $1/(2L)$, so every prime-power
phase changes by at most one. Its limiting cosine is at least
$\cos1>0$, while each sinc factor is at least $\sin1>0$.
The continuous term tends to zero by the Riemann--Lebesgue
lemma after $r=\log x$. Finally
$\sum_{n\le X}\Lambda(n)/\sqrt n\ge\psi(X)/\sqrt X\gg e^L$.
This proves the lower bound. The corresponding values of
$a(S_k)$ are unbounded; the argument proves no analogous
lower bound for $d_L-a$.
\end{proof}

\section{All-height zero sampling on the closed Weil form domain}
For $L\ge1$ write
\[
 \|f\|_{\log}^2=\frac1{2\pi}\int_\R
             \log(2+|t|)|\widehat f(t)|^2\dd t,
 \qquad z_\rho=-i(\rho-1/2).
\]
This norm is equivalent to the defining weighted form norm.
If $\rho=\beta+i\gamma$, then
$z_\rho=\gamma+i(1/2-\beta)$; functional-equation symmetry
gives $|\Im z_\rho|\le B$ without requiring a zero attaining $B$.

\begin{theorem}[Weighted sampling and the full zero-side identity]
\label{thm:zero-sampling}
There is a fixed constant $C$ such that, for every $L\ge1$
and every form-domain element supported on $[-L,L]$,
\[
 \sum_\rho m_\rho
  \bigl(|\widehat f(z_\rho)|^2+
              |\widehat f(\overline{z_\rho})|^2\bigr)
       \le C(1+L)^2e^{2BL}\|f\|_{\log}^2.               \tag{ZS}
\]
For all $f,g$ in that domain the exact identity is
\[
 q_{e^L}(f,g)=\sum_\rho m_\rho
    \overline{\widehat f(\overline{z_\rho})}
                            \widehat g(z_\rho),         \tag{ZF}
\]
with absolute convergence. This is the zero side of the
full form, including its poles and all prime powers.
Neither (ZS) nor (ZF) asserts its unconditional positivity.
\end{theorem}
\begin{proof}
The unconditional local counting estimate
\cite[Lemma 9.4]{Kedlaya}, together with conjugation and the
finite number of zeros on bounded sets, gives
\[
 \sum_{\gamma_\rho\in[j,j+1]}m_\rho
                       \le C_N\log(3+|j|)\quad(j\in\Z).
\]
Choose a fixed smooth cutoff equal to one on $(-\infty,0]$
and zero on $[1,\infty)$. The product of its translates
at the two endpoints gives $\chi_L=1$ on $[-L,L]$,
supported in $[-L-1,L+1]$, with uniformly bounded first
and second derivatives. For $|v|\le B\le1/2$, put
$k_{L,v}=\widehat{e^{vx}\chi_L(x)}$.
The zeroth and second derivative $L^1$ bounds, followed
by two integrations by parts, give
\[
 |k_{L,v}(u)|\le\kappa_L(u):=
       \frac{C_0(1+L)e^{B(L+1)}}{1+u^2}.                \tag{KC}
\]
The constant is uniform in $L,v,B$ in these ranges.

Set $F=\widehat f$. Since $\chi_L=1$ on its support,
Fourier multiplication gives pointwise
\[
 F(t+iv)=\frac1{2\pi}\int_\R F(s)k_{L,v}(t-s)\dd s.
\]
The integral converges by Cauchy--Schwarz; both sides are
continuous in $t$, so the usual $L^2$ convolution identity
also gives this pointwise identity. Another Cauchy--Schwarz
inequality with the common positive majorant yields
\[
 |F(t+iv)|^2\le\frac{\|\kappa_L\|_1}{4\pi^2}
               \int_\R|F(s)|^2\kappa_L(t-s)\dd s.
\]
Local counting now implies
\[
 \sum_\rho m_\rho\kappa_L(\gamma_\rho-s)
 \le C_1(1+L)e^{B(L+1)}\log(2+|s|).
\]
Indeed on each unit interval the kernel is at most a
fixed multiple of $(1+(j-s)^2)^{-1}$, and
$\log(3+|j|)\le\log(3+|s|)+\log(1+|j-s|)$.
The resulting two lattice sums are bounded uniformly in
$s$, since the kernel and its logarithmic weight are
summable. Tonelli and
$\|\kappa_L\|_1=O((1+L)e^{B(L+1)})$ prove (ZS) for
both signs of the imaginary evaluation coordinate.
The factor $e^{2B}$ is absorbed in the fixed constant.
This argument handles varying horizontal zero positions,
clusters, and multiplicities with one common kernel.

On $C_c^\infty((-L,L))$, (ZF) is the unconditional Weil
explicit formula in \cite[eqs.~(1.1), (2.2), (2.6),
(2.8)--(2.12), pp.~97, 101--103]{ConnesConsani}.
Its Fourier sign is $e^{-itx}$, and
$2\theta'(t)=a(t)$ in those formulas.
The multiplicative convolution $f^**g$, after $u=e^x$,
is exactly $C_{f,g}(\log u)$ used here. Consequently
the pole evaluations, $n^{-1/2}\Lambda(n)$ translations,
and the antilinear first variable agree with the displayed
definition of $q_\lambda$. Reflection
$\rho\mapsto1-\overline\rho$ preserves multiplicity and sends
$z_\rho$ to $\overline{z_\rho}$; reindexing by this symmetry
also verifies the displayed pairing and its Hermitian convention.
Every translated correlation
outside $[-2L,2L]$ vanishes; its endpoint values vanish
as well. No omitted endpoint or trivial-zero term is added.

For completeness the smooth interior functions are a core
for the weighted domain. If $0<b<1$, compress a vector by
$f_b(x)=b^{-1/2}f(x/b)$. Its Fourier transform is
$\sqrt bF(bt)$, and these weighted $L^2$ dilations converge
strongly as $b\uparrow1$. To see this, first use continuous
compactly supported Fourier functions, then their density;
the dilation norms are uniformly bounded for $1/2\le b\le1$
because $\log(2+|t|/b)\le C\log(2+|t|)$.
The compressed support has positive distance from both
endpoints. Convolution with a nonnegative smooth compact
approximate identity, of support smaller than that distance,
therefore gives interior smooth functions. Their Fourier
transforms converge in the weighted norm by dominated
convergence, using a multiplier bounded by one.

The arithmetic form is continuous on this weighted space
for each fixed $L$: the multiplier has logarithmic growth,
and its finite translations and pole evaluations are bounded
in $L^2([-L,L])$. The two zero-evaluation maps in (ZS) are
bounded from the weighted space to $\ell^2(m_\rho)$.
Their paired sum is therefore a continuous sesquilinear
form, absolutely convergent by Cauchy--Schwarz.
Approximation by the preceding core extends (ZF) to every
domain pair. This is a fixed-$L$ extension with the uniform
sampling estimate stated in (ZS), not an interchange of
ground-state or height limits.
\end{proof}

\begin{corollary}[A conditional sharp negative-ground growth rate]
\label{cor:off-line-count}
Suppose there is $M<\infty$ such that the off-line zeros obey
\[
 \sup_{j\in\Z}
   \sum_{\substack{\gamma_\rho\in[j,j+1]\cr \Re\rho\ne1/2}}
                      m_\rho\le M.                     \tag{OC}
\]
Then for every $L\ge1$,
\[
 \nu_{e^L}\le C M(1+L)^2e^{2BL},\qquad
 \lim_{\lambda\to\infty}
       \frac{\log(1+\nu_\lambda)}{\log\lambda}=2B.       \tag{SG}
\]
Condition (OC) holds if there are only finitely many off-line
zeros. Neither condition is asserted for zeta here, and
(SG) alone does not imply that $B=0$.
\end{corollary}
\begin{proof}
Repeat the common-kernel argument for the off-line subfamily.
Under (OC), its summed kernel is bounded by
$C M(1+L)e^{B(L+1)}$, without the logarithmic factor in $s$.
Thus each of its two sampling sums is at most
$C M(1+L)^2e^{2BL}\|f\|_2^2$.
In (ZF), every critical-line term is exactly
$m_\rho|\widehat f(\gamma_\rho)|^2\ge0$.
The absolute off-line paired sum is bounded by the same
constant using Cauchy--Schwarz. The variational principle
gives the first inequality in (SG). Its upper rate is
at most $2B$, and Corollary~\ref{cor:cofinal-ground-rate} gives
the reverse inequality. If $M=0$, the full form is
nonnegative, $\nu_\lambda=0$, and both rates are zero.
\end{proof}

The unconditional bound (ZS) is in $\|f\|_{\log}$, not
$\|f\|_2$. The ratio of those squared norms is unbounded:
modulating a fixed nonzero smooth packet by $e^{iTx}$
keeps its $L^2$ norm fixed and makes its weighted norm
grow like $\log(2+|T|)$. Thus (ZS) closes the all-height
domain extension, but does not close the global negative-ground
estimate. The extra off-line count (OC) explicitly supplies
the missing unweighted sampling in the conditional corollary.

\begin{corollary}[Compensation by quadratic horizontal mass]
\label{cor:horizontal-mass}
Suppose
\[
 M_2:=\sup_{j\in\Z}
  \sum_{\substack{\gamma_\rho\in[j,j+1)\cr \beta_\rho\ne1/2}}
          m_\rho(\beta_\rho-1/2)^2<\infty.             \tag{HM}
\]
Then, with a fixed $C$ and all $L\ge1$,
\[
 \nu_{e^L}\le C M_2(1+L)^4e^{2BL},\qquad
 \lim_{\lambda\to\infty}
       \frac{\log(1+\nu_\lambda)}{\log\lambda}=2B.     \tag{HGR}
\]
Condition (OC) implies (HM), but (HM) does not require a
uniform off-line count. No unconditional proof of (HM)
or of $B=0$ is supplied.
\end{corollary}
\begin{proof}
Write $\delta_\rho=1/2-\beta_\rho$ and $F=\widehat f$.
The difference of a reflected pair is exactly
\[
 D_\rho=F(\gamma_\rho+i\delta_\rho)
             -F(\gamma_\rho-i\delta_\rho)
     =\delta_\rho\,\widehat{b_{\delta_\rho}f}(\gamma_\rho),
 \qquad b_v(x)=\frac{2\sinh(vx)}v,
 \quad b_0(x)=2x.
\]
Use the cutoff $\chi_L$ from (KC). Uniformly for $|v|\le B$,
the zeroth and second derivative $L^1$ norms of
$b_v\chi_L$ are bounded by
$C(1+L)^2e^{B(L+1)}$. Indeed
$b_v=2x\int_0^1\cosh(tvx)\dd t$,
$b_v'=2\cosh(vx)$, and $b_v''=2v\sinh(vx)$;
the cutoff derivatives lie in boundary intervals of fixed length.
Thus all their Fourier kernels have the common majorant
\[
 \kappa_L^{(2)}(u)=
       \frac{C(1+L)^2e^{B(L+1)}}{1+u^2}.
\]
The convolution Cauchy--Schwarz argument of (ZS), now
weighted by $m_\rho\delta_\rho^2$, is simultaneous for
the different $\delta_\rho$. By (HM) and the summable lattice kernel,
\[
 \sum_{\rho\text{ off-line}}m_\rho\delta_\rho^2
           \kappa_L^{(2)}(\gamma_\rho-s)
       \le C M_2(1+L)^2e^{B(L+1)}\quad(s\in\R).
\]
Tonelli therefore gives
\[
 \sum_{\rho\text{ off-line}}m_\rho|D_\rho|^2
       \le C M_2(1+L)^4e^{2BL}\|f\|_2^2.
\]
The fixed factor $e^{2B}$ is absorbed in $C$.
For one orbit the exact pair identity in (ZF) is
\[
 \frac{m_\rho}{2}\left(
   |F(z_\rho)+F(\bar z_\rho)|^2-|D_\rho|^2\right).
\]
Thus the full form is a positive sum minus one half of
the difference-squared sum over one representative per
off-line orbit. Critical-line terms are nonnegative.
The displayed estimate bounds this
negative part and gives (HGR) by the variational principle
and the lower rate in (CL).
If $M_2=0$, every off-line term is absent and both rates are zero.
Finally (OC) implies $M_2\le B^2M$, since
$|\delta_\rho|\le B$ and a half-open unit interval is
contained in the corresponding closed interval.
\end{proof}

The unconditional local count only bounds each sum in
(HM) by $O(\log(3+|j|))$. Thus the horizontal factor
supplies a different sufficient sampling hypothesis,
not a proof that the hypothesis holds for zeta.

\section{Individual zero detectors and the full negative index}
Let $\mathcal Z$ be the set of distinct zeros of $\Xi$, with
their multiplicities $m_z$. Define
\[
 \mathcal H_{\mathcal Z}=\ell^2(\mathcal Z,m_z),\qquad
 \mathcal T f=(\widehat f(z))_{z\in\mathcal Z},\qquad
 (Jc)_z=c_{\bar z}.
\]
Conjugation preserves $\mathcal Z$ and its multiplicities, so
$J$ is a self-adjoint unitary on this Hilbert space.
The closed-domain identity (ZF) reads
$q_{e^L}(f,g)=\langle J\mathcal T f,\mathcal T g\rangle$.
The norm of this evaluation space is positive independently of
RH; the operator $J$ can have negative directions.

\begin{lemma}[All-height approximation of prescribed zero values]
\label{lem:zero-interpolation}
For every finite subset $S\subset\mathcal Z$ and every prescribed
vector $(\alpha_z)_{z\in S}$, there is a smooth function $f$ such that
\[
 \widehat f(z)=\begin{cases}\alpha_z,&z\in S,\\0,&z\in\mathcal Z\setminus S,
                         \end{cases}
 \qquad |f^{(r)}(x)|\le C_r e^{-c e^{2|x|}}\quad(r\ge0). \tag{ID}
\]
Here $c>0$ and the finite $C_r$ may depend on the selected
zeros, their exact orders, and the prescribed values.
For smooth cutoffs $\chi_R=1$ on $[-R,R]$ and supported on
$[-R-1,R+1]$, with uniform derivative bounds,
\[
 \mathcal T(\chi_R f)\longrightarrow\alpha
                       \quad\hbox{in }\mathcal H_{\mathcal Z}, \tag{IT}
\]
where $\alpha$ is extended by zero outside $S$.
In particular $\mathcal T(C_c^\infty(\R))$ is dense in
$\mathcal H_{\mathcal Z}$.
\end{lemma}
\begin{proof}
The Xi kernel already reconstructed in the partition section
satisfies, for every fixed derivative order,
\[
 |\Phi^{(r)}(x)|\le C_r e^{-c_0 e^{2|x|}}.
\]
To verify the derivative assertion, on $x\ge0$ put
$T_n=\pi n^2e^{2x}$. Each derivative of a summand is
$e^{x/2}$ times a fixed polynomial in $T_n$, times $e^{-T_n}$.
Absorb that polynomial in $e^{-T_n/2}$ and sum the remaining
Gaussian series. The resulting $e^{x/2}e^{-\pi e^{2x}/2}$
is bounded by $C_r e^{-\pi e^{2x}/4}$.
Smooth evenness gives the negative half-line bound.
The Fourier convention is $e^{-izx}$; evenness of $\Phi$
therefore makes $\widehat\Phi(z)=\Xi(z)$.

We first justify division by a zero as a physical operation.
Suppose $h$ has the preceding decay for every derivative and
$\widehat h(a)=0$. Set
\[
 u(x)=i e^{iax}\int_{-\infty}^x e^{-iat}h(t)\dd t
     =-i e^{iax}\int_x^\infty e^{-iat}h(t)\dd t.          \tag{OD}
\]
The second equality uses exactly the vanishing Fourier value
at $a$, and all integrals converge absolutely.
Use the right tail when $x>0$ and the left tail when $x<0$.
For example, writing $t=x+s$ in the right tail and using
$e^{2(x+s)}\ge e^{2x}(1+2s)$, its magnitude for large $x$
is at most
\[
 C e^{-c_0e^{2x}}\int_0^\infty
       e^{-(2c_0e^{2x}-|\Im a|)s}\dd s
       \le C' e^{-c_0e^{2x}}.
\]
The left tail has the same estimate. The bounded middle
interval is absorbed into its constant.
The equation $u'=ia u+ih$ gives the bounds for every fixed
derivative of $u$. Integration by parts, with vanishing
boundary terms on every fixed complex Fourier line, now gives
\[
 (-i\partial_x-a)u=h,\qquad
              (z-a)\widehat u(z)=\widehat h(z).
\]
Consequently $\widehat u=\widehat h/(z-a)$ as entire functions.
This proof also covers real $a$.

Put
\[
 D_S(z)=\prod_{a\in S}(z-a)^{m_a},\qquad
                   G_S(z)=\Xi(z)/D_S(z).
\]
All factors use the exact zero orders. Thus $G_S$ is entire,
is nonzero at every $a\in S$, and vanishes at every other
point of $\mathcal Z$. Repeated use of (OD) constructs a
smooth inverse Fourier function $g_S$ of $G_S$ with the
same type of double-exponential decay. At each division
the current Fourier function still vanishes at the factor
being removed: at a multiple zero this is exactly the
remaining-order condition. No simplicity premise is used.

The distinct interpolation nodes admit the polynomial
\[
 p(z)=\sum_{a\in S}\frac{\alpha_a}{G_S(a)}
                   \prod_{b\in S\setminus\{a\}}\frac{z-b}{a-b}.
\]
Then $F=G_Sp$ has precisely the values in (ID).
Multiplication by $p(z)$ corresponds to $p(-i\partial_x)g_S$;
this is the required $f$, with all the stated derivative bounds.
The factors $G_S(a)$ are retained separately: conjugate
values of $G_S$ need not be equal.

Here is an all-height error bound for the compact truncation.
Write $h_R=(1-\chi_R)f$ and define the finite quantity
\[
 E_R=2\max\left\{
  \sup_{|v|\le1/2}\|e^{vx}h_R\|_1,
  \sup_{|v|\le1/2}\|\partial_x^2(e^{vx}h_R)\|_1\right\}.
\]
The derivative bounds of the cutoff and of $f$ give
$E_R\le C e^{-c' e^{2R}}$ for $R\ge1$, after decreasing
$c'>0$ if necessary. The zeroth-order estimate and two
integrations by parts together yield
\[
 |\widehat h_R(\gamma+iv)|\le\frac{E_R}{1+\gamma^2}
                              \quad(|v|\le1/2).
\]
Local zero counting proves
\[
 S_4:=\sum_{z=\gamma+iv\in\mathcal Z}
                 \frac{m_z}{(1+\gamma^2)^2}<\infty,
 \qquad \|\mathcal T h_R\|^2\le S_4 E_R^2.             \tag{TE}
\]
Thus (IT) holds with all heights and multiplicities, without
interchanging a growing-$L$ sampling constant with a limit.
Every finitely supported vector has been approximated;
their density in $\ell^2(\mathcal Z,m_z)$ proves the last assertion.
\end{proof}

For the full operator let $n_-(A_{e^L})$ count its strictly
negative eigenvalues with spectral multiplicity. Its compact
resolvent and lower bound make this number finite at each fixed $L$.
Equivalently it is the maximal dimension of a subspace on which
the quadratic form is negative definite.

\begin{theorem}[The exact exhausted negative index]
\label{thm:negative-index}
Let $K$ be the number of distinct off-line reflection orbits
$\{z,\bar z\}$ of $\mathcal Z$, taking $K=\infty$ if necessary.
Negative and positive heights are counted as separate orbits.
Then
\[
 \sup_{L\ge1}n_-(A_{e^L})=K.                           \tag{NI}
\]
The index is nondecreasing in $L$. If $K<\infty$, it equals
$K$ for every sufficiently large $L$.
A zero order $m_z$ scales an orbit's contribution and does
not increase the number of its evaluation coordinates.
\end{theorem}
\begin{proof}
An off-line orbit contributes, for $a=\widehat f(z)$ and
$b=\widehat f(\bar z)$,
\[
 m_z(\bar b a+\bar a b)
              =\frac{m_z}{2}(|a+b|^2-|a-b|^2).
\]
Every real zero instead contributes $m_z|\widehat f(z)|^2$.
If $K<\infty$, the negative part is the sum of $K$ squared
linear functionals. A subspace of dimension $K+1$ contains
a nonzero vector in their common kernel, where the form is
nonnegative. Hence $n_-(A_{e^L})\le K$.

For the reverse inequality choose any $k$ off-line orbits.
Let $\alpha^{(j)}$ equal $1$ at $z_j$, $-1$ at $\bar z_j$,
and zero at all other zeros. These vectors satisfy
\[
 \langle J\alpha^{(i)},\alpha^{(j)}\rangle
                          =-2m_{z_j}\delta_{ij}.
\]
Lemma~\ref{lem:zero-interpolation} supplies compact cutoffs
$f_{j,R}$ with evaluation errors
$\eta_{j,R}=\|\mathcal T f_{j,R}-\alpha^{(j)}\|\to0$.
By (ZF), for $L=R+2$ their form matrix has entry error at most
\[
 \eta_{i,R}\sqrt{2m_{z_j}}+
 \eta_{j,R}\sqrt{2m_{z_i}}+\eta_{i,R}\eta_{j,R}.          \tag{ME}
\]
This follows from Cauchy--Schwarz and $\|J\|=1$.
If $k$ times the largest error is less than
$2\min_j m_{z_j}$, the matrix is strictly negative definite.
The bounds (TE) show that this occurs for sufficiently large
$R$. The cutoff functions then span a $k$-dimensional
negative subspace of the actual full form. Since every
finite set of orbits is allowed, this proves (NI), including
the infinite case.

Zero extension preserves the form from a smaller interval
inside a larger one, as in Theorem~\ref{thm:negative-ground}.
It therefore preserves any negative subspace and proves
index monotonicity. When $K$ is finite, the preceding construction
with $k=K$ and the upper bound give eventual equality.
\end{proof}

\begin{corollary}[One negative direction in each parity sector]
\label{cor:parity-index}
Let $K_+$ count distinct off-line reflection orbits at positive
height. Write $A_{e^L}^+$ and $A_{e^L}^-$ for the even and odd
restrictions of the full operator. Then
\[
 \sup_{L\ge1}n_-(A_{e^L}^+)
 =\sup_{L\ge1}n_-(A_{e^L}^-)=K_+,\qquad
 \sup_{L\ge1}n_-(A_{e^L})=2K_+.                       \tag{PI}
\]
If $K_+<\infty$, both sector indices equal $K_+$ for every
sufficiently large $L$. In particular, an unbounded sequence
of parameters with full negative index at most one would imply RH.
No such index bound is proved here.
\end{corollary}
\begin{proof}
There is no real zeta zero in $0<s<1$. Indeed
\[
 \eta(s)=\sum_{k\ge1}\bigl((2k-1)^{-s}-(2k)^{-s}\bigr)>0
                          \quad(0<s<1),
\]
whereas $\eta(s)=(1-2^{1-s})\zeta(s)$ has a negative
prefactor there. The identity first holds in $\Re s>1$
and continues to $\Re s>0$: the paired series converges
locally uniformly there, since each difference is
$O_K(k^{-1-\sigma_K})$ on a compact set with
$\sigma_K=\min_K\Re s>0$.
Thus every off-line zero belongs to four distinct points
$\{z,\bar z,-z,-\bar z\}$ with $\Re z>0$.

The full form is invariant under $Rf(x)=f(-x)$; hence its
even and odd sectors are orthogonal also for the form.
If $Rf=\varepsilon f$, $\varepsilon\in\{1,-1\}$, then
$F(-z)=\varepsilon F(z)$. The four-point contribution is
\[
 2m_z(\bar b a+\bar a b)
                  =m_z(|a+b|^2-|a-b|^2),
 \qquad a=F(z),\quad b=F(\bar z).
\]
It therefore has at most one negative direction in each
sector, independently of the common zero order. The
common-kernel argument of Theorem~\ref{thm:negative-index}
gives the upper bound $K_+$ when that number is finite.

Choose any $k$ positive-height orbits and the individual
tests used in that theorem, with values $1,-1$ on each
selected positive-height pair and zero at every other zero.
Their even and odd projections
\[
 f_{j,\varepsilon}=(f_j+\varepsilon Rf_j)/\sqrt2
\]
have zero-side form matrix $-2m_{z_j}\delta_{ij}$ in
each sector: at $z_j,\bar z_j$ the values are
$1/\sqrt2,-1/\sqrt2$, and at $-z_j,-\bar z_j$ they
are multiplied by $\varepsilon$.
Use even cutoffs in (IT). They preserve parity and the
same all-height convergence, so both sector matrices are
strictly negative for sufficiently large $L$.
This proves both lower bounds, including $K_+=\infty$.
Nested intervals give monotonicity and finite stabilization.
The full index is the sum of the two sector indices;
take a common sufficiently large parameter for each
finite selected set to obtain the last equality in (PI).
If any off-line orbit existed, both indices would be
at least one for all sufficiently large $L$. This
contradicts the stated cofinal index-at-most-one condition.
\end{proof}

The finite-exception negative-index phenomenon is already described
in the primary abstract of Bombieri~\cite{BombieriWeil}. That source's
full proof is not imported here. The explicit Fourier division and
all-height cutoff argument above supply the needed reconstruction
for the present full interval forms. No originality claim is made.
The detectors depend on selected zero positions and their orders;
they certify the implication from a hypothetical off-line zero to
a negative direction, not the existence of such a zero.
Neither (NI) nor the density statement bounds those negative
directions from arithmetic information, and RH remains open.

\begin{lemma}[Independent pole values and exact pole annihilation]
\label{lem:pole-free-interpolation}
For every finite $S\subset\mathcal Z$, prescribed values
$(\alpha_z)_{z\in S}$, and $\beta_+,\beta_-\in\C$, there is a
smooth, doubly exponentially decreasing function $f$, with the same
decay type for every fixed derivative, such that
\[
 \widehat f(z)=\alpha_z\ (z\in S),\qquad
 \widehat f(z)=0\ (z\in\mathcal Z\setminus S),\qquad
 \widehat f(\pm i/2)=\beta_\pm.
\]
More precisely, for every fixed derivative order $r$ there are
finite constants $C_r,M_r$ such that
\[
 |f^{(r)}(x)|\le C_r\Phi(x)e^{M_r|x|}.                 \tag{PNR}
\]
In particular, when $\beta_+=\beta_-=0$, there are
$w_R\in C_c^\infty(\R)$, supported on $[-R-1,R+1]$,
with both pole values exactly zero and
\[
 \|\mathcal T w_R-\alpha\|_{\mathcal H_{\mathcal Z}}
          \le C e^{-c e^{2R}},\qquad
 w_R\longrightarrow f\quad\hbox{in every fixed }H^j(\R).
                                                               \tag{PN}
\]
The constants depend on the selected nodes, their exact orders,
and their values. If the zero-value vector, extended by zero,
satisfies $\alpha_{-z}=\varepsilon\alpha_z$ and
$\beta_-=\varepsilon\beta_+$ for $\varepsilon\in\{1,-1\}$,
the construction preserves that parity.
The pole-free compact tests therefore have dense evaluation range,
also in each parity subspace of $\mathcal H_{\mathcal Z}$.
\end{lemma}
\begin{proof}
Use $D_S$ and $G_S=\Xi/D_S$ from
Lemma~\ref{lem:zero-interpolation}. Besides being nonzero at
every node of $S$, $G_S$ is nonzero at $\pm i/2$, since
$\Xi(i/2)=\xi(0)=1/2$ and $\Xi(-i/2)=\xi(1)=1/2$.
Interpolate a polynomial $p$ on $S\cup\{i/2,-i/2\}$ with
values $\alpha_z/G_S(z)$ and $\beta_\pm/G_S(\pm i/2)$.
The inverse Fourier function of $G_Sp$ is obtained by the
already proved Fourier division and constant-coefficient
differentiation. It has all the asserted values and derivative
decay, with no simplicity hypothesis.

We verify the stronger relative estimate (PNR), needed for the
weighted domains at zero heat. A positive theta summand on $x\ge0$
is $e^{x/2}(4T^2-6T)e^{-T}$, where $T=\pi n^2e^{2x}$.
Its derivative polynomial is $-8T^3+30T^2-15T$.
Since
\[
 (-8T^3+30T^2-15T)+T(4T^2-6T)
       =-T(4T^2-24T+15)<0\quad(T\ge6),
\]
the normally convergent positive series satisfies
$\Phi'(x)/\Phi(x)\le-\pi e^{2x}$ once $\pi e^{2x}\ge6$.
Its fixed derivatives satisfy
$|\Phi^{(r)}|\le C_r\Phi e^{M_r|x|}$: factor out the
$n=1$ exponential and sum the remaining polynomially weighted
Gaussian tails, while $4T^2-6T\ge2T^2$ on this half-line.
Evenness and a positive compact minimum cover the other points.
In (OD), if $|h(x)|\le C\Phi(x)e^{M|x|}$, the right-tail
bound now integrates
$\Phi(x+s)\le\Phi(x)e^{-\pi e^{2x}s}$ and gives
$|u(x)|\le C'\Phi(x)e^{(M-2)x}$ for sufficiently large $x$.
The left-tail version gives the corresponding bound at negative
$x$. The equation $u'=iau+ih$ gives all fixed derivative
estimates. Repeating these finitely many divisions and applying
the fixed interpolation polynomial proves (PNR).

For exact pole-free cutoffs, choose a fixed real, nonzero,
nonnegative even $\eta\in C_c^\infty((-1,1))$ and put
\[
 c_\eta=\int\eta(x)\cosh(x/2)\dd x>0,\qquad
 s_\eta=\int x\eta(x)\sinh(x/2)\dd x>0,\qquad
 \psi_0=\eta/c_\eta,\quad \psi_1=x\eta/s_\eta.
\]
Their pole-value pairs are respectively $(1,1)$ and $(1,-1)$.
With $p_{R,\pm}=\widehat{\chi_R f}(\pm i/2)$, define
\[
 w_R=\chi_R f-\frac{p_{R,+}+p_{R,-}}2\psi_0
                   -\frac{p_{R,+}-p_{R,-}}2\psi_1.
                                                               \tag{PC}
\]
Both pole values vanish exactly. Since the untruncated $f$
has zero pole values,
$|p_{R,\pm}|\le\|(1-\chi_R)f\,e^{|x|/2}\|_1
\le C e^{-c e^{2R}}$.
Two integrations by parts give
$|\widehat\psi_j(\gamma+iv)|\le C_j/(1+\gamma^2)$
uniformly for $|v|\le1/2$.
The finite sum $S_4$ in (TE) thus bounds
$\|\mathcal T\psi_j\|$ with every multiplicity retained.
Together with (TE), this proves the evaluation error in (PN).
The cutoff tails and the two small correction coefficients
also converge in every fixed $H^j$.
The polynomial is applied before truncation, so no growing
zero-height factor is introduced into the cutoff error.
Even cutoffs and the even/odd correction basis preserve parity;
equivalently one may project the prescribed interpolant onto
that parity. Density follows by approximating every finitely
supported evaluation vector.
\end{proof}

For $L\ge1$ define the closed pole-free form subspace
\[
 \mathcal D_L^0=\{f\in\mathcal D(q_{A_{e^L}}):
                  \widehat f(i/2)=\widehat f(-i/2)=0\}.
\]
Here the index of a restricted form is the maximal dimension
of a subspace on which that form is strictly negative.
In (BI), $\mathcal H_+$ and $\mathcal H_-$ denote the even and
odd subspaces of $\mathcal H_{e^L}$.
On $\mathcal D_L^0$ the forms of $A_{e^L}$ and $B_{e^L}$ agree.

\begin{theorem}[The exact pole-free exhausted indices]
\label{thm:pole-free-index}
With $K_+$ as in Corollary~\ref{cor:parity-index},
\[
 \begin{split}
 \sup_{L\ge1}n_-(q_A|_{\mathcal D_L^0})&=2K_+,\\
 \sup_{L\ge1}n_-(q_A|_{\mathcal D_L^0\cap\mathcal H_\pm})&=K_+,\\
 \sup_{L\ge1}n_-(B_{e^L}^+)&=K_++1,\\
 \sup_{L\ge1}n_-(B_{e^L}^-)&=K_+,\qquad
 \sup_{L\ge1}n_-(B_{e^L})=2K_++1 .
 \end{split}                                                 \tag{BI}
\]
The first two identities also hold for $q_B$ on the same subspaces.
The convention is $\infty+1=\infty$.
All these indices are nondecreasing in $L$ and attain their
displayed values for every sufficiently large $L$ if $K_+<\infty$.
In particular $B_{e^L}$ has an unconditional negative even
direction for all sufficiently large $L$.
\end{theorem}
\begin{proof}
The upper bounds on the pole-free subspaces follow by restriction
from (NI) and (PI). Lemma~\ref{lem:pole-free-interpolation}
approximates every finite negative evaluation basis used in those
proofs by exactly pole-free compact vectors. Their full form
matrices converge to the same negative diagonal matrices, by (ZF).
This proves the reverse bounds, in both parities and at all heights.

For the extra even direction put $\varphi=2\Phi$.
Its zero evaluations all vanish and its two pole values are one.
Use compact cutoffs to define its limiting form against the
rapidly decreasing interpolants. The convergence (TE), and
weighted $L^1$ convergence at the poles, give
\[
 q_A(\varphi,\varphi)=0,\qquad q_B(\varphi,\varphi)=-2.
\]
Every exactly pole-free negative interpolant $f_j$ satisfies
$q_B(\varphi,f_j)=q_A(\varphi,f_j)-q_P(\varphi,f_j)=0$.
Thus adjoining $\varphi$ to a finite even negative basis yields
one further negative block. Its compact-cutoff Gram matrix is
strictly negative for sufficiently large cutoff radius.
The even lower bound is therefore $K_++1$.
The odd lower bound $K_+$ already follows from the pole-free tests.

When $K_+$ is finite, the exact rank-one identities
\[
 B_+=A_+-2|c_\lambda\rangle\langle c_\lambda|,\qquad
 B_-=A_-+2|s_\lambda\rangle\langle s_\lambda|
\]
give $n_-(B_+)\le K_++1$ and $n_-(B_-)\le K_+$.
For the first inequality, a negative subspace of dimension
$K_++2$ would contain a $(K_++1)$-dimensional subspace in
$c_\lambda^\perp$, where $q_B=q_A$, contradicting (PI);
the second follows from adding a positive form.
Every finite selection gives the lower bounds in the infinite case.
Zero extension preserves the forms and both pole conditions,
so all the indices are monotone. A finite selection of all orbits
then proves eventual equality when $K_+$ is finite.
\end{proof}

This is an exact index relation, rather than a universal index bound.
An unbounded parameter sequence with $n_-(B_{e^L})\le1$ would
exclude every off-line orbit, but that bound is not established.
The simple positive pole-free ground vector does not imply it.

\begin{theorem}[A cofinal individual detector with bounded regularity cost]
\label{thm:cofinal-detector}
Suppose $a=\gamma+i\delta$ and $b=\gamma-i\delta$ are zeros of
$\Xi$, with $0<\delta<1/2$ and common order $m\ge1$.
There are selection-dependent constants $c_a,C_a>0$ and $L_a$
such that for every real $L\ge L_a$ there is
$g_L\in C_c^\infty((-L,L))\cap\mathcal D_L^0$ satisfying
\[
 q_{A_{e^L}}(g_L)=q_{B_{e^L}}(g_L)\le-m,\qquad
 \|g_L\|_2^2\le C_a e^{-2\delta L}(1+L)^\delta,
 \qquad
 \nu_{e^L}\ge c_a e^{2\delta L}(1+L)^{-\delta}.            \tag{CD}
\]
For every fixed $j\ge0$, the normalized vectors
$u_L=g_L/\|g_L\|_2$ have uniformly bounded $H^j$ norms,
and their log-Fourier norms are uniformly bounded.
No constant is asserted to be uniform in the chosen height,
horizontal distance, or zero order.
\end{theorem}
\begin{proof}
Choose the fixed pole-free interpolant $f$ from
Lemma~\ref{lem:pole-free-interpolation} with values $1,-1$
at $a,b$ and zero at all other distinct zeros.
Let $\alpha$ be that evaluation vector; then
$\|\alpha\|=\sqrt{2m}$ and $\langle J\alpha,\alpha\rangle=-2m$.
The corrected cutoffs $w_R$ have evaluation errors
$\eta_R\le C_*e^{-c_*e^{2R}}$.
For $v>R+1$ set
\[
 g_{v,R}(x)=
 \frac{e^{i\gamma v}w_R(x-v)+e^{-i\gamma v}w_R(x+v)}
                       {2\cosh(\delta v)},\qquad
 \widehat g_{v,R}(z)=
 \frac{\cos(v(z-\gamma))}{\cosh(\delta v)}\widehat w_R(z).
                                                               \tag{SC}
\]
The Fourier factor equals one at $a,b$.
At every Xi zero in the full strip $|\Im z|\le1/2$ its magnitude
is at most $2e^{(1/2-\delta)v}$, since
$|\cos(u+it)|^2=\cos^2u+\sinh^2t\le\cosh^2t$.
Hence the all-height evaluation error of $g_{v,R}$ is at most
\[
 2C_*\exp((1/2-\delta)v-c_*e^{2R}).
\]
Choose
\[
 D_*=\max\{1,\log(8C_*/\sqrt m)\},\quad
 R(L)=\max\left\{R_0,\frac12\log
                 \frac{(1/2-\delta)L+D_*}{c_*}\right\},\quad
 v(L)=L-R(L)-2,
\]
where $R_0\ge1$ is large enough that
$\|w_R\|_2^2\le2\|f\|_2^2$ for $R\ge R_0$.
For every sufficiently large real $L$ the two supports are
disjoint and lie in $[-L+1,L-1]$, and the error is at most
$\sqrt m/4$. Since $J$ is unitary, (ZF) bounds the form error by
\[
 2\sqrt{2m}\frac{\sqrt m}{4}+\frac m{16}
             =(\sqrt2/2+1/16)m<m.
\]
Thus $q_A(g_{v,R})\le-m$. Both pole values remain exactly zero,
because translation multiplies a zero Fourier value by a finite
exponential.

The disjoint copies give the exact identities
\[
 \|g_{v,R}^{(j)}\|_2^2
             =\frac{\|w_R^{(j)}\|_2^2}{2\cosh^2(\delta v)}
                       \qquad(j\ge0).
\]
Therefore Rayleigh variation gives
\[
 \nu_{e^L}\ge
 \frac{2m\cosh^2(\delta v)}{\|w_R\|_2^2}
       \ge\frac{m}{4\|f\|_2^2}e^{2\delta v}.
\]
Since $R(L)\le\frac12\log(1+L)+C_a$, this proves both norm
and ground-value bounds in (CD). Finally $w_R\to f$ in every
fixed $H^j$, with $\|f\|_2>0$, so the derivative ratios of the
normalized copies are uniformly bounded. The log-Fourier norm
is controlled by their $H^1$ norm.
\end{proof}

\begin{corollary}[A lower growth rate on the entire parameter tail]
\label{cor:cofinal-ground-rate}
With $B=\sup_\rho(\Re\rho-1/2)$ as above,
\[
 \liminf_{\lambda\to\infty}
       \frac{\log(1+\nu_\lambda)}{\log\lambda}\ge2B.       \tag{CL}
\]
In particular either additional hypothesis (OC) or (HM) makes
the corresponding upper rate $2B$ a genuine limit.
\end{corollary}
\begin{proof}
For each individual off-line distance $\delta>0$, (CD) holds
for every sufficiently large $L$, so its polynomial loss gives
the lower limit $2\delta$. Take the supremum over those distances,
using the functional-equation symmetries, to obtain $2B$.
If $B=0$ the inequality is immediate. The stated conditional
limits then follow from the already proved upper bounds.
\end{proof}

\begin{theorem}[The arithmetic signature of one prescribed zero pair]
\label{thm:individual-prime-signature}
Under the hypothesis of Theorem~\ref{thm:cofinal-detector}, let
$f$ be its fixed pole-free interpolant and
$C_f(r)=\int\overline{f(x-r)}f(x)\dd x$.
As $v\to+\infty$,
\[
 e^{2i\gamma v}\sum_{n\ge2}\frac{\Lambda(n)}{\sqrt n}
                             C_f(\log n-2v)
             =2m\cosh(2\delta v)+O_f(e^{-5v}).           \tag{IS}
\]
The same statement holds with the sum replaced by
\[
 \int_1^\infty x^{-1/2}C_f(\log x-2v)
                              \dd(\psi(x)-x),\qquad
 \psi(x)=\sum_{n\le x}\Lambda(n).
\]
All series and integrals converge absolutely for each fixed $v$.
Every prime power and the complete gamma contribution are retained.
\end{theorem}
\begin{proof}
The derivative decay of $f$ gives, for every fixed $j$,
\[
 |C_f^{(j)}(r)|\le C_j e^{-c_j e^{|r|}}.
                                                               \tag{IC}
\]
Indeed $\max(|x|,|x-r|)\ge|r|/2$; split the two exponential
majorants in halves, extract the displayed factor, and integrate
the remaining doubly exponentially decreasing majorant.
This also justifies differentiation under the correlation integral.
Together with $\Lambda(n)\le\log n$, (IC) proves the asserted
absolute convergence.

For each fixed $v$, use
$f_v^-(x)=e^{-i\gamma v}f(x+v)$ and
$f_v^+(x)=e^{i\gamma v}f(x-v)$.
Their Fourier values are $e^{\pm i(z-\gamma)v}F(z)$.
The two prescribed zero coordinates give exactly
\[
 q_A(f_v^-,f_v^+)=-m(e^{2\delta v}+e^{-2\delta v})
                         =-2m\cosh(2\delta v).
\]
This noncompact identity follows by compact cutoffs with $v$
fixed first: (TE) controls all zero coordinates, weighted $L^1$
controls the poles, $H^1$ convergence controls the logarithmic
multiplier, and the fixed-shift correlation bound
$C_v e^{-c_v e^{|r|}}$ sums all prime powers.
Only after this limit is justified do we let $v$ grow.
Both pole values of both shifted functions are zero.

The exact arithmetic cross-term has the two distinct channels
\[
 e^{2i\gamma v}
 \{C_f(\log n-2v)+C_f(-\log n-2v)\}.
\]
The second channel, including its infinite prime-power sum, is
$O_f(e^{-c e^{2v}})$ by (IC).
The archimedean cross-term is
\[
 \mathcal A_f(v)=\frac{e^{2i\gamma v}}{2\pi}
                      \int a(t)|F(t)|^2e^{-2itv}\dd t.
\]
Put $D(s)=C_f(-s)$ and $k(r)=e^{-r/2}/(1-e^{-2r})$.
The digamma jump identity gives exactly
\[
 e^{-2i\gamma v}\mathcal A_f(v)
 =a(0)D(2v)+\int_0^\infty k(r)
       \{2D(2v)-D(2v-r)-D(2v+r)\}\dd r.
\]
Fubini is justified near zero by the $O(r^2)$ second difference
and $k(r)=O(1/r)$, and on the tail by integrability of $k$.
The pole-free moment identities are
\[
 \int C_f(r)e^{\pm r/2}\dd r
           =\overline{F(\mp i/2)}F(\pm i/2)=0.
\]
On $0<r<1$, the second-difference bound and (IC) make the
archimedean contribution doubly exponentially small.
On $r\ge1$ write $k(r)=e^{-r/2}+k_1(r)$, with
$0\le k_1(r)\le C e^{-5r/2}$.
The only potentially larger convolution is
\[
 \int_1^\infty e^{-r/2}D(2v-r)\dd r
       =e^{-v}\int_{-\infty}^{2v-1}e^{s/2}D(s)\dd s.
\]
Its full moment is zero, so the complementary tail is doubly
exponentially small. The remaining convolution is at most
\[
 C e^{-5v}\int_\R e^{5s/2}|D(s)|\dd s=O_f(e^{-5v});
\]
the other channels have still faster decay.
Thus $\mathcal A_f(v)=O_f(e^{-5v})$, with no omitted gamma term.
The full form identity now proves (IS), including its sign and phase.

Finally
\[
 \int_1^\infty x^{-1/2}C_f(\log x-2v)\dd x
 =e^v\int_{-2v}^\infty e^{s/2}C_f(s)\dd s
 =-e^v\int_{-\infty}^{-2v}e^{s/2}C_f(s)\dd s
 =O_f(e^{-c e^{2v}}).
\]
Subtracting this exact continuous term proves the centered version.
\end{proof}

Theorems~\ref{thm:cofinal-detector} and
\ref{thm:individual-prime-signature} are conditional on the selected
off-line pair. They identify its individual negative direction and
arithmetic growth; they neither exhibit such a zero nor exclude it.
Their bounded normalized derivative costs do not establish a
uniform arithmetic bound for the moving pair of copies.

\begin{proposition}[The exact two-point detector cost]
\label{prop:two-point-cost}
Fix $L>0$, $\gamma\in\R$, and real $\delta\ne0$.
Among $L^2$ functions supported on $[-L,L]$ subject to
\[
 \widehat f(\gamma+i\delta)=1,\qquad
                       \widehat f(\gamma-i\delta)=-1,
\]
the minimum squared norm is
\[
 \min\|f\|_2^2=
 \left(\frac{\sinh(2\delta L)}{2\delta}-L\right)^{-1}.
                                                               \tag{DC}
\]
The unique minimizer is
$f(x)=e^{i\gamma x}\sinh(\delta x)\mathbf1_{[-L,L]}(x)/I$,
where $I=\sinh(2\delta L)/(2\delta)-L>0$.
It belongs to the logarithmic form domain. As $|\delta|L\to0$,
the minimum is
\[
 \frac{3}{2\delta^2L^3}\bigl(1+O((\delta L)^2)\bigr).
\]
\end{proposition}
\begin{proof}
Put $g(x)=e^{-i\gamma x}f(x)$.
The two constraints are exactly
\[
 \int_{-L}^L g(x)\cosh(\delta x)\dd x=0,
 \qquad \int_{-L}^L g(x)\sinh(\delta x)\dd x=1.
\]
The squared norm of the second, real evaluation vector is
$I=\int_{-L}^L\sinh^2(\delta x)\dd x$.
Cauchy--Schwarz gives $\|g\|_2^2\ge1/I$.
Equality forces $g=\sinh(\delta x)/I$, and that vector
also satisfies the first constraint by oddness.
Direct integration gives the displayed value of $I$.
The minimizer's Fourier transform is bounded near finite
frequencies and is $O((1+|t|)^{-1})$ on the tail, by one
integration by parts on its support with finite boundary terms.
Thus its logarithmic Fourier integral converges.
Expanding $\sinh(2\delta L)$ proves the asymptotic formula.
\end{proof}

The last proposition only enforces two evaluations; requiring
all other zero samples to vanish can increase the cost.
Consequently the qualitative detection and index theorems
do not supply a uniform norm bound as an off-line pair
approaches the critical line. This limitation is retained
when considering arithmetic estimates for all parameters.

\begin{proposition}[Fixed-support modulation and a sharp lattice control]
\label{prop:shifted-lattice}
For fixed $L$ and $f\in C_c^\infty((-L,L))$, set
$f_T(x)=e^{iTx}f(x)$. Then, as $T\to+\infty$,
\[
 \begin{split}
 q_{e^L}(f_T,f_T)
  ={}&\log(T/(2\pi))\|f\|_2^2\\
    &-2\sum_{n\le e^{2L}}\frac{\Lambda(n)}{\sqrt n}
          \Re\bigl(e^{iT\log n}C_{f,f}(\log n)\bigr)+o(1).
 \end{split}                                                \tag{FM}
\]
The remainder is for this fixed $L,f$, not for a moving test family.

For the separate model with $a>0$, $0<\delta<1/2$ and
$z_k^\pm=(k+1/2)a\pm i\delta$, define
\[
 \mathcal W_{a,\delta}(f)=\sum_{k\in\Z}
  \left(\overline{F(z_k^-)}F(z_k^+)
              +\overline{F(z_k^+)}F(z_k^-)\right),
 \quad F=\widehat f.
\]
For every smooth test supported on $[-L,L]$,
\[
 \mathcal W_{a,\delta}(f)=\frac{4\pi}{a}\|f\|_2^2
                                    \quad(L\le\pi/a). \tag{LP}
\]
For every $L>\pi/a$ a smooth test in that interval has
strictly negative $\mathcal W_{a,\delta}$. Thus $\pi/a$
is the exact support threshold for model positivity.
These model nodes are the zeros of the real even entire function
$\cos(\pi(z-i\delta)/a)\cos(\pi(z+i\delta)/a)$;
they are not asserted to be zeta zeros.
\end{proposition}
\begin{proof}
The modulated Fourier transform is $F(t-T)$.
The multiplier asymptotic is
$a(t)=\log(|t|/(2\pi))+o(1)$.
After subtracting $\log(T/(2\pi))$, its absolute value
at $t=T+u$ is bounded by
$C(1+\log(2+|u|))$ uniformly for $T\ge2$.
This follows by splitting at $|u|=T/2$ and using the
global logarithmic growth and lower bound of $a$.
Pointwise convergence and dominated convergence against
the Schwartz function $|F(u)|^2$ give the first term in (FM).
The two pole evaluations are Fourier transforms of
$e^{\pm x/2}f(x)$ at frequency $-T$ and tend to zero.
Finally $C_{f_T,f_T}(r)=e^{iTr}C_{f,f}(r)$ gives
the exact, finite prime-power sum in (FM).

For the model put $P=2\pi/a$, $g=e^{-\delta x}f$ and
$h=e^{\delta x}f$. The Schwartz function
$H(t)=\overline{\widehat g(t)}\widehat h(t)$ has
\[
 \widehat H(\xi)=2\pi\int\overline{g(x+\xi)}h(x)\dd x.
\]
Poisson summation on the half-shifted grid gives
\[
 \mathcal W_{a,\delta}(f)=\frac{2\pi}{a}
 \sum_{n\in\Z}(-1)^n
 \left(\int\overline{g(x+nP)}h(x)\dd x
            +\int\overline{h(x+nP)}g(x)\dd x\right).   \tag{LA}
\]
All evaluation sums are absolutely convergent.
Only finitely many integrals can be nonzero, by support.
When $2L\le P$, every $n\ne0$ term vanishes; at equality
the overlap has measure zero. The $n=0$ term proves (LP).

If $L>P/2$, choose a real norm-one smooth bump of support
$[-\epsilon,\epsilon]$, with
$0<\epsilon<\min\{L-P/2,P/4\}$, and translate it to
the two centers $\pm P/2$. Let $f$ be their sum divided
by $\sqrt2$. Then $\|f\|_2=1$,
$C_{f,f}(P)=C_{f,f}(-P)=1/2$, and correlations at
$nP$, $|n|\ge2$, vanish. The two weighted correlations
at each of $n=\pm1$ in (LA) sum to $\cosh(\delta P)$.
Consequently
\[
 \mathcal W_{a,\delta}(f)
                 =\frac{4\pi}{a}(1-\cosh(\delta P))<0.
\]
This proves the sharp threshold without any approximation.
\end{proof}

For each preassigned fixed $L$, the model becomes positive on
all those tests when $a$ is small enough, although every model
zero remains off-line and its unit-interval quadratic horizontal
mass grows like $2\delta^2/a$. The model changes with $a$;
it does not give one off-line zero set positive for all $L$.
Together with (FM), it prevents inferring (HM) from a
fixed-support high-frequency positivity argument alone.
The actual arithmetic identification and growing-support
quantifiers must be retained.

\begin{proposition}[Quantitative modulation with growing support]
\label{prop:joint-modulation}
There is a fixed $C$ with the following property. For every $L>0$,
$T\ge4$, integer $N\ge1$, and $f\in C_c^\infty((-L,L))$, set
\[
 A_N(f)=\sum_{j=0}^N\binom Nj2^{j-N}\|f^{(j)}\|_2,
 \qquad f_T(x)=e^{iTx}f(x).
\]
Then the full form satisfies the exact identity
\[
 \begin{split}
 q_{e^L}(f_T,f_T)
  ={}&\log(T/(2\pi))\|f\|_2^2\\
    &-2\sum_{n\le e^{2L}}\frac{\Lambda(n)}{\sqrt n}
          \Re\bigl(e^{iT\log n}C_{f,f}(\log n)\bigr)
       +\mathcal R_{L,T}(f),
 \end{split}
\]
with the uniform bound
\[
 \begin{split}
 |\mathcal R_{L,T}(f)|\le{}&
 \frac C T\bigl(\|f\|_2^2+\|f\|_2\|f'\|_2\bigr)
 +\frac{C\log(2+T)}{T^2}\|f'\|_2^2\\
 &+4Le^LT^{-2N} A_N(f)^2 .
 \end{split}                                               \tag{JM}
\]
Consequently, fix $c>0$ and $N$ with $2N>c$. For any family
$L=L_T\le c\log T$ whose functions $f=f^{[T]}$ satisfy
$\sup_T\| (f^{[T]})^{(j)}\|_2<\infty$ for $0\le j\le N$,
the remainder tends to zero uniformly as $T\to\infty$.
The prime-power sum is retained exactly; it is not part of the error.
\end{proposition}
\begin{proof}
Put $F=\widehat f$. The sectorial digamma expansion
\cite[eq.~(5.11.2), \S5.11(ii)]{DLMFGamma}, on
$z=1/4+it/2$, gives
\[
 a(t)=\log(|t|/(2\pi))+O(1/|t|)\quad(|t|\ge1),
 \qquad |a(t)|\le C+\log(2+|t|)\quad(t\in\R).
\]
The sector stays a fixed distance from the negative real axis;
its remainder constant is therefore independent of $t$.
The compact part of the second estimate uses continuity, since
the digamma argument has real part $1/4>0$.
For $|u|\le T/2$ these facts and
$|\log(1+u/T)|\le2|u|/T$ yield
\[
 |a(T+u)-\log(T/(2\pi))|\le C(1+|u|)/T.
\]
Plancherel and Cauchy--Schwarz therefore bound the integral
of this difference against $|F(u)|^2/(2\pi)$ by
$C(\|f\|_2^2+\|f\|_2\|f'\|_2)/T$ on that central interval.
Here $\int u^2|F(u)|^2\dd u/(2\pi)=\|f'\|_2^2$.

For $|u|>T/2$, the global logarithmic estimate and
$T<2|u|$ give
\[
 |a(T+u)-\log(T/(2\pi))|
     \le C(1+\log(2+|u|))
     \le\frac{C\log(2+T)}{T^2}u^2.
\]
For the last inequality, $(1+\log(2+r))/r^2$ decreases for
$r\ge2$; evaluate it at $T/2$. This proves the remaining
archimedean error in (JM), also near $u=-T$ where using
the large-argument expansion directly would be invalid.

The pole values are
\[
 I_\pm=\widehat f_T(\pm i/2)
             =\int e^{iTx}e^{\pm x/2}f(x)\dd x,
 \qquad q_P(f_T,f_T)=2\Re(\overline{I_+}I_-).
\]
$N$ integrations by parts have no boundary terms and give
\[
 |I_\pm|\le T^{-N}e^{L/2}
                    \|(D\pm1/2)^Nf\|_1
        \le T^{-N}e^{L/2}\sqrt{2L}\,A_N(f).
\]
Thus $|q_P(f_T,f_T)|\le4Le^LT^{-2N}A_N(f)^2$.
The prime correlations are exactly
$C_{f_T,f_T}(r)=e^{iTr}C_{f,f}(r)$, with every prime power
up to $e^{2L}$ included. Combining these computations gives (JM).
Finally the assumed derivative bounds control $A_N$ uniformly,
and $Le^LT^{-2N}\le c\log T\,T^{c-2N}\to0$.
The other two errors also tend to zero. All support, derivative
order, and height quantifiers in the claimed regime are explicit.
\end{proof}

For a fixed nonzero bump $b\in C_c^\infty((-d,d))$ and its
translations $b_y(x)=b(x-y)$ with $|y|+d<L$, even that support
restriction on $L$ is unnecessary. The Fourier density on the
real axis, $|\widehat b_y|^2=|\widehat b|^2$, is unchanged; the
two pole factors acquire $e^{y/2}$ and $e^{-y/2}$, which cancel
in their product; and each prime correlation is unchanged and
vanishes for $\log n\ge2d$. Hence, uniformly in all these
admissible $y,L$,
\[
 q_{e^L}(e^{iT\cdot}b_y,e^{iT\cdot}b_y)
       =\log(T/(2\pi))\|b\|_2^2+O_b(1).
\]
These individual bump diagonals are eventually positive.
Their separated cross terms and test shapes varying without
the derivative bounds are not controlled by this statement.
In particular, (JM) is not a lower bound on the entire
logarithmic form domain and proves neither (NP) nor an index bound.

\begin{corollary}[The joint modulation rest is independent of packet count]
\label{cor:dimension-free-rest}
For disjoint translates $b_j$ of one fixed norm-one
$b\in C_c^\infty((-d,d))$, the operator norm of the rest in (JM)
on their entire span is bounded by the scalar (JM) bound with
$\|b^{(r)}\|_2^2$, independently of the number $M$ of packets.
For $L\le\frac12\log T$, $t\in[T,2T]$, this rest is $O_b(T^{-1})$.
\end{corollary}
\begin{proof}
For each fixed derivative order,
$\|\sum v_jb_j^{(r)}\|_2^2=\|v\|_{\ell^2}^2\|b^{(r)}\|_2^2$.
Apply the quadratic estimate (JM) to the whole vector and take
the unit Rayleigh supremum of the Hermitian rest matrix.
With derivative order one, its pole term is
$O_b(Lt^{-2}e^L)=O_b((\log T)T^{-3/2})$.
The two gamma terms are $O_b(T^{-1})$ as well.
Disjoint support does not make the pole term zero.
\end{proof}

\section{Arithmetic mean values give growing positive packet compressions}
The height average below is in the modulation variable. Packet centers
and the test shape are chosen before taking that average.

\begin{lemma}[An elementary Dirichlet-polynomial mean bound]
\label{lem:elementary-mean}
For $X\ge2$, $H>0$, $U\in\R$, and arbitrary complex coefficients,
\[
 \int_U^{U+H}\left|\sum_{n\le X}a_n n^{it}\right|^2\dd t
 \le[H+4X(1+\log(2X))]\sum_{n\le X}|a_n|^2.       \tag{EM}
\]
There is an absolute $C_\psi$ such that, for $1\le A<B$,
\[
 \sum_{A<n\le B}\frac{\Lambda(n)^2}{n}
 \le C_\psi\log B\left(1+\log\frac BA\right).       \tag{LS}
\]
Both sums retain every prime power.
\end{lemma}
\begin{proof}
Expanding the square gives the diagonal $H\sum|a_n|^2$.
Each off-diagonal integral has magnitude at most
$2/|\log n-\log m|\le2X/|n-m|$.
Apply $2|a_na_m|\le|a_n|^2+|a_m|^2$ and sum the two
harmonic rows. Their total contribution is at most
$4XH_{\lfloor X\rfloor-1}\sum|a_n|^2$, which proves (EM).

For completeness the Chebyshev input $\psi(t)\le C_\psi t$
is unconditional and elementary. Write
$\vartheta(t)=\sum_{p\le t}\log p$.
Every prime in $(n,2n]$ divides $\binom{2n}{n}$, giving
$\vartheta(2n)-\vartheta(n)\le2n\log2$.
For real $t\ge2$, use $n=\lceil t/2\rceil$; at most one
integer is in $(t/2,n]$, so
$\vartheta(t)-\vartheta(t/2)\le(t+2)\log2+\log(t+2)$.
Summing over $t,t/2,t/4,\ldots$ proves $\vartheta(t)=O(t)$:
the errors are $O((1+\log t)^2)=O(t)$.
The identity $\psi(t)=\sum_{j\le\log_2t}\vartheta(t^{1/j})$
then gives $\psi(t)=O(t+\sqrt t\log t)=O(t)$.
Extend the bound to $1\le t<2$ using $\psi(t)=0$.
Partial summation now gives
\[
 \sum_{A<n\le B}\frac{\Lambda(n)}n
 =\frac{\psi(B)}B-\frac{\psi(A)}A
                   +\int_A^B\frac{\psi(u)}{u^2}\dd u
 \le C_\psi(1+\log(B/A)).
\]
Multiply by $\Lambda(n)\le\log B$ to prove (LS).
\end{proof}

\begin{theorem}[Positive finite compressions at most modulation heights]
\label{thm:mean-compressions}
Fix $d>0$, a real $b\in C_c^\infty((-d,d))$ with $\|b\|_2=1$,
and $0<c<1/2$. For each sufficiently large $T$, choose $L\ge1$
with $L\le c\log T$ and $M$ centers $x_j$ satisfying
\[
 |x_j|+d<L,\qquad |x_j-x_k|>2d\quad(j\ne k).
\]
Put $b_j(x)=b(x-x_j)$ and let the Hermitian matrix
\[
 \mathcal M(t)_{jk}=q_{e^L}(e^{it\cdot}b_j,e^{it\cdot}b_k),
 \qquad E(t)=\mathcal M(t)-\operatorname{diag}\mathcal M(t).
\]
Constants below depend only on $b,d,c$, not on $L,M,T$ or the centers.
Then
\[
 \frac1T\int_T^{2T}\|E(t)\|_{\rm HS}^2\dd t
                \le C_{b,d,c}M^2(1+L),                \tag{MC}
\]
and
\[
 \frac{|\{t\in[T,2T]:\mathcal M(t)\not\ge
                         \tfrac14\log T\,I_M\}|}{T}
       \le C_{b,d,c}\frac{M^2(1+L)}{(\log T)^2}.      \tag{MP}
\]
In particular if $M=M_T$ and
$M_T^2(1+L_T)=o((\log T)^2)$, these compressions are positive
at a set of heights of relative measure tending to one.
The exceptional set can depend on the selected centers.
\end{theorem}
\begin{proof}
Polarizing (JM), including its pole error, gives
\[
 \begin{split}
 q_{e^L}(f_t,g_t)={}&\log(t/(2\pi))\langle f,g\rangle\\
 &-\sum_{n\le e^{2L}}\frac{\Lambda(n)}{\sqrt n}
 \left[n^{it}C_{f,g}(\log n)+n^{-it}C_{f,g}(-\log n)\right]
 +\mathcal R_{L,t}(f,g),
 \end{split}                                         \tag{BM}
\]
where $f_t=e^{it\cdot}f$ and $g_t=e^{it\cdot}g$.
Indeed $C_{f_t,g_t}(r)=e^{itr}C_{f,g}(r)$; the two
phases in (BM) are different. For $f=b_j$, $g=b_k$,
polarization uses $b_j+i^r b_k$ ($0\le r\le3$), whose
derivative norms are bounded by twice those of $b$.
Use $N=1$ in (JM). Since $L\le c\log T$, it gives
\[
 \sup_{T\le t\le2T}|\mathcal R_{L,t}(b_j,b_k)|
 \le C_b\left(T^{-1}+T^{-2}\log(2+T)+Le^LT^{-2}\right)
 \le C_{b,c}T^{-1}.                                 \tag{BR}
\]
Support orthogonality cancels only the scalar main term.
It does not cancel the pole evaluations, which remain in (BR).

For $x_k>x_j$ set $s=x_k-x_j>2d$.
Then $C_{b_j,b_k}(r)=C_{b,b}(r-s)$, its magnitude is
at most one, and $C_{b_j,b_k}(-r)=0$ for $r>0$.
The nonzero coefficients in (BM) are thus supported on
the shell $A<n\le B$ with $A=e^{s-2d}>1$,
$B=e^{s+2d}<e^{2L}$. Endpoint coefficients vanish.
By (LS) their squared sum is at most
$C_\psi(s+2d)(1+4d)\le C_d(1+L)$.
With $X=e^{2L}$ and $H=T$, (EM) has normalized factor
\[
 1+4X(1+\log(2X))/T
 \le1+4T^{2c-1}(1+\log2+2c\log T),
\]
which is bounded for all sufficiently large $T$ because $c<1/2$.
Together with (BR), this proves
$T^{-1}\int_T^{2T}|\mathcal M(t)_{jk}|^2\dd t
\le C_{b,d,c}(1+L)$ for every $j\ne k$.
Reverse-ordered entries are complex conjugates.
Summing their squared magnitudes proves (MC).

The translation-invariant diagonal calculation following (JM)
gives, uniformly in all the centers,
$\mathcal M(t)_{jj}=\log(t/(2\pi))+O_b(1)$.
It is at least $\frac12\log T$ for all sufficiently large $T$.
Markov's inequality and (MC) show that
\[
 \frac{|\{t:\|E(t)\|_{\rm HS}>\tfrac14\log T\}|}{T}
 \le16C_{b,d,c}M^2(1+L)/(\log T)^2.
\]
On the complementary set $\|E(t)\|_{\rm op}\le\|E(t)\|_{\rm HS}$;
the diagonal lower bound therefore gives (MP).
\end{proof}

\begin{theorem}[Higher moments improve equally spaced compressions]
\label{thm:toeplitz-moments}
Keep the fixed $b,d$ above, fix $\ell>4d$, and take $M\ge2$ centers
$x_j=(j-(M-1)/2)\ell$, $0\le j<M$. Assume
$L> (M-1)\ell/2+d$, $L\ge1$, and put
\[
 X=e^{2L},\qquad Q_M=\lceil2\pi M^2\rceil,\qquad
 r=\lceil\log(8Q_M)\rceil.
\]
If $T$ is sufficiently large and
\[
 4X^r(1+\log(2X^r))\le T,                             \tag{ML}
\]
then outside a set of relative measure at most $2\cdot4^{-r}$
in $[T,2T]$, the full off-diagonal compression satisfies
\[
 \|E(t)\|_{\rm op}
 \le C\sqrt r\,\log(2X)+C(1+L)+C_bM/T.                \tag{TM}
\]
Here $C$ is absolute, and the last constant depends only on $b$;
all higher prime powers and the full pole term are included.
In particular, for fixed $\ell,d,b$, if $M_T\to\infty$,
$L_T\asymp M_T$, and
\[
 M_T=o\left(\frac{\log T}{\log\log T}\right),           \tag{TS}
\]
the full compression is at least $\frac14\log T\,I_{M_T}$
on a set of relative measure tending to one.
\end{theorem}
\begin{proof}
The logarithmic shells $[k\ell-2d,k\ell+2d]$, $1\le k<M$,
are disjoint. Split the exact prime-power matrix in (BM) into
its prime part and its higher-power part. For a prime in these
shells let $k(p)$ be its unique shell index and set
\[
 a_p=\frac{\log p}{\sqrt p}C_{b,b}(\log p-k(p)\ell),\quad
 A_X=\sum_p|a_p|^2,\quad
 S_\theta(t)=\sum_p a_pp^{it}e^{ik(p)\theta}.
\]
By (LS) and $|C_{b,b}|\le1$, $A_X\le C\log(2X)^2$.
For integer $r\ge1$, write
$S_\theta(t)^r=\sum_{n\le X^r}d_n(\theta)n^{it}$.
Unique prime factorization shows that a coefficient with
prime multiplicities $m_p$ is
\[
 \frac{r!}{\prod_p m_p!}\prod_p(a_pe^{ik(p)\theta})^{m_p}.
\]
Since $(r!/\prod m_p!)^2\le r!(r!/\prod m_p!)$, the
multinomial theorem proves
\[
 \sum_n|d_n(\theta)|^2\le r!A_X^r.
\]
Apply (EM) with length $X^r$, and then (ML), to obtain
\[
 T^{-1}\int_T^{2T}|S_\theta(t)|^{2r}\dd t\le2r!A_X^r
                                                          \tag{PM}
\]
uniformly in $\theta$. This argument concerns the prime part
only; the higher powers are retained separately below.

Here an elementary grid bound suffices, without an imported
Bernstein inequality. For a polynomial
$S(\theta)=\sum_{k=1}^{M-1}A_ke^{ik\theta}$, Fourier inversion
gives $|A_k|\le\|S\|_\infty$, hence
$\|S'\|_\infty\le M(M-1)\|S\|_\infty/2$.
The nearest point of the $Q_M$-point circular grid is at
distance at most $\pi/Q_M$. Therefore
$\|S\|_\infty\le2\max_{\theta\text{ on grid}}|S(\theta)|$.
If $A_X=0$ the prime part vanishes identically. Otherwise use (PM),
$r!\le r^r$, and Markov at threshold
$U=2\sqrt{rA_X}\,Q_M^{1/(2r)}$.
Each grid point has exceptional proportion at most
$2/(4^rQ_M)$. The union has proportion at most $2\cdot4^{-r}$,
and outside it
\[
 \sup_\theta|S_\theta(t)|\le4\sqrt e\,\sqrt{rA_X},
\]
since $r\ge\log Q_M$.

Let $P_k(t)$ be the coefficient of $e^{ik\theta}$ in $S_\theta(t)$.
The prime matrix entry with $j<j+k$ is $-P_k(t)$.
Its real Laurent symbol, with
$V(\theta)=\sum_jv_je^{ij\theta}$, is
$h(\theta,t)=-2\Re S_{-\theta}(t)$, and
\[
 v^*E_{\rm pr}(t)v
 =\frac1{2\pi}\int_0^{2\pi}h(\theta,t)|V(\theta)|^2\dd\theta.
\]
Thus $\|E_{\rm pr}(t)\|_{\rm op}\le2\sup_\theta|S_\theta(t)|$.

The higher-power symbol has norm at most twice
\[
 \sum_{p^m\le X,\,m\ge2}\frac{\log p}{p^{m/2}}
 \le \sum_{p\le\sqrt X}\frac{\log p}{p}
       +\sum_p\frac{\log p\,p^{-3/2}}{1-p^{-1/2}}
 \le C(1+\log X).
\]
The first term is bounded by partial summation of $\psi$;
the second converges by comparison with
$\sum_{n\ge2}(\log n)n^{-3/2}$.
Disjoint shells ensure that no prime power is counted twice.
This proves the second term in (TM).

Condition (ML) implies $L\le\frac12\log T$.
The $N=1$ polarized (JM) error is at most $C_b/T$ in
every entry for this range, by the calculation in (BR), now
with $c=1/2$. Its matrix norm is at most $C_bM/T$.
This factor $M$ is retained; entrywise convergence alone
would not justify growing-matrix convergence.
Together these estimates prove (TM).

Under (TS), $r=O(\log\log T)$ and $2Lr=o(\log T)$,
so (ML) holds for all sufficiently large $T$.
Moreover $(1+L)\sqrt r=o(\log T)$ and $M/T=o(1)$.
The bound (TM) is therefore $o(\log T)$, while the
diagonals are $\log(t/(2\pi))+O_b(1)$ uniformly.
This proves the positive lower bound. The exceptional
proportion tends to zero because $M_T\to\infty$.
The moment-length condition, rather than the weaker
$L\sqrt{\log M}=o(\log T)$ positivity condition alone,
limits the stated admissible scale.
\end{proof}

\begin{theorem}[Large compressions positive at every modulation height]
\label{thm:pointwise-dirichlet}
Fix $0<\kappa<\theta<1$, $T\ge4$, and $L\ge1$.
For $f\in H_0^1(-L,L)$, extended by zero, suppose
$\|f'\|_2\le\kappa T\|f\|_2$. For every $t\in[T,2T]$,
\[
 q_{e^L}(e^{it\cdot}f)\ge
 \left[(1-\kappa^2/\theta^2)\log T-C_\theta-4C_\psi e^L
           -4Le^L(\kappa+1/(2T))^2\right]\|f\|_2^2.     \tag{PD}
\]
Thus if $Le^L=o(\log T)$ the first
$M=\lfloor2\kappa LT/\pi\rfloor$ Dirichlet modes have positive
full-form compressions for every height in the band.
For $L=(1-\eta)\log\log T$, $0<\eta<1$, this gives
$M\sim2\kappa(1-\eta)T\log\log T/\pi$.
\end{theorem}
\begin{proof}
The gamma multiplier $a$ is globally bounded below by $-C_a$
and satisfies $a(t+u)\ge\log t-C_\theta$ for $|u|\le\theta t$;
enlarge $C_\theta\ge C_a$ using the already controlled digamma
asymptotic. The Fourier mass outside that band is at most
$\|f'\|_2^2/(\theta t)^2$, by Plancherel and zero extension.
The archimedean form is consequently at least
$[(1-\kappa^2/\theta^2)\log T-C_\theta]\|f\|_2^2$.

Partial summation of the elementary $\psi(X)\le C_\psi X$ bound gives
$\sum_{n\le X}\Lambda(n)n^{-1/2}\le2C_\psi\sqrt X$.
Every zero-extended translation is a contraction. The whole
prime-power form, including both channels, therefore has absolute
value at most $4C_\psi e^L\|f\|_2^2$.
One weak integration by parts, with zero boundary trace, gives
\[
 |\widehat{e^{it\cdot}f}(\pm i/2)|
 \le t^{-1}e^{L/2}\sqrt{2L}
                         (\|f'\|_2+\tfrac12\|f\|_2).
\]
Its two pole terms have absolute sum at most the last cost in (PD).
The zero extension really belongs to $H^1(\R)$; smooth approximation
therefore justifies every form step for $H_0^1$ inputs.
The normalized modes
$e_j=L^{-1/2}\sin(\pi j(x+L)/(2L))\mathbf1_{(-L,L)}$
have derivative norm on their first $M$-dimensional span bounded by
$\pi M/(2L)$. This proves the assertions.
\end{proof}

\begin{proposition}[The large positive blocks miss a fixed detector]
\label{prop:dirichlet-escape}
For $h\in C_c^\infty((-L,L))$, let $\Pi_{t,M}$ be the $L^2$ projection
onto the modulated Dirichlet span in the preceding theorem.
For every integer $N\ge1$, uniformly for $t\in[T,2T]$,
\[
 \|\Pi_{t,M}h\|_2^2
 \le \frac{2\kappa}{\pi(1-\kappa)^{2N}}\,
                 T^{1-2N}\|h^{(N)}\|_1^2.              \tag{DE}
\]
In particular positivity of these growing blocks does not supply
a projection approximation to an individual compact negative direction.
\end{proposition}
\begin{proof}
With $q_j=\pi j/(2L)$, direct Fourier integration gives
\[
 \langle e^{it\cdot}e_j,h\rangle
 =\frac{e^{iq_jL}\widehat h(t-q_j)
                 -e^{-iq_jL}\widehat h(t+q_j)}{2i\sqrt L}.
\]
Both arguments have modulus at least $(1-\kappa)T$.
Integration by parts gives
$|\widehat h(u)|\le\|h^{(N)}\|_1|u|^{-N}$.
Square, sum over $j\le M$, and use $M/L\le2\kappa T/\pi$.
\end{proof}

\begin{corollary}[A finite set of zeros disappears from a height average]
\label{cor:averaged-zero-loss}
For any finite set $S$ of distinct Xi zeros $z=\gamma+i\delta$,
any $f\in L^2([-L,L])$, and any real interval $I$ of length $H>0$,
\[
 \frac1H\int_I\sum_{z\in S}m_z
       |\widehat f(\gamma-t+i\delta)|^2\dd t
 \le\frac{2\pi}{H}\sum_{z\in S}m_z e^{2|\delta|L}\|f\|_2^2.
                                                               \tag{AZ}
\]
If $H=T$, $L\le c\log T$, and $\|f\|_2$ is uniformly bounded,
the right side is $O_S(T^{c-1})$. Thus for $c<1$ it tends to
zero even if $S$ contains an off-line zero.
\end{corollary}
\begin{proof}
Compact support makes $f$ integrable and its transform entire.
Plancherel on $e^{\delta x}f(x)$ gives
\[
 \int_\R|\widehat f(u+i\delta)|^2\dd u
       =2\pi\int_{-L}^Le^{2\delta x}|f(x)|^2\dd x
       \le2\pi e^{2|\delta|L}\|f\|_2^2.
\]
Each interval integral is bounded by this full real-line integral.
Sum over $S$ and use $|\delta|<1/2$ for nontrivial zeta zeros.
The finite-set qualification is essential; no infinite-zero limit
is interchanged with the height average.
\end{proof}

These arithmetic mean estimates are stronger than an absolute
prime-shell bound for the chosen packet families. They do not give
an upper bound on the full negative index or on $\nu_{e^L}$:
other vectors and exceptional heights remain uncontrolled.

\begin{proposition}[The critical-line sampling is not an unweighted frame]
\label{prop:critical-frame}
For $L>0$, define on smooth compact tests
\[
 P_{c,L}(f)=\sum_{z=\gamma\in\mathcal Z\cap\R}
                         m_z|\widehat f(\gamma)|^2.
\]
For no fixed $L>0$ is there a finite $B_L$ such that
$P_{c,L}(f)\le B_L\|f\|_2^2$ for all
$f\in C_c^\infty((-L,L))$.
If an off-line zeta zero exists, there are constants $C,c,L_0>0$
such that for every $L\ge L_0$,
\[
 \inf_{0\ne f\in C_c^\infty((-L,L))}
              \frac{P_{c,L}(f)}{\|f\|_2^2}
                  \le C e^{-c e^{2L}}.                 \tag{FC}
\]
The second assertion is conditional on that zero's existence.
\end{proposition}
\begin{proof}
Conrey's unconditional proportion theorem and zero-count asymptotic
\cite[eqs.~(8)--(12), printed p.~3; Theorem 1, printed p.~4]{Conrey}
give $N_c(T)\ge c_0T\log T$ for sufficiently large $T$, where
$N_c$ counts positive-height critical zeros with multiplicity.
There are $c_1>0$ and arbitrarily large dyadic $T$ such that
$N_c(2T)-N_c(T)\ge c_1T\log T$.
Indeed otherwise summing bounds $\epsilon2^j j$ over
dyadic intervals would give $N_c(2^R)\le C+2\epsilon R2^R$,
contradicting the cumulative lower bound for sufficiently small
$\epsilon<c_0\log2/4$.

Fix a nonzero $b\in C_c^\infty((-L,L))$ and $A>0$ with
$\mu_A=\int_{-A}^A|\widehat b(u)|^2\dd u>0$.
Removing the endpoint intervals of length $A$ removes only
$O_A(\log T)$ zeros by local counting. Thus a bulk of
$\gg T\log T$ critical zeros lies in $[T+A,2T-A]$
along the preceding sequence.
Tonelli and $\widehat{e^{it\cdot}b}(\gamma)=\widehat b(\gamma-t)$ give
\[
 \int_T^{2T}P_{c,L}(e^{it\cdot}b)\dd t
 \ge \mu_A\,[N_c(2T-A)-N_c(T+A)]\gg_b T\log T.
\]
An $L^2$ upper bound would instead make this at most
$B_LT\|b\|_2^2$, a contradiction.

For (FC), select one off-line zero and use (ID) with value one
there and zero at every other distinct zero. The resulting
nonzero smooth $f$ has $\|f\|_2>0$ and annihilates every
critical zero. Its cutoffs satisfy (TE).
For $L\ge3$ take $R=L-2$, so $f_R=\chi_Rf$ is compactly
supported in $[-L+1,L-1]\subset(-L,L)$.
Since the critical values of $f_R$ equal the negatives of
those of $(1-\chi_R)f$,
\[
 P_{c,L}(f_R)\le S_4 E_R^2\le C_f e^{-c_f e^{2R}}.
\]
The cutoff also converges in $L^2$, hence
$\|f_R\|_2\ge\|f\|_2/2$ for large $L$.
Divide and absorb $R=L-2$ into the constants to obtain (FC).
Exact zero orders in (ID) retain every multiplicity.
\end{proof}

The upper-frame failure holds regardless of RH and is compatible
with the log-weighted bound (ZS). The conditional lower-frame
collapse and (AZ) explain why positive critical-line sampling
and modulation averages alone cannot exclude each individual zero.

\section{Positive full compressions on multiplication meshes}
\label{sec:multiplication-mesh}
We prove positivity on growing families that retain every arithmetic
edge between their cells. The pole correction is included on the
whole compression, without a moment condition. The support measure
still tends to zero, so the result supplies no full-form exhaustion.

For an integer $N\ge3$ put
\[
 r_N=\frac1{100N(\log N)^2},\qquad
 I_n=\log n+[-r_N,r_N],\qquad E_N=\bigcup_{n=1}^N I_n,
 \qquad h_N=\sum_{n=1}^N\frac1n,
\]
and define
\[
 p_N=2\sinh(r_N)\sqrt{N(N+1)h_N/2}-2r_NN. \tag{MM1}
\]
Here $H_{\log}(\R)$ has the Fourier domain used above, and $q_A$
is the actual physical Weil form, with both poles and all prime powers.
For a compactly supported test its prime sum is finite by support,
rather than by an additional truncation.

\begin{theorem}[Every vector in the full multiplication compression]
\label{thm:full-multiplication-mesh}
For every integer $N\ge3$ and every complex $f\in H_{\log}(\R)$
supported in $E_N$,
\[
 q_A(f)\ge\{2\log\log N-6-p_N\}\|f\|_2^2. \tag{MM2}
\]
Moreover $0\le p_N<2/25$ and $p_N=O((\log N)^{-3/2})$.
The bound holds for every translate of $E_N$, with the same constants,
and at every modulation frequency. No local or global pole moment
is required.
\end{theorem}
\begin{proof}
We first establish the exact logarithmic support cost of the gamma
form. Write $\psi=\Gamma'/\Gamma$ and $C_0=-a(0)$.
The digamma series \cite[eq.~(5.7.6), away from the poles]{DLMFGamma}
gives, with $b_j=j+1/4$,
\[
 a(u)+C_0=\sum_{j=0}^{\infty}
     \frac{(u/2)^2}{b_j\{b_j^2+(u/2)^2\}}
 \ge\int_{1/4}^{\infty}
     \frac{(u/2)^2}{b\{b^2+(u/2)^2\}}\dd b
 =\frac12\log(1+4u^2)\ge\log|u|. \tag{MM3}
\]
The summand is nonnegative and decreasing in $b$, which justifies
the integral comparison. The last inequality is for $u\ne0$.
If $F$ is supported in an interval of radius $r$, translate its
center to zero and write $F(x)=r^{-1/2}g(x/r)$, with
$\operatorname{supp}g\subset[-1,1]$. Plancherel and (MM3) yield
\[
 q_{\rm arch}(F)\ge
 \{\log(1/r)-C_0\}\|g\|_2^2+
       \int_\R\log|v|\,|\widehat g(v)|^2\frac{\dd v}{2\pi}.
\]
On $|v|<1$, $|\widehat g(v)|\le\sqrt2\|g\|_2$, and
$\int_0^1-\log v\dd v=1$. Discarding the positive part of the
last integral therefore proves
\[
 q_{\rm arch}(F)\ge
       \{\log(1/r)-C_0-2/\pi\}\|F\|_2^2. \tag{MM4}
\]
All integrals converge: the logarithmic norm controls infinity,
and compact support bounds the Fourier transform near zero.

For completeness, the numerical constant needed here has an
elementary exact bound. Specialize
\cite[eq.~(5.4.19), $p=1,q=4$]{DLMFGamma} to obtain
$C_0=\gamma_E+\pi/2+3\log2+\log\pi$. The identity
\[
 H_m-\log m-\gamma_E
 =\sum_{j=m}^{\infty}
       \{\log(1+1/j)-1/(j+1)\}
 \ge\frac12\sum_{j=m}^{\infty}(j+1)^{-2}
 \ge\frac1{2(m+1)}
\]
follows by integrating
$1/t-1/(j+1)\ge(j+1-t)/(j+1)^2$ on $[j,j+1]$.
At $m=8$, since $H_8=761/280$ and $\log8=3\log2$, this gives
\[
 C_0\le\frac{6709}{2520}+\frac\pi2+\log\pi
 <\frac{6709}{2520}+\frac{11}7+\frac{23}{20}
 =\frac{13567}{2520}<\frac{27}5. \tag{MM5}
\]
The bounds used in this strict inequality follow from
$22/7-\pi=\int_0^1x^4(1-x)^4/(1+x^2)\dd x>0$ and
\[
 e^{23/20}>\sum_{j=0}^{6}\frac{(23/20)^j}{j!}
 =\frac{145501209049}{46080000000}>\frac{22}7.
\]
Also $\pi>3$, for example by integrating the lower geometric
polynomial through degree $14$:
$\pi>4\sum_{j=0}^7(-1)^j/(2j+1)=135904/45045>3$.

We next keep all gamma interactions between distinct cells.
The multiplier has the exact representation
\[
 a(u)=a(0)+2\int_0^\infty k(d)(1-\cos ud)\dd d,\qquad
 k(d)=\frac{e^{-d/2}}{1-e^{-2d}}. \tag{MM6}
\]
Expand its denominator geometrically and integrate each nonnegative
summand to recover (MM3); Tonelli proves the representation.
For separated supports the off-diagonal physical kernel is
$-k(|x-y|)$, with no diagonal finite-part ambiguity.
For all $d>0$, $k(d)\le4/d$. Indeed for $d\le1$,
$1-e^{-2d}\ge d(1-e^{-2})>d/2$, and for $d\ge1$ the
denominator exceeds $1/2$ while $e^{d/2}>d/2$.
The cell separation is
$|\log n-\log m|\ge|n-m|/N$, and $2r_N<1/(2N)$.
Consequently
\[
 \sup_{x\in I_n}\sum_{m\ne n}\int_{I_m}k(|x-y|)\dd y
 \le32r_NN\sum_{j=1}^{N-1}\frac1j
 \le\kappa_N:=32r_NN(1+\log N)\le16/25. \tag{MM7}
\]
The symmetric kernel Schur inequality, obtained by integrating
$2|f(x)f(y)|\le|f(x)|^2+|f(y)|^2$, bounds the whole
off-diagonal block by $\kappa_N$. Hence for $f_n=1_{I_n}f$,
\[
 q_{\rm arch}(f)\ge
       \sum_{n=1}^N q_{\rm arch}(f_n)-\kappa_N\|f\|_2^2
 \ge\{\log(1/r_N)-C_0-2/\pi-\kappa_N\}\|f\|_2^2. \tag{MM8}
\]
Each $f_n$ is in $H_{\log}$: disjoint smooth separators equal to
one on their cell and zero on the other cells give these restrictions
as smooth compact multipliers. Their boundedness on $H_{\log}$
follows from weighted Fourier Young; only finitely many are used
at each fixed $N$. Thus (MM8) applies to the stated whole domain.

Every prime-power edge is retained exactly. A shift $\log k$,
$k\ge2$, joins $I_m$ to $I_n$ precisely when $n=km$.
Otherwise
$|\log(km/n)|\ge\log(1+1/n)>1/(N+1)>2r_N$ when $km>n$;
the reverse case has the still larger separation
$\log(n/(n-1))$. This also treats $km>N$.
Under the unitary identification $f_n(t)=f(\log n+t)$,
$t\in[-r_N,r_N]$, the compression of $\Pi$ is $T_N\otimes I$,
where
\[
 (T_N)_{n,m}=
 \begin{cases}
   \Lambda(k)/\sqrt k,& n=km\ \hbox{or }m=kn,\quad k\ge2,\\
   0,&\hbox{otherwise}.
 \end{cases} \tag{MM9}
\]
We prove the needed row bound without a prime number theorem.
The factorial divisor identity gives, for every $X\ge1$,
\[
 \sum_{d\le X}\Lambda(d)\lfloor X/d\rfloor
   =\sum_{j\le X}\log j\le X\log X.
\]
For integers $M\ge1$, prime valuations of $\binom{2M}{M}$ give
$\psi(2M)-\psi(M)\le\log\binom{2M}{M}\le2M\log2$.
Dyadic summation and monotonicity imply
$\psi(X)<4X\log2<3X$ for $X>1$; the case $X=1$ is zero.
Here $\log2<3/4$ follows already from the first four positive
Taylor terms of $e^{3/4}$. It follows that
\[
 S(X):=\sum_{2\le d\le X}\frac{\Lambda(d)}d\le\log X+3.
\]
For Schur weights $w_n=n^{-1/2}$ the exact divisor sum now gives
\[
 \frac{(T_Nw)_n}{w_n}
   =S(N/n)+\sum_{\substack{k\mid n\\ k\ge2}}\Lambda(k)
   \le\log(N/n)+3+\log n=\log N+3. \tag{MM10}
\]
Applying $2|v_nv_m|\le(w_m/w_n)|v_n|^2+
(w_n/w_m)|v_m|^2$ and integrating over $t$ proves
$|\langle f,\Pi f\rangle|\le(\log N+3)\|f\|_2^2$.
The argument includes every power of every prime.
By (MM8), (MM10), $\log(1/r_N)=\log N+2\log\log N+\log100$,
$\log100>4$, (MM5), and $\pi>3$, we obtain
\[
 q_B(f)\ge
 \{2\log\log N+\log100-C_0-2/\pi-3-\kappa_N\}\|f\|_2^2
 \ge(2\log\log N-6)\|f\|_2^2. \tag{MM11}
\]
For the last inequality the total constant loss is strictly below
$27/5+2/3+3+16/25-4=428/75<6$.
The inequality $\log100>4$ follows from $e<3$.

Finally retain both poles together. On $L^2(E_N)$ put
$u(x)=e^{x/2}$ and $v(x)=e^{-x/2}$. Their full correction is
\[
 P_E=|u\rangle\langle v|+|v\rangle\langle u|
       =2|c_E\rangle\langle c_E|-2|s_E\rangle\langle s_E|.
\]
With $\alpha_0=\langle u,v\rangle=2r_NN$,
$U=\|u\|^2=\sinh(r_N)N(N+1)$ and
$V=\|v\|^2=2\sinh(r_N)h_N$, its two nonzero eigenvalues
are $\alpha_0\pm\sqrt{UV}$, by its matrix
$\left(\begin{smallmatrix}\alpha_0&V\\U&\alpha_0\end{smallmatrix}\right)$
on $\operatorname{span}(u,v)$. Cauchy--Schwarz gives
$\sqrt{UV}\ge\alpha_0$. Thus $q_{\rm pole}(f)\ge-p_N\|f\|^2$.
Moreover
\[
 p_N\le2r_Ne^{r_N}N\sqrt{1+\log N}
 =\frac{e^{r_N}}{50}\,
       \frac{\sqrt{1+\log N}}{(\log N)^2}<\frac2{25}. \tag{MM12}
\]
Use $\sinh r\le re^r$, $h_N\le1+\log N$ and
$(N+1)/(2N)\le1$. Since $r_N<1/300$, $e^{r_N}<2$,
and $\sqrt{1+x}/x^2<2$ for $x>1$, the displayed strict
bound holds. It also proves the stated asymptotic estimate.
Adding this exact pole bound to (MM11) proves (MM2).
Translation multiplies $U,V$ by reciprocal factors and leaves
$\alpha_0$ unchanged; the gamma and prime forms are translation
invariant. Modulation changes none of the support hypotheses.
\end{proof}

The compressions are positive for $2\log\log N>6+p_N$, but
$|E_N|=2r_NN=1/[50(\log N)^2]\to0$. For every fixed $F\in L^2(\R)$,
\[
 \sup_{a\in\R}\|1_{a+E_N}F\|_2\longrightarrow0. \tag{MM13}
\]
This is absolute continuity of the integral of $|F|^2$ on sets
of vanishing measure, uniformly in the set's location.
The projections therefore converge strongly to zero, including
all translated choices. A transfer from these positive compressions
to the full form must reconstruct the missing interactions and
their localization cost; (MM2) alone cannot supply that transfer.

\part{Complete arithmetic realizations and source transfers}

\section{An exact theta ground transformation of the full Weil form}
Here $\Phi$ is the complete positive smooth Xi density reconstructed
above, not a finite theta head. On the whole line put
\[
 c(x)=\cosh(x/2),\quad s(x)=\sinh(x/2),\quad
 \dd\mu(x)=c(x)\Phi(x)\dd x,\quad
 \dd\nu=2\dd\mu,\quad \kappa(x)=\tanh(x/2).
\]
The exact values $\Xi(\pm i/2)=1/2$ and evenness give
$\mu(\R)=1/2$ and $\int s\Phi=0$.
Thus $\nu$ is an even probability measure and
$\int\kappa\dd\nu=0$.
For $g\in C_c^\infty(\R)$ define the positive energy
\[
 \begin{split}
 \mathcal D_\Phi(g)={}&\frac12\iint_{\R^2}
       k(|x-y|)\Phi(x)\Phi(y)|g(x)-g(y)|^2\dd x\dd y\\
 &+\sum_{n\ge2}\frac{\Lambda(n)}{\sqrt n}
       \int_\R\Phi(x)\Phi(x-\log n)
                    |g(x)-g(x-\log n)|^2\dd x,\qquad
 k(r)=\frac{e^{-r/2}}{1-e^{-2r}}.
 \end{split}                                                 \tag{DPE}
\]
The sum retains every prime power. Its outgoing jumps extend
beyond the support of $g$; they are not cut off at $e^{2L}$.
When $g$ is compactly supported the energy is finite:
the small-distance difference is $O(|x-y|)$, and the
large-index tails are $O_g(e^{-c_gn^2})$ whenever at least
one endpoint is in that compact support.

\begin{theorem}[The exact theta ground identity]
\label{thm:theta-ground}
For $g\in C_c^\infty(\R)$ and $f=\Phi g$, in any interval
containing its support,
\[
 q_B(f)=\mathcal D_\Phi(g)-\int|g|^2\dd\mu,\qquad
 q_A(f)=\mathcal D_\Phi(g)
   -\frac12\left\{\operatorname{Var}_\nu(g)
               +\left|\int\kappa g\dd\nu\right|^2\right\},
                                                               \tag{TG}
\]
where
$\operatorname{Var}_\nu(g)=\int|g|^2\dd\nu-
|\int g\dd\nu|^2$. The identity includes the full gamma term,
both pole terms, and all prime powers.
\end{theorem}
\begin{proof}
First we prove the local, distributional all-prime identity
$\mathcal B\Phi=-c$. This does not put the untruncated $\Phi$
in the Hilbert space of a fixed finite-interval operator.
For a compact smooth test $h$, (TE) gives
$\mathcal T(\chi_R\Phi)\to0$, since $\widehat\Phi=\Xi$ vanishes
at every Xi zero. Therefore (ZF) gives
$q_A(\chi_R\Phi,h)\to0$.
The two pole values tend to $1/2$, so the pole cross-term
tends to $\int c(x)h(x)\dd x$.
The archimedean cross-term converges in the log-Fourier norm.
For the prime cross-terms, uniformly on the fixed support of $h$,
$|\Phi^{(j)}(x\pm\log n)|\le C_{h,j}e^{-c_{h,j}n^2}$
for every fixed $j$ and all large $n$. Consequently the all-prime
series is locally absolutely convergent, also after differentiation.
The digamma jump representation now yields
\[
 \begin{split}
 \mathcal B\Phi(x)={}&a(0)\Phi(x)
   +\int_0^\infty k(r)
        \{2\Phi(x)-\Phi(x-r)-\Phi(x+r)\}\dd r\\
 &-\sum_{n\ge2}\frac{\Lambda(n)}{\sqrt n}
        \{\Phi(x-\log n)+\Phi(x+\log n)\}=-c(x).
 \end{split}                                                 \tag{GDI}
\]
The integral is locally convergent: its second difference is
$O(r^2)$ at zero, and its tail is integrable.
Thus (GDI) is a well-defined local distributional, indeed smooth,
identity; no whole-line selfadjoint operator limit is assumed.

For the ground transformation, initially truncate the archimedean
distances to $\epsilon\le r\le R$ and the prime sum to $n\le N$.
For positive $\phi_x,\phi_y$ the elementary identity
\[
 |\phi_x g_x-\phi_y g_y|^2
 -\phi_x\phi_y|g_x-g_y|^2
 =(\phi_x-\phi_y)
                   (\phi_x|g_x|^2-\phi_y|g_y|^2)
\]
gives the archimedean potential after symmetric integration.
For each prime translation its negative cross-term, after
subtracting its contribution to
$\int(\mathcal B\Phi/\Phi)|f|^2$, gives exactly
\[
 \frac{\Lambda(n)}{\sqrt n}
 \int\Phi(x)\Phi(x-\log n)
                         |g(x)-g(x-\log n)|^2\dd x.
\]
The scalar $a(0)$ cancels in the difference.
Let the truncations disappear. Near zero use the $O(r^2)$
majorant. At large distances use the theta decay; the positive
prime energy is justified by Tonelli and its incoming potential
converges uniformly on the compact support of $g$.
The original $q_B(f)$ itself stabilizes once $N>e^{2L}$
if that support is inside $[-L,L]$, because larger translations
have no overlap. Hence the limiting algebra gives
\[
 q_B(f)-\int\frac{\mathcal B\Phi}{\Phi}|f|^2
                                     =\mathcal D_\Phi(g).
\]
Using (GDI) proves the first identity of (TG).
Restoring the exact pole term adds
$2|\int c\Phi g|^2-2|\int s\Phi g|^2$.
Since $\mu(\R)=1/2$, $\nu=2\mu$, and
$s\Phi\dd x=\kappa\dd\mu$, this is precisely the second identity.
\end{proof}

The remaining sharp global lower bound in this construction is
\[
 2\mathcal D_\Phi(g)\ge
       \operatorname{Var}_\nu(g)
                 +\left|\int\kappa g\dd\nu\right|^2
                  \qquad(g\in C_c^\infty(\R)).             \tag{GP}
\]
We have not proved (GP).
If it held, every compact smooth $f$ could be written
$f=\Phi g$, since $\Phi$ is positive on its compact support.
Then (TG) would make every such full Weil form nonnegative,
and the exact individual-zero detection theorem would exclude
every off-line zero. For even $g$ the $\kappa$ term vanishes;
for odd $g$ the mean vanishes, but the $\kappa$ term remains.
Discarding that negative odd pole correction would change the
required statement.

We next prove two unconditional structural results for this energy:
the ordinary sharp gap would already supply the full odd pole
correction, and the actual joint energy has a positive, possibly
smaller, gap. The value one remains a separate obligation.
Let $H_\mu=L^2(\mu)$, with inner product antilinear in its first
variable, and let $\mathcal Q_\mu$ be the closure of
$C_c^\infty(\R)$ for
$(\mathcal D_\Phi+\|\cdot\|_{H_\mu}^2)^{1/2}$.
This positive form is closable: an $H_\mu$-convergent,
form-Cauchy sequence has a subsequence converging almost everywhere;
the archimedean edge differences and all countably many translated
prime differences identify their edge-$L^2$ limits with that same
function. Let $L_\mu\ge0$ denote its selfadjoint operator and put
$Q_\mu=\mathcal D_\Phi-\langle\cdot,\cdot\rangle_{H_\mu}$.

\begin{proposition}[The score identity and the ordinary sharp-gap bridge]
\label{prop:theta-ordinary-gap}
The functions $1$ and $v=-2\Phi'/\Phi$ belong to $\mathcal Q_\mu$,
and
\[
 L_\mu1=0,\quad \int v\dd\mu=0,\quad
 \langle\kappa,v\rangle_{H_\mu}=\frac12,\quad
 Q_\mu(v,g)=\langle\kappa,g\rangle_{H_\mu},\quad
 Q_\mu(v,v)=\frac12.                                      \tag{SCI}
\]
In particular $v\in D(L_\mu)$ and $L_\mu v=v+\kappa$.
The full inequality (GP) is equivalent to
\[
 \mathcal D_\Phi(g)\ge\|g\|_{H_\mu}^2
 \quad\left(g\in\mathcal Q_\mu,\ \int g\dd\mu=0\right).
                                                               \tag{OGP}
\]
It suffices to require (OGP) on compact smooth mean-zero tests.
This proves the implication and equivalence; it does not prove
their common sharp-gap premise.
\end{proposition}
\begin{proof}
The positive theta series and its differentiated terms give, for
each fixed $j$, $|\Phi^{(j)}|/\Phi\le C_j(1+e^{2|x|})^j$.
On either tail the first term lower-bounds $\Phi$ and absorbs
the remaining Gaussian terms; on a compact use its positive minimum.
Thus $v,v'$ have at most exponential growth. The cutoff errors
for $1$ and $v$ in the archimedean energy are dominated by
\[
 C\Phi(x)\Phi(y)k(|x-y|)\min\{1,|x-y|^2\}
               (e^{A|x|}+e^{A|y|})
\]
for some fixed $A$, an integrable function. For prime shifts $r$
the integrated majorant is at most
$Ce^{A|r|}e^{-b e^{|r|}}$: split off half of each theta exponent
and use $\max\{|x|,|x-r|\}\ge|r|/2$ for the other half.
It sums every $\Lambda(n)/\sqrt n$ at $r=\log n$.
Dominated convergence therefore puts $1,v$ in the stated core
closure, and the energy of $1$ is zero.

Integration by parts, with these actual tails, gives the two
moments in (SCI). For $f=\Phi v=-2\Phi'$,
$\widehat f(z)=-2iz\Xi(z)$ vanishes at every zero, with each
full multiplicity, and its two pole values are $(1/2,-1/2)$.
The all-height cutoff bound (TE) and (ZF) imply
$q_A(f,\Phi g)=0$ for every compact smooth $g$.
The negative of its pole cross-term is
\[
 -q_P(f,\Phi g)=\frac12\{
       \widehat{\Phi g}(i/2)-\widehat{\Phi g}(-i/2)\}
                       =\langle\kappa,g\rangle_{H_\mu}.
\]
The polarized ground identity gives (SCI) on the core and
form continuity gives it on $\mathcal Q_\mu$.
Its right side is $H_\mu$-bounded, so $v$ is in the operator
domain with the asserted action.

Suppose (OGP). Then $Q_\mu$ is positive semidefinite on the
mean-zero form domain. Its Cauchy--Schwarz inequality and (SCI)
give
\[
 |\langle\kappa,g\rangle_{H_\mu}|^2
       \le Q_\mu(v,v)Q_\mu(g,g)=\tfrac12Q_\mu(g,g).
\]
For any compact smooth $g$ set $g_0=g-2\int g\dd\mu$.
The proved constant exhaustion puts $g_0$ in the form domain;
$\mathcal D_\Phi(g_0)=\mathcal D_\Phi(g)$ and
$\|g_0\|^2=\|g\|^2-2|\int g\dd\mu|^2$.
Also $\langle\kappa,g_0\rangle=\langle\kappa,g\rangle$.
The preceding inequality is exactly (GP), retaining its odd
pole correction. Conversely (GP) immediately gives (OGP) on
the mean-zero compact core. That core is form dense in the
mean-zero domain: correct the mean error of a core approximant
by a fixed compact smooth bump of $\mu$-integral one.
This proves the full converse and the stated core reduction.
\end{proof}

\begin{theorem}[An unconditional positive gap for the joint theta energy]
\label{thm:theta-positive-gap}
There is a kernel-defined constant $c_*>0$, without RH, such that
\[
 \mathcal D_\Phi(g)\ge c_*\|g\|_{H_\mu}^2
 \quad\left(g\in\mathcal Q_\mu,\ \int g\dd\mu=0\right).
                                                               \tag{JPG}
\]
The proof retains the joint archimedean and all-prime energies.
It supplies no assertion that $c_*=1$.
\end{theorem}
\begin{proof}
Choose $R\ge3$, put $b=R-1$, and let
$C=[-R,R]$, $I=[-b,b]$, $T=\R\setminus C$.
The joint degree from $x\ge R$ into $I$ is
\[
 V_R(x)=\int_I k(|x-y|)\Phi(y)\dd y
   +\sum_{x-\log n\in I}\frac{\Lambda(n)}{\sqrt n}
                                        \Phi(x-\log n).
\]
We first verify its uniform lower bound. Let $S_{\rm out}(x)$
be the complete prime degree whose other endpoint lies outside
$I$. Since the distance from $x$ to $I$ is at least one,
the exact local identity (GDI), split at $r=1$, gives
\[
\begin{split}
 V_R(x)-c(x)={}&a(0)\Phi(x)
 +\int_0^1k(r)\{2\Phi(x)-\Phi(x-r)-\Phi(x+r)\}\dd r\\
 &+2\Phi(x)\int_1^\infty k(r)\dd r
 -\int_{\substack{y\notin I\\|x-y|\ge1}}
                         k(|x-y|)\Phi(y)\dd y-S_{\rm out}(x).
\end{split}                                                    \tag{JRD}
\]
This preserves both tails. The local terms, divided by $c(x)$,
are bounded by $C_k\sup_{|u|\ge b}(|\Phi(u)|+|\Phi''(u)|)$;
the second difference supplies $r^2$ at zero.
Since $k(r)\le Ce^{-r/2}$ for $r\ge1$ and
$e^{x/2}/c(x)\le2$, the remaining archimedean term is at most
$C_k\int_{|y|\ge b}e^{|y|/2}\Phi(y)\dd y$.
The Chebyshev--variation estimate in the proof of
Lemma~\ref{lem:theta-prime-schur}, applied to
$u^{-1/2}\Phi(\log u)\mathbf1_{|\log u|\ge b}$ including
its endpoint atoms, gives
\[
 S_{\rm out}(x)/c(x)\le20C_\psi e^{b/2}\Phi(b).
\]
Consequently $|V_R(x)/c(x)-1|\le\varepsilon_R$, uniformly
for $x\ge R$, where one may take
\[
 \varepsilon_R=C_k\left\{
    \sup_{|u|\ge b}(|\Phi(u)|+|\Phi''(u)|)
     +\int_{|y|\ge b}e^{|y|/2}\Phi(y)\dd y\right\}
       +20C_\psi e^{b/2}\Phi(b).
\]
Here $C_k$ depends only on $a(0)$ and the displayed integrals
of $k$; all terms tend to zero doubly exponentially.
Reflection gives the same bound on $x\le-R$.
Fix one $R$ for which $a_R:=1-\varepsilon_R>0$.

The reverse degree from $y\in I$ into $T$ is at most
\[
 N_R:=k(1)\Xi(0)+C_{\rm pr}<\infty,
\]
where $C_{\rm pr}$ is the all-prime Schur constant in
Lemma~\ref{lem:theta-prime-schur}. The archimedean bound uses
the distance at least one and $c(y)\ge1$; the prime bound
retains its full positive row, including every prime power.
The detailed-balance conductance is
$k(|x-y|)\Phi(x)\Phi(y)\dd x\dd y$ for archimedean edges
and $\Lambda(n)n^{-1/2}\Phi(x)\Phi(x-\log n)\dd x$
for each prime edge. Thus these forward and reverse degree
bounds refer to the same energy and the same measure $\mu$.

For a compact smooth complex $g$ let
$m_C=\mu(C)^{-1}\int_Cg\dd\mu$ and
$V_C=\int_C|g-m_C|^2\dd\mu$.
The archimedean energy on $C\times C$ alone satisfies
\[
 D_{CC}(g)\ge b_R V_C,\qquad
 b_R=\frac{k(2R)\mu(C)}{\cosh(R/2)^2}>0:
\]
use $k(|x-y|)\ge k(2R)$,
$\Phi\dd x=\dd\mu/c$ and
$\tfrac12\iint_C|g(x)-g(y)|^2\dd\mu(x)\dd\mu(y)=\mu(C)V_C$.
The disjoint positive energy $E_{TI}$ of edges between $T$
and $I$, each counted once, satisfies
\[
 E_{TI}(g)\ge\frac{a_R}{2}\int_T|g-m_C|^2\dd\mu-N_RV_C.
\]
Indeed $|A-B|^2\ge\tfrac12|A|^2-|B|^2$, with
$A=g(x)-m_C$, $B=g(y)-m_C$; the two degree bounds apply.
Define
\[
 \theta_R=\min\{1,b_R/[2(N_R+1)]\},\qquad
 c_*:=\min\{b_R/2,\theta_Ra_R/2\}>0.
\]
Positivity allows retaining only the fraction $\theta_R$ of
$E_{TI}$, so
\[
 \mathcal D_\Phi(g)\ge D_{CC}(g)+\theta_RE_{TI}(g)
               \ge c_*\|g-m_C\|_{H_\mu}^2.
\]
The entire annulus $R-1<|x|\le R$ belongs to $C$.
For globally mean-zero $g$ the last squared norm is exactly
$\|g\|_{H_\mu}^2+|m_C|^2/2$, proving (JPG) on the core.
Form approximation with the mean correction used above proves
it on the entire stated domain. This is an explicit positive
kernel bound, without a numerical or optimal-value claim.
Increasing $R$ makes $b_R$ tend to zero, so this proof does
not bootstrap its constant to one.
\end{proof}

\begin{lemma}[A compact smooth operator core]
\label{lem:theta-operator-core}
Compact smooth functions form a graph-norm core for $L_\mu$
and for the operators of the closed forms $Q_\mu$ and
$q_A(\Phi\,\cdot)$. In particular
\[
 \|L_\mu g\|_{H_\mu}^2\ge\mathcal D_\Phi(g)
                       \qquad(g\in C_c^\infty(\R))          \tag{IGP}
\]
would imply (OGP). Unconditionally one has the same inequality
with $c_*\mathcal D_\Phi(g)$ on its right side, for every
$g\in D(L_\mu)$. The premise (IGP) is not proved here.
\end{lemma}
\begin{proof}
Under $Ug=\sqrt{c\Phi}\,g$ set $p=\sqrt{\Phi/c}$.
The ground identity is the exact form representation
\[
 UL_\mu U^{-1}=I+M_p a(D)M_p-\widetilde T_{\rm pr},
 \qquad \mathcal Q=\{u\in L^2:pu\in H_{\log}\}.
\]
The prime operator is bounded by
Lemma~\ref{lem:theta-prime-schur}; the remaining pole perturbation
is bounded and has rank at most two.
Closedness follows by applying the closed multiplier
$\sqrt{a(D)+C_a}$ to $pu$ and adding the bounded terms.
The form core is genuinely compact smooth: scaled smooth
cutoffs $\chi_R(x)=\chi(x/R)$ are uniformly bounded on
$H_{\log}$ by weighted Fourier--Young, then converge strongly
there by Schwartz density. On fixed support mollify $pu$ and
divide by the strictly positive smooth $p$.

The gamma commutator on compact smooth tests has kernel
\[
 [a(D),\chi_R]h(x)=\int_0^\infty k(r)
 \{[\chi_R(x)-\chi_R(x-r)]h(x-r)
    +[\chi_R(x)-\chi_R(x+r)]h(x+r)\}\dd r.
\]
Hence $\|[a(D),\chi_R]\|\le(C/R)\int_0^\infty rk(r)\dd r$.
Multiplication by bounded $p$ gives the same $O(1/R)$ bound
for the weighted gamma commutator. Each other commutator is
that of a bounded operator with $\chi_R\to I$ strongly,
and tends strongly to zero. Testing the distributional identity
$L(\chi_Ru)=\chi_RLu+[L,\chi_R]u$ on the proved form core
extends it to $u\in D(L)$. Its right side is $L^2$, so
$\chi_Ru\in D(L)$ and these cutoffs converge in graph norm.

For compactly supported $u$ set $h=pu$. The operator equation
gives $p\,a(D)h\in L^2$. On a fixed enlarged compact $p$
is bounded below. Outside its distance-one enlargement,
$a(D)h$ is the ordinary off-support convolution
$-\int k(|x-y|)h(y)\dd y$, whose truncated kernel is $L^1$.
It follows that $a(D)h\in L^2$ globally.
Mollification $h_\epsilon$ converges in its graph norm by
Fourier dominated convergence with weight $1+|a(t)|^2$.
Then $u_\epsilon=h_\epsilon/p$ are compact smooth and converge
in the graph norm of the weighted gamma operator; every other
term is bounded. This proves the operator-core assertion for
all three stated operators.

Thus (IGP), if proved on this core, extends to $D(L_\mu)$.
A nonzero spectral projection in any closed subinterval of
$(0,1)$ would make $\|L_\mu g\|^2-\mathcal D_\Phi(g)<0$.
The positive archimedean conductance shows that
$\ker L_\mu$ consists exactly of constants: zero energy makes
$g(x)=g(y)$ for almost every pair. Excluding $(0,1)$ would
therefore give (OGP). For the unconditional smaller constant,
subtract the global mean from $g\in D(L_\mu)$ and use (JPG):
\[
 \mathcal D_\Phi(g)=\langle L_\mu g,g_0\rangle
       \le\|L_\mu g\|\sqrt{\mathcal D_\Phi(g)/c_*}.
\]
Squaring proves the asserted integrated bound, also when the
energy is zero. No spectral interval below one is excluded by
the smaller constant alone.
\end{proof}

\begin{proposition}[Exact sharpness in both parity sectors]
\label{prop:theta-saturators}
For $j\ge1$ put
\[
 h_j^+=\Phi^{(2j)}-4^{-j}\Phi,\qquad
 h_j^-=\Phi^{(2j+1)}-4^{-j}\Phi',\qquad
 g_j^\pm=h_j^\pm/\Phi .
\]
These functions have finite $L^2(\nu)$ norms and finite energy
(DPE), and satisfy
$\int g_j^\pm\dd\nu=\int\kappa g_j^\pm\dd\nu=0$.
For every finite linear combination $g$ of these functions,
\[
             2\mathcal D_\Phi(g)=\|g\|_{L^2(\nu)}^2.       \tag{SAT}
\]
They are linearly independent within each parity sector.
Their weighted norms and energies are limits of compact smooth
tests with both moments tending to zero.
Consequently the coefficient one in (GP) cannot be replaced
by a uniformly larger coefficient on either moment-free sector.
\end{proposition}
\begin{proof}
The Fourier functions are
\[
 \widehat h_j^+(z)=\{(-1)^jz^{2j}-4^{-j}\}\Xi(z),\qquad
 \widehat h_j^-(z)=\{(iz)^{2j+1}-4^{-j}iz\}\Xi(z).
\]
They vanish at every Xi zero and at both pole points.
The latter values give both asserted moments.
The theta derivative series gives, for each fixed $r$,
$|\Phi^{(r)}(x)|/\Phi(x)\le C_r(1+e^{2|x|})^r$.
For large $|x|$ this follows by first-term dominance and
absorption of the remaining Gaussian series; on a compact
interval use the positive minimum of $\Phi$.
Thus all weighted norms converge.
The archimedean energy is finite near the diagonal by
differentiation and on its tail by this exponential-ratio bound.
For a large prime translation $r$, the product of the two
theta densities, multiplied by any of these fixed exponential
ratios, is bounded after integration by $C e^{-c e^{|r|}}$.
This sums every $\Lambda(n)/\sqrt n$ term.
The same majorants justify weighted-norm and energy convergence
under even cutoffs, so (TG) extends to the finite combinations.

Their compact cutoffs have zero-side evaluation norm tending
to zero by (TE), and their archimedean, pole, and prime forms
converge by the preceding bounds. Hence their limiting full
form is zero. In (TG) both moments are zero, which proves (SAT).
A finite linear dependence would give a nonzero polynomial
of distinct even or odd degrees times $\Xi$ identically zero
after Fourier transformation, an impossibility.
If exact compact moment conditions are desired, correct the
two small cutoff moments with a fixed even and odd compact
bump pair. Those correction coefficients tend to zero,
and their finite energy preserves the limit.
\end{proof}

\begin{proposition}[Fixed-support transfer to the theta energy]
\label{prop:theta-log-transfer}
For every fixed $L>0$ there is $C_L<\infty$ such that
every $e\in C_c^\infty((-L,L))$, with $h=e/\Phi$, satisfies
\[
 \mathcal D_\Phi(h)+\|h\|_{L^2(\nu)}^2
                         \le C_L\|e\|_{\log}^2.            \tag{TLT}
\]
Thus fixed-support log-Fourier convergence transfers to the
full theta energy and both moment terms. The constant is not
asserted to be uniform as $L$ grows.
\end{proposition}
\begin{proof}
The first identity in (TG) gives
\[
 \mathcal D_\Phi(h)
       =q_B(e)+\int_{-L}^L\frac{c(x)}{\Phi(x)}|e(x)|^2\dd x.
\]
The local potential is bounded on the fixed compact interval.
The full form $q_B$ is bounded in the log-Fourier norm there:
the gamma multiplier has logarithmic growth and the finite
sum of all translations $n\le e^{2L}$ is $L^2$ bounded.
Also
$\|h\|_{L^2(\nu)}^2=2\int(c/\Phi)|e|^2$.
The inequality $\|e\|_2^2\le(\log2)^{-1}\|e\|_{\log}^2$
proves (TLT). Cauchy--Schwarz for the positive energy, and for
$L^2(\nu)$, gives convergence of its polarized form and the
two bounded moment functionals.
This controls the squared energy error. Differences of energies
themselves need not have that same squared-error order.
Since $1/\Phi$ grows doubly exponentially on the tail, the
fixed-$L$ constant cannot be discarded in a cofinal argument.
\end{proof}

\section{Uniform theta-tail bounds and transfer of the full negative index}
The following estimates keep the full gamma form, every prime power,
and both pole moments. They apply to the complete density $\Phi$.
The norm throughout this section is $L^2(\nu)$, where
$\dd\nu=2c\Phi\dd x$; it is not the original physical $L^2$ norm.
Choose $C_a\ge0$ with $a(t)\ge-C_a$ for all real $t$, and put
\[
 p(x)=\sqrt{\Phi(x)/c(x)},\quad p_*:=\sup_xp(x)<\infty,\qquad
 P_R=\mathbf1_{|x|\ge R},\quad Q_R=1-P_R\quad(R\ge1).
\]
Define on $L^2(\nu)$ the all-prime operator
\[
 (\mathsf T_{\rm pr}g)(x)=
 \sum_{n\ge2}\frac{\Lambda(n)}{\sqrt n\,c(x)}
 \{\Phi(x-\log n)g(x-\log n)
                  +\Phi(x+\log n)g(x+\log n)\}.             \tag{TPS}
\]
It has a positive kernel; this does not assert positive
semidefiniteness of its quadratic form.

\begin{lemma}[Uniform all-prime Schur bounds]
\label{lem:theta-prime-schur}
Set
\[
 V_0=\int_\R e^{v/2}|\Phi'(v)-\Phi(v)/2|\dd v,\qquad
 C_{\rm pr}=4C_\psi V_0,\qquad
 \delta_R=20C_\psi e^{R/2}\Phi(R).
\]
Then $V_0<\infty$, the operator in (TPS) is bounded selfadjoint,
and
\[
 \|\mathsf T_{\rm pr}\|\le C_{\rm pr},\quad
 \|P_R\mathsf T_{\rm pr}P_R\|\le\delta_R,\quad
 \|\mathsf T_{\rm pr}P_R\|\le\sqrt{C_{\rm pr}\delta_R}.       \tag{TSS}
\]
These bounds are uniform over all test vectors and support sizes.
\end{lemma}
\begin{proof}
For a nonnegative bounded-variation function $W$ on $(0,\infty)$,
with zero endpoint limits and $V(W)=\int u|\dd W|(u)<\infty$,
\[
 W(u)\le\int_{[u,\infty)}|\dd W|(s),\qquad
 \frac1t\sum_{n\ge2}\Lambda(n)W(n/t)
       \le C_\psi V(W)\quad(t>0).                          \tag{BVC}
\]
Indeed Tonelli bounds the sum by
$\int t^{-1}\psi(ts)|\dd W|(s)$, and the elementary
$\psi(ts)\le C_\psi ts$ estimate applies. At a jump, any value
between the two one-sided limits, including the larger one,
is permitted: the displayed closed tail includes its jump atom.
Thus no coincident prime atom is lost.

For $W(u)=u^{-1/2}\Phi(\log u)$, $V(W)=V_0$.
Taking $t=e^x$ in (BVC) gives
\[
 \frac1{c(x)}\sum_n\frac{\Lambda(n)}{\sqrt n}\Phi(x-\log n)
 \le\frac{e^{x/2}}{c(x)}C_\psi V_0\le2C_\psi V_0.
\]
Evenness gives the same bound in the other direction.
The kernel's detailed-balance mass for each $n$ is
$2\Lambda(n)n^{-1/2}\Phi(x)\Phi(x-\log n)\dd x$.
The symmetric Schur test therefore gives boundedness and
selfadjointness with norm at most $C_{\rm pr}$.
Countable translations preserve null sets, and $\nu$ is equivalent
to Lebesgue measure, so the construction is well defined on
$L^2(\nu)$ equivalence classes.

For $x\ge1$, write $T=\pi n^2e^{2x}$ in each positive theta term.
Its logarithmic derivative is
$9/2+6/(2T-3)-2T\le-\pi e^{2x}$, since $T\ge\pi e^2$.
Termwise differentiation is locally dominated. Hence
\[
 \Phi'(x)\le-\pi e^{2x}\Phi(x),\qquad
 J_R:=\int_R^\infty e^{x/2}\Phi(x)\dd x
 \le\frac{e^{R/2}\Phi(R)}{\pi e^{2R}-1/2}.                 \tag{TJI}
\]
For $W_R(u)=W(u)\mathbf1_{|\log u|\ge R}$, the positive
continuous variation is $e^{R/2}\Phi(R)+J_R$.
The negative continuous variation is bounded by
$\int_R^\infty e^{-x/2}[-\Phi'(x)+\Phi(x)/2]\dd x
=e^{-R/2}\Phi(R)$. The two jump masses add
$(e^{R/2}+e^{-R/2})\Phi(R)$. Consequently
\[
 V(W_R)\le2(e^{R/2}+e^{-R/2})\Phi(R)+J_R
                                  \le5e^{R/2}\Phi(R).
\]
Applying (BVC) in both directions shows that, uniformly in every
starting point, the row mass into $|y|\ge R$ is at most
$4C_\psi V(W_R)\le\delta_R$. Detailed balance gives
the two Schur bounds in (TSS): the row constants are
$(\delta_R,\delta_R)$ and $(\delta_R,C_{\rm pr})$, respectively.
This includes connections between opposite tails.
\end{proof}

For compact smooth $g$, the exact physical prime form and (TG) give
\[
 2\mathcal D_\Phi(g)=\|g\|_\nu^2
        +2q_{\rm arch}(\Phi g)
                   -\langle g,\mathsf T_{\rm pr}g\rangle_\nu.\tag{TGE}
\]
\[
\begin{split}
 \mathfrak a(g):=2q_A(\Phi g)
    &=2q_{\rm arch}(\Phi g)
        -\langle g,\mathsf T_{\rm pr}g\rangle_\nu
        +\left|\int g\dd\nu\right|^2
        -\left|\int\kappa g\dd\nu\right|^2.
\end{split}                                                \tag{TFA}
\]
The factor two in these identities follows from the displayed
detailed-balance mass; it is not suppressed in the estimates below.

\begin{theorem}[Full theta-tail coercivity and touching-core control]
\label{thm:theta-tail-cross}
Write $\nu_R=\nu(|x|\ge R)$ and
\[
 \gamma_R=\sqrt2(1+\pi/2)p_*p(R),\quad
 \beta_R=\gamma_R+\sqrt{C_{\rm pr}\delta_R}
                     +\sqrt{(1-\nu_R)\nu_R},\quad
 e_R=\delta_R+C_ap(R)^2+\nu_R.
\]
For every compact smooth $g$ supported outside $[-R,R]$,
\[
 2\mathcal D_\Phi(g)\ge
        [1-\delta_R-C_ap(R)^2]\|g\|_\nu^2,\qquad
 \mathfrak a(g)\ge-e_R\|g\|_\nu^2.                       \tag{TTC}
\]
For core and tail form vectors $g_0=Q_Rg_0$, $g_1=P_Rg_1$,
\[
 |\mathfrak a(g_0,g_1)|\le\beta_R\|g_0\|_\nu\|g_1\|_\nu.
                                                               \tag{TCX}
\]
There is no restriction on modulation height, derivatives, mesh,
or packet dimension. The constants $e_R$ and $\beta_R$ tend to
zero doubly exponentially as $R\to\infty$.
\end{theorem}
\begin{proof}
Since $a\ge-C_a$, Parseval gives
\[
 2q_{\rm arch}(\Phi g)\ge
      -C_a\int(\Phi/c)|g|^2\dd\nu
      \ge-C_ap(R)^2\|g\|_\nu^2
\]
on the tail. Here $\Phi/c$ decreases for $x\ge1$ by (TJI).
Use (TSS) in (TGE) for the first bound in (TTC).
The negative pole moment in (TFA) is bounded by
$\nu_R\|g\|_\nu^2$, whereas the mean term is nonnegative.
This proves the second bound. Also $\nu_R\le4J_R$.

The gamma cross-block alone is
\[
 (\Gamma g)(x)=c(x)^{-1}\int k(|x-y|)\Phi(y)g(y)\dd y,
                  \qquad x\in\operatorname{core},\ y\in\operatorname{tail}.
\]
Only this cross-block is asserted to be bounded: the unseparated
diagonal integral of the full $\Gamma$ diverges.
Under the unitary $g\mapsto\sqrt{2c\Phi}\,g$, its kernel becomes
$p(x)p(y)k(|x-y|)$. For $-R\le x\le R$, $y\ge R$, put
$u=R-x$, $v=y-R$. The kernel is the Hankel kernel $k(u+v)$,
with $0\le u\le2R$ and $v\ge0$. The elementary bound
\[
 k(r)\le e^{-r/2}+\frac1{2r}\quad(r>0)
\]
follows from $e^{2r}-1\ge2r$.
The exponential Hankel kernel has rank one and norm
$\int_0^\infty e^{-u}\dd u=1$.
For the Carleman kernel, the Schur weight $u^{-1/2}$ gives
\[
 \int_0^\infty\frac{v^{-1/2}}{u+v}\dd v=\pi u^{-1/2}.
\]
Its norm is at most $\pi$, so the Hankel norm is at most
$1+\pi/2$. Restriction to $u\le2R$ cannot increase it.
Multiplication by the two weights has norms at most $p_*$ and
$p(R)$. The negative tail is identical. Combining the two
orthogonal tail input spaces gives
$\|Q_R\Gamma P_R\|\le\gamma_R$.

For disjoint supports the exact gamma cross-form is
$2q_{\rm arch}(\Phi g_0,\Phi g_1)
=-\langle g_0,\Gamma g_1\rangle_\nu$;
the scalar $a(0)$ term has no overlap. The all-prime cross-norm
is at most $\sqrt{C_{\rm pr}\delta_R}$ by (TSS).
The pole cross-block is
\[
 Q_R\{|1\rangle\langle1|-|\kappa\rangle\langle\kappa|\}P_R.
\]
Evenness makes $1$ and $\kappa$ orthogonal on each symmetric
piece separately. Its two possible singular values are
$\sqrt{(1-\nu_R)\nu_R}$ and
$\sqrt{\int_{\rm core}\kappa^2\dd\nu
               \int_{\rm tail}\kappa^2\dd\nu}$.
The second is no larger than the first. This proves (TCX),
with both pole terms present.

Cutting a compact smooth test by $Q_R$ and $P_R$ creates at most
two jumps. The resulting physical Fourier transforms are
$O(1/|t|)$ by piecewise integration by parts, hence belong to
the log-Fourier form domain. Mollification converges in that norm
by Fourier dominated convergence, and in $L^2(\nu)$ on the fixed
support. A tail mollification uses $R-\eta$ first; continuity of
the displayed constants lets $\eta\downarrow0$.
For separate general core and tail form vectors, truncate on a
fixed compact support, move the tail pieces slightly outward and
the core slightly inward, and mollify. Translation and dilation
are strongly continuous in the log-Fourier norm; the weighted and
physical norms are equivalent on that fixed support. The bounded
cross-kernel controls the limiting sesquilinear form.
Finally the theta-series bound
$\Phi(R)\le C e^{9R/2}e^{-\pi e^{2R}}$ and (TJI) give the
claimed decay of all constants.
\end{proof}

\begin{theorem}[A quantitative full-form transfer at fixed negative depth]
\label{thm:theta-core-index}
For $\epsilon>e_R$ set
\[
 \tau_R(\epsilon)=\epsilon-\frac{\beta_R^2}{\epsilon-e_R}.
\]
Let $\mathcal H_R$ consist of $g\in L^2(\nu)$ supported in
$[-R,R]$ such that $\Phi g$ belongs to the physical log-Fourier
form domain. Let $\mathcal I_R(\tau)$ be the supremum of dimensions
of subspaces of $\mathcal H_R$ on which
$\mathfrak a(g)+\tau\|g\|_\nu^2<0$ for every nonzero $g$.
On the left below take the index over all compact smooth tests,
with arbitrarily large supports. Then
\[
 n_-(\mathfrak a+\epsilon\|\cdot\|_\nu^2)
          \le\mathcal I_R(\tau_R(\epsilon))<\infty.        \tag{TCI}
\]
More explicitly, put
\[
 m_R=\min_{|x|\le R}\Phi(x)/c(x)>0,\qquad
 B_\tau=C_{\rm pr}+\|\kappa\|_\nu^2+\max\{-\tau,0\}.
\]
For any $H>0$ with $a(H)\ge2B_\tau/m_R+C_a$,
\[
                \mathcal I_R(\tau)\le\lceil4RH/\pi\rceil.
                                                               \tag{TCB}
\]
Consequently the dimension of any subspace satisfying
$\mathfrak a(g)\le-\epsilon\|g\|_\nu^2$ is uniformly finite
over all support parameters, for every fixed $\epsilon>0$.
\end{theorem}
\begin{proof}
For $g=g_0+g_1$, (TTC)--(TCX) and completion of the square give
\[
\begin{split}
 \mathfrak a(g)+\epsilon\|g\|_\nu^2
 \ge{}&\mathfrak a(g_0)+\tau_R(\epsilon)\|g_0\|_\nu^2\\
 &+(\epsilon-e_R)
       \left(\|g_1\|_\nu-
            \frac{\beta_R}{\epsilon-e_R}\|g_0\|_\nu\right)^2.
\end{split}                                                   \tag{TCS}
\]
On any strictly negative subspace of the left side, $Q_R$ is
injective: a pure tail has strictly positive lower bound.
Its image has the same dimension and is negative for the shifted
core form in (TCS). This proves the index inequality in (TCI).

The digamma series \cite[eq.~(5.7.6), away from its poles]{DLMFGamma}
gives, with $x_n=n+1/4$,
\[
 a(t)-a(0)=\sum_{n\ge0}
       \frac{(t/2)^2}{x_n\{x_n^2+(t/2)^2\}}.
\]
The locally uniformly convergent terms increase with $|t|$.
Thus $a(t)\ge a(H)$ for $|t|\ge H$; logarithmic growth guarantees
that an admissible $H$ exists.
For a negative core vector put $f=\Phi g$, $n=\|f\|_2^2$.
Then $\|g\|_\nu^2\le2n/m_R$ and (TFA) yields
\[
 \mathfrak a(g)+\tau\|g\|_\nu^2
                 \ge2q_{\rm arch}(f)-B_\tau\|g\|_\nu^2.
\]
Let $\Pi_H$ be the physical $L^2$ Fourier projection onto
$[-H,H]$. Parseval with factor $1/(2\pi)$ gives
\[
 q_{\rm arch}(f)\ge a(H)n-{a(H)+C_a}\|\Pi_Hf\|_2^2.
\]
Negativity therefore forces
\[
 \frac{\|\Pi_Hf\|_2^2}{n}
   >\frac{a(H)-B_\tau/m_R}{a(H)+C_a}\ge\frac12.
\]
The map $g\mapsto\Phi g$ is a linear bijection onto its core
image. Choose a physical orthonormal basis $f_1,\ldots,f_d$ of
that negative image, so this concentration bound applies to each.
The positive band-concentration operator on $[-R,R]$ has kernel
$\sin(H(x-y))/(\pi(x-y))$, diagonal $H/\pi$, and trace $2RH/\pi$.
Bessel gives
\[
 d/2<\sum_{j=1}^d\|\Pi_Hf_j\|_2^2\le2RH/\pi.
\]
This proves (TCB) and the finiteness in (TCI).
For the final non-strict depth assertion apply (TCI) with shift
$\epsilon/2$, choosing $R$ sufficiently large.
\end{proof}

For a fixed shifted depth one may choose $R$ with
$e_R\le\epsilon/2$ and $\beta_R^2\le\epsilon^2/4$;
then $\tau_R(\epsilon)\ge\epsilon/2$. Every strictly negative
direction below that depth therefore gives a genuine negative
direction on this central core. This transfer controls the full
form and has no exceptional set of modulation heights.
It does not prove core nonnegativity or a uniform total index bound.
The remaining condition (GP) is still open. In particular the small
negative remainder in (TTC) cannot be replaced by zero, and
$\|g\|_\nu^2=2\int c|\Phi g|^2/\Phi$ can be very large relative
to the physical norm. No subpower bound for the original
finite-interval negative eigenvalues is inferred from these estimates.

\section{A global weighted Weil operator and its exact essential spectrum}
This section constructs one actual operator on $L^2(\nu)$.
It does not take a limit of the physical finite-interval operators
and assumes no simple or uniformly isolated ground state.

\begin{theorem}[The full weighted operator and its kernel complement]
\label{thm:theta-global-spectrum}
The compact-test form $\mathfrak a(g)=2q_A(\Phi g)$ has a
semibounded closed form closure on $L^2(\nu)$, with domain
\[
 \mathcal Q_\Phi=\{g\in L^2(\nu):\Phi g\in H_{\log}(\R)\}.
\]
Let $\mathcal L_\Phi$ be its associated selfadjoint operator,
and $\mathcal E_\Phi$ the operator associated to the closure
of the positive form $2\mathcal D_\Phi$. Then
\[
 \sigma_{\rm ess}(\mathcal L_\Phi)=
 \sigma_{\rm ess}(\mathcal L_\Phi|_{\ker(\mathcal L_\Phi)^\perp})
 =[0,\infty),\qquad
 \sigma_{\rm ess}(\mathcal E_\Phi)=[1,\infty).             \tag{GES}
\]
The same assertions hold in each parity sector, removing the
entire sector kernel when taking its complement.
In particular no positive uniform gap is available even after
removing the complete operator nullspace. Negative discrete
eigenvalues are not excluded by these statements.
\end{theorem}
\begin{proof}
The unitary $Ug=\sqrt{2c\Phi}\,g$ satisfies
$\Phi g=pUg/\sqrt2$. If $v=Ug$, (TFA) becomes the exact form
\[
 \ell(v)=\langle pv,a(D)pv\rangle
       -\langle v,\widetilde{\mathsf T}_{\rm pr}v\rangle
       +\langle v,\widetilde Q_Pv\rangle,
 \quad Q_P=|1\rangle\langle1|-|\kappa\rangle\langle\kappa|,
 \quad \mathcal Q=\{v\in L^2:pv\in H_{\log}\},             \tag{GWF}
\]
where tildes denote conjugation by $U$.
The operator $S=\sqrt{a(D)+C_a}\,M_p$, with domain $\mathcal Q$,
is closed: convergence of $v_j$ and $Sv_j$ gives convergence of
$pv_j$, and the closed Fourier multiplier applies. The squared
norm of $S$ is a closed nonnegative form. The difference between
that form and (GWF) consists of the bounded forms
$-C_a\|pv\|_2^2$, $-\widetilde{\mathsf T}_{\rm pr}$ and the
rank-two pole form. Thus (GWF) is densely defined, closed and
semibounded. The graph domain is equivalent to the stated
log-Fourier domain after the underlying $L^2$ norm is included.

We verify that compact smooth vectors are a form core.
For fixed smooth compact cutoffs $\chi_R(x)=\chi(x/R)$,
the weighted Fourier $L^1$ norms with weight
$w(t)=\sqrt{\log(2+|t|)}$ are uniformly bounded for $R\ge1$.
The inequality $w(u+v)\le w(u)+w(v)$ and weighted Young
give uniform boundedness of multiplication by $\chi_R$ on
$H_{\log}$. Schwartz density, followed by this uniform bound,
gives $\chi_Rpv\to pv$ in $H_{\log}$, while $\chi_Rv\to v$
in $L^2$. On the resulting fixed compact support mollify
$\chi_Rpv$. Fourier dominated convergence gives log-norm
convergence. Division by the positive smooth $p$ on that fixed,
slightly enlarged support gives compact smooth $v$ approximants,
also converging in $L^2$. A diagonal choice proves the form-core
claim. Consequently (GWF) is precisely the compact-test form
closure, with no additional condition at infinity.

The closed-form representation theorem applies
\cite[Theorem 2.13, pp.~71--72]{TeschlQM2009}.
On $C_c^\infty(\R)$ the resulting physical conjugate acts as
\[
 \widetilde{\mathcal L}_\Phi v
   =p\,a(D)(pv)-\widetilde{\mathsf T}_{\rm pr}v
                                    +\widetilde Q_Pv.      \tag{GWA}
\]
Indeed $a(D)(pv)\in L^2$ by Fourier decay and $p$ is bounded.
Plancherel against every form-domain test gives this actual
operator representative, including operator-domain membership.
The closure of $2\mathcal D_\Phi$ has the same form core and
domain, since (TGE)--(TFA) give its bounded difference
$I-Q_P$ from $\mathfrak a$.

Here is an explicit operator Weyl sequence for every $\lambda>0$.
The positive theta-series on $x\ge1$ gives, for fixed $m$,
\[
 |\Phi^{(m)}(x)|/\Phi(x)\le C_me^{2mx},\qquad
 |p^{(m)}(x)|\le C_mp(x)e^{2mx}\quad(m=1,2).              \tag{GWP}
\]
To see the first bound write each term as
$2e^{x/2}T_n(2T_n-3)e^{-T_n}$, $T_n=\pi n^2e^{2x}$.
Its $m$th derivative is bounded by
$C_me^{x/2}T_n^{m+2}e^{-T_n}$. The first term lower-bounds
the denominator, and
$\sum_n n^{2m+4}e^{-\pi e^2(n^2-1)}<\infty$ controls all
others. Differentiation of $p=\sqrt{\Phi/c}$ proves the
second bound. These same majorants justify differentiation.

The actual digamma asymptotic and global logarithmic bound imply,
for $t\ge4$ and $b\in C_c^\infty$,
\[
 \|(a(D)-a(t))e^{itx}b\|_2
 \le Ct^{-1}(\|b\|_2+\|b'\|_2)
                  +C(\log t)t^{-2}\|b''\|_2.             \tag{GSE}
\]
For Fourier displacements $|u|\le t/2$, the multiplier
difference is at most $C(1+|u|)/t$, by the logarithm and both
asymptotic errors. On $|u|>t/2$, it is at most
$C\log(2+|u|)$; the supremum of
$\log(2+|u|)/u^2$ there is $O(\log t/t^2)$.
Parseval proves (GSE), including the Fourier tail where $t+u$
may approach zero. This controls the full operator remainder.

Fix smooth $\chi$ supported in $[-1,1]$ with $\|\chi\|_2=1$.
For large $R$ choose
\[
 \rho_R=e^{-3R},\quad p_R=p(R),\quad
 a(t_R)=\lambda/p_R^2,\qquad
 v_R(x)=\rho_R^{-1/2}\chi((x-R)/\rho_R)e^{it_Rx}.          \tag{GWV}
\]
Strict monotonicity of the digamma series and logarithmic growth
give a unique $t_R>0$, with
$t_R=2\pi\exp(\lambda/p_R^2+o(1))$.
On this support (GWP) gives $p(x)/p_R=1+O(e^{-R})$ and,
for $b_R=p\rho_R^{-1/2}\chi((x-R)/\rho_R)$,
$\|b_R^{(m)}\|_2\le C_mp_R\rho_R^{-m}$, $m=0,1,2$.
Thus (GSE) proves
\[
 \|p\,a(D)(pv_R)-\lambda v_R\|_2
 \le C\lambda e^{-R}
  +Cp_*p_R\left\{\frac{1+\rho_R^{-1}}{t_R}
             +\frac{\log t_R}{t_R^2\rho_R^2}\right\}\to0. \tag{GWR}
\]
The prime operator remainder is at most
$\sqrt{C_{\rm pr}\delta_{R-1}}$ by (TSS), and the pole
remainder at most $2\sqrt{\nu(|x|\ge R-1)}$.
Both tend to zero with all prime powers and pole terms present.
Taking $R_j=2j$ gives unit vectors with disjoint escaping
supports, hence an orthonormal weakly null sequence in the
operator domain. The Weyl criterion
\cite[Lemma 6.17, pp.~145--146]{TeschlQM2009} puts every
$\lambda>0$ in the essential spectrum.

There is no negative essential spectrum. Theorem
\ref{thm:theta-core-index} gives finite index below every fixed
negative weighted threshold on compact tests. If a spectral
projection below $-\epsilon$ had infinite rank, any fixed finite
dimension in its range could be approximated in form norm by
the proved compact core. The strictly negative matrix inequality
persists with margin $\epsilon/2$, and its Gram remains positive
definite. Arbitrarily large dimensions would contradict that
theorem. The essential-spectrum characterization
\cite[eq.~(6.28), p.~145]{TeschlQM2009} therefore excludes
negative essential points. Since the essential spectrum is
closed and contains every positive $\lambda$, it is $[0,\infty)$.

Let $P_0$ be the full kernel projection of
$\widetilde{\mathcal L}_\Phi$. For the preceding sequence,
\[
 \lambda\|P_0v_R\|_2
 \le\|(\widetilde{\mathcal L}_\Phi-\lambda)v_R\|_2\to0.
\]
Selfadjointness makes the kernel reducing and preserves the
operator domain under $P_0$. The normalized vectors
$(I-P_0)v_R$ remain weakly null and retain the Weyl remainder.
Hence every positive $\lambda$, and then zero by closedness,
remains essential on the entire kernel complement. Restriction
cannot add a negative infinite-rank spectral interval.
Finally $\mathcal E_\Phi=I+\mathcal L_\Phi-Q_P$;
finite-rank invariance \cite[Lemma 6.18, p.~146]{TeschlQM2009}
gives its asserted essential spectrum.

All operators preserve parity. The vectors
$v_R^\pm(x)=(v_R(x)\pm v_R(-x))/\sqrt2$ have unit norm
and at most $\sqrt2$ times the original remainder, since their
supports are disjoint and the operators commute with reflection.
They escape in each sector. The same negative-spectrum exclusion
and complete kernel projection then prove the sector statements.
\end{proof}

\begin{corollary}[Exact arithmetic index of the global weighted operator]
\label{cor:theta-global-index}
With $K_+$ as in Corollary~\ref{cor:parity-index},
\[
 \dim\mathbf1_{(-\infty,0)}(\mathcal L_\Phi^+)
 =\dim\mathbf1_{(-\infty,0)}(\mathcal L_\Phi^-)=K_+,\qquad
 \dim\mathbf1_{(-\infty,0)}(\mathcal L_\Phi)=2K_+.         \tag{GNI}
\]
Infinite $K_+$ gives infinite dimensions. The left sides count
operator eigenmultiplicities; $K_+$ counts distinct off-line
reflection orbits and assumes no simplicity of zeta zeros.
All negative eigenvalues are discrete, finite in number below
each fixed negative threshold, and can accumulate only at zero.
\end{corollary}
\begin{proof}
The bijection $f=\Phi g$ on compact smooth tests preserves the
full form up to its positive factor two. Their negative form
index is therefore the supremum in (PI). The proved form core
approximates every finite strictly negative form subspace by
compact tests without changing its dimension or sign.
Conversely any such negative subspace projects injectively into
the negative spectral subspace
\cite[Theorem 4.12(i), p.~119]{TeschlQM2009}.
Thus this index is exactly the dimension of the negative spectral
projection. Parity symmetrization preserves the form-core
approximation, so (PI) proves all three equalities, also when
$K_+$ is infinite. Discreteness and the accumulation statement
follow from (GES) and the fixed-threshold bound.
\end{proof}

\begin{proposition}[A relatively compact full arithmetic perturbation]
\label{prop:theta-relative-compact}
In the physical representation define by its closed form
\[
 H_\Phi=M_p\{a(D)+C_a\}M_p\ge0,\qquad
 B_\Phi=-C_aM_{p^2}-\widetilde{\mathsf T}_{\rm pr}
                                      +\widetilde Q_P.
\]
For every nonreal $z$, $B_\Phi(H_\Phi-z)^{-1}$ is compact and
\[
 (\widetilde{\mathcal L}_\Phi-z)^{-1}
                       -(H_\Phi-z)^{-1}\quad\hbox{is compact}. \tag{GRC}
\]
In particular the relatively compact perturbation retains the
complete all-prime operator, rather than a finite prime substitute.
\end{proposition}
\begin{proof}
For a bounded family $v_j$ in the $H_\Phi$ form norm,
$f_j=pv_j$ is bounded in $H_{\log}$ and
$\|\mathbf1_{|x|\ge R}f_j\|_2\le p(R)\|v_j\|_2\to0$
uniformly. Fourier truncation followed by spatial truncation
gives physical $L^2$ precompactness: the Fourier squared error
is $O(1/\log H)$, and the remaining finite-bandwidth kernel on
a compact spatial interval is Hilbert--Schmidt.
On a fixed core $1/p$ is bounded, so $Q_Rv_j$ is precompact too.
The spectral multiplier $\sqrt{1+u}/(u-z)$ for $u\ge0$ shows
that $(H_\Phi-z)^{-1}$ maps the unit ball into a bounded form
family. Consequently $M_p(H_\Phi-z)^{-1}$ and every
$Q_R(H_\Phi-z)^{-1}$ are compact. Now
\[
 \widetilde{\mathsf T}_{\rm pr}(H_\Phi-z)^{-1}
 =\widetilde{\mathsf T}_{\rm pr}Q_R(H_\Phi-z)^{-1}
  +\widetilde{\mathsf T}_{\rm pr}P_R(H_\Phi-z)^{-1}.
\]
The first term is compact; the second tends to zero in norm
by (TSS). Their norm limit is compact. The same conclusion
holds for $M_{p^2}(H_\Phi-z)^{-1}$, while the pole term has
finite rank. This proves relative compactness of $B_\Phi$.
Its bounded form difference also proves, by the representation
theorem, the actual domain equality
$D(\widetilde{\mathcal L}_\Phi)=D(H_\Phi)$ and operator equality
$\widetilde{\mathcal L}_\Phi=H_\Phi+B_\Phi$.
The resolvent identity
\[
 (\widetilde{\mathcal L}_\Phi-z)^{-1}-(H_\Phi-z)^{-1}
 =-(\widetilde{\mathcal L}_\Phi-z)^{-1}B_\Phi(H_\Phi-z)^{-1}
\]
proves (GRC). The essential-spectrum invariance theorem applies
\cite[Theorem 6.19, pp.~146--147]{TeschlQM2009}.
\end{proof}

These are statements about a specific closed weighted operator,
with an exact arithmetic negative index. The remaining RH step
is to exclude its discrete negative spectrum. Removing its
whole kernel supplies no positive spectral gap, and relatively
compact perturbations can create negative eigenvalues. Neither
property can replace the still unproved sign in (GP).

\subsection{Gaussian grounds, quantitative pole transfer, and exact defect flow}
\label{sec:heat-ground}
We prove the full Gaussian ground transformation, a quantitative
ordinary-gap bridge, and an exact scalar defect-flow theorem.
The scalar flow has a fixed physical variational limit.
A coherent family of detectors also detects a zero exactly at the
boundary of the corresponding exponential weight. These results
retain every prime power and both pole terms; they do not establish
the remaining sharp global gap.

For $t>0$ define
\[
 \begin{gathered}
 G_t(x)=(4\pi t)^{-1/2}e^{-x^2/(4t)},\qquad
 p_t=e^{-t/4}G_t*\Phi,\\
 \dd\mu_t=c p_t\dd x,\qquad H_t=L^2(\mu_t),\qquad
 a_t=\frac{e^{-t/4}}{\sqrt{4\pi t}},\qquad M_0=\Xi(0)>0.
 \end{gathered}
                                                               \tag{HGF1}
\]
Here $c=\cosh(x/2)$ and $\kappa=\tanh(x/2)$ have their previous
meaning. The function $p_t$ is positive, smooth, and even, and
\[
 \widehat p_t(z)=e^{-t(z^2+1/4)}\Xi(z),\qquad
 \widehat p_t(\pm i/2)=\frac12,\qquad \mu_t(\R)=\frac12.
                                                               \tag{HGF2}
\]
The entire multiplier is nonvanishing, so every Xi zero retains
its exact multiplicity. For each fixed $t$, every fixed derivative
of $p_t$ has a Gaussian upper tail. Indeed $\Phi(y)\le
C_\eta e^{-y^2/(4\eta)}$ for every $\eta>0$; convolution gives
$p_t(x)\le C_{t,\eta}e^{-x^2/[4(t+\eta)]}$, and differentiated
Gaussian kernels give the analogous bounds. A fixed positive
mass slice of $\Phi$ gives
$p_t(x)\ge c_t e^{-(|x|+M)^2/(4t)}$ for some $M,c_t>0$.
No estimate here is asserted uniformly as $t\to\infty$.

\begin{theorem}[The complete Gaussian ground identity]
\label{thm:heat-ground}
For every fixed $t>0$, the local smooth identity is
$\mathcal Bp_t=-c$. Define
\[
\begin{split}
 \mathcal D_t(g)={}&\frac12\iint k(|x-y|)p_t(x)p_t(y)
                              |g(x)-g(y)|^2\dd x\dd y\\
 &+\sum_{n\ge2}\frac{\Lambda(n)}{\sqrt n}
       \int p_t(x)p_t(x-\log n)
                          |g(x)-g(x-\log n)|^2\dd x .
\end{split}                                                   \tag{HGF3}
\]
This positive form is closable on $C_c^\infty(\R)$ in $H_t$.
Let $\mathcal V_t$ be its core closure and $L_t\ge0$ its
selfadjoint operator. For compact smooth $g$,
\[
\begin{split}
 q_B(p_tg)&=\mathcal D_t(g)-\|g\|_{H_t}^2,\\
 q_A(p_tg)&=\mathcal D_t(g)-\|g\|_{H_t}^2
       +2\left|\int g\dd\mu_t\right|^2
       -2\left|\int\kappa g\dd\mu_t\right|^2 .
\end{split}                                                   \tag{HGF4}
\]
The all-prime energy includes every outgoing edge, even when its
other endpoint lies outside the support of $g$.
\end{theorem}
\begin{proof}
The already proved identity (GDI) gives the positive prime degree
$S_\Phi=c+\mathcal A\Phi\le c+C$, where $\mathcal A=a(D)$ is
the full archimedean operator. Its second-difference expression
and $\int_0^\infty r^2k(r)\dd r<\infty$ bound
$\|\mathcal A\Phi\|_\infty$ using $\|\Phi\|_\infty$ and
$\|\Phi''\|_\infty$. Thus $G_t*S_\Phi$ is finite pointwise.
Tonelli commutes the positive all-prime sum with $G_t$.
For the archimedean second difference, use
$r^2G_t*|\Phi''|$ near zero and the integrable kernel at infinity.
Fubini therefore commutes that term too. Since
$G_t*c=e^{t/4}c$, convolution of (GDI) proves
$\mathcal Bp_t=-c$, without an undefined convolution of
arbitrary distributions.

On each fixed compact, the shifted prime terms and all fixed
derivatives converge absolutely: their tails are bounded by
$C_{t,K,j}(\log n)n^{-1/2}e^{-b_{t,K}(\log n)^2}$.
The small archimedean second difference cancels the singularity
of $k$. The pair algebra in Theorem~\ref{thm:theta-ground},
first with a distance cutoff and a finite prime cap, gives
\[
 q_B(p_tg)-\int\frac{\mathcal Bp_t}{p_t}|p_tg|^2\dd x
                                                    =\mathcal D_t(g).
\]
The displayed convergence and Tonelli remove the cutoffs.
The potential is exactly $-\int c p_t|g|^2$, and restoring
the actual pole form proves (HGF4).
The edge-$L^2$ subsequence argument used for the theta energy
proves closability: an $H_t$-convergent form-Cauchy sequence
identifies all its archimedean and countably many prime
edge-difference limits with the same almost everywhere limit.
This gives the asserted core closure. In particular no
fixed-support limit has been exchanged with $t\to\infty$.
\end{proof}

Put $Q_t=\mathcal D_t-\langle\cdot,\cdot\rangle_{H_t}$ and define
on $\mathcal V_t$ the closed full ground form
\[
 \mathcal W_t(g)=Q_t(g,g)
       +2\left|\int g\dd\mu_t\right|^2
       -2\left|\int\kappa g\dd\mu_t\right|^2 .
                                                               \tag{HGF5}
\]
On the compact core this equals $q_A(p_tg)$; the definition
does not presume a physical domain identity for every
noncompact vector.

\begin{proposition}[The score identity and uniform Fisher control]
\label{prop:heat-score}
The functions $1$ and $v_t=-2p_t'/p_t$ belong to $\mathcal V_t$.
Moreover $L_t1=0$ and
\[
\begin{gathered}
 \int v_t\dd\mu_t=0,\qquad
 \langle\kappa,v_t\rangle_{H_t}=\frac12,\qquad
 Q_t(v_t,g)=\langle\kappa,g\rangle_{H_t},\\
 Q_t(v_t,v_t)=\frac12,\qquad
 v_t\in D(L_t),\quad L_tv_t=v_t+\kappa,\qquad
 \|v_t\|_{H_t}^2\le\frac12+\frac1t .
\end{gathered}                                                \tag{HGP1}
\]
\end{proposition}
\begin{proof}
For fixed $t$, all ratios $p_t^{(j)}/p_t$ grow at most
polynomially. To verify this, use the conditional density
proportional to $\Phi(y)e^{xy/(2t)-y^2/(4t)}$.
Its $j$th absolute moment on $|y|\le4|x|+C$ is
$O_{t,j}((1+|x|)^j)$. The normalizing integral is at least
$c_te^{-C_t|x|}$ by a fixed mass slice, while the
double-exponential Phi tail outside that region bounds the
remaining normalized moments. Gaussian-Hermite
differentiation proves the ratio assertion.
The cutoff errors for $1,v_t$ are integrably dominated in
the small-distance energy by $r^2$ times polynomial factors
and Gaussian tails. For a prime shift $r$, integration of
the analogous majorant is
$C_t(1+|r|)^A e^{-b_tr^2}$, summable with every
$\Lambda(n)/\sqrt n$ at $r=\log n$.
Thus smooth cutoffs converge in the stated form norm.

For fixed $t,j$, two integrations by parts on the closed strip
$|\Im z|\le1/2$ give
\[
 \left|\widehat{(1-\chi_R)p_t^{(j)}}(z)\right|
 \le C_{t,j}e^{-b_{t,j}R^2}(1+|\Re z|^2)^{-1}.
                                                               \tag{HGP2}
\]
The all-height $S_4$ bound from (TE) sums these errors with
all multiplicities; the two pole errors converge in weighted
$L^1$. Now $\widehat{-2p_t'}(z)=-2iz e^{-t(z^2+1/4)}\Xi(z)$
vanishes at every zero, and its pole values are $(1/2,-1/2)$.
Apply (ZF), (HGF4), and the proved form exhaustion to obtain
the mixed identity in (HGP1). Integration by parts gives its
two moments. The right side is $H_t$-bounded, so the
representation theorem gives the stated operator action.

For the norm, mixture Cauchy--Schwarz gives
\[
 \frac{|p_t'(x)|^2}{p_t(x)}
 \le\frac{e^{-t/4}}{4t^2}
           \int G_t(x-y)\Phi(y)(x-y)^2\dd y.
\]
The exact Gaussian moment is
$\int c(y+z)G_t(z)z^2\dd z=(2t+t^2)e^{t/4}c(y)$.
Integrating and using $\int c\Phi=1/2$ yields
$\int c|p_t'|^2/p_t\le1/(4t)+1/8$; multiplication by four
proves the uniform bound.
\end{proof}

\begin{theorem}[Quantitative transfer of the complete pole correction]
\label{thm:heat-gap-transfer}
Let $H_t^0=\{g:\int g\dd\mu_t=0\}$ and
$S_tg=g-2\langle\kappa,g\rangle_{H_t}v_t$. On
$\mathcal V_t\cap H_t^0$,
\[
 \mathcal W_t(g)=Q_t(S_tg,S_tg),\qquad
 \|S_t\|_{H_t^0\to H_t^0}^2
       =4\|v_t\|_{H_t}^2\|\kappa\|_{H_t}^2\le1+2/t .
                                                               \tag{HGP3}
\]
If $\epsilon_t\ge0$ and
$\mathcal D_t(g)\ge(1-\epsilon_t)\|g\|_{H_t}^2$ for all
mean-zero form-domain vectors, then every compact smooth $g$
satisfies
\[
 q_A(p_tg)\ge-(1+2/t)\epsilon_t
       \left\|g-2\int g\dd\mu_t\right\|_{H_t}^2 .        \tag{HGP4}
\]
In particular the ordinary gap with exact constant one is
equivalent to full Weil positivity and RH. This is a proved
implication; its sharp-gap premise remains unproved.
\end{theorem}
\begin{proof}
The score identity gives
$Q_t(S_tg,S_tg)=Q_t(g,g)-2|\langle\kappa,g\rangle|^2$.
Both $\kappa$ and $v_t$ have mean zero, and
$\langle\kappa,v_t\rangle=1/2$, so $S_t$ is the projection
onto $\kappa^\perp$ along $v_t$. For
$A=\|v_t\|^2$, $B=\|\kappa\|^2$, write
$g=h+(z/B)\kappa$ with $h\perp\kappa$. Then
$S_tg=h+z(\kappa/B-2v_t)$, and
$\|\kappa/B-2v_t\|^2=4A-1/B$.
The row-map norm squared is $1+B(4A-1/B)=4AB$,
also on $H_t^0$, whose orthogonal complement to $\kappa$
is infinite dimensional. Since $B\le1/2$, (HGP1) gives
the bound in (HGP3).
Apply the ordinary gap to $S_tg$. For general $g$, subtract
its mean using the proved domain membership of $1$.
This proves (HGP4), including the negative odd pole correction.
If the gap is exact, every compact physical $f$ equals $p_tg$
for compact smooth $g$, and the full form is nonnegative;
(NI) excludes every off-line zero. Conversely RH makes
(ZF) a sum of positive real-node squares. On the mean-zero
core $q_B=q_A+2|\int\kappa g\dd\mu_t|^2\ge0$.
Correcting the mean of a core approximant with a fixed
compact bump proves density in $H_t^0\cap\mathcal V_t$,
and therefore the full converse.
\end{proof}

Define the ordinary gap
\[
 \beta_t=\inf_{0\ne g\in\mathcal V_t\cap H_t^0}
                       \frac{\mathcal D_t(g)}{\|g\|_{H_t}^2}.
                                                               \tag{HFL1}
\]
It satisfies $0\le\beta_t\le1$. For the upper bound the
nonzero vector $(p_t''-p_t/4)/p_t$ belongs to the proved
polynomial-ratio form domain, has both pole moments zero,
and has zero values at every Xi node. The mixed identity
and (HGP2) give $\mathcal D_t(g)=\|g\|^2$.

\begin{theorem}[The exact scalar heat-defect flow]
\label{thm:heat-defect-flow}
Set
\[
 C_-=\sup_{\substack{0\ne f\in C_c^\infty\\\int cf=0}}
                      \frac{[-q_B(f)]_+}{\int c|f|^2}
                         \in[0,\infty].
\]
Then $E_t=(1-\beta_t)/a_t$ is nondecreasing for $t>0$, and
\[
                   \lim_{t\to\infty}E_t=M_0C_- .         \tag{HFL2}
\]
The value infinity is permitted. A finite global bound
$q_B(f)\ge-C_0\int c|f|^2$ on this physical mean-zero core
is equivalent to $1-\beta_{t_j}\le K a_{t_j}$ on an
unbounded sequence, with constants $K=C_0M_0$ and
$C_0=K/M_0$ in the respective directions.
An estimate $1-\beta_{t_j}=o(a_{t_j})$ would prove RH;
no such uniform estimate is established here.
\end{theorem}
\begin{proof}
Put $r_t=p_t/a_t=\int\Phi(y)e^{-(x-y)^2/(4t)}\dd y$.
For each $x$, positivity gives $0<r_s(x)\le r_t(x)\le M_0$
when $s<t$, and $r_t(x)\uparrow M_0$.
The compact bijection $f=p_tg$ turns the mean condition
into the same physical constraint $\int cf=0$ for every $t$.
The mean-zero compact form core is dense, and $\beta_t\le1$,
so
\[
 E_t=\sup_{\substack{0\ne f\in C_c^\infty\\\int cf=0}}
             \frac{[-q_B(f)]_+}{\int c|f|^2/r_t}.        \tag{HFL3}
\]
Each ratio is nondecreasing. Its denominator converges on
the fixed compact support to $\int c|f|^2/M_0$.
The bound $r_t\le M_0$ gives $E_t\le M_0C_-$.
Taking the endpoint limit for each fixed $f$ and then its
supremum gives the reverse liminf, even if $C_-=\infty$.
This proves (HFL2) without a uniform exchange of vector,
support, or height limits. It also proves the finite-bound
equivalence. Directly, $p_t\le M_0a_t$ turns the physical
bound into $q_B(p_tg)\ge-C_0M_0a_t\|g\|_{H_t}^2$.

For completeness an individual off-line zero supplies,
by Lemma~\ref{lem:pole-free-interpolation} and (BI), one
fixed compact exact-pole-null $f$ with $q_B(f)=q_A(f)=-\eta<0$.
On its fixed support
$p_t=a_t[M_0+O_f(t^{-1})]$, because
$0\le1-e^{-u}\le u$ and $\Phi$ has a finite second moment.
Thus
\[
 1-\beta_t\ge\frac{\eta}{\|f/p_t\|_{H_t}^2}
   =a_t\frac{\eta M_0}{\int c|f|^2}[1+O_f(t^{-1})].      \tag{HFL4}
\]
Every individual hypothetical zero contradicts an $o(a_t)$
bound, whatever its height or multiplicity; no uniform
constant over zeros is required. Equivalently, the exact
nondecreasing flow makes that bound force $E_t=0$ for every
finite $t$, and Theorem~\ref{thm:heat-gap-transfer} applies.
\end{proof}

\begin{theorem}[Coherent detection on an exponential-weight boundary]
\label{thm:heat-boundary-detection}
For $\alpha>0$ put $c_\alpha(x)=\cosh(\alpha x)$ and
\[
 C_{-,\alpha}=\sup_{\substack{0\ne f\in C_c^\infty\\\int cf=0}}
                         \frac{[-q_B(f)]_+}{\int c_\alpha|f|^2}.
                                                               \tag{HBC1}
\]
The constraint uses the original $c=\cosh(x/2)$.
If $\Xi$ has a pair $z_\pm=\gamma\pm i\delta$, $\gamma,\delta>0$,
then $C_{-,\alpha}=\infty$ whenever $\delta\ge\alpha/2$.
This holds even when the supremum is restricted to either
parity sector and to exact pole-null tests.
Consequently $C_-<\infty$ would imply the open strip
$1/4<\Re\rho<3/4$ for every nontrivial zero.
Finiteness itself is not proved.
\end{theorem}
\begin{proof}
Choose the exact pole-null interpolant $f_0$ with values
$1,-1$ at $z_+,z_-$ and zero at every other distinct node,
using Lemma~\ref{lem:pole-free-interpolation}.
For either parity $\varepsilon\in\{1,-1\}$ put
$w(x)=f_0(x)+\varepsilon f_0(-x)$. Its values at
$z_+,z_-,-z_+,-z_-$ are $1,-1,\varepsilon,-\varepsilon$.
All other node values and both poles vanish. Xi symmetry
gives the same multiplicity $m\ge1$ at these four distinct
nodes. The function and every fixed derivative have
double-exponential tails.

Fix $D\ge1$ with $\gamma D\in2\pi\Z$, let $v_j=jD$, and put
\[
 h_j(x)=e^{-\delta v_j}\{w(x-v_j)+w(x+v_j)\},\quad
 H_J=\sum_{j=1}^Jh_j,\quad
 S_J=\sum_{j=1}^J(1+e^{-2\delta v_j})\ge J .
                                                               \tag{HBC2}
\]
These functions preserve parity and both zero pole values.
Phase locking gives $\widehat h_j(z_\pm)
=\pm(1+e^{-2\delta v_j})$ and the stated parity relations
at the negative nodes. The all-height identity therefore gives
\[
                         q_B(H_J)=-4mS_J^2\le-4mJ^2 .  \tag{HBC3}
\]
This can equivalently be read as its fixed-$J$ cutoff limit
using (TE); all unselected zero coefficients are exactly zero.

For $d\ge0$ define
$K_\alpha(d)=\int c_\alpha(u)|w(u-d/2)w(u+d/2)|\dd u$.
The tail bound $|w(x)|\le A e^{-b e^{2|x|}}$ implies
$K_\alpha(d)\le A_\alpha e^{-b_\alpha e^d}$.
Indeed the maximum of $|u-d/2|,|u+d/2|$ is at least
$|u|$ and $d/2$; split its exponent in halves to retain
both that decay in $d$ and an integrable weighted majorant
in $u$. For the two same-side overlaps, translating to
the midpoint uses
$c_\alpha(u+(v_i+v_j)/2)\le
e^{\alpha(v_i+v_j)/2}c_\alpha(u)$.
For each cross overlap the additional factor is
$e^{-\alpha\min(v_i,v_j)}$.
Since $\delta\ge\alpha/2$,
\[
 |\langle h_i,h_j\rangle_{L^2(c_\alpha\dd x)}|
 \le2K_\alpha(D|i-j|)+2K_\alpha(D(i+j)).                 \tag{HBC4}
\]
Its row sums are bounded by the finite constant
$M_\alpha=2\sum_{k\in\Z}K_\alpha(D|k|)
+2\sum_{k\ge2}K_\alpha(Dk)$, independent of $J$.
Schur's inequality gives
$\|\sum a_jh_j\|_{L^2(c_\alpha\dd x)}^2
\le M_\alpha\sum|a_j|^2$, hence
$\int c_\alpha|H_J|^2\le M_\alpha J$.

Now hold $J$ fixed. An even smooth cutoff preserves its
parity, and the fixed even or odd pole corrector in the proof
of Lemma~\ref{lem:pole-free-interpolation} restores both pools
exactly with a coefficient tending to zero.
Double-exponential tails give weighted norm convergence;
(TE) and (ZF) give convergence of the complete form, including
all multiplicities. The fixed centers up to $JD$ only change
constants in this fixed-$J$ limit.
There is consequently a compact same-parity, exact-pole-null
$f_J$ with $q_B(f_J)\le-2mJ^2$ and
$\int c_\alpha|f_J|^2\le(M_\alpha+1)J$.
Its negative quotient tends to infinity with $J$.
No uniform joint $J$/cutoff exhaustion was assumed.
This proves the theorem, including the boundary equality.
\end{proof}

\subsubsection{Optimal isolated-zero costs at every heat scale}
We now determine the optimal norm orders for an individual isolated
mirror pair. At and beyond the quarter boundary we also attain the
unrestricted two-evaluation lower bound to its exact leading constant,
while preserving every other zero value and both poles. The premise
that the selected off-line pair exists remains explicit.

Use the physical Hilbert space
\[
 \mathfrak H_t=\{f:f/p_t\in H_t\},\qquad
 \|f\|_{\mathfrak H_t}^2=\int\frac{c(x)}{p_t(x)}|f(x)|^2\dd x,
 \qquad N_t(\delta)=\int\frac{p_t(x)}{c(x)}e^{2\delta x}\dd x.
                                                               \tag{HOC1}
\]
The first variable of each Hilbert product is antilinear, as for
the full form. Weighted Gaussian tails make every complex Fourier
evaluation continuous on $\mathfrak H_t$.

\begin{theorem}[Exact evaluation costs and the quarter transition]
\label{thm:heat-evaluation-cost}
Let $t>0$, $\gamma\in\R$ and $0<\delta<1/2$, and put
$z_\pm=\gamma\pm i\delta$. Then
\[
 \inf_{\widehat f(z_+)=1,\ \widehat f(z_-)=-1}
        \|f\|_{\mathfrak H_t}^2=\frac{2}{N_t(\delta)-N_t(0)}.
                                                               \tag{HOC2}
\]
The unique minimizer is
\[
 f_*(x)=\frac{2p_t(x)e^{i\gamma x}\sinh(\delta x)}
                         {c(x)[N_t(\delta)-N_t(0)]}.
                                                               \tag{HOC3}
\]
It satisfies $f_*/p_t\in\mathcal V_t$. Its remaining zero and
pole values are not asserted to vanish. For each fixed $\delta$,
\[
\begin{aligned}
 N_t(\delta)-N_t(0)&\sim2\pi M_0a_t[\sec(2\pi\delta)-1],
                                      &&0<\delta<1/4,\\
 N_t(1/4)&=M_0e^{-t/4}, &&t>0,\\
 N_t(\delta)-N_t(0)&\sim2\Xi(i(2\delta-1/2))
                             e^{-(2\delta-4\delta^2)t},
                                      &&1/4<\delta<1/2.
\end{aligned}                                                \tag{HOC4}
\]
For any fixed $A<\infty$, uniformly for $|\eta|\le A$ and
$\delta_t=1/4+\eta/\sqrt t$,
\[
 N_t(\delta_t)=2M_0e^{-t/4+4\eta^2}
           [\Phi_{\rm N}(2\sqrt2\eta)+o_A(1)],\qquad
 N_t(0)=o(N_t(\delta_t)),                                    \tag{HOC5}
\]
where $\Phi_{\rm N}$ is the standard normal cumulative distribution,
distinct from the theta density $\Phi$.
\end{theorem}
\begin{proof}
The physical Riesz representer is
$e_{t,z}(x)=p_t(x)c(x)^{-1}e^{i\overline z x}$, since
$\widehat f(z)=\langle e_{t,z},f\rangle_{\mathfrak H_t}$.
Evenness gives the exact two-coordinate Gram matrix
\[
 \begin{pmatrix}N_t(\delta)&N_t(0)\\N_t(0)&N_t(\delta)\end{pmatrix}.
                                                               \tag{HOC6}
\]
Its difference eigenvalue is
$\int(p_t/c)[\cosh(2\delta x)-1]\dd x>0$.
Projection onto these two representers gives (HOC2) and (HOC3).
Additional zero or pole constraints can only increase the infimum.

For later domain estimates note, uniformly for $t\ge1$, that
$\|p_t^{(j)}\|_\infty\le C_j a_tt^{-j/2}$ for $j=0,1,2$.
The second-difference expression for $\mathcal A$ therefore gives
$\|\mathcal A p_t\|_\infty\le Ca_t$. The exact ground identity implies
\[
 S_t^{\rm pr}:=\sum_{n\ge2}\frac{\Lambda(n)}{\sqrt n}
         [p_t(\cdot-\log n)+p_t(\cdot+\log n)]
             =c+\mathcal A p_t\le Cc .                       \tag{HOC7}
\]
In particular the total prime conductance is finite, including every
prime power. The ratio $f_*/p_t$ and its first derivative are bounded
for $\delta<1/2$. Near-gamma energy uses
$\int_0^1r^2k(r)\dd r<\infty$; far-gamma and prime energies have
finite conductance. Cutoffs of this bounded ratio converge in the
actual form norm, proving its stated membership.

We have $p_t\le M_0a_t$ and $p_t/a_t\to M_0$ at every fixed point.
For $0\le\delta<1/4$, dominated convergence gives
\[
 N_t(\delta)/a_t\longrightarrow
 M_0\int e^{2\delta x}/c(x)\dd x=2\pi M_0\sec(2\pi\delta).
                                                               \tag{HOC8}
\]
Indeed $u=e^x$ makes the integral
$2\int_0^\infty u^{2\delta-1/2}(1+u)^{-1}\dd u$.
The beta-integral and reflection formulas apply with parameters
$\lambda=2\delta+1/2>0$, $1-\lambda>0$ and $\lambda\notin\Z$;
see \cite[eqs.~5.12.1, 5.12.3, 5.5.3]{DLMF}. At the boundary,
$e^{x/2}/c(x)=1+\kappa(x)$, so evenness proves the exact identity
$N_t(1/4)=\int p_t=M_0e^{-t/4}$.

Put $b=2\delta-1/2$. Gaussian completion and positive Tonelli give,
for every real $b$, the exact identity
\[
 N_t(\delta)=2e^{-t/4+tb^2}\Xi(ib)
                 \mathbb E[(1+e^{-X_{t,b}})^{-1}],\qquad
 X_{t,b}=2tb+Y_b+Z_t,                                      \tag{HOC9}
\]
where $Y_b$ has density $e^{by}\Phi(y)/\Xi(ib)$ and
$Z_t\sim N(0,2t)$ is independent. The Laplace integral
$\Xi(ib)=\int e^{by}\Phi(y)\dd y$ is positive. For fixed $b>0$,
$X_{t,b}\to+\infty$ in probability; bounded convergence of the
logistic factor and $1/4-b^2=2\delta-4\delta^2$ prove (HOC4).

For (HOC5), take $b=2\eta/\sqrt t$. The tilted variables $Y_b$
have uniformly bounded first absolute moments for $|\eta|\le A$,
and $\Xi(ib)\to M_0$ uniformly. Now
$X_{t,b}/\sqrt{2t}=2\sqrt2\eta+Z+Y_b/\sqrt{2t}$.
Replacing the logistic by $1_{\{X>0\}}$ costs at most
$e^{-B}+\mathbb P(|X|\le B)\le e^{-B}+CB/\sqrt t$,
uniformly by the conditional Gaussian density. Take $B=\log t$.
The indicator expectation differs from $\Phi_{\rm N}(2\sqrt2\eta)$
by $O_A(t^{-1/2})$, using the bounded first moment and the
Lipschitz bound for the normal cumulative distribution.
This proves (HOC5); (HOC8) makes $N_t(0)$ a factor $t^{-1/2}$ smaller.
\end{proof}

\begin{theorem}[All three optimal isolated-pair norm orders]
\label{thm:optimal-isolated-heat-orders}
Suppose $z_\pm=\gamma\pm i\delta$ are an actual Xi mirror pair,
with $\gamma>0$, $0<\delta<1/2$ and common multiplicity $m$.
For $\varepsilon\in\{1,-1\}$ let $\mathcal I_{t,\varepsilon}$
consist of $f\in\mathfrak H_t$ with $f/p_t\in\mathcal V_t$,
$f(-x)=\varepsilon f(x)$, values $1,-1$ at $z_+,z_-$,
the corresponding parity values at $-z_+,-z_-$, and zero values
at every other Xi node and at both poles. Then
\[
 \inf_{f\in\mathcal I_{t,\varepsilon}}\|f\|_{\mathfrak H_t}^2
 \asymp
 \begin{cases}
  a_t^{-1},&0<\delta<1/4,\\
  e^{t/4},&\delta=1/4,\\
  e^{(2\delta-4\delta^2)t},&1/4<\delta<1/2.
 \end{cases}                                                \tag{HIO1}
\]
Constants may depend on the selected pair and interpolant, but not
on $t$. The constructed global profiles have $q_A=q_B=-4m$.
Compact smooth, same-parity, exact-pole-null approximants retain
the displayed norm and negative-depth orders uniformly for growing $t$.
Their other zero values are approximate, with all-height error control.
\end{theorem}
\begin{proof}
The interpolant $w$ in the proof of
Theorem~\ref{thm:heat-boundary-detection} has all the required
four-node values. For $j=0,1,2$,
\[
 |w^{(j)}(x)|\le C_j e^{-c_we^{2|x|}}.                       \tag{HIO2}
\]
Fix $D\ge1$ with $\gamma D\in2\pi\Z$, and for positive centers
$v_j\in D\Z$ use
\[
 h_t=\sum_{j=1}^{J_t}e^{-\delta v_j}
             [w(\cdot-v_j)+w(\cdot+v_j)],\qquad
 S_t=\sum_{j=1}^{J_t}(1+e^{-2\delta v_j}).                   \tag{HIO3}
\]
Exact Fourier multiplication preserves every unselected value and
both poles. Phase locking gives values $S_t,-S_t$ at the positive
pair and the stated parity values at the negative pair. Thus
$J_t\le S_t\le2J_t$ and $q_A(h_t)=q_B(h_t)=-4mS_t^2$.

A fixed positive Phi mass slice $[-B,B]$ gives, for all $x\in\R$ and $t\ge1$,
\[
 p_t(x)\ge m_Ba_te^{-(|x|+B)^2/(4t)},\qquad
 \frac{c(x)}{p_t(x)}\le Ca_t^{-1}
              e^{|x|/2+x^2/(4t)+B|x|/(2t)}.                  \tag{HIO4}
\]
The coefficient $x^2/(4t)$ is essential. Put $b=2\delta-1/2$.
For signed centers $r,s$ with $|r|,|s|\le C_0t$, write
$z=(r+s)/2$, $d=r-s$. Midpoint substitution in the weighted overlap
and (HIO2) give
\[
\begin{split}
 &e^{-\delta(|r|+|s|)}
    |\langle w(\cdot-r),w(\cdot-s)\rangle_{\mathfrak H_t}|\\
 &\quad\le Ca_t^{-1}
    e^{-\delta(|r|+|s|)+|z|/2+z^2/(4t)}K(d),
 \qquad K(d)\le C_Ae^{-d^2/2-A|d|}. 
\end{split}                                                 \tag{HIO5}
\]
For the last bound, (HIO2) implies
$|w(y)|\le C_Ae^{-y^2-A|y|}$. The product at $u\pm d/2$
is bounded by $C_Ae^{-2u^2-d^2/2-A|d|}$; the remaining
shifted weight is at most $e^{C|u|+u^2/4}$. Its integral is
uniformly finite. The factor $B|z|/(2t)$ is uniformly bounded.
For $H(v)=-bv+v^2/(4t)$ the signed midpoint factor satisfies
\[
 -\delta(|r|+|s|)+|z|/2+z^2/(4t)
                     \le\tfrac12[H(|r|)+H(|s|)].             \tag{HIO6}
\]
The difference equals
$-[|r|+|s|-|r+s|]/4-(r-s)^2/(16t)\le0$.
In particular opposite lobes have not been discarded.

At $\delta=1/4$ choose $v_j=jD\le\sqrt t$. Above the boundary,
choose the grid points in $[2bt-\sqrt t,2bt+\sqrt t]$.
They are positive for large $t$ and $J_t\asymp\sqrt t$ in both cases.
The corresponding $H(v_j)$ are at most $1/4$ and $-b^2t+1/4$,
respectively. The signed centers are uniformly separated and the
Schur row sums of $K(r-s)$ are bounded. Hence
\[
 \|h_t\|_{\mathfrak H_t}^2\le CJ_ta_t^{-1}
       \begin{cases}1,&\delta=1/4,\\e^{-b^2t},&\delta>1/4.
       \end{cases}                                         \tag{HIO7}
\]
Normalize by $S_t\asymp J_t$. This gives the two claimed upper
orders without a polynomial loss. Below the boundary use $w$
itself: (HIO4) and its double-exponential tails give
$\|w\|_{\mathfrak H_t}^2\le Ca_t^{-1}$.
The matching lower orders follow from (HOC2)--(HOC4).

Here are the uniform domain and compactification details.
Put $V_t=\max_jv_j=O(t)$ and cut off at
$|x|\le V_t+2\log(t+2)$, with a fixed width-one transition.
Beyond this cutoff, $h_t$ and its first two derivatives are bounded
by $CJ_te^{-ce^{2(|x|-V_t)}}$. The lower mass slice and uniform
Gaussian derivative bounds give, for $|u|\le1$, the rough bounds
$|p_t'/p_t|,p_t(x+u)/p_t(x),p_t(x)/p_t(x+u)
\le Ce^{Cx^2+C|x|}$, uniformly for $t\ge1$.
The near-gamma cutoff energy is therefore bounded by
$C\int e^{Cx^2+C|x|}(|h_t|^2+|h_t'|^2)\dd x$ over that tail,
using the finite $r^2k(r)$ moment. It is at most
$CJ_t^2e^{CV_t^2-c(t+2)^4}$. The weighted norm tail follows
from (HIO4), far-gamma energy from its integrable kernel,
and all-prime energy from (HOC7). The two-derivative (TE) estimate
with weight $e^{|x|/2}$ has the same superexponential order and
retains all heights and multiplicities. Fixed even or odd pole
correctors restore the exact pole values with coefficients at most
$CJ_te^{V_t/2-c(t+2)^4}$; their norm and energy are $O(a_t^{-1})$.
Thus all errors are bounded by $Ce^{Ct^2-c(t+2)^4}$ after normalization.
They vanish in the actual form norm and preserve (HIO1).
This also proves membership of the noncompact profiles and the
stated complete-form values by (ZF) and core closure.
\end{proof}

\begin{theorem}[The exact leading cost of full isolation]
\label{thm:optimal-isolated-heat-constant}
Suppose $z_\pm=\gamma\pm i\delta$ are an actual Xi mirror pair
of common multiplicity $m$, with $\gamma>0$ and $1/4\le\delta<1/2$.
Let $\mathcal I_t$ be the analogous form-domain class with values
$1,-1$ at these two nodes and zero at every other Xi node and both
poles, without a parity constraint. Then
\[
 \inf_{f\in\mathcal I_t}\|f\|_{\mathfrak H_t}^2
                  \sim\frac{2}{N_t(\delta)-N_t(0)}.           \tag{HIC1}
\]
An explicit family $f_t$ attains this equivalent and satisfies
$q_A(f_t)=q_B(f_t)=-2m$. Its negative Rayleigh depth is
$m[N_t(\delta)-N_t(0)](1+o(1))$.
Exact-pole-null compact approximants preserve both equivalents.
\end{theorem}
\begin{proof}
Choose the fixed pole-null interpolant $w$ with values $1,-1$
only at $z_+,z_-$, by Lemma~\ref{lem:pole-free-interpolation}.
It has (HIO2). Put $b=2\delta-1/2\ge0$, $d=\delta-1/2<0$, and
\[
 K_t(v)=2p_t(v)e^{d|v|}e^{i\gamma v},\qquad
 A_t=2\int_0^\infty p_t(v)(e^{bv}+e^{-v/2})\dd v,
 \qquad f_t=(w*K_t)/A_t.                                    \tag{HIC2}
\]
Gaussian tails and (HIO2) justify absolute Fourier multiplication
at every complex point. Direct evaluation gives
$\widehat K_t(z_+)=\widehat K_t(z_-)=A_t$.
Thus all the required zero and pole values are exact. Gaussian
completion gives
$2\int_0^\infty p_t(v)e^{bv}\dd v
\sim2\Xi(ib)e^{-t/4+tb^2}$ when $b>0$; at $b=0$ that integral
is exactly $M_0e^{-t/4}$. The remaining term is $O(a_t)$.
Consequently $A_t\sim N_t(\delta)$ and $N_t(0)=o(N_t(\delta))$.

Let $R_t^{\rm ev}=e_{t,z_+}-e_{t,z_-}$ and
$Q_t^{\rm ev}(x)=\operatorname{sgn}(x)K_t(x)$.
Then $\|R_t^{\rm ev}\|_{\mathfrak H_t}^2=2[N_t(\delta)-N_t(0)]$,
and, off zero,
\[
 R_t^{\rm ev}(x)=Q_t^{\rm ev}(x)
                 \frac{1-e^{-2\delta|x|}}{1+e^{-|x|}}.
                                                               \tag{HIC3}
\]
Using $p_t\le M_0a_t$ gives
$\|R_t^{\rm ev}-Q_t^{\rm ev}\|_{\mathfrak H_t}^2=O(a_t)$:
the residual integral is dominated by
$Ce^{-(2\delta+1/2)|x|}$ after taking out $a_t$.
This is $o(N_t(\delta))$.

We prove the substantive convolution limit
\[
 \|w*K_t-Q_t^{\rm ev}\|_{\mathfrak H_t}
                                  =o(\sqrt{N_t(\delta)}).     \tag{HIC4}
\]
The two interpolation values imply exactly
\[
 Q_t^{\rm ev}(x)=K_t(x)
            \int w(u)e^{-i\gamma u+\delta u\operatorname{sgn}(x)}\dd u.
\]
Write
$p_t(x)=a_te^{-x^2/(4t)}M_t(x/(2t))$, where
$M_t(q)=\int\Phi(y)e^{qy-y^2/(4t)}\dd y$.
The fixed mass slice and $\Phi\le C_\epsilon G_\epsilon$ give,
for $t\ge1$,
\[
 ce^{-B|q|}\le M_t(q)\le C_\epsilon e^{\epsilon q^2},\qquad
 p_t(x)\le C_\epsilon a_te^{-x^2/[4(t+\epsilon)]}.              \tag{HIC5}
\]
After $v=x-u$ and $c(v+u)\le e^{|u|/2}c(v)$, these bounds give
\[
 \|K_t(\cdot-u)\|_{\mathfrak H_t}^2
 \le Ce^{C|u|+u^2/(4t)}\int p_t(v)
   e^{b|v|+(|u|+B)|v|/(2t)+\epsilon v^2/(4t^2)}\dd v.
                                                               \tag{HIC6}
\]
Take $\epsilon=1/10$. The combined Gaussian coefficient is
$\sigma_t=1/[4(t+\epsilon)]-\epsilon/(4t^2)>0$, with
$1/(4\sigma_t)=t+O(1)$ uniformly for $t\ge1$.
Completing its square proves
\[
 \|K_t(\cdot-u)\|_{\mathfrak H_t}^2
      \le Ca_t\sqrt t\,e^{b^2t}e^{C|u|+Cu^2}
      \le CN_t(\delta)e^{C|u|+Cu^2}.                         \tag{HIC7}
\]
In particular no loss proportional to $t$ occurs in the exponent.

For fixed $u$, (HIC6) puts the positive and negative norm lobes
at $\pm2bt+O_u(1)$ with variance $O(t)$. Outside windows of width
$M\sqrt t$ the squared norm is at most
$C_ue^{-cM^2}N_t(\delta)$, for both the shifted and unshifted kernel.
When $b=0$, the additional region $|x|\le\eta\sqrt t$ has norm at
most $C_u\eta N_t(\delta)$, since its density is bounded by $C_ua_t$.
On the retained bulk, $|x|\to\infty$ and
$x/(2t)\to b\operatorname{sgn}(x)$ uniformly.
Locally uniform dominated convergence gives $M_t(q)\to\Xi(iq)>0$,
and hence
\[
 \frac{p_t(x-u)}{p_t(x)}
 =e^{xu/(2t)-u^2/(4t)}
              \frac{M_t(x/(2t)-u/(2t))}{M_t(x/(2t))}
 \longrightarrow e^{bu\operatorname{sgn}(x)}.
\]
The absolute-value and phase factors in $K_t$ yield the full ratio
$e^{-i\gamma u+\delta u\operatorname{sgn}(x)}$, because $b-d=\delta$.
Let $t\to\infty$, then $M\to\infty$ and, at the boundary,
$\eta\downarrow0$. The displayed tail bounds prove fixed-$u$
weighted $L^2$ convergence of the difference in (HIC4).
Minkowski and (HIC7) dominate its normalized norm by
$Ce^{C|u|+Cu^2}$, which is integrable against $|w(u)|\dd u$ by (HIO2).
Dominated convergence over $u$ now proves (HIC4).
Together with (HIC3) and $A_t\sim N_t(\delta)$ this gives (HIC1),
with the matching lower bound (HOC2).

We give the domain and uniform cutoff bounds separately.
For fixed $t$, convolving (HIO2) with the Gaussian $p_t$ bounds the
physical numerator and its derivatives by Gaussians of width
$t+\epsilon$, with arbitrarily small fixed extra width.
The lower slice (HIO4) makes the weighted norm finite.
Near-gamma ratios have an arbitrarily small quadratic majorant at
fixed $t$ from (HIC5); far-gamma and prime energy tails are bounded
by the weighted norm using (HOC7). Thus $f_t/p_t\in\mathcal V_t$.
More explicitly, uniformly for large $t$ and $j=0,1,2$,
\[
 |f_t^{(j)}(x)|\le C_jt^{-1/2}e^{-x^2/[4(t+C_0)]}.             \tag{HIC8}
\]
Indeed $A_t\asymp e^{-t/4+tb^2}$ and the factor $e^{-tb^2}\le1$
can be discarded after convolution. Differentiating $M_t$ and
using its lower slice gives, for $|u|\le1$,
\[
 |p_t'/p_t|,\quad p_t(x+u)/p_t(x),\quad p_t(x)/p_t(x+u)
 \le C(1+|x|/t)e^{C\epsilon x^2/t^2+C|x|/t+C}.                \tag{HIC9}
\]
Cut off at $|x|\le t^2$ with fixed transition width.
Equations (HIO4), (HIC8)--(HIC9) bound the norm and all energy
tails by $Ce^{-ct^3+C't^2}$. The negative physical Gaussian
coefficient is of order $-x^2/t$, whereas the ratio loss is only
$O(\epsilon x^2/t^2+|x|/t)$; near-gamma uses the finite $r^2k(r)$
moment and the all-prime tail uses (HOC7). The two-derivative
all-height (TE) bound with $e^{|x|/2}$ has the same order.
Fixed compact pole correctors restore both pole values exactly;
their coefficients have that small order and their norm and energy
are $O(a_t^{-1})$. Thus the corrected compact profiles retain
(HIC1), and (ZF) gives $q_A=-2m+o(1)$ including all multiplicities.
Taking the actual form closure also proves the exact complete-form
value $q_A(f_t)=-2m$. No leading-constant claim is made below
$\delta=1/4$, where the central norm measure does not escape.
\end{proof}

\begin{theorem}[Full heat depth under a finite-exception hypothesis]
\label{thm:finite-off-heat-depth}
Suppose the set of distinct off-line Xi nodes is nonempty and finite.
Put
\[
 B=\max_{z\ {\rm off}}|\Im z|,\qquad
 m_* =\max_{|\Im z|=B}m_z .
\]
Let $d_t$ be the supremum of
$[-q_A(p_tg)]_+/\|g\|_{H_t}^2$ over nonzero compact smooth tests,
and $d_t^0$ the same supremum with both pole values zero.
The same suprema may be taken on the corresponding actual form
closures. If $B\ge1/4$, then
\[
 d_t^0\sim d_t\sim 1-\beta_t\sim m_*N_t(B).                  \tag{HFD1}
\]
If $0<B<1/4$, all three depths are $\asymp a_t$.
Finiteness of the off-line set is an unproved premise, and no
uniformity over varying finite zero configurations is asserted.
\end{theorem}
\begin{proof}
On the compact core (ZF) splits $q_A(f)$ into nonnegative real-node
squares and the finite off-line form
$\sum_{z\ {\rm off}}m_z\overline{\widehat f(\bar z)}\widehat f(z)$.
The latter is a bounded finite-rank selfadjoint form in $\mathfrak H_t$.
Its representer Gram is exactly
\[
 G_t(z,w)=\int\frac{p_t(x)}{c(x)}
          e^{(\Im z+\Im w)x}e^{i(\Re w-\Re z)x}\dd x.          \tag{HFD2}
\]
The diagonal is $N_t(|\Im z|)$. If $J$ swaps conjugate nodes and
$M=\operatorname{diag}(m_z)$, its nonzero eigenvalues are those of
$G_t^{1/2}JM G_t^{1/2}$ by the finite-rank polar decomposition.
No inverse of $G_t$ is needed, and $J,M$ commute.

For $B\ge1/4$, we claim $G_t/N_t(B)\to P_B$ in finite-dimensional
norm, where $P_B$ selects $|\Im z|=B$.
Equation (HOC4) and Cauchy--Schwarz handle any entry with a submaximal
index. Opposite imaginary signs at the top have absolute cross entry
at most $N_t(0)=o(N_t(B))$. Distinct same-sign top nodes have a
nonzero real-coordinate difference $\tau$. At $B=1/4$ their weight
is $p_t(1+\kappa)$ or its mirror, whose derivative satisfies
\[
 \|(p_t(1+\kappa))'\|_1
 \le2\|p_t'\|_1+\int p_t\kappa'\dd x
 \le Ce^{-t/4}/\sqrt t=o(N_t(1/4)).                          \tag{HFD3}
\]
Here $\|p_t'\|_1\le M_0e^{-t/4}\|G_t'\|_1$ and
$p_t\le M_0a_t$. Integrating by parts annihilates each $\tau$-entry.
For $B>1/4$ the weight is exactly
$2e^{-t/4+tb^2}\Xi(ib)\ell(x)\rho_{t,b}(x)$, with $b=2B-1/2$,
$\ell(x)=(1+e^{-x})^{-1}$ and $\rho_{t,b}$ the density in (HOC9).
Its derivative has normalized $L^1$ norm tending to zero:
$\|\rho_{t,b}'\|_1\le\|G_t'\|_1=O(t^{-1/2})$, and
$\int\ell'\rho_{t,b}\to0$ because $X_{t,b}\to+\infty$ in
probability and $\ell'$ is bounded and vanishes there.
Integration by parts again gives a vanishing cross entry.
Evenness treats the negative sign. This proves the Gram limit.

Continuity of the finite matrix square root shows that the least
off-form eigenvalue divided by $N_t(B)$ tends to the least eigenvalue
of $P_BJM P_B$, namely $-m_*$. The real-node squares can only raise
the form, so $d_t\le[m_*+o(1)]N_t(B)$.
Choose a top pair of multiplicity $m_*$. The exact isolated profiles
of Theorem~\ref{thm:optimal-isolated-heat-constant}, including their
compact pole-null approximants, give
$d_t^0\ge[m_*+o(1)]N_t(B)$. Thus $d_t^0\sim d_t\sim m_*N_t(B)$.
For $B<1/4$, the finite off-form is bounded below by
$-\sum_{z\ {\rm off}}m_zN_t(|\Im z|)\|f\|_{\mathfrak H_t}^2
\ge-Ca_t\|f\|_{\mathfrak H_t}^2$.
A fixed compact pole-null negative detector, whose norm costs are
$\asymp a_t^{-1}$, gives the matching lower depth.
These inequalities extend to the actual core closure, since all
finitely many off-line evaluations are norm continuous.

Finally $\mathcal W_t$ is unchanged by subtracting the mean:
its mixed value against $1$ is zero by (HGF5).
Such subtraction reduces the norm. On the mean-zero domain,
$Q_t(g,g)=\mathcal W_t(g)+2|\langle\kappa,g\rangle|^2$, and (HGP3)
gives
\[
 d_t^0\le1-\beta_t\le d_t\le(1+2/t)(1-\beta_t).              \tag{HFD4}
\]
Core density preserves the suprema and exact moment conditions.
The preceding matching bounds now prove (HFD1) and the central
orders. The finite-exception premise has not been established.
\end{proof}

These results give the sharp cost and attainable depth of complete
individual detection. They do not exclude any hypothetical pair:
a contradictory global lower estimate, still uniform over every
form vector, must separately be proved. The dependency on a finite exceptional set occurs only in the
final theorem of this subsection.

\begin{proposition}[The ordinary index survives every finite heat scale]
\label{prop:heat-index}
On $\mathcal V_t\cap H_t^0$, the negative form index of $Q_t$
is exactly $2K_+$ for every fixed $t>0$, including infinity.
\end{proposition}
\begin{proof}
On compact tests $f=p_tg$ identifies the constraint with
$\int cf=0$. In the even sector this annihilates both pole
values, so (BI) gives index $K_+$. In the odd sector the
constraint is automatic and the complete odd $B$ index is
$K_+$, again by (BI). Orthogonal parity decomposition adds
the two. Any finite negative domain basis can be approximated
by the mean-zero compact form core without changing the
negative matrix or its dimension. This proves equality
on the closed domain, and arbitrarily many finite pairs
give the infinite case.
\end{proof}

\begin{proposition}[Actual Loewner heat monotonicity fails]
\label{prop:heat-loewner-failure}
Transport $Q_t$ to $L^2(\dd x)$ by
$U_tg=\sqrt{cp_t}\,g$, and write $R_t=U_t(L_t-I)U_t^{-1}$.
For every $0<s<t$, $q_{R_t}-q_{R_s}$ has both signs on the
common odd compact smooth core.
\end{proposition}
\begin{proof}
Put $w_t=p_t/c$. On that core
$q_{R_t}(u)=q_B(\sqrt{w_t}\,u)$.
The weighted masses of $p_t,p_s$ agree, whereas their
unweighted integrals are $e^{-t/4}M_0,e^{-s/4}M_0$.
Thus the even continuous function $w_t-w_s$ has a positive
and a negative open interval in $(0,\infty)$.
For a real nonzero bump $\eta$ in either interval take
$u_q=e^{iqx}\eta(x)-e^{-iqx}\eta(-x)$.
It is odd and hence mean zero for both grounds.
All prime terms are finite on this fixed support and
uniformly bounded in $q$. The actual digamma multiplier has
$a(q)=\log q+O(1)$, by the same sectorial expansion
\cite[eq.~5.11.2]{DLMF} used in (JM).
Fourier modulation and the rapid transform tail give
$q_{\rm arch}(e^{iqx}h)=a(q)\|h\|_2^2+O_h(1)$ for a fixed
compact smooth $h$. The two supports have positive separation,
so their archimedean cross-kernel is bounded too. Therefore
\[
 q_{R_t}(u_q)-q_{R_s}(u_q)
   =2\log q\int(w_t-w_s)|\eta|^2\dd x+O_{s,t,\eta}(1).    \tag{HFL5}
\]
The two choices of interval give opposite signs for large $q$.
No individual eigenvalue monotonicity is inferred or denied.
\end{proof}

The exact densities also satisfy
$\partial_tw_t=w_t''+\kappa w_t'$ and
$\partial_t(cp_t)=(cp_t)''-(\kappa cp_t)'$.
These equations do not supply the missing sharp gap.
The actual global target remains a bound uniform over all
mean-zero form vectors,
$\mathcal D_t(g)\ge[1-o(a_t)]\|g\|_{H_t}^2$ on an unbounded
sequence. The exact scalar flow shows that its normalized
negative leading coefficient is $M_0C_-$ and cannot decrease
under heat. A favorable finite compression or a fixed-support
limit is insufficient. The proof obligation is still to show
$C_-=0$, or otherwise to exclude every negative direction.

\section{Shifted heat grounds expose every positive displacement}
\label{sec:shifted-heat}

The quarter-strip threshold above belongs to its particular limiting
norm. A moving, positive convolution changes that norm without
changing any nonreal zero. We retain the full ground equation,
the score correction, and the costs of the moving detector.
Throughout this section $c(x)=\cosh(x/2)$, $s(x)=\sinh(x/2)$,
$\kappa=s/c$, and $M(b)=\int e^{by}\Phi(y)\dd y=\Xi(ib)>0$.

\begin{theorem}[Shifted heat and uniform pole costs]
\label{thm:shifted-heat-ground}
For $t>0$ and $0\le b\le1/2$ put
\[
 p_{t,b}(x)=\frac{p_t(x-2bt)+p_t(x+2bt)}{2\cosh(bt)},
 \qquad
 \sigma_{t,b}=\frac{e^{-(b^2+1/4)t}}
                    {\sqrt{4\pi t}\cosh(bt)},\qquad
 \dd\mu_{t,b}=cp_{t,b}\dd x.                         \tag{SHG1}
\]
These are smooth, even, strictly positive grounds, with
$Bp_{t,b}=-c$, $\int cp_{t,b}=1/2$ and $Ap_{t,b}=0$.
Their transforms are exactly
\[
 \widehat p_{t,b}(z)
 =e^{-t(z^2+1/4)}\Xi(z)\frac{\cos(2btz)}{\cosh(bt)}. \tag{SHG2}
\]
Every original nonreal zero has exactly its original order;
at a real zero the order is at least its original order.
All added zeros are real.

The closed jump energy $\mathcal D_{t,b}$ uses the same gamma
kernel and every $\Lambda(n)$ as the preceding ground identity.
On its compact physical core,
\[
 q_A(p_{t,b}g)=\mathcal D_{t,b}(g)-\|g\|_\mu^2+
             2|\langle1,g\rangle_\mu|^2-
             2|\langle\kappa,g\rangle_\mu|^2.       \tag{SHG3}
\]
The score $v=-2p_{t,b}'/p_{t,b}$ belongs to the ordinary
operator domain and satisfies
\[
 \langle1,v\rangle_\mu=0,\quad
 \langle\kappa,v\rangle_\mu=1/2,\quad
 (\mathcal D_{t,b}-I)(v,g)=\langle\kappa,g\rangle_\mu,\quad
 \|v\|_\mu^2\le1/2+1/t.                            \tag{SHG4}
\]
On the mean-zero space the exact pole correction
$Sg=g-2\langle\kappa,g\rangle_\mu v$ has
$\|S\|^2\le1+2/t$, uniformly in $b$, including $b=1/2$.
\end{theorem}
\begin{proof}
The identities
$c(x-a)+c(x+a)=2\cosh(a/2)c(x)$ and the analogous identity
for $s$ give the ground equation and the two pole values.
Finite translations commute with gamma and every finite prime
cap. For fixed $t,b$ their Gaussian bounds dominate the full
positive prime degree, including every higher power. Tonelli
and the cancelled near-diagonal gamma second difference
therefore justify the full identity, rather than only a capped
one. The preceding ground-algebra proof gives (SHG3) and
its compact core closure. Translation gives (SHG2).
For $z=\gamma+i\delta$, $\delta\ne0$, $b>0$,
\[
 |\cos(2btz)|^2=\cos^2(2bt\gamma)+\sinh^2(2bt\delta)>0.
\]
At $b=0$ the mask is one. This proves the stated order claims.

For the uniform score bound the probability density
$\rho_+(x)=2e^{x/2}p_{t,b}(x)$ is exactly the law of
three independent variables
\[
 X=Y+Z+H,\qquad
 \Pr(Y\in\dd y)=2e^{y/2}\Phi(y)\dd y,\quad
 Z\sim N(t,2t),\quad
 \Pr(H=\pm2bt)=\frac{e^{\pm bt}}{2\cosh(bt)}.       \tag{SHG5}
\]
This follows by exponential tilting of the actual convolution.
All three measures have mass one. Writing $Z=t+U$, the score
of $\rho_+$ is the conditional average of $-U/(2t)$.
Differentiation is legitimate under Gaussian convolution.
Conditional Cauchy--Schwarz and Tonelli prove
\[
 I(\rho_+)=\int\frac{|\rho_+'|^2}{\rho_+}
            \le \frac{\mathbb E U^2}{4t^2}=\frac1{2t}.
\]
If $h=p'/p$, integration by parts gives
$\int\rho_+h=-1/2$ and
$\int\rho_+h^2=I(\rho_+)+1/4$.
Evenness gives the same integral for
$\rho_-=2e^{-x/2}p$, and
$\mu=(\rho_++\rho_-)/4$. Hence
$\|v\|_\mu^2=1/2+2I(\rho_+)\le1/2+1/t$.

The physical score is $-2p'$.
For each fixed $t,b$ its derivative and pole tails allow
smooth cutoffs in every fixed exponential $H^1$ norm.
Their full prime forms converge because
$\sum\Lambda(n)n^{-1/2-\eta}<\infty$ for $\eta>1/2$;
gamma and the two pole functionals converge as well.
The mixed differentiated equation $B(-2p')=s$, with the
two integrations by parts, proves the first three identities
in (SHG4) and places the score in the operator domain.
The exact rank-one algebra of the preceding score theorem
then gives the asserted projector norm.
Only this norm bound is uniform in $b$; the cutoff statements
here are for each fixed $t,b$.
\end{proof}

\begin{theorem}[Exact shifted defect flow, including the endpoint]
\label{thm:shifted-defect-flow}
For fixed $0\le b\le1/2$ define
\[
 d_{t,b}=
 \sup_{\substack{0\ne f\in C_c^\infty\\\int cf=0}}
 \frac{[-q_B(f)]_+}{\int c|f|^2/p_{t,b}},
 \qquad
 \mathfrak C_b=
 \sup_{\substack{0\ne f\in C_c^\infty\\\int cf=0}}
 \frac{[-q_B(f)]_+}{\int[c/\cosh(bx)]|f|^2}.        \tag{SHD1}
\]
Let $\eta_{t,b}$ be the negative depth of the full closed
ground form. Then, with extended limits allowed,
\[
 d_{t,b}\le\eta_{t,b}\le(1+2/t)d_{t,b},\qquad
 \frac{d_{t,b}}{\sigma_{t,b}}\nearrow M(b)\mathfrak C_b,
 \qquad
 \frac{\eta_{t,b}}{\sigma_{t,b}}\longrightarrow
                       M(b)\mathfrak C_b.           \tag{SHD2}
\]
In particular
\[
 \mathfrak C_{1/2}<\infty\ \Longleftrightarrow\ {\rm RH},
 \qquad
 {\rm RH}\ \Longrightarrow\ \mathfrak C_{1/2}=0.     \tag{SHD3}
\]
Thus a full signed $O(\sigma_{t,1/2})$ lower bound on one
unbounded sequence would prove RH. Such a bound is not
established here.
\end{theorem}
\begin{proof}
On the whole physical line Gaussian completion gives
\[
 \frac{p_{t,b}(x)}{\sigma_{t,b}}=
 \int\Phi(y)\cosh[b(x-y)]e^{-(x-y)^2/(4t)}\dd y
       \nearrow M(b)\cosh(bx).                     \tag{SHD4}
\]
The limit is uniform on each compact, by the double-exponential
theta tails. For every fixed compact $f$, its normalized
denominator decreases to
$M(b)^{-1}\int[c/\cosh(bx)]|f|^2$.
Its variational quotient increases and is bounded by
$M(b)\mathfrak C_b$. Taking the supremum of these fixed-test
limits gives the reverse lower limit, also if it is infinite.
This proves the scalar assertion without interchanging a
growing test and a nonuniform limit. Compact approximation
with a fixed mean corrector identifies $d_{t,b}$ with the
closed ordinary depth. The score correction proves the
full-depth sandwich; removing the constant component preserves
the full form and decreases its norm.

At $b=1/2$ the limiting physical weight is exactly one and
$M(1/2)=1/2$. If an original off-line pair exists, take its
pole-free fast interpolant $w$ from
Lemma~\ref{lem:pole-free-interpolation}, symmetrized into
one parity as in the coherent-boundary proof.
Choose $D\ge1$ with $\gamma D\in2\pi\mathbb Z$ and put
\[
 H_J(x)=\sum_{j=1}^J e^{-\delta jD}
                         [w(x-jD)+w(x+jD)] .
\]
All other zero values and both poles are zero. The full
formula with multiplicity $m$ gives $q_B(H_J)\le-4mJ^2$,
whereas
\[
 \|H_J\|_2\le2\|w\|_2\sum_{j\ge1}e^{-\delta jD}<\infty.
\]
For each fixed $J$ cut off first, then repair the two pools
with fixed compact correctors. The all-height strip tail
and weighted norm convergence permit norm error at most one
and $q_B(f_J)\le-2mJ^2$. Consequently
$\mathfrak C_{1/2}=\infty$ for every $\delta>0$.
There is no assumed uniform cutoff in $J$.
Conversely, under RH the proved full zero formula gives
$q_B(f)=q_A(f)+2|\int sf|^2\ge0$ on $\int cf=0$.
This proves (SHD3) as an implication and its converse.
\end{proof}

\begin{theorem}[Joint endpoint costs at every positive displacement]
\label{thm:endpoint-detection-cost}
Write $p=p_{t,1/2}$, $\sigma_t=\sigma_{t,1/2}
\sim e^{-t}/\sqrt{\pi t}$, and
\[
 N_t(\delta)=\int\frac{p(x)}{c(x)}e^{2\delta x}\dd x .
\]
For every fixed actual pair $z_\pm=\gamma\pm i\delta$,
$\gamma>0$, $0<\delta<1/2$, of multiplicity $m$,
\[
 N_t(\delta)\sim\sigma_t\sqrt{4\pi t}\,
              e^{4\delta^2t}M(2\delta-1/2),\qquad
 N_t(0)\sim\tfrac12\sigma_t\sqrt{4\pi t}.           \tag{EHC1}
\]
For the physical fast-decay class with the two pole values
zero, $F_f(z_+)=1$, $F_f(z_-)=-1$ and every other original
zero value zero, its least norm $I_t(f)=\int c|f|^2/p$ satisfies
\[
 \frac2{N_t(\delta)-N_t(0)}
 \le\inf I_t(f)\le
       \frac{K_{\gamma,\delta,m}}
                  {\sigma_t\sqrt t\,e^{4\delta^2t}}. \tag{EHC2}
\]
Moreover there are compact tests with both poles exactly
zero, support radius
$4\delta t+\sqrt t+\tfrac12\log t+B_0+2$, and strictly
negative full form, yielding
\[
 d_{t,1/2}\ge c_{\gamma,\delta,m}
                     \sigma_t\sqrt t\,e^{4\delta^2t}. \tag{EHC3}
\]
All constants may depend on the selected pair and its fixed
interpolant. No rightmost zero or uniform isolation is assumed.
In particular a bound
$d_{t_j,1/2}\le\sigma_{t_j}\exp(o(t_j))$ for one unbounded
sequence, uniformly over all ordinary mean-zero form vectors,
would exclude every off-line zero.
\end{theorem}
\begin{proof}
Put $M_t(q)=\int\Phi(y)e^{qy-y^2/(4t)}\dd y$.
On the whole line the exact moving-density identity is
\[
 \frac{p(x)}{c(x)}=\sigma_t e^{-x^2/(4t)}
 \frac{e^{x/2}M_t(x/(2t)-1/2)+
       e^{-x/2}M_t(x/(2t)+1/2)}
                  {e^{x/2}+e^{-x/2}}.              \tag{EHC4}
\]
For $t\ge1$, $m_1=\int_{-1}^1\Phi>0$ gives
$M_t(q)\ge m_1e^{-|q|-1/4}$ and consequently
\[
 \frac{c(x)}{p(x)}\le
       \frac{e^{3/4}}{m_1\sigma_t}
                         e^{x^2/(4t)+|x|/(2t)}.    \tag{EHC5}
\]
For every fixed $\epsilon>0$, the theta tails also give
$\Phi\le C_\epsilon G_\epsilon$ and
$M_t(q)\le C_\epsilon e^{\epsilon q^2}$.
Complete the square in $N_t(\delta)$ and set
$x=4\delta t+y\sqrt t$. For each fixed $y$ and $\delta>0$
the first weight in (EHC4) tends to one and its argument
to $2\delta-1/2$. With $\epsilon=1/16$ the entire integrand
after this substitution is dominated by
$C_\delta e^{-7y^2/32}$, uniformly for $t\ge1$.
Dominated convergence proves the first part of (EHC1).
For $\delta=0$, at $x=y\sqrt t$, either dominant weight
gives $M(1/2)=1/2$. The same dominator proves the second.

In $L^2(cp\dd x)$ the representer of $F_{pg}(z)$ is
$r_z=e^{i\bar z x}/c$. The exact two-point Gram is
\[
 \begin{pmatrix}N_t(\delta)&N_t(0)\\
                 N_t(0)&N_t(\delta)\end{pmatrix}.
\]
The crossphase cancels exactly. Cauchy--Schwarz applied
to $r_{z_+}-r_{z_-}$ gives the first inequality of (EHC2).
It holds before imposing other zero or pole constraints,
so remains valid afterward. On this class the complete
zero formula gives exactly $q_A(f)=-2m$.

For the upper bound use one fixed complex pole-free interpolant
$w$ whose only original zero values are $1,-1$ at $z_+,z_-$.
It satisfies, for $0\le k\le4$,
\[
 |w^{(k)}(y)|\le A e^{-cE(y)},\qquad
 E(y)=e^{2y}+e^{-2y},\qquad A,c>0.                 \tag{EHC6}
\]
These are the previously proved theta interpolation bounds,
with constants for this selected pair.
For $a=4\delta$, let
$v_j=at-\sqrt t+j$, $0\le j<J_t=\lfloor2\sqrt t\rfloor+1$,
and, for sufficiently large $t$, define
\[
 h_j(x)=e^{-\delta v_j}
       [e^{i\gamma v_j}w(x-v_j)+
        e^{-i\gamma v_j}w(x+v_j)],\quad
 H_t=\sum_jh_j,\quad
 S_t=\sum_j(1+e^{-2\delta v_j})\ge J_t.             \tag{EHC7}
\]
The explicitly chosen phases give the selected values
$\pm S_t$. Every other zero and both poles are still exactly
zero. Hence $q_A(H_t)=-2mS_t^2$.

For the two signed centers of an overlap write
$z=(\varepsilon v_j+\varepsilon'v_k)/2$,
$h=(\varepsilon v_j-\varepsilon'v_k)/2$, and $x=z+y$.
By (EHC5) its absolute weighted overlap is at most
\[
 K\sigma_t^{-1}e^{z^2/(4t)-\delta(v_j+v_k)}
   \int e^{2|y|+y^2/4}|w(y-h)w(y+h)|\dd y.        \tag{EHC8}
\]
Here $|z|/t\le a+1\le3$; all omitted factors are bounded.
For equal signs the exponent in front is at most
$-4\delta^2t+1/4$. For opposite signs it is at most
$-4\delta^2t+1/2$, since
$-4\delta^2t+2\delta\sqrt t\le1/4$.
The exact identity
$E(y-h)+E(y+h)=2\cosh(2h)E(y)$ bounds the integral
in (EHC8) by $K e^{-c_1e^{2|h|}}$.
The signed grid has uniformly bounded counts per unit interval.
Every absolute Gram row is therefore at most
$K\sigma_t^{-1}e^{-4\delta^2t}$.
Schur's inequality gives
\[
 I_t(H_t)\le KJ_t\sigma_t^{-1}e^{-4\delta^2t}.       \tag{EHC9}
\]
The fast vector $H_t/S_t$ proves the upper bound in (EHC2).
It satisfies the exact all-other-zero constraints.

For a genuine compact detector choose a smooth cutoff equal
to one up to radius $L-1$ and supported up to $L$, with
four uniformly bounded derivatives, and set
$L=at+\sqrt t+B+2$. Bounds (EHC5)--(EHC8) give
for $E_{t,L}=(1-\chi_L)H_t$, with fixed $K,c_*>0$,
\[
 \begin{split}
 I_t(E_{t,L})&\le KJ_t^2\sigma_t^{-1}e^{-4\delta^2t}
                                             e^{-c_*e^{2B}},\\
 |F_{E_{t,L}}(\gamma'+i\delta')|
 &\le KJ_t(1+|\gamma'|)^{-4}
       e^{(1/2-\delta)(at+\sqrt t)}e^{-c_*e^{2B}},
                         \qquad|\delta'|\le1/2.
 \end{split}                                      \tag{EHC10}
\]
The second bound uses four actual physical integrations
by parts, including the cutoff derivatives; the translated
theta envelope absorbs every fixed polynomial factor.
For example take $c_*=c/16$ and then $B_*$ sufficiently
large that
$(c/16)E(y)\ge2|y|+y^2/4+|y|/2+2$ for $|y|\ge B_*$.
These are explicit finite thresholds from the fixed
source envelope.

Choose $B=\tfrac12\log t+B_0$, with $B_0\ge B_*$,
$c_*e^{2B_0}>10$, and covering the two fixed pole
correctors. Since
$(1/2-\delta)(at+\sqrt t)\le3t/2$ for $t\ge1$,
the complete strip error is
$O(J_t e^{-8t}(1+|\gamma'|)^{-4})$.
The actual counting bound, including multiplicities,
gives
$\sum_{\Xi(z)=0}m_z(1+|\Re z|)^{-8}<\infty$.
Thus the newly introduced zero values have weighted
$\ell^2$ error $O(J_t e^{-8t})$.
The same estimate gives both pole errors.
Repair them exactly with two fixed compact correctors.
Their weighted squared cost is
$O(J_t^2e^{-16t}\sigma_t^{-1})$, negligible relative
to (EHC9), since $4\delta^2<1$.
The full zero formula and Cauchy--Schwarz then give
\[
 q_A(\widetilde H_t)=-2mS_t^2+
                              O(J_t^2e^{-8t})<-mJ_t^2 .
\]
These compact tests have exactly the asserted radius
and both exact pole values zero. Their other zero values
are small with the written all-height bound; exact
annihilation is only claimed for the preceding fast tests.
Dividing this strict negative bound by (EHC9) proves (EHC3).
A subexponential normalized upper bound on any unbounded
sequence contradicts it for each fixed $\delta>0$.
\end{proof}

\begin{theorem}[The endpoint interpolation cost has an attained leading constant]
\label{thm:endpoint-leading-cost}
Under the selected-pair hypotheses of
Theorem~\ref{thm:endpoint-detection-cost}, the full pole-free,
all-other-zero interpolation class has
\[
 \inf I_t(f)=\frac2{N_t(\delta)-N_t(0)}(1+o(1)).
                                                        \tag{EHL1}
\]
There are actual compact tests $\widehat f_t$, with both
poles exactly zero, supported in radius
\[
 L_t=4\delta t+4\sqrt{t\log t}
                         +\tfrac12\log t+B_0+2,
\]
for which
\[
 \begin{split}
 I_t(\widehat f_t)&=\frac2{N_t(\delta)-N_t(0)}(1+o(1)),\\
 q_A(\widehat f_t)=q_B(\widehat f_t)&=-2m+O(e^{-8t}),\\
 E_\Gamma(\widehat f_t)&\le
      K e^{-8\delta^2t+8\delta\sqrt{t\log t}}+K e^{-16t},\\
 \Pi(\widehat f_t)&=2m+o(1).
 \end{split}                                             \tag{EHL2}
\]
Here $\Pi=q_{\rm arch}-q_B$ is the complete prime-power
form and
$E_\Gamma(f)=\int[a(u)+C_a+1]|F_f(u)|^2\dd u/(2\pi)$,
with $C_a$ fixed so that the multiplier is positive and
dominates $|a|$.
In particular
$d_{t,1/2}\ge m[N_t(\delta)-N_t(0)](1+o(1))$.
Leading optimality is asserted for this interpolation
class, rather than for the full global negative depth.
\end{theorem}
\begin{proof}
Use the fixed interpolant $w$ in (EHC6) and put
$a=4\delta$. The convergent Gaussian superposition
\[
 f_t(x)=\int_\R G_t(v-at)e^{-\delta v}
       [e^{i\gamma v}w(x-v)+e^{-i\gamma v}w(x+v)]\dd v
                                                        \tag{EHL3}
\]
has the exact transform
\[
 F_{f_t}(z)=F_w(z)e^{-4\delta^2t}
       [e^{-t(z-\gamma+i\delta)^2}
                        +e^{-t(z-\gamma-i\delta)^2}].
                                                        \tag{EHL4}
\]
The theta envelope permits every strip transform and pole
evaluation to pass through the Gaussian integral.
The selected values are $\pm(1+e^{-4\delta^2t})$.
Every other original zero value and both pole values are
exactly zero. Thus
$q_A(f_t)=q_B(f_t)=-2m(1+e^{-4\delta^2t})^2$.

Define $h_\pm(y)=w(y)e^{-i\gamma y\pm\delta y}$.
Their integrals are $1,-1$. The physical lobes of (EHL3) are
\[
 f_{t,\pm}(x)=e^{i\gamma x\mp\delta x}
                         (G_t*h_\pm)(x\mp at).
\]
Let $A_t(x)$ denote the bracket after the Gaussian in (EHC4).
At $x=at+u\sqrt t$ put
$Q_{t,+}(u)=\int h_+(y)e^{uy/(2\sqrt t)-y^2/(4t)}\dd y$.
Direct substitution, including the moving center, gives
\[
 I_t(f_{t,+})=
 \frac1{4\pi\sigma_t\sqrt t\,e^{4\delta^2t}}
 \int_\R\frac{e^{-u^2/4}|Q_{t,+}(u)|^2}
                          {A_t(at+u\sqrt t)}\dd u .       \tag{EHL5}
\]
For fixed $u$, $Q_{t,+}(u)\to1$ and
$A_t(at+u\sqrt t)\to M(2\delta-1/2)$.
The theta envelope implies $|h_\pm|\le C_\epsilon G_\epsilon$
for every fixed $\epsilon>0$. Taking $\epsilon=1/16$ gives
$|Q_{t,\pm}(u)|\le C e^{u^2/(64t)}$.
The actual lower bound on $M_t$ in the proof of (EHC5)
also gives
$A_t(\pm at+u\sqrt t)^{-1}\le C_\delta e^{|u|/(2\sqrt t)}$.
The integrand is therefore dominated by
$C_\delta e^{-u^2/8}$ for all sufficiently large $t$.
Dominated convergence and (EHC1) yield
$I_t(f_{t,+})\sim1/N_t(\delta)$.
The negative lobe has $Q_{t,-}\to-1$ and the same limit,
since $A_t$ is even.

The cross term is negligible, with a genuine global bound:
\[
 |f_{t,+}(x)f_{t,-}(x)|
 \le\frac K{t+\epsilon}
      \exp\!\left[-\frac{x^2+a^2t^2}{2(t+\epsilon)}\right].
\]
Multiplying by (EHC5) and integrating shows
that its weighted integral is
$O(\sigma_t^{-1}t^{-1/2}e^{-8\delta^2t})
 =o(N_t(\delta)^{-1})$.
Hence $I_t(f_t)\sim2/N_t(\delta)$.
Dividing by $1+e^{-4\delta^2t}$ achieves the exact selected
values. Since $N_t(0)/N_t(\delta)\to0$, the Riesz lower
bound (EHC2) proves (EHL1).

For completeness these fast tests are in the actual closed
ground domain. For each fixed $t$, physical cutoffs converge
in $I_t$ by the preceding Gaussian majorants and in
$H^1_\rho$ for any fixed $\rho>1/2$.
On this latter space the gamma form is continuous and
the complete prime form is continuous, since a shift by
$\log n$ has overlap at most a fixed weighted norm times
$n^{-\rho}$ and
$\sum\Lambda(n)n^{-1/2-\rho}<\infty$.
The compact identity $D_p(f/p)=I_t(f)+q_B(f)$ makes
these cutoffs Cauchy in the positive energy norm.
Thus its closure retains the identity for $f_t/p$.
No generic physical-domain assertion is required.

We next construct the stated compact tests, retaining
the sharper constant. Restrict (EHL3) to
\[
 V_t=[at-4\sqrt{t\log t},\,at+4\sqrt{t\log t}].
\]
Its selected scalar is exactly
\[
 S_t^V=\int_{V_t}G_t(v-at)(1+e^{-2\delta v})\dd v
                    =1+o(1),\qquad S_t^V\ge1/2.          \tag{EHL6}
\]
All other zero and pole values are still exactly zero.
The omitted first Gaussian mass is $O(t^{-4})$; the second
term is at most
$e^{-8\delta^2t+8\delta\sqrt{t\log t}}$.
In (EHL5) the truncated-error integrand has the additional
indicator
$\boldsymbol1_{\{|u-y/\sqrt t|>4\sqrt{\log t}\}}$.
It tends to zero at every fixed $u,y$, while the same
Gaussian dominator remains valid. Applying dominated
convergence to both lobes gives
$I_t(f_t-f_t^V)=o(N_t(\delta)^{-1})$.

Take the physical cutoff at radius $L_t$ and
$B_t=\tfrac12\log t+B_0$, with $B_0$ chosen as in (EHC10).
On either truncated lobe tail, $|x\mp v|\ge B_t+1$.
Weighted Minkowski, (EHC5), and the exact square
$v^2/(4t)-2\delta v=-4\delta^2t+(v-at)^2/(4t)$ give
\[
 I_t((1-\chi_{L_t})f_t^V)
 \le K\sigma_t^{-1}e^{-4\delta^2t}t^4
                              e^{-c_*e^{2B_0}t}
 =o(N_t(\delta)^{-1}).                                \tag{EHL7}
\]
Four physical integrations by parts, using coefficient
mass $\int_{V_t}G_t\le1$, give on the entire closed strip
\[
 |F_{\rm error}(\gamma'+i\delta')|
 \le K(1+|\gamma'|)^{-4}
       e^{(1/2-\delta)(at+4\sqrt{t\log t})}
                              e^{-c_*e^{2B_0}t}
 =O(e^{-8t}(1+|\gamma'|)^{-4}).                         \tag{EHL8}
\]
Here $c_*e^{2B_0}>10$, and
$4\delta(1/2-\delta)\le1/4$.
The original counting bound makes the multiplicity-weighted
square sum finite. Fixed compact correctors repair both
pole values exactly with squared cost
$O(\sigma_t^{-1}e^{-16t})=o(N_t(\delta)^{-1})$.
Normalize this corrected compact vector by the exact $S_t^V$.
The full zero formula and its Cauchy--Schwarz bound prove
the first two assertions of (EHL2), including every height
and multiplicity. Compact tests have the written small
other-zero error; the fast tests preceding them have exact
other-zero annihilation.

Finally, on the whole real Fourier line,
\[
 |F_{f_t^V}(u)|
 \le2e^{-\delta(at-4\sqrt{t\log t})}|F_w(u)|.              \tag{EHL9}
\]
The positive gamma energy of $w$ is finite.
Equation (EHL8) gives gamma energy error $O(e^{-16t})$,
since $a(u)=O(\log(2+|u|))$.
This proves the gamma bound in (EHL2).
The exact pole annihilation gives
$q_A=q_B=q_{\rm arch}-\Pi$; therefore the sum of
all prime-power correlations is $2m+o(1)$.
Dividing the actual negative form by the attained norm
proves the stated lower bound on $d_{t,1/2}$.
\end{proof}

\subsection{The principal prime and first gamma terms cancel exactly}
For compact smooth physical $f$ put
$C_f(r)=\int\overline{f(x)}f(x+r)\dd x$,
$E(r)=\psi(e^r)-e^r$, and
$k_1(r)=e^{-5r/2}/(1-e^{-2r})$.
The full arithmetic form has the exact identities
\[
 \begin{split}
 q_A(f)={}&[a(0)+4]\|f\|_2^2+
       2\int_0^\infty k_1(r)[\|f\|_2^2-\Re C_f(r)]\dd r
       -2\int_0^\infty e^{-r/2}\Re C_f(r)\dd E(r),\\
 ={}&[a(0)+2]\|f\|_2^2+
       2\int_0^\infty k_1(r)[\|f\|_2^2-\Re C_f(r)]\dd r\\
 &\hspace{15mm}
       +2\int_0^\infty E(r)e^{-r/2}
                    [\Re C_f'(r)-\tfrac12\Re C_f(r)]\dd r .
 \end{split}                                            \tag{EPA1}
\]
Indeed the two physical poles give
$4\int_0^\infty\cosh(r/2)\Re C_f(r)\dd r$.
The principal part of the complete prime Stieltjes measure
gives $2\int e^{r/2}\Re C_f$. Their difference cancels
the correlation part of the whole first gamma atom.
That atom also contributes $4\|f\|_2^2$.
Finally $E(0)=-1$ gives the boundary shift $4$ to $2$
under Stieltjes integration by parts.
The remaining $1/r$ singularity is cancelled by
$\|f\|_2^2-\Re C_f(r)=O(r^2)$. All correlations have
compact support and every prime power remains present.

\begin{proposition}[Absolute prime error is not a gamma-relative budget]
\label{prop:absolute-prime-error-obstruction}
Unconditionally,
\[
 \int_0^\infty e^{-r/2}|E(r)|\dd r=\infty .
\]
For
$\mathcal B_E(f)=2\int e^{-r/2}|E(r)|
                   |\Re C_f'(r)-\tfrac12\Re C_f(r)|\dd r$
there are compact smooth tests with both poles exactly zero
such that
\[
 \frac{\mathcal B_E(f_R)}{\|f_R\|_2^2}\longrightarrow\infty,
 \qquad
 \frac{q_{\rm arch}(f_R)}{\|f_R\|_2^2}\longrightarrow a(0).
                                                        \tag{EPA2}
\]
In particular no uniform $O(\sigma_t)$ endpoint budget
of this absolute form holds, even if RH is true.
\end{proposition}
\begin{proof}
If $\int_1^\infty|\psi(x)-x|x^{-3/2}\dd x$ were finite,
Euler differentiation and Abel integration would give
\[
 s\int_1^\infty(\psi(x)-x)x^{-s-1}\dd x
       =-\frac{\zeta'}{\zeta}(s)-\frac{s}{s-1}
\]
with a holomorphic left side on $\Re s>1/2$, bounded on
approach from the right to each fixed critical ordinate.
There is unconditionally a critical zero, for example
by \cite[Theorem~1, p.~4]{Conrey}.
At that zero the right side has residue $-m\ne0$.
Meromorphic continuation and dominated convergence give
a contradiction. This uses no simplicity or RH premise.

Take any nonzero real $\eta\in C_c^\infty$ and
$f_R=(-\partial_x^2+1/4)\eta(x/R)$.
Integration by parts makes both poles zero exactly.
Its norm is asymptotic to $R\|\eta\|_2^2/16$.
For every fixed $r$, $\Re C_{f_R}(r)/\|f_R\|_2^2\to1$,
and uniformly in $r$,
$|C_{f_R}'(r)|/\|f_R\|_2^2\le\|f_R'\|_2/\|f_R\|_2=O(1/R)$.
Fatou proves the first part of (EPA2).
Fourier scaling, the Schwartz transform of $\eta$, and
$|a(u)|\le C(1+u^2)$ prove its second part.
Thus no fixed positive gamma-plus-$L^2$ budget bounds
$\mathcal B_E$. If a uniform endpoint
$\mathcal B_E(f)\le C\sigma_t\int c|f|^2/p_{t,1/2}$
held, first fixing $R$ and then letting $t\to\infty$
in (SHD4) would give
$\mathcal B_E(f_R)\le2C\|f_R\|_2^2$, another contradiction.
The signed term in (EPA1) has not been bounded in this argument.
\end{proof}

\section{A canonical Gaussian family exhausts the full negative index}
The preceding individual detectors need no most-right zero.
We now give a predetermined finite family whose full negative
index eventually detects each finite collection of pairs.
The family is not nested; its cofinal transfer is proved directly.

\begin{theorem}[Canonical Gaussian index exhaustion]
\label{thm:canonical-gaussian-index}
For each integer $n\ge1$ put
\[
 \mathcal V_n=\operatorname{span}_{\C}
 \left\{(-i\partial_x)^j(-\partial_x^2+1/4)
       [p_n(x)e^{ikx/n^2}]:
          0\le j\le n,\quad k\in\Z,\quad |k|\le n^3\right\}.
                                                        \tag{CDX1}
\]
Every vector is smooth, Gaussian decreasing, and has both
physical pole moments zero. Its complete gamma and prime-power
Gram is well defined. With $K_+$ as in the exact index theorem,
\[
 \lim_{n\to\infty}n_-(q_A|_{\mathcal V_n})=2K_+.            \tag{CDX2}
\]
If $K_+$ is finite, equality holds for every sufficiently
large $n$. If it is infinite, the left side tends to infinity.
Consequently positivity of the full Gram along any unbounded
subsequence would imply RH; this positivity is not proved here.
\end{theorem}
\begin{proof}
The factor $-\partial_x^2+1/4$ makes the transform vanish
at both poles. For fixed $n$, every derivative has Gaussian
physical tails. The gamma integral and the complete
prime-power correlation form converge absolutely.
They agree with the full zero formula by physical cutoff,
the strip counting bound, and fixed compact pole repair.

Choose any finite collection of distinct off-line mirror
pairs $z_{\ell,\pm}=\gamma_\ell\pm i\delta_\ell$.
Pairs at both positive and negative heights can be selected.
For each $\ell$ choose a fixed real polynomial $P_\ell$,
divisible by $z^2+1/4$, which vanishes on every other selected
pair and every other original zero with
$|\Re z-\gamma_\ell|\le1$.
These constraints are a finite conjugation-invariant set;
$P_\ell(z_{\ell,+})\ne0$.
Orders of zeros need not be repeated in the polynomial,
since the full form uses their values with multiplicity.

For $T$ real, the exact Fourier transform of the trial vector is
\[
 F_{\ell,n,T}(z)=
 P_\ell(z)e^{-n[(z-T)^2+1/4]}\Xi(z-T).                    \tag{CDX3}
\]
This is a polynomial derivative of $p_n e^{iTx}$.
Write $\theta_\ell=\arg P_\ell(z_{\ell,+})$ and choose a
fixed real $u_\ell$ with
$\theta_\ell+2\delta_\ell u_\ell=\pi/2$ modulo $2\pi$.
The nearest grid point
$T_{\ell,n}=k_{\ell,n}/n^2$ to $\gamma_\ell+u_\ell/n$
has error at most $1/(2n^2)$.
For all sufficiently large $n$ its index obeys
$|k_{\ell,n}|\le n^3$, and the fixed degree of
$P_\ell/(z^2+1/4)$ is at most $n$.
Thus these vectors belong to $\mathcal V_n$.

Multiply the $\ell$th vector by the positive scalar
\[
 \frac{e^{n(1/4-\delta_\ell^2)}}
                {|P_\ell(z_{\ell,+})|\Xi(i\delta_\ell)}.
                                                        \tag{CDX4}
\]
The denominator is positive: $\Xi(i\delta)=
\int\Phi(x)\cosh(\delta x)\dd x>0$.
At the selected two nodes (CDX3) then tends to $i,-i$.
The grid phase error is $O(1/n)$.
Every other locally constrained value is exactly zero.
For the remaining nodes $|\Re z-\gamma_\ell|>1$,
$|\Im z|<1/2$.
After (CDX4), their Gaussian exponent is at most
$-n[\,|\Re z-T_{\ell,n}|^2+\delta_\ell^2-1/4\,]$.
For large $n$ this has a fixed strictly negative gap.
The fixed polynomial, the strip bound for $\Xi$,
and the multiplicity counting bound dominate the entire
weighted square sum by a convergent Gaussian tail.
It therefore tends to zero. This is an all-height
weighted $\ell^2$ assertion, rather than a finite-zero
approximation.

The finite target vectors have disjoint pair supports.
Polarizing the full zero formula, their Gram is diagonal
with entries $-2m_\ell$. The actual Gram converges to it
in matrix norm, so is strictly negative definite for
every sufficiently large $n$. We obtain as many negative
dimensions as the chosen finite set of mirror pairs.
The potentially large rescaling in (CDX4) does not affect
negative definiteness; no uniform gap estimate is inferred.

For the reverse inequality, any negative finite-dimensional
subspace of $\mathcal V_n$ can be approximated by a common
physical compact cutoff and the two fixed pole correctors.
Gaussian tails give convergence of its entire finite
gamma, all-prime, and zero Grams. A sufficiently large
cutoff retains its strict negative dimension.
The already proved exact compact index bounds it by $2K_+$.
Selecting all pairs when this number is finite proves
eventual equality. Selecting arbitrarily large finite sets
when it is infinite proves the extended limit (CDX2).
No monotonicity of $\mathcal V_n$ is used.
\end{proof}

\section{A positive rest regularizer and uniform global norm decay}
\label{sec:rest-regularizer-global}

We use the Fourier convention
\[
 \widehat f(z)=\int_{\mathbb R}f(x)e^{-izx}\,dx,
 \qquad f(x)=\frac1{2\pi}\int_{\mathbb R}\widehat f(u)e^{iux}\,du.
 \tag{RGC1}
\]
Write \(c(x)=\cosh(x/2)\), \(s(x)=\sinh(x/2)\), and
\(\kappa=s/c\). Recall the already proved theta-ground facts:
\(\Phi\) is positive, even, and smooth;
\(\int c\Phi=1/2\); \(\Phi,\Phi''\) are bounded;
\(\sup_x e^{|x|/2}\Phi(x)<\infty\); and
\[
 B\Phi=-c,\qquad A=B+P,\qquad B=\Gamma-\Pi.
 \tag{RGC2}
\]
Here the exact arithmetic and pole kernels are
\[
 \Pi=\sum_{n\ge2}\frac{\Lambda(n)}{\sqrt n}
       \bigl(\delta_{\log n}+\delta_{-\log n}\bigr),
 \qquad P(x-y)=2c(x-y).
 \tag{RGC3}
\]
Every prime power is included. The archimedean multiplier and its
second-difference formula are
\[
 a(z)=\frac12\left[
 \psi_{\rm dig}\!\left(\frac14+\frac{iz}{2}\right)+
 \psi_{\rm dig}\!\left(\frac14-\frac{iz}{2}\right)\right]-\log\pi,
 \quad
 \Gamma f=a(0)f+\int_0^\infty
 \frac{e^{-r/2}}{1-e^{-2r}}
 \bigl(2f-f(\cdot-r)-f(\cdot+r)\bigr)\,dr.
 \tag{RGC4}
\]
All physical identities below initially hold on smooth compactly
supported tests. Their indicated bounded weighted extensions are
then proved, rather than assumed.

\begin{lemma}[Explicit positive rest construction]
\label{lem:rest-regularizer-construction}
Fix \(C>\max(0,-a(0))\), put \(q=C+a(0)>0\),
\(x_n=n+1/4\), and \(\lambda_n=4x_n^2\). Define
\[
 \varphi_C(\zeta)=q+
   \sum_{n\ge0}\frac1{x_n}\frac{\zeta}{\zeta+\lambda_n}.
 \tag{RGC5}
\]
There are simple numbers
\[
 0<\mu_0<\lambda_0<\mu_1<\lambda_1<\mu_2<\lambda_2<\cdots,
 \qquad A_j=\varphi_C'(-\mu_j)^{-1}>0,
 \tag{RGC6}
\]
such that, away from its poles,
\[
 \varphi_C(\zeta)^{-1}
   =\sum_{j\ge0}\frac{A_j}{\zeta+\mu_j},
 \qquad \sum_{j\ge0}\frac{A_j}{\mu_j}=\frac1q.
 \tag{RGC7}
\]
In fact \(\mu_1>1\). The functions
\[
 H_C(z)=\sum_{j\ge1}\frac{A_j}{z^2+\mu_j},\qquad
 G_C(r)=\sum_{j\ge1}\frac{A_j}{2\sqrt{\mu_j}}
                   e^{-\sqrt{\mu_j}|r|}
 \tag{RGC8}
\]
are respectively the Fourier transform and kernel of a positive
even \(L^1\) convolution operator \(R_C\). Its exact pool moment is
\[
 r_C=\int_{\mathbb R}c(r)G_C(r)\,dr
     =H_C(i/2)=\frac{A_0}{1/4-\mu_0}>0.
 \tag{RGC9}
\]
With
\[
 B_n=\frac{A_0\lambda_n}{x_n(\lambda_n-\mu_0)},\qquad
 U_C^{\rm hi}(r)
   =C G_C(r)+\sum_{n\ge1}\frac{B_n}{2\sqrt{\lambda_n}}
                    e^{-\sqrt{\lambda_n}|r|},
 \tag{RGC10}
\]
the exact multiplier identity is
\[
 1-a(z)H_C(z)=
 C H_C(z)+\sum_{n\ge0}\frac{B_n}{z^2+\lambda_n}.
 \tag{RGC11}
\]
In particular \(B_0=r_C\), and the inverse kernel on the right is
\(U_C^{\rm hi}+r_Ce^{-|r|/2}\).
The number
\[
 H_U=\int_{\mathbb R}e^{|r|/2}U_C^{\rm hi}(r)\,dr
 \tag{RGC12}
\]
is finite. For \(b_1=\sqrt{\mu_1}>1\), \(G_C\) has every
exponential \(L^1\) moment of order less than \(b_1\).
\end{lemma}

\begin{proof}
The digamma partial-fraction formula
\href{https://dlmf.nist.gov/5.7\#E6}{NIST DLMF 5.7.6},
valid off the nonpositive integers, gives
\(a(z)+C=\varphi_C(z^2)\). The displayed series converges locally
uniformly off the negative poles: its tail is \(O(n^{-3})\)
on each such compact set.

First truncate at \(n=N\). On every pole-free real interval,
\[
 \varphi_{C,N}'(\zeta)=
 \sum_{n=0}^N\frac{\lambda_n/x_n}{(\zeta+\lambda_n)^2}>0.
 \tag{RGC13}
\]
The one-sided pole signs give precisely \(N+1\) simple negative
zeros with the stated interlacing. Partial fractions therefore give
\[
 \varphi_{C,N}(\zeta)^{-1}
   =b_N+\sum_{j=0}^N\frac{A_{N,j}}{\zeta+\mu_{N,j}},
 \quad
 b_N=\left(q+\sum_{n=0}^Nx_n^{-1}\right)^{-1},
 \quad A_{N,j}=\varphi_{C,N}'(-\mu_{N,j})^{-1}>0.
 \tag{RGC14}
\]
At zero,
\(b_N+\sum_jA_{N,j}/\mu_{N,j}=1/q\), and \(b_N\to0\).
There is a uniform positive lower bound \(\mu_*\) on all
\(\mu_{N,0}\): for \(0<\mu\le\lambda_0/2\),
\[
 \varphi_{C,N}(-\mu)\ge
 q-2\mu\sum_{n\ge0}(x_n\lambda_n)^{-1}.
 \tag{RGC15}
\]
Choose \(\mu_*<\min(\lambda_0/2,
 q/[4\sum_n(x_n\lambda_n)^{-1}])\).

The finite measures
\(\eta_N=\sum_j(A_{N,j}/\mu_{N,j})\delta_{\mu_{N,j}}\)
have uniformly bounded mass on the compactification
\([\mu_*,\infty]\). Pass to a weakly convergent subsequence.
For real \(\zeta>0\), the continuous function
\(\mu/(\zeta+\mu)\), with value \(1\) at infinity, gives
\[
 \varphi_C(\zeta)^{-1}
   =\int\frac{\mu}{\zeta+\mu}\,d\eta(\mu).
 \tag{RGC16}
\]
Every fixed finite partial sum shows
\(\varphi_C(\zeta)\to\infty\) as \(\zeta\to\infty\);
thus dominated convergence rules out mass at infinity.
Letting \(\zeta\downarrow0\) gives total mass \(1/q\).
On any fixed bounded interval of \(\mu\), interlacing bounds the
number of atoms uniformly. Local uniform convergence of
\(\varphi_{C,N}\) and its derivative makes each such root and its
residue converge to the unique root of \(\varphi_C\) in the same
interval. Consequently the limiting measure is exactly
\(\sum_j(A_j/\mu_j)\delta_{\mu_j}\), proving the series.
For nonreal \(\zeta\), the imaginary part of
\(\varphi_C(\zeta)\) has the strict sign of
\(\operatorname{Im}\zeta\); hence no additional zeros are omitted.
Analytic continuation supplies the asserted identity.

At \(\zeta=-1\), which lies between the first two poles,
\[
 \varphi_C(-1)=q+\frac{16}{3}
       -\sum_{n\ge1}\frac1{x_n(\lambda_n-1)}>0.
 \tag{RGC17}
\]
Indeed the sum is less than
\(\frac13\sum_{n\ge1}n^{-3}<1/2\). Monotonicity on that interval
therefore gives \(\mu_1>1\).

The inverse transform of \((u^2+\mu)^{-1}\) is
\((2\sqrt\mu)^{-1}e^{-\sqrt\mu|r|}\).
Tonelli proves the \(L^1\) statements. Removing the \(j=0\)
atom defines the displayed \(H_C,G_C\). Because
\(\varphi_C^{-1}(-1/4)=0\) at the first digamma pole,
\[
 H_C(i/2)=-\frac{A_0}{\mu_0-1/4}=r_C.
 \tag{RGC18}
\]
The moment identities converge absolutely since \(\mu_j\ge\mu_1>1\)
for \(j\ge1\) and \(\sum A_j/\mu_j<\infty\).

Finally \(\varphi_C(-\mu_0)=0\) yields
\[
 \frac{\varphi_C(\zeta)}{\zeta+\mu_0}
   =\sum_{n\ge0}
      \frac{\lambda_n}{x_n(\lambda_n-\mu_0)(\zeta+\lambda_n)}.
 \tag{RGC19}
\]
Insert \(H_C=\varphi_C(z^2)^{-1}-A_0/(z^2+\mu_0)\)
to obtain the exact \(U\)-identity. At \(n=0\), it gives \(B_0=r_C\).
The higher \(\lambda_n\)-rates are at least \(5/2\);
\(\sum B_n/\lambda_n<\infty\). For \(G_C\), all rates exceed \(b_1>1\).
The ratio of the exponentially weighted to unweighted mass of
each exponential of rate \(b\) is \(b/(b-1/2)\), uniformly bounded
over these rates. This proves \(H_U<\infty\) and the remaining
exponential moments. Notice that the low term
\(r_Ce^{-|r|/2}\) is not included in \(H_U\).
\end{proof}

Define the signed and positive kernels
\[
 S_C=U_C^{\rm hi}+G_C*\Pi-r_Ce^{|r|/2},
 \qquad
 S_C^+=S_C+2r_Cc
       =U_C^{\rm hi}+G_C*\Pi+r_Ce^{-|r|/2}.
 \tag{RGC20}
\]
The all-prime sum defining \(G_C*\Pi\) is locally integrable and
positive. On compact intervals its distant terms are bounded by
a constant times \(\Lambda(n)n^{-b-1/2}\) for any
\(1/2<b<b_1\); the finitely many local singularities are integrable.
Indeed the positive series and \(\sum A_j/\mu_j<\infty\) give
\(G_C(r)\le K_be^{-b|r|}\) for \(|r|\ge1\): after multiplying
by \(e^{b|r|}\), each term is bounded by \(A_j/\mu_j\) times
\(\sup_{\mu\ge\mu_1}(\sqrt\mu/2)e^{-(\sqrt\mu-b)}<\infty\).
The multiplier identity and the exact pool moment now give
\[
 R_CB=I-S_C^+,\qquad R_CA=I-S_C.
 \tag{RGC21}
\]
These identities retain both pools on every test, not only on
pool-null tests.

\begin{theorem}[Uniform norm decay on the entire weighted space]
\label{thm:rest-shifted-global-norm}
For \(t\ge0\), let
\[
 p_t=e^{-t/4}G_t*\Phi,\qquad
 G_t(x)=(4\pi t)^{-1/2}e^{-x^2/(4t)}\quad(t>0),\qquad p_0=\Phi,
 \tag{RGC22}
\]
and for \(v\ge0\) set
\[
 \phi_{v,t}(x)=
 \frac{p_t(x-v)+p_t(x+v)}{2\cosh(v/2)},\qquad
 d\mu_{v,t}=c\phi_{v,t}\,dx,\qquad
 \omega_{v,t}=\sqrt{\phi_{v,t}/c}.
 \tag{RGC23}
\]
Then \(B\phi_{v,t}=-c\), \(\mu_{v,t}(\mathbb R)=1/2\).
The physical form \(q_{R_CA}(\phi_{v,t}g)\) extends to a bounded
self-adjoint form on all of \(L^2(\mu_{v,t})\). Under the unitary
map \(g\mapsto\sqrt{c\phi_{v,t}}\,g\), its operator is
\[
 L_{C;v,t}=M_{\omega_{v,t}^2}
     -\omega_{v,t} S_C(D)\omega_{v,t}.
 \tag{RGC24}
\]
For this fixed \(C\),
\[
 \sup_{t\ge0}\|L_{C;v,t}\|\longrightarrow0
                  \qquad(v\to\infty).
 \tag{RGC25}
\]
Let
\[
 D_{C;v,t}^+(g)=\frac12\iint
 S_C^+(x-y)\phi_{v,t}(x)\phi_{v,t}(y)
                  |g(x)-g(y)|^2\,dx\,dy .
 \tag{RGC26}
\]
This is finite on the entire weighted space. Its normalized
ordinary mean-zero gap
\[
 \beta_{C;v,t}=
 \inf_{\substack{g\ne0\\ \int g\,d\mu_{v,t}=0}}
 \frac{D_{C;v,t}^+(g)}
      {r_C\|g\|_{L^2(\mu_{v,t})}^2}
 \tag{RGC27}
\]
satisfies
\(\sup_{t\ge0}|\beta_{C;v,t}-1|\to0\).
\end{theorem}

\begin{proof}
Put \(q_0(x)=e^{-|x|/2}\) and
\(A_\Phi=\sup_x e^{|x|/2}\Phi(x)\).
Distributionally \((\partial_x^2-1/4)q_0=-\delta_0\);
hence the exact positive heat identity is
\[
 e^{-t/4}G_t*q_0
   =q_0-\int_0^t e^{-s/4}G_s\,ds\le q_0 .
 \tag{RGC28}
\]
Thus \(p_t\le A_\Phi q_0\) for every \(t\ge0\).
For \(|x|\le v\),
\((q_0(x-v)+q_0(x+v))/(2c(v)c(x))=e^{-v/2}/c(v)\);
for \(|x|\ge v\) the quotient equals \(e^{-|x|/2}/c(x)\)
and is no larger. Therefore, uniformly in all \(t\),
\[
 M_{v,t}:=\sup_x\frac{\phi_{v,t}(x)}{c(x)}
       \le \frac{A_\Phi e^{-v/2}}{c(v)}
       \le 2A_\Phi e^{-v}.
 \tag{RGC29}
\]
Convolution and translation preserve the exact \(B\)-ground
identity, while \(c(x-v)+c(x+v)=2c(x)c(v)\).
The normalization follows from evenness and
\(\int c\Phi=1/2\).

There is also a uniform all-prime degree bound. The exact
archimedean formula gives
\[
 \|\Gamma\Phi\|_\infty\le M_\Gamma:=
 |a(0)|\|\Phi\|_\infty
 +\|\Phi''\|_\infty\int_0^1r^2k(r)\,dr
 +4\|\Phi\|_\infty\int_1^\infty k(r)\,dr<\infty,
 \quad k(r)=\frac{e^{-r/2}}{1-e^{-2r}}.
 \tag{RGC30}
\]
Its positive heat convolution and the two translations imply
\(\|\Gamma\phi_{v,t}\|_\infty\le M_\Gamma\).
The ground equation, with all prime powers retained, implies
\[
 0\le\Pi\phi_{v,t}=c+\Gamma\phi_{v,t}
        \le K_\Pi c,\qquad K_\Pi=1+M_\Gamma .
 \tag{RGC31}
\]
All these constants are independent of \(v,t\).

For a positive even kernel \(Q\), the row and column Schur test
with weight \(w=\sqrt{c\phi_{v,t}}\) gives
\[
 \|\omega_{v,t}Q(D)\omega_{v,t}\|
      \le \sup_x\frac{(Q*\phi_{v,t})(x)}{c(x)}.
 \tag{RGC32}
\]
For signed kernels use \(|Q|\). In particular a kernel with
finite \(1/2\)-exponentially weighted \(L^1\) mass costs at most
\(M_{v,t}\int e^{|r|/2}|Q(r)|\,dr\), since
\(c(x-r)\le c(x)e^{|r|/2}\).
Also \(2c*\phi_{v,t}=c\), so a kernel bounded in absolute
value by \(2\epsilon c\) costs at most \(\epsilon\).

The exact row of \(S_C^+\) is
\[
 S_C^+*\phi_{v,t}=\phi_{v,t}+r_Cc .
 \tag{RGC33}
\]
Consequently its weighted norm is at most \(M_{v,t}+r_C\).
This proves that \(D^+\) is finite and bounded on the entire
weighted space. Polarization of its positive difference integral,
followed by the exact two-pool formula, gives
\[
 q_{R_CA}(\phi_{v,t}g)=D_{C;v,t}^+(g)
       -r_C\|g\|_{\mu_{v,t}}^2
       +2r_C\left|\int g\,d\mu_{v,t}\right|^2
       -2r_C\left|\int\kappa g\,d\mu_{v,t}\right|^2 .
 \tag{RGC34}
\]
Thus the asserted bounded operator really extends the physical
form and keeps both pools.

It remains to prove decay without estimating singular prime
spikes pointwise. Use the finite positive Stieltjes head
\[
 G_J(r)=\sum_{j=1}^J
        \frac{A_j}{2\sqrt{\mu_j}}e^{-\sqrt{\mu_j}|r|},
 \quad r_J=\int cG_J,\quad
 r_{{\rm tail},J}=r_C-r_J\longrightarrow0,
 \quad W_J=G_J*\Pi-r_Je^{|r|/2}.
 \tag{RGC35}
\]
We claim \(W_J(r)=o(e^{|r|/2})\) for each fixed \(J\).
This uses only the unconditional prime number theorem
\(\Psi(X):=\sum_{n\le X}\Lambda(n)=X+o(X)\);
see \href{https://kskedlaya.org/ant/chap-pnt.html}
{Kedlaya, Chapter 1, Sections 1.3--1.4}.
Its assertion for \(\vartheta(X)=\sum_{p\le X}\log p\) gives
the displayed all-prime-power version: the same chapter's
Chebyshev bound gives
\(\Psi(X)-\vartheta(X)=\sum_{k\ge2}\vartheta(X^{1/k})
=O(\sqrt X\log X)=o(X)\).
For \(r>0\), set
\(a_r(X)=X^{-1/2}G_J(r-\log X)\).
Stieltjes integration by parts with \(E(X)=\Psi(X)-X\)
gives an \(O_J(e^{-b r})\) boundary, where
\(1/2<b<\min_{j\le J}\sqrt{\mu_j}\), and error
\(-\int_1^\infty E(X)a_r'(X)\,dX\).
On \(X\ge X_0\) with \(|E(X)|\le\epsilon X\), its magnitude is
at most
\[
 \epsilon e^{r/2}
 \int_{\mathbb R}e^{-s/2}
             |G_J'(s)+G_J(s)/2|\,ds .
 \tag{RGC36}
\]
The integral is finite. The remaining bounded \(X\)-interval
costs \(O_{J,X_0}(e^{-b r})\). The main integral satisfies
\[
 e^{-r/2}\int_1^\infty a_r(X)\,dX
   =\int_{-\infty}^r e^{-s/2}G_J(s)\,ds\longrightarrow r_J .
 \tag{RGC37}
\]
The negative prime translations cost \(O_J(e^{-b r})\)
by \(\sum_n\Lambda(n)n^{-b-1/2}<\infty\).
The finite exponential cusps are continuous, and create no
omitted jump in this integration by parts. First let \(r\to\infty\)
and then \(\epsilon\to0\). Evenness handles negative \(r\).
This proves the claim with every prime power.

Given \(J\) and \(\epsilon>0\), choose \(R\) such that
\(|W_J(r)|\le\epsilon e^{|r|/2}\) for \(|r|>R\), and put
\[
 B_{J,R}=\int_{|r|\le R}e^{|r|/2}|W_J(r)|\,dr<\infty .
 \tag{RGC38}
\]
The compact part of \(W_J\) costs \(M_{v,t}B_{J,R}\);
the exterior part costs at most \(\epsilon\).
The \(U_C^{\rm hi}\) part costs \(M_{v,t}H_U\).
For the remaining positive Stieltjes tail, Tonelli and the
uniform genuine prime degree give
\[
 ((G_C-G_J)*\Pi)*\phi_{v,t}
      \le K_\Pi (G_C-G_J)*c
      =K_\Pi r_{{\rm tail},J}c .
 \tag{RGC39}
\]
Its norm therefore costs at most \(K_\Pi r_{{\rm tail},J}\).
The matching growing-pool term costs at most
\(r_{{\rm tail},J}\), because \(e^{|r|/2}\le2c(r)\).
We have proved the explicit whole-space budgets
\[
 \|\omega_{v,t}S_C(D)\omega_{v,t}\|
 \le 2A_\Phi e^{-v}(B_{J,R}+H_U)
       +\epsilon+(K_\Pi+1)r_{{\rm tail},J},
 \tag{RGC40}
\]
\[
 \|L_{C;v,t}\|
 \le 2A_\Phi e^{-v}(1+B_{J,R}+H_U)
       +\epsilon+(K_\Pi+1)r_{{\rm tail},J}.
 \tag{RGC41}
\]
For a desired error, first choose \(J\), then \(\epsilon,R\),
and finally \(v\). The final choice is independent of all
\(t\ge0\) and of every form vector. This proves the uniform
operator norm limit.

On the mean-zero subspace, the displayed exact energy identity
gives \(\beta_{C;v,t}\ge1-\|L_{C;v,t}\|/r_C\).
Choose a nonzero vector in the common kernel of the two
bounded moments \(\int g\,d\mu\) and \(\int\kappa g\,d\mu\);
the infinite-dimensional weighted space has such vectors.
Its quotient is at most \(1+\|L_{C;v,t}\|/r_C\).
The opposite bound follows, proving the uniform gap limit.
\end{proof}

This theorem is an unconditional global, all-vector estimate.
Its proof alone provides no exponential rate: its constants
\(B_{J,R}\), \(R\), and \(r_{{\rm tail},J}\) are chosen in the
stated order. No bound at the scale \(e^{-v+o(v)}\), and no
exclusion of an off-line zero, follows by exchanging these
choices with \(v\).

\subsection{A quantitative norm limit and its positive-tail budget}

\begin{theorem}[A quantitative, uniform all-vector norm limit]
\label{thm:quantitative-rest-norm}
Fix \(C>\max(0,-a(0))\). There are constants \(K_C,V_C>0\)
such that, for every \(v\ge V_C\),
\[
 \sup_{t\ge0}\|L_{C;v,t}\|\le K_Cv^{-1/2},\qquad
 \sup_{t\ge0}|\beta_{C;v,t}-1|\le K_C/(r_C\sqrt v). \tag{RQ1}
\]
The form retains every prime power and both pool terms of
Theorem~\ref{thm:rest-shifted-global-norm}.
\end{theorem}
\begin{proof}
Use the unconditional classical prime number theorem
\[
 |\Psi(X)-X|\le BXe^{-c_0\sqrt{\log X}}\quad(X\ge1),
 \qquad\Psi(X)=\sum_{n\le X}\Lambda(n),             \tag{RQ2}
\]
with absolute \(B,c_0>0\), enlarging \(B\) on a bounded interval.
This is \href{https://people.maths.ox.ac.uk/greenbj/papers/primenumbers.pdf}
{Ben Green, \emph{The Prime Number Theorem}, Theorem 9.2,
Section 9.2, printed page 55}. Its separate RH Section 9.3 is
not used. In particular (RQ2) retains all prime powers.

Set \(b_j=\sqrt{\mu_j}\), and use \(G_J,r_J,W_J\) from (RGC35).
Interlacing and the positive resolvent mass give
\[
 1<b_1\le b_j<2j+\tfrac12,\quad
 \sum_{j\le J}A_j/b_j\le K_CJ,\quad
 G_J(r)\le K_CJ e^{-b_1|r|}.                       \tag{RQ3}
\]
The finite cusps are continuous. Integration on both half-lines
gives the exact derivative moment
\[
 I_J:=\int e^{-s/2}|G'_J(s)+G_J(s)/2|\,ds
 =\sum_{j\le J}\frac{A_j}{b_j}
                   \frac{b_j^2+1/4}{b_j^2-1/4}\le K_CJ. \tag{RQ4}
\]
On the positive half-line its tail beyond \(R\) is at most
\(K_CJ e^{-(b_1+1/2)R}\).
With \(S=\lambda_J=4(J+1/4)^2\), interlacing gives \(\mu_j>S\)
for \(j>J\). Therefore
\[
 r_{{\rm tail},J}\le\frac43\sum_{j>J}\frac{A_j}{\mu_j}
       \le\frac83H_C(\sqrt S)\le\frac{K_C}{\log(2+J)}. \tag{RQ5}
\]
The last estimate follows from the real-axis digamma expansion
\(a(x)=\log(x/(2\pi))+O(x^{-1})\); it pays the actual singular tail.

For \(r\ge1\), repeat the Stieltjes integration by parts in
(RGC36)--(RGC37), splitting \(X\) at \(e^{r/2}\).
The upper part costs
\(Be^{r/2}e^{-c_0\sqrt{r/2}}I_J\).
On the lower part use \(|\Psi(X)-X|\le BX\) and the positive
tail of (RQ4). The missing main integral, lower boundary, and
negative prime translations are paid by (RQ3) and
\(\sum_n\Lambda(n)n^{-b_1-1/2}<\infty\). Hence, uniformly in \(J\),
\[
 |W_J(r)|e^{-|r|/2}
 \le K_CJ\{e^{-c_0\sqrt{|r|/2}}+e^{-\alpha_C|r|}\}
 \quad(|r|\ge1),\qquad \alpha_C=(b_1+1/2)/2>0.     \tag{RQ6}
\]
There is no omitted jump at \(X=e^r\). For \(|r|\le1\), split
the actual prime sum at \(n=e^2\) and use the same convergent
Dirichlet sum. Integration gives \(B_{J,R}\le K_CJe^R\).
Insert these bounds into the entire-space estimate (RGC41):
\[
 \sup_{t\ge0}\|L_{C;v,t}\|
 \le K_C\left[e^{-v}+Je^{R-v}+Je^{-c_0\sqrt{R/2}}
                   +Je^{-\alpha_CR}+\frac1{\log(2+J)}\right]. \tag{RQ7}
\]
Take \(R=v/2\) and \(J=\lfloor e^{(c_0/4)\sqrt v}\rfloor\).
The first four terms are \(O_C(v^{-1/2})\), and the last has
that order. All constants are independent of the heat parameter and
form vector. The gap comparison following (RGC41) completes (RQ1).
\end{proof}

The positive-tail budget has a precise limitation. For these actual
atoms and the moment (RQ4), we prove
\[
 (\log J)r_{{\rm tail},J}\longrightarrow1,\qquad
             \frac{\log I_J}{\log J}\longrightarrow1. \tag{RQ8}
\]
Here is a direct proof. Put
\(\eta=\sum_{j\ge1}(A_j/\mu_j)\delta_{\mu_j}\),
\(h(S)=H_C(\sqrt S)\), and \(T(Y)=\eta((Y,\infty))\).
The finite positive measure satisfies
\[
 h(S)=\int\frac{\mu}{S+\mu}\,d\eta(\mu)
 =[\tfrac12\log S+C-\log(2\pi)+O(S^{-1/2})]^{-1}
                    -\frac{A_0}{S+\mu_0}\sim\frac2{\log S}. \tag{RQ9}
\]
Splitting at \(Y\), positivity gives
\(T(Y)\ge h(Y^\alpha)-\eta(\mathbb R_+)Y^{1-\alpha}\) for
\(\alpha>1\), and
\(T(Y)\le(1+Y^{\alpha-1})h(Y^\alpha)\) for \(0<\alpha<1\).
First \(Y\to\infty\), then \(\alpha\to1\), gives
\(T(Y)\sim2/\log Y\).
Interlacing identifies the actual tail with \(\mu>\lambda_J\),
where \(\mu/(\mu-1/4)=1+O(\lambda_J^{-1})\) uniformly.
Since \(\log\lambda_J\sim2\log J\), the first limit follows.

For the second limit, the full expansion in (RQ9), rather than
differentiation of its equivalence symbol, gives
\[
 h(S)-h(2S)
 =\int\frac{S\mu}{(S+\mu)(2S+\mu)}\,d\eta(\mu)
                  \sim\frac{2\log2}{(\log S)^2}.   \tag{RQ10}
\]
For \(0<\epsilon<1\), this positive kernel is at most
\(S^{-\epsilon}\) outside
\([S^{1-\epsilon},S^{1+\epsilon}]\); the mass in that band is
therefore at least \(c_\epsilon/(\log S)^2\).
Choose \(S=\lambda_J^{1/(1+\epsilon)}\), placing the entire band
in the actual finite head. Formula (RQ4) yields
\[
 I_J\ge c_\epsilon
 \frac{\lambda_J^{(1-\epsilon)/(2(1+\epsilon))}}
                             {(\log\lambda_J)^2}. \tag{RQ11}
\]
First \(J\to\infty\), then \(\epsilon\downarrow0\), together with
\(I_J\le K_CJ\), proves the second limit in (RQ8).

Requiring this positive-tail budget itself to be \(e^{-v+o(v)}\)
forces \(\log\log J\ge v-o(v)\). The corresponding sufficient
absolute PNT budget \(I_Je^{-c_0\sqrt{R/2}}\), with \(R=O(v)\),
then diverges. This is a limitation of that absolute error
budget, not a lower bound on the signed arithmetic error.
Polynomial smoothing or a signed long-block estimate is a
different construction; increasing this head order does not
supply the missing exclusion estimate.


\subsection{A sharp Gaussian comparison for the entire zero form}
\label{sec:sharp-gaussian-rest}

\begin{lemma}[The full regularized zero identity and its weights]
\label{lem:rest-zero-identity}
The regularizer just constructed satisfies
\[
 |H_C(z)|\le\frac{K_C}{\log(2+|\Re z|)},\qquad
 \Re H_C(z)>0\quad (|\Im z|\le1/2).                 \tag{RGS1}
\]
For every $f\in C_c^\infty(\R)$, with $F=\widehat f$, the exact
full arithmetic identity is
\[
 q_{R_CA}(f)=\sum_{\Xi(z)=0}m_z H_C(z)
                         \overline{F(\bar z)}F(z). \tag{RGS2}
\]
The sum runs over all distinct nodes with their multiplicities
and converges absolutely. It includes every prime power and
both poles through the stated left-hand form.
\end{lemma}
\begin{proof}
For $z=u+iv$ in the closed strip, the arguments
$1/4\pm iz/2$ have nonnegative real parts. When $|u|$ is large
they stay uniformly in $|\arg w|\le\pi/2$ and away from the
digamma poles. The sector expansion
$\psi_{\rm dig}(w)=\log w-1/(2w)+O(|w|^{-2})$
in \cite[5.11.2 and Section 5.11(ii)]{DLMFGamma} therefore gives
\[
 a(u+iv)=\log(|u|/(2\pi))+O(|u|^{-1})
\]
uniformly on this strip. Since
$H_C(z)=(C+a(z))^{-1}-A_0/(z^2+\mu_0)$, this proves the
first bound at large $|u|$. On the remaining compact set use
the normally convergent positive rest series.
Each denominator there, and indeed throughout the strip, has
$\Re(z^2+\mu_j)=u^2-v^2+\mu_j>0$ for $j\ge1$.
The positive coefficients give $\Re H_C>0$ and a bounded
compact remainder. The same bound applies to all multiplicities.

To justify the new use of the explicit formula, put
$C_f(r)=\int\overline{f(x)}f(x+r)\,dx$ and $h=G_C*C_f$.
For every $1<b<b_1$ and $|r|\ge1$ the positive atom expansion gives
$G_C(r)\le K_b e^{-b|r|}$: write each summand as
$(A_j/\mu_j)(\sqrt{\mu_j}/2)e^{-\sqrt{\mu_j}|r|}$
and bound $\sqrt{\mu_j}e^{-(\sqrt{\mu_j}-b)|r|}$ uniformly.
The summable coefficients $A_j/\mu_j$ pay the entire series.
Consequently $h$ is smooth, and all its fixed derivatives have
the same exponential bound outside a compact set; differentiating
the compact smooth factor $C_f$ justifies this assertion even
at a singularity of $G_C$.

The multiplicative test $u^{-1/2}h(\log u)$ is therefore in the
class $\mathcal W$ of \cite[Section 2.1.1, eqs.~(2.1)--(2.7),
pp.~101--102]{ConnesConsani}: at zero it is
$O(u^{b-1/2})$, and at infinity it is
$O(u^{-b-1/2})$, with $b-1/2>0$.
Its derivative has the required continuity, and its Mellin
transform is $\widehat h$ under the same change of variables as
in (ZF). Fourier convolution gives exactly
\[
 \widehat h(z)=H_C(z)\overline{F(\bar z)}F(z).
\]
The two pole evaluations are multiplied by $r_C=H_C(i/2)$;
the archimedean integral has multiplier $aH_C$; and every
prime translation uses the actual values $h(\pm\log n)$.
Tonelli or absolute convergence pays all these terms:
$b>1/2$ gives $\sum\Lambda(n)n^{-b-1/2}<\infty$,
and the smooth even second difference pays the gamma singularity.
Thus the linear Weil formula is exactly (RGS2).
The rapid Fourier decay of the compact smooth $f$, uniformly
in $|\Im z|\le1/2$, and the unit-height count prove absolute
convergence of its complete zero sum. No larger raw physical
form domain has been assumed.
\end{proof}

\begin{theorem}[Sharp full-vector endpoint upper and exact negative rate]
\label{thm:sharp-gaussian-rest-rate}
Keep $C>\max(0,-a(0))$ fixed. For $t\ge1$ set
\[
 \phi_t=\phi_{t,t},\qquad
 \sigma_t=\frac{e^{-t/2}}{\sqrt{4\pi t}\cosh(t/2)},\qquad
 \eta_{C,t}=\max\{0,-\inf\operatorname{spec}(L_{C;t,t})\}.
                                                        \tag{RGS3}
\]
If every actual zero satisfies $|\Im z|\le d$, $0\le d\le1/2$,
then, with a constant independent of $t,d$ and every vector,
\[
 \|L_{C;t,t}\|\le K_C\sigma_t\sqrt t\,e^{4d^2t},
 \qquad \eta_{C,t}\le K_C\sigma_t\sqrt t\,e^{4d^2t}.
                                                        \tag{RGS4}
\]
Writing $d_*=\sup_{\Xi(z)=0}|\Im z|$, we also have
\[
 \lim_{t\to\infty}\frac1t
           \log(1+\eta_{C,t}/\sigma_t)=4d_*^2.          \tag{RGS5}
\]
This supremum need not be attained, and the off-line set may
be infinite. No simplicity assumption is made.
\end{theorem}
\begin{proof}
Put $\omega_t=\sqrt{\phi_t/c}$. Completing the two shifted
Gaussians in the exact ground gives
\[
 \frac{\omega_t(x)^2}{\sigma_t}=
 \int_\R\Phi(y)\frac{c(x-y)}{c(x)}
                    e^{-(x-y)^2/(4t)}\,dy.            \tag{RGS6}
\]
The actual theta tails imply
$A_1=\sup_y\Phi(y)e^{|y|/2+y^2/4}<\infty$.
As $c(x-y)\le c(x)e^{|y|/2}$, the full Gaussian integral gives
\[
 \omega_t(x)^2\le W_s(x):=
 A_G\sigma_t e^{-x^2/(4s)},\qquad
 s=t+1,\quad A_G=\sqrt{4\pi}A_1.                    \tag{RGS7}
\]
This is a positive multiplier comparison; no Fourier decay
of the non-Gaussian weight is inferred from it.

Let $w_z=m_z|H_C(z)|$. By (RGS1) and the unconditional
unit-interval count with multiplicities,
\[
 M_C:=\sup_{a\in\R}
     \sum_{\Re z\in[a,a+1)}w_z<\infty .              \tag{RGS8}
\]
An arbitrary real unit interval costs at most two integer ones.
For $R_z(x)=\sqrt{W_s(x)}e^{i\bar z x}$ in $L^2(dx)$, Gaussian
integration in the antilinear-first convention gives
\[
 N_z:=\|R_z\|_2^2=A_G\sigma_t\sqrt{4\pi s}
                                      e^{4s(\Im z)^2},
 \qquad
 \frac{|\langle R_z,R_{z'}\rangle|}{\sqrt{N_zN_{z'}}}
                         =e^{-s|z-z'|^2}.             \tag{RGS9}
\]
First take any finite subset of actual nodes. Apply the
matrix Schur test to the vectors $\sqrt{w_z}R_z$ with
coordinate weight $\sqrt{w_z}$. Its row bound is at most
\[
 N_d\sum_{z'}w_{z'}e^{-s(\Re z-\Re z')^2}
 \le N_dM_CC_0,\qquad
 N_d=A_G\sigma_t\sqrt{4\pi s}e^{4sd^2},
 \quad C_0=\sum_{k\in\Z}e^{-\max(|k|-1,0)^2}<\infty .
                                                        \tag{RGS10}
\]
The column bound agrees by symmetry. The constant is uniform
over all finite subsets. Passing to all nodes by monotone
convergence of sums of squares proves the full Gaussian
Bessel inequality.

For $g\in L^2(c\phi_tdx)$ write
$u=\sqrt{c\phi_t}\,g$. The factor
$r_t=\omega_t/\sqrt{W_s}$ satisfies $0<r_t\le1$.
Thus each actual Fourier evaluation is exactly the model
evaluation of $r_tu$, and the Bessel inequality yields
\[
 \sum_zm_z|H_C(z)|\,|F_{\phi_tg}(z)|^2
                   \le M_CC_0N_d\|g\|_{L^2(c\phi_tdx)}^2.
                                                        \tag{RGS11}
\]
Each individual integral exists by the same Gaussian
Cauchy--Schwarz bound. The test map is therefore defined and
bounded on the entire weighted Hilbert space.

On the smooth compact core (RGS2), conjugate-node pairing,
and $2|ab|\le|a|^2+|b|^2$ give
$|q_{R_CA}(\phi_tg)|\le\sum_z w_z|F_{\phi_tg}(z)|^2$.
Polarization of the same sample bound defines a bounded form;
it coincides with the already bounded whole operator on a
dense core and hence throughout the space. This proves both
the full zero identity in this weighted realization and its
norm bound. Since $\sqrt{s}\le\sqrt{2t}$ and
$e^{4sd^2}\le e\,e^{4td^2}$, (RGS4) follows.
All prime powers and pools remain in the matched arithmetic
form; none has been removed to obtain the comparison.

For the lower rate fix any one actual off-line pair
$a=\gamma+i\delta,\bar a$, $\delta>0$, of multiplicity $m$.
Theorem~\ref{thm:endpoint-leading-cost} provides actual fast
tests with only the values $1,-1$ on this pair and both
poles zero, whose squared weighted norm is
$2/(N_t(\delta)-N_t(0))(1+o(1))$.
The full weighted identity just proved gives their value
$-2m\Re H_C(a)<0$. Consequently
\[
 \eta_{C,t}\ge (1+o(1))m\Re H_C(a)
       \Xi(i(2\delta-1/2))\,
               \sigma_t\sqrt{4\pi t}\,e^{4\delta^2t}.
                                                        \tag{RGS12}
\]
The real part is strictly positive by (RGS1), and the
Laplace moment of the original positive theta density is
strictly positive. This is a fixed-pair bound, with no
uniform isolation hypothesis for a growing zero family.
It gives the lower logarithmic rate $4\delta^2$ for EVERY
such pair. Taking their supremum gives $4d_*^2$ without
selecting a most extreme zero. The full upper with $d=d_*$
gives the opposite inequality. If no off-line pair exists,
all terms in (RGS2) are nonnegative real-node terms;
bounded extension gives $\eta_{C,t}=0$, so the rate is zero.
This completes all cases of (RGS5).
\end{proof}

The unconditional choice $d=1/2$ is permitted in (RGS4).
The theorem proves the exact cost of any actual imaginary
displacement, rather than proving that displacement absent.
A cofinal bound $\eta_{C,t}\le\sigma_t e^{o(t)}$ would now
contradict (RGS12) for every off-line pair individually;
that signed arithmetic bound remains unproved.

\subsection{A fixed extra resolvent gives a complete trace-class operator}
\label{sec:smoothed-rest-trace}

\begin{theorem}[Uniform nuclear decay with every zero still detected]
\label{thm:smoothed-rest-trace}
Fix the same admissible \(C\), and put
\[
 E_1(x)=\tfrac12e^{-|x|},\quad G_1=E_1*G_C,\quad
 U_1=E_1*U_C^{\rm hi},\quad
 h_1(z)=\frac{H_C(z)}{1+z^2},\quad r_1=\tfrac43r_C .
 \tag{RT1}
\]
The full physical kernel and its weighted operator are
\[
 Q_1=(1+r_1)E_1-U_1-G_1*\Pi+r_1e^{|x|/2},\qquad
 L^{(1)}_{v,t}(x,y)=\omega_{v,t}(x)Q_1(x-y)\omega_{v,t}(y).
 \tag{RT2}
\]
For every fixed \(v,t\ge0\) this is a self-adjoint trace-class
operator on the entire \(L^2(\mathbb R)\). It retains every prime
power, both pools, and the full negative index \(2K_+\), with
index \(K_+\) in each parity and infinity allowed.
There are fixed constants \(K,c>0\) such that
\[
 \sup_{t\ge0}\|L^{(1)}_{v,t}\|_1\le Ke^{-c\sqrt v}\quad(v\ge1).
 \tag{RT3}
\]
The trace, rather than the trace norm, has the exact formula
\[
 \operatorname{tr}L^{(1)}_{v,t}
 =N_0(v,t)\tau_C,\quad 0<N_0(v,t)\le2e^{-v},\qquad
 \tau_C=\tfrac12+2r_C-U_1(0)
       -2\sum_{n\ge2}\frac{\Lambda(n)}{\sqrt n}G_1(\log n)>0 .
 \tag{RT4}
\]
\end{theorem}
\begin{proof}
Direct integration gives
\(E_1*e^{|x|/2}=\frac43e^{|x|/2}+\frac43E_1\),
\(E_1*e^{-|x|/2}=\frac43e^{-|x|/2}-\frac43E_1\), and
\(E_1*c=\frac43c\). Thus (RT2) is precisely the full
\((1-\partial_x^2)^{-1}R_CA\), including its nonlocal first term.
The positive kernel \(S_1^+=E_1*S_C^+\) has the exact row
\(S_1^+*\phi=E_1*\phi+r_1c\).
Writing \(D_K^\phi=\frac12\iint K(x-y)\phi(x)\phi(y)
|g(x)-g(y)|^2\,dx\,dy\), its complete ground form is
\[
 D_{S_1^+}^\phi-D_{E_1}^\phi-r_1\|g\|_\mu^2
       +2r_1\left|\int g\,d\mu\right|^2
       -2r_1\left|\int\kappa g\,d\mu\right|^2 .     \tag{RT5}
\]
Both energies are bounded on all \(L^2(\mu)\): use
\(E_1*\phi\le1/4\), the positive row, and the square difference
inequality. Tonelli pays every prime power; the pool moments
are bounded. Smooth compact functions form a norm core, so
(RT5) proves the actual entire-space physical extension.

Since \(b_1>1\), the atom sums give
\(\int e^{|x|}G_C(x)\,dx<\infty\), and likewise for
\(U_C^{\rm hi}\). Hence
\(G_1\le Ke^{-|x|}\), \(|G'_1|\le G_1\),
\(G_1''=G_1-G_C\in L^1\), and \(U_1\le Ke^{-|x|}\).
Chebyshev--Abel and the convergent sum
\(\sum\Lambda(n)n^{-3/2}\) make \(G_1*\Pi\) continuous and
bounded by \(Ke^{|x|/2}\).
Apply the linear Weil formula used in (RGS2) to
\((E_1*G_C)*C_f\). Its smooth derivatives have exponential
decay of every fixed rate below \(1\), and in particular above
\(1/2\); it is in the same multiplicative test class.
Polarization gives
\[
 q_{(1-\partial^2)^{-1}R_CA}(f,h)
 =\sum_zm_z h_1(z)\overline{F_f(\bar z)}F_h(z).       \tag{RT6}
\]
This also follows by compact cutoffs of \(E_1*h\);
their exponential derivative norms and uniform Fourier
integration by parts pay the entire node sum.

On \(|\Im z|\le1/2\), each denominator \(z^2+\mu\), \(\mu\ge1\),
has argument at most \(\pi/6\) in magnitude, since
\(|u|/(u^2+3/4)\le1/\sqrt3\).
The inverse atom summands have the same argument sign. Therefore
\[
 \Re h_1(z)\ge\tfrac12|h_1(z)|>0,\quad
 |h_1(z)|\le\frac{K_C}{(1+|\Re z|^2)\log(2+|\Re z|)},\quad
 \sum_{|\Re z|>T}m_z|h_1(z)|\le K_C/T.             \tag{RT7}
\]
The unconditional unit-height zero count, including multiplicity,
justifies both sums. Equality of compact product pairings in
(RT6) identifies the continuous kernel:
\(Q_1(r)=\sum_zm_z h_1(z)e^{izr}\), normally on each real compact.

Define, for \(0\le\delta\le1/2\),
\[
 N_\delta(v,t)=\int_{\R}\omega_{v,t}(x)^2e^{2\delta x}\,dx,
 \qquad u_z=\omega e^{izx},\quad v_z=\omega e^{i\bar z x}.
\]
Evenness and the exact ground moment give
\(\|u_z\|\|v_z\|=N_{|\Im z|}\le N_{1/2}=1-N_0\le1\).
Consequently
\[
 L^{(1)}_{v,t}=\sum_zm_z h_1(z)|u_z\rangle\langle v_z|
                                                        \tag{RT8}
\]
converges in trace norm on the whole space.
Real nodes are positive rank-one forms; each nonreal conjugate
pair has one negative direction in its off-diagonal value matrix.
The pole-free fast interpolation of
Lemma~\ref{lem:pole-free-interpolation} prescribes any finite
set of pair values and annihilates the others.
Estimate (PNR) gives finite norms for each fixed shifted heat
ground. At zero heat and \(v>0\), on the right tail the
denominator is at least \(\Phi(x-v)/(2c(v))\), so the
integrand is bounded by a constant times
\(c(x)\Phi(x)^2e^{M|x|}/\Phi(x-v)\); its theta
exponents make it integrable. Evenness treats the left tail;
at \(v=t=0\) use \(\phi=\Phi\), and positive heat has a
Gaussian mass-slice lower bound. Compact weighted cutoffs
therefore preserve the negative matrices by bounded form
convergence. The parity values \((1,-1,\varepsilon,-\varepsilon)\)
on \((z,\bar z,-z,-\bar z)\) give
\(-4m_z\Re h_1(z)<0\) per selected quartet in each parity.
The finite upper count and arbitrary finite lower selections
prove the full index, including infinity.
There are no height-zero off-pairs, since \(\Xi(i\delta)>0\).

Each rank-one trace is exactly
\(\langle v_z,u_z\rangle=N_0\), so (RT8) yields (RT4) and
\(\tau_C=\sum_zm_zh_1(z)>0\).
The positive real parts and the nonempty zero set prove strict
positivity. Formula (RT2) at zero gives its arithmetic expression.
For \(N_0\), use the two shifted copies and
\(c(x)\ge e^{\pm x/2}/2\). Their exact moments
\(\int p_t e^{\pm x/2}=1/2\) give \(N_0\le2e^{-v}\).

More generally set \(\beta=1-2\delta\), \(\theta=2\delta-1/2\).
On \(x>0\), \(1/c(x)\le2e^{-x/2}\).
The first copy has bound
\(e^{\theta v}\Xi(i\theta)e^{-2\delta\beta t}/c(v)\).
For the second use \(e^{-\theta v}\le e^{\theta v}\) if
\(\theta\ge0\); if \(\theta\le0\), write \(y=x+v\ge v\)
and use \(e^{\theta y}\le e^{2\theta v}e^{-\theta y}\).
Evenness and \(\Xi(i\theta)\le1/2\) prove
\[
 N_\delta(v,t)\le4e^{-(1-2\delta)v-2\delta(1-2\delta)t}.
                                                        \tag{RT9}
\]
The unconditional classical zero-free region
\href{https://people.maths.ox.ac.uk/greenbj/papers/primenumbers.pdf}
{Green, Theorem 9.1 and its bounded-height remark, printed
pages 53--55}, and the functional symmetries give
\(1-2|\Im z|\ge2b_0/\log(2+|\Re z|)\).
Split (RT8) at \(T=e^{\sqrt v}\).
The lower heights cost \(Ke^{-2b_0v/\log(T+2)}\) by (RT9);
the upper heights cost \(K/T\) by (RT7) and \(N_\delta\le1\).
This proves (RT3), uniformly in every heat parameter.
\end{proof}

The positive trace (RT4) does not bound the negative trace or
exclude a negative eigenvalue. Opposite members of an off-line
pair can cancel in the net trace. The extra resolvent removes
the previous singular-kernel obstruction; the missing signed
arithmetic estimate remains a separate obligation.
For \(x,y>0\), put
\[
 W^{(1)}(x,y)=\omega(x)
 [(G_1*\Pi)(x+y)-r_1e^{(x+y)/2}]\omega(y),\quad
 d\nu(x)=\omega(x)^2e^x\,dx .
\]
The exact whole-vector energy is
\[
 \|W^{(1)}\|_{\rm HS}^2
 =\iint\left|e^{-(x+y)/2}(G_1*\Pi)(x+y)-r_1\right|^2
                                                    d\nu(x)d\nu(y).
 \tag{RT10}
\]
The required bound \(e^{-2v+o(v)}\) on a cofinal sequence
\(t(v)\to\infty,\ t(v)=o(v)\) is unproved.
It would contradict each fixed off-node's detector depth
\(c_z t^{-1/2}e^{-t/4}e^{-(1-2|\Im z|)v}\).
Here the depth follows directly from (PNR). Select a pole-free
interpolant \(f\) with values one at \(z=\gamma+i\delta\)
and \(\bar z\), zero at all other nodes, and set
\(u_v=\tau_vf-e^{-2i\gamma v}\tau_{-v}f\).
The two values and the antilinear first slot in (RT6) give
\[
 \begin{gathered}
 q_{(1-\partial^2)^{-1}R_CA}(u_v)
       =-8m_z\Re h_1(z)\sinh^2(\delta v),\\
 \|u_v\|_\phi^2\le8e^vJ_{f,t},\qquad
 J_{f,t}=\int c|f|^2/p_t\le K_f/a_t\quad(t\ge1).
 \end{gathered}
\]
For the last inequality integrate \(\Phi\) on a fixed positive
mass interval to obtain
\(p_t(x)\ge c a_t e^{-x^2/4-B|x|/2}\) for every \(t\ge1\),
and then use (PNR). The elementary same-side kernel bound is
\(K e^{-v+t}\); relative to this negative depth it is
\(O(\sqrt t\,e^{-2\delta v+5t/4})\), tending to zero when
\(t=o(v)\). Thus the whole-vector cross-energy bound stated
after (RT10), if proved, would exclude each such fixed node.
Smoothing leaks support, so the raw finite Möbius cutoff cannot
be reused without its extra tail or outgoing-mode calculation.


\subsection{Four exact moments preserve support and the entire negative index}
\label{sec:support-resolvent}

Write \(F_f(z)=\widehat f(z)\), with Fourier phase \(e^{-izx}\),
and \(\|f\|_\phi^2=\int c|f|^2/\phi_{v,t}\).
The next regularizer acts on the original complete \(A\), rather
than on the earlier gamma-rest operator. Put
\[
 \begin{gathered}
 J_*=(1-\partial^2)^{-1}(4-\partial^2)^{-1},\qquad
 j_*(z)=\frac1{(1+z^2)(4+z^2)},\\
 P_*(z)=(1+z^2)(4+z^2),\qquad
 K_*(x)=\frac{2e^{-|x|}-e^{-2|x|}}{12}.
 \end{gathered}                                      \tag{SC1}
\]

\begin{theorem}[Support-preserving complete regularization]
\label{thm:support-resolvent}
For a compact interval \(I\) and \(h\in L^2\) supported in \(I\),
\[
 \operatorname{supp}(J_*h)\subset I
 \quad\Longleftrightarrow\quad
 F_h(i)=F_h(-i)=F_h(2i)=F_h(-2i)=0.                  \tag{SC2}
\]
In that case \(J_*h\in H^4\). Its complete continuous kernel is
\[
 \begin{split}
 Q_*&=\gamma_*^{\rm hi}-K_* *\Pi+\frac{16}{45}e^{|x|/2},\\
 \gamma_*^{\rm hi}(r)
 &=\frac{a(i)}6e^{-|r|}-\frac{a(2i)}{12}e^{-2|r|}
   -\sum_{n\ge1}\frac{e^{-(2n+1/2)|r|}}
                         {(\lambda_n-1)(\lambda_n-4)},\quad
 \lambda_n=4(n+1/4)^2.
 \end{split}                                        \tag{SC3}
\]
For every \(v,t\ge0\), \(\omega(x)Q_*(x-y)\omega(y)\)
defines a trace-class operator on the entire canonical space.
Its negative index is \(2K_+\), and each parity index is \(K_+\),
also after restricting to the four moments in (SC2) and to both
original pole moments \(F_f(\pm i/2)=0\).
Here infinity is allowed and multiplicities are arbitrary.
\end{theorem}
\begin{proof}
Partial fractions give (SC1). Outside \(I=[a,b]\) its convolution
is respectively
\[
 12(J_*h)(x)=2e^{-x}F_h(i)-e^{-2x}F_h(2i)\quad(x>b),
 \qquad
 12(J_*h)(x)=2e^xF_h(-i)-e^{2x}F_h(-2i)\quad(x<a).
\]
Linear independence proves both directions of (SC2); Plancherel
gives \(H^4\) and \(P_*(-i\partial)J_*h=h\).
For smooth tests the constrained range is exactly
\(P_*(-i\partial)C_c^\infty(I)\), where compact support is
inside the interval's interior.

The normally convergent digamma series in (RGC4) gives
\(a(z)=a(0)+\sum_{n\ge0}(n+1/4)^{-1}z^2/(z^2+\lambda_n)\).
The residues of \(a(z)j_*(z)\), as a function of \(z^2\),
at \(-1,-4,-\lambda_n\) are
\(a(i)/3,-a(2i)/3,-\lambda_n/[(n+1/4)(1-\lambda_n)(4-\lambda_n)]\).
Inverse Fourier transformation divides the last by \(2(2n+1/2)\).
The \(n=0\) term is exactly \(-16e^{-|r|/2}/45\).
The remaining coefficients are \(O(n^{-4})\), while the two
fixed-pole residue series converge absolutely. Thus all partial
fractions and their inverse transforms are justified; no delta
term is omitted. The original pool becomes
\(2j_*(i/2)c=32c/45\), and its combination with that low gamma
term is exactly the last term of (SC3).
Chebyshev--Abel gives \(K_* *\Pi\le Ke^{|r|/2}\), locally
uniformly with the summable majorant \(\Lambda(n)n^{-3/2}\).

The smoothing Weil formula, justified by exponential cutoffs of
\(J_*h\) of any fixed decay rate between \(1/2\) and \(1\), is
\[
 q_{J_*A}(f,h)=\sum_{z\in\mathcal Z}m_zj_*(z)
                         \overline{F_f(\bar z)}F_h(z). \tag{SC4}
\]
The same local zero count as in (RT7) makes
\(\sum m_z|j_*(z)|<\infty\), and each weighted rank-one
norm is \(N_{|\Im z|}\le1\). This identifies the whole
trace-class operator. On the closed strip both denominator
arguments have the same sign and magnitude at most \(\pi/6\),
so \(\Re j_*>0\) and \(\Re P_*>0\).
Take the interpolants of (PNR) and apply \(P_*(-i\partial)\).
The selected conjugate values \((1,-1)\) give
\(-2m_z\Re P_*(z)\), since the new coefficient is
\(j_*(z)\overline{P_*(\bar z)}P_*(z)=P_*(z)\),
rather than \(j_*|P_*|^2\).
Parity values \((1,-1,\varepsilon,-\varepsilon)\) give
\(-4m_z\Re P_*(z)\). Apply the polynomial to compact
cutoffs and to the fixed two pole correctors in (PC).
All four extra moments remain exactly zero; the pole matrix
is multiplied by \(P_*(i/2)=45/16\), so remains invertible.
Estimate (PNR) gives convergence in each fixed weighted norm,
including zero heat, as in (RT8). Arbitrary finite selections
give the lower indices, also for infinitely many pairs.
The positive real-node blocks and the one negative direction
per nonreal conjugate block give the matching upper indices.
All moment functionals are continuous on each fixed weighted
space by Cauchy--Schwarz and the Gaussian or theta tails.
Thus the asserted constrained spaces are closed.
\end{proof}

\begin{theorem}[An ordinary finite Möbius source for every constrained pair]
\label{thm:ordinary-finite-source}
Let \(f,h\in L^2\) have ordered compact support intervals
\(I_f,I_h\), with \(\sup I_f<\inf I_h\), and put
\(B=\sup I_h-\inf I_f\). Suppose (SC2) holds for \(h\) and
\(F_h(i/2)\overline{F_f(-i/2)}=0\).
For every integer \(N\ge e^B\),
\[
 q_{J_*A}(f,h)=\frac{i}{2\pi}\int_{\R}
 j_*(u)F_h(u)\overline{F_f(u)}
 \frac{\Xi'(u)}{\mathfrak G(1/2-iu)}D_N(1/2-iu)\,du,
                                                               \tag{SC5}
\]
where
\(\mathfrak G(s)=\tfrac12s(s-1)\pi^{-s/2}\Gamma(s/2)\) and
\(D_N(s)=\sum_{d\le N}\mu(d)d^{-s}\).
This is an absolutely convergent ordinary integral on the whole
stated \(L^2\) class. No limit \(D_N\to1/\zeta\) is used.
The inward pole product may be nonzero.
\end{theorem}
\begin{proof}
We give the causal finite-source calculation. Initially take
smooth ordered compact tests \(f,k\), and set
\(C(r)=\int\overline{f(x)}k(x+r)\,dx\),
\(W(z)=F_k(z)\overline{F_f(\bar z)}\), and
\(g=\mathfrak G'/\mathfrak G\).
Then \(\operatorname{supp}C\subset(0,B]\) and \(W(i/2)=0\).
The inverse Fourier transform of \(g(1/2-iu)W(u)\) is
\[
 H(r)=-\int_0^\infty e^{-s/2}C(r-s)\,ds
 -\frac{\log\pi+\gamma_E}{2}C(r)
 +\int_0^\infty
 \frac{e^{-2s}C(r)-e^{-5s/2}C(r+s)}{1-e^{-2s}}\,ds.
                                                               \tag{SC6}
\]
This follows from the digamma partial fractions and the recurrence
\(\mathfrak G(s)=(s-1)\pi^{-s/2}\Gamma(1+s/2)\).
At zero the numerator is \(O(s)\); its derivatives have the same
bound, and all large-\(s\) tails converge absolutely.
Thus Fourier inversion is justified first on
\([\epsilon,M]\), then by its integrable limiting bounds.
The sole right tail vanishes by \(W(i/2)=0\);
\(H\) has right support at most \(B\) and a left tail
\(O(e^{5r/2})\). The inverse of \(W'\) is \(-irC(r)\).
Their exponential moments are respectively
\(iW'(i/2)\) and \(W'(i/2)\), since \(g(1/2-iu)\)
has principal part \(i/(u-i/2)\).

For any of these inverse functions \(H_T\), Stieltjes
integration with \(Q(x)=\lfloor x\rfloor-x\) gives
\[
 \sum_{k\ge1}(dk)^{-1/2}H_T(\log(dk))
 =\frac{T(i/2)}d+
 \frac{d^{-1/2}}{2\pi}\int T(u)\zeta(1/2-iu)e^{iu\log d}\,du.
\]
Both endpoints vanish; \(|Q(x)|\le\min(x,1)\), the just
proved left tail, and Schwartz decay pay absolute Fubini.
All terms with \(dk>e^B\) vanish. Finite divisor inversion
\(\mu*1=\epsilon\), with
\(S_N=\sum_{d\le N}\mu(d)/d\), therefore yields
\[
 \int gW(\zeta D_N-1)=-2\pi iW'(i/2)S_N,\qquad
 \int W'(\zeta D_N-1)=-2\pi W'(i/2)S_N.             \tag{SC7}
\]
The zero-principal original \(C\) and
\(\Lambda=-(\mu\log)*1\) give the complete forward prime
half. With \(B_N(s)=\sum_{d\le N}\mu(d)\log d\,d^{-s}\),
it is \((2\pi)^{-1}\int[gW+W\zeta B_N]\).
Use \(\zeta B_N=\zeta'D_N-(\zeta D_N)'\) and
\(\zeta'=i\Xi'/\mathfrak G-g\zeta\).
Integration by parts and (SC7) cancel the two finite principal
corrections exactly, leaving \(i(2\pi)^{-1}\int
W(\Xi'/\mathfrak G)D_N\).
The opposite prime orientation is zero by ordering. Its gamma
half is \(-\int_0^\infty e^{-r/2}C(r)\,dr=-W(-i/2)\)
by (SC6) with reversed shifts; the retained inward pole is
\(+W(-i/2)\). They cancel. The outward pole is zero by
hypothesis. This proves the identity for the full \(A\).

Now use \(k=J_*h\). Equation (SC2) preserves its support,
and \(F_k=j_*F_h\) preserves the outward zero. This proves
(SC5) on the smooth core. The elementary floor continuation
\(\zeta(s)=s/(s-1)-s\int_1^\infty\{x\}x^{-s-1}dx\)
and its derivative give \(|\zeta|+|\zeta'|\le K(1+|u|)\)
on \(s=1/2-iu\); the digamma series gives
\(|g|\le K\log(2+|u|)\). Since \(|D_N|\le2\sqrt N\),
the complete multiplier in (SC5) is bounded by \(K\sqrt N\).
Plancherel gives \(\int|F_fF_h|\le2\pi\|f\|_2\|h\|_2\).
Finally approximate inside the fixed intervals and restore the
four moments and whichever outward factor vanishes with fixed
smooth correctors. The relevant distinct exponentials have a
positive definite Gram on an interval with interior. Both the
kernel form and the ordinary integral are \(L^2\)-continuous.
Their limits give the assertion for every stated pair.
\end{proof}

\begin{lemma}[Joint compact detectors with a logarithmic radius]
\label{lem:joint-compact-detector}
Fix an actual hypothetical node \(a=\gamma+i\delta\),
\(0<\delta<1/2\), of order \(m_a\).
There are constants \(c_a,C_a,R_0>0\) and smooth functions
\(f_R\) supported on \([-R-1,R+1]\), with all six moments
\(F_{f_R}(\pm i)=F_{f_R}(\pm2i)=F_{f_R}(\pm i/2)=0\), such
that the translated twins, for
\(R=R_0+\log(v+t+2)\), satisfy
\[
 \frac{-q_{J_*A}(\tau_vf_R-e^{-2i\gamma v}\tau_{-v}f_R)}
 {\|\tau_vf_R-e^{-2i\gamma v}\tau_{-v}f_R\|_\phi^2}
 \ge c_a a_t e^{-(1-2\delta)v}                       \tag{SC8}
\]
for all sufficiently large \(v\) along every regime
\(t\to\infty\), \(t=o(v)\).
Their actual finite source cap is
\(\lceil e^{2v+2R+2}\rceil=e^{2v+o(v)}\).
\end{lemma}
\begin{proof}
Select a fixed pole-free \(w\) of (PNR) with values one at
\(a,\bar a\) and zero at every other node. Restore the two
pole values of \(\chi_Rw\) as in (PC), obtaining \(z_R\),
and put \(f_R=P_*(-i\partial)z_R\), \(f=P_*(-i\partial)w\).
All six moments vanish exactly. A fixed positive mass slice
of \(\Phi\) gives, uniformly for \(t\ge1\),
\(p_t(x)\ge c a_t e^{-x^2/4-B|x|/2}\).
Leibniz and (PNR) hence yield
\[
 \int c|f_R-f|^2/p_t\le C a_t^{-1}e^{-c e^{2R}},\qquad
 \int c(|f|^2+|f_R|^2)/p_t\le C a_t^{-1}.          \tag{SC9}
\]
The denominator copy and \(c(v+x)\le2c(v)c(x)\) give
\(\|\tau_{\pm v}h\|_\phi^2\le4c(v)^2\int c|h|^2/p_t\).
The uniformly bounded full operator in (SC4) therefore makes
the twin form error at most
\(C c(v)^2a_t^{-1}e^{-c'e^{2R}}\).
The ideal twin form is exactly
\(-8m_a\Re P_*(a)\sinh^2(\delta v)\), with the conjugation
already checked in (SC4). For the stated radius that error is
\(o(e^{2\delta v})\), whereas the squared norm is at most
\(C c(v)^2/a_t\). This proves (SC8), with all joint quantifiers.
The two support intervals give the stated integer cap directly.
\end{proof}

The new support transfer makes (SC5) applicable to complete
constrained moving windows. It supplies no signed cancellation
estimate for its finite Möbius multiplier. Such an estimate of
\(e^{-v+o(v)}\) in canonical norm, together with the elementary
same-side cost \(Ke^{-v+t}\), would contradict (SC8) for each
fixed \(\delta>0\). This is a separate, still open obligation.

\subsection{An actual arithmetic limit of growing complete resolvents}
\label{sec:growing-resolvents}

For integers \(m\ge1\), define
\[
 Q_m(z)=\prod_{k=1}^m(1+z^2/k^2),\quad j_m=Q_m^{-1},\quad
 J_m=\prod_{k=1}^m k^2(k^2-\partial^2)^{-1},\quad
 K_m=E_1*\cdots*E_m,\quad E_k(x)=\frac{k}{2}e^{-k|x|}.
                                                               \tag{GL1}
\]
These act on the full original \(A\); they are distinct from
the earlier resolvents multiplied by \(R_C\).

\begin{theorem}[Uniform arithmetic identification and complete index]
\label{thm:growing-resolvent-limit}
The complete weighted operators \(L_m(v,t)\) are self-adjoint
and trace-class for every \(m\ge1\), \(v,t\ge0\).
They retain every prime power, both pools, the exact negative
index \(2K_+\), and parity indices \(K_+\), including infinity.
There is an actual arithmetic limit \(L_\infty(v,t)\) with
\[
 j_\infty(z)=\frac{\pi z}{\sinh(\pi z)},\qquad
 K_\infty(x)=\frac1{4\cosh^2(x/2)},\qquad
 \sup_{v,t\ge0}\|L_m-L_\infty\|_1\le K/m.            \tag{GL2}
\]
The value at \(z=0\) is one. The limit has the same complete
indices, and for fixed positive constants \(K,b\),
\[
 \sup_{m,v,t}\|L_m(v,t)\|_1\le K,\qquad
 \sup_{t\ge0}\|L_\infty(v,t)\|_1
          \le Ke^{-b v/\log(v+e)}\quad(v\ge2).         \tag{GL3}
\]
\end{theorem}
\begin{proof}
The primary product identity
\href{https://dlmf.nist.gov/4.36.E1}{DLMF 4.36.1} gives normally
on compact sets \(Q_m\to Q_\infty=\sinh(\pi z)/(\pi z)\).
For \(z=x+iy\), \(|y|\le1/2\),
\[
 \Re Q_\infty(z)=
 \frac{x\sinh(\pi x)\cos(\pi y)+y\cosh(\pi x)\sin(\pi y)}
      {\pi(x^2+y^2)}>0,                               \tag{GL4}
\]
with the removable value at zero. The factor arguments have
the sign of \(xy\). Their convergent sum is continuous, equals
zero on \(y=0\), and cannot cross \(\pm\pi/2\) by (GL4).
Consequently every partial sum lies between zero and the total;
\(\Re Q_m,\Re j_m>0\) on the whole closed strip. For each
fixed \(a\) there, the product bound
\(\prod(1-|\Im a|^2/k^2)\ge2/\pi\), normal convergence,
and those arguments give constants independent of \(m\) with
\[
 c_a\le|Q_m(a)|\le C_a,\quad \Re Q_m(a)\ge c_a,
             \quad\Re[j_\infty(a)Q_m(a)^2]\ge c_a.    \tag{GL5}
\]
For the last assertion, the angle is
\(2\theta_m-\theta_\infty\), whose magnitude is at most
\(|\theta_\infty|<\pi/2\). No margin uniform in growing
node height is claimed.

We next identify the full physical operators. For every \(m\),
\(K_m\le2e^{-|x|}\). Indeed, write
\(K_m=E_1*\kappa_m\), \(\kappa_m=E_2*\cdots*E_m\),
and use
\(\int e^{|x|}\kappa_m\le2\prod_{k=2}^m k^2/(k^2-1)\le4\).
Let \(\gamma_m\) be the ordinary inverse transform of
\(a(u)j_m(u)\). It is well defined since
\(|a(u)|\le K\log(2+|u|)\) and \(j_m(u)\le(1+u^2)^{-1}\).
Put \(r_m=j_m(i/2)\), \(1\le r_m\le\pi/2\), and
\(\gamma_m^{\rm hi}=\gamma_m+r_me^{-|x|/2}\).
The full physical kernel is
\[
 \mathcal Q_m=\gamma_m^{\rm hi}-K_m*\Pi+r_me^{|x|/2}.
                                                               \tag{GL6}
\]
The low gamma residue and the two original pools combine exactly
as in (SC3). To verify uniform tails, the first partial fraction is
\[
 \gamma_1^{\rm hi}(r)=\tfrac12a(i)e^{-|r|}
       +\sum_{n\ge1}\frac{e^{-(2n+1/2)|r|}}{\lambda_n-1},
 \qquad \gamma_1=\gamma_1^{\rm hi}-\tfrac43e^{-|r|/2}.
\]
Its coefficients are absolutely summable, giving the stated
high-kernel bound directly when \(m=1\). For \(m\ge2\),
\(\kappa_m\le12e^{-2|r|}\), by
\(2\prod_{k=3}^m k^2/(k^2-4)\le12\).
For \(r\ge0\), the exact convolution defect is
\[
 \begin{split}
 D_m(r)&=(\kappa_m*e^{-|\cdot|/2})(r)
                  -j_{\kappa_m}(i/2)e^{-r/2}\\
 &=-e^{-r/2}\int_r^\infty\kappa_m(y)e^{y/2}dy
        +e^{r/2}\int_r^\infty\kappa_m(y)e^{-y/2}dy.
 \end{split}
\]
Thus \(|D_m|\le Ke^{-2r}\) and
\(\gamma_m^{\rm hi}=\kappa_m*\gamma_1^{\rm hi}-\frac43D_m\)
is bounded by \(Ke^{-|r|}\), uniformly in \(m\).
Chebyshev--Abel with the same uniform kernel bound makes
\(K_m*\Pi\le Ke^{|r|/2}\) with locally absolute convergence.
All \(\Lambda(n)\) and both pool orientations are retained.

The smoothing Weil formula, with the same exponentially convergent
cutoffs as in (SC4), gives
\[
 q_{J_mA}(f,h)=\sum_zm_zj_m(z)\overline{F_f(\bar z)}F_h(z).
                                                               \tag{GL7}
\]
On the closed strip one has
\[
 |j_m(x+iy)|\le\frac\pi2\prod_{k=1}^m(1+x^2/k^2)^{-1}.\tag{GL8}
\]
This follows from
\(|k^2+z^2|\ge(k^2-y^2)(1+x^2/k^2)\).
The unconditional unit-height count makes
\(\sum m_z|j_m(z)|\le K\) uniformly in \(m\).
Since each weighted rank-one norm in (RT8) is \(N_{|\Im z|}\le1\),
the sum is trace-norm convergent on the entire space. Compact
product pairings identify it with (GL6), proving the first bound
in (GL3). Real nodes are positive; each off-line conjugate block
has one negative value direction. Lemma (PNR), multiplication by
the fixed polynomial \(Q_m\), and compact weighted cutoffs realize
arbitrary finite selections, with either parity and both pole
values zero. This gives the complete indices, including infinity.

The actual limiting physical kernel is identified as follows.
Substitute \(u=e^x\) into the transform of \(K_\infty\):
\[
 \int_\R K_\infty(x)e^{-izx}dx
 =\int_0^\infty\frac{u^{-iz}}{(1+u)^2}du
 =\Gamma(1-iz)\Gamma(1+iz)=\frac{\pi z}{\sinh\pi z}.
\]
The primary beta integral
\href{https://dlmf.nist.gov/5.12.E3}{DLMF 5.12.3}, its
\href{https://dlmf.nist.gov/5.12.E1}{gamma quotient 5.12.1}, and
\href{https://dlmf.nist.gov/5.5.E3}{reflection 5.5.3}, with
the recurrence 5.5.1, apply for \(|\Im z|<1\); the value at
zero is removable. Dominated Fourier inversion with \(j_1\)
gives uniform real-kernel convergence. The uniform exponential
tails give convergence in \(L^1(e^{b|x|}dx)\) for every \(b<1\).
The majorant \(\Lambda(n)n^{-3/2}\) passes every prime power on
compact physical intervals. Dominated inversion of \(aj_m\) and
\(r_m\to\pi/2\) also pass the gamma and pool terms. Hence (GL6)
has a genuine arithmetic physical limit, retaining all terms.

For the uniform trace error split the nodes at \(|\Re z|=\sqrt m\).
When \(m\ge4\) and the node is below this height,
\(\sum_{k>m}|\log(1+z^2/k^2)|\le K|z|^2/m\), a bounded
quantity, so
\(|j_m-j_\infty|\le K|z|^2|j_m|/m\).
Equation (GL8) with \(m\ge2\) gives
\(\sum m_z|z|^2|j_m(z)|\le K\), proving the low cost \(K/m\).
The high cost is at most
\(K\log(2+m)/m^{3/2}\le K'/m\), using the same \(j_2\)
majorant and local zero count. Each rank-one norm is at most
one, independently of every \(v,t\); small \(m\) are absorbed
in \(K\). This proves (GL2) on the whole spaces. The physical
and spectral limits have identical compact product pairings,
so are the same operator. Nonvanishing \(j_\infty\), its positive
real-node values, and the same finite interpolation argument
prove its exact indices.

Finally \(|j_\infty(x+iy)|\le K(1+|x|)e^{-\pi|x|}\) on the
closed critical strip. Its node tail above \(T\) costs
\(Ke^{-\pi T/2}\). The unconditional zero-free input used after
(RT9) gives \(1-2|\Im z|\ge2b_0/\log(2+|\Re z|)\).
Below \(T\), (RT9) costs \(Ke^{-2b_0v/\log(T+2)}\), uniformly
in every \(t\ge0\). Take \(T=v/\log(v+e)\), adjusting fixed
constants for small \(v\), to prove the second bound in (GL3).
\end{proof}

\begin{lemma}[Growing moments with an explicitly paid detector norm]
\label{lem:growing-moment-detectors}
Fix one actual hypothetical node \(a=\gamma+i\delta\),
\(0<\delta<1/2\), of arbitrary order \(m_a\).
There are constants \(C,c,R_0>0\), independent of \(m\ge1\),
\(R\ge R_0\), \(t\ge1\), and smooth compact tests \(f_{m,R}\)
supported on \([-R-1,R+1]\), such that
\[
 F_{f_{m,R}}(\pm ik)=0\ (1\le k\le m),\qquad
 F_{f_{m,R}}(\pm i/2)=0,
 \quad \int c|f_{m,R}|^2/p_t\le C a_t^{-1}e^{Cm}.    \tag{GL9}
\]
For \(R=R_0+\log(v+t+m+2)\), along every
\(t\to\infty\), \(t=o(v)\), \(m=o(v)\), the translated
twins have negative Rayleigh depth at least
\[
 c_a a_t e^{-(1-2\delta)v-Cm}                         \tag{GL10}
\]
for both full operators \(J_mA\) and \(J_\infty A\).
The ordinary source cap \(\lceil e^{2v+2R+2}\rceil\) applies
directly to the finite \(J_m\) family. The generally noncompact
\(J_\infty f_{m,R}\) requires a separate outgoing-tail estimate.
\end{lemma}
\begin{proof}
Choose a fixed pole-free PN interpolant \(w\) with values one
at \(a,\bar a\), zero at all other nodes. The full theta series
is analytic on \(|\Im x|<\pi/4\); on every smaller closed
strip it is bounded by \(Ce^{-c e^{2|\Re x|}}\).
In the finite Fourier divisions (OD), use horizontal tail
integrals. Their total integrals vanish on every such line by
rectangular contour shifts with vanishing vertical tails.
Thus all finitely many divisions and differentiations preserve
the strip and its decay. Fixed-radius Cauchy circles give
\[
 |w^{(n)}(x)|\le C n!\rho^{-n}B(x),\quad
 B(x)=e^{-c_\rho e^{2|x|}},\quad 0<\rho<\pi/4.       \tag{GL11}
\]
The positive coefficients \(e_{m,j}\) of \(Q_m\) satisfy
\(0\le e_{m,j}\le\pi^{2j}/(2j+1)!\), by the normally
convergent product and its positive coefficients. Consequently
\(Q_m(-i\partial)w\) has bound \(Ce^{Cm}B\).

To control growing cutoff order, convolve a fixed smooth
probability bump supported on \([-1/4,1/4]\) with \(2m\)
uniform probability densities on \([-1/(8m),1/(8m)]\),
obtaining \(\eta_m\). Set
\(\chi_{R,m}=\mathbf1_{[-R-1/2,R+1/2]}*\eta_m\).
It equals one on \([-R,R]\), vanishes outside \([-R-1,R+1]\),
and \(\|\chi_{R,m}^{(r)}\|_\infty\le(8m)^r\) for
\(0\le r\le2m\). Indeed, place each derivative on a different
uniform factor, whose derivative has total variation \(8m\).
Leibniz and (GL11) then give, for \(n\le2m\),
\[
 |[(\chi_{R,m}-1)w]^{(n)}|
   \le C n!\rho^{-n}e^{8\rho m}B(x)\mathbf1_{|x|\ge R}.
\]
The factorial coefficient bound makes its \(Q_m(-i\partial)\)
sum cost at most \(e^{Cm}\).
Restore the two residual pole values with the compact functions
\(\chi_{1,m}e^{-x^2}\) and \(x\chi_{1,m}e^{-x^2}\), normalized
as in (PC). Their pole matrix has uniform positive lower and
upper bounds, since they equal the fixed Gaussian on \([-1,1]\)
and vanish outside \([-2,2]\). Their derivatives through \(2m\)
have the same Cauchy--Leibniz bound; arbitrary fixed smooth
correctors would not provide this growing-order control.
The residual coefficients are \(O(e^{-c e^{2R}})\).
Call the corrected function \(z_{m,R}\), and set
\(f_{m,R}=Q_m(-i\partial)z_{m,R}\).
All stated moments vanish exactly. The same mass-slice lower
bound used in (SC9) proves, uniformly for every \(t\ge1\),
\[
 \int c|f_{m,R}-Q_m(-i\partial)w|^2/p_t
       \le C a_t^{-1}e^{Cm-c e^{2R}}.                 \tag{GL12}
\]
It also proves (GL9). Twin shifts and the uniform full operator
norm in (GL3) make both form errors at most
\(C c(v)^2a_t^{-1}e^{Cm-c'e^{2R}}\).
For the ideal twins the exact forms are respectively
\[
 -8m_a\Re Q_m(a)\sinh^2(\delta v),\qquad
 -8m_a\Re[j_\infty(a)Q_m(a)^2]\sinh^2(\delta v).
\]
Their coefficients have the uniform positive margin (GL5).
The twin squared norm is at most \(C c(v)^2a_t^{-1}e^{Cm}\).
The stated radius makes the form errors negligible and proves
(GL10). For finite \(J_m\), its differential inverse gives
\(J_m f_{m,R}=z_{m,R}\) with the same compact support, hence
the stated cap by the causal source calculation (SC6)--(SC7).
\end{proof}

Requiring all infinitely many integer moments on one compact
test would give the zero space. Indeed, put \(u=e^x\) on its
compact positive interval. The positive powers have zero integrals
against \(f(\log u)du/u\); polynomials multiplied by \(u\) are
dense among continuous functions there, by approximating \(1/u\).
The measure, and therefore \(f\), is zero. The preceding lemma
uses only finitely many moments at each finite scale, with all
growing-order costs paid. Neither this exact arithmetic limit nor
its unconditional decay reaches the exclusion rate
\(e^{-v+o(v)}\); every fixed off-node remains a separate detectable
negative direction. The missing signed estimate is still open.

\subsection{Paid logistic tails and a linear source-frequency window}
\label{sec:logistic-source-window}

\begin{theorem}[A cofinal ordinary finite-source transfer]
\label{thm:logistic-finite-source}
Let \(m\ge2\), \(R\ge1\), \(\eta\ge1\), and
\(2v>2R+\eta+1\). Let \(f,h\in L^2\) be supported respectively
in \([-v-R,-v+R]\) and \([v-R,v+R]\), with
\(F_f(\pm i/2)=0\) and \(F_h(\pm ik)=0\) for \(1\le k\le m\).
For the actual integer cap
\(N=\lceil e^{2v+2R+\eta+1}\rceil\), set
\[
 \mathcal I_N(f,h)=\frac{i}{2\pi}\int_\R
 j_\infty(u)F_h(u)\overline{F_f(u)}
          \frac{\Xi'(u)}{\mathfrak G(1/2-iu)}D_N(1/2-iu)du.
                                                               \tag{LS1}
\]
It is an absolutely convergent ordinary integral. Uniformly in
all these parameters and every \(t\ge0\),
\[
 |q_{J_\infty A}(f,h)-\mathcal I_N(f,h)|
 \le K(R+1)(m+1)^3e^{R-(m+1/2)\eta}\|f\|_\phi\|h\|_\phi.
                                                               \tag{LS2}
\]
Choose fixed sufficiently large \(C_0,C_1\) and
\[
 \begin{split}
 \eta&=\max\left(1,
  \frac{v+R+4\log(m+1)+2\log(R+1)+C_0}{m+1/2}\right),\\
 H&=\max\left(2,
  \frac{v+R+(\eta+1)/2+5\log(v+R+\eta+2)+C_1}{\pi}\right).
 \end{split}                                                   \tag{LS3}
\]
Whenever the stated ordering holds, (LS1) can then be restricted
to \(|u|\le H\), with total error at most
\(2e^{-v}\|f\|_\phi\|h\|_\phi\) from the full form.
For \(m=\lfloor v/\log^2(v+e^2)\rfloor\), \(R=O(\log v)\),
ordering holds eventually, \(\eta=O(\log^2v)=o(v)\),
\(N=e^{2v+o(v)}\), and \(H=(v+o(v))/\pi\).
\end{theorem}
\begin{proof}
Put \(A_\Phi=\sup e^{|x|/2}\Phi(x)<\infty\).
Gaussian convolution gives \(p_t\le2A_\Phi e^{-|x|/2}\)
for every \(t\ge0\). Splitting \(|x|\le v\) and \(|x|\ge v\)
then gives
\[
 \phi_{v,t}/c\le4A_\Phi e^{-v},\qquad
 \|f\|_2\|h\|_2\le4A_\Phi e^{-v}\|f\|_\phi\|h\|_\phi.
                                                               \tag{LS4}
\]
For \(r\ne0\), the exact absolutely convergent expansion is
\(K_\infty(r)=\sum_{k\ge1}(-1)^{k-1}ke^{-k|r|}\).
Outside \(I_h=[v-R,v+R]\), the first \(m\) convolution modes
vanish by the given moments. Thus, with \(k=J_\infty h\) and
\(d(x)=\dist(x,I_h)\ge1\),
\[
 |k^{(j)}(x)|\le C_j\sqrt R(m+1)^{j+1}e^{-(m+1)d(x)}\|h\|_2
                         \quad(j=0,1,2).              \tag{LS5}
\]
Here \(\|h\|_1\le\sqrt{2R}\|h\|_2\), and factoring out
\(e^{-(m+1)d}\) leaves a uniformly bounded geometric polynomial
sum. No factor \(e^{mR}\) occurs.
Take a cutoff equal to one on
\([v-R-\eta,v+R+\eta]\), vanishing outside its unit enlargement,
with fixed derivative bounds. Denote the cut function by \(k_c\)
and the error by \(e=k-k_c\). Then
\[
 \|e\|_{H^2}\le K\sqrt R(m+1)^3e^{-(m+1)\eta}\|h\|_2.
                                                               \tag{LS6}
\]
Both physical pool products are exactly zero because both moments
of \(f\) are zero. No pool value of \(k_c\) is assumed zero.
The gamma error is at most \(K\|f\|_2\|e\|_{H^2}\), since
\(|a(u)|\le K(1+u^2)\).

For clarity we pay all prime orientations. With \(a=m+1\),
Chebyshev--Abel gives the forward right-tail bound, for \(x\in I_f\),
\[
 \sum_{\log n\ge v+R+\eta-x}\frac{\Lambda(n)}{\sqrt n}
 e^{-a(x+\log n-v-R)}
 \le Ke^{(v+R-x)/2}e^{-(a-1/2)\eta}.                 \tag{LS7}
\]
The forward left-tail sum is empty if its upper logarithmic
threshold is negative; otherwise it is bounded by
\(Ke^{(v-R-x)/2}e^{-(a+1/2)\eta}\).
The backward orientation is at most
\(Ke^{-a(v-R-x)}\), since \(v-R-x>\eta+1\) by ordering and
\(\sum\Lambda(n)n^{-a-1/2}\) is bounded uniformly for \(a\ge3\).
It is also bounded by (LS7). These are genuine Stieltjes bounds,
including endpoint atoms; their exponent ratios are uniformly
bounded. Integrating (LS5) against \(|f(x)|\), and using
\(\|f\|_1\le\sqrt{2R}\|f\|_2\), therefore gives
\[
 |q_\Pi(f,e)|\le KR(m+1)e^{v+R-(m+1/2)\eta}\|f\|_2\|h\|_2.
\]
Together with (LS4) and (LS6), this pays the entire physical
error by the right side of (LS2), with all prime powers retained.

The raw causal calculation (SC6)--(SC7) applies to \(f,k_c\).
Its outward pole is zero because \(F_f(-i/2)=0\); its support
span is at most \(2v+2R+\eta+1\). Thus it gives (LS1) with
\(F_{k_c}\) in place of \(j_\infty F_h\), exactly at the cap
stated above. For \(L^2\) tests use smooth approximants for
\(f\), restore its two moments, and use the \(H^2\) bound on
\(k_c\). The ordinary source and physical mixed form are both
continuous under this approximation. The source difference is
at most
\(K\sqrt N\|f\|_2\|e\|_{H^2}\), using the floor bound
\(|\Xi'/\mathfrak G|\le K(1+u^2)\) and \(|D_N|\le2\sqrt N\).
Since \(\sqrt N\le\sqrt2e^{v+R+(\eta+1)/2}\), (LS4)--(LS6)
give the same bound as (LS2). Absorb both constants into \(K\).

On \(|u|>H\ge2\), \(|j_\infty(u)|\le K(1+|u|)e^{-\pi|u|}\).
The sharper floor bound \(|\Xi'/\mathfrak G|\le
K(1+|u|)\log(2+|u|)\) and Plancherel make this source tail
at most
\[
 K e^{R+(\eta+1)/2}(1+H)^2\log(2+H)e^{-\pi H}
                                             \|f\|_\phi\|h\|_\phi.
\]
The choices (LS3) pay both tails by \(e^{-v}\) each. The original
tests have undergone no frequency projection, so every exact
moment is retained. The last regime follows directly from (LS3).
\end{proof}

The compact detectors of (GL9)--(GL10) belong to this class,
using their actual support radius \(R+1\). All their errors are
paid on a cofinal scale. What remains is the signed norm estimate
for the concrete finite Möbius integral on \(|u|\le(v+o(v))/\pi\).
Neither positivity on most modulation heights nor positivity of a
chosen smaller compression proves that estimate.

\subsection{Complete source matrices with an all-vector error}
\label{sec:complete-source-matrices}

We now replace the remaining infinite-dimensional cross norm by
explicit complete matrices. This replacement preserves every test
in Theorem~\ref{thm:logistic-finite-source}; it is not a choice of
favorable packets.

\begin{theorem}[A polynomial-dimensional complete norm transfer]
\label{thm:complete-source-matrices}
Use the parameters, cap \(N\), and choices of \(\eta,H\) in
Theorem~\ref{thm:logistic-finite-source}. Write
\[
 M_N(u)=j_\infty(u)\frac{\Xi'(u)}{\mathfrak G(1/2-iu)}
                         D_N(1/2-iu),\quad Q=v+R+\eta+2,
\]
and choose an integer
\[
 k\ge 2eHR+4(v+R+\eta+\log Q).
                                                               \tag{CM1}
\]
Set \(a_L(z)=a_{v,t}(-v+z)\), \(a_R(z)=a_{v,t}(v+z)\)
for \(|z|\le R\), and use the exact moment sets
\[
 \mathcal B_L=\{-1/2,1/2\},\qquad
 \mathcal B_R=\{-m,\ldots,-1,1,\ldots,m\}.
                                                               \tag{CM2}
\]
For \(X=L,R\) define the real matrices and vectors
\[
 \begin{split}
 C^X_{\beta\gamma}
   &=\int_{-R}^R a_X(z)^2e^{(\beta+\gamma)z}\,dz,\\
 d^X_{p,\beta}
   &=\int_{-R}^R a_X(z)^2z^pe^{\beta z}\,dz,\\
 G^X_{pq}
   &=\int_{-R}^R a_X(z)^2z^{p+q}\,dz
                    -(d^X_p)^T(C^X)^{-1}d^X_q,\quad 0\le p,q<k,
 \end{split}                                                   \tag{CM3}
\]
and the ordinary finite-integral matrix
\[
 B_{pq}=\frac{i}{2\pi}\frac{i^p(-i)^q}{p!q!}
       \int_{-H}^H u^{p+q}M_N(u)e^{-2ivu}\,du,\qquad
 \mathcal M=(G^L)^{1/2}B(G^R)^{1/2}.                \tag{CM4}
\]
Both \(C^X\) and \(G^X\) are positive definite. If \(\mathcal S\)
is the supremum of
\(|q_{J_\infty A}(f,h)|/(\|f\|_\phi\|h\|_\phi)\)
over all nonzero pairs in that theorem, then, uniformly in
every \(t\ge0\),
\[
                  |\mathcal S-\|\mathcal M\||\le K e^{-v}.
                                                               \tag{CM5}
\]
For its cofinal regime \(R=O(\log v)\),
\(m=\lfloor v/\log^2(v+e^2)\rfloor\), one may take
\(k=O(v\log v)\). The underlying Möbius cap remains
\(N=e^{2v+o(v)}\).
\end{theorem}
\begin{proof}
The two errors already paid in (LS2)--(LS3) are all-vector bounds.
It remains to approximate the complete band kernel in canonical
\(L^2([-R,R])\), which is
\[
 T_H(z,w)=a_L(z)a_R(w)\frac{i}{2\pi}
           \int_{-H}^H M_N(u)e^{-2ivu}e^{iuz-iuw}\,du.
                                                               \tag{CM6}
\]
For real \(s\), the integral Taylor remainder of \(e^{is}\)
after \(k\) terms is at most \(|s|^k/k!\): its exponential
along the real integration segment has modulus one.
On this band the two remainders are each at most
\((HR)^k/k!\le2^{-k}\). Each partial sum has modulus at
most \(1+2^{-k}\); the product remainder is at most \(3\,2^{-k}\).
The real-line floor estimate and exponential \(j_\infty\)
give \(\int_\R|M_N|\le K\sqrt N\). By (LS4),
\(\|a_L\|_2\|a_R\|_2\le KRe^{-v}\). The complete kernel
therefore differs in operator norm by at most
\[
            KRe^{R+(\eta+1)/2}2^{-k}\le K'e^{-v}       \tag{CM7}
\]
from
\(T^{[k]}(z,w)=a_L(z)a_R(w)\sum_{p,q<k}B_{pq}z^pw^q\).
All coefficients are finite integrals; no Taylor series is
interchanged with an uncontrolled infinite frequency integral.

On side \(X\), let \(P_X\) be the orthogonal projection onto the
complement of the real functions \(a_Xe^{\beta z}\),
\(\beta\in\mathcal B_X\). These are precisely the moment
conditions in (CM2); the translation factors \(e^{\pm\beta v}\)
are nonzero scalars. The moment Gram \(C^X\) is positive
definite by independence of the distinct exponentials.
The projected functions \(u_p^X=P_X(a_Xz^p)\) have exactly
the Gram \(G^X\) in (CM3). They are independent: if a
polynomial of degree less than \(k\) equals a combination
of the distinct nonzero exponentials on an open interval,
its \(k\)-th derivative forces every exponential coefficient
to vanish and then the polynomial itself vanishes.
Consequently \(G^X>0\). No bound on a raw moment-matrix
condition number is used.

Let \(U_X:\mathbb C^k\longrightarrow L^2([-R,R])\) be synthesis
by the \(u_p^X\). Then
\(P_LT^{[k]}P_R=U_LBU_R^*\), with \(U_X^*U_X=G^X\).
The finite polar decompositions give isometries \(V_X\) such that
\[
 U_LBU_R^*=V_L(G^L)^{1/2}B(G^R)^{1/2}V_R^*.
\]
Its operator norm is exactly \(\|\mathcal M\|\). The bases and
Grams are real, and the \(i^p(-i)^q\) in (CM4) retains the
antilinear first-slot convention. Compression is contractive;
combining (CM7) with the two preceding all-vector errors proves
(CM5). The stated dimension follows from \(H=O(v)\) and
\(R=O(\log v)\). This is a polynomial vector-space dimension,
not a claim of polynomial total arithmetic computation.
\end{proof}

\begin{corollary}[An exact arithmetic norm target that excludes every off-node]
\label{cor:complete-matrix-rh-transfer}
Suppose, for all sufficiently large \(v\), the complete matrices
(CM4), with \(t=\log(v+e)\), the cofinal moment order above and
each fixed sufficiently large \(R_0\), satisfy
\[
 \|\mathcal M\|\le \exp[-v+o(v)],\qquad
 R=R_0+1+\log(v+t+m+2).                             \tag{CM8}
\]
Then RH holds. The hypothesis (CM8) is not proved here.
\end{corollary}
\begin{proof}
Fix any hypothetical actual node \(a=\gamma+i\delta\),
\(0<\delta<1/2\), of arbitrary multiplicity. Choose its fixed
PN-seed in the proof of Lemma~\ref{lem:growing-moment-detectors}.
Its ideal \(Q_m\)-modified seed has only the two conjugate
node values, both pole values zero. Put
\(A_m=j_\infty(a)Q_m(a)^2\), so \(\Re A_m\ge c_a>0\)
by (GL5). For the ordered copies
\(f(x)=e^{-2i\gamma v}f_{m,R-1}(x+v)\) and
\(h(x)=f_{m,R-1}(x-v)\), the phase is in the antilinear
left slot. The real part of their ideal ordered cross is
\[
                  2m_a\Re A_m\cosh(2\delta v).
                                                               \tag{CM9}
\]
This follows directly from the two terms in (GL7); it uses no
assumption of a rightmost or isolated zero. The corrected
compact copies in (GL12) change this cross by \(o(1)\), since
their whole-operator norm and cutoff errors are uniformly paid.
Their radius is exactly \(R\) above, their left pole moments
and right first \(m\) high moments are zero, and their product
of canonical norms is at most
\(Kc(v)^2a_t^{-1}e^{Cm}=\exp[v+o(v)]\).
The premise and (CM5) would bound that same cross by
\(\exp[o(v)]\). This contradicts (CM9), which is at least
\(c_a e^{2\delta v}\). Every height, every \(\delta>0\), and
every multiplicity was arbitrary. The functional equation
then gives the stated RH conclusion.
\end{proof}

\subsection{Cofinal signature of one complete real arithmetic matrix}
\label{sec:cofinal-real-source-signature}

Retaining the same-side blocks allows a stronger statement than
the preceding cross-norm transfer: one finite real matrix detects
the entire negative index cofinally.

\begin{theorem}[Exact significant counts for the complete source matrix]
\label{thm:cofinal-real-source-signature}
Fix \(R_0\ge2\), set \(t=\log(v+e)\), and take
\[
 m=\lfloor v/\log^2(v+e^2)\rfloor,\quad
 R=R_0+1+\log(v+t+m+2),                             \tag{CS1}
\]
with \(\eta,N,H,B\) as in (LS3),(CM4), and \(k\) the least integer
satisfying (CM1).
On both windows use the symmetric moment set
\(\mathcal B=\{-m,\ldots,-1,-1/2,1/2,1,\ldots,m\}\).
Let \(G\) be (CM3) on the right window for this set, and define
\[
 E_k=\operatorname{diag}((-1)^p)_{0\le p<k},\qquad
                 \mathcal H_v=G^{1/2}E_kBG^{1/2}.     \tag{CS2}
\]
This is a real symmetric \(k\)-by-\(k\) matrix. There are
fixed constants \(C_D,C_E>0\) such that, with
\[
 \tau_v=C_D(R+1)e^{R-v}+2C_Ee^{-v}
                       =e^{-v+O(\log v)},            \tag{CS3}
\]
one has at every sufficiently large finite \(v\)
\[
 n_+(\mathcal H_v-\tau_v I)\le K_+,\qquad
 n_-(\mathcal H_v+\tau_v I)\le K_+,
 \quad
 \lim_{v\to\infty}n_+(\mathcal H_v-\tau_v I)
  =\lim_{v\to\infty}n_-(\mathcal H_v+\tau_v I)=K_+.   \tag{CS4}
\]
For finite \(K_+\), both counts eventually equal \(K_+\) exactly.
For \(K_+=\infty\), each is eventually at least any prescribed
integer; monotonicity in \(v\) is not asserted.
The matrix has \(O(v\log v)\) rows and retains the ordinary cap
\(N=e^{2v+o(v)}\).
\end{theorem}
\begin{proof}
First we verify the exact real fold. On the real line,
\(M_N(-u)=-\overline{M_N(u)}\), since \(\Xi'\) is real and odd
and \(\mathfrak G,D_N\) have real type. If
\(I_n=\int_{-H}^H u^nM_N(u)e^{-2ivu}du\), reflection gives
\(\overline{I_n}=(-1)^{n+1}I_n\).
The coefficient \(i\,i^p(-i)^q\) in (CM4) therefore makes
every \(B_{pq}\) real and gives \(B^T=E_kBE_k\).
Evenness of the actual weight implies \(a_L(z)=a_R(-z)\).
The symmetric moment set is reflection-invariant, and the
left polynomial index \(p\) acquires exactly \((-1)^p\).
Thus the Gram-normalized folded matrix is precisely (CS2).

Let \(X_v\) be the right canonical moment complement and \(P_v\)
its orthogonal projector. With the complete physical kernel
\(\mathcal Q_\infty\) in (GL6), define the two full compressed blocks
\[
 \begin{split}
 S_v&=P_v[a_R(z)\mathcal Q_\infty(2v+z+w)a_R(w)]P_v,\\
 D_v&=P_v[a_R(z)\mathcal Q_\infty(z-w)a_R(w)]P_v.
 \end{split}                                                   \tag{CS5}
\]
The proof of (CM5) applies on every stronger exact moment class:
the kernel error was paid before orthogonal compression. If
\(U:\mathbb C^k\to X_v\) synthesizes the projected polynomials
and \(W=UG^{-1/2}\), then \(W\) is isometric and
\[
 \|S_v-W\mathcal H_vW^*\|\le C_Ee^{-v}.              \tag{CS6}
\]
The finite matrix is extended by zero on the complement of its
polynomial range. This is a bound on the entire \(X_v\).

The full arithmetic bound
\(|\mathcal Q_\infty(r)|\le C_Qe^{|r|/2}\), proved in (GL6),
and (LS4) give
\(\|D_v\|\le C_D(R+1)e^{R-v}\): on a single window
\(|z-w|\le2R\), and the absolute rank-one majorant costs
\(C_Qe^R\|a_R\|_2^2\le C_DRe^{R-v}\).
Every gamma term, prime power, and both original pools remain
in this bound. The even and odd unitary two-window embeddings
are \((+\mathcal R\xi,\xi)/\sqrt2\) and
\((-\mathcal R\xi,\xi)/\sqrt2\), with physical reflection
\(\mathcal R\). Their full forms are exactly \(D_v+S_v\) and
\(D_v-S_v\). Hence eigenvalues of \(\mathcal H_v\) above
\(\tau_v\) give a strictly negative odd subspace; those below
\(-\tau_v\) give a strictly negative even subspace.
The whole parity indices in Theorem~\ref{thm:growing-resolvent-limit}
give the two upper bounds in (CS4).

For the lower bounds select any fixed \(q\) distinct actual
off-line mirrorpairs \(a_j=\gamma_j+i\delta_j\), \(b_j=\bar a_j\),
\(\delta_j>0\), with arbitrary orders \(m_j\).
For each parity \(\epsilon=\pm1\), finite PN interpolation gives
a fixed basis with arbitrary selected samples, all other nodes
and both pools zero. Choose its coefficient data
\[
 F_w(a_j)=
 \frac{\eta_\epsilon c_j e^{\delta_jv}}
      {2\cos(a_jv)Q_m(a_j)},\qquad
 F_w(b_j)=
 \frac{-\eta_\epsilon c_j e^{\delta_jv}}
      {2\cos(b_jv)Q_m(b_j)},\quad
 \eta_+=i,\quad\eta_-=1,                            \tag{CS7}
\]
and their \(\epsilon\)-parity mirror values.
For real \(c_j\) these have the required physical real type;
complex linear extension preserves the Hermitian negative
subspaces. Their basis coefficients are uniformly bounded,
because (GL5) gives a fixed-node lower for \(|Q_m(a_j)|\) and
\[
 |\cos((\gamma_j+i\delta_j)v)|^2
   =\cos^2(\gamma_jv)+\sinh^2(\delta_jv)
                         \ge\sinh^2(\delta_jv).       \tag{CS8}
\]
There is no exceptional phase subsequence.

Put \(f=Q_m(-i\partial)w\) and \(u=f(\,\cdot-v)+f(\,\cdot+v)\).
Then \(F_u(z)=2\cos(zv)Q_m(z)F_w(z)\), and the full node formula
(GL7) gives exactly
\[
 q_{J_\infty A}(u)=
       -4\sum_{j=1}^q m_j\Re j_\infty(a_j)
                         e^{2\delta_jv}|c_j|^2<0.    \tag{CS9}
\]
The coefficient four counts the entire mirrorquartet.
Prepare the bounded basis combination by the cutoff in
(GL12) with inner radius \(R-1\); its parity, both pools and
first \(m\) high moments remain exact. Linearity of that
construction and its fixed-basis estimates give
\[
 \begin{split}
 \|u_c\|_\phi^2
   &\le C_qc(v)^2a_t^{-1}e^{C_qm}\sum_j|c_j|^2,\\
 |q(u_c)-q(u)|
   &\le C_qc(v)^2a_t^{-1}
         e^{C_qm-c_qe^{2(R-1)}}\sum_j|c_j|^2.
 \end{split}                                                   \tag{CS10}
\]
The second is negligible, since \(e^{2(R-1)}\) grows quadratically
in \(v+t+m\). Thus each parity has a \(q\)-dimensional compact
subspace of negative depth at least
\[
 \Delta_{v,q}\ge c_q a_t
       e^{-(1-2\delta_{\min,q})v-C_qm},\qquad
                       \Delta_{v,q}/\tau_v\longrightarrow\infty.
                                                               \tag{CS11}
\]
Only the minimum within this fixed finite family is used.
The negative bound makes the coefficient map injective.

Via (CS5)--(CS6) and the retained \(D_v\) bound, min--max
transports these even and odd spaces to at least \(q\) finite
matrix eigenvalues respectively below \(-\tau_v\) and above
\(\tau_v\). It applies even outside the polynomial range,
because (CS6) holds on the entire space; extension by zero
has the same nonzero eigenvalues as the finite matrix.
For finite \(K_+\), take \(q=K_+\) and use the upper counts;
for \(K_+=0\) those uppers already vanish. In the infinite
case take each \(q\) separately. This proves (CS4) with all
heights and multiplicities retained.
\end{proof}

Consequently either one-sided bound
\(\mathcal H_v\preceq\tau_v I\) or
\(\mathcal H_v\succeq-\tau_v I\) along an unbounded sequence of
scales would prove RH. This is a fully proved implication:
both significant counts have the cofinal limit (CS4).
Neither one-sided arithmetic bound is established here.

\subsection{Prepared infinite moment ranges with a polynomial cofactor}
\label{sec:prepared-nearlinear}

We next reduce the arithmetic cost on each fixed finite detector
basis. The regularizer is fixed throughout; its differential order
does not grow with the support scale. Put
\[
 \alpha_\ell=8\ell\log^2(\ell+1),\qquad
 P_{\mathrm n}(z)=\prod_{\ell\ge1}(1+z^2/\alpha_\ell^2)
       =\sum_{r\ge0}c_rz^{2r},\qquad j_{\mathrm n}=P_{\mathrm n}^{-1}.
                                                               \tag{NL1}
\]
The norm remains the actual two-copy norm
\(\|f\|_\phi^2=\int c|f|^2/\phi_{v,t}\), with \(p_t,a_t\)
as in the preceding sections. All constants below may depend on
a fixed finite PN basis, including its selected heights and
multiplicities. No uniform bound in a growing selected-node family
is asserted.

\begin{lemma}[A fixed local regularizer with an infinite compact range]
\label{lem:nearlinear-compact-range}
The product in (NL1) is entire and has exactly the simple zeros
\(\pm i\alpha_\ell\). For one fixed \(B_0\), every \(r\ge0\) satisfies
\[
 0<c_r\le
 \frac{B_0^{2r}}{(2r+1)!\log^{4r}(2r+e)}.                       \tag{NL2}
\]
There is a positive, smooth probability kernel \(K_{\mathrm n}\),
with Fourier transform \(j_{\mathrm n}\), such that
\[
 |K_{\mathrm n}^{(q)}(x)|\le C_qe^{-\alpha_1|x|}\quad(q\ge0),
 \qquad r_{\mathrm n}=\int e^{x/2}K_{\mathrm n}(x)\dd x
                        =j_{\mathrm n}(i/2).
                                                               \tag{NL3}
\]
For every compact interval \(I\) and \(h\in L^2\) supported in \(I\),
\[
 \operatorname{supp}(J_{\mathrm n}h)\subset I
 \quad\Longleftrightarrow\quad F_h(\pm i\alpha_\ell)=0
                 \text{ for every }\ell,\qquad
 J_{\mathrm n}h=K_{\mathrm n}*h.                                \tag{NL4}
\]
This compact range is nonzero. On the complete original form \(A\),
the regularized form has the zero expansion
\[
 q_{J_{\mathrm n}A}(f,h)
   =\sum_{z\in\mathcal Z}m_zj_{\mathrm n}(z)
                  \overline{F_f(\bar z)}F_h(z).                 \tag{NL5}
\]
Its whole negative index is \(2K_+\), its parity indices are \(K_+\),
and compact tests satisfying (NL4) and both original pole moments
attain the same suprema, including infinity.
\end{lemma}
\begin{proof}
Both \(\sum\alpha_\ell^{-1}\) and \(\sum\alpha_\ell^{-2}\)
converge, and
\(\sum\alpha_\ell^{-1}<1\), by integral comparison with
\(1/(x\log^2x)\). For an \(r\)-element index set, at most
\(\sqrt r\) indices are below \(\sqrt r\). For \(r\ge4\), therefore,
with \(L=\tfrac12\log r\) and \(L_r=\log(2r+e)\),
\[
 \prod_{\ell\in I}\log^{-4}(\ell+1)
 \le L^{-4r}(L/\log2)^{4\sqrt r}
 \le L_r^{-4r}4^{4r}e^{4r}.
\]
Here \(L_r/L\le4\) and
\(\log\max(1,L/\log2)\le\sqrt r\). The primary sinh product
\cite[4.36.1]{DLMF} gives
\(\sum_{|I|=r}\prod_{\ell\in I}\ell^{-2}
 =\pi^{2r}/(2r+1)!\).
Together with the factor \(8^{-2r}\) this proves (NL2);
\(B_0=64\pi\) also covers \(r=0,1,2,3\).
All sums first concern finite subsets and then increase monotonically.

Let \(E_\alpha(x)=(\alpha/2)e^{-\alpha|x|}\).
Independent variables with these densities have an almost surely
absolutely convergent sum, since
\(\sum\mathbb E|X_\ell|=\sum\alpha_\ell^{-1}<\infty\).
Its characteristic function is \(j_{\mathrm n}\).
Convolving a fixed finite prefix with the remaining probability
law proves smoothness to every fixed order and (NL3):
the remaining law has a finite exponential moment at \(\alpha_1\).
Strict positivity follows from the first density.
A fixed exponential moment between \(1/2\) and \(\alpha_1\)
also pays uniform integrability in the halfmoment identity.

The finite kernels have residues
\[
 A_{\ell,M}=\frac{\alpha_\ell}{2}
       \prod_{\substack{k\le M\\k\ne\ell}}
                     (1-\alpha_\ell^2/\alpha_k^2)^{-1}.
\]
Comparison with the sine product
\cite[4.22.1]{DLMF}, differentiated at its simple integer zeros,
gives \(0<|A_{\ell,M}|\le|A_\ell|\le\alpha_\ell\).
Indeed \(\alpha_\ell/\alpha_k\ge\ell/k\) for \(k<\ell\),
and the reverse inequality holds for \(k>\ell\); the linear
absolute reciprocal product is exactly two.
Thus on \(|x|\ge d>0\),
\[
 K_{\mathrm n}(x)=\sum_{\ell\ge1}A_\ell e^{-\alpha_\ell|x|}
\]
converges absolutely and uniformly. No absolute residue sum at
zero is required. Outside \(I=[a,b]\), the convolution on the right
is \(\sum A_\ell e^{-\alpha_\ell x}F_h(i\alpha_\ell)\).
Its majorant is \(\|h\|_1\sum\alpha_\ell
 e^{-\alpha_\ell(x-b)}\).
Moment vanishing makes it zero; conversely, multiplication by
the successive \(e^{\alpha_\ell x}\) and \(x\to\infty\)
isolates each nonzero coefficient. The left side proves the
negative moments. This proves (NL4).

For completeness, the local class is explicit. Set \(\beta=3/2\),
\[
 S=\sum_{\ell\ge1}\frac1{\ell\log^\beta(\ell+e)},\qquad
 b_\ell=\frac1{2S\ell\log^\beta(\ell+e)}.
\]
The infinite convolution of uniform densities on
\([-b_\ell/2,b_\ell/2]\) has compact support and derivatives
\[
 \|\eta^{(n)}\|_\infty\le CB_1^n n!\log^{\beta n}(n+e).          \tag{NL6}
\]
Distributional differentiation of the first \(n\) factors costs
\(\prod_{\ell\le n}2/b_\ell\); the next density supplies
\(b_{n+1}^{-1}\), and the remaining convolution cannot enlarge
the supremum. The extra factor is absorbed in \(CB_1^n\).
The next derivative bound makes each derivative Lipschitz,
so this is an actual smooth representative.
For every compact function \(g\) with (NL6),
\[
 P_{\mathrm n}(-i\partial)g
     =\sum c_r(-1)^rg^{(2r)},\qquad
 F_{P_{\mathrm n}(-i\partial)g}=P_{\mathrm n}F_g,\qquad
 J_{\mathrm n}P_{\mathrm n}(-i\partial)g=g.                     \tag{NL7}
\]
The derivative series converge normally, including every fixed
extra derivative, by
\(\sum[C/\log^{2-\beta}(n+e)]^n<\infty\).
Every summand has the same support. Taking \(g=\eta\)
proves that the compact range is nonzero.

On \(|\Im z|\le1/2\), the product arguments have the same sign
and total absolute value at most
\(2\sum\arcsin(|\Im z|/\alpha_\ell)<\pi/2\).
Consequently
\[
 \Re P_{\mathrm n}(z)\ge c_*|P_{\mathrm n}(z)|,\quad
 \Re j_{\mathrm n}(z)\ge c_*|j_{\mathrm n}(z)|,\quad
 |j_{\mathrm n}(x+iy)|\le Cj_{\mathrm n}(x).                    \tag{NL8}
\]
The modulus estimate follows from
\(\Re(1+z^2/\alpha_\ell^2)\ge
 (1-y^2/\alpha_\ell^2)(1+x^2/\alpha_\ell^2)\)
and the convergent reciprocal product.
For large \(|u|\), each of the first
\(\lfloor |u|/(32\log^2(|u|+e))\rfloor\) real factors
is at least 17. Hence
\[
 j_{\mathrm n}(u)\le C
       \exp[-c|u|/\log^2(|u|+e)].                              \tag{NL9}
\]
The smoothed gamma kernel, in the normalization of (RGC4), is
\[
 \gamma_{\mathrm n}=a(0)K_{\mathrm n}
 +\int_0^\infty\frac{e^{-s/2}}{1-e^{-2s}}
       [2K_{\mathrm n}-K_{\mathrm n}(\cdot+s)
                         -K_{\mathrm n}(\cdot-s)]\dd s.
\]
Its second difference pays the singularity at zero. Removing
the exact first gamma branch gives
\(\gamma_{\mathrm n}^{\rm hi}
 =\gamma_{\mathrm n}+r_{\mathrm n}e^{-|\cdot|/2}
 =O(e^{-2|\cdot|})\). Thus the complete physical kernel is
\[
 Q_{\mathrm n}=\gamma_{\mathrm n}^{\rm hi}
       -K_{\mathrm n}*\Pi+r_{\mathrm n}e^{|\cdot|/2}.            \tag{NL10}
\]
Chebyshev--Abel pays its prime tail by \(Ce^{|x|/2}\),
with every \(\Lambda(n)\) retained.
The same cutoff proof of the complete smoothing formula as in
(SC4), now using (NL3), proves (NL5).
The multiplicity-weighted height count and (NL8)--(NL9) pay
absolute trace-class convergence on the whole canonical space.

Finally the PN interpolants constructed earlier are analytic in
a fixed horizontal strip. On \(|z-x|\le4/100\), the full theta
series has bound
\[
 Ce^{M|x|-c_\epsilon e^{2|x|}},\qquad
 c_\epsilon=\pi e^{-8/100}\cos(8/100)>\pi/2 .
\]
Finite Fourier division solves the local first-order analytic ODE;
its integral solution preserves this bound, using the already
proved real theta tail. Fixed polynomial derivatives are paid on
nested circles before applying Cauchy. Therefore, on each fixed
finite PN basis,
\[
 |w^{(n)}(x)|\le Cn!\rho^{-n}
                      e^{M|x|-c_\epsilon e^{2|x|}},\qquad
                         \rho=1/100.                          \tag{NL11}
\]
Equations (NL2) and (NL11) make \(P_{\mathrm n}(-i\partial)w\)
converge with this same tail and with all fixed derivatives.
It preserves every selected node and every other zero value.
Apply the compact cutoffs (NL6), correct both pole moments by
an even bump and its product with \(x\), and use (NL7).
Their error has the same theta tail outside the cutoff and
size \(Ce^{-d e^{2R}}\) on the fixed correction support.
A positive fixed mass slice gives
\(p_t^{-1}\le Ca_t^{-1}e^{x^2/4+C|x|}\) for \(t\ge1\).
For each fixed \(v,t=0\), the reciprocal theta bound costs
at most \(e^{\pi e^{2|x|}+C_v|x|}\), paid by \(2c_\epsilon>\pi\).
Thus the compact tests converge in every fixed ground norm.
The node formula then gives every finite negative detector
subspace in this exact moment class. Real nodes are positive
rank ones; each non-real conjugate pair contributes one
negative direction, and the mirror gives one per parity.
Taking arbitrary finite families proves the stated whole
and compact indices without assuming finitely many off-line zeros.
\end{proof}

\begin{theorem}[Prepared ordinary sources with \(K=O(v^4)\) and \(H=O(v)\)]
\label{thm:prepared-nearlinear-source}
Fix a finite PN basis \(w_1,\ldots,w_d\) with independent selected
actual-node data, zero data at every other Xi node, and both
original pole moments zero. For every fixed \(A>0\) there are
fixed constants \(B,C_A\) and linear compact prepared functions
\(f_{R,p}[a]\), for coefficient vectors \(a\in\C^d\), such that
\[
 p=\lfloor v/2\rfloor,\quad
 R=\tfrac12\log(v+p+t+2)+B,\quad W=R+1,\quad 1\le t\le v,        \tag{NL12}
\]
and all sufficiently large \(v\) have the following properties.
Every function satisfies all moments (NL4) and both pole
moments exactly. Its limiting main function
\(f_\infty[a]=P_{\mathrm n}(-i\partial)\sum a_jw_j\) is fixed,
and the translated norms satisfy
\[
 c\,c(v)a_t^{-1/2}|a|
 \le\|\tau_{\pm v}f_{R,p}[a]\|_\phi
 \le C\,c(v)a_t^{-1/2}|a|.                                   \tag{NL13}
\]
For \(f=\tau_{-v}f_{R,p}[a]\), \(h=\tau_vf_{R,p}[b]\), set
\(C(r)=\int\overline{f(x)}(J_{\mathrm n}h)(x+r)\dd x\) and
\[
 \begin{split}
 N&=\lceil e^{2v+2W}\rceil,\\
 E_p&=\max\{1,C_0\,2^{p-2}(p+1)!(C_1p)^p\},\\
 D&=e^{(2-1/p)(v-W)}
                  [E_pW\log(2N)]^{-1/p}.
 \end{split}                                                   \tag{NL14}
\]
Then \(1<D<N\), \(K_{\rm short}=\lfloor N/D\rfloor\le C_Bv^4\),
and the actual divisor head has
\[
 |T_{\le D}(C)|\le e^{-v}\|f\|_\phi\|h\|_\phi.                 \tag{NL15}
\]
Both complete ordinary finite sources
\[
 \begin{split}
 q_{J_{\mathrm n}A}(f,h)
  &=\frac{i}{2\pi}\int_\R j_{\mathrm n}(u)F_h(u)
          \overline{F_f(u)}
          \frac{\Xi'(u)}{\mathfrak G(1/2-iu)}
                         D_N(1/2-iu)\dd u,\\
 T_{>D}(C)
  &=-\frac1{2\pi}\int_\R j_{\mathrm n}(u)F_h(u)
        \overline{F_f(u)}
        \zeta(1/2-iu)M^{\log}_{D,N}(1/2-iu)\dd u
 \end{split}                                                   \tag{NL16}
\]
can be restricted simultaneously to \(|u|\le H=C_Av\), with
error at most \(e^{-Av}\|f\|_\phi\|h\|_\phi\).
Here \(D_N(s)=\sum_{d\le N}\mu(d)d^{-s}\) and
\(M^{\log}_{D,N}(s)=\sum_{D<d\le N}\mu(d)\log d\,d^{-s}\).
The long signed sum, including its cofactor-one channel and
every prime power, is retained.
\end{theorem}
\begin{proof}
Rescale the fixed bump (NL6) into \([-1/8,1/8]\) with positive
variance. Let \(U_p\) be uniform on \([-1/(8p),1/(8p)]\) and
\(\eta_p=U_p^{*p}*\eta\). It has support in \([-1/4,1/4]\) and
\[
 \|\eta_p^{(q)}\|_\infty\le C(Cp)^q\ (q\le p),\qquad
 \|\eta_p^{(p+s)}\|_\infty
       \le C(Cp)^pB_1^ss!\log^{\beta s}(s+e).                  \tag{NL17}
\]
Differentiate separate uniform factors first, each at total
variation cost \(8p\), then the fixed bump. The integrated
switches define \(\chi_{R,p}=1\) on \([-R,R]\), zero outside
\([-R-1,R+1]\). Their transitions are disjoint, so (NL17)
also bounds their derivatives. The correctors
\(\eta_p,x\eta_p\) obey these first-\(p\) bounds and the
all-order bound \(Ce^{Cp}B_2^nn!\log^{\beta n}(n+e)\).
For \(n=p+s\), use \(n!\ge p!s!\) and
\((Cp)^p/p!\le e^{C'p}\); for \(n\le p\), use
\((Cp)^n/n!\le e^{Cp}\).

Their two halfmoment matrix is uniformly invertible:
\(\int c\eta_p\ge1\) and
\(\int x\sinh(x/2)\eta_p\ge\frac12\operatorname{Var}(\eta)>0\);
the off-parity entries vanish. Set
\[
 z_{R,p}[a]=\chi_{R,p}w[a]-b_0[a]\eta_p-b_1[a]x\eta_p,\qquad
 f_{R,p}[a]=P_{\mathrm n}(-i\partial)z_{R,p}[a].                \tag{NL18}
\]
The coefficients enforce both original moments of \(z_{R,p}\)
and satisfy \(|b_j[a]|\le C|a|e^{-d e^{2R}}\).
Leibniz, (NL11), and (NL17) give, for every derivative order \(n\),
\[
 |[(\chi_{R,p}-1)w[a]]^{(n)}|
 \le C|a|e^{Cp}B_3^nn!\log^{\beta n}(n+e)
       e^{M|x|-c_\epsilon e^{2|x|}}1_{\{|x|\ge R\}} .
\]
Thus (NL2) pays the complete local series, with every fixed
derivative. Its error is \(e^{Cp-d e^{2R}}\) in \(L^1,L^2\)
and, with factor \(a_t^{-1/2}\), in the central weighted norm.
The large constant \(B\) in (NL12) pays any prescribed exponential
error, also after translation and application of the bounded
whole operator. The \(e^{Cp}\) factor occurs only in this error;
the main function is unchanged.

The independent main functions have a strictly positive Gram on
some fixed slice \([-T,T]\), by convergence to their full
positive-definite \(L^2\) Gram. The paid error preserves that
lower bound. Since \(p_t\le a_t\Xi(0)\),
the translated slice gives the lower bound in (NL13).
The upper follows from the fixed mass-slice bound and
\(\phi_{v,t}(x\pm v)\ge p_t(x)/(2c(v))\).
This proves (NL13) for every coefficient vector, rather than
only separately for the basis elements.

Now \(J_{\mathrm n}h=\tau_vz_{R,p}[b]\) exactly.
Put correlation derivatives on this latter function.
Its first \(p\) jets have bound \(C|b|(C_1p)^q\), by (NL11),
(NL17), and the adaptive correctors. Consequently
\[
 \max_{q\le p}\|C^{(q)}\|_\infty
 \le C(C_1p)^p|a||b|
 \le C(C_1p)^pe^{-v}\|f\|_\phi\|h\|_\phi.                    \tag{NL19}
\]
The last step uses the genuine lower norm (NL13).
The correlation has support in \([2v-2W,2v+2W]\), all endpoint
jets vanish, and its outward halfmoment is zero.
Repeated Euler integration on this compact smooth function,
using \(\Lambda=-(\mu\cdot\log)*1\), bounds its divisor head by
\[
 2^{p-2}(p+1)!(4W)D^p\log(2D)
 e^{-(p-1/2)(2v-2W)}\max_{q\le p}\|C^{(q)}\|_\infty .
\]
This is the same complete Euler head identity used above,
with both endpoints retained before their proven vanishing.
Choose \(C_0\) in (NL14) sufficiently large. Inserting \(D\)
cancels its full prefactor and the support exponent, proving
(NL15). Moreover \(E_p^{1/p}\le Cp^2\), and
\[
 N/D\le 2e^{4W+(v-W)/p}
                  [E_pW\log(2N)]^{1/p}\le C_Bv^4.
                                                               \tag{NL20}
\]
The positive main exponent \(2v-O(\log v)\) also proves \(1<D<N\).

It remains to pay the complete multiplier in the Fourier tails.
On the real axis the sinh product, split at \(\ell=\sqrt{|u|}\),
gives
\[
 1\le P_{\mathrm n}(u)\le
       C\exp[C|u|/\log^2(|u|+e)].
\]
The low factors cost \(C\sqrt{|u|}\log(|u|+e)\); the remaining
factors are bounded by the linear sinh product with parameter
\(|u|/(2\log^2|u|)\). Fixed-strip contour shift in (NL11)
gives \(|F_{w[a]}(u)|\le C|a|e^{-\rho_1|u|}\).
Every derivative of \(e_{R,p}=z_{R,p}-w\) has \(L^1\) bound
\(C|a|\mathcal E_v B_4^nn!\log^{\beta n}(n+e)\), where
\(\mathcal E_v=e^{Cp-d e^{2R}}\).
Integration by parts at
\(n=\lfloor |u|/(4B_4\log^\beta(|u|+e))\rfloor\) proves
\[
 |F_{e_{R,p}[a]}(u)|\le
 C|a|\mathcal E_v e^{-c|u|/\log^\beta(|u|+e)} .
                                                               \tag{NL21}
\]
Because \(\beta<2\), this dominates the growth of \(P_{\mathrm n}\).
For \(w_0(u)=(1+|u|)\log(2+|u|)\),
\[
 \int_{|u|>H}w_0(u)P_{\mathrm n}(u)
    |F_{z[a]}(u)F_{z[b]}(u)|\dd u
 \le C|a||b|\,[e^{-\rho_2H}+\mathcal E_v+\mathcal E_v^2].
                                                               \tag{NL22}
\]
The mixed and error terms are paid on their entire frequency
line by the explicit integrable majorants in (NL21).

The exact compact ordinary completion gives (NL16); the original
pool product vanishes by the two separately imposed moments.
On the real line the unconditional one-step Euler continuation
and completion bounds give
\[
 |\zeta(1/2-iu)|\le C(1+|u|),\quad
 |\Xi'/\mathfrak G|\le Cw_0(u),\quad
 |D_N|\le2\sqrt N,\quad |M^{\log}_{D,N}|\le2\sqrt N\log N.
\]
The exact prepared product is
\(j_{\mathrm n}F_h\overline{F_f}
 =e^{-2iuv}P_{\mathrm n}F_{z[b]}\overline{F_{z[a]}}\).
After (NL13), the remaining normalized prefactor is at most
\(Ca_te^W\log(2N)=O_B(v^{3/2})\).
First enlarge \(B\) to make \(\mathcal E_v\le e^{-(A+3)v}\),
then take \(C_A>(A+3)/\rho_2\). Equation (NL22) pays both
tails at \(H=C_Av\). Every Dirichlet sum remains finite, with
the same exact cap \(N\); no outgoing arithmetic tail is discarded.
\end{proof}

To use the theorem, choose any individual hypothetical pair
\(b=\gamma+i\delta,\bar b\), \(0<\delta<1/2\), of order \(m\).
PN interpolation gives \(F_w(b)=F_w(\bar b)=1\), all other
node values and both pools zero. By (NL5),
\[
 q_{J_{\mathrm n}A}
   (\tau_vf_\infty-e^{-2i\gamma v}\tau_{-v}f_\infty)
       =-8m\Re P_{\mathrm n}(b)\sinh^2(\delta v).                \tag{NL23}
\]
The two separate diagonals stay bounded; the phased ordered cross
has real part \(2m\Re P_{\mathrm n}(b)\cosh(2\delta v)\).
The prepared errors are exponentially small. The high gamma
cross from (NL10), the zero outward pool product, and (NL15)
therefore force the actual long signed cross to have normalized
size at least
\[
 c_ba_te^{-(1-2\delta)v},\qquad
                 t=1\ \text{or }t=o(v),\ t\ge1.                \tag{NL24}
\]
This uses an arbitrary fixed node and arbitrary multiplicity,
without a rightmost zero, a finite exceptional set, or a
growing main-seed norm.

Fix the prepared construction with \(A=2\), and its sufficiently
large constants \(B,C_2\). Consequently the following implication
is proved: if, for each
fixed finite PN basis, its actual low-frequency long integral
in (NL16) satisfies
\[
 |T_{>D}^{\,|u|\le C_2v}(f,h)|
       \le e^{-v+o(v)}\|f\|_\phi\|h\|_\phi,\qquad t=1,           \tag{NL25}
\]
uniformly in its coefficient vectors, then RH holds.
Indeed (NL13) makes the assumed unnormalized upper \(e^{o(v)}\),
whereas every fixed \(\delta>0\) forces \(e^{2\delta v}\).
The condition (NL25) remains unproved. The polynomial cofactor,
linear frequency window, exact infinite moment range, and bounded
main-seed cost are proved theorems; they do not estimate the
remaining signed long Möbius sum.

\subsection{Uniform Fourier division in the actual central heat graph}
\label{sec:central-heat-division}

The next theorem changes the domain on which division is controlled.
It uses the single central heat ground, rather than assuming convexity
of the two-copy ground in its intervening valley.

\begin{theorem}[Global central heat division and the full logarithmic derivative]
\label{thm:central-heat-division}
There is a finite constant \(T_\Phi\ge1\), depending only on the
complete theta density, such that for every \(t\ge T_\Phi\)
\[
 (\log p_t)''(x)\le-\frac1{4t}\quad(x\in\R).
                                                               \tag{HD1}
\]
Put \(W_t=c/p_t\), \(\mathcal H_t=L^2(W_t\dd x)\), and let
\(\mathcal D_t=-i\partial\) on its maximal domain
\[
 \operatorname{Dom}(\mathcal D_t)
 =\{u\in\mathcal H_t:\ u\text{ is locally absolutely continuous},
                               \ u'\in\mathcal H_t\}.
\]
This operator is closed and \(C_c^\infty(\R)\) is a graph core.
For every \(a\in\C\), without a restriction on its height or
imaginary part,
\[
 \|(\mathcal D_t-a)u\|_{\mathcal H_t}
          \ge\frac{\|u\|_{\mathcal H_t}}{2\sqrt t},\qquad
 \operatorname{ran}(\mathcal D_t-a)=\ker F(a).                  \tag{HD2}
\]
Hence each \(f\in\mathcal H_t\) with \(F_f(a)=0\) has unique
division \(B_af\in\operatorname{Dom}(\mathcal D_t)\), satisfying
\[
 F_{B_af}(z)=\frac{F_f(z)}{z-a},\qquad
              \|B_af\|_{\mathcal H_t}\le2\sqrt t\|f\|_{\mathcal H_t}.
                                                               \tag{HD3}
\]
The Fourier identity is entire, with its removable value at \(a\).

Let \(\mathcal Z_+\) contain all distinct Xi nodes with positive
real part, including both conjugate representatives when different,
and let \(C_2=\sum_{a\in\mathcal Z_+}m_a/|a|^2<\infty\).
For \(f\in\operatorname{Dom}(\mathcal D_t^2)\) vanishing at every
Xi node, the series
\[
 \mathcal M_\Xi f=\sum_{a\in\mathcal Z_+}m_a(B_a+B_{-a})f
                                                               \tag{HD4}
\]
converges absolutely in \(\mathcal H_t\), with
\[
 F_{\mathcal M_\Xi f}=(\Xi'/\Xi)F_f,\qquad
 \|\mathcal M_\Xi f\|_{\mathcal H_t}
 \le2C_2\|\mathcal D_tf\|_{\mathcal H_t}
          +4\sqrt t\,C_2\|\mathcal D_t^2f\|_{\mathcal H_t}.
                                                               \tag{HD5}
\]
Every height and multiplicity is included. Vanishing at all
nodes is a zero-ideal hypothesis; it places no nodes on the real axis.
\end{theorem}
\begin{proof}
We first prove the required global heat curvature from the actual
theta series, without a premise of global logconcavity of \(\Phi\).
For \(x\ge0\), \(T=e^{2x}\), factor it as
\(\Phi(x)=e^{9x/2-\pi T}A(T)\), where
\[
 A(T)=4\pi^2-6\pi/T+
 \sum_{n\ge2}(4\pi^2n^4-6\pi n^2/T)e^{-\pi(n^2-1)T}.
\]
The remainder and its first two \(x\) derivatives are
\(O(T^je^{-3\pi T})\), \(j=0,1,2\): factor \(e^{-3\pi T}\)
and bound the remaining polynomially weighted series at \(T=1\).
For \(V=-\log\Phi\), smooth and even by the Jacobi identity,
\[
 V'(x)=2\pi e^{2x}-9/2+O(e^{-2x}),\qquad
 V''(x)=4\pi e^{2x}+O(e^{-2x}) .
                                                               \tag{HD6}
\]
Choose fixed \(R>2\) such that \(V''\ge1\) for \(|x|\ge R-1\)
and \(V'(R-1)>R\). Choose an even smooth switch \(\chi\),
zero on \(|x|\le R-1\), one on \(|x|\ge R-1/2\), and set
\(q_0=1+\chi(V''-1)\). It is at least one, equals \(V''\)
outside the transition, and
\[
 \int_0^Rq_0\le R-1+V'(R)-V'(R-1)<V'(R).
\]
Add a positive multiple of a fixed even nonnegative bump supported
in \((-R+1,R-1)\) to make this integral \(V'(R)\).
Integrating twice with derivative zero at zero and value \(V(R)\)
at \(R\) produces a smooth even \(U\), with \(U''\ge1\)
everywhere and \(U=V\) on \(|x|\ge R\).
Define the finite constants
\[
 B_\Phi=\sup(V-U)-\inf(V-U),\qquad C_\Phi=e^{B_\Phi},\qquad
                         T_\Phi=\max(1,C_\Phi).
\]

For any real tilt \(\theta\) and \(t>0\), compare the posterior
density proportional to
\(\Phi(u)e^{-u^2/(4t)+\theta u}\) with the density proportional
to \(e^{-U(u)-u^2/(4t)+\theta u}\).
The latter potential \(H\) satisfies \(H''\ge\alpha=1+1/(2t)\).
If \(m\) is its mean, strong monotonicity and integration by parts give
\[
 \alpha\operatorname{Var}(u)
 \le\int(u-m)[H'(u)-H'(m)]\rho(u)\dd u=1 .
\]
All boundary terms vanish by the proven tails.
The actual/model normalized density ratio lies between
\(e^{-B_\Phi}\) and \(e^{B_\Phi}\).
Using the model mean as a comparison constant in the actual
variance shows that the actual variance is at most \(C_\Phi\),
uniformly in \(\theta,t\).

Differentiating the positive heat integral is justified locally
by the complete double-exponential theta majorant and yields
\[
 (\log p_t)''(x)=-\frac1{2t}
    +\frac{\operatorname{Var}_{t,x/(2t)}(u)}{4t^2}.
\]
This proves (HD1). If \(V_t=\log W_t\) and \(\kappa_t=1/(4t)\),
evenness and \((\log c)''\ge0\) give
\[
 V_t''\ge\kappa_t,\qquad
               W_t(x)\ge W_t(0)e^{\kappa_t x^2/2}.              \tag{HD7}
\]
Every complex Fourier evaluation is therefore bounded:
its squared norm is \(\int e^{2\Im(a)x}/W_t(x)\dd x<\infty\).
The same majorants make each Fourier transform entire.

Closedness of the maximal derivative follows by local
distributional convergence. Multiplying a domain vector by
cutoffs \(\chi_R\) with \(\|\chi_R'\|_\infty\le C/R\)
converges in its graph norm; subsequent compact mollification
does so by the positive local upper and lower weight bounds.
This proves the asserted core.
Under the unitary map \(u\mapsto h=\sqrt{W_t}u\), for
\(a=\gamma+i\delta\), the operator on that core becomes
\[
 A=-i\partial-\gamma+iq,\qquad q=V_t'/2-\delta .
\]
Its formal adjoint on the core is
\(A^\#=-i\partial-\gamma-iq\), and direct integration gives
\[
 \|Ah\|_2^2-\|A^\#h\|_2^2
       =\int V_t''|h|^2\dd x\ge\kappa_t\|h\|_2^2.
\]
The graph core extends this lower bound to the maximal domain.
It is uniform in both \(\gamma\) and \(\delta\), and makes
the range closed and the operator injective.
Its range annihilator is obtained distributionally: if \(g\)
annihilates it, \(b=\bar gW_t\) solves \(ib'-ab=0\).
Thus the annihilator is precisely the one-dimensional span of
\(g_a=W_t^{-1}e^{i\bar ax}\), which belongs to \(\mathcal H_t\)
by (HD7). This proves (HD2)--(HD3).
Fourier integration by parts is legal because every derivative
in the domain has all exponential \(L^1\) moments by (HD7);
cutoff derivative errors vanish in each such moment.
It gives the entire division identity, rather than only a
real-frequency equality.

The unconditional paired Hadamard formula, already used in
(EUL10), is
\[
 \frac{\Xi'(z)}{\Xi(z)}
  =\sum_{a\in\mathcal Z_+}m_a
                  \left(\frac1{z-a}+\frac1{z+a}\right).
                                                               \tag{HD8}
\]
The order-at-most-one factorization \cite[Theorem 7.7]{Petrow},
the evenness, \(\Xi(0)>0\), and the full height count give
\(C_2<\infty\); pairing removes the residual linear exponential.
No imaginary-axis node is omitted: on the real critical strip,
the alternating eta series is strictly positive and
\(\zeta(s)=\eta(s)/(1-2^{1-s})<0\), \(0<s<1\).
Thus \(\xi(s)>0\) there. Formula (HD8) converges normally away
from its actual poles, with all orders retained.

For an input in the stated zero ideal, both derivatives also
vanish at every node. The entire Fourier identities verify
the exact weighted-domain relation
\[
 B_af=-f/a-\mathcal D_tf/a^2+B_a\mathcal D_t^2f/a^2 .
\]
Pairing \(a,-a\) cancels the first term, while (HD3) gives
\[
 \|(B_a+B_{-a})f\|_{\mathcal H_t}
 \le |a|^{-2}
       [2\|\mathcal D_tf\|_{\mathcal H_t}
                  +4\sqrt t\|\mathcal D_t^2f\|_{\mathcal H_t}].
\]
Summing proves absolute Hilbert-space convergence in (HD4).
Bounded complex evaluations pass this limit locally uniformly;
(HD8) identifies its entire transform, including all removable
values. This proves (HD5).
\end{proof}

These division costs apply directly to fixed PN seeds.
For a fixed polynomial \(H\), the genuine theta profiles and
the positive mass slice give
\(\|\mathcal D_t^q H(\partial)\Phi\|_{\mathcal H_t}
 \le C_{H,q}/\sqrt{a_t}\) for every fixed \(q\), uniformly
for \(t\ge T_\Phi\).
Legal division commutes with those derivatives: its divided
Fourier transform and the uniqueness in (HD3) prove the claim,
and \(\mathcal D_tB_af=f+aB_af\) places it in the next domain.
A fixed chain of \(m\) divisions thus costs at most
\((2\sqrt t)^m\) in each fixed graph norm.
At \(t=\log^3v\) this extra fixed-factor cost is \(v^{o(1)}\).
This is a product cost; the separate finite logarithmic-derivative
sum costs only its number of terms.

The source-transfer condition remains distinct. To apply (HD5)
to a completed finite M\"obius cofactor, one must establish its
membership in this central domain, its zero-ideal property,
and the required two derivative norms there.
Local signed theta inversion or a bounded projected scalar pairing
does not establish these global weighted norms.
For example, translations of \(\Phi\) preserve its real Fourier
modulus and all zero values, while (HD7) makes their central
norms grow with the translation. This theorem proves division
and its graph bound, not a uniform M\"obius source estimate.

\subsection{A compact realization of the complete arithmetic spectrum}
\label{sec:compact-complete-arithmetic-spectrum}

We prove an exact spectral identification on the whole closed arithmetic
range. The operator retains the original gamma term, every prime power
and both pole terms. Its realization is compact and may be
nonselfadjoint. Realness of its spectrum remains a separate obligation.

We keep the Fourier convention
\(F_f(z)=\int_{\R}f(x)e^{-izx}\dd x\), with inverse factor
\(1/(2\pi)\). Let \(\Phi\) be the complete positive even theta density
with \(F_\Phi=\Xi\). For \(t\ge T_\Phi\), put
\[
 \phi=p_t=e^{-t/4}G_t*\Phi,\quad
 G_t(x)=(4\pi t)^{-1/2}e^{-x^2/(4t)},\quad
 c(x)=\cosh(x/2),\quad \kappa(x)=\tanh(x/2),
\]
\[
 H_t=L^2(c\phi\dd x),\qquad
 \mathcal H_t=L^2((c/\phi)\dd x),\qquad U_tg=\phi g.
                                                        \tag{CAS1}
\]
Thus \(U_t:H_t\to\mathcal H_t\) is unitary,
\(F_\phi(z)=e^{-t(z^2+1/4)}\Xi(z)\), and
\(\int c\phi=1/2\). Inner products are antilinear in the first
variable. For each integer \(p\ge1\), define
\[
 K_p(x)=\frac{\Gamma(2p)}{\Gamma(p)^2[4\cosh^2(x/2)]^p},\qquad
 J_p(z)=\frac{\Gamma(p+iz)\Gamma(p-iz)}{\Gamma(p)^2},\qquad
 r_p=J_p(i/2)>0,
                                                        \tag{CAS2}
\]
\[
 \Pi_{\rm ar}=\sum_{n\ge2}\frac{\Lambda(n)}{\sqrt n}
             (\delta_{\log n}+\delta_{-\log n}),\quad
 a_\Gamma(u)=\Re\psi(1/4+iu/2)-\log\pi,
\]
\[
 \gamma_p=F^{-1}(a_\Gamma J_p),\quad
 S_p=K_p*\Pi_{\rm ar}-\gamma_p,
 \qquad
 (P_pg)(x)=\frac1{r_pc(x)}\int_\R S_p(x-y)\phi(y)g(y)\dd y.
                                                        \tag{CAS3}
\]
Every term \(\Lambda(n)\) is retained, including \(n=\ell^k\) for
every prime \(\ell\) and \(k\ge1\). The completed explicit formula is
\[
 2r_pc(x)-S_p(x)
      =\sum_{z:\,\Xi(z)=0}m_zJ_p(z)e^{izx}.              \tag{CAS4}
\]
The sum runs over all distinct nodes, at both heights and with every
conjugate and mirror present; \(m_z\) is the actual zero order.
The exponential decay of \(J_p\) on \(|\Im z|<1/2\) and the
unconditional local zero count pay normal convergence of this series
and all its fixed derivatives. In particular,
\( |S_p^{(j)}(x)|\le C_{p,j}e^{|x|/2}\).
The term \(2r_pc\) in (CAS4) is the original two-pole kernel.

These bounds give bounded selfadjoint operators \(P_p\) on \(H_t\).
For example, their kernels against \(c\phi\dd x\) are
\(S_p(x-y)/(r_pc(x)c(y))\); absolute Fubini follows from
\(e^{|x-y|/2}/(c(x)c(y))\le4\).
Equivalently, (CAS4) is an absolutely summable rank-one expansion,
so \(P_p\) is trace class. Its original row identities are
\[
 P_p1=1,\qquad P_p(\phi'/\phi)=\kappa/2.                 \tag{CAS5}
\]
The common closed range is
\[
 \mathcal C_t=(\ker P_1)^{\perp}
      =\overline{\operatorname{ran}P_j}\quad(j\ge1).
                                                        \tag{CAS6}
\]
It is defined by the full arithmetic operators, without assuming a
real zeroset. The restrictions \(P_j|_{\mathcal C_t}\) are injective.

\begin{theorem}[Complete compact arithmetic realization]
\label{thm:compact-complete-arithmetic-spectrum}
For every fixed integer \(p\ge1\) and \(t\ge T_\Phi\), there is a bounded
compact operator \(R_{p,t}:\mathcal C_t\to\mathcal C_t\) satisfying
\[
 R_{p,t}P_{p+1}=P_p.                                    \tag{CAS7}
\]
It is injective and has dense range. Its entire spectrum is
\[
 \boxed{\ \sigma(R_{p,t})=\{0,1\}\cup
       \left\{\frac{p^2-1/4}{p^2+z^2}:\Xi(z)=0\right\}.\ } \tag{CAS8}
\]
Zero belongs to the spectrum and is not an eigenvalue. The eigenspace
at \(1\) is \(\operatorname{span}\{1,\kappa\}\). For any other nonzero
spectral value associated with a node \(z\), the eigenspace is exactly
\[
 \operatorname{span}\{e^{izx}/c(x),e^{-izx}/c(x)\}.        \tag{CAS9}
\]
Every nonzero eigenvalue is semisimple, with algebraic and geometric
dimension two. This conclusion does not assume simple Xi zeros.
Their original orders remain the weights \(m_z\) in (CAS4), and do
not produce operator Jordan chains.

Compactness is asserted for each fixed \(p,t\). No uniform
compactness, trace-class estimate for \(R_{p,t}\), or unitary
identification of the different spaces \(\mathcal C_t\) is asserted.
\end{theorem}

\begin{proof}
We first construct the arithmetic realization with its actual
domains. Set
\[
 \lambda_p=p^2-1/4,\quad q_p=\lambda_p/p^2,\quad
 E_p(x)=\tfrac p2e^{-p|x|},\quad
 \mathcal N_p=\operatorname{span}\{e^{px}/c,e^{-px}/c\},\quad
 \mathcal M_p=\mathcal N_p^\perp,
\]
and let \(\Pi_p\) be the orthogonal projection onto \(\mathcal M_p\).
The raw expression
\[
 \mathcal R_pg=q_pc^{-1}E_p*(cg)                         \tag{CAS10}
\]
is first used only where the displayed convolution is ordinary and
absolute. It is not a bounded operator on all of \(H_t\).
The exact highmoment condition is
\[
 U_t\mathcal M_p=\ker F(ip)\cap\ker F(-ip).              \tag{CAS11}
\]

The ordinary kernel identities retain all terms in (CAS3) and give
\[
 \mathcal R_pP_{p+1}=P_p,
 \qquad
 (D_0+\lambda_p)P_p=\lambda_pP_{p+1},\qquad
 D_0=c^{-1}(-\partial_x^2+1/4)c.
                                                        \tag{CAS12}
\]
Indeed, \(J_{p+1}(z)=(1+z^2/p^2)J_p(z)\),
\(r_{p+1}=q_pr_p\), and
\(S_p''=p^2(S_p-S_{p+1})\).
All the differentiated integrals are absolute by the bound following
(CAS4). These identities also imply
\(\ker P_p=\ker P_{p+1}\), giving (CAS6).

We pay the missing two-dimensional angle before extending (CAS10).
For any \(0<s<2t\), the vectors
\[
 k_s=(p_s-p_t)/p_t,\qquad l_s=(p_s'-p_t')/p_t             \tag{CAS13}
\]
belong to \(H_t\) and to every \(\ker P_j\).
Their Gaussian integrability follows, for example, from
\[
\begin{aligned}
 \int\frac{c p_s^2}{p_t}\dd x
 &\le\frac12\frac{t}{\sqrt{s(2t-s)}}
       \exp\!\left(\frac{(s-t)^2}{2(2t-s)}\right),
 \\
 \int\frac{c (p_s')^2}{p_t}\dd x
 &\le\frac12\frac{t}{\sqrt{s(2t-s)}}
       \exp\!\left(\frac{(s-t)^2}{2(2t-s)}\right)
       \left(\frac{t}{2s(2t-s)}+
                         \frac{t^2}{4(2t-s)^2}\right).
\end{aligned}
                                                        \tag{CAS14}
\]
obtained by inserting the Gaussian score before Cauchy--Schwarz.
Their membership in the common kernel follows either from the
completed formula (CAS4), since every nodal Fourier value vanishes,
or from the full ground identity and its derivative. Both original
pole moments are independent of \(s\) and cancel in (CAS13).
For \(n_e=\cosh(px)/c\), \(n_o=\sinh(px)/c\), write
\[
 b_s=\Xi(ip)\bigl(e^{s\lambda_p}-e^{t\lambda_p}\bigr).
\]
Then the ordinary tilted Gaussian integrals give
\[
 \langle k_s,n_e\rangle=b_s,\qquad
 \langle l_s,n_o\rangle=-p b_s.                          \tag{CAS15}
\]
Choose \(s\ne t\). Both numbers are nonzero because \(\Xi(ip)>0\).
Parity makes the two sectors orthogonal. Since \(k_s,l_s\in
\mathcal C_t^\perp\), their nonzero pairings with the two homogeneous
modes imply an actual \(\delta_{p,t}>0\) such that
\[
 \|y\|_{H_t}\le\delta_{p,t}^{-1}\|\Pi_py\|_{H_t}
                     \quad(y\in\mathcal C_t).            \tag{CAS16}
\]
For example, one may take the minimum of the two normalized
pairings in (CAS15) as the lower angle. This argument pays a fixed
parameter bound; growing-parameter refinements are separate.

We next prove compactness in the actual physical space.
Put \(W=c/p_t\), \(V=\log W\), and
\[
 \mathcal D=\{f\in\mathcal H_t:
                         -i\partial_xf\in\mathcal H_t\}.
\]
The central heat division theorem gives
\(V''\ge\alpha=1/(4t)>0\) on the whole real line.
The maximal derivative operator \(D=-i\partial_x\) is closed on
\(\mathcal D\). Smooth compact functions form a graph core:
cutoffs \(\chi_R\) with \(\|\chi_R'\|_\infty\le C/R\) cost the
two graph tails and \((C/R)\|f\|\), and subsequent mollification is
ordinary \(H^1\) approximation on a compact where \(W,W^{-1}\) are
bounded.

For \(a=A+iB\), write \(\sqrt W f=e^{iAx}v\). On this core,
integration by parts gives the positive gauge identity
\[
 \|(D-a)f\|_{\mathcal H_t}^2
    =\|v'\|_2^2+\|(V'/2-B)v\|_2^2
                                  +\frac12\int V''|v|^2\dd x.
                                                        \tag{CAS17}
\]
Core approximation and weak lower semicontinuity extend the
required inequality to the maximal graph domain.
There is a unique \(x_B\) with \(V'(x_B)/2=B\), and
\[
 |V'(x)/2-B|\ge\tfrac\alpha2|x-x_B|.
\]
Thus a graph-bounded sequence has bounded ordinary \(H^1\) norm
after gauging and a bounded second spatial moment.
Outside \([x_B-L,x_B+L]\) its squared norm is uniformly \(O(L^{-2})\).
Inside that interval, piecewise averages on a mesh of size \(h\)
cost at most \(h\|v'\|_2\), and lie in a fixed finite-dimensional
space. First choosing \(L\), then \(h\), and then a finite net proves
that \(\mathcal D\hookrightarrow\mathcal H_t\) is compact.

The proved division
\(B_a:\ker F(a)\to\mathcal D\) satisfies
\(\|B_a\|\le2\sqrt t\) and \(DB_af=aB_af+f\).
Its graph image of a bounded set is bounded, so \(B_a\) is compact.
For \(f=U_tg\), \(g\in\mathcal M_p\), division first by \(z+ip\)
leaves its \(+ip\) zero. Division next by \(z-ip\) is consequently
legal. Entire Fourier division and ordinary \(L^1\) Fourier
uniqueness give
\[
 E_p*f=p^2B_{ip}B_{-ip}f.                               \tag{CAS18}
\]
Consequently \(A_p=q_pU_t^{-1}(E_p*\cdot)U_t:
\mathcal M_p\to H_t\) is compact. Its Hilbert adjoint is exactly
the bounded extension of \(\Pi_p\mathcal R_p\).
First pair \(h\in\mathcal M_p\) with \(y\in C_c^\infty\).
The two highmoments belong to \(h\); the compact support of \(y\)
makes the ordinary convolution Fubini integral absolute. It gives
\(\langle A_ph,y\rangle=\langle h,\Pi_p\mathcal R_py\rangle\).
Density of these \(y\) extends the adjoint to all of \(H_t\).
For \(y=P_{p+1}g\), its bounded ordinary output and exponential
convolution tails justify the same pairing directly.
Adjoint compactness follows directly by approximating
\(A_p\) in norm with its finite-rank image projections and taking
adjoints.

By (CAS16), \(\Pi_p\mathcal C_t\) is closed, and its inverse
\(J_C:\Pi_p\mathcal C_t\to\mathcal C_t\) is bounded. On the dense
actual range \(\operatorname{ran}P_{p+1}\), (CAS12) gives
\(A_p^*y=\Pi_p\mathcal R_py\in\Pi_p\mathcal C_t\).
Closedness and continuity give the same membership on all of
\(\mathcal C_t\). Define
\[
 R_{p,t}=J_CA_p^*|_{\mathcal C_t}.                       \tag{CAS19}
\]
This is compact and satisfies (CAS7). Local \(L^2\) convergence in
\(H_t\) passes the ordinary Green identity on the dense actual
range to all inputs distributionally:
\[
 (D_0+\lambda_p)R_{p,t}y=\lambda_py
                                  \quad(y\in\mathcal C_t). \tag{CAS20}
\]
It proves injectivity. The range is dense because (CAS7) makes it
contain \(\operatorname{ran}P_p\), which is dense in \(\mathcal C_t\).
No selfadjoint domain for \(D_0\), or raw convolution domain for
every completion input, has been assumed.

Suppose \(R_{p,t}y=\lambda y\ne0\), \(\lambda\ne0\).
Equation (CAS20) gives
\[
 D_0y=d_\lambda y,
 \quad d_\lambda=\lambda_p(\lambda^{-1}-1),
 \quad a^2=1/4-d_\lambda.
\]
Hence \((cy)''=a^2cy\) distributionally. Factoring the first-order
equations proves their ordinary smooth solutions. For \(a\ne0\),
\[
 y=(b_+e^{ax}+b_-e^{-ax})/c,\qquad
                       \lambda=\frac{\lambda_p}{p^2-a^2}. \tag{CAS21}
\]
All such candidate modes are \(H_t\)-integrable; the arithmetic
kernel must exclude the incorrect ones. Orthogonality to the two
actual vectors in (CAS13), for every \(0<s<2t\), gives
\[
 (b_++b_-)\Xi(ia)
      [e^{s(a^2-1/4)}-e^{t(a^2-1/4)}]=0,
\]
\[
 -a(b_+-b_-)\Xi(ia)
      [e^{s(a^2-1/4)}-e^{t(a^2-1/4)}]=0.                \tag{CAS22}
\]
If \(a^2\ne1/4\), the bracket is not identically zero on this
real \(s\)-interval. Thus a nonzero vector requires \(\Xi(ia)=0\)
or \(a^2=1/4\), the latter giving \(\lambda=1\).
For \(a=0\), \(y=(b_0+b_1x)/c\); the same two tests give
\(b_0\Xi(0)(e^{-s/4}-e^{-t/4})=0\) and
\(-b_1\Xi(0)(e^{-s/4}-e^{-t/4})=0\). Choosing \(s\ne t\) excludes
that repeated-root mode.

Every original node is actually present, not merely permitted by
the ODE. If \(z\) has order \(m\), take the physical isolator
\[
 F_{f_z}(\zeta)=
 \frac{m!}{(z^2+1/4)\Xi^{(m)}(z)}
       (\zeta^2+1/4)\frac{\Xi(\zeta)}{(\zeta-z)^m}.      \tag{CAS23}
\]
It is obtained by the genuine causal \(m\)-fold division of
\(\Phi\), followed by the displayed polynomial in \(D\).
For one division the physical solution is
\(i\int_{-\infty}^x e^{iz(x-u)}\Phi(u)\dd u\); higher divisions
insert \(i^{m-1}(x-u)^{m-1}/(m-1)!\).
The full vanished moments through order \(m-1\) replace the right
tail by the negative integral from \(x\) to infinity. The original
double-exponential theta tails pay both tails, all finite
derivatives and the complex phase. A fixed positive theta-mass
slice gives a Gaussian lower bound on \(p_t\), so
\(f_z/p_t\in H_t\).
Its selected Fourier value is one, all other node values are zero,
and both original pole values are zero. Formula (CAS4) gives
\[
 P_p(f_z/p_t)=-\frac{mJ_p(z)}{r_p}\,\frac{e^{izx}}{c(x)}.\tag{CAS24}
\]
Both gamma arguments in \(J_p(z)\) have positive real part.
They are finite and nonzero by the gamma-function analytic
properties~\cite{DLMFGamma}. Hence both \(e^{\pm izx}/c\) are
genuine actual-range vectors. The ordinary absolute Laplace
integral yields
\[
 R_{p,t}(e^{izx}/c)=\frac{\lambda_p}{p^2+z^2}\,e^{izx}/c.\tag{CAS25}
\]
The pools give \(1,\kappa\) at eigenvalue one.
Positivity of \(\Phi\) gives \(\Xi(iu)>0\) for real \(u\), so a
node is never zero or purely imaginary and never coincides with
\(\pm i/2\). Two nodes give the same value in (CAS25) precisely
when they are opposites. The ODE bound and actual presence therefore
prove the exact two-dimensional eigenspaces.

We exclude every operator Jordan jet without assuming simple nodes.
If \((R_{p,t}-\lambda)y=h\ne0\),
\(R_{p,t}h=\lambda h\), then (CAS20) gives
\[
 (D_0-d_\lambda)y=-\lambda_p\lambda^{-2}h.
\]
For \(h=(b_+e^{ax}+b_-e^{-ax})/c\), \(a\ne0\), every solution has
\[
 cy=(A_+x+B_+)e^{ax}+(A_-x+B_-)e^{-ax},\quad
 A_+=\frac{\lambda_pb_+}{2a\lambda^2},\quad
 A_-=-\frac{\lambda_pb_-}{2a\lambda^2}.                  \tag{CAS26}
\]
At least one \(A_\pm\) is nonzero. For a node eigenvalue put
\(w=i\overline a\), an original node of its actual order \(m_w\),
and construct the physical test
\[
 F_f(\zeta)=(\zeta^2+1/4)
       \frac{\Xi(\zeta)}{(\zeta-w)^{m_w-1}}
                         \frac{\zeta+w}{2w}.            \tag{CAS27}
\]
It is a genuine finite causal division and polynomial derivative
of \(\Phi\), with double-exponential tails. Every original node
value and both pool values vanish, while
\(F_f'(w)\ne0\) and \(F_f'(-w)=0\). Thus
\(f/p_t\in\ker P_j\) for every \(j\), and first-slot conjugation
gives exactly
\[
 \langle f/p_t,y\rangle_{H_t}
   =A_+\overline{iF_f'(w)}+A_-\overline{iF_f'(-w)}
       +B_+\overline{F_f(w)}+B_-\overline{F_f(-w)}.       \tag{CAS28}
\]
Orthogonality of \(y\in\mathcal C_t\) forces \(A_+=0\).
Interchanging \(w\) and \(-w\) forces \(A_-=0\), a contradiction.
For the groundvalue replace (CAS27) by
\(F_f(\zeta)=(\zeta^2+1/4)\Xi(\zeta)(\zeta+w)/(2w)\),
\(w=\pm i/2\). The selected derivative is
\(2w\Xi(w)=w\ne0\), while the opposite derivative vanishes.
The same argument excludes both groundjets. No chain of length
two exists, hence no longer chain exists.

For clarity, compactness closes the entire spectral set by an
elementary argument. If \(K\) is compact, \(\lambda\ne0\), and
\(A=K-\lambda I\) is injective, then \(A\) is bounded below.
Otherwise compact images of unitvectors with \(Ax_n\to0\) give a
nonzero kernelvector. Its ranges \(\operatorname{ran}A^n\) are
closed. If \(A\) were not onto, these ranges would be strictly
decreasing. Choose unit \(x_n\in\operatorname{ran}A^{n-1}\)
orthogonal to \(\operatorname{ran}A^n\). For \(n>m\),
\(Kx_n\in\operatorname{ran}A^m\) whereas
\(Kx_m=\lambda x_m+Ax_m\), so
\(\|Kx_m-Kx_n\|\ge|\lambda|\), contradicting compactness.
Thus injectivity implies a bounded inverse. Every nonzero spectral
point is consequently an eigenvalue, and there are no extra
nonzero spectral points beyond (CAS25) and the pools.

The absence of jets also gives algebraic dimension two directly.
For \(A=R_{p,t}-\lambda I\) and \(E=\ker A\), compactness bounds
\(A\) below on \(E^\perp\), so \(\operatorname{ran}A\) is closed.
The jet exclusion gives \(E\cap\operatorname{ran}A=0\).
The compact operator \(K'=R_{p,t}-\lambda Q_E\) has
\(K'-\lambda I=A-\lambda Q_E\) injective: for \(x=e+v\),
\(e\in E\), \(v\perp E\), its kernel equation says
\(Av=\lambda e\in E\cap\operatorname{ran}A\).
The proved compact argument makes it onto. Therefore
\[
 \mathcal C_t=E\dotplus\operatorname{ran}A.               \tag{CAS29}
\]
On the second closed invariant summand, \(A\) is boundedly
invertible; on \(E\), \(R_{p,t}=\lambda I\). This pays the
semisimple spectral summand without a basis assumption.

There are infinitely many distinct critical nodes, for example
by Conrey's unconditional simple-critical count~\cite{Conrey}.
Their actual-range exponential modes are linearly independent
by the finite Vandermonde argument. Thus \(\mathcal C_t\) is
infinite-dimensional. A compact operator on that space cannot
be boundedly invertible, so \(0\in\sigma(R_{p,t})\).
Injectivity and dense range established above show that zero
is neither an eigenvalue nor a hidden cokernel mode. This proves
all of (CAS8)--(CAS9).
\end{proof}

For an actual node \(z\), \(\Re z\ne0\), since \(\Xi(iu)>0\).
The value \((p^2-1/4)/(p^2+z^2)\) is therefore real precisely when
\(\Im z=0\). Realness of the complete nonzero spectrum would
prove RH for all heights and all original multiplicities.
The compact realization and the exact spectral identification
prove no such realness assertion by themselves.

\subsection{Fractional complete arithmetic and full Dirichlet--Hankel positivity}
\label{sec:fractional-complete-arithmetic}

We next identify a continuous smoothing family, prove positivity of its
entire Dirichlet block, and compute exactly what this positivity retains.
Throughout this subsection, \(p=1\), \(r=\pi/2\), \(q=3/4\),
\(J(z)=\Gamma(1+iz)\Gamma(1-iz)\), and \(S=S_1\) in (CAS3).
The already proved source and ground identities give constants
\(\epsilon>0,C<\infty\), independent of normalized heat, such that
\[
 2\epsilon r c\le S\le Cr c,\qquad S*\phi=rc,
 \qquad \int c\phi=1/2,
 \quad \phi=p_t,\quad t\ge T_\Phi .                       \tag{FH1}
\]
In particular, the full zero expansion is
\(S(x)=2rc(x)-\sum_{a\in Z}m_aJ(a)e^{iax}\), with every distinct
actual Xi node, both height signs and all original orders retained.

\begin{theorem}[Fractional arithmetic, full block positivity, and persistent index]
\label{thm:fractional-complete-arithmetic}
For \(\Re\alpha>0\), set
\[
 E_\alpha=\frac1{\Gamma(\alpha)}
    \int_0^\infty e^{-s}s^{\alpha-1}G_s\,\dd s,
 \qquad S_\alpha=q^\alpha E_\alpha*S,
 \quad \rho(a)=\frac q{1+a^2}.                            \tag{FH2}
\]
The integral defines an \(L^1\) kernel, and its full arithmetic
identification is
\[
 S_\alpha=(q^\alpha E_\alpha*K_1)*\Pi
             -F^{-1}(a_\Gamma J_\alpha),\qquad
 J_\alpha(z)=q^\alpha J(z)(1+z^2)^{-\alpha}.              \tag{FH3}
\]
Every prime power in \(\Pi\), the full gamma multiplier and both poles
remain present. On \(H_t=L^2(c\phi\,\dd x)\), define
\[
 (P_\alpha g)(x)=\frac1{rc(x)}
             \int S_\alpha(x-y)\phi(y)g(y)\,\dd y.       \tag{FH4}
\]
For real \(\alpha>0\), this is a selfadjoint, positive-preserving
Markov contraction. The family is holomorphic in actual operator norm
on \(\Re\alpha>0\), with the exact bound
\[
 \|P_{\sigma+i\tau}\|\le
             \frac{\Gamma(\sigma)}{|\Gamma(\sigma+i\tau)|}.
                                                               \tag{FH5}
\]
For every \(b>0\), integer \(d\ge0\) and all \(d+1\) form vectors,
the entire block satisfies
\[
 \begin{split}
 \mathcal D_{b,d}
   &:=[\Gamma(b+j+k)q^{-(b+j+k)}(I-P_{b+j+k})]_{j,k=0}^d\ge0,\\
 \mathcal D_{b,d}
   &\ge\epsilon[\Gamma(b+j+k)q^{-(b+j+k)}
     (I-2|1\rangle\langle1|+2|\kappa\rangle\langle\kappa|)]_{j,k=0}^d .
 \end{split}                                                   \tag{FH6}
\]
There are constants \(C_Z<\infty\) and \(\beta<1\), independent of
every normalized positive heat ground, such that for real
\(\alpha\ge0\), with \(P_0\) the original logistic operator,
\[
 \|P_\alpha-P_{\rm pool}\|_1\le C_Z\beta^\alpha,
 \qquad \|QP_\alpha Q\|_1\le C_Z\beta^\alpha,
 \quad P_{\rm pool}=2|1\rangle\langle1|-2|\kappa\rangle\langle\kappa| .
                                                               \tag{FH7}
\]
Here \(Q\) is the orthogonal projection off both pool vectors, and
each trace norm is the actual norm in its own \(H_t\). Nevertheless,
at every finite real \(\alpha\ge0\),
\[
 n_+(P_\alpha)=1+2K_+,\qquad n_+(QP_\alpha Q)=2K_+,       \tag{FH8}
\]
including \(K_+=\infty\) and arbitrary original multiplicities.
\end{theorem}

\begin{proof}
Write \(\alpha=\sigma+i\tau\). The absolute integral in (FH2)
has majorant \(\Gamma(\sigma)E_\sigma/|\Gamma(\alpha)|\).
For \(|\Im z|<1\), Fourier--Laplace Fubini is paid by
\(\int_0^\infty e^{-[1-(\Im z)^2]s}s^{\sigma-1}\,\dd s<\infty\).
The Gaussian integral and the ordinary gamma Laplace integral give
\[
 F_{E_\alpha}(z)=(1+z^2)^{-\alpha},\qquad
 E_\alpha*c=q^{-\alpha}c .                              \tag{FH9}
\]
The powers are principal: \(\Re(1+z^2)>0\) in this strip.
The gamma integral is used with \(\Re\alpha>0\) and
\(\Re(1+z^2)>0\), precisely the hypotheses of
\href{https://dlmf.nist.gov/5.9.E1}{DLMF 5.9.1} with \(\mu=1\).
For real \(\alpha>0\), positivity, Fubini and (FH1) give
\(S_\alpha*\phi=rc\) and \(2\epsilon rc\le S_\alpha\le Cr c\).
For (FH3), fix \(1/2<\eta<1\). Since \(K_1(x)\le e^{-|x|}\)
and \(E_\sigma\) has a finite exponential \(\eta\)-moment,
\(|q^\alpha E_\alpha*K_1(x)|\le C_{\alpha,\eta}e^{-\eta|x|}\).
Every compact prime series is dominated by
\(C\sum_{n\ge2}\Lambda(n)n^{-1/2-\eta}<\infty\), by
Chebyshev--Abel. The gamma multiplier and all fixed derivatives
are integrable by the exponential decay of \(J\). Thus convolution
retains all terms and Fourier uniqueness proves (FH3).

The symmetric positive kernel in (FH4), against \(c\phi\,\dd x\),
has row integral one. Jensen and its finite invariant mass prove
selfadjointness and the contraction bound. Its Dirichlet form is
\[
 \langle g,(I-P_\alpha)g\rangle
   =\frac1{2r}\iint S_\alpha(x-y)\phi(x)\phi(y)
                  |g(x)-g(y)|^2\,\dd x\dd y .           \tag{FH10}
\]
The row identity pays this integral for every \(g\in H_t\).
For \(s>0\), let \(B_s\) have kernel
\((G_s*S)(x-y)/(rc(x)c(y))\). Its row integral is \(e^{s/4}\),
so \(\|B_s\|\le e^{s/4}\). The actual operator identity
\[
 P_\alpha=\frac{q^\alpha}{\Gamma(\alpha)}
       \int_0^\infty e^{-s}s^{\alpha-1}B_s\,\dd s          \tag{FH11}
\]
is absolutely convergent in operator norm. Its norm is at most
\(q^\sigma\Gamma(\sigma)q^{-\sigma}/|\Gamma(\alpha)|\),
which proves (FH5). Powers of \(|\log s|\) from differentiation
are locally uniformly integrable at both endpoints. This proves
operator-norm holomorphy without a false Markov semigroup identity.

For arbitrary \(g_0,\ldots,g_d\in H_t\), put
\(g(s,x)=\sum_{j=0}^d s^jg_j(x)\). The quadratic value of (FH6) is
exactly
\[
 \frac1{2r}\int_0^\infty e^{-s}s^{b-1}
   \iint(G_s*S)(x-y)\phi(x)\phi(y)
                 |g(s,x)-g(s,y)|^2\,\dd x\dd y\dd s .    \tag{FH12}
\]
The finite polynomial and \(\|B_s\|\le e^{s/4}\) pay every absolute
cross integral. The final integrand is nonnegative. Tonelli and
polarization identify all block entries and prove positivity.
Using \(G_s*c=e^{s/4}c\), the lower bound in (FH1), and
\(c(x-y)=c(x)c(y)-\sinh(x/2)\sinh(y/2)\) gives the lower in (FH6).
In particular, the odd pool has the displayed positive sign.
The row matrix is strictly positive:
\[
 \sum_{j,k=0}^d\overline{z_j}\Gamma(b+j+k)q^{-(b+j+k)}z_k
  =\int_0^\infty e^{-qs}s^{b-1}|\sum_{j=0}^d z_js^j|^2\,\dd s>0
                                                               \tag{FH13}
\]
for every nonzero finite vector. The Hilbert-valued version also
vanishes only when all coefficients vanish, by pairing the polynomial
with its coefficients. A single row constraint therefore cannot
remove this entire positive row cost.

The complete node expansion, convolved absolutely in (FH2), is
\[
 S_\alpha(x)=2rc(x)-\sum_{a\in Z}m_aJ(a)\rho(a)^\alpha e^{iax}.
                                                               \tag{FH14}
\]
Since \(\Re(1+a^2)>q\), every \(0<|\rho(a)|<1\).
The principal powers have the compact-parameter bound
\(|\rho^\alpha|\le e^{\pi|\tau|/2}|\rho|^\sigma\); the full series
and all fixed derivatives converge normally. As \(|\Re a|\to\infty\),
\(\rho(a)\to0\) uniformly in the zero strip. Inside a bounded
rectangle there are only finitely many zeros of the nonzero entire
Xi function. Their finite maximum, together with a sufficiently
small outer bound, supplies one \(\beta<1\); no rightmost zero is
assumed. For \(h_a=e^{i\overline a x}/c(x)\),
\(|h_a|\le2\) and \(\|h_a\|_{H_t}^2\le2\).
The rank-one series in (FH14) is absolutely nuclear, and hence
\[
 \|P_\alpha-P_{\rm pool}\|_1
   \le\frac2r\sum_{a\in Z}m_a|J(a)||\rho(a)|^\alpha
   \le\left(\frac2r\sum_am_a|J(a)|\right)\beta^\alpha .    \tag{FH15}
\]
The constants do not depend on heat. Projection off both pools proves
the second bound in (FH7).

For the exact index, evaluate \(f=\phi g\) at every node and both
pools in (FH14). Each critical-node coefficient
\(-m_aJ(a)\rho(a)^\alpha/r\) is strictly negative real.
Each nonreal conjugate pair gives a nonzero Hermitian off-diagonal
two-coordinate block, with one positive and one negative value.
The positive-height mirror pair and its negative-height mirror
therefore contribute two positive coordinates. The pool block has
exactly one positive and one negative coordinate. These facts give
the upper bound when \(K_+\) is finite.

For the matching lower, any finite selected collection is realized
by the actual theta isolators in (CAS23): divide \(\Xi(z)\) by
\((z-a)^m\) at an order-\(m\) node, multiply by \(z^2+1/4\),
and normalize the selected value. All other node values and both
pools vanish. Their finite causal divisions have double exponential
tails, since every required lower moment vanishes; their derivatives
remain admissible after dividing by \(\phi\), using its Gaussian
lower bound. The transforms \(2(1/2\mp iz)\Xi(z)\) similarly realize
the two separate pool coordinates and vanish at all nodes. Finite
linear combinations therefore realize every required block. If
\(K_+=\infty\), arbitrarily large finite positive subspaces result.
On \(QH_t\) the pool-null node isolators give the entire lower directly.
This proves (FH8), without continuity of an index at zero.
\end{proof}

The correct unweighted row factorization uses the bounded complete
range factor, not factorial division. Put
\(\widehat R=R_{1,t}\operatorname{Proj}_{\mathcal C_t}\).
The full node series gives \(\widehat RP_b=P_{b+1}=P_b\widehat R^*\)
for real \(b>0\). Consequently
\[
 [P_{b+j+k}]_{j,k=0}^d=V P_bV^*,\qquad
 Vg=(g,\widehat Rg,\ldots,\widehat R^dg).                 \tag{FH16}
\]
At \(b=0\), the same identity holds for the original integer family.
Its signed middle factor remains; (FH13) cannot be removed by
entrywise division by factorials.

\begin{corollary}[The paid complex-parameter bound does not force realness]
\label{cor:fractional-complex-cost}
The consequence of (FH5) at an individually isolated node, along all
fixed rays \(\alpha=\sigma(1+ik)\), is
\[
 \log|\rho|-k\arg\rho
   \le k\arctan k-\tfrac12\log(1+k^2),\qquad
 |\rho|\le\cos(\arg\rho).                               \tag{FH17}
\]
Equivalently, \((\Re a)^2-(\Im a)^2\ge-1/4\). This inequality is
already strict in the original zero strip and supplies no exclusion.
\end{corollary}

\begin{proof}
A pool-null actual isolator makes \(P_\alpha g_a\) a fixed nonzero
vector times \(\rho(a)^\alpha\). Apply (FH5), divide logarithms by
\(\sigma\), and let \(\sigma\to\infty\). Stirling's formula on this
fixed right-half-plane ray gives the first inequality in (FH17);
\href{https://dlmf.nist.gov/5.11.E1}{DLMF 5.11.1} is used only in
its fixed sector of validity. Minimizing the right side plus
\(k\arg\rho\) gives \(k=-\tan(\arg\rho)\) and the second inequality.
It says \(\Re(1/\rho)\ge1\), which is exactly the stated strip
inequality. No uniform-ray Stirling estimate is required.
\end{proof}

Thus an arithmetically identified positive family, whole-vector
Dirichlet--Hankel positivity and even heat-uniform trace-norm decay
can retain every off-line index at all finite parameters. The missing
arithmetic step must control the signed interaction, rather than its
limit or its positive row kernel alone.

\subsection{Exact smoothed inertia on every interval}
\label{sec:local-smoothed-index}

Keep the complete logistic source \(S\), \(J\) and \(r=\pi/2\) from
(FH1), and define
\[
 \mathcal W_J(x)=2c(x)-S(x)/r
      =\sum_{a\in\mathcal Z}\frac{m_aJ(a)}r e^{iax}.
 \tag{LI1}
\]
This definition retains the complete arithmetic source
\(S=K*\Pi-F^{-1}(a_\Gamma J)\); in particular, smoothing does not
truncate the prime-power measure. For a bounded nonempty open
interval \(I\), let \(\mathcal V\) be the span of any fixed finite
collection of functions \(x^je^{\alpha x}\), with real \(\alpha\)
and nonnegative integer \(j\). Put
\[
 \mathcal D_I=\{f\in L^2(I):\int_I f(x)v(x)\,\dd x=0
                                  \text{ for every }v\in\mathcal V\},
 \quad \mathcal D_I^\infty=\mathcal D_I\cap C_c^\infty(I),
 \tag{LI2}
\]
and
\[
 \mathfrak q_J(f,g)=
   \int_I\int_I\overline{f(x)}\mathcal W_J(x-y)g(y)\,\dd y\,\dd x .
 \tag{LI3}
\]

\begin{theorem}[The full smoothed index on every bounded interval]
\label{thm:local-smoothed-index}
For every such \(I\) and every fixed finite \(\mathcal V\),
\[
 n_-(\mathfrak q_J|_{\mathcal D_I})
   =n_-(\mathfrak q_J|_{\mathcal D_I^\infty})=2K_+ .
 \tag{LI4}
\]
The assertion includes infinitely many off-line pairs and arbitrary
original Xi multiplicities. It asserts no uniform lower bound
for negative eigenvalues or for the norm of interpolating vectors
as the interval or the moment collection changes.
\end{theorem}

\begin{proof}
Set \(w_a=m_a|J(a)|/r>0\) and \(\theta_a=J(a)/|J(a)|\).
On \(\ell^2(\mathcal Z)\), define
\[
 (T_I\xi)(x)=\sum_a\sqrt{w_a}\xi_a e^{i\overline a x},
 \qquad B e_a=\theta_a e_{\overline a}.
 \tag{LI5}
\]
Since \(\theta_{\overline a}=\overline{\theta_a}\), \(B\) is a
selfadjoint unitary. At a real zero, \(J(a)>0\), so its coordinate
is positive. Each nonreal conjugate pair has one positive and one
negative coordinate. The two height signs therefore give
\(n_-(B)=2K_+\), with the original \(m_a\) retained in \(w_a\).

The uniform bound
\[
 |J(\gamma+i\delta)|\le C(1+|\gamma|)e^{-\pi|\gamma|},
 \qquad |\delta|<1/2,
 \tag{LI6}
\]
and the unconditional local zero count make \(T_I\)
Hilbert--Schmidt. For every \(\xi\in\ell^2(\mathcal Z)\), its
synthesis is holomorphic on \(|\Im x|<\pi/2\). Indeed, on each
compact subset, Cauchy--Schwarz gives the summable squared
majorant \(w_a\exp(2|\Re a||\Im x|+|\Re x|)\); the same argument
pays every fixed derivative. The adjoint has coordinates
\((T_I^*f)_a=\sqrt{w_a}F_f(a)\), and direct absolute reconstruction
gives
\[
 \mathfrak q_J(f,g)=\langle T_I^*f,BT_I^*g\rangle .
 \tag{LI7}
\]
The coefficient at \(a\), after relabeling \(a\) and \(\overline a\),
is \(m_aJ(a)\overline{F_f(\overline a)}F_g(a)/r\), as required by
the kernel in (LI3).

We next prove joint independence of this synthesis and the
exponential moment functions. Choose a basis
\(v_\nu=x^{j_\nu}e^{\alpha_\nu x}\) of \(\mathcal V\). Suppose
\[
 T_I\xi+\sum_\nu b_\nu v_\nu=0 \quad\text{a.e. on }I .
 \tag{LI8}
\]
The analytic representative vanishes on \(I\), hence on its whole
strip of analyticity and on the entire real axis. There is no
purely imaginary Xi zero. To see this directly, \(\xi(s)>0\) for
real \(s>1\); for \(0<s<1\), the positive alternating eta series
and \(1-2^{1-s}<0\) give \(\zeta(s)<0\), which is compensated by
\(s(s-1)<0\). The functional equation covers \(s<0\), and
\(\xi(0)=\xi(1)=1/2\).

Because \(\sum_a\sqrt{w_a}|\xi_a|<\infty\), the right Laplace
transform of (LI8) is absolutely justified for
\(\Re z>\max(1/2,\max_\nu\alpha_\nu)\). It is
\[
 \sum_a\frac{\sqrt{w_a}\xi_a}{z-i\overline a}
       +\sum_\nu\frac{b_\nu j_\nu!}
                          {(z-\alpha_\nu)^{j_\nu+1}}=0 .
 \tag{LI9}
\]
The first series extends meromorphically to the whole plane,
normally on compact sets outside its discrete poles: for large
\(|\Re a|\), its denominator is at least \(|\Re a|/2\) on a fixed
compact. The nodal poles \(i\overline a\) are distinct and
nonreal, so none coincides with a real \(\alpha_\nu\). Meromorphic
continuation of the identity makes every nodal residue zero.
Then the principal parts at the finitely many \(\alpha_\nu\),
in descending jet order, make all \(b_\nu\) zero. Thus (LI8)
implies \(\xi=0\) and \(b_\nu=0\).

If \(\xi\) is orthogonal to \(T_I^*\mathcal D_I\), then
\(T_I\xi\in\mathcal V\). The preceding independence proves
\[
 \overline{T_I^*\mathcal D_I}=\ell^2(\mathcal Z).
 \tag{LI10}
\]
The smooth exactly constrained space is dense in \(\mathcal D_I\).
In fact the finite map
\(M_\nu(f)=\int_I f v_\nu\), from \(C_c^\infty(I)\) to
\(\mathbb C^{\dim\mathcal V}\), is onto; otherwise a nonzero
linear combination of the \(v_\nu\) would vanish as a distribution
on \(I\). Choose fixed smooth correctors \(\chi_\nu\) with
\(M_\mu(\chi_\nu)=\delta_{\mu\nu}\). For a smooth approximation
\(f_n\to f\in\mathcal D_I\), subtract
\(\sum_\nu M_\nu(f_n)\chi_\nu\). This correction tends to zero in
\(L^2(I)\) and makes all constraints exact. Therefore
\[
 \overline{T_I^*\mathcal D_I^\infty}=\ell^2(\mathcal Z).
 \tag{LI11}
\]
No uniform growing-moment norm for the correctors is assumed.

Approximate any fixed finite collection of orthonormal negative
coordinates of \(B\) by vectors \(T_I^*f_j\), with
\(f_j\in\mathcal D_I^\infty\). By continuity and \(\|B\|=1\),
their Gram matrix in (LI7) is strictly negative for sufficiently
accurate approximation. This supplies every finite lower index,
also when \(K_+=\infty\). Conversely, on a strictly negative
subspace the projection of \(T_I^*f\) onto the negative
coordinates of \(B\) is injective. Its dimension is at most
\(2K_+\) when that number is finite. This proves (LI4).
\end{proof}

If the constraints include \(e^{x/2}\) and \(e^{-x/2}\), the
rank-two form of \(2rc\) vanishes. With \(Uf=f/p_t\), the actual
heat-space form satisfies, for compact \(f,h\in\mathcal D_I\),
\[
 \langle Uf,P_0 Uh\rangle_{H_t}
  =\frac1r\int_I\int_I\overline{f(x)}S(x-y)h(y)\,\dd y\,\dd x
  =-\mathfrak q_J(f,h).
 \tag{LI12}
\]
On each fixed compact \(\overline I\), \(U\) is bounded and
invertible onto the supported subspace of \(H_t\): the positive
continuous weights \(p_t,c\) have finite maxima and positive
minima there. Thus the corresponding supported, pool-null
arithmetic compression has positive index exactly \(2K_+\).
All distant prime powers remain in \(K*\Pi\). Small test support
does not remove the infinite logistic tail.

\begin{proposition}[The explicit cost of moment depth on a fixed interval]
\label{prop:local-moment-price}
Let \(I=(l,r_I)\), \(A=e^l\), \(U=e^{r_I}\),
\(u_c=(A+U)/2\) and \(D=(U-A)/2\).
Fix \(D<r_0<u_c\), set \(\eta=r_0/D>1\), and fix a Xi node \(a\).
If \(f\in L^2(I)\) satisfies
\(\int_I f(x)e^{jx}\,\dd x=0\) for \(1\le j\le m\), then
\[
 |F_f(a)|\le C_{a,I,r_0}\eta^{-m}\|f\|_{L^2(I)},\qquad
 C_{a,I,r_0}
   =\frac{\sqrt{|I|}\,U}{1-D/r_0}
                \max_{|u-u_c|=r_0}|u^{-1-ia}|.
 \tag{LI13}
\]
In particular, \(|(T_I^*f)_a|\ge1\) requires
\(\|f\|\ge\eta^m/(\sqrt{w_a}C_{a,I,r_0})\).
\end{proposition}
\begin{proof}
The principal function \(H(u)=u^{-1-ia}\) is analytic on
\(|u-u_c|<u_c\). Cauchy's coefficient estimate bounds its Taylor
remainder after degree \(m-1\), on \([A,U]\), by
\[
 \frac{M(D/r_0)^m}{1-D/r_0},
 \qquad M=\max_{|u-u_c|=r_0}|H(u)|.
\]
Multiplying the Taylor polynomial by \(u\) gives only powers
\(u,\ldots,u^m\), whose integrals against \(f(x)\), \(u=e^x\),
are exactly zero. Since \(uH(u)=e^{-iax}\), the remainder and
\(\int_I|f|\le\sqrt{|I|}\|f\|\) prove (LI13).
For \(a=\gamma+i\delta\), an explicit finite bound is
\[
 M\le (u_c-r_0)^{-1+\delta}
          \exp\!\left(|\gamma|\arcsin(r_0/u_c)\right).
 \tag{LI14}
\]
Every finite depth therefore retains the full index but carries
a genuine selected-coordinate price. If all \(j\ge1\) moments
vanish on this fixed interval, \(f=0\): polynomials with no
constant term are dense on \([A,U]\), because \(A>0\) and one
can first approximate the continuous function divided by \(u\).
The change of variables has a bounded positive Jacobian. This
last infinite constraint is not the finite-depth statement.
\end{proof}

\subsection{Exact theta--M\"obius cancellation and the selected-node mode}
\label{sec:theta-mobius-mode}

We can now prove signed cancellation on a concrete family and
compute exactly what removing a selected zero changes.
The arithmetic function is the ordinary \(\mu\), throughout.
Define the first theta summand on the entire real line by
\[
 A_1(x)=(4\pi^2e^{9x/2}-6\pi e^{5x/2})e^{-\pi e^{2x}} .
                                                               \tag{TMO1}
\]
It is different from the even finite head used in the Gaussian
head certificates. Write
\(\mathfrak G(s)=\tfrac12s(s-1)\pi^{-s/2}\Gamma(s/2)\).

\begin{theorem}[Exact radical cancellation and a nonzero selected-node mode]
\label{thm:theta-mobius-mode}
Normally on compact real sets with every fixed derivative,
\[
 \Phi(x)=\sum_{n\ge1}n^{-1/2}A_1(x+\log n),\qquad
 \sum_{n\ge1}\mu(n)n^{-1/2}\Phi(x+\log n)=A_1(x).                \tag{TMO2}
\]
Put \(\lambda_j=5/2+2j\),
\(a_j=-2\pi(2j+3)(-\pi)^j/j!\), and
\[
 H_J(X)=(X^2-1/4)\prod_{0\le j<J}(X^2-\lambda_j^2),\quad
 f_J=H_J(\partial)\Phi,\quad
 C_J(s)=\int f_J(y)f_J(y+s)\dd y .
\]
For every fixed \(J\ge0\), as \(T\to+\infty\),
\[
 \sum_{n\ge1}\mu(n)n^{-1/2}C_J(T-\log n)
 =a_JH_J(\lambda_J)^2\xi(2J+3)e^{-\lambda_JT}
                              +O_J(e^{-(\lambda_J+2)T}).       \tag{TMO3}
\]
Its coefficient is nonzero; for \(J=0\) it is
\(-216\pi\xi(3)\).

In contrast, fix a hypothetical actual pair
\(a=\gamma+i\delta\), \(b=\bar a\), \(0<\delta<1/2\), of
exact multiplicity \(m\), and its PN interpolant
\[
 F_w(z)=R(z)\frac{\Xi(z)}{D(z)},\qquad
 D(z)=(z-a)^m(z-b)^m,\qquad F_w(a)=F_w(b)=1,                   \tag{TMO4}
\]
with polynomial \(R\), all other node values zero, and both
poles canceled. Use the fixed \(P_{\mathrm n}\) in (NL1) and set
\(f=P_{\mathrm n}(-i\partial)w\),
\(C(s)=\int\overline{f(y)}w(y+s)\dd y\).
Then there are polynomials \(B_a,B_b\), of degrees at most
\(m-1\), such that
\[
 \begin{split}
 \sum_{n\ge1}\mu(n)n^{-1/2}C(T-\log n)
   &=e^{ibT}B_b(T)+e^{iaT}B_a(T)+O(e^{-5T/2}),\\
 [T^{m-1}]B_b
   &=P_{\mathrm n}(b)
          \frac{i^m m\,\mathfrak G(1/2+ib)}{\Xi^{(m)}(b)}
                       \ne0 .
 \end{split}                                                   \tag{TMO5}
\]
Thus its modulus is at least \(c_wT^{m-1}e^{\delta T}\)
for every sufficiently large \(T\), without a favorable
phase subsequence or a rightmost-zero assumption.
Every sum in (TMO3) and (TMO5) is absolutely convergent for
each fixed \(T\); its signed value uses all integers.
\end{theorem}
\begin{proof}
For \(x\ge0\) the first series in (TMO2) is literally the full
theta series. Both sides are holomorphic on the connected strip
\(|\Im x|<\pi/4\); the series is normal on its compact sets.
The analytic even continuation therefore proves the identity
on the whole real line. Substituting this series in its
M\"obius sum gives an absolutely convergent double series:
after grouping \(n=dk\), its majorant is
\(\sum\tau(n)n^Me^{-c_xn^2}\), also after each fixed derivative.
The exact divisor identity \(\sum_{d\mid n}\mu(d)=1_{n=1}\)
proves the second equality. Evenness gives the reflected version
\(\sum\mu(n)n^{-1/2}\Phi(x-\log n)=A_1(-x)\).

For a fixed even polynomial \(H\), let \(f_H=H(\partial)\Phi\)
and \(A_H=H(\partial)A_1\).
The true theta profiles imply
\[
 \left|\int f_H(y)f_H(y+s)\dd y\right|
              \le Ce^{M|s|-ce^{|s|}} .
                                                               \tag{TMO6}
\]
Use \(e^{2|y|}+e^{2|y+s|}\ge2e^{|s|}\), and also
\(\ge e^{2|y|}\), assigning half to each bound;
the remaining polynomially weighted \(y\) integral is finite.
This majorant pays absolute Fubini in
\[
 \sum\mu(n)n^{-1/2}C_H(T-\log n)
                        =\int f_H(y)A_H(-y-T)\dd y.            \tag{TMO7}
\]
The expansion of (TMO1) at the left is
\(A_1(x)=\sum_{j\ge0}a_je^{\lambda_jx}\).
Applying \(H_J(\partial)\) kills its first \(J\) modes.
Its remainder has the global bound
\[
 |A_{H_J}(x)-a_JH_J(\lambda_J)e^{\lambda_Jx}|
                                    \le C_Je^{(\lambda_J+2)x}.
\]
On \(x\le0\) this follows from the convergent Taylor expansion;
on \(x\ge0\) its growing right side also pays the actual
double-exponential tail. Integration by parts and evenness give
\[
 \int f_J(y)e^{-\lambda_Jy}\dd y
        =H_J(\lambda_J)\Xi(i\lambda_J)
        =H_J(\lambda_J)\xi(2J+3).
\]
The latter number is positive from the defining factors at
the real point \(2J+3>1\).
The next exponential moment pays the remainder in (TMO7).
This proves (TMO3), including its nonzero coefficient.

For the quotient in (TMO4), define the actual causal sum
\(Y(x)=\sum\mu(n)n^{-1/2}w(x-\log n)\).
Its true theta profile makes it normal with all fixed derivatives
on compact sets and double-exponentially small as \(x\to-\infty\).
Termwise application of the finite polynomial operators and
(TMO2) therefore give the exact equation
\[
 D(-i\partial)Y=R(-i\partial)A_1(-x)=q(x).                    \tag{TMO8}
\]
The right side has a superexponential left tail and right
tail \(O(e^{-5x/2})\), with its fixed derivatives.
The causal inverse of \((-i\partial-r)^j\) has kernel
\[
 \frac{i^j}{(j-1)!}1_{x\ge0}x^{j-1}e^{irx}.
\]
This follows by the first-order integral and repeated convolution.
Partial fractions of \(1/D\) have top coefficient
\(c_{b,m}=(b-a)^{-m}\).
Their finite causal solution has the same left boundary as \(Y\);
the difference is a finite homogeneous exponential polynomial.
No nonzero such solution can be double-exponentially small at
the left, by the confluent Vandermonde independence.
Thus the causal solution is exactly \(Y\).

Replacing the upper endpoint of each causal moment by infinity
is legal since \(|\Im r|<1/2<5/2\).
Its omitted part is \(O(e^{-5x/2})\), including all polynomial
factors, by integration after \(s=x+u\).
Consequently
\[
 Y(x)=e^{ibx}P_b(x)+e^{iax}P_a(x)+O(e^{-5x/2}),\qquad
 [x^{m-1}]P_b
       =\frac{i^mF_q(b)}{(m-1)!(b-a)^m}.                      \tag{TMO9}
\]
The bound can be made global on the real line; on the left
the exponential polynomials are dominated by \(Ce^{-5x/2}\).

The substitution \(u=\pi e^{2x}\) in (TMO1) and the primary
Gamma recurrence \cite[5.2.1, 5.5.1]{DLMFGamma} give, for
\(s=1/2-iz\) with \(\Re s>0\),
\[
 F_{A_1}(z)=\tfrac12\pi^{-s/2}
                  [4\Gamma(s/2+2)-6\Gamma(s/2+1)]
                          =\mathfrak G(s).                   \tag{TMO10}
\]
Both selected Gamma arguments have strictly positive real part.
Reflection then gives \(F_q(b)=R(b)\mathfrak G(1/2+ib)\).
The selected value in (TMO4) says exactly
\(R(b)/(b-a)^m=m!/\Xi^{(m)}(b)\).
This fixes the top coefficient in (TMO9) without assuming simplicity.

The same profile (TMO6) pays absolute Fubini for
\(\sum\mu(n)n^{-1/2}C(T-\log n)
 =\int\bar f(y)Y(y+T)\dd y\).
The global remainder in (TMO9) costs \(O(e^{-5T/2})\).
The top coefficient of its \(b\) mode is multiplied by
\[
 \int\bar f(y)e^{iby}\dd y
      =\overline{F_f(\bar b)}
      =P_{\mathrm n}(b)\overline{F_w(a)}=P_{\mathrm n}(b).
\]
This proves (TMO5), including its first-slot conjugation.
The exact-order derivative, the nonvanishing Gamma factor,
and (NL8) make its top coefficient nonzero.
The \(b\) mode grows as \(e^{\delta T}\);
the \(a\) mode decays as \(e^{-\delta T}\).
Their leading-degree comparison proves the all-large-\(T\)
modulus bound.
\end{proof}

There is also a completely paid compact consequence.
For every fixed \(A>0\), choose \(J\) with \(4J+6>A+2\).
Cut off \(f_J\) at
\(R=\frac12\log(v+t+2)+B\), correct both pole moments exactly,
and write the resulting real function as \(z_R\).
Use \(\sqrt{a_t}e^{-v/2}\tau_{-v}z_R\) on the left and
\(\sqrt{a_t}e^{-v/2}\tau_vP_{\mathrm n}(-i\partial)z_R\)
on the right, with \(t=\log^3v\).
Their canonical norm product stays between fixed positive
constants, their actual correlation is
\(a_te^{-v}C_R(r-2v)\), and
\(\sup_s|C_R(s)-C_J(s)|\le Ce^{-d e^{2R}}\).
The right vector has all high moments (NL4);
both vectors have the two original pole moments.
The left vector is not claimed to have every high moment.

Let \(N=\lceil e^{2v+2R+2}\rceil\),
\(D=e^{2v}/h_v\), with
\(L_Av\le h_v\le v^{C_0}\) for fixed \(C_0,L_A\).
The complete unlogged finite fringe satisfies
\[
 \left|a_te^{-v}\sum_{D<n\le N}
          \mu(n)n^{-1/2}C_R(\log n-2v)\right|
                                \le C_Aa_te^{-Av}.            \tag{TMO11}
\]
Indeed (TMO3) at \(T=2v\) costs \(a_te^{-(4J+6)v}\).
For \(n\le D\), (TMO6) bounds the discarded sum by
\(Ca_th_v^{M-1/2}e^{-ch_v}\).
For \(n>N\), its integral majorant is
\[
 Ca_t\int_{N/e^{2v}}^\infty q^{M-1/2}e^{-cq}\dd q
                           \le Ca_te^{-cN/(2e^{2v})}.
\]
The finite compactification error is at most
\(Ca_te^{R+1}e^{-d e^{2R}}\).
Choose the fixed \(L_A,B\) sufficiently large; these three
errors prove (TMO11) uniformly in the stated polynomial cap range.
If \(D=e^{2v}/(L_Av)\), the inverse cofactor is \(O_A(v^2)\).
This optional inverse cap does not assert a Type-I-head estimate
at that \(D\); using the separately paid Type-I cap retains
its actual larger cofactor.

For the selected pair (TMO4), use \(p=\lfloor v/2\rfloor\),
\(R=\frac12\log(v+p+t+2)+B\), and the paid adaptive
preparation (NL18). Its correlation error is bounded by
\(Ce^{Cp-d e^{2R}}\); the whole finite fringe multiplies it
by at most \(2a_te^{R+1}\). Increasing the fixed \(B\)
pays this cost with any fixed margin \(A\ge2\).
The same exact lower and upper truncation argument then
gives the opposite finite-fringe forcing:
\[
 \left|a_te^{-v}\sum_{D<n\le N}
       \mu(n)n^{-1/2}C_R(\log n-2v)\right|
                    \ge c_wa_tv^{m-1}e^{-(1-2\delta)v}.        \tag{TMO12}
\]
Here \(C_R=\int\overline{P_{\mathrm n}(-i\partial)z_R(y)}
 z_R(y+s)\dd y\), and the compact errors from the prepared
construction are paid by increasing \(B\).
This correlation is Hermitian; replacing \(2v-\log n\)
by \(\log n-2v\) conjugates its signed sum and does not change
the modulus. The profile and cutoff bounds pay the same
three truncation costs used in (TMO11), now smaller than the
nonzero mode in (TMO5). No endpoint or principal term is dropped.

Adding a fixed theta-polynomial radical to the selected seed
replaces \(R\) by \(R+DH\) in (TMO8) and leaves the top
coefficient in (TMO5) unchanged.
The radical cancellation is therefore a proved arithmetic result,
but it cannot remove this selected-node mode by such additions.
Nor is radical density available: its Fourier value at that node
is zero, and evaluation is bounded in each fixed ground norm.
An upper bound excluding the mode in (TMO12) is still required.

\subsection{All heat times and the exact cost of escaping completion jets}
\label{sec:normalized-heat-flux}

Large heat removes a local Gaussian weight cost. It does not
automatically bound a completed arithmetic source.
Write
\[
 \widetilde p_t=p_t/a_t
  =\int\Phi(y)e^{-(x-y)^2/(4t)}\dd y,\quad
 \widetilde\phi_{v,t}
  =\frac{\widetilde p_t(\cdot-v)+\widetilde p_t(\cdot+v)}{2c(v)},
 \quad\widehat W_t=c/\widetilde p_t,\quad
 \widehat{\mathcal H}_t=L^2(\widehat W_t\dd x).                 \tag{NHA1}
\]
The norm identities are exact:
\(\|f\|_{\widetilde\phi}^2=a_t\|f\|_\phi^2\) and
\(\|f\|_{\widehat{\mathcal H}_t}^2=a_t\|f\|_{\mathcal H_t}^2\).

\begin{theorem}[Prepared sources uniformly over every heat time]
\label{thm:all-heat-prepared}
In Theorem~\ref{thm:prepared-nearlinear-source}, replace the radius
by \(R_v=\frac12\log(v+p+3)+B\) and the norm by the explicitly
normalized norm in (NHA1). Then the prepared construction is
independent of \(t\), and for all \(t\ge1\), simultaneously,
\[
 c\,c(v)|a|\le\|\tau_{\pm v}f_v[a]\|_{\widetilde\phi}
       \le C\,c(v)|a|,\quad
 K_{\rm short}\le C_Bv^4,\quad H=C_Av.                        \tag{NHA2}
\]
Its actual head has upper
\(e^{-v}\|f\|_{\widetilde\phi}\|h\|_{\widetilde\phi}\);
both complete finite source UV errors have upper
\(C_Ae^{-Av}\|f\|_{\widetilde\phi}\|h\|_{\widetilde\phi}\).
The exact \(D,N\), pole cancellations, and prime powers are retained.
If \(t/v^2\to\infty\), for every fixed main seed
\[
 \frac{\|\tau_{\pm v}f_\infty\|_{\widetilde\phi}^2}{c(v)^2}
   \longrightarrow
       \frac1{\Xi(0)}\int e^{\pm x/2}|f_\infty(x)|^2\dd x .
                                                               \tag{NHA3}
\]
The same limit holds for its prepared compact versions and their
fixed-basis Grams.
\end{theorem}
\begin{proof}
Put \(M_0=\Xi(0)\), \(M_2=\int x^2\Phi\).
Positivity, evenness, and \(1-e^{-u}\le u\) give everywhere
\[
 0\le M_0-\widetilde p_t(x)\le(M_0x^2+M_2)/(4t),\quad
 \widehat W_t(x)\ge c(x)/M_0,\quad
 \widetilde p_t(x)\ge m_1e^{-(|x|+1)^2/4}\ (t\ge1),            \tag{NHA4}
\]
where \(m_1=\int_{-1}^1\Phi>0\).
The last fixed mass-slice bound pays the prepared central errors
\(e^{Cp-d e^{2R_v}}\) from (NL18) uniformly in all \(t\).
Translation uses
\(\widetilde\phi(x\pm v)\ge\widetilde p_t(x)/(2c(v))\).
The fixed-slice Gram and \(\widetilde p_t\le M_0\) give the
lower norm for all coefficient vectors. No growing main-seed
price is introduced.

The first-\(p\) correlation derivatives are physical bounds,
so (NL19) now holds with the normalized norm and without a
heat restriction. Use the same full \(E_p\) and exact cap
\[
 D=e^{(2-1/p)(v-R_v-1)}
        [E_p(R_v+1)\log(2N)]^{-1/p},\quad
                 N=\lceil e^{2v+2R_v+2}\rceil .
\]
The head cancellation and \(O(v^4)\) calculation in (NL20)
are unchanged. The normalized UV prefactor is now at most
\(C\sqrt N\log(2N)e^{-v}=O_B(v^{3/2})\).
The fixed Fourier main and all-order Denjoy error therefore
pay both complete finite source tails at \(H=C_Av\),
uniformly over \(t\ge1\).

For \(t/v^2\to\infty\), both
\(\widetilde p_t(x)\) and \(\widetilde p_t(2v+x)\) tend to
\(M_0\) on each fixed slice. The actual normalized right
integrand is
\[
 \frac{2c(v+x)}{c(v)}
   \frac{|f_\infty(x)|^2}
        {\widetilde p_t(x)+\widetilde p_t(2v+x)} .
\]
The mass slice gives a common integrable theta/Gaussian
majorant, so dominated convergence proves (NHA3).
Both copies contribute; retaining only one would insert an
incorrect factor two. The uniformly paid prepared norm error
then proves the compact and Gram versions.

Finally no bound uniform in \(t\) for the whole operator
divided by \(a_t\) has been assumed. The detector cutoff
form error is paid directly by the complete physical bound
\(|Q_{\mathrm n}(r)|\le Ce^{|r|/2}\) in (NL10), integrated
against the two true theta profiles. Translation costs at
most \(Ce^v\); increasing the fixed \(B\) pays it by the
same \(e^{Cp-d e^{2R_v}}\). Thus the individual normalized
off-line forcing is \(c_be^{-(1-2\delta)v}\) at every
heat time, also in the new regime.
\end{proof}

\begin{theorem}[The actual finite completion must pay its escaping zero jets]
\label{thm:escaping-completion-cost}
Fix an actual Xi node \(b=\gamma-i\delta\), \(\gamma>0\),
\(0\le\delta<1/2\), of order \(m\ge1\).
Choose a fixed polynomial \(R\) with \(R(\pm i/2)=0\),
\(R(b)\ne0\), and set
\[
 \psi=P_{\mathrm n}(-i\partial)R(-i\partial)\Phi,\qquad
 S_N(x)=\sum_{n\le N}\mu(n)n^{-1/2}\psi(x-\log n).              \tag{NHA5}
\]
Take \(L=2v\) and the actual cap \(N=\lceil e^{L+2R_v+2}\rceil\)
from the preceding theorem. For every fixed \(q\ge0\),
\[
 \|(-i\partial)^qS_N\|_{\widehat{\mathcal H}_t}
           \ge c_qL^{m-1}e^{(1/4+\delta)L},
                   \qquad t\ge1,                             \tag{NHA6}
\]
for all sufficiently large \(v\), uniformly in \(t\).
Every \(S_N\) retains all original Xi zero jets, both poles,
all high moments, and every fixed derivative graph domain.
The bound already holds at an actual critical node \(\delta=0\).
\end{theorem}
\begin{proof}
The finite entire Fourier identity is
\[
 F_{S_N}(z)=P_{\mathrm n}(z)R(z)\Xi(z)
                  \sum_{n\le N}\mu(n)n^{-1/2-iz}.
\]
Its exact zero jets and moments follow immediately.
Every fixed derivative has a finite sum of translated
theta tails, which pays the mass-slice Gaussian majorant
in (NHA4), so its graph membership is genuine.

The normal theta--M\"obius identity (TMO2), on a fixed
complex circle, permits \(P_{\mathrm n}(-i\partial)R(-i\partial)\)
termwise: (NL2) pays the full derivative sum after Cauchy,
and the theta series has a compact-normal divisor majorant.
Thus the actual local limit is
\[
 g=P_{\mathrm n}(-i\partial)R(-i\partial)A_1(-x),\qquad
 G_b=F_g(b)=P_{\mathrm n}(b)R(b)\mathfrak G(1/2+ib)\ne0.
                                                               \tag{NHA7}
\]
It has a double-exponential left tail and \(O_q(e^{-5x/2})\)
right derivatives. It has no growing selected-node mode.

For \(x\le L+1\), the true omitted finite-sum tail is bounded by
\[
 |(S_N-g)^{(q)}(x)|
        \le C_qe^{x/2}e^{-d_qN^2e^{-2x}}.                    \tag{NHA8}
\]
Indeed substitute \(y=ne^{-x}\) in the decreasing tail
integral starting at \((N-1)e^{-x}\), and absorb its fixed
polynomial prefactor into \(d_q>0\).
Since \(Ne^{-L}\asymp_Bv\), integration against
\((1+|x|^j)e^{-\delta x}\) on this left half-line gives error
\(C_j(1+L^j)e^{(1/2-\delta)L-c_jv^2}\).
The integrability at the left uses precisely \(\delta<1/2\).

Choose a fixed smooth switch \(0\le\theta\le1\), zero on
\((-\infty,-1]\) and one on \([0,\infty)\), and split
\(S_N=I+T\), \(T=\theta(x-L)S_N\).
The omitted Gamma tail and (NHA8) give
\(F_I^{(j)}(b)=F_g^{(j)}(b)+O_j(\operatorname{poly}(L)
e^{-\epsilon_jL})\), \(j<m\).
All those full source jets are zero, hence \(F_T^{(j)}=-F_I^{(j)}\).
For \(Q(s)=T(L+s)\), supported on \(s\ge-1\), the exact
centered highest jet is
\[
 F_Q^{(m-1)}(b)
  =-e^{ibL}\sum_{j<m}\binom{m-1}{j}(iL)^{m-1-j}F_I^{(j)}(b),
 \quad
 |F_Q^{(m-1)}(b)|\sim|G_b|L^{m-1}e^{\delta L}.                 \tag{NHA9}
\]
The \(j=0\) coefficient is nonzero; the other terms have
lower powers and the errors have an exponential margin.
On the whole actual right half-line, (NHA4) and Cauchy--Schwarz give
\[
 |F_Q^{(m-1)}(b)|
 \le C_{m,\delta}e^{-L/4}\|T\|_{\widehat{\mathcal H}_t},\qquad
 C_{m,\delta}^2
  =2M_0\int_{-1}^\infty
               |s|^{2m-2}e^{-(2\delta+1/2)s}\dd s<\infty .
\]
This evaluation decays on the right because \(b=\gamma-i\delta\).
It is valid for every \(0\le\delta<1/2\), not only for
\(\delta<1/4\). Since \(\|S_N\|\ge\|T\|\), it proves
(NHA6) for \(q=0\). For fixed \(q\), the local transform
is \(b^qG_b\ne0\), and all zero jets remain; the identical
argument proves the derivative bound.
\end{proof}

The full source is therefore eligible for the division theorem
(HD5), but its actual derivative costs cannot be replaced by
the small local prefix. In fact, when \(t/L^2\to\infty\), that
prefix has bounded graph limits
\(\|D^qI\|_{\widehat{\mathcal H}_t}^2
 \to M_0^{-1}\int c|D^qg|^2\), by (NHA8), flatness on
\([0,L]\), and the uniform left theta majorant.
It also has \(F_I(b)\to G_b\ne0\), so it is outside the
zero ideal on which (HD5) is proved.
The missing jets are precisely in the expensive right tail.
A signed bilinear estimate retaining that flux, or a different
completion construction, remains necessary for RH.

\subsection{Every heat regime and the exact multiplicity cost of a finite shell}
\label{sec:bimodal-multiplicity-flux}
The next bound concerns the original two-copy norm (NL11),
not an inference from a lower bound in a stronger single-copy
space. We prove it directly and retain the actual finite
M\"obius endpoints. Its premise is one actual off-line zero;
it establishes the cost of that premise without assuming it occurs.

\begin{theorem}[Heat-independent finite-shell multiplicity cost]
\label{thm:bimodal-multiplicity-flux}
Fix \(q>1\), a nonzero \(h_0\in C_c^\infty(\R)\) with
\(F_{h_0}(\pm i/2)=0\), and \(Z=\Phi*h_0\).
Put \(M_0=\int\Phi\), \(D=-i\partial\), and, for \(X>1\),
\[
 \begin{split}
 E&=L^2(e^{-y/2}\dd y),\qquad
 U_X(y)=\sum_{X/q<n\le X}\mu(n)n^{-1/2}
                         Z(y+\log(n/X)),\\
 A_{q,r}(w;y)&=\int_0^{\log q}e^{-ws}(D^rZ)(y-s)\dd s,
                   \quad r=0,1,2,\ldots .
 \end{split}                                                   \tag{BMF1}
\]
Each \(A_{q,r}\) is entire with values in \(E\), and is nonzero
at every \(w\in\C\). For \(v,t>0\) take exactly
\[
 f_v(x)=\sqrt{a_t}e^{-v/2}U_{e^{2v}}(x+v),\qquad
 \phi_{v,t}(x)=\frac{p_t(x-v)+p_t(x+v)}{2\cosh(v/2)},\quad
 \|f\|_{\phi_{v,t}}^2=\int\frac{c(x)|f(x)|^2}{\phi_{v,t}(x)}\dd x.
                                                               \tag{BMF2}
\]
Every fixed derivative of this finite source belongs to the
actual weighted graph domain. It retains every Xi zero to
at least its original order and both original pole moments.

If \(s_0=1/2+\delta-i\gamma\) is an actual zeta zero of
order \(m\ge1\), with \(0<\delta<1/2\), put
\(w_0=\delta-i\gamma\). Then, for every fixed \(r\ge0\),
\[
 \limsup_{X\to\infty}
 \frac{\|D^rU_X\|_E}{X^\delta(\log X)^{m-1}}
 \ge\frac{m\,\|A_{q,r}(w_0)\|_E}{|\zeta^{(m)}(s_0)|}>0 .       \tag{BMF3}
\]
For every positive function \(t=t(v)\), without a continuity
or size restriction, this implies the original bimodal bound
\[
 \limsup_{v\to\infty}
 \frac{\|D^rf_v\|_{\phi_{v,t(v)}}}
      {e^{2\delta v}(2v)^{m-1}}
 \ge\frac{m\,\|A_{q,r}(w_0)\|_E}
             {2\sqrt{M_0}\,|\zeta^{(m)}(s_0)|}>0 .              \tag{BMF4}
\]
The constants may depend on the selected zero, its order,
the fixed seed and \(q,r\), but not on the heat choice.
No rightmost-zero or simplicity premise is used.
\end{theorem}
\begin{proof}
The full theta tails and compact convolution pay every fixed
derivative and exponential moment of \(Z\).
Translations by \(s\in[0,\log q]\) have \(E\)-norm
\(e^{-s/4}\|Z\|_E\). Thus every actual finite profile
is in \(E\), with \(\|U_X\|_E\le C_q\sqrt X\).
It is strongly measurable in \(X\): on compact \(X\)-intervals
there are finitely many continuous translation terms and endpoints.
The fixed \(s\)-interval pays all \(w\)-derivatives of \(A_{q,r}\)
locally uniformly, proving its entire \(E\)-valued dependence.

It is never the zero vector. Indeed its ordinary absolute
Fourier transform on the real axis is
\[
 F_{A_{q,r}(w)}(u)
  =u^r\Xi(u)F_{h_0}(u)
                \frac{1-e^{-(w+iu)\log q}}{w+iu}.             \tag{BMF5}
\]
The last factor has the nonzero removable value \(\log q\)
at zero and only discrete zeros. The first analytic factor
is not identically zero. Their product cannot vanish at every
real \(u\). If the \(E\)-vector were zero, its smooth
representative and then this transform would vanish everywhere,
which is impossible. This argument retains even those actual
nodes at which \(\Xi\) itself vanishes.

For \(\Re w>1/2\), absolute Bochner Fubini gives the exact identity
\[
 \int_1^\infty U_X X^{-w-1}\dd X
                       =\frac{A_{q,0}(w)}{\zeta(1/2+w)}.      \tag{BMF6}
\]
For each integer \(n\), its exact integration range is
\(n\le X<qn\); substitute \(X=ne^s\).
The remaining coefficient is \(\mu(n)n^{-1/2-w}\),
whose sum is absolutely convergent in this half-plane.
All integers and both finite endpoints are therefore retained.
The same identity with \(D^rU_X\) has numerator \(A_{q,r}\).
At \(w_0\) the right side has a genuine vector pole of order \(m\),
with leading coefficient \(m!A_{q,r}(w_0)/\zeta^{(m)}(s_0)\).

Let \(M\) be the limsup on the left of (BMF3).
If it is infinite there is nothing to prove.
Otherwise, for each \(\eta>0\), above a fixed \(X_0\),
\(\|D^rU_X\|_E\le(M+\eta)X^\delta(\log X)^{m-1}\).
This pays the actual Mellin integral and its derivatives
locally uniformly on \(\Re w>\delta\).
Meromorphic uniqueness identifies it there with (BMF6).
Any pole further right would already contradict this holomorphy;
its absence has not been assumed.
At \(w=w_0+\epsilon\), \(\epsilon>0\), triangle inequality gives
\[
 \left\|\frac{A_{q,r}(w_0+\epsilon)}
                   {\zeta(s_0+\epsilon)}\right\|_E
 \le C_{X_0}+(M+\eta)\int_1^\infty
                 X^{-1-\epsilon}(\log X)^{m-1}\dd X
 =C_{X_0}+(M+\eta)(m-1)!\epsilon^{-m}.                       \tag{BMF7}
\]
The finite-interval term is uniformly bounded as
\(\epsilon\downarrow0\). Multiply by \(\epsilon^m\), use the
nonzero leading Laurent coefficient, and then let
\(\eta\downarrow0\). The ratio \(m!/(m-1)!=m\) proves (BMF3).

For the bimodal transfer set \(Q_t=p_t/a_t\), so
\(0<Q_t\le M_0\) everywhere. Direct substitution \(y=x+v\) gives
\[
 \begin{split}
 \|D^rf_v\|_{\phi_{v,t}}^2
 &=\frac{1+e^{-v}}2
  \int\frac{e^{-y/2}+e^{-v}e^{y/2}}
            {Q_t(y)+Q_t(y-2v)}|D^rU_{e^{2v}}(y)|^2\dd y\\
 &\ge\frac1{4M_0}\|D^rU_{e^{2v}}\|_E^2 .
 \end{split}                                                   \tag{BMF8}
\]
Both heat copies stay in the denominator. Their sum is at most
\(2M_0\); only the positive second numerator is discarded.
The true finite double tails pay this entire integral and every
fixed derivative graph norm by the positive theta mass slice,
for every fixed \(v,t>0\).
The exact finite transform is
\[
 F_{f_v}(z)=\sqrt{a_t}e^{-v/2-izv}\Xi(z)F_{h_0}(z)
                   \sum_{e^{2v}/q<n\le e^{2v}}
                              \mu(n)n^{-1/2+iz}.
\]
It proves the claimed original zero jets and both pole moments.
Finally (BMF8) holds pointwise at every \(v,t(v)>0\).
As \(X=e^{2v}\) covers the entire positive tail,
(BMF3) yields (BMF4) for every such heat function.
\end{proof}

The sources in (BMF2) remain in the full zero ideal.
Their large norm alone is not a negative Weil-form value.
For \(t\ge\log^3v\) the actual input
\(\sqrt{a_t}e^{-v/2}\tau_vh_0\) has bimodal norm bounded
above and below by fixed positive constants: the nearby
\(Q_t\)-copy tends to \(M_0\) on its fixed support and the other
copy lies between zero and \(M_0\).
Thus the theorem gives a genuine source cost in that whole
regime. A subexponential upper for these same finite sources
in the same full norm would exclude every off-line zero.
That upper remains an arithmetic proof obligation.

\subsection{Adaptive completion with a paid growing support radius}
\label{sec:adaptive-completion-radius}

The fixed-profile Mellin theorem also detects genuinely adaptive
cofactors without factoring a vector-valued Mellin transform. Put
\(E=L^2(e^{-y/2}\,\dd y)\), \(J_Ef=e^{-y/4}f\), and, for fixed \(q>1\),
\[
 L_Xh(y)=\sum_{X/q<n\le X}\mu(n)n^{-1/2}h(y+\log(n/X)),
 \quad\Delta_X(s)=\sum_{X/q<n\le X}\mu(n)n^{-s},
 \quad v=\tfrac12\log X.                                  \tag{AC1}
\]
All shifts lie in \([-\log q,0]\). This window differs from a full cap
whose shifts extend to \(\log X\).

\begin{theorem}[All-radius inverse completion and adaptive node visibility]
\label{thm:adaptive-completion-radius}
Fix an actual Xi zero \(a=\gamma+i\delta\), \(|\delta|<1/2\), of order
\(m_a\ge1\). Set
\(b_a=\max(2|\delta|-1/2,0)\), \(\kappa_a=5/8+b_a^2/2<3/4\),
and \(p_a=m_a/2+1/4\). For every fixed \(A>0\) there are constants
\(B,C\), depending on the selected node, its order, \(A\) and \(\Phi\),
such that for all \(v\ge1,R\ge0\) and \(k\in C_c^\infty([-R,R])\),
\[
 \begin{split}
 Y&=R+\tfrac12\log(v+R+e)+B,\qquad
 K_a(v,R)=C(1+Y)^{p_a}e^{\kappa_aY},\\
 |F_k(a)|&\le K_a(v,R)
          [\|\Phi*k\|_E+e^{-Av}\|k\|_E].
 \end{split}                                               \tag{AC2}
\]
In particular, \(\log K_a(v,R(v))=o(v)\) whenever \(R(v)=o(v)\).
Suppose \(\delta>0\) and the actual adaptive inputs satisfy eventually
\[
 \operatorname{supp}h_X\subset[-R_X,R_X],\quad R_X=o(\log X),
 \quad\|h_X\|_E=X^{o(1)},\quad |F_{h_X}(a)|\ge X^{-o(1)}.
                                                               \tag{AC3}
\]
Then
\[
 \limsup_{X\to\infty}
   \frac{\log(1+\|\Phi*L_Xh_X\|_E)}{\log X}\ge\delta .       \tag{AC4}
\]
For \(f_v=\sqrt{a_t}e^{-v/2}(\Phi*L_{e^{2v}}h_{e^{2v}})(\cdot+v)\),
the original bimodal norm in (BMF2) satisfies
\[
 \limsup_{v\to\infty}v^{-1}\log(1+\|f_v\|_{\phi_{v,t(v)}})
     \ge2\delta\qquad\text{for every }t(v)>0.              \tag{AC5}
\]
The corresponding input has subexponential norm if, in addition,
\((R_X+B_0)^2/t(X)=o(\log X)\). Output (AC5) requires no such input
condition.
\end{theorem}

\begin{proof}
Choose \(B_0>0\) with \(m_0=\int_{|u|\le B_0}\Phi(u)\,\dd u>0\).
The actual Gaussian ground has the global lower bound
\[
 p_\tau(y)\ge m_0a_\tau e^{-(|y|+B_0)^2/(4\tau)},
 \quad\tau=Y+B_0+T_\Phi .                               \tag{AC6}
\]
For \(f=\Phi*k\), weighted Cauchy--Schwarz gives
\(\int|k|\le2e^{R/4}\|k\|_E\). The full theta envelope therefore gives
\(|f(y)|\le Ce^{CR+M|y|}e^{-ce^{2(|y|-R)}}\|k\|_E\) outside \(R+1\).
It also applies to every fixed derivative of \(f\), so the following
heat and graph domains are actual domains. In
\(\mathcal H_\tau=L^2(c/p_\tau\,\dd y)\), the squared normalized tail
is bounded by
\[
 Ce^{CR}\|k\|_E^2\int_{|y|>Y}
  e^{C|y|+(|y|+B_0)^2/(4\tau)-2ce^{2(|y|-R)}}\,\dd y.
                                                               \tag{AC7}
\]
Writing \(|y|=Y+s\), the negative exponent is
\(-2c(v+R+e)e^{2B}e^{2s}\); the positive quadratic exponent is at
most \(CY+Cs+Cs^2\). Since \(\sup_{s\ge0}(s+s^2)e^{-2s}<\infty\),
choose fixed \(B\) sufficiently large to pay these terms, \(CR+CY\),
and \(2Av\). This works for all \(R\), including \(R\gg v\), and
pays both infinite endpoints. On the core, (AC6) and
\(c(y)e^{y/2}\le e^Y\) give
\[
 \sqrt{a_\tau}\|f\|_{\mathcal H_\tau}
   \le Ce^{5Y/8}\|f\|_E+Ce^{-Av}\|k\|_E .                \tag{AC8}
\]
The support of \(k\) has not been used as the support of \(f\).

Since \(F_f=\Xi F_k\), all \(m_a\) zero jets vanish. The point-division
theorem (HD3)--(HD5), applied successively on its actual ranges, costs
\((2\sqrt\tau)^{m_a}\) and gives the selected value
\(\Xi^{(m_a)}(a)F_k(a)/m_a!\). Its coefficient is nonzero by the
actual order. Evaluation in \(\mathcal H_\tau\) has norm bounded by
\[
 \left(\int\frac{p_\tau}{c}e^{2\delta y}\,\dd y\right)^{1/2}
   \le[4\Xi(ib_a)e^{(b_a^2-1/4)\tau}]^{1/2}.
\]
Here \(1/c\le2e^{-|y|/2}\) and the two exact heat moments at
\(\pm b_a\) pay both signs. After division by \(\sqrt{a_\tau}\),
the evaluation price is \(C_a\tau^{1/4}e^{b_a^2\tau/2}\).
Combine this with (AC8) and \(\tau=Y+O(1)\) to obtain (AC2).

Now \(k_X=L_Xh_X\) has radius at most \(R_X+\log q\),
\(\|k_X\|_E\le C_q\sqrt X\|h_X\|_E\), and exactly
\[
 F_{k_X}(a)=X^{-ia}\Delta_X(1/2-ia)F_{h_X}(a).             \tag{AC9}
\]
For \(s_a=1/2-ia\), the fixed scalar identity
\(\int_1^\infty\Delta_X(s_a)X^{-w-1}\,\dd X
 =(1-q^{-w})/[w\zeta(s_a+w)]\)
has an order-\(m_a\) pole at \(w=0\), with nonzero numerator limit
\(\log q\). Thus \(\Delta_X(s_a)\) cannot satisfy any eventual
\(O(X^{-\eta})\) bound, \(\eta>0\). Along arbitrarily large \(X\),
combine this fact, (AC3), \(|X^{-ia}|=X^\delta\), and
\(\log K_a=o(\log X)\). Choosing \(A=4\) makes the remainder in
(AC2) at most \(X^{-3/2+o(1)}\). This proves (AC4), using no adaptive
vector Mellin factorization. The exact two-copy lower (BMF8) gives
\(\|f_v\|_{\phi_{v,t}}^2\ge(4M_0)^{-1}\|\Phi*L_Xh_X\|_E^2\)
at every original \(t>0\), proving (AC5).

Finally, the nearby copy in (AC6), with the original heat \(t\), pays
the actual preparation bound
\[
 \|\sqrt{a_t}e^{-v/2}h_X(\cdot-v)\|_{\phi_{v,t}}
   \le C e^{R_X/2+(R_X+B_0)^2/(8t)}\|h_X\|_E .           \tag{AC10}
\]
The asserted additional condition makes its logarithm \(o(\log X)\).
Auxiliary \(\tau\) and original \(t\) are distinct parameters.
\end{proof}

The theorem permits the cofactor to respond to the actual arithmetic
coefficients while retaining visibility of every selected off-line node.
An eventual subpower upper for the same adaptive completed windows
would contradict (AC4). A favorable subsequence alone does not
contradict its limsup. The required arithmetic upper remains open.

\subsection{Every zero is visible in a fixed Fourier band}
\label{sec:fixed-band-source}

The logarithmic window admits an exact Mellin identity on the full
weighted space. This gives a cumulative energy statement for one
fixed input, independently of the adaptive preparation theorem.
Throughout this subsection fix \(q>1\), \(Q=\log q\), and write
\[
 E=L^2(\R,e^{-y/2}\dd y),\qquad J_Eh=e^{-y/4}h,\qquad
 T_rh(y)=h(y+r),\qquad C_\Phi h=\Phi*h .
\]
Thus \(J_E\) is unitary onto ordinary \(L^2\),
\(J_ET_r=e^{r/4}T_rJ_E\), and
\(\|T_r\|_{E\to E}=e^{r/4}\).
The full theta tails imply boundedness of \(C_\Phi\) and its
first derivative on \(E\). Let \(\Pi_H\) denote the orthogonal
output projection defined by
\[
 \widehat{J_E\Pi_Hf}(u)=\mathbf1_{[-H,H]}(u)\widehat{J_Ef}(u),
 \qquad H>0 .
\]
This projection is on \(E\); it imposes no new compact support
condition on a Weil-form input. For \(X\ge1\), set
\[
 \mathcal U_Xh=
 \sum_{X/q<n\le X}\frac{\mu(n)\log(X/n)}{\sqrt n}
                     T_{\log(n/X)}C_\Phi h,\qquad
 \mathcal A_q(z)h=\int_0^Q t e^{-zt}T_{-t}C_\Phi h\dd t .
 \tag{FB1}\label{eq:fixed-band-source}
\]
All indices are integers, with the displayed endpoint conventions.

\begin{lemma}[The exact Mellin product and nonvanishing numerator]
\label{lem:log-source-mellin}
The map \(z\mapsto\mathcal A_q(z)\) is entire in operator norm.
For every \(h\in E\) and \(\Re z>1/2\), the following Bochner
integral is absolutely convergent:
\[
 \int_1^\infty X^{-z-1}\mathcal U_Xh\dd X
       =\frac{\mathcal A_q(z)h}{\zeta(1/2+z)} .
 \tag{FB2}\label{eq:log-source-mellin}
\]
For each \(z\in\C\) and \(h\ne0\), \(\mathcal A_q(z)h\ne0\).
If \(\Pi_Hh\ne0\), then also \(\Pi_H\mathcal A_q(z)h\ne0\).
In particular, the latter conclusion holds for every \(H>0\)
when \(h\) is a nonzero compact smooth function.
\end{lemma}
\begin{proof}
The weighted convolution identity gives
\(\widehat{J_EC_\Phi h}(u)=\Xi(u-i/4)\widehat{J_Eh}(u)\).
Boundedness of the derivative convolution and the fundamental
theorem of calculus give
\(\|T_{-t}C_\Phi-T_{-s}C_\Phi\|\le C_Q|t-s|\) on \([0,Q]\).
The integral in (FB1) is consequently an operator-norm Bochner
integral. Its Taylor series in \(z\), including every derivative,
converges uniformly on compact sets.

Each integer \(n\) contributes precisely for \(n\le X<qn\).
Substitute \(X=ne^t\). The sum of the integrals of the \(E\)-norms
is at most
\[
 \|C_\Phi\|\|h\|_E
 \sum_{n\ge1}n^{-1/2-\Re z}
       \int_0^Q t e^{-(\Re z+1/4)t}\dd t<\infty .
\]
Absolute Fubini and the reciprocal Dirichlet series for zeta
therefore give (FB2). The source is strongly measurable: on a
compact \(X\)-interval it has finitely many continuous translation
terms and finitely many endpoint changes.

Put
\[
 G_Q(\lambda)=\int_0^Q t e^{-\lambda t}\dd t
   =\frac{1-(1+Q\lambda)e^{-Q\lambda}}{\lambda^2},
 \qquad G_Q(0)=Q^2/2 .
\]
The exact numerator multiplier is
\[
 \widehat{J_E\mathcal A_q(z)h}(u)
   =\Xi(u-i/4)\widehat{J_Eh}(u)G_Q(z+1/4+iu).
 \tag{FB3}\label{eq:fixed-band-numerator}
\]
The two entire multipliers are not identically zero and have
only discrete real-\(u\) zero sets. The Fourier transform of a
nonzero \(L^2\) function is nonzero on a set of positive measure;
it cannot be supported on those discrete sets. The same argument
applies inside \([-H,H]\) when \(\Pi_Hh\ne0\). A nonzero compact
smooth input has an entire Fourier transform, which cannot vanish
on any nonempty real interval. This proves all assertions.
\end{proof}

\begin{theorem}[Fixed-band detection and cumulative multiplicity cost]
\label{thm:fixed-band-cumulative}
Fix a nonzero \(h\in C_c^\infty(\R)\) and \(H>0\). Set
\[
 F_h(X)=\|\Pi_H\mathcal U_Xh\|_E,\qquad
 \mathcal E_h(Y)=\int_1^Y F_h(X)^2\,\frac{\dd X}{X}.
 \tag{FB4}\label{eq:cumulative-log-source}
\]
If \(\rho=1/2+\delta+i\gamma\) is a nontrivial zeta zero of
exact order \(m\), with \(\delta>0\), define
\[
 z_\rho=\delta+i\gamma,\qquad
 N_\rho=\frac{m!\,\Pi_H\mathcal A_q(z_\rho)h}
                      {\zeta^{(m)}(\rho)} .
 \tag{FB5}\label{eq:fixed-band-laurent}
\]
Then \(N_\rho\ne0\), and
\[
 \limsup_{Y\to\infty}
 \frac{\mathcal E_h(Y)}{Y^{2\delta}(\log Y)^{2m-2}}
 \ \ge\ \frac{\|N_\rho\|_E^2}{2\delta(2m-2)!}>0 .
 \tag{FB6}\label{eq:cumulative-multiplicity-lower}
\]
In particular, the limsup exponential rate in \(V=\log Y\)
is at least \(2\delta\).

There is one compact input canceling both original poles and
one \(H>0\), independent of \(\rho\) and \(m\), for which
\[
 \|\Pi_H\mathcal A_q(\delta+i\gamma)h\|_E
       \ge \frac{c_{h,q,H}}{1+|\gamma|}
       \quad(0\le\delta\le1/2,\ \gamma\in\R),
 \qquad c_{h,q,H}>0 .
 \tag{FB7}\label{eq:uniform-height-numerator}
\]
\end{theorem}
\begin{proof}
Lemma~\ref{lem:log-source-mellin} gives a genuine vector pole at
\(z_\rho\) of exact order \(m\), with leading Laurent vector
(FB5). No ordering or isolation uniform in the zeros is needed.

Write \(A(V)=\mathcal E_h(e^V)\) and let \(M\) be the limsup
on the left of (FB6). If \(M=\infty\), the result follows.
Otherwise, for each \(\eta>0\), above a finite \(V_0\),
\[
 A(V)\le(M+\eta)e^{2\delta V}V^{2m-2}.
\]
Weighted Cauchy--Schwarz and Stieltjes integration show that the
actual vector Mellin integral is holomorphic on \(\Re z>\delta\).
For a compact subset of this half-plane, choose
\(2\delta<s<2\min\Re z\); the displayed bound gives
\(\int_0^\infty e^{-sv}\dd A(v)<\infty\), and the remaining
exponential margin pays every derivative in \(z\).
Meromorphic uniqueness identifies this integral with (FB2).
A pole further right would already contradict this conclusion;
absence of such a pole has not been assumed.

At \(z=z_\rho+\epsilon\), \(\epsilon>0\), Cauchy--Schwarz gives
\[
 \left\|\int_0^\infty e^{-(z_\rho+\epsilon)v}
                    \Pi_H\mathcal U_{e^v}h\dd v\right\|_E^2
 \le \epsilon^{-1}\int_0^\infty
                   e^{-(2\delta+\epsilon)v}\dd A(v).
 \tag{FB8}
\]
Since \(A(0)=0\), integration by parts rewrites the last
integral as
\((2\delta+\epsilon)\int_0^\infty
 e^{-(2\delta+\epsilon)v}A(v)\dd v\).
Its initial finite interval is uniformly bounded as
\(\epsilon\downarrow0\); its remainder is at most
\[
 (2\delta+\epsilon)(M+\eta)(2m-2)!\,
                                      \epsilon^{-(2m-1)} .
\]
Multiplying (FB8) by \(\epsilon^{2m}\) and taking
\(\epsilon\downarrow0\), then \(\eta\downarrow0\), gives
\(\|N_\rho\|_E^2\le2\delta M(2m-2)!\). This proves (FB6).

For (FB7), choose any nonnegative, nonzero real
\(g\in C_c^\infty(\R)\), and put \(h=(D^2+1/4)g\),
\(D=-i\partial\). Then \(F_h(\pm i/2)=0\), and
\[
 W_h(0):=\Xi(-i/4)F_h(-i/4)
       =\frac3{16}\Xi(-i/4)\int g(y)e^{-y/4}\dd y>0 .
\]
The positive theta integral proves the last inequality.
Choose \(H>0\) and \(b>0\) so that
\(|W_h(u)|\ge b\) on \([-H,H]\).
For \(\lambda=r+i(\gamma+u)\), \(1/4\le r\le3/4\),
\[
 G_Q(\lambda)=-\frac{Qe^{-Q\lambda}}{\lambda}
                      +\frac{1-e^{-Q\lambda}}{\lambda^2}.
\]
Consequently
\[
 |G_Q(\lambda)|\ge\frac{Qe^{-3Q/4}}{2|\lambda|}
 \quad\text{if}\quad
 |\lambda|\ge 2(e^{3Q/4}+1)/Q .
\]
For sufficiently large \(|\gamma|\), this holds throughout
the band, and \(|\lambda|\le|\gamma|+H+3/4\). Plancherel
in (FB3) proves (FB7) there. On the remaining compact
\((\delta,\gamma)\)-rectangle, the continuous function
\[
 \int_{-H}^H
    |G_Q(\delta+1/4+i(\gamma+u))|^2\dd u
\]
is strictly positive: otherwise the entire function \(G_Q\)
would vanish on an interval. Its compact minimum is positive.
Using \(|W_h|\ge b\) and Plancherel completes the proof.
\end{proof}

\begin{corollary}[An exceptional-set lower bound in the source parameter]
\label{cor:source-exception-measure}
For the input and band in Theorem~\ref{thm:fixed-band-cumulative},
take \(0<\epsilon<\delta\) and \(X_0\ge1\), and set
\[
 B_\epsilon(Y)=
 \big|\{X\in[X_0,Y]:F_h(X)>X^\epsilon\}\big|.
\]
There is \(C_{h,q}>0\), independent of \(X\), such that
\[
 \limsup_{Y\to\infty}
 \frac{B_\epsilon(Y)}{Y^{2\delta}(\log Y)^{2m-2}}
 \ge
 \frac{\|N_\rho\|_E^2}{2\delta(2m-2)!\,C_{h,q}}>0 .
 \tag{FB9}
\]
\end{corollary}
\begin{proof}
The absolute finite-window estimate gives
\(\|\mathcal U_Xh\|_E\le C_q\sqrt X\|h\|_E\), since every
shift lies in \([-Q,0]\). Thus \(F_h(X)^2\le C_{h,q}X\).
Split (FB4) into the good and exceptional sets to obtain
\[
 \mathcal E_h(Y)\le C_0+
       \frac{Y^{2\epsilon}}{2\epsilon}+C_{h,q}B_\epsilon(Y).
\]
The good contribution is negligible on the scale in (FB6),
because \(\epsilon<\delta\). Theorem~
\ref{thm:fixed-band-cumulative} proves (FB9).
\end{proof}

\begin{proposition}[A quantitative dyadic exceptional-set implication]
\label{prop:dyadic-source-exceptions}
Fix a nonzero compact smooth input \(h\), a band \(H>0\), and
\(0\le\theta\le1\). Suppose that for each \(\epsilon>0\)
there is a measurable set \(\mathcal B_\epsilon\) and a constant
\(C_\epsilon\) such that, eventually,
\[
 F_h(X)\le C_\epsilon X^\epsilon
       \quad(X\notin\mathcal B_\epsilon),\qquad
 |\mathcal B_\epsilon\cap[T,2T]|
       \le C_\epsilon T^\theta .
 \tag{FB10}
\]
Then no nontrivial zero satisfies
\(\Re\rho>\max(1/2,\theta)\). In particular, (FB10) with
\(\theta\le1/2\) would imply RH.
\end{proposition}
\begin{proof}
On a dyadic interval, the exceptional part of the vector
Mellin norm is at most
\(C T^{\theta-1/2-\Re z}\), by the absolute
\(O(\sqrt X)\) source bound. The good part is at most
\(C_\epsilon T^{\epsilon-\Re z}\).
Summing dyadically pays holomorphy, including derivatives,
for \(\Re z>\max(\epsilon,\theta-1/2)\).
If \(\Re\rho>\max(1/2,\theta)\), choose
\(\epsilon<\Re\rho-1/2\). Holomorphy then contradicts the
genuine pole from Lemma~\ref{lem:log-source-mellin}.
Conjugation and the functional equation give the stated RH
implication. The sets here are measured in \(X\); no
modulation-height exceptional set is substituted for them.
\end{proof}

\subsection{Integer gamma completion and the exact logarithmic covariance}
\label{sec:integer-gamma-completion}

The preceding source has a finite gamma realization with a cap
that works simultaneously for every input in \(E\).
For \(X\ge q\), define
\[
 \begin{split}
 A_\Gamma(y)&=(4\pi^2e^{9y/2}-6\pi e^{5y/2})
                                      e^{-\pi e^{2y}},\\
 B_\Gamma(y)&=-2\pi e^{5y/2}e^{-\pi e^{2y}},\qquad
 A_\Gamma=B_\Gamma'+B_\Gamma/2 .
 \end{split}
 \tag{GC1}
\]
The theta identity proved above gives on the entire real axis
\(\Phi=\sum_{k\ge1}k^{-1/2}T_{\log k}A_\Gamma\), with normal
convergence on physical compact sets. For \(j=0,1\), put
\[
 \begin{split}
 \nu_0(d)&=\mu(d),\qquad
 \nu_1(d)=\mu(d)\log(X/d),\qquad X/q<d\le X,\\
 a_j(k)&=\sum_{\substack{d\mid k\\X/q<d\le X}}\nu_j(d),\qquad
 \kappa_j=\sum_{X/q<d\le X}\nu_j(d)/d,\qquad
 m_j=\sum_{X/q<d\le X}\nu_j(d),\\
 \widetilde D_j(u)&=-\sum_{X/q<d\le X}
                         \nu_j(d)(\{u/d\}-1/2).
 \end{split}
 \tag{GC2}
\]
Values at integers are the actual after-jump values. For an
integer \(M\ge X\), define bounded operators on \(E\) by
\[
 \begin{split}
 G_{j,X,M}h&=\sum_{k\le M}a_j(k)k^{-1/2}
                         T_{\log(k/X)}(A_\Gamma*h),\\
 \mathcal B_{j,X,M}h&=\kappa_j\sqrt M
                         T_{\log(M/X)}(B_\Gamma*h),\\
 \mathcal C_{j,X,M}&=G_{j,X,M}-\mathcal B_{j,X,M},\\
 \mathcal U_X^{(j)}h&=\sum_{X/q<d\le X}\nu_j(d)d^{-1/2}
                                      T_{\log(d/X)}C_\Phi h .
 \end{split}
 \tag{GC3}
\]
Thus \(\mathcal U_X^{(1)}=\mathcal U_X\).
Every integer \(k\), including every prime power, is retained
with all its divisor columns.

\begin{theorem}[Joint integer caps and absolute energy transfer]
\label{thm:joint-integer-gamma-completion}
For every integer \(T\ge\lceil X\rceil\), choose \(M=M_*(X,T)\)
as the smallest minimizer, among \(T\le k<2T\), of
\[
 |\widetilde D_0(k)|^2+Q^{-2}|\widetilde D_1(k)|^2 .
 \tag{GC4}
\]
This cap is independent of \(h\). For \(j=0,1\), uniformly on
the full space \(E\),
\[
 \|\mathcal C_{j,X,M}-\mathcal U_X^{(j)}\|_{E\to E}
 \le C_q\bigl(X^{1/4}T^{-1/4}+X^{5/4}T^{-3/4}\bigr).
 \tag{GC5}\label{eq:joint-gamma-error}
\]
At \(T=\lceil X^2\rceil\) the error is \(O_q(X^{-1/4})\).
There is an unconditional constant \(c_\mu>0\) such that,
for \(T=\lceil X^3\rceil\), every \(h\in E\), and all
sufficiently large \(X\),
\[
 \left|\|\mathcal C_{1,X,M}h\|_E^2-\|\mathcal U_Xh\|_E^2\right|
 \le C_q\bigl(e^{-c_\mu f_0(\log X)}+X^{-1}\bigr)\|h\|_E^2,
 \quad f_0(L)=L^{3/5}(\log L)^{-1/5}.
 \tag{GC6}\label{eq:joint-gamma-energy}
\]
In particular, this is absolute \(o(1)\) for every family with
\(\log(2+\|h_X\|_E)=O((\log X)^{1/3})\).
For a compact input with \(F_h(\pm i/2)=0\), the finite
head and boundary in (GC3) retain both of these pole moments.
\end{theorem}
\begin{proof}
We first prove the floor-field estimate. For \(r\ge1\), let
\[
 B_r(Y)=\sum_{\substack{n\le Y\\(n,r)=1}}\mu(n)/n .
\]
The finite identity
\(\sum_{n\le Y,(n,r)=1}\mu(n)\lfloor Y/n\rfloor\)
counts the positive integers at most \(Y\) having all prime
factors among the divisors of \(r\). It lies in \([0,Y]\).
The fractional remainder is at most \(\lfloor Y\rfloor\)
in absolute value, so \(-1\le B_r(Y)\le2\).
For \(A_r^{(j)}=\sum_{r\mid d}\nu_j(d)/d\), nonsquarefree
\(r\) contributes zero. For squarefree \(r\), set \(Y=X/r\).
Finite Abel summation gives exactly
\[
 \begin{split}
 A_r^{(0)}&=\frac{\mu(r)}r[B_r(Y)-B_r(Y/q)],\\
 A_r^{(1)}&=\frac{\mu(r)}r\int_{Y/q}^{Y}
                   [B_r(t)-B_r(Y/q)]\,\frac{\dd t}{t}.
 \end{split}
 \tag{GC7}
\]
Hence \(|A_r^{(0)}|\le3/r\) and
\(|A_r^{(1)}|\le3Q/r\).

The centered sawtooth has cycles-Fourier coefficients
\(-1/(2\pi i k)\). Its exact covariance is
\((d,e)^2/(12de)\). On the common period
\(P=\operatorname{lcm}(1,\ldots,\lfloor X\rfloor)\), this yields
\[
 V_j=P^{-1}\int_0^P|\widetilde D_j(u)|^2\dd u
       =\frac1{12}\sum_{r\le X}J_2(r)|A_r^{(j)}|^2,\qquad
 V_0\le3X/4,\quad V_1\le3Q^2X/4,
 \tag{GC8}
\]
where \(J_2(r)=r^2\prod_{p\mid r}(1-p^{-2})\) and
\(\sum_{r\mid n}J_2(r)=n^2\).
The grouped nonzero frequencies are reduced rationals \(a/b\),
\(b\le X\), with coefficient
\(A_b^{(j)}/(2\pi i(a/b))\).
Distinct frequencies differ by at least \(X^{-2}\).

For an interval \(I\) of length \(\ell\), set
\(\eta=(2X^2)^{-1}\), \(K_\eta(u)=\eta
(\sin(\pi\eta u)/(\pi\eta u))^2\), and let \(I^+\) extend
\(I\) by \(1/\eta\) at both ends. Then
\(W_I=(\pi^2/4)\mathbf1_{I^+}*K_\eta\ge1\) on \(I\),
\(\int W_I=(\pi^2/4)(\ell+4X^2)\), and its cycles-Fourier
transform is supported in \([-\eta,\eta]\).
Expanding a finite Fourier truncation annihilates all distinct
frequency differences in its \(W_I\)-integral.
The periodic series converges in \(L^2(I)\), so passage to
that limit proves
\[
 \int_I|\widetilde D_j|^2\dd u
       \le\frac{\pi^2}{4}(\ell+4X^2)V_j .
 \tag{GC9}
\]
No pointwise convergence at a jump has been used.

On the true right unit interval,
\(\widetilde D_j(k+s)=\widetilde D_j(k)-\kappa_js\),
\(0\le s<1\). Therefore
\[
 |\widetilde D_j(k)|^2
 \le2\int_k^{k+1}|\widetilde D_j(u)|^2\dd u
                                      +2|\kappa_j|^2/3 .
\]
The elementary bounds
\(|\kappa_0|\le Q+1\), \(|\kappa_1|\le Q(Q+1)\),
(GC9), and averaging the \(T\) integer endpoints in
(GC4) give simultaneously
\[
 |\widetilde D_j(M)|\le C_q
                  \bigl(\sqrt X+X^{3/2}T^{-1/2}\bigr).
 \tag{GC10}
\]

For a compact smooth input set
\(g_h=A_\Gamma*h\), \(b_h=B_\Gamma*h\), and
\(f(u)=u^{-1/2}T_{\log(u/X)}g_h\).
The full physical theta series can be grouped by \(k=d n\).
Abel summation above the integer \(M\), with
\(D_j(u)=\widetilde D_j(u)-m_j/2\), gives
\[
 \mathcal U_X^{(j)}h
 =\mathcal C_{j,X,M}h+R_{j,X,M}h,\qquad
 R_{j,X,M}h=-D_j(M)f(M)-\int_M^\infty D_j(u)f'(u)\dd u .
 \tag{GC11}
\]
To verify the moving boundary, (GC1) implies pointwise
\[
 \int_M^\infty f(u)\dd u
                 =-\sqrt M\,T_{\log(M/X)}b_h .
\]
The derivative of \(\sqrt u\,T_{\log(u/X)}b_h\) is \(f(u)\);
its upper endpoint vanishes by the physical gamma tail.
This calculation is pointwise for compact inputs; it does
not presume absolute \(E\)-integrability of the mean integral.

Weighted Young bounds for \(A_\Gamma,B_\Gamma\) and their
derivatives give, for all \(h\in E\),
\[
 \|f(u)\|_E\le C X^{-1/4}u^{-1/4}\|h\|_E,\qquad
 \|f'(u)\|_E\le C X^{-1/4}u^{-5/4}\|h\|_E .
\]
Thus \(f(\infty)=0\) strongly in \(E\), and \(f'\) is
absolutely integrable there. The constant \(m_j/2\)
cancels exactly between the endpoint and the integral in
(GC11). We may replace \(D_j\) by \(\widetilde D_j\).
By (GC9) and Cauchy--Schwarz on each \([Y,2Y]\),
\[
 \int_Y^{2Y}|\widetilde D_j(u)|u^{-5/4}\dd u
 \le C_q\bigl(\sqrt X\,Y^{-1/4}
                         +X^{3/2}Y^{-3/4}\bigr).
\]
Summing over \(Y=2^\ell M\) and using (GC10) pays the
absolute \(E\)-valued remainder and proves (GC5).
The identity and bound extend by density to the whole \(E\).

For the energy estimate we use only the unconditional Mertens
bound from \cite[Theorem 2.2, (23)]{LeeLeong}. Fix one
admissible parameter \(W\) in that theorem and reduce its
positive exponential constant if necessary. Its bounded
prefactor and the factor \(\log x\) are absorbed for all
large \(x\), giving
\[
 \left|\sum_{n\le x}\mu(n)\right|
 \le Cx e^{-c_0 f_0(\log x)}
 \quad(x\ge x_0),\qquad c_0>0 .
\]
Since \(f_0'(L)\to0\), one fixed \(c_\mu>0\) works uniformly
on \([X/q,X]\), with \(q\)-dependent constants and threshold.
Vector Abel summation for
\(\log(X/u)u^{-1/2}T_{\log(u/X)}C_\Phi h\),
using bounded derivative convolution, gives
\(\|\mathcal U_X\|_{E\to E}
 \le C_q\sqrt X e^{-c_\mu f_0(\log X)}\).
At \(T=\lceil X^3\rceil\), (GC5) gives an operator error
at most \(C_qX^{-1/2}\).
The inequality
\[
 \big|\|U-r\|^2-\|U\|^2\big|\le2\|U\|\|r\|+\|r\|^2
\]
now proves (GC6). Finally, the exponential moments of
\(A_\Gamma,B_\Gamma\) at \(\pm i/2\) are absolute.
Their finite translates preserve a compact input's zero
pole moments by convolution factorization. Higher moments
remain conditions on the original input; no corresponding
absolute moments of a separate finite gamma head are required.
\end{proof}

For a fixed input, (GC5) and the Abel bound prove (GC6)
also for the projected vectors, by contractivity of \(\Pi_H\).
They therefore transfer the cumulative
lower bound (FB6) to
\(\int_q^Y\|\Pi_H\mathcal C_{1,X,M_*(X,\lceil X^3\rceil)}
h\|_E^2\dd X/X\). On a bounded \(X\)-interval there are
finitely many candidate caps and divisor indices. Their scores
are measurable, and the specified finite minimum with its tie
rule makes \(M_*(X,\lceil X^3\rceil)\) measurable. The finite
completed source is therefore strongly measurable.
The energy error is \(o(1)\) per \(X\), hence \(o(\log Y)\)
after integration. This is negligible on the scale in (FB6).

\begin{theorem}[The exact logarithmic primitive-sector remainder]
\label{thm:log-primitive-remainder}
For \(h\in E\), set \(\psi=C_\Phi h\),
\(f=J_E\psi\), and
\[
 C_h(r)=\int f(y)\overline{f(y+r)}\dd y,\qquad
 \rho(M)=\zeta(2)^{-1}\prod_{p\mid M}\frac p{p+1}.
\]
Let \(E_{\le H}(X)\) be the sector of the exact expansion of
\(\|\mathcal U_Xh\|_E^2\) obtained by writing \(n=d a\),
\(k=d b\), \((a,b)=1\), and imposing \(\max(a,b)\le H\).
For \(1\le H\le X\), it has the main term
\[
 \begin{split}
 M_{\le H}(X)&=
 \sum_{\substack{a,b\le H\\(a,b)=1}}
 \frac{\mu(a)\mu(b)\rho(ab)}{\sqrt{ab}}\,
                         w_Q(\log(b/a))C_h(\log(b/a)),\\
 w_Q(r)&=\int_{|r|/2}^{Q-|r|/2}
                 e^{-s/2}(s^2-r^2/4)\dd s\quad(|r|<Q),\\
 w_Q(r)&=0\quad(|r|\ge Q).
 \end{split}
 \tag{GC12}
\]
Uniformly over \(h\in E\),
\[
 |E_{\le H}(X)-M_{\le H}(X)|
 \le C_q C_h(0)X^{-1/2}H^{3/2}\log^2(2H).
 \tag{GC13}\label{eq:log-primitive-error}
\]
If \(C_{h_X}(0)\le C\exp(B(\log X)^{1/3})\), then
\[
 H_X=X^{1/3}e^{-A(\log X)^{1/3}},\qquad A>2B/3,
 \tag{GC14}
\]
gives an absolute \(o(1)\) remainder.
\end{theorem}
\begin{proof}
Expansion of the finite source, with the linear first
inner-product slot, gives exactly
\[
 \|\mathcal U_Xh\|_E^2=X^{-1/2}
 \sum_{X/q<n,k\le X}
 \frac{\mu(n)\mu(k)}{(nk)^{1/4}}
 \log(X/n)\log(X/k)C_h(\log(k/n)).
 \tag{GC15}
\]
In a coprime slot \(a,b\), its remaining sum is
\[
 \sum_{\substack{X/[q\min(a,b)]<d\le X/\max(a,b)\\(d,ab)=1}}
 \frac{\mu(d)^2}{\sqrt d}
                      \log\frac X{da}\log\frac X{db}.
 \tag{GC16}
\]
The square-divisor expansion of \(\mu(d)^2\), followed by
finite inclusion--exclusion for \((d,M)=1\), proves
\[
 \sum_{\substack{d\le Y\\(d,M)=1}}\mu(d)^2
       =\rho(M)Y+O(2^{\omega(M)}\sqrt Y)
 \quad(Y>0).
 \tag{GC17}
\]
Indeed the inner coprime count is
\((\varphi(M)/M)Y/k^2+O(2^{\omega(M)})\);
there are \(O(\sqrt Y)\) square divisors, and the tail of
\(\sum k^{-2}\) costs \(O(\sqrt Y)\). The Euler product
of the main constant is precisely \(\rho(M)\).
For \(Y<1\) the same estimate follows from \(Y\le\sqrt Y\).

The weight in (GC16) has modulus \(O_q(d^{-1/2})\)
and derivative \(O_q(d^{-3/2})\) throughout the actual
nonempty cap interval, whose endpoint ratio is at most \(q\).
Stieltjes integration with (GC17) costs
\(O_q(2^{\omega(ab)})\) per slot, including both endpoints.
Replacing just the count main by \(\rho(ab)\dd d\) and
substituting \(d=Xe^{-s}/\sqrt{ab}\) proves (GC12).

For a coprime slot,
\(2^{\omega(ab)}=2^{\omega(a)}2^{\omega(b)}\), and
\[
 \sum_{n\le Y}2^{\omega(n)}
 \le\sum_{n\le Y}\tau(n)
 \le Y(1+\log Y).
\]
In a dyadic block \(A<\max(a,b)\le2A\) the support
condition forces \(\min(a,b)>A/q\). Thus
\((ab)^{-1/4}\le C_qA^{-1/2}\).
Using \(|C_h(r)|\le C_h(0)\), the total error in this
block is at most
\(C_q C_h(0)X^{-1/2}A^{3/2}\log^2(2A)\).
The geometric dyadic sum, with the initial singleton
included, proves (GC13). Substitution of (GC14) gives
\(\exp[-(3A/2-B)(\log X)^{1/3}]O_q((\log X)^2)=o(1)\).
\end{proof}

Combining (GC6) and (GC13) for the stated input family yields
the genuine finite signed identity
\[
 \|\mathcal C_{1,X,M}h_X\|_E^2
       =M_{\le H_X}(X)+E_{>H_X}(X)+o(1),\qquad
 E_{>H_X}(X)=\|\mathcal U_Xh_X\|_E^2-E_{\le H_X}(X).
 \tag{GC18}
\]
The complete boundary cross term in
\(\|G_{1,X,M}h_X-\mathcal B_{1,X,M}h_X\|_E^2\)
is included. Neither primitive sector is asserted to be
positive separately, and neither signed term is bounded
subpolynomially by this identity.

\subsection{The staged derivative estimate gives a cubic short cofactor}
\label{sec:cubic-prepared-source}
The conservative maximum in (NL19) loses one factor of the
order in the cofactor budget. We retain the separately paid
derivative orders and prove a sharper complete head estimate.

\begin{theorem}[Prepared cubic cofactor, uniformly at every heat time]
\label{thm:cubic-prepared-source}
Fix a finite PN basis as in Theorem~\ref{thm:prepared-nearlinear-source}
and a desired \(A>0\) before choosing the preparation and its \(B\).
Use the all-heat prepared family (NHA2) in the normalized
ground \(\widetilde\phi_{v,t}=\phi_{v,t}/a_t\), \(t\ge1\).
Put
\[
 p=\lfloor v/2\rfloor\ge2,\quad
 W=\tfrac12\log(v+p+3)+B+1,\quad
 a_-=2v-2W,\quad a_+=2v+2W,\quad N=\lceil e^{a_+}\rceil.
                                                               \tag{CPH1}
\]
For \(f=\tau_{-v}f_v[a]\), \(h=\tau_vf_v[b]\), write
\(C(r)=\int\overline{f(x)}(J_{\mathrm n}h)(x+r)\dd x\).
There are fixed constants \(C_1,J_0>0\), depending only on
the fixed basis and the preparation, such that the actual cap
\[
 \begin{split}
 E_p^\sharp&=\max\left\{1,\frac{4J_0}{p-1/2}
                 \left(\frac{(C_1+1)p}{2\pi}\right)^p\right\},\\
 D^\sharp&=e^{(2-1/p)(v-W)}
                          [E_p^\sharp\log(2N)]^{-1/p}
 \end{split}                                                   \tag{CPH2}
\]
satisfies, for all sufficiently large \(v\),
\[
 1<D^\sharp<N,\qquad
 K_{\rm short}=\lfloor N/D^\sharp\rfloor\le C_Bv^3,\qquad
 |T_{\le D^\sharp}(C)|
      \le e^{-v}\|f\|_{\widetilde\phi}\|h\|_{\widetilde\phi}.
                                                               \tag{CPH3}
\]
All original moments, every prime power and the cofactor-one
channel of the remaining long sum are retained.
Both ordinary sources (NL16), with \(D^\sharp\) in place of \(D\),
have the simultaneous paid window \(H=C_Av\) and error
\(C_Ae^{-Av}\|f\|_{\widetilde\phi}\|h\|_{\widetilde\phi}\),
uniformly for all \(t\ge1\), for this chosen \(A\).
Uniformity in a growing PN basis is not asserted.
\end{theorem}
\begin{proof}
The actual preparation has support
\(C\in C_c^\infty([a_-,a_+])\), all endpoint jets vanish,
and \(\int e^{r/2}C(r)\dd r=0\).
Put every derivative on \(J_{\mathrm n}h=\tau_vz_v[b]\).
The proof before the maximum in (NL19) gives, individually,
\[
 \|C^{(j)}\|_\infty
  \le C_{\rm jet}(C_1p)^j|a||b|\quad(0\le j\le p),\qquad
 |a||b|\le C_{\rm norm}e^{-v}
                 \|f\|_{\widetilde\phi}\|h\|_{\widetilde\phi}.
                                                               \tag{CPH4}
\]
Indeed \(\|f_v[a]\|_1\le C|a|\); the analytic main jets use
\(j!\le p^j\), while (NL17) pays each first-\(p\) cutoff and
corrector jet. The second inequality is the actual fixed-basis
norm lower in (NHA2). Set \(J_0=C_{\rm jet}C_{\rm norm}\).

For an integer \(d\ge1\), extend
\(F_d(k)=k^{-1/2}C(\log(dk))\) by zero at \(k\le0\).
It is genuinely smooth and compactly supported away from zero.
Induction gives the exact derivative, with commuting factors,
\[
 F_d^{(p)}(k)=k^{-p-1/2}
          \prod_{j=0}^{p-1}(\partial_r-j-1/2)C(r),
                         \quad r=\log(dk).                    \tag{CPH5}
\]
The absolute coefficient polynomial is
\(\prod_{j=0}^{p-1}(X+j+1/2)\). Insert each order in (CPH4)
with \(X=C_1p\), rather than using a common derivative maximum.
It is at most \([(C_1+1)p]^p\), so
\[
 \|F_d^{(p)}\|_1
 \le\frac{C_{\rm jet}[(C_1+1)p]^p|a||b|}{p-1/2}
                              (e^{a_-}/d)^{1/2-p}.            \tag{CPH6}
\]
The last bound integrates \(k^{-p-1/2}\) from its actual
lower support endpoint to infinity. It needs no support-length
factor \(4W\) and no second factorial.

For completeness, periodize any \(F\in C_c^\infty(\R)\).
The resulting smooth one-periodic function has Fourier coefficients
\(\widehat F(2\pi\ell)\), where the hat is our ordinary
Fourier integral. Its absolutely convergent Fourier series at zero,
followed by \(p\) integrations by parts, gives for every integer
\(p\ge2\), both odd and even,
\[
 \left|\sum_{n\in\Z}F(n)-\int F\right|
 \le\frac{2\zeta(p)}{(2\pi)^p}\|F^{(p)}\|_1
 \le\frac4{(2\pi)^p}\|F^{(p)}\|_1 .                          \tag{CPH7}
\]
Only the absolutely convergent series
\(\zeta(p)=\sum_{\ell\ge1}\ell^{-p}<2\) is used here.
For \(F_d\) the complete integral is exactly
\(d^{-1/2}\int e^{r/2}C(r)\dd r=0\); every nonpositive
integer term vanishes. The actual identity
\(\Lambda=-(\mu\cdot\log)*1\), with all prime powers, therefore yields
\[
 \begin{split}
 T_{\le D}(C)
 &=-\sum_{d\le D}\mu(d)\log d\,d^{-1/2}
                         \sum_{k\ge1}F_d(k),\\
 |T_{\le D}(C)|
 &\le E_p^\sharp D^p\log(2D)e^{-(p-1/2)a_-}e^{-v}
                         \|f\|_{\widetilde\phi}\|h\|_{\widetilde\phi}.
 \end{split}                                                   \tag{CPH8}
\]
Here \(\sum_{d\le D}d^{p-1}\log d\le D^p\log(2D)\)
is a plain absolute bound after the zero continuum term.
No M\"obius cancellation is a premise.
At (CPH2) its complete prefactor is
\(\log(2D^\sharp)/\log(2N)\le1\), proving the head in (CPH3).
The positive leading \(2v-O(\log v)\) gives \(D^\sharp>1\),
while \(E_p^\sharp\ge1\), \(v>W\), and \(\log(2N)>1\)
give \(D^\sharp<N\).
Since \((E_p^\sharp)^{1/p}=O(p)\),
\[
 \frac{N}{D^\sharp}
 \le2e^{4W+(v-W)/p}[E_p^\sharp\log(2N)]^{1/p}
 \le C_Bv^3.                                                 \tag{CPH9}
\]
This includes the actual ceiling: \(N\le2e^{a_+}\),
\(e^{4W}=O_B(v^2)\), and the other exponential and logarithmic
factors are bounded after taking the \(p\)-th root.

The finite paired completion (NL16) is valid for every actual
\(1<D<N\), so it applies at \(D^\sharp\) without alteration.
Its absolute UV bound \(|M^{\log}_{D,N}|\le2\sqrt N\log N\)
does not depend on the lower cap. The first ordinary source is
unchanged. Thus the already paid fixed-strip main decay,
all-order cutoff decay and normalized prefactor in (NL21)--(NL22)
and (NHA2) give both claimed \(H=C_Av\) errors, uniformly in heat.
\end{proof}

Taking \(A=2\) and \(t=1\), the same individual forcing (NL24)
now lies in the actual smaller retained range
\(D^\sharp<d\le N\), \(k\le O(v^3)\).
Consequently the conditional exclusion (NL25) remains valid
with this cap and the normalized norm.
The cubic head bound proves no signed upper on that long range;
its cofactor-one channel is still present.

\part{Theta approximation and certified zero control}

\section{Circular Euler heat: a global zero theorem for a modified Xi}
The next result concerns a genuinely changed entire function.
Its heat averages rotations of the argument, rather than translations
of the Fourier density. We prove its global zero statement and the
precise obstruction to interpreting the result as a positive
Fourier-measure construction.

\begin{lemma}[A rational zero-free rectangle]\label{lem:euler-low-box}
For the complete Riemann density $\Phi$ in the normalization above,
\[
  \Re\Xi(x+iy)>\frac{41}{2916}
       \qquad (|x|\le5/3,\ |y|\le1/2).                    \tag{EUL1}
\]
Consequently every zero $z=\gamma+i\delta$ of $\Xi$ satisfies
$|\gamma|>5/3$, $|\delta|<1/2$, and
\[
 |\arg z^2|<\phi_0=2\arctan(3/10)<\pi/2,\qquad
             \cos\phi_0=\frac{91}{109}.                 \tag{EUL2}
\]
\end{lemma}
\begin{proof}
Write $M(b)=\int_{\R}\Phi(u)\cosh(bu)\dd u=\Xi(ib)$.
The exact value of $\xi$ at its removable pole gives $M(1/2)=1/2$.
For $u\ge0$, every term of the full theta density is positive, and
\[
 0<\Phi(u)
 \le4\pi^2e^{9u/2}\sum_{n\ge1}n^4e^{-\pi n^2e^{2u}}
 <80e^{9u/2}e^{-3e^{2u}}.                               \tag{EUL3}
\]
Here all theta terms have been retained. The constants can be justified
elementarily: polynomial division gives
$\int_0^1x^4(1-x)^4/(1+x^2)\dd x=22/7-\pi>0$, hence
$\pi^2<10$; the lower eight-term geometric head for $1/(1+x^2)$
gives $4\sum_{j=0}^7(-1)^j/(2j+1)=135904/45045>3$.
Convexity of $1/x$ on $[1,2]$ gives $\log2<3/4$.
The derivative of $v^2-1-4\log v$ is positive for $v\ge2$,
so $4\log n<n^2-1$ for $n\ge2$. Therefore
\[
 \sum_{n\ge1}n^4e^{-3(n^2-1)}
 \le1+\sum_{n\ge2}e^{-2(n^2-1)}
 \le1+\sum_{n\ge2}e^{-6(n-1)}<2,
\]
which proves (EUL3) on the full half-line.

Put $w=e^{2u}$ and use $\cosh(u/2)\le e^{u/2}$,
$\log w\le w-1$, and $w^{3/2}\le w^2$. Evenness yields
\[
 \begin{split}
 M''(1/2)
 &<20\int_1^\infty(\log w)^2w^{3/2}e^{-3w}\dd w\\
 &\le20e^{-3}\int_0^\infty v^2(1+v)^2e^{-3v}\dd v\\
 &=20e^{-3}\left(\frac2{27}+\frac{12}{81}+\frac{24}{243}\right)
 <\frac{26}{81}.
 \end{split}                                             \tag{EUL4}
\]
The last strict inequality follows from
$\sum_{j=0}^8 3^j/j!=89641/4480>20$.
These differentiations and integrals converge absolutely by (EUL3).
Twice integrating the second derivative of $\cosh$ gives
$\cosh a-1\le a^2\cosh a/2$ and hence
$M(0)\ge1/2-M''(1/2)/8$.
For $|y|\le1/2$ we have $M(y)\ge M(0)$ and
$M''(y)\le M''(1/2)$.
The inequality $\cos(xu)\ge1-x^2u^2/2$ now gives
\[
 \Re\Xi(x+iy)\ge M(y)-\tfrac12x^2M''(y)
 >\frac12-\left(\frac18+\frac{25}{18}\right)\frac{26}{81}
 =\frac{41}{2916}.
\]
Every nontrivial zero lies in the open critical strip.
Its square has the argument and cosine bound (EUL2).
\end{proof}

For any entire $F(z)=\sum_{j\ge0}c_jz^j$ and $\epsilon>0$ define
\[
 \mathcal C_\epsilon F(z)=
 \frac1{\sqrt{4\pi\epsilon}}\int_{\R}
    e^{-u^2/(4\epsilon)}F(ze^{iu})\dd u
 =\sum_{j\ge0}c_je^{-\epsilon j^2}z^j,
 \qquad \mathcal C_0F=F.                                \tag{EUL5}
\]
The integral is bounded on $|z|\le R$ by $\|F\|_R$ times a
Gaussian density. Normal Taylor integration is valid there, and
$\|\mathcal C_\epsilon F\|_R\le\|F\|_R$.
The operators form a semigroup with generator
$-(z\partial_z)^2$. In particular they preserve entire functions
and their values at zero.

\begin{lemma}[The complete polynomial sector flow]
\label{lem:euler-sector-flow}
Let $Q$ be a real polynomial with $Q(0)\ne0$, all of whose roots
belong to $|\arg w|\le\theta<\pi/2$. Apply the coefficient multipliers
$e^{-4\epsilon n^2}$ to $Q(w)=\sum_nq_nw^n$.
For every $\epsilon\ge-\frac14\log\cos\theta$, all roots of
the resulting polynomial are positive real. Initial and final
multiplicities are allowed.
\end{lemma}
\begin{proof}
Write $Q_\epsilon=\mathcal C_{4\epsilon}Q$.
Its equation is
$\partial_\epsilon Q_\epsilon=-4(w\partial_w)^2Q_\epsilon$.
For a simple root $w_j$ implicit differentiation gives the exact identity
\[
 \frac{w_j'}{w_j}
     =4+8\sum_{\ell\ne j}\frac{w_j}{w_j-w_\ell}.           \tag{EUL6}
\]
If $w_j$ has the largest positive angle $\phi\in(0,\pi/2)$,
then $\arg(w_\ell/w_j)\in[-2\phi,0]$ for every $\ell$.
Thus each term of the sum has nonpositive imaginary part:
$\Im[1/(1-v)]=\Im v/|1-v|^2\le0$.
The conjugate root specifically contributes
$\Im[w_j/(w_j-\overline w_j)]=-\tfrac12\cot\phi$.
At every simple-root time, each branch attaining the angle maximum obeys
\[
       \phi'\le-4\cot\phi,\qquad
       \frac{\dd}{\dd\epsilon}\cos\phi\ge4\cos\phi.        \tag{EUL7}
\]

We justify integration through all collisions.
At an initial nonzero root $a$ of order $m$, write
$Q(w)=c(w-a)^m+O((w-a)^{m+1})$.
In the exact rotation integral for $Q_\epsilon$, set
$w=a+2a\sqrt\epsilon\,v$ and $u=2\sqrt\epsilon\,s$.
After division by $c(2a\sqrt\epsilon)^m$, the integral tends
normally on compact $v$-sets to
\[
 H_m(v)=\frac1{\sqrt{4\pi}}\int_{\R}(v+is)^me^{-s^2/4}\dd s.
                                                               \tag{EUL8}
\]
For $|u|\le u_0$ the Taylor majorant and Gaussian moments
justify this limit. The remaining circle is bounded and contributes
$O(e^{-u_0^2/(32\epsilon)}\epsilon^{-m/2})$.
The polynomial $H_m$ is the monic Hermite polynomial for
the weight $e^{-v^2/4}$; its Rodrigues formula proves orthogonality
to lower-degree polynomials. The usual sign-change argument for
a positive weight then proves that its $m$ roots are distinct and real.
Separate Rouch\'e circles around these roots give $m$ distinct
nearby roots of $Q_\epsilon$. This holds also when $a$ is complex.

The discriminant is consequently analytic in $\epsilon$ and not
identically zero. On any compact time interval it has only finitely
many zeros. The leading and constant coefficients remain nonzero;
Cauchy bounds for the polynomial and its reciprocal keep every
root bounded and separated from zero on that interval.
The root multiset and its maximum absolute angle are continuous.
Between collisions the maximum is locally Lipschitz, and (EUL7)
holds for its upper derivative. Integration over the finitely many
pieces, using continuity at their ends, gives
\[
 \log\cos\phi(\epsilon)-\log\cos\phi(0)\ge4\epsilon
\]
while a nonreal root remains. A first exit from the initial sector
is impossible. The inequality forces $\phi=0$ by
$\epsilon=-\frac14\log\cos\theta$.
Real roots cannot subsequently leave the positive axis: on any
first component of positive angle starting at angle zero,
the same inequality would give $\cos\phi>1$.
Continuity includes the exact time boundary and its multiplicities.
\end{proof}

The general sector reduction in this lemma is already present in
\cite[Theorem 2.5 and Corollary 2.6, p.~5]{SectorReducing}.
That source uses multipliers $e^{-a^2n^2/2}$ and the sector rule
$\cos\gamma=\min\{1,e^{a^2/2}\cos\theta\}$.
Here $a^2/2=4\epsilon$. Its entire-function version requires
normal approximation by polynomials in the same fixed sector;
we establish that hypothesis for the complete Xi function next.

\begin{theorem}[Global reality after fixed circular Euler heat]
\label{thm:euler-xi-real}
For every
\[
                  \epsilon\ge\frac14\log\frac{109}{91},  \tag{EUL9}
\]
the actual entire function $\mathcal C_\epsilon\Xi$ has only real
zeros, at every height and with arbitrary multiplicity.
It is a normal limit of real-rooted polynomials.
\end{theorem}
\begin{proof}
The even Taylor series defines the entire function
$H(w)=\Xi(\sqrt w)$ without a square-root branch.
For $R\ge0$ and $a=R/2+9/4$, (EUL3) gives
\[
 \max_{|z|\le R}|\Xi(z)|
 \le80\,3^{-a}\int_3^\infty t^{a-1}e^{-t}\dd t
 \le80(\lceil a\rceil-1)!.
\]
Thus $\log M_\Xi(R)\le C(1+R)\log(2+R)$ and $\Xi$ has
order at most one. The order of $H$ is at most $1/2<1$.
Hadamard factorization for a nonzero finite-order entire function
gives an exponential polynomial factor of degree at most the
integer part of its order; see \cite[theorem (1), p.~1]{ElkiesHadamard}.
There is no zero at the origin by Lemma~\ref{lem:euler-low-box}.
The factor for $H$ is therefore constant. Equivalently one may pair
the genus-one $\Xi$ factors in \cite[Theorem 7.7, p.~58]{Petrow}.
We obtain the exact normal product
\[
 H(w)=\Xi(0)\prod_{\substack{z=\gamma+i\delta:\ \Xi(z)=0\\
                          \gamma>0}}
                  \left(1-\frac{w}{z^2}\right)^{m_z}.    \tag{EUL10}
\]
Each positive-height representative occurs with its multiplicity;
both conjugate representatives occur when they differ.
Their negative-height partners are accounted for by evenness.
The unconditional count $N(T)=O(T\log(2+T))$ gives
$\sum_zm_z/|z|^2<\infty$, so these products converge normally.
Every product root lies in the fixed sector (EUL2).
Take finite products closed under conjugation. Their corresponding
even polynomials converge normally to $\Xi$.
Lemma~\ref{lem:euler-sector-flow} makes all their heated roots
real at every time (EUL9).
The disk contraction in (EUL5) preserves normal convergence.
The limiting function is not identically zero, since its value
at zero is $\Xi(0)>0$. Hurwitz on every disk disjoint from the
real axis excludes each nonreal limit zero.
\end{proof}

\begin{theorem}[No positive Fourier measure after circular heat]
\label{thm:euler-no-positive-measure}
For every $\epsilon>0$, the restriction of
$\mathcal C_\epsilon\Xi$ to $\R$ is not the Fourier transform
of any finite positive measure on $\R$.
\end{theorem}
\begin{proof}
Let $\mu_{2n}=\int_{\R}u^{2n}\Phi(u)\dd u>0$.
The complete superexponential theta tail implies
$\Phi(u)\le Ke^{-u^2}$ for one finite $K$ and all real $u$.
Hence
$\mu_{2n}\le K\Gamma(n+1/2)\le K\sqrt\pi\,n!$.
The Taylor series of the actual heated function $f$ gives
\[
 (-1)^nf^{(2n)}(0)=\widetilde\mu_{2n}
        =e^{-4\epsilon n^2}\mu_{2n},\qquad
 \widetilde\mu_{2n}^{1/(2n)}
        \le C\sqrt n\,e^{-2\epsilon n}\longrightarrow0. \tag{EUL11}
\]
Suppose $f(t)=\int e^{-itx}\dd\nu(x)$ for a finite positive
measure, without any initial moment assumption.
With $\Delta_h^cf(t)=f(t+h)-2f(t)+f(t-h)$ we have exactly
\[
 \frac{(-\Delta_h^c)^nf(0)}{h^{2n}}
       =\int\left(\frac{2\sin(hx/2)}h\right)^{2n}\dd\nu(x).
                                                               \tag{EUL12}
\]
Its Taylor limit is $\widetilde\mu_{2n}<\infty$.
Fatou first proves that $\int x^{2n}\dd\nu$ is finite.
The bound $|2\sin(hx/2)/h|\le|x|$ then permits dominated
convergence in (EUL12), giving
$\int x^{2n}\dd\nu=\widetilde\mu_{2n}$.
For every $r>0$,
$\widetilde\mu_{2n}\ge r^{2n}\nu(\{|x|\ge r\})$.
The root limit in (EUL11) forces $\nu(\{|x|\ge r\})=0$
for every $r>0$, so $\nu$ is supported at zero.
Its second moment is zero, contradicting
$\widetilde\mu_2=e^{-4\epsilon}\int u^2\Phi(u)\dd u>0$.
\end{proof}

\subsection*{Every positive Euler time has a controlled complete high tail}
The next coefficient argument removes the fixed lower time threshold
from the high tail, with its low core retained in the quantifiers.

\begin{lemma}[The full theta moments and their curvature]
\label{lem:euler-moment-curvature}
Put $\mu_{2n}=\int_\R u^{2n}\Phi(u)\dd u$ and
$b_n=\mu_{2n}/(2n)!$. Then
\[
 \frac{\mu_{2n+2r}}{\mu_{2n}(\tfrac12\log n)^{2r}}\longrightarrow1
       \quad(r=1,2),\qquad
 \frac{b_{n+1}}{b_n}\sim\frac{(\log n)^2}{16n^2}.         \tag{EUT1}
\]
There are fixed $C,N_0$ such that
\[
 \left|\log\frac{b_{n+2}b_n}{b_{n+1}^2}\right|
           \le C\frac{\log n}{n}\qquad(n\ge N_0).        \tag{EUT2}
\]
\end{lemma}
\begin{proof}
On $u\ge1$ the first positive theta term and the complete
sum respectively give
\[
 18e^{9u/2-\pi e^{2u}}\le\Phi(u)
                         \le80e^{9u/2-\pi e^{2u}}.      \tag{EUT3}
\]
The upper bound uses the full summation bound in (EUL3).
The lower bound follows from $4\pi^2-6\pi e^{-2u}>24>18$,
using $3<\pi<4$ and $e^2>2$.

Let $s_n=(\log n)/2$ and
$v_n=(\log n-\log\log n)/2$.
Integration on $[v_n,v_n+1/n]$ gives
\[
 \mu_{2n}\ge v_n^{2n}
       \exp[-C_0n/\log n-C_0\log n]                     \tag{EUT4}
\]
for all large $n$. Fix $0<\eta<1$.
The contribution of $0\le u\le(1-\eta)s_n$, normalized by
$\mu_{2n}s_n^{2r}$, has logarithm at most
$2n\log[(1-\eta)s_n/v_n]+C_0n/\log n+O(\log n)$,
which tends to $-\infty$, for $r=0,1,2$.
On $u\ge u_0=(1+\eta)s_n$ use the coarser full bound (EUL3).
The logarithmic derivative of
$u^{2n+2r}e^{9u/2-3e^{2u}}$ is at most
$-3e^{2u_0}$ throughout this half-line for large $n$.
Its integral is bounded by its value at $u_0$ divided
by $3e^{2u_0}$. The normalized logarithm is at most
$2n\log[(1+\eta)s_n/v_n]-3n^{1+\eta}
+C_0n/\log n+O(\log n)$ and tends to $-\infty$.
On the remaining interval $u/s_n\in[1-\eta,1+\eta]$.
A limsup and liminf squeeze, first in $n$ and then in $\eta$,
proves the moment limits; the factorial ratio proves (EUT1).

For the quantitative curvature put
$\varphi_n(u)=2n\log u+9u/2-\pi e^{2u}$.
For $n\ge e^8$ its unique maximum $u_n$ satisfies
\[
 \tfrac14\log n<u_n<\log n,\quad u_n\ge2,\qquad
 2\pi e^{2u_n}=2n/u_n+9/2.
\]
Its derivative has the required signs at the two displayed bounds
and its second derivative is negative.
On $u\ge1$, $-\varphi_n''(u)\ge\lambda_n=2n/(\log n)^2$.
On $|u-u_n|\le1/(10\sqrt n)$ it is at most $106n$.
The lower density in (EUT3) therefore gives
$\mu_{2n}\ge2n^{-1/2}e^{\varphi_n(u_n)}$.
The upper density, the global Gaussian potential bound, and
$(u^2-u_n^2)^2\le8u_n^2(u-u_n)^2+2(u-u_n)^4$ give
\[
 \frac{\mathbb E[(U^2-u_n^2)^2]}{u_n^4}
                    \le C_1\frac{\log n}{n},            \tag{EUT5}
\]
where $U$ has density $u^{2n}\Phi(u)/\mu_{2n}$ on $\R$.
The two Gaussian integrals are
$\sqrt{2\pi}\lambda_n^{-3/2}$ and
$3\sqrt{2\pi}\lambda_n^{-5/2}$.
After division by the mass lower bound and $u_n^4$,
they cost $O(\log n/n)$ and $O(\log n/n^2)$.
The $|u|<1$ part is exponentially smaller, since
$\varphi_n(u_n)\ge n/2$.
Cauchy--Schwarz in (EUT5) gives
$\mathbb E U^2\ge u_n^2/2$ for all large $n$. Hence
\[
 1\le\frac{\mu_{2n+4}\mu_{2n}}{\mu_{2n+2}^2}
 =1+\frac{\operatorname{Var}(U^2)}{(\mathbb E U^2)^2}
 \le1+C_2\frac{\log n}{n}.
\]
The logarithm of the exact factorial quotient is $O(1/n)$.
This proves (EUT2).
\end{proof}

\begin{lemma}[All high roots from a limiting second quotient]
\label{lem:euler-second-quotient}
Let $a_j>0$ and $a_{j+2}a_j/a_{j+1}^2\to q\in(0,1)$.
Then $f(w)=\sum_{j\ge0}(-1)^ja_jw^j$ is entire.
With $W_n=a_n/a_{n+1}$ its high roots are positive simple,
$w_n/W_n\to1$, and $w_{n+1}/w_n\to1/q$.
Every sufficiently large modulus root belongs to this list.
\end{lemma}
\begin{proof}
Put $\rho_j=a_{j+1}/a_j$. Its ratio tends to $q<1$,
so the series is entire. On $\C\setminus\{0\}$ normalize it:
\[
 \Psi_n(t)=\frac{f(W_nt)}{(-1)^na_nW_n^nt^n}
 =\sum_{k=-n}^\infty(-1)^k A_{n,k}t^k,\quad
             A_{n,k}=\frac{a_{n+k}W_n^k}{a_n}.           \tag{EUT6}
\]
For fixed integer $k$, $A_{n,k}\to q^{k(k-1)/2}$.
Choose $q<Q<1$ and $N$ with $\rho_{j+1}/\rho_j\le Q$
for $j\ge N$. For $n\ge N,k\ge0$,
$A_{n,k}\le Q^{k(k-1)/2}$.
For $k=-\ell$ its logarithm is exactly
$\sum_{j=n-\ell}^{n-1}(\log\rho_n-\log\rho_j)$.
Indices $j\ge N$ obey the corresponding geometric bound.
The early indices add at most
$\sum_{j=0}^{N-1}\max\{D_j,0\}$, where
$D_j=\sum_{i=j}^{N-1}
[\log(\rho_{i+1}/\rho_i)-\log Q]$.
Thus $A_{n,-\ell}\le C_NQ^{\ell(\ell+1)/2}$
with a constant independent of $n,\ell$.
These summable bounds prove normal convergence on compact annuli to
\[
 \Theta_q(t)=\sum_{k\in\Z}(-1)^kq^{k(k-1)/2}t^k
 =\prod_{j\ge1}(1-q^j)(1-tq^{j-1})(1-q^j/t).             \tag{EUT7}
\]
This is \cite[eq.~20.5.9]{DLMF} with nome $\sqrt q$
and $p^2=-t/\sqrt q$. Its normal single product has exactly
the simple positive zeros $q^m$, $m\in\Z$.

Fix $1/2<\beta<1$. The annulus
$q^\beta<|t|<q^{-\beta}$ contains exactly the zero1.
Normal convergence keeps both boundary winding numbers unchanged.
The annular argument principle gives exactly one zero for $\Psi_n$.
Rouch\'e on a small disk around1 gives the same count.
Conjugation makes the single zero positive and simple.
Shrinking this fixed disk proves $w_n/W_n\to1$.
Since $W_{n+1}/W_n\to1/q$ and $\beta>1/2$, the physical
annuli $q^\beta W_n<|w|<q^{-\beta}W_n$ overlap for all
large $n$ and cover all large radii.
Small fixed relative disks around $W_n$ are disjoint for large $n$.
This proves exhaustion and unambiguous high-root labeling.
\end{proof}

\begin{theorem}[The complete Euler-Xi tail at every positive time]
\label{thm:euler-xi-tail}
For every fixed $\epsilon>0$ there exists a finite $R_\epsilon$
such that every zero of $\mathcal C_\epsilon\Xi$ outside this
disk is real and simple. Its positive high zeros can be labeled
by coefficient indices with
\[
 z_n\sim\frac{4n}{\log n}e^{4\epsilon n+2\epsilon},
 \qquad N_{\epsilon,+}(R)\sim\frac{\log R}{4\epsilon}.    \tag{EUT8}
\]
The nonreal zeros of this modified function form a finite set;
no reality claim is made for that set.
\end{theorem}
\begin{proof}
Its even coefficients are $a_n=b_ne^{-4\epsilon n^2}$.
By (EUT2) their second quotient tends to
$q=e^{-8\epsilon}\in(0,1)$.
Apply Lemma~\ref{lem:euler-second-quotient} to
$H_\epsilon(w)=\mathcal C_\epsilon\Xi(\sqrt w)$.
Its high $w$-zeros are positive simple, hence their two
$z$-preimages are real simple.
Their scales satisfy
$W_n\sim16n^2(\log n)^{-2}e^{4\epsilon(2n+1)}$.
Square roots give (EUT8) and its count.
A finite shift of the global numbering may differ from the
coefficient labels.
\end{proof}

\begin{theorem}[A uniform tail as Euler time tends to zero]
\label{thm:euler-xi-uniform-tail}
For every fixed $A>\pi^2/16$ there exist $\bar\epsilon>0$
and $N_A<\infty$ such that the following holds.
For all $0<\epsilon\le\bar\epsilon$ and all $n\ge N_A$
with $\epsilon\log n\ge A$, put
$W_n=(b_n/b_{n+1})e^{4\epsilon(2n+1)}$.
The actual $H_\epsilon$ has exactly one positive simple zero in
\[
                   |w/W_n-1|<\epsilon/\log n.           \tag{EUT9}
\]
All sufficiently large modulus zeros are these zeros.
For $n_\epsilon=\max\{N_A,\lceil e^{A/\epsilon}\rceil\}$,
the real simple $z$-tail may begin at a radius bounded by
\[
 R_\epsilon\le C_A n_\epsilon e^{4\epsilon n_\epsilon},
 \qquad\log R_\epsilon=O_A(\epsilon e^{A/\epsilon})
                      \quad(\epsilon\downarrow0).       \tag{EUT10}
\]
These analytic bounds include no numerical low-core certificate.
\end{theorem}
\begin{proof}
Set $h=8\epsilon,q=e^{-h}$ and normalize by (EUT6).
For $|k|\le n/2$, telescoping (EUT2) bounds the logarithm
of $b_{n+k}(b_n/b_{n+1})^k/b_n$ by $C(\log n/n)k^2$.
On $q^\beta\le|t|\le q^{-\beta}$, fixed $1/2<\beta<1$,
the remaining factors are bounded by
$e^{-4\epsilon k^2+O(\epsilon|k|)}$.
The condition makes $\log n/n=o(\epsilon)$ uniformly.
Using $|e^v-1|\le|v|e^{|v|}$ bounds the complete local
difference by $C(\log n/n)\epsilon^{-3/2}$.

For $|k|>n/2$, $\Phi\le Ke^{-u^2}$ gives
$b_j\le K\sqrt\pi/(4^jj!)\le K\sqrt\pi$ for every $j$.
(EUT4) and a factorial bound give $\log b_n\ge-Cn\log n$;
(EUT1) gives $|\log(b_n/b_{n+1})|\le C\log n$.
Thus the logarithm of the coefficient correction is at most
$C(n+|k|)\log n\le C'|k|\log n$.
The negative quadratic heat exponent absorbs it uniformly,
since $\epsilon\ge A/\log n$ and $n/(\log n)^2\to\infty$.
The complete far tails of both series are at most
$C\epsilon^{-1/2}e^{-c\epsilon n^2}$. Therefore
\[
 \|\Psi_{\epsilon,n}-\Theta_q\|
 \le C(\log n/n)\epsilon^{-3/2}
                  +C\epsilon^{-1/2}e^{-c\epsilon n^2}. \tag{EUT11}
\]

The genuine theta minimum must be retained.
Factorwise (EUT7) gives
$\min_{|t|=r}|\Theta_q(t)|\ge|\Theta_q(r)|$.
For $a=1/2+\log r/h$, square completion and
Gaussian Poisson summation give exactly
\[
 \Theta_q(r)=e^{ha^2/2}\sqrt{2\pi/h}
 \sum_{m\in\Z}e^{-2\pi^2(m+1/2)^2/h}
                         e^{2\pi i(m+1/2)a}.            \tag{EUT12}
\]
On $r=q^{\pm\beta}$ the leading coefficient is $2\cos(\pi a)$,
of modulus $2\sin(\pi\beta)>0$.
The remote terms cost relatively $O(e^{-4\pi^2/h})$.
Both boundary minima are consequently at least
$c_\beta\epsilon^{-1/2}e^{-\pi^2/(16\epsilon)}$.

On $|t-1|=\epsilon/\log n$, use the holomorphic identity
(EUT12) with $d=\log t/h$.
Its leading pair is $-2\sin(\pi d)$, of modulus at least
$c/\log n$. Each remote paired term also contains
$\sin((2m+1)\pi d)$, bounded for $|d|\le1/10$ by
$C(2m+1)|d|e^{Cm|d|}$.
The entire remainder therefore shares the leading vanishing
$|d|$ factor and is uniformly relatively $O(e^{-c/h})$.
The small-circle minimum is at least
$c\epsilon^{-1/2}e^{-\pi^2/(16\epsilon)}/\log n$;
an absolute fixed-time Poisson remainder would not suffice.

The first ratio of (EUT11) to the annulus minimum is at most
\[
 C\frac{\log n}{n\epsilon}e^{\pi^2/(16\epsilon)}
 \le C_A(\log n)^2n^{-1+\pi^2/(16A)}\longrightarrow0
\]
uniformly. Division by the small-circle minimum adds one $\log n$.
The strict $A>\pi^2/16$ makes that ratio vanish too.
The far-error exponent
$-c\epsilon n^2+\pi^2/(16\epsilon)$ tends uniformly to $-\infty$.
Rouch\'e and both winding numbers give the exact counts1,
and conjugation gives a positive simple root.
Finally (EUT2) gives
$\log(W_{n+1}/W_n)=8\epsilon+O(\log n/n)
=8\epsilon+o(\epsilon)$ throughout the regime.
Every consecutive annulus overlaps its predecessor and exhausts
all subsequent radii; the small disks are disjoint.
(EUT1) gives
$\sqrt{W_{n_\epsilon}}\le C_A n_\epsilon e^{4\epsilon n_\epsilon}$,
which proves (EUT10).
\end{proof}

The operation (EUL5) does not preserve the original Xi zeros.
The fixed positive time (EUL9) therefore proves a global theorem
about a modified function, with no assertion of RH or of a
realifying times sequence tending to zero.
Theorem~\ref{thm:euler-no-positive-measure} also precludes a direct
identification through finite positive Fourier measures.
An arithmetic transfer or a zero-time realification still requires
a separately proved mechanism.

\section{A theta summand does not have termwise Fourier positivity}
This result is about a summand of the exact density, not about
the sign of the complete Xi defect. Set
\[
 \phi_1(u)=2\pi e^{5|u|/2}(2\pi e^{2|u|}-3)
              e^{-\pi e^{2|u|}},
 \qquad
 f_p(q)=\phi_1((p+q)/2)\phi_1((p-q)/2),\quad p>0.
\]

\begin{proposition}\label{prop:cusp}
For every fixed $p>0$, the Fourier transform of $f_p$
has infinitely many positive and negative values. More precisely,
\[
 \widehat f_p(x)=-\frac{2J_p}{x^2}\cos(px)+o(x^{-2}),\qquad
 J_p=\phi_1(p)\pi e^{-\pi}(30\pi-15-8\pi^2)>0.
\]
\end{proposition}
\begin{proof}
Direct differentiation on the positive half-line gives
$\phi_1'(0+)=\pi e^{-\pi}(30\pi-15-8\pi^2)$.
The quadratic $30v-15-8v^2$ is decreasing for $v>15/8$.
The elementary bounds $3<\pi<22/7$ therefore give
\[
 30\pi-15-8\pi^2>
 30(22/7)-15-8(22/7)^2=13/49>0.
\]
Also $\phi_1(p)>0$. The function $f_p$ is continuous and
even, smooth except at $q=\pm p$, and its first derivative
has jump $J_p$ at each of those two points.
All its piecewise derivatives have integrable double
exponential tails, so in distributions
\[
 D^2f_p=g_p+J_p(\delta_p+\delta_{-p}),\qquad g_p\in L^1(\R).
\]
Fourier transformation gives
$-x^2\widehat f_p(x)=\widehat g_p(x)+2J_p\cos(px)$.
The Riemann--Lebesgue lemma yields the stated error.
Evaluation at $x_n=2\pi n/p$ and
$x'_n=(2n+1)\pi/p$ proves eventual negative and positive
values, respectively.
\end{proof}

Multiplication by $p^2\sinhc(yp)>0$ does not change these
signs, so even this single fixed-$p$ summand cannot support
a termwise positive-slice proof. The full theta series is
smoothly even and its derivatives cancel at the origin:
$\sum_{n\ge1}\phi_n'(0+)=\Phi'(0)=0$.
The cancellation is justified by uniform derivative
convergence near zero from the Gaussian series.
Thus the cusp argument does not give a sign obstruction
for the complete $\Phi$ or for $C_y$, and provides no
off-line zeta zero.

\begin{proposition}[Integrating the first theta slice does not restore its sign]
\label{prop:integrated-slice}
Set
\[
 C_y^{11}(x)=\iint (u+v)^2\sinhc(y(u+v))
                    \phi_1(u)\phi_1(v)\cos(x(u-v))\dd u\dd v.
\]
For every $0\le y\le1/2$ this is positive definite as a function
of $x$, but for each fixed $y$ it is strictly negative for
all sufficiently large $|x|$. With $a_1=\phi_1'(0+)>0$,
\[
 C_y^{11}(x)=-\frac{16a_1^2}{(x^2+y^2)^3}
                       +O_y(|x|^{-7})\quad(0<y\le1/2),
 \qquad
 C_0^{11}(x)=-16a_1^2x^{-6}+O(|x|^{-8}).
\]
\end{proposition}
\begin{proof}
The change of variables $p=u+v$, $q=u-v$ has Jacobian $1/2$.
Thus $C_y^{11}$ is the Fourier transform of the nonnegative
integrable density
\[
 H_y^{11}(q)=\tfrac12\int_\R p^2\sinhc(yp)
                  \phi_1((p+q)/2)\phi_1((p-q)/2)\dd p.
\]
For any finite $z_j,x_j$ its Gram sum is
$\int H_y^{11}(q)|\sum_jz_je^{ix_jq}|^2\dd q\ge0$,
which proves positive definiteness without claiming pointwise sign.

Put $f=\phi_1|_{[0,\infty)}$ and
$F_1(\delta)=\int_\R\phi_1(u)e^{\delta u}\dd u$.
All half-line derivatives of $f$ have integrable
double exponential tails. Define finite weighted norms
$\|v\|_{1/2}=\int_0^\infty e^{u/2}|v(u)|\dd u$ and
\[
 \begin{split}
 K_0&=2\bigl(|f'''(0)|+\|f^{(4)}\|_{1/2}\bigr),\\
 K_1&=8|f'''(0)|+2\|(uf)^{(5)}\|_{1/2},\\
 K_2&=40|f'''(0)|+2\|(u^2f)^{(6)}\|_{1/2}.
 \end{split}
\]
For $I_v(\delta)=\int_0^\infty v(u)e^{\delta u}\dd u$,
integration by parts $N$ times gives the exact finite formula
\[
 I_v(\delta)=\sum_{k=0}^{N-1}
      (-1)^{k+1}\frac{v^{(k)}(0)}{\delta^{k+1}}
             +(-1)^N\delta^{-N}I_{v^{(N)}}(\delta).
\]
Apply it separately to $\delta$ and $-\delta$ for $v=f,uf,u^2f$,
with $N=4,5,6$, respectively. The cancellations of boundary
terms then prove, uniformly for $|\Re\delta|\le1/2$, $|\delta|\ge1$,
\[
 \left|F_1(\delta)-2a_1\delta^{-2}\right|\le K_0|\delta|^{-4},
 \quad
 \left|F_1'(\delta)+4a_1\delta^{-3}\right|\le K_1|\delta|^{-5},
 \quad
 \left|F_1''(\delta)-12a_1\delta^{-4}\right|\le K_2|\delta|^{-6}.
\]
These are separate derivative proofs, not differentiation of an
unspecified asymptotic error.

Fubini, evenness and the defect identity give, for $y>0$,
\[
 C_y^{11}(x)=-\frac1y\partial_s
   [F_1(s-1/2+ix)F_1(s-1/2-ix)]_{s=1/2-y}.
\]
Indeed the derivative is the integral of
$p e^{-yp}\phi_1(u)\phi_1(v)e^{ixq}$;
even reflection of $p$ makes it
$-\int p\sinh(yp)\phi_1(u)\phi_1(v)\cos(xq)$.
Dividing by $-y$ gives exactly the stated $p^2\sinhc(yp)$.
At $y=0$ the continuous extension is the second $s$-derivative
of the same product at $s=1/2$.

Let $\delta_\pm=-y\pm ix$, $r=(x^2+y^2)^{1/2}$.
Since $\delta_++\delta_-=-2y$ and $\delta_+\delta_-=r^2$,
the two leading product-derivative terms equal $16ya_1^2r^{-6}$.
For $|x|\ge1$ the remaining terms give the explicit bound
\[
 \left|C_y^{11}(x)+16a_1^2r^{-6}\right|
 \le \frac{K_P}{y}r^{-7},\qquad
 K_P=8a_1K_0+4a_1K_1+2K_0K_1.
\]
Thus $r\ge K_P/(8a_1^2y)$ is a sufficient negative-sign threshold
for each fixed $y>0$, together with $|x|\ge1$.
At $y=0$, applying the independent $F_1''$ bound directly to
the second product derivative gives
\[
 |C_0^{11}(x)+16a_1^2x^{-6}|\le K_3|x|^{-8},
 \quad
 K_3=24a_1K_0+4a_1K_2+16a_1K_1+2K_0K_2+2K_1^2.
\]
For example $|x|\ge\max\{1,\sqrt{K_3/(8a_1^2)}\}$ is sufficient.
This proves both negative tails, with finite analytic error bounds.
\end{proof}

The proposition concerns only the $(n,m)=(1,1)$ term.
The complete theta density has the cusp cancellations already
described, so this result cannot be extended termwise to the complete
Xi defect. It excludes that proposed sign mechanism, and gives
no off-line zeta zero.

\begin{proposition}[Every fixed finite theta head has a negative tail]
\label{prop:finite-theta-head}
For $N\ge1$ let $\Phi_N=\sum_{n=1}^N\phi_n$, where
\[
 \phi_n(u)=(4\pi^2n^4e^{9|u|/2}-6\pi n^2e^{5|u|/2})
                                      e^{-\pi n^2e^{2|u|}},
\]
and let $C_{N,y}$ be the Xi-defect integral with both factors
$\Phi$ replaced by $\Phi_N$, retaining every cross term.
Then
\[
 a_N:=\Phi_N'(0+)>0,\qquad
 a_N\sim8\pi^3(N+1)^6e^{-\pi(N+1)^2}\quad(N\to\infty).
                                                               \tag{TH}
\]
For each fixed $N$ and $0<y\le1/2$, writing
$r=(x^2+y^2)^{1/2}$, there are explicit finite constants
$K_{P,N},K_{3,N}$ such that, for $|x|\ge1$,
\[
 \left|C_{N,y}(x)+16a_N^2r^{-6}\right|
                       \le K_{P,N}y^{-1}r^{-7},
 \qquad
 \left|C_{N,0}(x)+16a_N^2x^{-6}\right|
                       \le K_{3,N}|x|^{-8}.             \tag{TN}
\]
In particular $C_{N,y}(x)$ is strictly negative for all
sufficiently large $|x|$, for every fixed $N,y$.
For the full defect and every such fixed $N,y$,
\[
 \lim_{|x|\to\infty}
       \frac{|C_y(x)-C_{N,y}(x)|}{|C_{N,y}(x)|}=1.       \tag{TR}
\]
Thus no fixed finite head admits a global relative-error
bound with a factor strictly less than one.
\end{proposition}
\begin{proof}
Direct differentiation gives
\[
 \phi_n'(0+)=\pi n^2e^{-\pi n^2}
              (30\pi n^2-15-8\pi^2n^4).
\]
The polynomial $30t-15-8t^2$ is decreasing for $t>15/8$
and negative for $t\ge12$. Since $\pi n^2>12$ for $n\ge2$,
all these derivatives are negative when $n\ge2$.
The Gaussian series and its derivatives converge uniformly
near zero. The full smooth even kernel has derivative zero
there, so
\[
 a_N=-\sum_{n>N}\phi_n'(0+)>0.
\]
For large $n$ the summand on the right is
$8\pi^3n^6e^{-\pi n^2}(1+O(n^{-2}))$.
The ratio of successive terms is bounded by
$C e^{-\pi(2n+1)}$ for all sufficiently large $n$.
Their tail is therefore asymptotic to its first term,
proving (TH).

For the quantitative Fourier calculation put
$f_N=\Phi_N|_{[0,\infty)}$ and
$F_N(\delta)=\int_\R\Phi_N(u)e^{\delta u}\dd u$.
All half-line derivatives have finite weighted $L^1$ norms.
Repeat the exact finite integration-by-parts formula in
Proposition~\ref{prop:integrated-slice} with $f_N$ in place
of $f$. Its constants $K_0,K_1,K_2$, defined there from
$f,uf,u^2f$ and their derivatives, give $K_{0,N},K_{1,N},K_{2,N}$.
The resulting independent derivative estimates are
\[
 \begin{split}
 |F_N(\delta)-2a_N\delta^{-2}|&\le K_{0,N}|\delta|^{-4},\\
 |F_N'(\delta)+4a_N\delta^{-3}|&\le K_{1,N}|\delta|^{-5},\\
 |F_N''(\delta)-12a_N\delta^{-4}|&\le K_{2,N}|\delta|^{-6},
 \end{split}
\]
for $|\Re\delta|\le1/2$, $|\delta|\ge1$.
The proof uses only even reflection and weighted
half-line derivative integrability, both valid for $\Phi_N$;
it is not a differentiation of an unspecified remainder.
For $y>0$, the exact product identity is
\[
 C_{N,y}(x)=-\frac1y\partial_s
    [F_N(s-1/2+ix)F_N(s-1/2-ix)]_{s=1/2-y}.
\]
At $y=0$ use its second $s$ derivative at $s=1/2$.
The same multiplication of the three derivative estimates
proves (TN) with
\[
 \begin{split}
 K_{P,N}&=8a_NK_{0,N}+4a_NK_{1,N}+2K_{0,N}K_{1,N},\\
 K_{3,N}&=24a_NK_{0,N}+4a_NK_{2,N}+16a_NK_{1,N}
                    +2K_{0,N}K_{2,N}+2K_{1,N}^2.
 \end{split}
\]
For fixed $y>0$, $r\ge K_{P,N}/(8a_N^2y)$ and $|x|\ge1$
suffice for a negative sign. At zero use
$|x|\ge\max\{1,\sqrt{K_{3,N}/(8a_N^2)}\}$.

The complete $\Phi$ has none of these cusp derivatives.
Its defect density is
\[
 H_y(q)=\frac12\int_\R p^2\sinhc(yp)
       \Phi((p+q)/2)\Phi((p-q)/2)\dd p.
\]
For every derivative order its $L^1$ norm is bounded
uniformly for $0\le y\le1/2$. Indeed differentiate the two
smooth factors in $q$ and use
$p^2\sinhc(yp)\le p^2e^{|p|/2}$ together with their
double-exponential derivative bounds. After the linear
change to $u,v$, every term is bounded by a product of
finite weighted derivative moments. This also justifies
the differentiations and integrations by parts.
Consequently $C_y(x)=O_m(|x|^{-2m})$ for every $m$,
uniformly in that closed $y$ interval.
For each fixed $N,y$, (TN) has the nonzero leading
coefficient $-16a_N^2x^{-6}$. Hence
$C_y(x)/C_{N,y}(x)\to0$, proving (TR).
\end{proof}

This obstruction concerns fixed finite theta heads. It
does not give the sign of the full defect. Letting
$N$ vary with $x$ requires a new uniform comparison
with its actual positive margin; the fixed-head argument
does not rule out that different mechanism.

\begin{corollary}[Uniform absolute theta-tail transfer]
\label{cor:theta-tail-transfer}
Set $R_N=\Phi-\Phi_N\ge0$ and, for $k=0,2$, define
\[
 M_k=\int_\R |u|^k e^{|u|/2}\Phi(u)\dd u,\qquad
 R_{k,N}=\int_\R |u|^k e^{|u|/2}R_N(u)\dd u.
\]
For all $N\ge1$, all real $x$, and $0\le y\le1/2$,
\[
 |C_y(x)-C_{N,y}(x)|
       \le4(R_{2,N}M_0+R_{0,N}M_2),\qquad
 R_{k,N}\le C_k(N+1)^4e^{-\pi(N+1)^2/2},              \tag{AT}
\]
where the explicit finite constants
\[
 C_k=\frac{8\pi^2}{2^k}\int_1^\infty
             (\log t)^k t^{3/2}e^{-\pi t/2}\dd t
\]
are independent of $N,x,y$.
This is an absolute bound, compatible with the relative
obstruction (TR).
\end{corollary}
\begin{proof}
The positive product difference is at most
$R_N(u)\Phi(v)+\Phi(u)R_N(v)$.
Use
\[
 (u+v)^2\sinhc(y(u+v))
       \le2(u^2+v^2)e^{|u|/2}e^{|v|/2}
\]
and $|\cos(x(u-v))|\le1$ to obtain the first estimate.
The moments are finite by the full-kernel bound.
For the tail, the upper bound
$\phi_n(u)\le4\pi^2n^4e^{9u/2}e^{-\pi n^2e^{2u}}$
on $u\ge0$ and $t=e^{2u}$ give
\[
 R_{k,N}\le\frac{4\pi^2}{2^k}\sum_{n>N}n^4
       \int_1^\infty(\log t)^k t^{3/2}e^{-\pi n^2t}\dd t.
\]
For $n,t\ge1$, $n^2t\ge(n^2+t)/2$. Extract
$e^{-\pi n^2/2}$ and use the displayed integral defining
$C_k/2$. For $n\ge2$ the ratio of consecutive terms of
$n^4e^{-\pi n^2/2}$ is less than $1/2$.
Its tail is bounded by twice its first term, proving (AT).
\end{proof}

\section{Gaussian smoothing changes the finite-head tail mechanism}
We keep the Laplace transform $F_N(\delta)$ used above separate from
the Fourier transform
\[
 \mathcal F_N(z)=\widehat\Phi_N(z)=F_N(-iz).
\]
Evenness makes the sign of the Fourier exponent immaterial here.
For any real even kernel $D$ with the required exponential moments,
write $C_{D,y}$ for its defect integral. Direct differentiation and
even reflection of one integration variable give
\[
 C_{D,y}(x)=\frac1y\partial_y|\widehat D(x+iy)|^2\quad(y>0),
 \qquad C_{D,0}(x)=\partial_y^2|\widehat D(x+iy)|^2\big|_{y=0}.
                                                               \tag{DI}
\]
The double integrals justify these derivatives by their weighted moments.

\begin{theorem}[Exact Gaussian correction and a uniform positive tail]
\label{thm:gaussian-head}
For every $N\ge1$ and $\tau>0$, put
\[
 g_\tau(u)=(4\pi\tau)^{-1/2}e^{-u^2/(4\tau)},\qquad
 D_{N,\tau}=\Phi_N*g_\tau.
\]
This is a smooth, positive, even kernel, and for every complex $z$
\[
 \widehat D_{N,\tau}(z)=e^{-\tau z^2}\mathcal F_N(z).       \tag{GS}
\]
Its zeros and their multiplicities are exactly those of $\mathcal F_N$.
For all real $x$ and $0\le y\le1/2$,
\[
 C_{D_{N,\tau},y}(x)=e^{-2\tau(x^2-y^2)}
       \left[C_{N,y}(x)+4\tau|\mathcal F_N(x+iy)|^2\right]. \tag{GD}
\]
For fixed $N$, uniformly over that closed $y$ interval,
\[
 C_{N,y}(x)=-16a_N^2x^{-6}+O_N(|x|^{-8}),\qquad
 |\mathcal F_N(x+iy)|^2=4a_N^2x^{-4}+O_N(|x|^{-6}).       \tag{UT}
\]
Consequently for every fixed $N,\tau$ the smoothed defect is
strictly positive for all sufficiently large $|x|$, uniformly
in $0\le y\le1/2$. Every $\mathcal F_N$ nevertheless has
at least one nonreal zero somewhere in $\C$.
No location within the critical strip is asserted for that zero.
\end{theorem}
\begin{proof}
Fubini and the Gaussian integral give (GS), with absolute convergence
on every complex Fourier line. Derivatives of the Gaussian are
integrable against the bounded, rapidly decreasing $\Phi_N$, proving
smoothness of the convolution. Positivity and evenness follow directly.
Since the Gaussian transform never vanishes, the zero assertion
includes every multiplicity. Taking squared moduli in (GS) gives
$e^{-2\tau(x^2-y^2)}|\mathcal F_N(x+iy)|^2$.
Differentiate using (DI); the quotient at zero is obtained by evenness
in $y$. This proves (GD), including $y=0$.

The three independent derivative bounds in the finite-head proof,
after $\delta=-iz$, read
\[
 \mathcal F_N(z)=-2a_Nz^{-2}+O_N(|z|^{-4}),\quad
 \mathcal F_N'(z)=4a_Nz^{-3}+O_N(|z|^{-5}),\quad
 \mathcal F_N''(z)=-12a_Nz^{-4}+O_N(|z|^{-6}),
\]
uniformly for $|\Im z|\le1/2$, $|\Re z|\ge1$.
These are separately proved derivative estimates, not derivatives
of an unspecified error term. Set
$H_x(y)=\mathcal F_N(x+iy)\mathcal F_N(x-iy)$.
It is even in $y$, so
\[
 C_{N,y}(x)=\int_0^1H_x''(ty)\dd t
\]
also at zero. Its second derivative is
\[
 -\mathcal F_N''(x+iy)\mathcal F_N(x-iy)
 +2\mathcal F_N'(x+iy)\mathcal F_N'(x-iy)
 -\mathcal F_N(x+iy)\mathcal F_N''(x-iy).
\]
Products containing an error are $O_N(|x|^{-8})$ uniformly in $y$.
The leading product is $4a_N^2/(x^2+y^2)^2$; its first
$y$ derivative divided by $y$ is $-16a_N^2/(x^2+y^2)^3$.
Expanding this expression uniformly for $|y|\le1/2$ proves the
first part of (UT). The zeroth derivative bound proves the second.
In (GD) the bracket therefore equals
\[
 16\tau a_N^2x^{-4}+O_{N,\tau}(|x|^{-6}).
\]
Its leading coefficient is positive, proving the uniform tail assertion.
The threshold is not claimed uniform as $N\to\infty$ or $\tau\downarrow0$.

Finally $\mathcal F_N$ is real even entire and $\mathcal F_N(0)>0$.
The elementary bound $\Phi_N(u)\le C_Ne^{9|u|/2}e^{-\pi e^{2|u|}}$
implies, with $b=|\Im z|$,
\[
 |\mathcal F_N(z)|\le \exp(C_N(1+b)\log(2+b)).
\]
To see this, maximize $(b+9/2)u-(\pi/2)e^{2u}$ on $u\ge0$
and integrate the remaining $e^{-(\pi/2)e^{2u}}$.
Thus its order is at most one. If all its zeros were real,
Hadamard factorization \cite[Theorem 7.7, printed p.~58]{Petrow}
and evenness would give the normally convergent paired product
\[
 \mathcal F_N(z)=\mathcal F_N(0)
                     \prod_{a>0}(1-z^2/a^2)^{m_a}.
\]
Indeed the linear exponential coefficient is
$\mathcal F_N'(0)/\mathcal F_N(0)=0$, and pairing cancels
the genus-one exponential factors. The real-rooted finite products
satisfy $P'^2-PP''\ge0$ on $\R$: off their zeros, divide by $P^2$
and differentiate $P'/P$ to obtain $\sum_a m_a/(x-a)^2$;
at zeros use continuity. Normal convergence implies convergence
of the first two derivatives, preserving this inequality.
But (DI) gives $C_{N,0}=2(\mathcal F_N'^2-\mathcal F_N\mathcal F_N'')$,
which is negative in (UT) for large $|x|$. This contradiction
proves existence of a nonreal finite-head zero.
\end{proof}

\begin{proposition}[Uniform absolute convergence of smoothed heads]
\label{prop:gaussian-transfer}
Define $w_0(u)=e^{|u|/2}$, $w_2(u)=(1+u^2)e^{|u|/2}$ and
$\|h\|_{w_k}=\int w_k|h|$. Put $\mathsf M_k=\|\Phi\|_{w_k}$,
$E=\Phi-D_{N,\tau}$, and
\[
 \mu_j(\tau)=\int_\R |u|^j e^{|u|/2}g_\tau(u)\dd u.
\]
The following explicit bounds hold:
\[
 \begin{split}
 \|E\|_{w_0}&\le R_{0,N}+\mu_1(\tau)\|\Phi_N'\|_{w_0},\\
 \|E\|_{w_2}&\le R_{0,N}+R_{2,N}
       +2(\mu_1(\tau)+\mu_3(\tau))\|\Phi_N'\|_{w_2},\\
 \mu_j(\tau)&\le e^{\tau/2}
   \frac{8^{(j+1)/2}\Gamma((j+1)/2)}{2\sqrt\pi}\tau^{j/2},\\
 \sup_{x\in\R,\,0\le y\le1/2}|C_y(x)-C_{D_{N,\tau},y}(x)|
 &\le4(\|E\|_{w_2}\mathsf M_0+\|E\|_{w_0}\mathsf M_2)
                       +4\|E\|_{w_2}\|E\|_{w_0}.
 \end{split}                                               \tag{GT}
\]
The constants $R_{k,N}$ are bounded explicitly in (AT).
For any sequence $N\to\infty$, $\tau_N\downarrow0$, the last
right side tends to zero. This is absolute convergence, not a
relative estimate against the vanishing defect tail.
\end{proposition}
\begin{proof}
The finite head is absolutely continuous, including its cusp.
For any absolutely continuous $f$ with the stated weighted
derivative integrability,
$f(u-t)-f(u)=-\int_0^t f'(u-r)\dd r$.
The weight inequalities
\[
 w_0(v+r)\le e^{|r|/2}w_0(v),\qquad
 w_2(v+r)\le2e^{|r|/2}(1+r^2)w_2(v)
\]
give weighted translation errors at most
$|t|e^{|t|/2}\|f'\|_{w_0}$ and
$2|t|e^{|t|/2}(1+t^2)\|f'\|_{w_2}$, respectively.
Integrate against $g_\tau$ and add the tail $\Phi-\Phi_N$.
This proves the first two inequalities. The third follows from
$|u|/2\le u^2/(8\tau)+\tau/2$ and direct Gaussian integration.

Expand $\Phi(u)\Phi(v)-D(u)D(v)$ as
$E(u)\Phi(v)+\Phi(u)E(v)-E(u)E(v)$.
The bound $(u+v)^2\sinhc(y(u+v))\le
2(u^2+v^2)e^{|u|/2}e^{|v|/2}$ proves the last inequality,
uniformly also at $y=0$. The weights $w_2$ dominate the
required second moments.
Lastly $\sup_N\|\Phi_N'\|_{w_k}<\infty$ for $k=0,2$.
For $u\ge0$ each first derivative is $e^{u/2}$ times a
fixed polynomial in $\pi n^2e^{2u}$, times $e^{-\pi n^2e^{2u}}$.
Absorb the polynomial in half the exponential and sum the
remaining Gaussian series, as in the proof of (ID).
Its weighted integral is finite; even reflection handles $u<0$.
The tail bounds (AT) and $\mu_1=O(\sqrt\tau)$,
$\mu_3=O(\tau^{3/2})$ now prove the claimed diagonal convergence.
\end{proof}

\begin{corollary}[Exact obstruction to a finite Gaussian boost]
\label{cor:gaussian-threshold}
For a fixed $N$, a finite $\tau>0$ makes
$C_{D_{N,\tau},y}(x)\ge0$ for every real $x$, $0\le y\le1/2$
if and only if $\mathcal F_N$ has no zero with $0<\Im z\le1/2$.
More precisely, off its zeros let
\[
 \tau_N^*=\sup_{\substack{x\in\R,\,0\le y\le1/2\\
                         \mathcal F_N(x+iy)\ne0}}
       \frac{[-C_{N,y}(x)]_+}{4|\mathcal F_N(x+iy)|^2}.    \tag{GB}
\]
Under that zero exclusion $\tau_N^*<\infty$ and every
$\tau>\tau_N^*$ gives strict positivity for $y>0$ and
nonnegativity for $y=0$. If a zero $x_0+i\delta$ of order $m$
has $0<\delta\le1/2$, then along $x=x_0$, $y\uparrow\delta$,
\[
 \frac{-C_{N,y}(x_0)}{4|\mathcal F_N(x_0+iy)|^2}
                    \sim\frac{m}{2\delta(\delta-y)}.      \tag{GO}
\]
No finite boost then works on the closed strip.
\end{corollary}
\begin{proof}
Near that nonreal zero, write
$\mathcal F_N(z)=(z-x_0-i\delta)^mG(z)$ with $G$ nonzero.
The logarithmic derivative in (DI) gives
\[
 \frac{C_{N,y}(x_0)}{|\mathcal F_N(x_0+iy)|^2}
          =\frac{2m}{y(y-\delta)}+O(1)
\]
as $y\uparrow\delta$, proving (GO). A zero itself has
$C_{N,\delta}(x_0)=0$; its nearby negative defect,
rather than its value at the zero, is the obstruction.

For sufficiency assume the stated nonreal-zero exclusion.
Near a real zero $a$ of order $m$, divide by $(z-a)^m$.
Reality gives $|G(x+iy)|^2$ even in $y$ and bounded away from
zero in a sufficiently small neighborhood. Thus
\[
 \frac{C_{N,y}(x)}{|\mathcal F_N(x+iy)|^2}
           =\frac{2m}{(x-a)^2+y^2}+O(1).
\]
The error has a bounded continuous extension to $y=0$.
The leading term dominates, so no negative ratio arises close
to a real zero. At $(a,0)$ the defect is $2|\mathcal F_N'(a)|^2$
for a simple zero and zero for a multiple zero; both are allowed.
By (UT), the ratio in (GB) is $x^{-2}+O_N(|x|^{-4})$ for large
$|x|$, uniformly in $y$. On the remaining compact region there
are only finitely many real zeros; after removing their preceding
neighborhoods, continuity gives a finite supremum.
This proves finiteness of (GB). Formula (GD) proves the positivity
assertion for $\tau>\tau_N^*$, including every real multiplicity.
\end{proof}

\begin{corollary}[A positive head sequence needs no uniform Gaussian scale]
\label{cor:head-exhaustion}
Suppose there are integers $N_j\to\infty$ and numbers $\tau_j>0$
such that $C_{D_{N_j,\tau_j},y}(x)\ge0$ for every real $x$
and $0\le y\le1/2$. Then RH holds.
No upper bound on $\tau_j$, and no limit $\tau_j\to0$, is required
for this implication. Its positivity hypothesis is not established.
\end{corollary}
\begin{proof}
Corollary~\ref{cor:gaussian-threshold} excludes every zero of
$\mathcal F_{N_j}$ in $0<\Im z\le1/2$.
Independently of $\tau_j$, the unsmoothed transforms satisfy
\[
 \sup_{|\Im z|\le1/2}|\Xi(z)-\mathcal F_{N_j}(z)|
              \le R_{0,N_j}\longrightarrow0
\]
by the integral formula and (AT).
If $\Xi$ had an off-line zero, conjugation would give one
$z_0$ with $0<\Im z_0<1/2$.
Choose a closed disk about $z_0$ strictly inside that upper strip,
with $\Xi$ nonzero on its boundary. Its boundary minimum modulus
is positive. For sufficiently large $j$ the displayed error is
smaller than this minimum. The straight homotopy from $\Xi$ to
$\mathcal F_{N_j}$ is then nonzero on the boundary, so their
argument-principle zero counts in the disk agree, including
multiplicities. The first count is positive and the second is
zero, a contradiction.
\end{proof}

A finite-head nonreal zero can lie outside this closed strip.
No unbounded sequence of heads satisfying the preceding positivity
hypothesis has been established. This is the exact missing step of
the zero-preserving route; uniformity of the Gaussian scale is unnecessary.
If a stronger sequence with $\tau_j\to0$ were found, (GT) would also
transfer nonnegativity to the full defect, without a uniform strict margin.
Indeed a hypothetical off-line zero would make
$y\mapsto|\Xi(x_0+iy)|^2$ nondecreasing and zero at $y=\delta$;
it would vanish on $[0,\delta]$, contradicting analyticity.
The proved positive tails alone provide neither positivity sequence.

\subsection{Two independently reconstructed head certificates}

\begin{lemma}[Exact head evaluation and derivative majorants]
\label{lem:head-gamma}
Writing $a_n=\pi n^2$, one has for every complex $z$
\[
 \mathcal F_N(z)=\sum_{n=1}^N\sum_{\varepsilon=\pm1}
 a_n^{-1/4-\varepsilon iz/2}
 \left[2\Gamma(9/4+\varepsilon iz/2,a_n)
       -3\Gamma(5/4+\varepsilon iz/2,a_n)\right].        \tag{HG}
\]
For every integer $j\ge0$,
\[
 \sup_{|\Im z|\le1/2}|\mathcal F_N^{(j)}(z)|
 \le M_{N,j}:=
 8\pi^2j!\sum_{n=1}^N
             \frac{n^4e^{-\pi n^2}}{(2\pi n^2-5)^{j+1}}. \tag{HMJ}
\]
\end{lemma}
\begin{proof}
On $u\ge0$, put $t=a_ne^{2u}$ in
$\phi_n(u)=e^{u/2-t}(4t^2-6t)$.
The two exponential terms from $2\cos(zu)$ give
$a_n^{-1/4-\varepsilon iz/2}t^{1/4+\varepsilon iz/2}(2t-3)e^{-t}\dd t$.
Their integrals prove (HG). The powers use the real logarithm of $a_n>0$;
the upper incomplete gamma integral is entire in its first parameter.
Further, $\phi_n(u)>0$ and
\[
 \phi_n(u)e^{u/2}
 \le4\pi^2n^4e^{5u}e^{-\pi n^2e^{2u}}
 \le4\pi^2n^4e^{-\pi n^2}e^{-(2\pi n^2-5)u}.
\]
Since $2\pi n^2-5>0$, integration of $u^j$ and even reflection prove
(HMJ), justifying every differentiated integral.
\end{proof}

\begin{theorem}[Global strip reality and positivity for the first head]
\label{thm:first-head-certificate}
The first head $\mathcal F_1$ has exactly two zeros in
$|\Im z|\le1/2$, counted with multiplicity. They are real, simple,
and opposite. Consequently $\tau_1^*<\infty$, and for every
$\tau>\tau_1^*$ the Gaussian defect is strictly positive for
all $x\in\R$ and $0\le y\le1/2$.
\end{theorem}
\begin{proof}
We specify the analytic enclosures and the finite certificate used.
For $f=\phi_1|_{[0,\infty)}$, $T=\pi e^{2u}$, write
$f^{(j)}(u)=e^{u/2-T}P_j(T)$, where
\[
 P_0=4T^2-6T,\qquad
 P_{j+1}=(1/2-2T)P_j+2TP_j'.
\]
In particular $P_1=-15T+30T^2-8T^3$ and
\[
 P_4=-\tfrac{1875}{8}T+\tfrac{15465}{4}T^2-8512T^3
                                      +5176T^4-1056T^5+64T^6.
\]
The finite-head integration-by-parts bound gives
\[
 |\mathcal F_1(z)+2a_1/z^2|\le K_0|z|^{-4},\quad
 K_0=2\bigl(|f'''(0+)|+\|f^{(4)}\|_{1/2}\bigr),
 \quad |\Im z|\le1/2,\ |z|\ge1 .
\]
The weighted norm is exactly
\[
 \|f^{(4)}\|_{1/2}
 =\frac1{2\sqrt\pi}\int_\pi^\infty
                  T^{-1/2}|P_4(T)|e^{-T}\dd T.
\]
Replacing the lower endpoint by $3$ gives an upper bound.
On $[3,100]$ the certificate uses $100000$ rational cells,
interval evaluation of $P_4$, and the decreasing factors
$T^{-1/2}$ and $e^{-T}$ at each cell's left endpoint.
The remaining integral is bounded by
$\sum_j|[T^j]P_4|\Gamma(j+1/2,100)$.
The outward-rounded 80-decimal-digit computation proves
\[
 \begin{gathered}
 a_1>0.0394987652677,\qquad |f'''(0+)|<10.538986,\\
 \|f^{(4)}\|_{1/2}<331.276011,\qquad K_0<683.630,\\
 K_0/(2a_1)<8653.815<94^2.
 \end{gathered}
\]
Thus $\mathcal F_1(z)\ne0$ on $|\Re z|\ge94$, $|\Im z|\le1/2$.

For the remaining rectangle use its counterclockwise boundary
$[-94,94]\times[-1/2,1/2]$ and (HG) at exact dyadic endpoints.
Derivatives at segment midpoints are integrated on $[0,2]$ using
Acb; all callbacks are entire. For $c_0=1$, $c_j=(2j)^j$ if $j\ge1$,
the omitted derivative tail is at most
\[
 2c_j\pi^{-3/4}
       [2\Gamma(11/4,\pi e^4)+3\Gamma(7/4,\pi e^4)] .
\]
Here $u^j\le c_je^{u/2}$ and
$\Gamma(s,T)\le5T^{s-1}e^{-T}$ for $1\le s\le3$, $T\ge1$:
substitute $t=T+v$ and bound $(1+v/T)^{s-1}\le(1+v)^2$.
These also bound the global fourth-derivative moment.
On a segment of half-length $r$ centered at $z_c$, enclose the
degree-three Taylor polynomial and enlarge its image by a complex disk
of radius $M_4r^4/4!$, where $M_4$ is that computed moment bound.
The integral Taylor remainder along the segment proves the enclosure.
Midpoints, radii, powers, endpoint inclusion and cyclic adjacency
are all calculated exactly or with outward-rounded intervals.

The 140-bit certificate accepts all $4128$ dyadic segments.
Every full segment image lies in an open half-plane
$\Re(\overline m\,w)>0$, for its certified image midpoint $m$.
It therefore avoids zero and has a continuous argument branch
of width less than $\pi$. The principal argument of the independently
enclosed endpoint ratio is the true argument increment.
Their summed winding interval is
\[
 2+[-4.79\cdot10^{-34},\,4.79\cdot10^{-34}].
\]
It contains just the integer $2$, so the argument principle gives
exactly two zeros, with multiplicity. The certificate also proves
$\mathcal F_1(0)>0$ and $\mathcal F_1(94)<0$.
Continuity and evenness give at least one positive and one negative
real zero. The count forces precisely these zeros, both simple;
boundary and tail exclusion extend the result to the closed strip.

The contour and tail calculations use outward-rounded Arb/Acb
arithmetic through python-flint 0.9.0 \cite{FlintDocs}, with the
exact geometry, subdivision and enclosure bounds specified above.
Corollary~\ref{cor:gaussian-threshold} gives $\tau_1^*<\infty$.
Away from a real zero, the extra
$4(\tau-\tau_1^*)|\mathcal F_1|^2$ makes the defect strict.
At either real zero its value is $2|\mathcal F_1'|^2>0$.
This proves strictness also at $y=0$, without imposing simplicity on
any zero of $\Xi$.
\end{proof}

\begin{theorem}[A certified third-head obstruction inside the strip]
\label{thm:third-head-certificate}
The third head $\mathcal F_3$ has exactly one zero, with multiplicity,
in the disk $|z-c|<2^{-85}$, where
\[
 c=\frac{
 90228248448811850205073154815756697372+
 i\,634498798097303045007033058879710098}{2^{120}}.
\]
The disk lies in $0<\Im z<1/2$.
This zero is nonreal and simple. In particular $\tau_3^*=\infty$:
no finite Gaussian boost makes the third-head defect nonnegative
on the whole closed strip.
\end{theorem}
\begin{proof}
All geometry is dyadic. Use (HG), $h=2^{-120}$ and
$r=2^{-85}$, and set
$D_h=(\mathcal F_3(c+h)-\mathcal F_3(c-h))/(2h)$.
The global bounds (HMJ) prove
\[
 |\mathcal F_3'(c)-D_h|\le M_{3,3}h^2/6,\qquad
 |\mathcal F_3(z)-\mathcal F_3(c)-\mathcal F_3'(c)(z-c)|
          \le M_{3,2}r^2/2 \quad(|z-c|\le r).
\]
Both follow from integral Taylor remainders on straight segments
inside the strip; the first uses the symmetric pair of segments.
The independently reproduced 640-bit computation gives
\[
 \begin{split}
 M_{3,2}&<3.230,\qquad M_{3,3}<7.552,\\
 |\mathcal F_3(c)|&<1.36\cdot10^{-56},\\
 |D_h|-M_{3,3}h^2/6&>1.7775\cdot10^{-20},\\
 |\mathcal F_3(c)|+M_{3,2}r^2/2&<1.080\cdot10^{-51},\\
 (|D_h|-M_{3,3}h^2/6)r&>4.594\cdot10^{-46}.
 \end{split}
\]
The strict Rouch\'e inequality has ratio below $2.35\cdot10^{-6}$.
Comparison with $\mathcal F_3'(c)(z-c)$ gives exactly one zero
in the disk. Its multiplicity is therefore one.
The exact center satisfies $0<\Im c-r<\Im c+r<1/2$.
The displayed intervals are outward-rounded enclosures for the
specified disk, computed with the preceding series and tail bound.
Corollary~\ref{cor:gaussian-threshold} now proves $\tau_3^*=\infty$.
This is a zero of an approximating head, not a zero of $\Xi$.
\end{proof}

\subsection{Uniform positive-real tails and a quadratic contour radius}

\begin{theorem}[A positive-real omitted tail, uniformly in the head index]
\label{thm:positive-real-tail}
Put $A=\pi(N+1)^2$, $Y_N=A-7$, and
$G_{N,t}(z)=\widehat{e^{tu^2}R_N}(z)$.
For every $N\ge1$, $0\le t\le A$ and $|\Im z|\le Y_N$,
\[
 \Re G_{N,t}(z)>0 .
\]
For $|\Im z|<Y_N$,
\[
 \left|\frac{G_{N,t}'(z)}{G_{N,t}(z)}\right|
                   \le\frac{2}{Y_N-|\Im z|}.            \tag{PT}
\]
In particular the true omitted theta tail has no zero in this strip.
For $G_N=G_{N,0}$ and $|y|\le1/2$, $|x|>A$, one has
\[
 \Re G_N(x+iy)\ge
 \frac{12e^{-4}A^3e^{-A}}{x^2}\left(1-\frac A{|x|}\right). \tag{TL}
\]
\end{theorem}
\begin{proof}
On the positive half-line write
$f_n(u)=e^{u/2-T}P_0(T)$, $T=\pi n^2e^{2u}$.
The recurrence already used gives
\[
 P_1=-8T^3+30T^2-15T,\qquad
 P_2=16T^4-112T^3+165T^2-75T/2.
\]
For $T\ge12$ the exact identities
\[
 P_1+(T-7)P_0=T(27-4T-4T^2)<0,\qquad
 P_2+2(T-7)P_1=T(60T^2-285T+345/2)>0
\]
hold, and $P_1<0$.
For $n>N$, $T\ge A>12$.
By evenness take $0\le y\le A-7$ and set
$h_n=e^{tu^2}f_n\cosh(yu)$ and $S=2tu+y\tanh(yu)$.
Here
$0\le S\le2Au+A-7\le Ae^{2u}-7\le T-7$.
Thus
\[
 h_n'=e^{tu^2}\cosh(yu)(f_n'+Sf_n)<0,\qquad
 h_n''\ge e^{tu^2}\cosh(yu)(f_n''+2(T-7)f_n')>0.
\]
In the second inequality the other coefficient of $f_n$ is
$2t+4t^2u^2+4tuy\tanh(yu)+y^2\ge0$.
Consequently $q(u)=e^{tu^2}R_N(u)\cosh(yu)$ is strictly
decreasing and strictly convex. Polynomial factors in $T$ and
$e^{tu^2+Y_Nu}$ are dominated by the double-exponential theta decay.
This justifies summation, both derivatives, their integrability,
and vanishing of $q,q'$ at infinity.
For $x\ne0$, two integrations by parts give
\[
 \Re G_{N,t}(x+iy)
       =\frac2{x^2}\int_0^\infty q''(u)(1-\cos xu)\dd u>0.
\]
The strictness follows from $q''>0$ and the discrete zero set
of $1-\cos xu$. At $x=0$ the original positive integral applies.
All closed strip endpoints are included.

On the disk about $z$ of radius $d=Y_N-|\Im z|$, the Cayley function
$(G(w)-G(z))/(G(w)+\overline{G(z)})$ maps to the unit disk and vanishes
at $z$. Schwarz's inequality gives
$|G'(z)|\le2\Re G(z)/d$, proving (PT).

For (TL), retain just $n=N+1$ and $0\le u\le1/A$, with $t=0$.
Since $2|y|\le1$ and $f_n'<0$,
\[
 (f_n\cosh(yu))''\ge f_n''+f_n'
       \ge6T^4e^{u/2-T}\ge6A^4e^{-A-4}.
\]
Indeed $P_2+P_1\ge6T^4$ for $T\ge12$, and
$T=Ae^{2u}\le A+4$ on this interval.
The other slices have positive second derivatives.
Finally $\int_0^{1/A}(1-\cos xu)\dd u
=1/A-\sin(x/A)/x\ge(1-A/|x|)/A$.
This proves (TL).
\end{proof}

\begin{lemma}[An explicit unconditional Xi envelope]
\label{lem:xi-envelope}
For $|x|\ge2$ and $|y|\le1/2$,
\[
 |\Xi(x+iy)|\le6|x|^3e^{-\pi|x|/4}.                     \tag{XE}
\]
\end{lemma}
\begin{proof}
Take $x\ge2$, $s=\sigma+ix$ with $0\le\sigma\le1$, and
$M=\lceil x\rceil$.
Euler--Maclaurin, obtained by twice integrating with the periodic
Bernoulli polynomials, gives
\[
 \zeta(s)=\sum_{n<M}n^{-s}+\frac{M^{1-s}}{s-1}
 +\frac12M^{-s}+\frac{s}{12}M^{-s-1}
 -\frac{s(s+1)}2\int_M^\infty B_2(\{u\})u^{-s-2}\dd u .
\]
The integral converges for $\sigma>-1$; equality first holds for
$\sigma>1$ and extends meromorphically. Since $|B_2|\le1/6$,
the last term is at most
$|s(s+1)|M^{-\sigma-1}/[12(\sigma+1)]$.
The five displayed terms are bounded, respectively, by
$x$, $3/2$, $1/2$, $3/32$ and $9x/64$.
These follow from $x\le M\le3x/2$,
$|s|\le9x/8$ and $|s+1|\le3x/2$.
Their sum is below $3x$.

Use the explicit log-Gamma Stirling remainder
\cite[\S5.11(i)--(ii), equation (5.11.1)]{DLMFGamma}:
at truncation before its Bernoulli terms,
$|R(w)|\le\sec^2(\arg w/2)/(12|w|)$.
For $w=s/2=b+ix/2$, $0\le b\le1/2$,
$|\arg w|\le\pi/2$, $|w|\ge1$, so $|R(w)|\le1/6$.
Moreover $(x/2)\arctan(2b/x)\le b$ and $|w|^{b-1/2}\le1$.
The real part of the Stirling formula therefore gives
$|\Gamma(s/2)|\le\sqrt{2\pi}e^{1/6}e^{-\pi x/4}
<3e^{-\pi x/4}$.
The last numerical inequality follows from
$e^{1/3}<17/12$ and $\pi<22/7$.
Since $|\pi^{-s/2}|\le1$ and
$|s(s-1)|\le(5/4)x^2$, the definition of $\Xi$ gives
$|\Xi|\le(45/8)x^3e^{-\pi x/4}<6x^3e^{-\pi x/4}$.
Conjugation treats negative $x$.
\end{proof}

\begin{corollary}[Every head has an explicit quadratic strip radius]
\label{cor:quadratic-head-radius}
For every $N\ge1$, every zero of $\mathcal F_N$ in
$|\Im z|\le1/2$ satisfies $|\Re z|<X_N$, where
\[
 X_N=4(N+1)^2+\frac8\pi\log\!\bigl(\pi(N+1)^2\bigr)+16. \tag{QR}
\]
In fact $\Re\mathcal F_N(x+iy)<0$ for $|x|\ge X_N$,
$|y|\le1/2$.
For each $\epsilon>0$, all sufficiently large $N$ have no strip
zero with $|\Re z|\ge(4+\epsilon)(N+1)^2$.
\end{corollary}
\begin{proof}
Again put $A=\pi(N+1)^2>12$, $x_0=X_N$.
One has $x_0/A<10/3$, using $\pi>3$, $\log A/A\le1/4$ and
$16/A<4/3$.
For $x\ge x_0$, $1-A/x>1-\pi/4>3/14$.
By (TL) and (XE),
\[
 \frac{|\Xi(x+iy)|}{\Re G_N(x+iy)}
 <\frac73 e^4\frac{x^5}{A^3}e^{-\pi x/4+A}.
\]
The right side decreases for $x\ge x_0$.
At $x_0$ it is at most
$\frac73(10/3)^5e^{4-4\pi}
<\frac{700000}{729}e^{-8}<1$.
The final inequality follows from $e>8/3$ and
$6300000<8^8$.
Thus $\Re(\Xi-G_N)<0$ there.
Evenness and conjugation give the negative side.
Finally (QR) is $4(N+1)^2+O(\log(N+1))$, proving the
stated quantified asymptotic consequence.
\end{proof}

\begin{proposition}[The on-axis multiplicity obstruction is exact]
\label{prop:head-multiplicity}
Suppose $\Xi$ has a real zero $\gamma$ of order $m\ge1$, and
let $c=\Xi^{(m)}(\gamma)/m!\in\R\setminus\{0\}$.
Put $r_N=(G_N(\gamma)/|c|)^{1/m}$.
Then $r_N\to0$ and, locally uniformly in $w\in\C$,
\[
 \frac{\mathcal F_N(\gamma+r_Nw)}{G_N(\gamma)}
                  \longrightarrow \operatorname{sgn}(c)w^m-1. \tag{MS}
\]
For every sufficiently large $N$, the $m$ headzeros in a fixed
sufficiently small disk about $\gamma$ are simple. Their real count is
one for odd $m$, two for even $m$ with $c>0$, and zero for even $m$
with $c<0$. All the others are nonreal stripzeros.
Consequently a globally positive unbounded head sequence would also
exclude $m\ge3$ and double zeros with $c<0$. These restrictions are
not part of RH; no such multiple $\Xi$-zero is asserted to exist.
\end{proposition}
\begin{proof}
Theorem~\ref{thm:positive-real-tail} gives $G_N(\gamma)>0$,
and the positive integrable theta tail gives $G_N(\gamma)\to0$.
On each fixed scaled compact set, (PT), integrated on straight segments,
gives
$\log(G_N(\gamma+r_Nw)/G_N(\gamma))\to0$ uniformly.
The true Taylor expansion of $\Xi$ gives the other term in (MS),
with error $O_{\gamma,m,w}(r_N)$ uniformly on that compact.
Every root of the limiting polynomial is simple.
Disjoint small disks about its roots, Rouch\'e's theorem, and conjugation
give one simple headzero per disk; roots in the disks with real centers
are exactly real because each such disk is invariant under conjugation.
All other disks stay away from the real axis.
Uniform convergence $\mathcal F_N\to\Xi$ on the boundary of the fixed
original disk counts precisely $m$ zeros there.
The scaled disks account for them all.
The limiting real polynomial has exactly the claimed real-root count,
which proves every assertion.
\end{proof}

\subsection{Shrinking physical heat removes that local obstruction}

\begin{theorem}[Arbitrary critical multiplicities split into real headzeros]
\label{thm:heated-head-local}
Put $A_N=\pi(N+1)^2$, $t_N=A_N^{-1}$ and
\[
 F_t(z)=\widehat{e^{tu^2}\Phi}(z),\qquad
 H_{N,t}(z)=\widehat{e^{tu^2}\Phi_N}(z).
\]
For every fixed real $\Xi$-zero $\gamma$ of arbitrary finite order $m$,
with $c=\Xi^{(m)}(\gamma)/m!$, one has locally uniformly
\[
 t_N^{-m/2}H_{N,t_N}(\gamma+\sqrt{t_N}\,w)
 \longrightarrow cP_m(w),\quad
 P_m(w)=m!\sum_{k=0}^{\lfloor m/2\rfloor}
              \frac{(-1)^kw^{m-2k}}{k!(m-2k)!}.          \tag{HH}
\]
For every sufficiently large $N$, the $m$ zeros in a fixed small disk
about $\gamma$ are all simple and real, at distance $O_\gamma(\sqrt{t_N})$.
The threshold depends on $\gamma,m$; no uniformity over all zeros is
asserted. For every fixed $Y<\infty$,
\[
 \sup_{|\Im z|\le Y}|H_{N,t_N}(z)-\Xi(z)|\longrightarrow0. \tag{HC}
\]
Physical heat moves zeros; subsequent Gaussian convolution preserves
the new zeros. This operation differs from the zero-preserving
Gaussian convolution alone.
\end{theorem}
\begin{proof}
The double-exponential decay makes $F_t(z)$ entire jointly in complex
$t,z$, with all derivatives justified by integrable compact majorants.
For $0\le t\le1$, $Y\ge0$, set
\[
 C_Y=\exp\sup_{u\ge0}\{u^2+(Y+9/2)u-(\pi/2)e^{2u}\},
 \quad
 B_Y=8\pi C_Y\sum_{k\ge0}(1+k)^2e^{-\pi k-\pi k^2/2}<\infty .
\]
The exact positive omitted tail and $e^{2u}\ge1+2u$ give
\[
 \sup_{|\Im z|\le Y}|F_t(z)-H_{N,t}(z)|
       \le B_Y(N+1)^2e^{-\pi(N+1)^2/2}.                 \tag{HT}
\]
Indeed bound each summand by
$4\pi^2 C_Yn^4e^{-\pi n^2e^{2u}/2}$, integrate both halves to obtain
$8\pi C_Yn^2e^{-\pi n^2/2}$, and write $n=N+1+k$.

Joint Taylor convergence gives
$F_t(z)=\sum_{k\ge0}(-t)^k\Xi^{(2k)}(z)/k!$.
Set $t=v^2$, $z=\gamma+vw$.
All coefficients of $v^\ell$ with $\ell<m$ vanish, and the coefficient
of $v^m$ is $cP_m(w)$.
The error is $O_{\gamma,m,R}(v^{m+1})$ on $|w|\le R$.
At $v=\sqrt{t_N}$, (HT) divided by $t_N^{m/2}$ tends to zero,
proving (HH).

For completeness, $P_m$ is monic, real, and satisfies
$P_m(w)=(-2)^me^{w^2/4}D_w^me^{-w^2/4}$.
Integrating by parts $m$ times makes it orthogonal to every real
polynomial of degree below $m$ against $e^{-w^2/4}\dd w$.
If its real sign-changing roots were fewer than $m$, multiply their
linear factors to get a polynomial of degree below $m$ whose product
with $P_m$ has a constant sign and a nonzero weighted integral.
This contradicts orthogonality. Thus all $m$ roots are distinct and real.
Disjoint conjugation-invariant disks about these roots and (HH)
give exactly one simple real headzero in each.

Lastly dominated convergence gives
$\sup_{|\Im z|\le Y}|F_t(z)-\Xi(z)|
\le\int\Phi(u)(e^{tu^2}-1)e^{Y|u|}\dd u\to0$.
Together with (HT) this proves (HC).
Rouch\'e on the fixed disk around $\gamma$ counts exactly $m$ zeros;
the preceding real clusters give all of them.
The same proof works for any $t_N\downarrow0$ for which the right side
of (HT) is $o(t_N^{m/2})$ for every fixed integer $m$,
including $t_N=(1+\log(N+1))^{-1}$.
\end{proof}

\begin{theorem}[A uniform contour radius for every heated head]
\label{thm:heated-head-radius}
For $N\ge1$, $A=\pi(N+1)^2$, $0\le t\le1$, and $|y|\le1/2$,
\[
 |F_t(x+iy)|<132e^{-\pi|x|/8},\qquad
 \Re G_{N,t}(x+iy)\ge
       5e^{-4}A^3e^{-A}/x^2\quad(|x|\ge2A).              \tag{WR}
\]
In particular $\Re H_{N,t}(x+iy)<0$ for $|x|\ge5A$.
Every heated-head zero in the closed strip therefore lies in $|x|<5A$.
All constants are uniform in $N,t$ and the entire closed strip.
\end{theorem}
\begin{proof}
For the omitted tail put $q=e^{tu^2}R_N(u)\cosh(yu)$ on $u\ge0$.
The proof of (PT) gives $q''>0$ and
$\Re G_{N,t}=2x^{-2}\int q''(1-\cos xu)\dd u$.
On $0\le u\le1/A$, the logarithmic derivative
$b=2tu+y\tanh(yu)$ lies in $[0,2/3]$.
Since $f_n'<0$, retain the first omitted slice to obtain
\[
 q''\ge e^{tu^2}\cosh(yu)\{f_{N+1}''+\tfrac43f_{N+1}'\}
       \ge5A^4e^{-A-4}.
\]
Indeed $P_2+(4/3)P_1
=16T^4-(368/3)T^3+205T^2-(115/2)T\ge5T^4$ for $T\ge12$,
and $A\le T=Ae^{2u}\le A+4$ here.
For $|x|\ge2A$,
$\int_0^{1/A}(1-\cos xu)\dd u
=A^{-1}-\sin(x/A)/x\ge(2A)^{-1}$.
This proves the lower bound.

For the upper bound use the full, unreflected theta series for complex $w$.
It is holomorphic on $|\Im w|<\pi/4$ by normal convergence.
Jacobi modularity makes it even on the real line; the identity
theorem therefore gives $\Phi(w)=\Phi(-w)$ on this complex strip.
For $x\ge0$ shift the Fourier contour to $w=u-i\pi/8$.
The two vertical sides vanish: their theta factors are bounded by
$C e^{9|u|/2-c e^{2|u|}}$, which dominates
$e^{u^2+|u|/2}$ uniformly on those sides.
The shifted exponential has modulus at most
$e^{u^2+|u|/2-\pi x/8}$.
Writing $c_0=\pi/\sqrt2$, positive majorants on the two halves give
\[
 |F_t(x+iy)|\le e^{-\pi x/8}
 2e^{5/4}\sum_{n\ge1}(4\pi\sqrt2\,n^2+6\sqrt2)
                         e^{-\pi n^2/(2\sqrt2)} .
\]
Here $u^2+5u-(c_0/2)e^{2u}\le-u^2+3u-1\le5/4$,
and the remaining integral is at most
$e^{-c_0n^2/2}/(c_0n^2)$. All majorants are nonnegative.
Elementary bounds $\pi<4$, $\sqrt2<3/2$, $e^{5/4}<4$ and
\[
 \sum_{n\ge1}n^2e^{-n^2}
 <3/8+\sum_{n\ge2}n^2\,7^{-n}
 =3/8+22/189<1/2
\]
make the displayed constant strictly less than $132$.
Evenness gives $x<0$ as well.
For $|x|\ge5A$, the ratio of these two bounds is at most
\[
 \frac{660e^4}{A}e^{-(5\pi/8-1)A}<55e^{-13/2}<1.
\]
Use $A>12$, $\pi>3$, and $e^{13/2}>64$.
Thus $\Re H_{N,t}\le|F_t|-\Re G_{N,t}<0$.
\end{proof}

\begin{theorem}[The shrinking-heat construction also fails at the third head]
\label{thm:heated-third-certificate}
For $N=3$ and $t=1/(16\pi)=A_3^{-1}$ the heated head $H_{3,t}$ has
exactly one zero, simple and nonreal, in the disk with
\[
 \begin{aligned}
 \Re c&=\frac{90204354118549976947605343860520658130}{2^{120}},\\
 \Im c&=\frac{570243309641866082381721850049555627}{2^{120}},
 \qquad r=2^{-90}.
 \end{aligned}                                         \tag{WD}
\]
The disk lies strictly in $0<\Im z<1/2$.
This is a finite-head zero, not a zero of $\Xi$ or $\zeta$.
\end{theorem}
\begin{proof}
For derivative orders $j=0,1,2,3$, $u^j\le e^{u^2}$ on $u\ge0$.
Also $2u^2+5u\le(\pi/2)e^{2u}$, from
$e^{2u}\ge1+2u+2u^2$ and $\pi>3$.
The absolute derivative integrands on the whole closed strip are
therefore bounded by $4\pi^2n^4e^{-\pi n^2e^{2u}/2}$.
Including both halves gives the uniform moment bound
\[
 M=8\pi\sum_{n=1}^3 n^2e^{-\pi n^2/2}<6.
\]
For a quadrature cutoff $U$, the error in each complex derivative
has modulus at most
$T_U=8\pi e^{-2U}\sum_{n=1}^3n^2e^{-\pi n^2e^{2U}/2}$.
The saved $600$-bit complex interval integrals on $[0,3]$,
enlarged in both coordinates by $[-T_3,T_3]$, give
\[
 |H_{3,t}(c)|+Mr^2/2<1.778\cdot10^{-54}
             <1.320\cdot10^{-47}<|H_{3,t}'(c)|r .
\]
The centre and radius are exact dyadics and the disk lies within the strip.
Taylor's integral remainder and Rouch\'e count exactly one zero.
A separate $500$-bit evaluation with cutoff $5/2$, three quadrature intervals
and the coarser bound $M=6$ gives a strict Rouch\'e ratio below
$1.491\cdot10^{-7}$ for the same disk.
\end{proof}

The heated omitted tail obeys (PT) and has the uniform radius (WR),
but its positivity does not exclude zeros of the heated head.
Global strip reality on an unbounded
heated-head sequence remains a separate obligation. Any such sequence
would imply RH by (HC) and the same small-contour argument as
Corollary~\ref{cor:head-exhaustion}, with every multiplicity retained.

\subsection{The growing cusp creates nonreal zeros after shrinking heat}

We prove that the preceding local Hermite mechanism does not extend to
global strip reality for the two following heat regimes. The zeros
constructed here belong to finite heated heads, not to $\zeta$.
No assumption about the reality or simplicity of the zeros of $\Xi$
is made.

\begin{theorem}[Cusp zeros at every fixed positive logarithmic heat]
\label{thm:cusp-fixed-logheat}
Let $K>0$ and
\[
 \max\{0,(1-K)/2\}<\delta<1/2,
 \quad A_N=\pi(N+1)^2,\quad t_N=K/\log A_N.
\]
For every sufficiently large integer $N$, $H_{N,t_N}$ has a simple zero
$z_N$ satisfying
\[
 \Im z_N=\delta+O_{K,\delta}((\log A_N)^{-1}),
 \qquad
 \Re z_N=\frac{4A_N}{\pi}
  +\frac4\pi\left(\frac34+\frac\delta2+\frac K4\right)\log A_N
  +O_{K,\delta}(1).                                    \tag{CZ1}
\]
In particular $0<\Im z_N<1/2$ eventually.
\end{theorem}

\begin{theorem}[Cusp zeros with logarithmic heat tending to zero]
\label{thm:cusp-small-logheat}
Let $t_N>0$, put $A_N=\pi(N+1)^2$, $L_N=\log A_N$,
$\eta_N=t_NL_N$, and suppose
\[
 \eta_N\longrightarrow0,\qquad
 \frac{\eta_NL_N}{\log(2+1/\eta_N)}\longrightarrow\infty. \tag{CZ2}
\]
For every sufficiently large $N$, $H_{N,t_N}$ has a simple zero $z_N$
with
\[
 \Im z_N=\frac12-\frac{\eta_N}4+O(L_N^{-1}),
\]
\[
 \Re z_N=\frac{4A_N}{\pi}+\frac4\pi L_N
       +\frac{\eta_N}{2\pi}L_N
       +O(\log(2+1/\eta_N)).                           \tag{CZ3}
\]
The zero lies strictly inside $0<\Im z<1/2$.
In particular this holds for $t_N=L_N^{-\beta}$, $1<\beta<2$,
and for $t_N=(\log L_N)^2/L_N^2$.
\end{theorem}

We give the complete joint-parameter argument, starting from the exact
primary normalization. Polymath~\cite[Theorem 1.3, printed p.~6,
eqs.~(13)--(24)]{PolymathHeat} uses a cosine transform on $[0,\infty)$
with $\Phi_{\rm ours}(u)=2\Phi_{\rm PM}(u/2)$. Consequently
\[
 F_t(z)=8H^{\rm PM}_{4t}(2z).                            \tag{PNH}
\]
Set $s=1/2-iz$, $T=4t$, $X=2x$, $Y=2y$, $Q=X/(4\pi)=x/(2\pi)$,
and define the explicit nonvanishing factors
\[
 \begin{split}
 M_0(s)&=\frac{s(s-1)}{16}\pi^{-s/2}\sqrt{2\pi}
  \exp\left[(s/2-1/2)\operatorname{Log}(s/2)-s/2\right],\\
 \alpha(s)&=\frac1{2s}+\frac1{s-1}
                         +\frac12\operatorname{Log}\frac{s}{2\pi},\\
 M_T(s)&=e^{T\alpha(s)^2/4}M_0(s),\quad
 D_t(z)=8M_{4t}(s),\quad s_*(z)=s+2t\alpha(s).
 \end{split}                                          \tag{PMG}
\]
The branches are those in the primary definitions (6)--(10):
\[
 \log M_0(s)=\operatorname{Log}s+\operatorname{Log}(s-1)
 -\frac s2\log\pi+\log\frac{\sqrt{2\pi}}{16}
 +(s/2-1/2)\operatorname{Log}(s/2)-s/2.
\]
All regions below stay away from the branch cut $(-\infty,1]$.
The primary theorem holds unconditionally in precisely
\[
 0<T\le1/2,\qquad X\ge200,\qquad 0\le Y\le1.
\]
Put
\[
 J=\left\lfloor\sqrt{Q+T/16}\right\rfloor,
 \quad b_n=e^{(T/4)\log^2 n},
 \quad
 \gamma=\frac{M_T(1/2-y+ix)}{M_T(1/2+y-ix)},
\]
\[
 \kappa_{\rm PM}=\frac T2
       [\alpha(1/2-y+ix)-\alpha(1/2+y+ix)].
\]
Its full approximation, retaining the conjugation and both sums, is
\[
 \frac{F_t(z)}{D_t(z)}
  =\sum_{n\le J}\frac{b_n}{n^{s_*}}
       +\gamma\sum_{n\le J}\frac{n^Yb_n}{n^{\overline{s_*}+\kappa_{\rm PM}}}
       +O_{\le}(e_A+e_B+e_{C,0}).                       \tag{PRS}
\]
The needed primary bounds are
\[
 |\gamma|\le e^{0.02Y}Q^{-Y/2},\qquad
 |\kappa_{\rm PM}|\le\frac{TY}{2(X-6)},
\]
\[
 e_A+e_B\le\sum_{n\le J}(1+|\gamma|J^{|\kappa_{\rm PM}|}n^Y)
  \frac{b_n}{n^{\Re s_*}}
 \left\{\exp\left(\frac{(T^2/16)\log^2(Q/n^2)+0.626}{X-6.66}\right)-1\right\},
\]
\[
 e_{C,0}\le Q^{-(1+Y)/4}
 \exp\left[-\frac T{16}\log^2Q
 +\frac{1.24(3^Y+3^{-Y})}{J-0.125}
 +\frac{3|\log Q+i\pi/2|+10.44}{X-12}\right].          \tag{PRE}
\]
In particular the constant $0.626$ is retained. Directly from (PMG),
uniformly for bounded $y$ and large $x$,
\[
 \Re s_*=(1+Y)/2+(T/4)\log Q+O(T/X),
\]
\[
 \Re s_*-(T/4)\log n
       \ge(1+Y)/2+(T/8)\log Q-O(T/X)\quad(n\le J).    \tag{PNE}
\]
The $T/2$ coefficient printed in the preliminary asymptotic proof on
p.~57 of the primary version is a factor-two error;
its definition (17) and its Theorem 1.3, eq.~(21), have the correct
$T/4$ coefficient. We use (PRS)--(PRE) directly, rather than that
incorrect intermediate estimate. Neither the extended $|Y|\le10$
region of Proposition 9.1 nor its high-zero Theorem 1.5 is imported
as a statement about finite heads.

\begin{lemma}[A joint Lorentz limit for the heated omitted tail]
\label{lem:heated-tail-crossover}
Let $M=N+1$, $A=\pi M^2$, and
\[
 \mathcal L_M(z)=\frac{16A^3e^{-A}}{4A^2+z^2}.
\]
Uniformly for $0\le t\le1$, $|y|\le1/2$ and $x/A$ in any fixed
compact real interval,
\[
 G_{N,t}(z)/\mathcal L_M(z)\longrightarrow1.            \tag{PLC}
\]
\end{lemma}
\begin{proof}
For the first omitted slice scale $v=Au$. Its even Fourier integral is
\[
 8Ae^{-A}\int_0^\infty f_{A,t}(v)\cos(zv/A)\dd v,
\]
\[
 f_{A,t}(v)=e^{tv^2/A^2+(9/2)v/A-A(e^{2v/A}-1)}
              \left(1-\frac3{2A}e^{-2v/A}\right).
\]
It tends locally uniformly to $e^{-2v}$. With $u=v/A$,
$tu^2+5u\le(A/2)(e^{2u}-1)$ for $A\ge5$, by
$e^{2u}-1\ge2u+2u^2$. Thus
$|f_{A,t}(v)\cos(zv/A)|\le e^{-v}$, including the
$e^{u/2}$ cost of the complex cosine. Splitting at a fixed $v$ cutoff
and then sending that cutoff to infinity proves uniform convergence.
The limiting integral is $2/(4+(z/A)^2)$, uniformly bounded away
from zero in the stated region, and gives exactly $\mathcal L_M$.
The same physical majorant gives for every $n>M$
\[
 |\widehat{e^{tu^2}\phi_n}(z)|\le8A_ne^{-A_n},
 \qquad A_n=\pi n^2.
\]
Their sum divided by $Ae^{-A}$ is at most
\[
 8\sum_{j\ge1}(1+j/M)^2e^{-2\pi Mj-\pi j^2}\longrightarrow0.
\]
This follows by dominated convergence with $M\ge2$; every tail index
is retained. This proves (PLC).
\end{proof}

\begin{proof}[Proof of Theorem~\ref{thm:cusp-fixed-logheat}]
Write $A=A_N$, $t=K/\log A$. Choose a fixed small $\varepsilon>0$
such that $y_\pm=\delta\pm\varepsilon$ lie in $(0,1/2)$ and
$p_0=1/2+y_-+K/2>1$. Work on
\[
 y_-\le y\le y_+,\qquad x=4A/\pi+O_{K,\delta}(\log A).
\]
All original primary hypotheses hold eventually. Choose $1<a<p_0$.
The corrected net exponent (PNE) gives
$b_n n^{-\Re s_*}\le n^{-a}$ for $n\le J$ and $\Re s_*\ge a$.
Splitting at a fixed integer $P$ therefore proves uniformly
\[
 \sum_{n\le J}b_n n^{-s_*}=\zeta(s_*)+o(1):
\]
on $n\le P$ the weights tend to one; both remaining tails are
bounded by $\sum_{n>P}n^{-a}$, independently of the large imaginary
part of $s_*$. For the second sum,
$|\gamma|J^YJ^{|\kappa_{\rm PM}|}=O(1)$ gives the same summable
majorant, and $\gamma\to0$ on each fixed finite prefix since $Y>0$.
Hence it tends to zero. The exponential bracket in (PRE) is
$O_K(A^{-1})$, because $T^2\log^2(Q/n^2)=O_K(1)$.
Thus $e_A+e_B=O_{K,\delta}(A^{-1})$ and $e_{C,0}=O(A^{-1/4})$.
We have proved
\[
 F_t(z)/D_t(z)=\zeta(s_*(z))+o(1).                      \tag{PZA}
\]
The arithmetic factor is retained, not replaced by one.

The entire $s_*$-image lies in a fixed half-plane $\Re s_*>b>1$.
There the absolute Euler logarithm
\[
 Z(s)=\sum_p\sum_{r\ge1}p^{-rs}/r=\log\zeta(s)          \tag{EZH}
\]
and its first two derivatives are uniformly bounded: each series
is majorized by $\sum_{n\ge2}(\log n)^j n^{-b}$ for $j=0,1,2$.
The identity $e^Z=\zeta$ follows by expanding the absolutely
convergent Euler product and unique prime factorization.
Every prime power is included. In particular $\zeta(s_*)$ is bounded
away from zero, so (PZA) proves
$F_t/(D_t\zeta(s_*))\to1$ uniformly. This quotient and the ratio in
(PLC) eventually have holomorphic logarithms tending uniformly to zero.

Use the exact analytic model
\[
 \ell_N(z)=\log D_t(z)+Z(s_*(z))-\log\mathcal L_M(z).    \tag{PMO}
\]
The denominator $4A^2+z^2$ has positive real part in this rectangle,
so its principal log suffices. The explicit gamma logarithm gives
\[
 \Re\ell_N(x+i\delta)
  =A-\frac\pi4x+(7/4+\delta/2)\log x-\log A
     +\frac t4\log^2\frac{x}{2\pi}
     +\Re Z(s_*(x+i\delta))+O(1+t).                    \tag{PAM}
\]
Indeed $s(s-1)$ contributes $2\log x$, the Stirling power contributes
$(\Re s/2-1/2)\log x$, the gamma exponential contributes $-\pi x/4$,
and $t\Re\alpha^2=(t/4)\log^2(x/(2\pi))+O(t)$.
The Lorentz log contributes $A-\log A+O(1)$.
At
\[
 X_N^*=4A/\pi+(4/\pi)(3/4+\delta/2+K/4)\log A
\]
and bounded real $q$, (PAM) is $-\pi q/4+O_{K,\delta}(1)$
at $x=X_N^*+q$. Choosing two fixed endpoints with opposite signs
therefore gives $x_N=X_N^*+O_{K,\delta}(1)$ with
$\Re\ell_N(x_N+i\delta)=0$. No amplitude monotonicity is assumed.

Derivatives are taken only of explicit holomorphic formulas.
Since $\alpha'=O(A^{-1})$, $\alpha''=O(A^{-2})$, the Euler-log
bounds and chain rule give throughout the rectangle
\[
 \ell_N'(z)=-\frac\pi4-\frac i2\log\frac{x}{2\pi}+O_{K,\delta}(1),
 \qquad \ell_N''(z)=O_{K,\delta}(1).                   \tag{PHD}
\]
Let $z_0=x_N+i\delta$, $v_N=\ell_N'(z_0)$; then
$|v_N|\sim(\log A)/2$ and $\Re v_N=O_{K,\delta}(1)$.
Choose $j_N\in\mathbb Z$ and $|\theta_N|\le\pi$ with
$\ell_N(z_0)-2\pi i j_N=i\theta_N$, and set
\[
 z_c=z_0-i\theta_N/v_N,\qquad r_N=1/|v_N|.
\]
The real center shift is $O(1/\log A)$ and its imaginary shift
is $O_{K,\delta}(1/\log^2 A)$; the closed disk lies in the rectangle.
Taylor's formula gives uniformly on that disk
\[
 \ell_N(z)-2\pi i j_N=v_N(z-z_c)+O_{K,\delta}(1/\log^2 A).
\]
The actual analytic log-ratio is
\[
 \mathcal L_N=\ell_N+
       \log\frac{F_t}{D_t\zeta(s_*)}
       -\log\frac{G_{N,t}}{\mathcal L_M}.
\]
The last two terms are $o(1)$ uniformly. On the circle the linear
term has modulus one. Rouch\'e gives exactly one zero, counting
multiplicity, of $\mathcal L_N-2\pi i j_N$ in the disk.
On the whole disk $|\mathcal L_N-2\pi i j_N|\le1+o(1)<2<2\pi$,
so no other exponential integer level creates an additional zero.
Since $e^{\mathcal L_N}=F_t/G_{N,t}$ and $G_{N,t}$ is nonzero there,
these zeros and their orders are exactly those of $H_{N,t}$.
The unique zero is simple and has the location (CZ1).
\end{proof}

\begin{proof}[Proof of Theorem~\ref{thm:cusp-small-logheat}]
Write $A=A_N$, $L=\log A$, $\eta=tL$, $h=\eta L$.
We first establish the needed Euler-log bounds uniformly as the real
part approaches one. For $1<\sigma\le2$, the absolute prime-power
series and elementary real integral comparisons give
\[
 |Z(s)|\le\log(1+1/(\sigma-1)),\quad
 |Z'(s)|\le\frac2{\sigma-1},\quad
 |Z''(s)|\le\frac C{(\sigma-1)^2},\qquad \Re s=\sigma. \tag{EZS}
\]
Here $\zeta(\sigma)\ge1/(\sigma-1)$ and
$-\zeta'(\sigma)\le2/(\sigma-1)^2$.
For the latter bound retain $n=2,3$ separately: their sum is less
than $\log2/2+\log3/3<1$; the remaining decreasing sum is bounded
by the integral from $3$, and the full integral is $(\sigma-1)^{-2}$.
The absolute prime-power series for $Z'$ equals
$-\zeta'(\sigma)/\zeta(\sigma)$.
For $Z''$ it equals $Z''(\sigma)\le\zeta''(\sigma)/\zeta(\sigma)$;
the finite prefix $n<8$ and the remaining decreasing integral give
$\zeta''(\sigma)\le C(\sigma-1)^{-3}$.
This proves (EZS) without an interior-strip zero-free premise.

Consider the shrinking rectangle
\[
 |y-(1/2-\eta/4)|\le\eta/32,
 \qquad x=4A/\pi+O(L).                                \tag{SRH}
\]
All original primary hypotheses hold eventually, including $Y<1$.
From (PMG), uniformly on this rectangle,
\[
 1+11\eta/16\le\sigma_*:=\Re s_*\le1+13\eta/16,
 \qquad t\log J\le9\eta/16.                           \tag{SNE}
\]
Indeed $\sigma_*=(1+Y)/2+t\log Q+O(t/X)$ and
$t\log Q=\eta+o(\eta)$.
Thus $b_n n^{-\sigma_*}\le n^{-1-\eta/8}$ for $n\le J$,
and $|\zeta(s_*)|\ge c\eta$ by the absolute inverse Euler product.

The near-one relative approximation requires a different argument.
Let $V$ be a centered real Gaussian of variance $2t$, so exactly
$e^{t\log^2n}=\mathbb E e^{V\log n}$.
Put $P=\lfloor A^{1/16}\rfloor<J$. The high indices satisfy
\[
 \sum_{P<n\le J}e^{t\log^2n}n^{-\sigma_*}
       \le C\eta^{-1}e^{-h/128}.                       \tag{GH}
\]
For $n\le P$ tilt the Gaussian and retain the tails explicitly:
\[
 \mathbb E[e^{V\log n}\mathbf1_{|V|>\eta/4}]
 =e^{t\log^2n}\Pr(|W+2t\log n|>\eta/4)
\le2e^{t\log^2n}e^{-h/256},                           \tag{GWT}
\]
where $W$ has the same centered Gaussian law. The last inequality
uses $0\le2t\log n\le\eta/8$ and the Gaussian Chernoff bound.
The low weighted sum is $O(1/\eta)$ by (SNE).
Its noncentral contribution is therefore
$O(\eta^{-1}e^{-h/256})$.
On $|V|\le\eta/4$, $\Re(s_*-V)\ge1+7\eta/16$.
The low unweighted Dirichlet sum can be replaced by $\zeta(s_*-V)$,
with tail cost $C\eta^{-1}e^{-7h/256}$.
After dividing all these errors by $|\zeta(s_*)|\ge c\eta$,
the main term is
\[
 \mathbb E\left[
  \frac{\zeta(s_*-V)}{\zeta(s_*)}\mathbf1_{|V|\le\eta/4}\right].
\]
By (EZS) along the straight real segment,
$|Z(s_*-v)-Z(s_*)|\le5|v|/\eta$, and consequently
\[
 \left|\frac{\zeta(s_*-v)}{\zeta(s_*)}-1\right|
       \le(5|v|/\eta)e^{5|v|/\eta}.
\]
Cauchy--Schwarz and the Gaussian moment formula give
\[
 \mathbb E[|V|e^{a|V|}]
       \le\sqrt{2t}\sqrt{2e^{4a^2t}}=2\sqrt t\,e^{2a^2t}.
\]
At $a=5/\eta$ the resulting relative cost is at most
$10h^{-1/2}e^{50/h}$. The omitted Gaussian mass is at most
$2e^{-h/64}$. We have therefore proved uniformly
\[
 \frac{\sum_{n\le J}b_n n^{-s_*}}{\zeta(s_*)}
     =1+O(h^{-1/2})+O(\eta^{-2}e^{-h/256}).             \tag{GZA}
\]
This proof keeps all integer indices and all prime powers in (EZH).

For the second sum in (PRS), keep the correlation with $Y$:
\[
 r=\sigma_*-Y-t\log J=(1-Y)/2+\eta/2+o(\eta).
\]
Eventually $0<r<1/2$ and $J^{|\kappa_{\rm PM}|}=1+o(1)$.
Thus its absolute value is at most
\[
 C Q^{-Y/2}J^{1-r}
 =C\exp[((1-Y)/4-\eta/4+o(\eta))\log Q]
       \le Ce^{-h/16},                                \tag{GMS}
\]
since $1-Y\le9\eta/16$ on (SRH).
The sum bound $\sum_{n\le J}n^{-r}\le CJ^{1-r}$ has a uniform
constant for $r<1/2$.
In (PRE) the exponential bracket is $O(A^{-1})$:
$T^2\log^2(Q/n^2)=O(\eta^2)$ and the fixed $0.626$ is retained.
The outside sum is $O(1/\eta)$, so
$e_A+e_B=O((\eta A)^{-1})$.
Because $Y\to1$, the other error is $O(A^{-1/3})$; its
negative heat exponent is retained or discarded for this upper bound,
and both remaining positive exponents tend to zero.
Together with (GZA) and (GMS) this proves
\[
 \frac{F_t}{D_t\zeta(s_*)}=1+O(E_N),\quad
 E_N=h^{-1/2}+\eta^{-2}e^{-h/256}
            +(\eta^2 A)^{-1}+(\eta A^{1/3})^{-1}\to0. \tag{GER}
\]
Condition (CZ2) dominates the inverse-$\eta$ exponential cost.
Also $\eta\ge1/L$ eventually, so both algebraic $A$ costs vanish.

The Lorentz limit (PLC) applies uniformly on (SRH).
Use the same exact model (PMO), now at $\delta_N=1/2-\eta/4$.
Formula (PAM) remains uniform, while (EZS) bounds the Euler log by
$O(\log(2+1/\eta))$. Evaluating its affine amplitude on two sides
gives a crossing at
\[
 x_N=4A/\pi+(4/\pi)L+(\eta/(2\pi))L
                         +O(\log(2+1/\eta)).
\]
No monotonicity is used. The explicit chain rule and (EZS) give
\[
 \ell_N'=-\pi/4-(i/2)\log(x/(2\pi))+O(1/\eta),
 \qquad \ell_N''=O(1/\eta^2).                         \tag{GPD}
\]
Since $h=\eta L\to\infty$, the gamma phase dominates:
$|v_N|\sim L/2$ and $\Re v_N=O(1/\eta)$.
Wrap the phase and choose $z_c,r_N$ exactly as in the preceding
proof. The imaginary center shift is $O((\eta L^2)^{-1})$,
and $r_N=O(L^{-1})$. Both are $o(\eta)$, so the disk stays
inside (SRH). The Taylor remainder is $O((\eta L)^{-2})=o(1)$.
The actual log-ratio errors tend to zero by (GER) and (PLC).
The same modulus-one Rouch\'e argument gives exactly one simple head
zero; the whole-disk bound below $2<2\pi$ excludes other integer
levels. Its location is (CZ3), and $\eta L\to\infty$ makes its
distance $\eta/4$ below $1/2$ strictly larger than the location error.
\end{proof}

For both theorems a subsequent Gaussian convolution, or multiplication
of the transform by $e^{-\tau z^2}$, retains these zeros for every
$\tau\ge0$. It cannot make the defect nonnegative throughout the
open strip: if $|F(x+iy)|^2$ were nondecreasing in $y>0$ and vanished
at an interior $y$, it would vanish on the entire lower vertical
segment, contradicting analyticity of a nonzero entire function.
The remaining smaller heat scales, including $t_N=1/\log^2A_N$
and $t_N=A_N^{-1}$, are not decided by these two estimates.
Their Gaussian/Dirichlet margin fails at the specific boundary
identified in (CZ2); the local arbitrary-multiplicity theorem is
unaffected. A successful global construction must address that smaller
arithmetic regime or remove the physical cusp itself.

\subsection{Cofinal physical windows with a completely real high spectrum}
\label{subsec:physical-windows}

The preceding obstruction concerns reflected finite theta heads. We now
retain the complete modular kernel at the origin and truncate its physical
variable instead. For $U>0$ put
\[
 J_{U,t}(z)=2\int_0^U e^{tu^2}\Phi(u)\cos(zu)\dd u,
 \qquad A=\pi e^{2U},\qquad L=\log A.                 \tag{PW1}
\]
This is a real even entire function of exponential type at most $U$.
There is an endpoint jump at $\pm U$, while every modular jet at the
origin is retained. The parameter $t$ is fixed when differentiating in
$z$. All results below allow arbitrary multiplicities in the part of the
spectrum that is not asserted to be simple.

\begin{theorem}[Cofinal Euler--endpoint windows]
\label{thm:physical-windows}
Let $t(U)>0$ be continuous for all sufficiently large $U$ and satisfy
\[
 t(U)\longrightarrow0,
 \qquad t(U)\log(\pi e^{2U})\ge16.                    \tag{PW2}
\]
There exists a strictly increasing unbounded sequence $U_j$ such that
every zero of $J_{U_j,t(U_j)}$ in
\[
 |\Re z|\ge4\pi e^{4/t(U_j)},\qquad |\Im z|\le1/2    \tag{PW3}
\]
is real and simple. These windows converge to $\Xi$ uniformly on every
complete closed horizontal strip of fixed width. Thus $t=L^{-1/2}$
gives the high-spectrum conclusion from $4\pi e^{4\sqrt L}$, while
$t=16/L$ gives it from $4\pi A^{1/4}$.
\end{theorem}

\begin{proof}
We give the uniform endpoint, arithmetic, phase, and zero-transfer
estimates. The claim concerns every zero in (PW3); no sample of model
zeros or favorable finite compression substitutes for that assertion.

\emph{The exact arithmetic factors.}
Use the explicit $M_T$ and $\alpha$ in (PMG). The normalization is
$F_t(z)=8H^{\rm PM}_{4t}(2z)$, with the primary variables
$T=4t$, $X=2x$, $Y=2y$, and $Q=x/(2\pi)$. Theorem~1.3 of
\cite[printed p.~6, eqs.~(13)--(24)]{PolymathHeat} assumes
$0<T\le1/2$, $X\ge200$, and $0\le Y\le1$; these hypotheses
hold eventually in the part of the strip with $y\ge0$.
Set
\[
 \begin{gathered}
 s_1=1/2-iz,\quad s_2=1/2+iz,\quad
 s_j^*=s_j+2t\alpha(s_j),\quad D_j=8M_{4t}(s_j),\\
 J=\lfloor\sqrt{Q+t/4}\rfloor,
 \quad P_{t,J}(s)=\sum_{n\le J}e^{t\log^2 n}n^{-s}.
 \end{gathered}                                      \tag{PW4}
\]
For $G^\#(z)=\overline{G(\bar z)}$ we have $D_2=D_1^\#$
and $s_2^*=(s_1^*)^\#$. The conjugated denominator in the second
primary sum is essential: exactly
\[
 \overline{s_1^*}+\kappa_{\rm PM}-Y
 =\overline{s_1}+2t\alpha(\overline{s_1})
  +2t\{\alpha(s_2)-\alpha(\overline{s_1})\}-Y=s_2^*.
                                                            \tag{PW5}
\]
The two complete sums in (PRS) are therefore
$D_1P_{t,J}(s_1^*)+D_2P_{t,J}(s_2^*)$.

Put $X_0=2\pi e^{4/t}$. Uniformly for $x\ge X_0$ and
$|y|\le1/2$, the actual exponent, obtained from the explicit $\alpha$, is
\[
 \Re s_j^*\ge t\log Q-O(t/x)\ge4-o(1),\qquad
 \Re s_j^*-t\log n\ge(t/2)\log Q-O(t/x)
                    \ge2-o(1)\quad(n\le J).          \tag{PW6}
\]
Consequently, for all large $U$,
\[
 |P_{t,J}(s_j^*)-\zeta(s_j^*)|
 \le t\sum_{n\le J}(\log n)^2n^{-3/2}
       +\sum_{n>J}n^{-3}
 \le Ct+C/x.                                        \tag{PW7}
\]
Here $e^v-1\le ve^v$ controls the finite weighted sum. The omitted
ordinary zeta tail is absolutely convergent. There is no estimate for
a divergent infinite heat-weighted Dirichlet series.
The complete primary errors (PRE), including their fixed constants,
are bounded by
$C(1+t^2\log^2x)/x+Cx^{-1/4}$ on this range. Indeed
$|\gamma|J^YJ^{|\kappa_{\rm PM}|}\le C$, the outside weighted sum
has exponent at least $3/2$, and the exponential bracket in (PRE)
is $O((1+t^2\log^2x)/x)$. The supremum of both displayed errors
on $x\ge X_0$ is $o(t)$. Thus
\[
 F_t=D_1\zeta(s_1^*)+D_2\zeta(s_2^*)
              +O(t)(|D_1|+|D_2|)                    \tag{PW8}
\]
uniformly on the whole high strip. Negative $y$ follows by conjugation,
within the original source hypotheses.

In the half-plane $\Re s\ge3$ use the absolutely convergent Euler log
\[
 Z(s)=\log\zeta(s)
     =\sum_p\sum_{m\ge1}\frac{p^{-ms}}m
     =\sum_{n\ge2}\frac{\Lambda(n)}{\log n}n^{-s}.   \tag{PW9}
\]
All prime powers remain present. The functions $Z,Z',Z''$ are uniformly
bounded there, and $1/\zeta(3)\le|\zeta(s)|\le\zeta(3)$.
On $\Re s\ge15$ they satisfy the explicit bounds
\[
 |Z(s)|,\ |Z'(s)|,\ |Z''(s)|<10^{-3}.              \tag{PW10}
\]
For example $\Lambda(n)\le\log n\le n$ and
$\Lambda(n)\log n\le(\log n)^2\le n^2$ bound the last two sums
by $\sum_{n\ge2}n^{-14}$ and $\sum_{n\ge2}n^{-13}$.
The elementary bounds $2^{-a}+2^{1-a}/(a-1)$ are below $10^{-3}$
for both $a=14,13$; the first sum is smaller.

\emph{The complete physical tail at every frequency.}
Write $f_t(u)=e^{tu^2}\Phi(u)$ and
\[
 E(z)=\frac{f_t(U)e^{-iUz}}{2A+iz},\qquad
 T_{U,t}=F_t-J_{U,t}.
\]
For $0\le t\le1$, $|y|\le1/2$, and every real $x$, we claim
\[
 T_{U,t}=E+E^\#
       +O((1+tU)/A)(|E|+|E^\#|).                  \tag{PW11}
\]
To prove the componentwise relative error, scale $v=A(u-U)$ and put
$q(v)=f_t(U+v/A)/f_t(U)$, so $q(0)=1$. The first theta slice has
the exact expression
\[
 \begin{gathered}
 q_1(v)=e^{-2v}e^{w(v)}r(v),\\
 w(v)=(2tU+9/2)v/A+tv^2/A^2
                 -A(e^{2v/A}-1-2v/A),\\
 r(v)=\frac{1-(3/(2A))e^{-2v/A}}{1-3/(2A)}.
 \end{gathered}                                      \tag{PW12}
\]
On $v\le A^{1/4}$, Taylor bounds give
$|w|\le C((1+tU)v+v^2)/A$,
$|w'|\le C(1+tU+v)/A$,
$|r-1|\le Cv/A^2$, and $|r'|\le C/A^2$.
Multiplication by $e^{-2v}$ and integration give error
$C(1+tU)/A$ both for $q_1-e^{-2v}$ and its derivative, with
weight $e^{v/(2A)}$. Beyond this prefix,
$A\ge4U+10$ and $e^{2r}-1\ge2r+2r^2$ give
$q_1\le2e^{-v}$ and $|q_1'|\le Ce^{-v/2}$, so the remaining
weighted error is $O(e^{-A^{1/4}/4})$.
Every higher theta index contributes relatively
$O(n^4e^{-A(n^2-1)e^{2v/A}})$; the scaled derivative has at most
the additional factor $(n^2-1)e^{2v/A}$. Both complete sums have
weighted integral $O(e^{-3A})$. Thus, including the exact
normalization by the whole $\Phi$,
\[
 \|q-e^{-2v}\|_{L^1(e^{v/(2A)})}
  +\|q'+2e^{-2v}\|_{L^1(e^{v/(2A)})}
                   \le C(1+tU)/A.                 \tag{PW13}
\]
The difference vanishes at $v=0$. Its Fourier--Laplace integral
has error at most the first norm, and a separate integration by
parts bounds it by the second norm divided by $|z/A|$.
The main integral is $(2\pm iz/A)^{-1}$, and
$|2\pm iz/A|\le2+|z/A|$. Combining the direct and integrated
bounds proves a uniform relative error, even when $|x|$ tends to
infinity. Restoring the factors $f_t(U)e^{\mp iUz}/A$ proves (PW11).

\emph{A cofinal phase choice for the entire crossover.}
Let
\[
 \mathcal D=D_1\zeta(s_1^*),\quad
 \mathcal B=\mathcal D-E,\quad
 \ell=\log D_1-\log f_t(U)+iUz+\operatorname{Log}(2A+iz)+Z(s_1^*).
                                                            \tag{PW14}
\]
The primary gamma log and the Euler log fix the branches;
$\Re(2A+iz)>0$ fixes the denominator log. They are holomorphic on
$x\ge X_0$, $|y|\le1/2$. On $A\le x\le2A$,
\[
 \begin{gathered}
 \ell'=-\pi/4-(i/2)\log(x/(2A))+\epsilon_1,
                    \quad|\epsilon_1|<0.002,\quad |\ell''|\le C,\\
 \Re\ell=A-\pi x/4+L/2+(tL/2)\log(x/(2A))+O(1).
 \end{gathered}                                      \tag{PW15}
\]
These are explicit derivative estimates, rather than derivatives of
an approximation error. Indeed $U=(L-\log\pi)/2$ cancels the large
gamma phase derivative. The exact endpoint satisfies
$\log f_t(U)=-A+tU^2+2\log A+U/2+\log4+O(A^{-1})$.
The gamma amplitude is
$-\pi x/4+(7/4+y/2)\log x+(t/4)\log^2(x/(2\pi))+O(1+t)$;
both $yL$ terms and the leading $tL^2/4$ terms cancel in their ratio.
Moreover $\Re s_1^*\ge16-o(1)\ge15$ on the crossover,
$(s_1^*)'=-i+O(t/x)$, and $(s_1^*)''=O(t/x^2)$.
Hence (PW10) bounds the added Euler derivatives as asserted.

For each $|y|\le1/2$ the amplitude has a unique zero
$x=X(U,y)$ in $[A,2A]$, since $\Re\ell'<-1/2$, and
\[
 X(U,y)=4A/\pi+(2/\pi)L
                +(2tL/\pi)\log(2/\pi)+O(1).       \tag{PW16}
\]
Set $\Theta(U,y)=\Im\ell(X(U,y)+iy)$ with the unwrapped branch
above. It is continuous in $U$ and
\[
 \Theta(U,0)=\frac{2A}{\pi}(1-\log(2/\pi))+O(L)\longrightarrow\infty.
                                                            \tag{PW17}
\]
The leading phase follows from
$\Im\log D_1=-(x/2)\log(x/(2\pi))+x/2+O(1+t\log x)$
and $\Im\log E=-Ux+O(1)$. Choose the first hitting points of the
successive odd multiples of $\pi$ for this continuous unbounded
phase. They form a strictly increasing cofinal sequence $U_j$.
Continuity excludes repeated hitting points for different targets,
and compact $U$-intervals have bounded phase.

If $\ell'=a+ib$ on the amplitude curve, the Cauchy--Riemann equations
give $X_y=b/a$ and $\Theta_y=(a^2+b^2)/a$. For all large $U$,
$|a+\pi/4|<0.002$ and $|b|<0.23$ there. Therefore
\[
 |\Theta_y|=|a|+b^2/|a|
            <0.789+0.23^2/0.78<1.                 \tag{PW18}
\]
The elementary bounds $3.14<\pi<22/7$ and $\log(\pi/2)<0.456$
suffice for these strict constants. At a chosen $U_j$, the phase
on the whole amplitude curve is within $1/2$ of its central odd
multiple of $\pi$. Whenever $|\Re\ell|\le1$, the horizontal
offset from that curve is at most $2$, and $|\Im\ell'|\le1/2$
on this interval adds at most $1$ to the phase variation.
Thus $\mathcal D/E$ lies in the left half-plane:
the total variation is below $3/2<\pi/2$. Outside this amplitude
band its modulus is separated from $1$. We obtain
\[
 |\mathcal B|\ge\tfrac14(|\mathcal D|+|E|).         \tag{PW19}
\]
This holds on the entire high strip. On $[X_0,A]$ the log ratio
is at least $cA$, and on $[2A,\infty)$ it is at most $-cA$ and
decreases further. The explicit gamma log proves these bounds:
$\log x+t\log^2x=o(x)$ uniformly beyond $2A$, whereas all
nonlinear costs below $A$ are $O(L^2)=o(A)$.
The Euler log is bounded throughout. Finally
$\log X_0\le L/4+\log(2\pi)$, so $X_0/A\to0$.
There is no unaccounted intermediate range.

\emph{A uniform carrier and every actual zero.}
The nonvanishing carrier (PW19) satisfies
\[
 \begin{gathered}
 c/t\le\Re(i\mathcal B'/\mathcal B)\le CL,\qquad
 |\Re(\mathcal B'/\mathcal B)|\le C,\qquad
 |(\log\mathcal B)''|\le C,\\
 |J_{U,t}-(\mathcal B+\mathcal B^\#)|
       \le\varepsilon_U(|\mathcal B|+|\mathcal B^\#|),\qquad
 \varepsilon_U\le C\{t+(1+tU)/A\}\longrightarrow0.
 \end{gathered}                                      \tag{PW20}
\]
On $[X_0,A]$ the logarithmic derivative is
$-\pi/4-(i/2)\log(x/(2\pi))+Z'(s_1^*)(s_1^*)'+o(1)$.
Both Euler derivatives are bounded by absolute summation.
On $[A,2A]$ it is
$E'/E+\ell'e^\ell/(e^\ell-1)=-iU+O(1)$, with bounded
second log derivative by (PW19). Beyond $2A$ it is $-iU+o(1)$;
the exponential ratio absorbs every logarithmic gamma derivative
even on the unbounded $x$-axis. This proves the first line of
(PW20). The second line follows from (PW8), (PW11), (PW19), and
the lower bound for the Euler factor.

Put $M=\mathcal B+\mathcal B^\#$. Vertical integration of the
positive first bound in (PW20) gives
$|\mathcal B(x+iy)|>|\mathcal B^\#(x+iy)|$ when $y>0$.
Thus $M$ has no nonreal high-strip zero. On the real axis write
$\mathcal B=|\mathcal B|e^{i\theta(x)}$ and
$m(x)=-\theta'(x)$; then $c/t\le m\le CL$ and $|m'|\le C$.
The phase tends to $-\infty$, since eventually the endpoint dominates.
Every model zero $x_n$ satisfies $\theta(x_n)=\pi/2\pmod\pi$
and is simple.

Let $z=x+iy$ be an actual window zero with $x\ge2X_0$.
The relative function error in (PW20) implies
\[
 \frac{\big||\mathcal B|-|\mathcal B^\#|\big|}
      {|\mathcal B|+|\mathcal B^\#|}\le\varepsilon_U.
\]
The vertical log ratio has magnitude at least
$2|y|(m(x)-C/2)$. Hence $|y|\le C\varepsilon_U/m(x)$.
Taylor expansion on that segment next gives
$|\cos\theta(x)|\le C\varepsilon_U$.
The nearest real model zero is at distance
$C\varepsilon_U/m(x)$, and $m(x_n)/m(x)\to1$ uniformly.
Every actual high-strip zero is therefore in a disk
$C\varepsilon_U/m(x_n)$ about some model zero.

Choose a fixed small $r>0$. The larger disks of radius $r/m(x_n)$
are disjoint for all large $U$, by the phase gap $\pi$ and $|m'|\le C$.
Write $\mathcal B(x_n)=ia_n$, with $a_n\in\R\setminus\{0\}$,
and $\mathcal B'/\mathcal B(x_n)=b_n-im_n$. Uniform log Taylor
expansion on $|w|=r/m_n$ gives
\[
 M(x_n+w)=2a_ne^{b_nw}\sin(m_nw)+O(|a_n|/m_n^2).   \tag{PW21}
\]
Here $|\sin(m_nw)|\ge r/2$. The model has exactly one zero in
the disk, and (PW20) gives actual error $O(\varepsilon_U|a_n|)$
on its boundary. Both errors are strictly smaller than the fixed
sinusoidal minimum, uniformly for every high model zero.
Rouch\'e gives exactly one window zero including multiplicity
in each disk. Reality of the function and conjugation symmetry
of the disk make this zero real and simple.
The prior localization accounts for every actual high-strip zero;
the possible leftmost disk stays above $X_0$ because $x\ge2X_0$.
Evenness treats negative $x$. This proves (PW3).

For convergence, $e^{tu^2}\Phi(u)$ converges to $\Phi$ in every
weighted norm $L^1(e^{Y|u|})$ by double-exponential decay and
$t\to0$. The omitted physical tails tend to zero in these same
norms as $U\to\infty$. The Fourier--Laplace bound is uniform in
all real parts. This proves the complete-strip convergence.
\end{proof}

\begin{proposition}[Physical-tail control and multiplicity-preserving low counts]
\label{prop:physical-tail-low-count}
For $U\ge2$, $0\le t\le1$, $A=\pi e^{2U}$, and $|\Im z|\le1/2$,
\[
 |F_t-J_{U,t}|\le128A^{3/2}e^{-A+tU^2},\qquad
 \frac{|\partial_z^m(F_t-J_{U,t})|}{m!}
       \le128A^2e^{-A+tU^2}\quad(m=0,1,\ldots).     \tag{PW22}
\]
If a bounded domain has a rectifiable zero-free boundary $C$ in this
strip and $\min_C|F_t|>128A^{3/2}e^{-A+tU^2}$, then $F_t$ and
$J_{U,t}$ have exactly the same enclosed zero count, including
every multiplicity. In particular every fixed hypothetical nonreal
zero of $\Xi$ is retained locally by both members of a sequence in
Theorem~\ref{thm:physical-windows} for all sufficiently large indices.
\end{proposition}
\begin{proof}
The scaled physical density above satisfies $q(v)\le4e^{-v}$ for
$U\ge2$, and the full theta sum gives
$f_t(U)\le8A^2e^{U/2-A+tU^2}$. Both tail components are bounded by
\[
 \frac{2f_t(U)e^{U/2}}A
       \int_0^\infty q(v)e^{v/(2A)}\dd v
       \le128A^{3/2}e^{-A+tU^2}.
\]
The integral is at most $8$ and $e^U=(A/\pi)^{1/2}$.
For every derivative order simultaneously, $|u|^m/m!\le e^{|u|}$
changes the weight $1/2$ to $3/2$; the scaled integral remains
at most $8$, and $e^{2U}=A/\pi$ gives the second bound.
Double-exponential decay justifies all differentiations.
Rouch\'e proves the asserted count identity. For a fixed nonreal
$\Xi$ zero, choose a small zero-free disk boundary separated from
the real axis. Uniform $F_t\to\Xi$ makes its boundary minimum
positive eventually, while the right side of (PW22) tends to zero.
Both local counts are then the original multiplicity.
\end{proof}

The high theorem does not exclude these low nonreal clusters. Its
cutoff still tends to infinity. The exact remaining obligation for
this construction is to prove absence of every nonreal window zero
below (PW3) along a cofinal sequence, or to prove a transfer that
forces any hypothetical nonreal $\Xi$ cluster into the proved high
range before it reaches the real axis. Gaussian convolution of the
physical window multiplies its transform by $e^{-\tau z^2}$ and
preserves every zero; it cannot supply that missing low exclusion.

\section{Positive modular Mehler approximants without an endpoint cusp}
\label{sec:mehler}
We prove a different theta approximation whose physical kernel is
analytic at the origin. It retains exact Jacobi modularity and converges
uniformly on the whole closed critical strip. Its transforms are
holomorphic on a larger strip, rather than entire; this domain distinction
is essential.

Put $L_{\rm th}=D_u^2-1/4$ and
\[
 \Theta(u)=\tfrac12e^{u/2}\theta(e^{2u}),\qquad
 \mathcal R(u)=\Theta(u)-\cosh(u/2).
\]
The already proved Jacobi identity makes $\Theta$ and $\mathcal R$ even,
and differentiating the normally convergent theta series gives
$L_{\rm th}\mathcal R=\Phi$. On $u\ge0$,
\[
 \mathcal R(u)=e^{u/2}\sum_{n\ge1}e^{-\pi n^2e^{2u}}
                                      -\tfrac12e^{-u/2}.
\]
Consequently $\mathcal R$ and its first derivative tend to zero at both
ends. Since $L_{\rm th}e^{-|u|/2}=-\delta_0$ distributionally, the
decaying Green solution is exactly
\[
 \mathcal R(u)=-\int_\R e^{-|u-v|/2}\Phi(v)\dd v<0. \tag{MG}
\]
Indeed the convolution exists absolutely, solves the same equation, and
its difference from $\mathcal R$ is a linear combination of
$e^{u/2},e^{-u/2}$ vanishing at both ends. Both coefficients are zero.
In particular the existing Fourier identity and
$\xi(0)=\xi(1)=1/2$ give
\[
 \int_\R e^{v/2}\Phi(v)\dd v
 =\int_\R\cosh(v/2)\Phi(v)\dd v=\tfrac12.
\]

For $\alpha>1$ define
\[
 H=\tfrac12\log\alpha,\qquad
 h_\alpha(u)=\tfrac12\log\frac{\alpha e^{2u}+1}{e^{2u}+\alpha},
\]
\[
 A_\alpha(u)=\sqrt\alpha
       [ (\alpha+1)^2+4\alpha\sinh^2u]^{-1/4},\qquad
 \mathcal R_\alpha=A_\alpha\mathcal R(h_\alpha),\qquad
 \Phi_\alpha=L_{\rm th}\mathcal R_\alpha.
\]
All real branches are positive where appropriate. The continuous
definition at $\alpha=1$ is
$h_1=0$, $A_1=2^{-1/2}\operatorname{sech}^{1/2}u$.

\begin{theorem}[Positive modular approximants and full-strip convergence]
\label{thm:mehler-approximation}
For every $\alpha\ge1$, $\Phi_\alpha$ is strictly positive, even, and
real analytic. It extends holomorphically to $|\Im u|<\pi/4$.
Its Fourier transform
$\mathcal M_\alpha(z)=\int_\R\Phi_\alpha(u)e^{-izu}\dd u$
is absolutely convergent and holomorphic on $|\Im z|<5/2$.
For each fixed $\alpha>1$ there is $C_\alpha>0$ such that
\[
 \Phi_\alpha(u)=C_\alpha e^{-5u/2}
                              +O_\alpha(e^{-9u/2})\quad(u\to+\infty).
\]
The exact boundary moment and the global convergence are
\[
 \int_\R\cosh(u/2)\Phi_\alpha(u)\dd u
 =\tfrac12-\sqrt\alpha\sum_{n\ge1}e^{-\pi\alpha n^2},
\]
\[
 \int_\R e^{|u|/2}|\Phi_\alpha(u)-\Phi(u)|\dd u\longrightarrow0,
 \qquad
 \sup_{|\Im z|\le1/2}|\mathcal M_\alpha(z)-\Xi(z)|
                                                    \longrightarrow0.
 \tag{MCV}
\]
The supremum includes every real height. Moreover
$\mathcal M_1$ is zero-free on $|\Im z|<5/2$.
\end{theorem}
\begin{proof}
The warp $h=h_\alpha$ is odd, $A=A_\alpha$ is even,
$h'>0$ and $h(\R)=(-H,H)$. Its two exact branches are
\[
 W_+(u)=Ae^{h/2}=\frac{\alpha^{1/4}}{\sqrt2}
                  \operatorname{sech}^{1/2}(u-H),\qquad
 W_-(u)=Ae^{-h/2}=\frac{\alpha^{1/4}}{\sqrt2}
                  \operatorname{sech}^{1/2}(u+H).
\]
For each fixed $v$, apply $L_{\rm th}$ to
$W_v(u)=A(u)e^{-|h(u)-v|/2}$ in (MG).
Away from $h(u)=v$,
$L_{\rm th}W_\pm=-(3/4)\operatorname{sech}^2(u\mp H)W_\pm$.
At its crossing the derivative jump of $W_v$ is $-Ah'$.
After integration in $v$ this produces $A(h')^2\Phi(h)$:
the second $h'$ is the delta-function Jacobian. Thus exactly
\[
\begin{split}
 \Phi_\alpha(u)={}&A(h')^2\Phi(h(u))\\
 &+\tfrac34 W_-(u)\operatorname{sech}^2(u+H)
                  \int_{-\infty}^{h(u)}e^{v/2}\Phi(v)\dd v\\
 &+\tfrac34 W_+(u)\operatorname{sech}^2(u-H)
                  \int_{h(u)}^\infty e^{-v/2}\Phi(v)\dd v.
\end{split} \tag{MK}
\]
All terms are positive. The double exponential bound for $\Phi$ and
its derivatives justifies the weighted integrals and distributional
calculation. At $\alpha=1$ direct differentiation instead gives
\[
 \Phi_1(u)=-\frac{3\mathcal R(0)}{4\sqrt2}
                             \operatorname{sech}^{5/2}u>0.
\]

For complex $u$ with $|\Im u|<\pi/4$, put $v=e^{2u}$.
Then $\Re v>0$, and
$w=(\alpha v+1)/(v+\alpha)$ also has positive real part: after
rationalization its real numerator is
$\alpha|v|^2+(\alpha^2+1)\Re v+\alpha>0$.
Use the quarter-root expression
$A=\sqrt\alpha\,v^{1/4}[(\alpha v+1)(v+\alpha)]^{-1/4}$
with the separate right-half-plane branches. Normal theta convergence
there proves the claimed physical holomorphy.
As $u\to+\infty$, $h=H+O_\alpha(e^{-2u})$,
$h'=O_\alpha(e^{-2u})$ and
$A=\alpha^{1/4}e^{-u/2}(1+O_\alpha(e^{-2u}))$.
Formula (MK) gives
\[
 C_\alpha=3\left[
 \alpha^{-1}\int_{-\infty}^{H}e^{v/2}\Phi(v)\dd v
 +\alpha^{3/2}\int_H^\infty e^{-v/2}\Phi(v)\dd v\right]>0.
\]
The main term in (MK) is $O_\alpha(e^{-9u/2})$.
The same expansions hold after real differentiation. Evenness and
these tails prove the stated Fourier-integral domain. Its Laplace
integral diverges at $z=\pm5i/2$; no entire-integral assertion is made.

The first tail of $\mathcal R_\alpha$ is
$\alpha^{1/4}\mathcal R(H)e^{-u/2}+O_\alpha(e^{-5u/2})$.
The Green identity on $[-T,T]$ with test $e^{u/2}$ reads
\[
 \int_{-T}^{T}e^{u/2}L_{\rm th}\mathcal R_\alpha\dd u
 =\left[e^{u/2}\mathcal R_\alpha'
               -\tfrac12e^{u/2}\mathcal R_\alpha\right]_{-T}^{T}.
\]
The lower boundary tends to zero; the upper tends to
$-\alpha^{1/4}\mathcal R(H)$.
The theta expression for $\mathcal R(H)$ gives the claimed moment.
The warp and all its real derivatives converge locally to the identity
and to $1$ as $\alpha\to\infty$, so $\Phi_\alpha\to\Phi$ pointwise.
Put $f_\alpha=\cosh(u/2)\Phi_\alpha$ and $f=\cosh(u/2)\Phi$.
Both are nonnegative, their integrals converge by the moment identity,
and $\min(f_\alpha,f)\to f$ is dominated by the integrable $f$.
Therefore
\[
 \int|f_\alpha-f|=\int f_\alpha+\int f-2\int\min(f_\alpha,f)
                                                     \longrightarrow0.
\]
Since $e^{|u|/2}\le2\cosh(u/2)$, this proves both limits in (MCV).

Finally, for $a>0$ and $|\Im z|<a$, substitution $v=e^{2u}$ gives
\[
 G_a(z):=\int_\R\operatorname{sech}^{a}u\,e^{-izu}\dd u
 =\frac{2^{a-1}}{\Gamma(a)}
       \Gamma((a+iz)/2)\Gamma((a-iz)/2).
\]
The beta integral here is absolutely convergent because both arguments
have positive real parts; the beta--gamma identity follows by the
change of variables $x+y=t$, $x/t=r$ in the two absolutely convergent
Euler gamma integrals. Gamma has no zeros and no poles with positive
real argument \cite[\S5.2(i)]{DLMFGamma}. Hence $G_a$ is zero-free
on that open strip, which proves the assertion for $\mathcal M_1$.
\end{proof}

The modularity is exact at finite $\alpha$. To see its origin, use the
separate convention
$\mathcal F f(\eta)=\int_\R f(x)e^{-2\pi ix\eta}\dd x$.
Poisson summation makes
$\mathcal D=\sum_{n\in\Z}\delta_n-\delta_0-1$ Fourier-selfdual.
The Mehler kernel
\[
 K_\alpha(x,y)=\sqrt\alpha
 e^{-\pi[\alpha(x^2+y^2)-2\sqrt{\alpha^2-1}\,xy]}
\]
commutes with $\mathcal F$: Gaussian integration in either coordinate
gives the same transformed kernel, using
$\alpha-(\alpha^2-1)/\alpha=1/\alpha$.
Its application to $\mathcal D$ is the signed selfdual function
\[
 f_\alpha(x)=\sqrt\alpha\sum_{n\ne0}
 e^{-\pi[\alpha(x^2+n^2)-2\sqrt{\alpha^2-1}\,nx]}
                                  -e^{-\pi x^2/\alpha}.
\]
For fixed $\alpha$ this is Schwartz. Completing the square in $n$
leaves a Gaussian $e^{-\pi x^2/\alpha}$ times a uniformly bounded
periodized Gaussian; derivatives only add polynomial factors.
Thus all subsequent Gaussian integrations are justified and give
\[
 \vartheta_\alpha(v)=\int_\R f_\alpha(x)e^{-\pi vx^2}\dd x
 =\sqrt\alpha\left[
 \frac{\theta((\alpha v+1)/(v+\alpha))-1}{\sqrt{v+\alpha}}
                           -\frac1{\sqrt{\alpha v+1}}\right],
\]
\[
 \vartheta_\alpha(1/v)=v^{1/2}\vartheta_\alpha(v),\qquad
 \tfrac12v^{1/4}\vartheta_\alpha(v)
                      =\mathcal R_\alpha(\tfrac12\log v)\quad(v>0).
\]
The first identity follows from Fourier selfduality and the Gaussian
transform, or directly from Jacobi. Signed $f_\alpha$ is not used as
a positive measure; positivity of $\Phi_\alpha$ has the separate proof
(MK).

\begin{proposition}[Entire factor, composition law, and exact evolution]
\label{prop:mehler-entire}
Set $k=(\alpha-1)/(\alpha+1)$, $c(k)=\sqrt{(1+k)/2}$, and
\[
 h_k(u)=\operatorname{atanh}(k\tanh u),\qquad
 B_k(u)=\sqrt{\frac{\cosh h_k(u)}{\cosh u}},\qquad
 T_k\mathcal R=B_k\mathcal R(h_k).
\]
Then $A_\alpha=c(k)B_k$ and $T_kT_l=T_{kl}$ for $0<k,l\le1$.
With $r=-\log k$, on every real compact the generator is
\[
 \partial_rT_{e^{-r}}\mathcal R=\mathcal G(T_{e^{-r}}\mathcal R),
 \qquad
 \mathcal G=-\tfrac12\sinh(2u)D_u-\tfrac12\sinh^2u.
\]
Define the even analytic disk function
\[
 Q(x)=(1-x)^{-1/2}\sum_{n\ge1}e^{-\pi n^2(1+x)/(1-x)}
                           -\tfrac12(1+x)^{-1/2}
       =\sum_{j\ge0}q_{2j}x^{2j},\quad |x|<1.
\]
For $a>0$ let $p_0^{(a)}=1$, $p_1^{(a)}=z$, and
\[
 p_{n+1}^{(a)}=zp_n^{(a)}-n(a+n-1)p_{n-1}^{(a)}.
\]
Each $p_n^{(a)}$ has simple real zeros, strictly interlacing those
of $p_{n+1}^{(a)}$. For every $0\le k<1$ the series
\[
 E_k(z)=\sum_{j\ge0}q_{2j}k^{2j}(-1)^j
                  \frac{p_{2j}^{(1/2)}(z)}{(1/2)_{2j}}
\]
converges normally on every complex compact and is entire and even.
It gives the exact factorization
\[
 \mathcal M_\alpha(z)=-\tfrac34c(k)G_{5/2}(z)E_k(z)
                      \qquad(|\Im z|<5/2), \tag{MFE}
\]
and the exact entire-function evolution
\[
 k\partial_kE_k(z)=\tfrac14\left[
 (\tfrac12-iz)E_k(z+2i)+(\tfrac12+iz)E_k(z-2i)-E_k(z)\right]
                   \quad(0<k<1,\ z\in\C). \tag{MD}
\]
The scalar $c(k)$ is not part of the composition semigroup.
\end{proposition}
\begin{proof}
The identity $\tanh h_k=k\tanh u$ proves composition of the warps;
the factors $\sqrt{\cosh h_k/\cosh u}$ telescope. Differentiating
at $r=0$ gives $\partial_rh=-\sinh u\cosh u$ and
$\partial_r\log B=-\sinh^2u/2$, proving the generator formula.
No globally bounded-generator assertion is needed.
With $x=\tanh u$, the displayed theta formula for $\mathcal R$ gives
$Q(x)=\sqrt{\cosh u}\mathcal R(u)$ and
\[
 \mathcal R_\alpha(u)
       =c(k)\operatorname{sech}^{1/2}u\,Q(k\tanh u).
\]
The map $(1+x)/(1-x)$ has positive real part on the unit disk,
so $Q$ is analytic there by normal theta convergence. Jacobi makes
it even.

For $I_n=\int_\R\operatorname{sech}^{a}u\tanh^n u\,e^{-izu}\dd u$,
absolute convergence on $|\Im z|<a$ permits integration by parts.
Differentiation of the physical factor gives
$izI_n=nI_{n-1}-(a+n)I_{n+1}$. Induction from $I_0=G_a$ yields
\[
 I_n=(-i)^nG_a\,p_n^{(a)}/(a)_n.
\]
The polynomial recurrence is the characteristic-determinant recurrence
for the symmetric Jacobi matrix with zero diagonal and positive
off-diagonal entries $\sqrt{j(a+j-1)}$, $1\le j<n$.
Its eigenvalues are real. The eigenvector recurrence determines every
entry from its first, which cannot vanish, so each eigenspace is
one-dimensional and every eigenvalue is simple. Principal-matrix
interlacing follows from the finite-dimensional min--max principle.
Common zeros of consecutive characteristic polynomials would propagate
backwards to $p_0=0$, and hence interlacing is strict.

The same recurrence gives the analytic generating function
\[
 \sum_{n\ge0}p_n^{(a)}(z)\frac{s^n}{n!}
            =(1+s^2)^{-a/2}e^{z\arctan s},\qquad |s|<1,
\]
by the equation $(1+s^2)G_s=(z-as)G$, $G(0)=1$.
Cauchy on $|s|=b<1$ bounds $|p_n(z)/n!|$ by $C_{b,K}b^{-n}$
for $z$ in a fixed compact $K$. Also
$n!/(1/2)_n\le2n+1$: its consecutive ratio is
$(2n+2)/(2n+1)$, which proves this bound by induction.
Disk Cauchy gives $|q_n|\le M_\rho\rho^{-n}$ for $\rho<1$.
Choose $b,\rho$ with $k<b\rho<1$; the entire series is then normally
dominated by $C(2n+1)(k/(b\rho))^n$. The domination is uniform on
closed $k$-intervals inside $(0,1)$ and allows parameter derivatives.

On $|\Im z|<1/2$, termwise physical integration of the $Q$ series
is legitimate because $\sum|q_n|k^n<\infty$ and the factor
$\operatorname{sech}^{1/2}u\,e^{|\Im z||u|}$ is integrable.
The Fourier multiplier of $L_{\rm th}$ is $-(z^2+1/4)$.
Gamma recurrence gives
$(z^2+1/4)G_{1/2}=(3/4)G_{5/2}$, proving (MFE) first there.
Analytic continuation proves it on $|\Im z|<5/2$; outside that
strip it describes a meromorphic continuation, not an absolutely
convergent physical integral.

Finally the generating function, together with
$e^{2i\arctan s}=(1+is)/(1-is)$, gives
\[
 (a-iz)p_n(z+2i)+(a+iz)p_n(z-2i)-2ap_n(z)=4np_n(z).
\]
Indeed its multiplier on the generating function is
$4s(z-as)/(1+s^2)=4sG_s/G$. Comparing coefficients proves the
identity. Normal convergence on shifted complex compacts and termwise
$k$ differentiation now prove (MD). These shifts use entire $E_k$
values, without extending a convergent $\mathcal R$ integral beyond
its domain.
\end{proof}

\begin{theorem}[An effective error on the entire closed critical strip]
\label{thm:mehler-effective-error}
For every real $\alpha\ge2$ the preceding exact theta approximants obey
\[
 \int_\R e^{|u|/2}|\Phi_\alpha(u)-\Phi(u)|\dd u
 \le C_*/\alpha<1500/\alpha,\qquad
 \sup_{|\Im z|\le1/2}|\mathcal M_\alpha(z)-\Xi(z)|<1500/\alpha.
 \tag{MEB}
\]
Here $C_*$ is the explicit positive elementary constant below.
The supremum includes all real heights.
\end{theorem}
\begin{proof}
Put $\delta=\alpha^{-1}\le1/2$, $v=e^{2u}\ge1$ for $u\ge0$,
$S(w)=\sum_{n\ge1}e^{-\pi n^2w}$, and $d=1-\delta^2$.
The exact branches of $\mathcal R_\alpha$ give
\[
 w=\frac{v+\delta}{1+\delta v},\qquad
 W=v^{1/4}(1+\delta v)^{-1/2},\qquad
 \mathcal R_\alpha=WS(w)-\tfrac12e^{-u/2}
                                   (1+\delta e^{-2u})^{-1/2}.
\]
Applying $L_{\rm th}$ gives $\Phi_\alpha=F_\delta+P_\delta$ with
\[
 F_\delta=W[C_2S''(w)+C_1S'(w)+C_0S(w)],\qquad
 P_\delta=\tfrac32\delta e^{-5u/2}
                              (1+\delta e^{-2u})^{-5/2},
\]
\[
 C_2=\frac{4v^2d^2}{(1+\delta v)^4},\quad
 C_1=\frac{6vd(1-\delta v)}{(1+\delta v)^3},\quad
 C_0=-\frac{3\delta v}{(1+\delta v)^2}. \tag{MEC}
\]
For example $w'=2vd/(1+\delta v)^2$ and
$2(W'/W)w'+w''=C_1$; these verify the potentially mixed derivative
coefficient. At $\delta=0$, $F_0=\Phi$ and $P_0=0$.

For $j=0,1,2,3$ and real $w\ge1$, the inequality
$n^2-1\ge3(n-1)$ gives $|S^{(j)}(w)|\le B_je^{-\pi w}$, where
\[
 B_j=\pi^j\sum_{n\ge1}n^{2j}e^{-3\pi(n-1)}.
\]
Writing $q=e^{-3\pi}$, repeated differentiation of $q/(1-q)$ yields
\[
 B_0=(1-q)^{-1},\quad B_1=\pi(1+q)(1-q)^{-3},
\]
\[
 B_2=\pi^2(1+11q+11q^2+q^3)(1-q)^{-5},
\]
\[
 B_3=\pi^3(1+57q+302q^2+302q^3+57q^4+q^5)(1-q)^{-7}.
\]
Set $V=1/\delta$. On $1\le v\le V$, for every
$0\le\eta\le\delta$, we have $\eta v\le1$, $w_\eta\ge v/2$,
$W_\eta\le v^{1/4}$, $|\partial_\eta w_\eta|\le v^2$ and
$|\partial_\eta W_\eta|\le vW_\eta/2$.
Directly differentiating (MEC) gives
\[
 |C_2|\le4v^2,\quad |\partial_\eta C_2|\le24v^3,
 \quad |C_1|\le6v,\quad |\partial_\eta C_1|\le30v^2,
 \quad |C_0|\le3,\quad |\partial_\eta C_0|\le3v.
\]
Consequently
\[
 |\partial_\eta F_\eta|
 \le v^{1/4}e^{-\pi v/2}
       [4B_3v^4+32B_2v^3+36B_1v^2+(9/2)B_0v].
\]
The mean-value integral in $\eta$ and
$e^{u/2}\dd u=v^{-3/4}\dd v/2$ now prove
\[
 \int_{1\le v\le V}e^{u/2}|F_\delta-F_0|\dd u\le\delta C_I,
 \qquad
 C_I=24(2/\pi)^5[2B_3+16B_2+18B_1+(9/4)B_0].
\]
Indeed the resulting powers $7/2,5/2,3/2,1/2$ are at most $v^4$
for $v\ge1$, and $\int_0^\infty v^4e^{-\pi v/2}\dd v=24(2/\pi)^5$.

For the complete tail $v\ge V$, we have $w\ge1/(2\delta)$ and
\[
 e^{u/2}W\dd u\le\tfrac12\delta^{-1/2}v^{-1}\dd v,
 \quad |C_2|\le4\delta^{-4}v^{-2},\quad
 |C_1|\le6\delta^{-2}v^{-1},\quad |C_0|\le3\delta^{-1}v^{-1}.
\]
Thus
\[
 \int_{v\ge V}e^{u/2}|F_\delta|\dd u
 \le e^{-\pi/(2\delta)}
       [B_2\delta^{-5/2}+3B_1\delta^{-3/2}+(3/2)B_0\delta^{-1/2}].
\]
The original theta tail also satisfies
\[
 \int_{v\ge V}e^{u/2}|F_0|\dd u
 \le C_0^*e^{-\pi/(2\delta)},\qquad
 C_0^*=16(2B_2+3B_1)/\pi^3.
\]
To check this, its two $v$-integrand powers $3/2,1/2$ are bounded
by $v^2$; split $e^{-\pi v}$ into two halves and integrate one half.
Differentiation gives the elementary maximum
\[
 \sup_{0<\delta\le1/2}\delta^{-7/2}e^{-\pi/(2\delta)}
     =S_*=(7/\pi)^{7/2}e^{-7/2}.
\]
It occurs at $\delta=\pi/7<1/2$. Since all smaller inverse powers
are dominated by $\delta^{-7/2}$, the combined tail is at most
$\delta C_TS_*$, where $C_T=B_2+3B_1+(3/2)B_0+C_0^*$.
Finally the actual positive branch in (MEC) gives
$\int_0^\infty e^{u/2}P_\delta\dd u\le3\delta/4$.
Doubling by evenness proves the weighted bound with
\[
 C_*=2C_I+2C_TS_*+\tfrac32.
\]
Outward-rounded evaluation of the displayed elementary constants
at 256 bits gives
\[
 1439.0019386987023213173153476<C_*<
 1439.0019386987023213173153477<1500.
\]
Only this final elementary constant uses interval arithmetic; the
estimates above hold for every real $\alpha\ge2$.
Since $|e^{-izu}|\le e^{|u|/2}$ on the closed strip, the transform
bound follows without a height truncation.
\end{proof}

\begin{theorem}[A nonreal zero of an actual positive modular approximant]
\label{thm:mehler-strip-obstruction}
At $k=177/200$, $\alpha=377/23$, the exact function
$\mathcal M_\alpha$ has exactly one zero, counted with multiplicity,
in the disk $|z-c|<r$, where
\[
 \begin{aligned}
 \Re c&=\frac{30561398818003719921715546903735782215}{2^{120}},\\
 \Im c&=\frac{597079250754110273230973026105322206}{2^{120}},
 \qquad r=2^{-60}.
 \end{aligned} \tag{MCert}
\]
It is simple, and the closed disk is contained in $0<\Im z<1/2$.
Therefore reality of all critical-strip zeros fails for the actual
theta approximants at this parameter. This zero is not a zero of
$\Xi$ or $\zeta$, and a favorable cofinal parameter sequence remains
unexcluded.
\end{theorem}
\begin{proof}
We supply both infinite tails used in the outward-rounded calculation.
Define $l_0(X)=1$, $l_1(X)=1/2-X$ and
\[
 (m+1)l_{m+1}(X)=(2m+1/2-X)l_m(X)-(m-1/2)l_{m-1}(X).
\]
Differentiation of
$(1-x)^{-1/2}e^{-Xx/(1-x)}$ proves that it equals
$\sum_{m\ge0}l_m(X)x^m$. The exact coefficients of the disk function
are therefore
\[
 q_m=\sum_{n\ge1}e^{-\pi n^2}l_m(2\pi n^2)
                    -\tfrac12(-1)^m(1/2)_m/m!.
\]
Odd coefficients vanish exactly by the already proved Jacobi evenness.
The certificate retains every even $m\le M=1500$ and every theta
term $n\le50$. On $|x|=1/2$ we have
$\Re[(1+x)/(1-x)]\ge1/3$ and $|1-x|^{-1/2}\le\sqrt2$.
Successive $n^2$ increments for $n\ge51$ are at least $103$.
Cauchy's estimate therefore bounds all omitted theta contributions by
\[
 |q_m-q_m^{[50]}|\le2^m\varepsilon_\theta,\qquad
 \varepsilon_\theta=
 \frac{\sqrt2e^{-\pi51^2/3}}{1-e^{-\pi103/3}}
                                      <1.723\,10^{-1183}.
\]
These errors are included in every coefficient interval; no discarded
theta term is assumed to vanish.

Choose the exact generating-circle radii $b=\rho=97/100$ and
$t=k/(b\rho)<1$. On $|x|=\rho$ the disk formula for $Q$ gives
\[
 |Q(x)|\le M_\rho=
 \frac{\sqrt{1+\rho}}{2(1-\rho)}+\frac1{2\sqrt{1-\rho}}.
\]
Here $\Re[(1+x)/(1-x)]\ge(1-\rho)/(1+\rho)$ and the decreasing
Gaussian has $\sum_{n\ge1}e^{-\pi a n^2}\le
\int_0^\infty e^{-\pi a v^2}\dd v=1/(2\sqrt a)$.
On $|s|=b$, elementary real and imaginary part formulas for
$\arctan s$ give
$|\Re\arctan s|\le\arctan b$ and
$|\Im\arctan s|\le\operatorname{atanh}b$; their maxima occur at
$s=\pm b$ and $s=\pm ib$, respectively.
Thus the generating function in Proposition~\ref{prop:mehler-entire}
and $n!/(1/2)_n\le2n+1$ give, on
$|\Re z|\le X$, $|\Im z|\le Y$,
\[
 \left|\frac{p_n^{(1/2)}(z)}{(1/2)_n}\right|
 \le(2n+1)C_b(X,Y)b^{-n},\quad
 C_b(X,Y)=(1-b^2)^{-1/4}
                  e^{X\arctan b+Y\operatorname{atanh}b}.
\]
Writing $E_M$ for the polynomial using the exact even $q_m$, $m\le M$,
the complete remaining entire-series tail satisfies
\[
 |E_k(z)-E_M(z)|\le T_M(X,Y),\quad
 T_M=M_\rho C_b(X,Y)t^{M+1}
       \left[\frac{2M+3}{1-t}+\frac{2t}{(1-t)^2}\right].
\]
This sums the safe majorant over all $n>M$, even and odd. Cauchy on
the unit disk about $c$ bounds the derivative tail there by
$T_M(|\Re c|+1,|\Im c|+1)$. On the radius two disk the entire series
has bound
\[
 C_F=M_\rho C_b(|\Re c|+2,|\Im c|+2)\frac{1+t}{(1-t)^2}.
\]
Every unit disk about a point in $|z-c|\le r$ stays inside that
radius two disk, so $|E_k''|\le2C_F$ on the small disk.

The finite computation uses
$r_n=p_n^{(1/2)}/(1/2)_n$, with
\[
 r_0=1,\quad r_1=2z,\quad
 r_{n+1}=(zr_n-nr_{n-1})/(n+1/2).
\]
Its differentiated recurrence computes $E_M'(c)$.
Evaluation of this recurrence in outward-rounded Arb/Acb
arithmetic, using python-flint 0.9.0 at 4096 bits and the exact
center and rational parameters above, gives
\[
 T_M(|\Re c|,|\Im c|)<4.042\,10^{-26},\quad
 T_M(|\Re c|+1,|\Im c|+1)<7.075\,10^{-25},
\]
\[
 C_F<1.127\,10^{15},\qquad |E_M'(c)|>1.202\,10^6,
\]
\[
 \begin{aligned}
 &|E_M(c)|+T_M(|\Re c|,|\Im c|)\\
 &\qquad{}+rT_M(|\Re c|+1,|\Im c|+1)+C_Fr^2\\
 &\qquad{}<8.476\,10^{-22}<1.042\,10^{-12}<r|E_M'(c)|.
 \end{aligned}
\]
The actual ratio is below $8.129\,10^{-10}$. Taylor's formula and
these complete tails show on $|z-c|=r$ that $E_k$ differs from
$E_M'(c)(z-c)$ by less than the latter's modulus.
Rouch\'e's theorem gives exactly one zero counted with multiplicity;
in particular it is simple and the contour is zero-free.
Exact rational comparisons place the entire closed disk in
$0<\Im z<1/2$. On it the two gamma arguments in $G_{5/2}$ have
positive real part, so $G_{5/2}$ is holomorphic and zero-free.
The exact factorization (MFE) transfers the same count to
$\mathcal M_{377/23}$. This certifies a zero of the stated approximant,
without a claim about zeros of the limiting $\Xi$.
\end{proof}

\begin{proposition}[Exact Kummer--Dirichlet representation]
\label{prop:mehler-kummer}
For every fixed $\alpha\ge1$ and every $z\in\C$, put
$s=1/4-iz/2$ and $X_n=\pi(\alpha-\alpha^{-1})n^2$. Then
\[
 E_k(z)=\frac1{\sqrt2c(k)}
 \left[\alpha^s\sum_{n\ge1}e^{-\pi n^2/\alpha}
   {}_1F_1(s;1/2;-X_n)-\tfrac12\alpha^{1/2-s}\right]. \tag{MKD}
\]
The sum converges normally on every complex compact. Here
\[
 {}_1F_1(a;b;w)=\sum_{m\ge0}\frac{(a)_m w^m}{(b)_m m!},
 \qquad b=1/2.
\]
The displayed subtraction term is part of the exact identity.
\end{proposition}
\begin{proof}
First take $0<\Re s<1/2$, equivalently $|\Im z|<1/2$.
The Gaussian Mellin transform of the already defined selfdual
$f_\alpha$ gives, by $v=e^{2u}$,
\[
 \widehat{\mathcal R_\alpha}(z)
 =\tfrac14\int_0^\infty\vartheta_\alpha(v)v^{s-1}\dd v
 =\frac{\Gamma(s)\pi^{-s}}4
                     \int_\R f_\alpha(x)|x|^{-2s}\dd x .
\]
Fubini is absolute: $f_\alpha$ is Schwartz, $|x|^{-2\Re s}$ is
locally integrable, and $\Re s>0$ justifies the $v$ integral.
Writing $\beta=(\alpha^2-1)^{1/2}$, even Gaussian integration yields
\[
 \int_\R e^{-\pi\alpha x^2+2\pi\beta nx}|x|^{-2s}\dd x
 =(\pi\alpha)^{s-1/2}\Gamma(1/2-s)
       {}_1F_1(1/2-s;1/2;\pi\beta^2n^2/\alpha).
\]
This follows by expanding the even $\cosh$ factor and using Gaussian
moments and $(2m)!=4^m m!(1/2)_m$; its absolute Gaussian majorant
also justifies the complex-$s$ computation. Completing squares
in the original Gaussian sum justifies summing the moments.
Thus
\[
 \widehat{\mathcal R_\alpha}
 =\frac{\Gamma(s)\Gamma(1/2-s)}{4\sqrt\pi}
 \left[2\alpha^s\sum_{n\ge1}e^{-\pi\alpha n^2}
       {}_1F_1(1/2-s;1/2;X_n)-\alpha^{1/2-s}\right].
\]
Kummer's identity
\({}_1F_1(a;b;w)=e^w{}_1F_1(b-a;b;-w)\) follows directly from
$wF''+(b-w)F'-aF=0$ and uniqueness of its power-series solution
with $F(0)=1$, $F'(0)=a/b$. Dividing by
$cG_{1/2}=c\Gamma(s)\Gamma(1/2-s)/\sqrt{2\pi}$ proves (MKD)
on the initial strip.

We verify normal convergence on all complex compacts rather than
assuming a parameter-uniform hypergeometric asymptotic.
Initially $0<\Re s<1/2$, beta integration of the defining series gives
\[
 {}_1F_1(s;1/2;-X)
 =\frac{\sqrt\pi}{\Gamma(s)\Gamma(1/2-s)}
       \int_0^1e^{-Xt}t^{s-1}(1-t)^{-s-1/2}\dd t,\quad X\ge0.
\]
Use a fixed smooth partition into neighborhoods of $0$, $1$ and
the interior. At $0$ subtract the Taylor polynomial of
$e^{-Xt}(1-t)^{-s-1/2}$ through degree $m-1$; its residual is
integrable when $\Re s>-m$. The added $j$th power has only the
simple pole $s=-j$. For example, for a cutoff equal to $1$ on
$[0,1/4]$ and zero after $1/2$, its continued integral is
$(1/4)^{s+j}/(s+j)+\int_{1/4}^{1/2}\chi(t)t^{s+j-1}\dd t$.
At $1$ instead subtract the Taylor polynomial of
$e^{-Xt}t^{s-1}$ in $1-t$; its residual is integrable for
$\Re s<m+1/2$, and the added powers have only simple poles
$s=j+1/2$. The factors $1/\Gamma(s)$ and $1/\Gamma(1/2-s)$
cancel these poles exactly.
For a prescribed compact $K$, choose $m$ covering an open
neighborhood of its real parts. All coefficient and remainder
derivatives are bounded by $C_K(1+X)^m$, since at most $m$
derivatives of $e^{-Xt}$ occur and $e^{-Xt}\le1$ on $[0,1]$.
The canceled poles satisfy the same bound by the maximum principle
on fixed small surrounding circles. Consequently
\[
 \sup_{s\in K}|{}_1F_1(s;1/2;-X)|\le C_K(1+X)^{m_K}.
\]
For fixed $\alpha$, the summable majorant is therefore
$C_K(1+n^2)^{m_K}e^{-\pi n^2/\alpha}$. The identity theorem
extends (MKD) to all $z$.
At $\alpha=1$ it gives exactly
$E_0=\sum_{n\ge1}e^{-\pi n^2}-1/2=\mathcal R(0)$, checking
the normalization against the Gamma starting member.
\end{proof}

\begin{theorem}[Joint parameter and height envelopes]
\label{thm:mehler-joint-envelope}
For every $\alpha>1$ the physical kernel is holomorphic on
$|\Im u|<\pi/4+\arctan(1/\alpha)$.
For all $\alpha\ge2$, $x\in\R$ and $|y|\le1/2$,
\[
 |\mathcal M_\alpha(x+iy)|
 \le2304\alpha^3e^{-[\pi/4+\arctan(1/(2\alpha))]|x|}. \tag{MJH}
\]
There is an explicitly defined absolute constant $C_\sharp$ such that
\[
 |\mathcal M_\alpha(x+iy)-\Xi(x+iy)|
 \le \frac{C_\sharp}{\alpha}(1+|x|)^{9/2}e^{-\pi|x|/4}. \tag{MJD}
\]
Both bounds cover all heights and indicated parameters; neither
is a nonvanishing margin.
\end{theorem}
\begin{proof}
For $u=t+iv$,
$|\tanh u|^2=(\sinh^2t+\sin^2v)/(\sinh^2t+\cos^2v)$.
If $|v|<\arctan(1/k)=\pi/4+\arctan(1/\alpha)$, then
$|k\tanh u|<1$ and $\cosh u\ne0$. The disk representation proves
holomorphy and, writing $x_u=\tanh u$, gives exactly
\[
 \Phi_\alpha(u)=c\,\operatorname{sech}^{5/2}u
 [k^2(1-x_u^2)Q''(kx_u)-3kx_uQ'(kx_u)-3Q(kx_u)/4].
\]
The bracket and derivatives are bounded on fixed closed substrips,
with tails $e^{-5|\Re u|/2}$.
Set $\delta=1/\alpha\le1/2$, $\eta=\pi/4+\arctan(\delta/2)$.
The disk radius on $|\Im u|\le\eta$ is
$r=(2-\delta-\delta^2)/(2+\delta-\delta^2)$, $1-r\ge8\delta/9$.
For $\rho=(1+r)/2$, both gaps are at least $4\delta/9$ and the
full disk bound gives $M_\rho<3/\delta$. Cauchy bounds the bracket by
\[
 M_\rho[2(k^2+r^2)/(\rho-r)^2+3r/(\rho-r)+3/4]
 \le(381/16)M_\rho\delta^{-2}<24M_\rho\delta^{-2}.
\]
Also $|\operatorname{sech}(t+i\eta)|\le\sec\eta\operatorname{sech}t$,
$\sec\eta<2$ and
$\int_\R e^{|t|/2}\operatorname{sech}^{5/2}t\dd t\le2^{5/2}$.
The complete shifted norm is at most $2304\delta^{-3}$.
Shifting down for $x\ge0$ and up for $x\le0$ proves (MJH);
the vertical ends vanish by the displayed tails.

For (MJD), use the exact rational formula (MEC). For
$a=\Re w>0$ and $j=0,1,2,3$, Gaussian splitting into exponent
factors $1/2,1/4,1/4$ gives
\[
 |S^{(j)}(w)|\le D_ja^{-j-1/2}e^{-\pi a/2},
 \quad D_0=1,\quad D_j=(4j/e)^j\ (j>0).
\]
The middle factor absorbs $n^{2j}$ by its maximum; the last sum
is at most $\int_0^\infty e^{-\pi at^2/4}\dd t=a^{-1/2}$.
Put $K_j=4^{j+1/2}D_j$.
Choose $\epsilon\in(0,\pi/4)$, $\beta=\pi/4-\epsilon$ and
$c_0=\sin(2\epsilon)$. On $u=r-i\beta$, $r\ge0$, set
$v=e^{2u}$ and $q=|v|\ge1$, so $\Re v=c_0q$.
For $q\le1/\delta$, every $0\le\eta\le\delta$ satisfies
\[
 \Re w_\eta=
 \frac{c_0q(1+\eta^2)+\eta(1+q^2)}{|1+\eta v|^2}\ge c_0q/4.
\]
Moreover $|W_\eta|\le q^{1/4}$,
$|\partial_\eta W_\eta|\le q|W_\eta|/2$,
$|\partial_\eta w_\eta|\le2q^2$, and
\[
 (|C_2|,|\partial_\eta C_2|)\le(4q^2,24q^3),\quad
 (|C_1|,|\partial_\eta C_1|)\le(12q,48q^2),\quad
 (|C_0|,|\partial_\eta C_0|)\le(3,6q).
\]
Product differentiation gives
$|\partial_\eta F_\eta|\le q^{1/4}
[8q^4|S'''|+50q^3|S''|+60q^2|S'|+(15/2)q|S|]$.
After weighting and $\dd r=\dd q/(2q)$ all four powers cancel,
giving interior error at most $\delta C_{I,c}c_0^{-9/2}$ with
\[
 C_{I,c}=4[8K_3+50K_2+60K_1+(15/2)K_0]/\pi.
\]
For $q\ge1/\delta$, $\Re w_\delta\ge1/(4\delta)$,
$|W_\delta|\le\delta^{-1/2}q^{-1/4}$, and the three coefficient
bounds are $4\delta^{-4}q^{-2}$, $12\delta^{-2}q^{-1}$,
$3\delta^{-1}q^{-1}$. Integration gives error
\[
 [K_2+6K_1+(3/2)K_0]e^{-\pi/(8\delta)}
 \le\delta C_{T,c},\quad
 C_{T,c}=8[K_2+6K_1+(3/2)K_0]/(\pi e).
\]
The maximum of $\delta^{-1}e^{-\pi/(8\delta)}$ is at $\delta=\pi/8$.
For $F_0$, Gaussian integration with $q^{-1}\le\delta$ gives
$\delta C_{0,c}c_0^{-7/2}$, $C_{0,c}=(4D_2+6D_1)/\pi$.
The actual pole branch has shifted weighted norm at most $3\delta/4$.
Evenness and conjugation double the half-line bounds.
The full shifted error norm is at most
$\delta C_cc_0^{-9/2}$, $C_c=2(C_{I,c}+C_{T,c}+C_{0,c}+3/4)$.
The contour shift within $|\Im u|\le\beta$ is valid by the true
exponential approximant tails and double exponential theta majorant.
It gives $\delta C_cc_0^{-9/2}e^{-\beta|x|}$.
Choose $\epsilon=\pi/[8(1+|x|)]$; then $c_0\ge1/[2(1+|x|)]$
and $\epsilon|x|\le\pi/8$. This proves (MJD) with
$C_\sharp=2^{9/2}e^{\pi/8}C_c$.
\end{proof}

\begin{theorem}[An arithmetic tangent with a uniform quadratic remainder]
\label{thm:mehler-tangent}
There is an even real analytic kernel $J$, integrable with every
exponential weight of rate less than $5/2$, such that for $\alpha\ge2$
\[
 \int_\R e^{|u|/2}
       |\Phi_\alpha-\Phi-\alpha^{-1}J|\dd u\le C_2\alpha^{-2}.
 \tag{MTQ}
\]
Here $C_2$ is the explicit constant below. Its exact transform is
\[
 \widehat J(z)=-(z^2+1/4)
 \left[\frac{\Xi(z+2i)}{5-2iz}+\frac{\Xi(z-2i)}{5+2iz}\right],
                      \qquad|\Im z|<5/2. \tag{MTS}
\]
Thus the transform remainder is at most $C_2\alpha^{-2}$ on the
entire closed critical strip. On this strip, functional symmetry
expresses the two shifted terms through $\xi(s)$ with $\Re s\in[2,3]$.
\end{theorem}
\begin{proof}
In (MEC) put $\delta=\alpha^{-1}$ and differentiate at $\delta=0$.
For $u\ge0$, $v=e^{2u}$, the exact result is
\[
\begin{split}
 J(u)={}&v^{1/4}\bigl[
 4v^2(1-v^2)S'''(v)+(6v-24v^3)S''(v)\\
 &\hspace{25mm}-27v^2S'(v)-3vS(v)\bigr]
                         +(3/2)e^{-5u/2}.
\end{split}
\]
Indeed $W_\delta=-vW/2$, $w_\delta=1-v^2$,
$(C_2)_\delta=-16v^3$, $(C_1)_\delta=-24v^2$,
$(C_0)_\delta=-3v$ at zero. The generator also gives
$J=2L_{\rm th}\mathcal G\mathcal R-\Phi/2$, proving even real
analyticity at the origin without a reflected cusp.

We supply the full second-derivative majorant, including its moving
tail. Set $B_j=\pi^j\sum_{n\ge1}n^{2j}e^{-3\pi(n-1)}$, $0\le j\le4$.
For $0<\eta\le1/2$, $1\le v\le1/\eta$,
$|W_\eta|/W\le v/2$, $|W_{\eta\eta}|/W\le3v^2/4$,
$|w_\eta|\le2v^2$ and $|w_{\eta\eta}|\le4v^3$.
The coefficient bounds for derivative orders $0,1,2$ are
\[
 \begin{array}{c|ccc}
 &|C_j|&|(C_j)_\eta|&|(C_j)_{\eta\eta}|\\ \hline
 j=2&4v^2&24v^3&160v^4\\
 j=1&12v&54v^2&288v^3\\
 j=0&3&9v&30v^2
 \end{array}
\]
These follow by differentiating (MEC). For example, for
$f=(1-\eta v)(1+\eta v)^{-3}$,
$f'=-2v(2-\eta v)(1+\eta v)^{-4}$ and
$f''=6v^2(3-\eta v)(1+\eta v)^{-5}$.
The second product rule therefore gives
\[
 |\partial_\eta^2F_\eta|
 \le v^{1/4}e^{-\pi v/2}
 [16B_4v^6+168B_3v^5+487B_2v^4+405B_1v^3+(165/4)B_0v^2].
\]
The coefficients are the respective sums
$16$, $8+96+16+48$, $3+24+160+24+216+48+12$,
$9+54+288+6+36+12$, and $9/4+9+30$.
After weighting, all powers are at most $v^6$, giving half-line
bound $360(2/\pi)^7D_I$, where
$D_I=16B_4+168B_3+487B_2+405B_1+(165/4)B_0$.

For $v\ge1/\eta$, instead
$W\le\eta^{-1/2}v^{-1/4}$,
$|W_\eta|/W\le1/(2\eta)$, $|W_{\eta\eta}|/W\le3/(4\eta^2)$,
$|w_\eta|\le\eta^{-2}$ and $|w_{\eta\eta}|\le2\eta^{-3}$.
The coefficient bounds are
\[
 \begin{array}{c|ccc}
 &|C_j|&|(C_j)_\eta|&|(C_j)_{\eta\eta}|\\ \hline
 j=2&4\eta^{-4}v^{-2}&24\eta^{-5}v^{-2}&160\eta^{-6}v^{-2}\\
 j=1&6\eta^{-2}v^{-1}&30\eta^{-3}v^{-1}&168\eta^{-4}v^{-1}\\
 j=0&3\eta^{-1}v^{-1}&9\eta^{-2}v^{-1}&30\eta^{-3}v^{-1}
 \end{array}
\]
Using $w\ge1/(2\eta)$ and
$e^{u/2}W\dd u\le\eta^{-1/2}\dd v/(2v)$, the complete tail is at most
\[
 e^{-\pi/(2\eta)}
 [B_4\eta^{-13/2}+18B_3\eta^{-11/2}
 +(349/4)B_2\eta^{-9/2}+(459/4)B_1\eta^{-7/2}
                                      +(165/8)B_0\eta^{-5/2}].
\]
For clarity the coefficients before integration are
$4\eta^{-8}v^{-2}$,
$60\eta^{-7}v^{-2}+6\eta^{-6}v^{-1}$,
$187\eta^{-6}v^{-2}+81\eta^{-5}v^{-1}$,
$(459/2)\eta^{-4}v^{-1}$ and $(165/4)\eta^{-3}v^{-1}$.
All inverse powers are bounded by $\eta^{-13/2}$; its exponential
product has maximum $(13/\pi)^{13/2}e^{-13/2}$ at $\eta=\pi/13$.
The separate pole branch has
\[
 \partial_\eta^2P_\eta=(3/2)e^{-5u/2}
 [-5t(1+\eta t)^{-7/2}+(35/4)\eta t^2(1+\eta t)^{-9/2}],
 \quad t=e^{-2u}.
\]
Its weighted half-line norm is at most $225/64$.
The pointwise Taylor remainder
$\int_0^\delta(\delta-\eta)\partial_\eta^2\Phi_\eta\dd\eta$
is now dominated uniformly, so Tonelli applies. The factor
$\delta^2/2$ cancels the doubling of the half-lines. This proves
(MTQ) with
\[
\begin{split}
 C_2={}&360(2/\pi)^7D_I+
             D_T(13/\pi)^{13/2}e^{-13/2}+225/64,\\
 D_T={}&B_4+18B_3+(349/4)B_2+(459/4)B_1+(165/8)B_0 .
\end{split}
\]

To prove (MTS) without invalid shifts of $\widehat{\mathcal R}$,
split $\mathcal R=B_0^{\rm phys}+U$ with
$B_0^{\rm phys}(u)=-e^{-|u|/2}/2$ and
$U=e^{|u|/2}S(e^{2|u|})$. The latter is continuous,
double exponentially decreasing, and has piecewise integrable first
derivative. Its ordinary shifted transforms are entire and
\[
 \widehat U=(1/2-\Xi)/(z^2+1/4);
\]
the apparent poles are removable by $\Xi(\pm i/2)=1/2$.
Integration by parts has canceling boundary values at zero and gives
\[
 \widehat{\mathcal GU}
 =\frac{3-2iz}{8}\widehat U(z+2i)
  +\frac{3+2iz}{8}\widehat U(z-2i)+\tfrac14\widehat U(z).
\]
Independently
$\mathcal GB_0^{\rm phys}=-e^{-|u|/2}/8+e^{-5|u|/2}/8$.
Factor
$(z+2i)^2+1/4=-(3/2-iz)(5/2-iz)$ and its conjugate.
The shifted constants are $-1/[4(5-2iz)]-1/[4(5+2iz)]
=-5/[8(z^2+25/4)]$, canceling the second physical $B$ term;
the center constant cancels its first. Hence
\[
 \widehat{\mathcal G\mathcal R}
 =\frac{\Xi(z+2i)}{2(5-2iz)}
  +\frac{\Xi(z-2i)}{2(5+2iz)}-\frac{\Xi(z)}{4(z^2+1/4)}.
\]
Apply $2L_{\rm th}$ and subtract $\Phi/2$; the center term cancels.
This first proves (MTS) in $|\Im z|<1/2$. Both sides are holomorphic
on $|\Im z|<5/2$, so continuation proves the larger domain.
For $z=x+iy$ on the closed critical strip, functional symmetry
identifies the two shifted values with
$\xi(5/2+y-ix)$ and $\xi(5/2-y+ix)$, whose real parts lie in $[2,3]$.
\end{proof}

\begin{corollary}[A precise unheated multiplicity obstruction]
If $\gamma$ is a real Xi zero of order $m$, $a=\Xi^{(m)}(\gamma)/m!$
and $j=\widehat J(\gamma)\ne0$, then with
$r_\alpha=(|j|/(|a|\alpha))^{1/m}$,
\[
 \frac{\mathcal M_\alpha(\gamma+r_\alpha w)}{|j|/\alpha}
          \longrightarrow\operatorname{sgn}(a)w^m+\operatorname{sgn}(j)
\]
normally on compact $w$-sets. For $m\ge3$ this gives nonreal
critical-strip zeros of every sufficiently large unheated approximant.
For $m=2$ it does so exactly when $j/a>0$.
\end{corollary}
\begin{proof}
Taylor expansion, (MTQ), and holomorphy of $\widehat J$ give the limit.
Its roots are distinct, so disjoint small circles and Rouch\'e's
theorem give one simple zero near each. A nonreal limiting root
remains nonreal; a conjugation-invariant circle containing a real
root gives a real zero. The condition $j\ne0$ remains a condition:
nonvanishing of the two shifted Xi values does not establish
nonvanishing of their specified real-part combination.
\end{proof}

\begin{theorem}[An admissible entire heat flow and all-multiplicity exhaustion]
\label{thm:mehler-heated-exhaustion}
For $\alpha\ge2$, $k=(\alpha-1)/(\alpha+1)$, $\tau>0$, and
$H_\alpha=\tfrac12\log(4\alpha+3)$, define
\[
 \mathcal E_{\alpha,\tau}(z)=\frac1{\sqrt{4\pi\tau}}
       \int_\R E_k(z+iu)e^{-u^2/(4\tau)}\dd u,\qquad
 \mathcal B_{\alpha,\tau}(z)
       =-\tfrac34c(k)G_{5/2}(z)\mathcal E_{\alpha,\tau}(z).
 \tag{MHeat}
\]
The first function is entire, even, and of real type, and has the
normally convergent derivative expansion
$\mathcal E_{\alpha,\tau}=\sum_{j\ge0}(-\tau)^jE_k^{(2j)}/j!$.
The second is holomorphic on $|\Im z|<5/2$. There are finite positive
constants $C_1,C_2,C_3$, independent of $\alpha,\tau$, and height,
such that, whenever $8\tau H_\alpha\le1$,
\[
 \sup_{|\Im z|\le1/2}|\mathcal B_{\alpha,\tau}(z)-\Xi(z)|
 \le \frac{C_1}{\alpha}+C_2\tau+
                   C_3(\alpha+1)^{17/4}e^{-1/(16\tau)}. \tag{MHU}
\]
Consequently every diagonal $\tau(\alpha)\to0$ satisfying
$\tau(\alpha)\log\alpha\to0$ converges uniformly on the whole closed
critical strip. For the particular diagonal
$\tau_\alpha=(\log\alpha)^{-3}$, every fixed real Xi zero of arbitrary
finite multiplicity $m$ has exactly $m$ simple real zeros of
$\mathcal E_{\alpha,\tau_\alpha}$ in a fixed sufficiently small
neighborhood, for all sufficiently large $\alpha$.
\end{theorem}
\begin{proof}
Choose $b=\rho=(2\alpha+1)/(2\alpha+2)$ and $t=k/b^2$ in the
coefficient bound used for (MCert). Its generating function gives
\[
 |E_k(x+iy)|\le
 M_b(1-b^2)^{-1/4}\frac{1+t}{(1-t)^2}
               e^{|x|\arctan b+|y|\operatorname{atanh}b},
 \qquad
 M_b=\frac{\sqrt{1+b}}{2(1-b)}
                  +\frac1{2\sqrt{1-b}}.
\]
For $A=\alpha+1$, the exact expression
$1-t=(4A+1)/(2A-1)^2\ge1/A$ gives
$M_b\le3A$, $(1-b^2)^{-1/4}\le(2A)^{1/4}$, and
$\operatorname{atanh}b=H_\alpha$. Therefore
\[
 |E_k(x+iy)|\le8(\alpha+1)^4
                    e^{\pi|x|/4+H_\alpha|y|},\qquad
 |\mathcal E_{\alpha,\tau}(x+iy)|\le16(\alpha+1)^4
                    e^{\pi|x|/4+H_\alpha|y|+\tau H_\alpha^2}.
 \tag{MHG}
\]
The first bound proves normal convergence of the Gaussian integral
on every $z$-compact. Completing the square on the two $u$-halflines
proves the second. Evenness and complex conjugation, with $u\mapsto-u$
where needed, prove both symmetries. Cauchy's estimate with radius
$2j/(1+H_\alpha)$ on a fixed compact bounds the $2j$th derivative by
$C_{\alpha,K}(2j)![e(1+H_\alpha)/(2j)]^{2j}$.
Multiplying by $\tau^j/j!$ yields a summable sequence for every
fixed $\alpha,\tau$. The same Cauchy bounds sum the absolute Taylor
series after integration against the Gaussian; its even moments
are $(2j)!\tau^j/j!$ and its odd moments vanish. Thus its normally
convergent derivative expansion is the asserted one.

We first strengthen the weighted estimate in (MEB) to
\[
 \int_\R e^{3|u|/2}|\Phi_\alpha-\Phi|\dd u
           \le C_{3/2}/\alpha,\qquad
 C_{3/2}=2C_I+2C_{T,3/2}S_4+3,
 \tag{MHW}
\]
where $C_I,B_j,C_0^*$ are the explicit constants in that proof and
\[
 C_{T,3/2}=\tfrac43B_2+6B_1+3B_0+C_0^*,\qquad
 S_4=(8/\pi)^4e^{-4}.
\]
Indeed on $1\le v\le1/\delta$, the same mean-value derivative bound,
now multiplied by $e^{3u/2}\dd u$, has precisely the powers
$v^4,v^3,v^2,v$, so the same $C_I$ works. On $v\ge1/\delta$,
the complete rational bounds from (MEC) instead give
\[
 \int e^{3u/2}|F_\delta|\dd u\le e^{-\pi/(2\delta)}
          [\tfrac43B_2\delta^{-3}
                         +6B_1\delta^{-2}+3B_0\delta^{-1}].
\]
For $F_0$, its resulting powers $v^2,v$ again give
$C_0^*e^{-\pi/(2\delta)}$. The maximum of
$\delta^{-4}e^{-\pi/(2\delta)}$ is $S_4$ at $\delta=\pi/8$.
Finally
$\int_0^\infty e^{3u/2}P_\delta\dd u\le3\delta/2$.
Evenness proves (MHW), including every real height in its transform
consequence.

Write $G=G_{5/2}$. For each fixed $Y<5/2$, the vertical gamma
asymptotic, uniform for bounded real parts
\cite[eq.~(5.11.9)]{DLMFGamma}, gives constants $0<g_Y<G_Y<\infty$ with
\[
 g_Y(1+|x|)^{3/2}e^{-\pi|x|/2}
 \le |G(x+iy)|\le
 G_Y(1+|x|)^{3/2}e^{-\pi|x|/2},\qquad |y|\le Y.
 \tag{MHGa}
\]
The two gamma arguments have real parts $(5/2\pm y)/2$ in a fixed
compact positive interval, so their powers sum to $3/2$ and their
exponential factors to $e^{-\pi|x|/2}$. On bounded $x$-intervals,
continuity and gamma nonvanishing extend both bounds.
Thus $G(z)/G(z+iu)$ is uniformly bounded for
$|\Im z|\le1/2$, $|u|\le1$, at every real height.
The logarithmic derivative
\[
 g(z)=G'(z)/G(z)=\tfrac i2
 [\psi((5/2+iz)/2)-\psi((5/2-iz)/2)]
\]
is bounded on every fixed smaller strip by the digamma asymptotic
\cite[eq.~(5.11.2)]{DLMFGamma}: its two $\log|x|$ terms cancel.
Cauchy's estimate on the fixed strip $|\Im z|\le2$ bounds $g'$
on $|\Im z|\le3/2$.

On the central interval $|u|\le1$, (MFE) identifies the integrand
in $\mathcal B_{\alpha,\tau}$ exactly with
$G(z)\mathcal M_\alpha(z+iu)/G(z+iu)$.
Equation (MHW) replaces $\mathcal M_\alpha$ by $\Xi$ at error
at most a height-independent constant divided by $\alpha$.
Set $S(z,u)=G(z)\Xi(z+iu)/G(z+iu)$. Its second $u$-derivative is
\[
 -\frac{G(z)}{G(z+iu)}
       [\Xi''-2g\Xi'+(g^2-g')\Xi](z+iu).
\]
The ordinary theta integral bounds $\Xi,\Xi',\Xi''$ on the entire
closed strip $3/2$ by the finite moments
$\int_\R |v|^j e^{3|v|/2}\Phi(v)\dd v$, $j=0,1,2$.
Hence this derivative is uniformly bounded. Taylor's formula and
the vanishing first moment on the symmetric central interval give
an error at most $C_2\tau$.

For $|u|\ge1$, the hypothesis $8\tau H_\alpha\le1$ implies
$H_\alpha|u|\le u^2/(8\tau)$. Equation (MHG) then gives
\[
 \frac1{\sqrt{4\pi\tau}}\int_{|u|\ge1}
    |E_k(z+iu)|e^{-u^2/(4\tau)}\dd u
 \le16(\alpha+1)^4
      e^{\pi|\Re z|/4+H_\alpha|\Im z|}e^{-1/(16\tau)}.
\]
Multiply by $(3/4)c(k)|G(z)|$ and use (MHGa) on the critical strip.
The remaining height factor
$(1+|x|)^{3/2}e^{-\pi|x|/4}$ has a finite supremum, and
$e^{H_\alpha/2}\le C(\alpha+1)^{1/4}$. This gives the third
term of (MHU). The omitted Gaussian mass multiplying $\Xi(z)$
is uniformly bounded by its theta-integral majorant and
$e^{-1/(4\tau)}$, which is absorbed into the same term.
If $\tau\log\alpha\to0$, $1/(16\tau)$ dominates every fixed
multiple of $\log\alpha$, proving the claimed whole-strip limit.

To retain every multiplicity, let $\gamma$ be a real Xi zero of
exact order $m$ and put $F=-4\Xi/(3G)$. On a sufficiently small
fixed disk around $\gamma$,
$F(\gamma+w)=a_mw^m+O(w^{m+1})$, $a_m\in\R\setminus\{0\}$.
Equations (MFE) and (MHW), and $c(k)=1+O(\alpha^{-1})$, give
$E_k=F+O(\alpha^{-1})$ uniformly on this disk.
In (MHeat), put $z=\gamma+\sqrt\tau\,w$. On $|u|\le d$,
where $d>0$ is fixed and small, this Taylor expansion gives an
error $O(\alpha^{-1})+O(\tau^{(m+1)/2})$ after Gaussian integration,
uniformly on every fixed compact $w$-set. On the remote interval
$|u|\ge d$ we explicitly require $8\tau H_\alpha\le d$,
which holds eventually when $\tau\log\alpha\to0$.
The same calculation as above then bounds the contribution by
$C\alpha^C e^{-d^2/(16\tau)}$; after division by $\tau^{m/2}$
it tends to zero. The polynomial terms omitted outside the
central interval have the same Gaussian tail property.
Consequently, if additionally $\alpha^{-1}=o(\tau^{m/2})$,
\[
 \tau^{-m/2}\mathcal E_{\alpha,\tau}
        (\gamma+\sqrt\tau\,w)\longrightarrow a_mP_m(w),\qquad
 P_m(w)=\sum_{j=0}^{\lfloor m/2\rfloor}
       \frac{(-1)^jm!}{j!(m-2j)!}w^{m-2j}. \tag{MHH}
\]
Convergence is normal on every compact $w$-set.
The Gaussian moments give exactly these coefficients.
The Rodrigues identity
$P_m=(-2)^me^{w^2/4}\partial_w^me^{-w^2/4}$ and $m$ integrations
by parts prove orthogonality to every real polynomial of degree
less than $m$ for the positive weight $e^{-w^2/4}$.
If $P_m$ had fewer than $m$ real sign-changing roots, multiplying
it by the polynomial of those roots would give a nonzero function
of constant sign whose weighted integral is both nonzero and,
by orthogonality, zero. Thus its $m$ roots are distinct and real.
Disjoint conjugation-invariant circles around them and Rouch\'e's
theorem give one simple real zero of the heated approximant in
each circle. A fixed enclosing contour around $\gamma$ has zero
count $m$ by compact convergence, so there are no further zeros
in that fixed neighborhood. The diagonal $(\log\alpha)^{-3}$
satisfies $\alpha^{-1}=o(\tau_\alpha^{m/2})$ for every fixed $m$,
completing the proof without assuming Xi simplicity.
\end{proof}

The heat in (MHeat) acts on the entire factor $E_k$; multiplication
of the physical $\Phi_\alpha$ by $e^{\tau u^2}$ would diverge for
$\tau>0$ because its positive exponential tail is retained.
A hypothetical fixed off-line Xi zero persists in
$\mathcal B_{\alpha,\tau_\alpha}$ by (MHU) and Rouch\'e's theorem.
Hence local multiplicity repair does not remove that zero.
An unbounded parameter sequence on which this particular heated
family has only real zeros in $|\Im z|\le1/2$ would imply RH.
No such sequence is proved here.

Meixner--Pollaczek expansions and Poisson flows already occur in
\cite[\S3.4, Theorems 3.6--3.8 and Proposition 3.9]{Romik2019}.
That work uses the parameter $3/4$, whereas the above monic recurrence
has parameter $1/4$ and includes the Gamma factor in (MFE).
It supplies context, not a premise of the proofs here; no priority claim for the present construction is intended.
Neither positivity in (MK), the real-rooted basis, nor (MD) proves
strip reality for the specific theta combinations $E_k$.
An unbounded sequence of $\alpha$ for which $\mathcal M_\alpha$ has
no nonreal critical-strip zeros would imply RH: on a sufficiently
small contour around any hypothetical off-line Xi zero, (MCV) and
Rouch\'e's theorem would preserve its positive zero count for every
sufficiently large $\alpha$. Such a sequence has not been proved.

\subsection{Tail matching and every fixed logarithmic heat diagonal}
The restriction $\tau\log\alpha\to0$ can be broadened by matching
the actual analytic tails before applying the entire-factor heat.

\begin{theorem}[Whole-strip exhaustion on bounded logarithmic diagonals]
\label{thm:mehler-logarithmic-heat}
For every fixed $d>0$ there are finite constants $C_1(d),C_2(d),C_3$
such that, for $\alpha\ge2$, $\tau>0$, and
$8\tau H_\alpha\le d$, $H_\alpha=\tfrac12\log\alpha$,
\[
 \sup_{|\Im z|\le1/2}|\mathcal B_{\alpha,\tau}(z)-\Xi(z)|
 \le\frac{C_1(d)}{\alpha}+C_2(d)\tau+
               C_3(\alpha+1)^{17/4}e^{-d^2/(16\tau)}.    \tag{TM1}
\]
Consequently the whole-strip convergence in (MHU) holds whenever
$\tau\to0$ and $\tau\log\alpha=O(1)$.
For every fixed $C>0$, the diagonal $\tau=C/\log\alpha$ is
admissible. On each such diagonal, every fixed real Xi zero of
arbitrary order $m$ produces exactly $m$ simple real zeros in
a sufficiently small fixed neighborhood, for all sufficiently
large $\alpha$. The thresholds need not be uniform in that zero.
\end{theorem}
\begin{proof}
Put $a=5/2$, $\delta=\alpha^{-1}$, $r=e^{-2u}$ for $u\ge0$,
and $S(w)=\sum_{n\ge1}e^{-\pi n^2w}$.
The exact theta and pool branches of $\Phi_\alpha$ are
$r^{a/2}\mathfrak A_\delta(r)$ and
$(3/2)\delta r^{a/2}(1+\delta r)^{-a}$, respectively, where
\[
 \begin{split}
 \mathfrak A_\delta(r)&=
 \frac{4(1-\delta^2)^2r}{(r+\delta)^{9/2}}S''(w)
 +\frac{6(1-\delta^2)(r-\delta)}{(r+\delta)^{7/2}}S'(w)
 -\frac{3\delta}{(r+\delta)^{5/2}}S(w),\\
 w&=\frac{1+\delta r}{r+\delta}.
 \end{split}                                                   \tag{TM2}
\]
These follow by substitution in the exact Mehler branches used
in (MEB). On $|r|\le\delta/2$ they are analytic continuations
from positive $r$, and
\[
 \Re w=\frac{\delta(1+|r|^2)+(1+\delta^2)\Re r}
                         {|r+\delta|^2}\ge\frac1{6\delta}.
\]
The full theta moment estimate used in (MHW) gives
\[
 |\mathfrak A_\delta(r)|
 \le C_A\delta^{-1}e^{-\pi/(12\delta)},\qquad
 C_A=2^{11/2}6^{5/2}D_2+
            9\,2^{7/2}6^{3/2}D_1+3\,2^{5/2}6^{1/2},     \tag{TM3}
\]
where $D_0=1$ and $D_j=(4j/e)^j$ for $j>0$.
Indeed, the three rational factors cost
$2^{11/2}\delta^{-7/2}$, $9\,2^{7/2}\delta^{-5/2}$
and $3\,2^{5/2}\delta^{-3/2}$; their respective theta moments
supply $(6\delta)^{5/2}D_2$, $(6\delta)^{3/2}D_1$ and
$(6\delta)^{1/2}D_0$, with the displayed exponential.
All theta indices remain in these bounds.

For a fixed integer $J\ge1$, let $h_j(\delta)$, $j<J$, be the
Taylor coefficients of
$\mathfrak A_\delta(r)+(3/2)\delta(1+\delta r)^{-a}$ at zero.
The theta contribution is bounded by
$C_A\delta^{-1}(2/\delta)^j e^{-\pi/(12\delta)}$;
the pool contribution is exactly
$(3/2)\delta(-\delta)^j(a)_j/j!$.
Thus $h_j=O_J(\delta)$. Here the uniform bound follows from
the finite maximum of $\delta^{-p}e^{-q/\delta}$ on $(0,1/2]$
for every fixed $p,q>0$.
Define
\[
 \begin{aligned}
 \beta_0&=h_0,&
 \beta_j&=h_j-\sum_{\ell<j}\beta_\ell
           \frac{(-1)^{j-\ell}(a+2\ell)_{j-\ell}}{(j-\ell)!},\\
 b_j&=2^{-a-2j}\beta_j,&
 W_{\delta,J}(u)&=\sum_{j<J}b_j\operatorname{sech}^{a+2j}u .
 \end{aligned}                                          \tag{TM4}
\]
The finite triangular system has diagonal one because
$\operatorname{sech}^{a+2j}u
 =2^{a+2j}r^{a/2+j}(1+r)^{-a-2j}$.
It matches exactly the first $J$ tails, with $b_j=O_J(\delta)$.

For every fixed $Y\ge0$ and $J$ with $a+2J>Y$ we claim
\[
 \int_{\mathbb R}e^{Y|u|}
       |\Phi_\alpha(u)-\Phi(u)-W_{\delta,J}(u)|\,du
                                      \le C_{Y,J}\delta. \tag{TM5}
\]
The following decomposition proves the claim over the entire line.
Split the coefficients in (TM4) linearly into theta and pool parts.
The pool analytic remainder, after removal of $r^{a/2}$,
is bounded by $C_J\delta$ on $|r|\le3/4$ and vanishes to order $J$.
Cauchy bounds it by $C'_J\delta r^J$ for $0\le r\le1/2$;
on $[1/2,1]$ its positive-real formulas give $O_J(\delta)$.
The measure is $dr/(2r)$ and the weight $r^{-Y/2}$,
so precisely $a+2J>Y$ makes the first integral finite.

For the theta part use the three consecutive regions
$r\ge\delta$, $\delta/4\le r\le\delta$, and $r\le\delta/4$.
In the first, $0\le u\le H_\alpha$, the mean-value calculation
underlying (MEB), with the additional weight, gives a bound
\[
 \frac{\delta}{2}\int_1^\infty
 v^{Y/2-3/4}e^{-\pi v/2}
 [4B_3v^4+32B_2v^3+36B_1v^2+(9/2)B_0v]\,dv .
\]
The $B_j$ are the finite full-theta derivative bounds from that
calculation; this positive moment is finite for every fixed $Y$.
The matched theta coefficients on this region contribute at most
a fixed power of $\delta^{-1}$ times $e^{-\pi/(12\delta)}$,
which is $O_{Y,J}(\delta)$.
In the second region $w\ge1/(2\delta)$, so all rational factors
cost only fixed powers, whereas the theta moments supply
$e^{-\pi/(2\delta)}$. Its $u$-length is $\log2$.
The true $\Phi$ tail and the matched theta tails satisfy the
same required $O_{Y,J}(\delta)$ conclusion.
In the third region Cauchy on the disk of (TM3) gives
\[
 \left|\mathfrak A_\delta(r)-\sum_{j<J}h_j^\theta r^j\right|
 \le2^{J+1}C_A\delta^{-J-1}e^{-\pi/(12\delta)}r^J.
\]
The matched analytic basis has the same coefficients through $J-1$;
Cauchy on its fixed disk $|r|\le1/2$ bounds its remaining factor
by $C_J\delta^{-J}e^{-\pi/(12\delta)}r^J$.
The weighted remote integral is therefore at most
\[
 \frac{C_J}{a+2J-Y}\delta^{-J-1}e^{-\pi/(12\delta)}
                       (\delta/4)^{(a+2J-Y)/2}
                  =O_{Y,J}(\delta).
\]
The omitted true $\Phi$ tail has its double exponential bound.
There is no gap between these regions; evenness completes (TM5).

Set $G=G_a$ and $F=-4\Xi/(3G)$.
Gamma recurrence yields the polynomial identity
\[
 \Pi_j(z)=\frac{G_{a+2j}(z)}{G_a(z)}
       =\frac{\Gamma(a)}{\Gamma(a+2j)}
                         \prod_{\ell<j}[(a+2\ell)^2+z^2]. \tag{TM6}
\]
The powers $2^{2j}$ and $4^{-j}$ cancel.
If $T_{\delta,J}$ denotes the transform of the remainder in (TM5),
physical integration first on $|\Im z|<a$ gives
\[
 c(k)E_k(z)=F(z)-\tfrac43\sum_{j<J}b_j\Pi_j(z)
                   -\frac{4T_{\delta,J}(z)}{3G_a(z)}.     \tag{TM7}
\]
The remainder is holomorphic for $|\Im z|<a+2J$.
Since $1/G_a$ is entire, every fixed Gamma pole is canceled;
analytic continuation extends (TM7) to that larger strip.

The vertical Gamma asymptotic \cite[5.11.9]{DLMFGamma}, uniformly
for bounded real parts of its arguments, and continuity of $1/G_a$
on bounded regions, imply
$|1/G_a(x+iy)|\le C_Y(1+|x|)^{-3/2}e^{\pi|x|/2}$
for every fixed $Y$, $|y|\le Y$.
The true theta integral bounds $\Xi$ uniformly on every fixed
horizontal strip. The finite polynomials in (TM7) are dominated
by the same exponential envelope. With $c(k)=1+O(\delta)$ we obtain
\[
 |E_k(x+iy)-F(x+iy)|
 \le K_Y\alpha^{-1}(1+|x|)^{-3/2}e^{\pi|x|/2},
           \qquad |y|\le Y,\quad\alpha\ge2.               \tag{TM8}
\]
This is valid also across the fixed Gamma poles.
It is not a bound uniform in growing $Y$.

In (MHeat), use (TM8) on $|u|\le d$, with $Y=d+1/2$.
Multiplication by $G_a(z)$ cancels both the power $3/2$ and
the exponential, giving an error $O_d(\alpha^{-1})$ at every
real height. The center integrand
$S(z,u)=G_a(z)\Xi(z+iu)/G_a(z+iu)$ is uniformly bounded in
a fixed complex $u$-neighborhood of $[-d,d]$.
Indeed, the reciprocal is entire, the two Gamma envelopes have
a bounded ratio when the real coordinate shifts by a bounded
amount, and $\Xi$ is bounded on the corresponding larger strip.
A fixed-radius Cauchy estimate bounds the second $u$ derivative
uniformly. Symmetric Gaussian integration therefore gives
$\Xi(z)+O_d(\tau)$, including its missing central mass.
For $|u|\ge d$, the majorant (MHG) and
$H_\alpha|u|\le u^2/(8\tau)$, valid under $8\tau H_\alpha\le d$,
give the full remote contribution
$C_3(\alpha+1)^{17/4}e^{-d^2/(16\tau)}$.
The omitted Xi mass is bounded by a smaller Gaussian tail.
This proves (TM1).

If $\tau\log\alpha\le B$ eventually, choose, for $B>0$,
a fixed $d=8B+10$. Then the window condition holds and
$d^2/(16B)>17/4$, so the remote term tends to zero.
For $B=0$ the previously proved regime applies.

Finally fix a real zero $\gamma$ of order $m$ and
$\tau=C/\log\alpha$. Use a small Taylor window
$|u|\le\varepsilon$ and a fixed large $D>4C+1$ with
$D^2/(16C)>17/4+1$.
In the small window (TM8) on compact sets and the Taylor
expansion of $F$ give the Hermite scaling (MHH), with errors
$O(\alpha^{-1})+O(\tau^{(m+1)/2})$.
On $\varepsilon\le|u|\le D$, (TM8) bounds $E_k$ uniformly
on the fixed larger compact: the contribution is
$O(e^{-\varepsilon^2/(8\tau)})$.
On $|u|\ge D$ it is
$O(\alpha^{17/4}e^{-D^2/(16\tau)})$.
Both remote terms vanish after division by $\tau^{m/2}$;
so does $\alpha^{-1}/\tau^{m/2}$.
The simple real Hermite roots, conjugation-invariant
Rouch\'e disks and the enclosing count $m$ give exactly the
stated local real zeros, without assuming simplicity of $\Xi$.
\end{proof}

The preceding proof uses a fixed larger window. We next control the
window dependence and remove its restriction on the heat diagonal.

\subsection{Uniform vertical costs and every vanishing heat diagonal}

\begin{theorem}[Unrestricted vanishing Mehler heat]
\label{thm:unrestricted-mehler-heat}
There are absolute finite $C,K>0$ such that for every $\alpha\ge2$,
$Y\ge0$, $x\in\R$ and $|y|\le Y$,
\[
 |c(k)E_k(x+iy)-F(x+iy)|
 \le\frac K\alpha e^{C(Y+1)\log(Y+2)}
              (1+|x|)^{-3/2}e^{\pi|x|/2},
 \qquad F=-\frac{4\Xi}{3G_{5/2}}.                 \tag{UVH1}
\]
The same bound holds for $E_k-F$ after increasing the absolute
constants. For each fixed $T>0$ there is $K_T<\infty$ such that
for every $\alpha\ge2$ and $0<\tau\le T$,
\[
 \sup_{|\Im z|\le1/2}
       |\mathcal B_{\alpha,\tau}(z)-\Xi(z)|
                         \le K_T(\alpha^{-1}+\tau).       \tag{UVH2}
\]
Here $D=\partial_z$ and $\mathcal B_{\alpha,\tau}=-3c(k)G_{5/2}
e^{-\tau D^2}E_k/4$ is the actual entire-factor heat with its
gamma factor restored. It is not physical heat of $\Xi$.
Consequently every diagonal $\tau(\alpha)\to0$ is admissible.
Along such a vanishing diagonal, for a fixed real Xi zero of order $m$, if also
$\alpha^{-1}=o(\tau^{m/2})$, its small fixed neighborhood
contains exactly $m$ simple real zeros of
$\mathcal B_{\alpha,\tau}$ for all sufficiently large $\alpha$.
No threshold uniform in the Xi zero or its multiplicity is claimed.
\end{theorem}
\begin{proof}
Retain $a=5/2$, $\delta=\alpha^{-1}$ and the exact tail matching
(TM2)--(TM7). The additional obligation is a controlled constant
when the matching order grows with the desired strip width.
For $t=r/(1+r)^2$ its analytic inverse is
\[
 r(t)=\frac{1-2t-\sqrt{1-4t}}{2t},\qquad r(0)=0 .
\]
It is analytic on $|t|<1/4$, with nonnegative Taylor
coefficients, as follows from $r=t(1+r)^2$ recursively.
In particular $|r(t)|\le r(|t|)$,
$r(1/16)=7-4\sqrt3<1/8$, and $|r(t)|\le2|t|$ there.
Put $Z^\theta_\delta=\mathfrak A_\delta$ from (TM2) and
$Z^{\rm p}_\delta=(3/2)\delta(1+\delta r)^{-a}$.
The coefficients already defined in (TM4) satisfy exactly
\[
 \sum_{j\ge0}\beta_j^b t^j
       =(1+r(t))^a Z^b_\delta(r(t)),
                         \qquad b\in\{\theta,{\rm p}\}.  \tag{UVH3}
\]
Substitution of $t=r/(1+r)^2$ gives the analytic basis
$r^j(1+r)^{-a-2j}$, whose triangular coefficient matrix
has diagonal one. This proves equality of the coefficients
at every finite order without an infinite matching assumption.

Cauchy on $|t|=1/16$ bounds
$|\beta_j^{\rm p}|\le3\delta16^j$.
Indeed $(3/2)(9/8)^a(15/16)^{-a}<3$.
Cauchy on $|t|=\delta/8$, using (TM3), gives
\[
 |\beta_j^\theta|
 \le2C_A\delta^{-1}e^{-q/\delta}(8/\delta)^j,
 \qquad q=\pi/12 .
\]
For $p>0$ define $M_p(q)=(p/(eq))^p$. Maximizing
$v^pe^{-qv}$ for $v>0$ gives
$\delta^{-p}e^{-q/\delta}\le M_p(q)$. Hence
\[
 |\beta_j|\le\delta\mathfrak B_j,\quad
 \mathfrak B_j=3\,16^j+2C_A8^jM_{j+2}(q),\quad
 |b_j|\le2^{-a-2j}\delta\mathfrak B_j .            \tag{UVH4}
\]
In particular $\log(1+\mathfrak B_j)=O((j+1)\log(j+2))$
with an absolute constant.

We make the remainder estimate (TM5) uniform. Take
$J=\lceil Y/2\rceil+2$ and $g=a+2J-Y\ge a+4$.
Then
\[
 \int_\R e^{Y|u|}
       |\Phi_\alpha-\Phi-W_{\delta,J}|\,\dd u
       \le\delta\mathcal C(Y),\qquad
 \mathcal C(Y)\le e^{C_0(Y+1)\log(Y+2)}              \tag{UVH5}
\]
for an absolute $C_0$. Here is a finite expression and the
estimates proving it, retaining the whole theta series.

For the pole remainder after removing $r^{a/2}$, use the
circle $|r|=1/32$. On it $|t|\le32/961$; (UVH3) and
$|\beta_j^{\rm p}|\le3\delta16^j$ bound the finite sum
by $8\delta$. The analytic remainder vanishes to order $J$,
so on $0\le r\le1/64$ it is at most
$20\delta(32r)^J$. On $1/64\le r\le1$,
$t\le1/4$ bounds it by $2\delta4^J$.
With $u_0=(\log64)/2$ its weighted integral on both
half-lines is bounded by $\delta\mathcal C_{\rm p}(Y)$, where
\[
 \mathcal C_{\rm p}(Y)=4u_0\,4^J64^{Y/2}
                   +40\,32^J64^{-g/2}/g .
\]

For the theta part split the whole positive half-line at
$H=(\log(1/\delta))/2$ and $U=H+\log2$.
The inner mean-value integral displayed in the proof of
(TM5) is at most
$\delta B_*(\lceil Y/2\rceil+4)!$, where
$B_*=2B_3+16B_2+18B_1+(9/4)B_0$;
these are the four fixed derivative constants of (MEB).
Every power in that integral is at most
$v^{\lceil Y/2\rceil+4}$ for $v\ge1$, and $\pi/2>1$,
which proves the factorial estimate directly.
The matched theta terms on $0\le u\le U$ contribute at most
\[
 4C_A2^{J+Y}\delta\,M_{J+Y/2+3}(q).
\]
This follows from
$|W^\theta|\le C_A\delta^{-1}e^{-q/\delta}(2/\delta)^J$,
$U\le2/\delta$ and $e^{YU}=2^Y\delta^{-Y/2}$.

On $H\le u\le U$, the exact argument $w$ is at least
$1/(2\delta)$. The three rational factors in (TM2)
cost respectively $4\delta^{-7/2}$,
$6\delta^{-5/2}$ and $3\delta^{-3/2}$, before
the factor $r^{a/2}$. The complete theta moment bounds
give $|\Phi_\alpha^\theta|\le
C_{\rm tr}\delta^{1/4}e^{-\pi/(4\delta)}$, where
$C_{\rm tr}=4\,2^{5/2}D_2+6\,2^{3/2}D_1+3\sqrt2$.
Its weighted integral on both half-lines is at most
$2C_{\rm tr}(\log2)2^Y\delta M_{Y/2+3/4}(\pi/4)$.

For $u\ge U$, Cauchy on $|r|=\delta/2$, applied to
the matched analytic theta remainder, bounds it by
$6C_A8^J\delta^{-1}e^{-q/\delta}(2r/\delta)^J$.
The circle bound follows from (TM3), (UVH3) and
$|(8/\delta)t|\le64/9<8$ there. Integrating the
remaining $e^{-gu}$ exactly gives
$12C_A4^J2^Y\delta M_{Y/2+3/4}(q)/g$.
Finally the omitted original $\Phi$ tail from $H$ has
weighted integral at most
$\delta(4D_2+6D_1)(\lceil Y/2\rceil+1)!$:
use $|\Phi|\le(4D_2+6D_1)v^{-1/4}e^{-\pi v/2}$
and $1\le\delta v$ on that tail.
Thus (UVH5) holds with the sum of the seven positive
expressions just given, allowing twice the inner factorial
term for evenness. Since $J\le Y/2+3$,
$\log M_p(q)=p(\log p-1-\log q)$ and $n!\le n^n$,
the logarithm of every summand is
$O((Y+1)\log(Y+2))$. Their finite sum has the same bound.
This controls the growing order before any limit is taken.

We also need a uniform reciprocal-gamma estimate.
Binet's formula \cite[eq.~(5.9.10), $\Re w>0$]{DLMFGamma}
uses the holomorphic logarithm of $\Gamma$ continued from the
positive real axis in the right half-plane, the principal inverse
tangent, and gives
\[
 \log\Gamma(w)=(w-\tfrac12)\log w-w
        +\tfrac12\log(2\pi)+
             2\int_0^\infty\frac{\arctan(t/w)}{e^{2\pi t}-1}\dd t.
\]
For $\Re w\ge1$, the remainder has absolute value at most
$1/(12\Re w)$: integrate the bound
$|\arctan(t/w)|\le t/\Re w$.
For $w=\sigma+it$, $\sigma\ge1$, this implies
\[
 |1/\Gamma(w)|\le
 \frac{e^{\sigma+1/12}}{\sqrt{2\pi}}\,
       2^{(\sigma-1/2)/2}
       (1+|t|)^{1/2-\sigma}e^{\pi|t|/2}.           \tag{UVH6}
\]
Use $|w|\ge(1+|t|)/\sqrt2$ and $|t\arg w|\le\pi|t|/2$.
Choose $N_Y=\lceil Y/2\rceil+2$ and apply exact recurrence
to $w_\pm=(a\mp y\pm ix)/2$:
$1/\Gamma(w_\pm)=(w_\pm)_{N_Y}/\Gamma(w_\pm+N_Y)$.
The shifted real parts exceed one, and
$|(w_\pm)_{N_Y}|\le(Y+5)^{N_Y}(1+|x|/2)^{N_Y}$.
The powers $N_Y$ cancel those in (UVH6), giving
\[
 |1/G_a(x+iy)|\le D_\Gamma(Y)
             (1+|x|)^{1-a}e^{\pi|x|/2},\quad
 D_\Gamma(Y)=10e^{a+2N_Y+1}2^{a+2N_Y}(Y+5)^{2N_Y}. \tag{UVH7}
\]
Reciprocal recurrence is valid even at every gamma pole.
$\log D_\Gamma(Y)=O((Y+1)\log(Y+2))$.
The same Binet bound on $|y|\le1/2$ gives
$|G_a(x+iy)|\le K_G(1+|x|)^{a-1}e^{-\pi|x|/2}$,
with a fixed $K_G$.

In (TM7), (UVH5) bounds $T_{\delta,J}$ by
$\delta\mathcal C(Y)$. The polynomials (TM6) satisfy, with
$p_j=2j+a-1$ and $B_j^*=a+2j+Y$,
\[
 \sup_{x\ge0}(1+x)^{a-1}e^{-\pi x/2}
                |\Pi_j(x+iy)|
 \le 2^{p_j+1}(B_j^*+2p_j/\pi)^{p_j}.
\]
Indeed $|\Pi_j|\le(x+B_j^*)^{2j}$ and one maximizes
$x^{p_j}e^{-\pi x/2}$. Equations (UVH4)--(UVH7)
now prove (UVH1). The theta moment bound for $\Xi$
has factorial growth in $Y$, so $F$ has the same envelope.
Since $c(k)^{-1}\le2$ and $|c(k)^{-1}-1|\le\delta/2$,
this also proves the estimate for $E_k-F$.

Define the exact limiting conjugated heat
\[
 \mathcal B_{\infty,\tau}(z)=
 G_a(z)(4\pi\tau)^{-1/2}\int_\R
       \frac{\Xi(z+iu)}{G_a(z+iu)}e^{-u^2/(4\tau)}\dd u .
\]
The reciprocal is entire; the integral is normally convergent.
Use (UVH1) with $Y=1/2+|u|$, and cancel its height envelope
with the fixed-strip bound for $G_a(z)$. On the whole closed
critical strip this bounds
$|\mathcal B_{\alpha,\tau}-\mathcal B_{\infty,\tau}|$
by a constant times
\[
 \frac{\delta}{\sqrt\pi}
 \int_\R e^{-v^2}
       e^{C(2+2\sqrt\tau|v|)\log(3+2\sqrt\tau|v|)}\dd v .
                                                               \tag{UVH8}
\]
For $0<\tau\le T$ the integrand is dominated by its value
with $\tau=T$, which is integrable. The error is $O_T(\delta)$,
without an $\alpha$-dependent remote prefactor.

Set $S(z,u)=G_a(z)\Xi(z+iu)/G_a(z+iu)$.
It is analytic in a fixed complex neighborhood of $u=0$,
uniformly bounded there at every real height by (UVH7)
and the theta moment bound. Cauchy gives
$|S(z,u)+S(z,-u)-2\Xi(z)|\le K u^2$ for $|u|\le1$.
For $|u|\ge1$ the same expression is bounded by
$K u^2e^{C(2+|u|)\log(3+|u|)}$.
Gaussian integration and $u=2\sqrt\tau v$ therefore give
$|\mathcal B_{\infty,\tau}-\Xi|\le K_T\tau$,
again uniformly on the closed strip. This proves (UVH2).

Finally fix a real zero $\gamma$ of order $m$.
The Taylor expansion of $F$ at $\gamma$, the all-$u$ Gaussian
majorant just proved, and the compact version of (UVH1)
give the Hermite scaling with errors
$O(\alpha^{-1})+O(\tau^{(m+1)/2})$; the remote Gaussian
contribution is smaller than every power of $\tau$.
Under $\alpha^{-1}=o(\tau^{m/2})$ the scaled limit has
$m$ simple real Hermite roots. Conjugation-invariant
Rouch\'e disks each contain one actual zero counted with
multiplicity; each is real and simple.
Whole-strip convergence on an enclosing fixed circle,
chosen to contain no other Xi zero, gives total count $m$.
This proves the local assertion for arbitrary multiplicity.
\end{proof}

A fixed off-line Xi zero would still persist under (UVH2).
The larger class of heat diagonals is available for an arithmetic
or phase transfer; global strip reality along a cofinal sequence
remains a separate obligation.

\section{A polynomial cutoff for every high-strip theta zero}
We now prove a uniform tail theorem for the actual entire factor
in (MFE), retaining the complete theta series and the subtraction.
Write $D=\partial_z$ and use exactly
$\mathcal E_{\alpha,\tau}=e^{-\tau D^2}E_k$ from (MHeat).
For $\alpha\ge3$ put
\[
 \begin{aligned}
 \sigma&=\arctan k,& d&=\frac{\pi\alpha}{\alpha^2+1},
 &\ell&=\frac{\pi(\alpha^2-1)}{\alpha^2+1},\\
 H&=\tfrac12\log\alpha,&J&=\lceil32/d\rceil,\\
 A&=\frac{\sqrt\alpha}{\sqrt2c(k)(\alpha^2+1)^{1/4}},
 &B&=\frac{\alpha^{1/4}}{\sqrt2c(k)} .
 \end{aligned}                                          \tag{PHT1}
\]
Thus $0<d<1$ and $\ell>0$. For finite $T>0$ define
\[
 \begin{split}
 W_C&=128\cdot18\cdot256^3\ell^2J^2d^{-3},\\
 W_R&=8\ell^2\sqrt J\,d^{-3/2},\\
 W_B&=\max\left\{1,\frac{2Hd}{\sigma},
 \sqrt{\frac{8\ell}{\sigma}
 \left[\max\{0,\log(256\pi B/(Ad))\}+\frac d4+\frac4d\right]}\right\},\\
 W&=2\max\{2,W_C,W_R,W_B,\,
 65536\ell\sqrt T\,d^{-3/2},\,16384\ell TH/d,\,4\ell\sqrt{T/d}\},\\
 X(\alpha,T,Y)&=2T\sigma+
                    \frac{(W+1+16\ell Y/d)^2}{2\ell}.
 \end{split}                                                   \tag{PHT2}
\]
These are elementary constants independent of any Xi zero.

\begin{theorem}[All high-strip zeros with a polynomial cutoff]
\label{thm:polynomial-theta-tail}
For $\alpha\ge3$, finite $T>0,Y\ge0$, and every $0\le\tau\le T$,
every zero of $\mathcal E_{\alpha,\tau}$ in
\[
       |\Re z|\ge X(\alpha,T,Y),\qquad |\Im z|\le Y
\]
is real and simple. For $Y<5/2$ the same assertion, with the same
multiplicities, holds for
$\mathcal B_{\alpha,\tau}=-(3/4)c(k)G_{5/2}\mathcal E_{\alpha,\tau}$.
For fixed $T,Y$,
\[
                X(\alpha,T,Y)=O_{T,Y}(\alpha^{10}).       \tag{PHT3}
\]
The statement includes every height above this cutoff.
\end{theorem}

\begin{proof}
First we derive the exact sampling identity used in the estimates.
Let $f_z(t)$ be the even entire solution
\[
 f_z''(t)+(t^2+2z)f_z(t)=0,\qquad f_z(0)=1,\quad f_z'(0)=0,
 \quad \nu_n=n+\tfrac12,\quad \epsilon_n=(-1)^{n(n+1)/2}.
\]
Then
\[
 E_k(z)=Ae^{\sigma z}\sum_{n\ge0}\epsilon_ne^{-d\nu_n^2}
                   f_z(\sqrt\ell\,\nu_n)-B\cos(Hz).        \tag{PHT4}
\]
Here is a derivation including the phases and convergence.
With $s=1/4-iz/2$, initially $0<\Re s<1/2$, set
$v=e^{2u}$ and $w=(\alpha v+1)/(v+\alpha)$.
Poisson summation on the alternating Gaussian gives
\[
 \theta(w)=(w+i)^{-1/2}
       \sum_{\nu\in\mathbb Z+1/2}e^{-\pi\nu^2/(w+i)}.
\]
Use the right-half-plane square root and put
$a=\alpha+i$, $b=1+i\alpha$, $\lambda=b/a=e^{2i\sigma}$.
Then
\[
 w+i=\frac{a(v+\lambda)}{v+\alpha},\qquad
 \frac1{w+i}=\frac1a+\frac{\alpha^2-1}{a^2(v+\lambda)}.
\]
The theta part of $\mathcal R_\alpha$ is consequently
\[
 \tfrac12\sqrt\alpha\,a^{-1/2}v^{1/4}
 \sum_\nu e^{-\pi\nu^2/a}(v+\lambda)^{-1/2}
       e^{-\pi\nu^2(\alpha^2-1)/(a^2(v+\lambda))}.
\]
Rotate its Mellin contour to $v=\lambda r$, $r>0$.
The singularity $-\lambda$ is outside the sector; the two arcs
vanish by $0<\Re s<1/2$. On this ray the coefficient of $\pi\nu^2$
in the exponent has real part $\alpha/(\alpha^2+1)$,
independent of $r$. Thus the full Gaussian sum is integrable
before and after rotation. Expanding the last exponential and
integrating by the beta integral gives
\[
 \int_0^\infty v^{s-1}(v+\lambda)^{-1/2}
 e^{-\pi\nu^2(\alpha^2-1)/(a^2(v+\lambda))}\,dv
 =\lambda^{s-1/2}\frac{\Gamma(s)\Gamma(1/2-s)}{\sqrt\pi}
                  M(1/2-s;1/2;i\ell\nu^2).
\]
The subtraction terms of $\mathcal R_\alpha$ give both
$-\alpha^s$ and $-\alpha^{1/2-s}$.
By (MFE), $\widehat{\mathcal R_\alpha}=c(k)G_{1/2}E_k$.
Division and folding the half-integer sum give
\[
 E_k(z)=\frac{\sqrt\alpha\,a^{-1/2}\lambda^{s-1/2}}{\sqrt2c(k)}
 \sum_{n\ge0}e^{-\pi\nu_n^2/a}
        M(1/2-s;1/2;i\ell\nu_n^2)
       -\frac{\alpha^s+\alpha^{1/2-s}}{2\sqrt2c(k)}.
\]
Direct substitution in Kummer's differential equation yields
$f_z(t)=e^{-it^2/2}M(1/4+iz/2;1/2;it^2)$.
The phase of
$e^{-\pi\nu_n^2/a+i\ell\nu_n^2/2}$ is
$e^{i\pi/8}\epsilon_n$, and the prefactor has the opposite
$e^{-i\pi/8}$ phase and modulus $Ae^{\sigma z}$.
This proves (PHT4) on the initial strip.
For real $t\ge0$, the energy of
$f''+(1+t^2)f=(1-2z)f$, divided by $1+t^2$, gives
$|f_z(t)|\le\exp(\tfrac12|1-2z|\operatorname{asinh}t)$.
This is a polynomial bound in $t$ on every compact $z$-set.
The Gaussian sum is therefore normally entire, and the identity
theorem proves (PHT4) everywhere.

Put
\[
 T_d(w)=\sum_{n\ge0}\epsilon_ne^{-d\nu_n^2}\cos(\nu_nw),
 \qquad a_m=(-1)^me^{-2dm}.
\]
The theta Fourier series and Jacobi product
\cite[20.2.2, 20.5.2]{DLMF}, at nome $q=ie^{-d}$ and argument $w/2$,
give the normally convergent product
\[
 T_d(w)=e^{-d/4}\cos(w/2)
       \prod_{m\ge1}(1-a_m)(1+a_me^{iw})(1+a_me^{-iw}).    \tag{PHT5}
\]
In $|\Im w|<2d$ its only zeros are the simple points
$w_j=(2j+1)\pi$. On $|\Im w|\le d$,
\[
 |T_d(w)|\ge c_d\operatorname{dist}(w,\pi+2\pi\mathbb Z),
             \qquad c_d=\pi^{-1}e^{-d/4-4/d}.             \tag{PHT6}
\]
Indeed, the three factor families are bounded below by
$[\prod_{j\ge1}(1-e^{-jd})]^2$; the possible inverse even
product is at least one. Monotone integral comparison bounds the
logarithm of the single product below by
$\int_0^\infty\log(1-e^{-dt})dt>-2/d$.
Also $|\cos(w/2)|\ge\operatorname{dist}(w,\pi+2\pi\mathbb Z)/\pi$
by the sine-concavity bound and the complex modulus identity.

We need relative, rather than absolute, control of all corrections.
For $|\Im w|\le d/2$ and distance at least $d/32$ from these zeros,
\[
                         |T_d'(w)/T_d(w)|\le128/d.       \tag{PHT7}
\]
The $m$th product derivative is
$-(-1)^m\sin w/[\cosh(2dm)+(-1)^m\cos w]$.
For each parity let $r\in[0,\pi]$ be the distance of $\Re w$
from its minimizing real angle, and put $v=\Im w,t=2dm$.
Since $|v|\le t/4$,
\[
 \Re(\cosh t-\cos r\cosh v)\ge(t^2+r^2)/5,\qquad
 |\sin w|\le2(r+|v|).
\]
For each parity the sums of
$r/[(2dm)^2+r^2]$ and $|v|/(2dm)^2$ are at most
$\pi/(4d)$ and $1/(4d)$, respectively.
Both parity families cost less than $21/d$; the tangent term
costs at most $32\pi/d<101/d$. This proves (PHT7).
Cauchy on a circle of radius $d/256$, integrating this logarithmic
derivative along its radii, now gives
\[
 |T_d^{(p)}(w)|\le2p!(256/d)^p|T_d(w)|,\quad p\ge0,
 \quad |\Im w|\le d/4,\quad
 \operatorname{dist}(w,\pi+2\pi\mathbb Z)\ge d/16.         \tag{PHT8}
\]
The circle remains at distance at least $15d/256>d/32$ and
within $|\Im w|<d/2$. The logarithmic ratio is at most $1/2$.
No exponential factor $c_d^{-1}$ enters (PHT8).

For $z=x+iy$, $x>0$, oscillator energy gives
\[
 |f_z(t)|\le e^{|y|\operatorname{asinh}(t/\sqrt{2x})}
                         \le e^{ht},\qquad h=|y|/\sqrt{2x}.
\]
With $\lambda_z=\sqrt{2z}$, $w=\sqrt{2\ell z}$, define
$Vg(t)=\lambda_z^{-1}\int_0^t\sin(\lambda_z(t-s))s^2g(s)\,ds$.
Exact Volterra iteration gives
\[
 f_z=\sum_{j=0}^{J-1}(-V)^j\cos(\lambda_z t)+(-V)^Jf_z,\qquad
 |V^Jf_z(t)|\le\frac{t^{3J}e^{ht}}{3^JJ!|\lambda_z|^J}.
\]
After summing every half-integer, the remainder is bounded by
\[
 |R_J(z)|\le3d^{-1/2}e^{b_0^2/d}
       \left[\frac{2\ell^2\sqrt J}{|w|d^{3/2}}\right]^J,
                  \qquad b_0=h\sqrt\ell.                \tag{PHT9}
\]
To check the bound, use
$b_0\nu\le d\nu^2/4+b_0^2/d$,
$\max_{\nu\ge0}\nu^{3J}e^{-d\nu^2/2}=(3J/(de))^{3J/2}$,
$\sum_{n\ge0}e^{-d\nu_n^2/4}\le3/\sqrt d$ and
$J!\ge(J/e)^J$. For $|\Im w|\le d/4$, $\Re w\ge2$,
$b_0\le d/2$; if $|w|\ge W_R$ this is at most
$6d^{-1/2}4^{-J}$. Dividing by the lower bound $c_dd/16$
in (PHT6) gives at most
$96\pi d^{-3/2}e^{d/4+4/d}4^{-J}
 <400d^{-3/2}e^{-40/d}<1/16$.

Each explicit correction retains the signed theta sum.
Define $F_0(u)=\cos u$ and
$F_j''+F_j=-u^2F_{j-1}$, with zero initial data for $j\ge1$.
Then $(-V)^j\cos(\lambda_z t)=\lambda_z^{-4j}F_j(\lambda_z t)$,
and $F_j=P_j\cos+Q_j\sin$, with degree at most $3j$,
$P_j$ even and $Q_j$ odd.
Uniqueness of the polynomial cosine/sine representation and the even
ODE prove the parities, so only the corresponding derivatives
of $T_d$ occur. For an explicit construction put
\[
 \begin{split}
 A_0(u)&=1,\\
 B_j(u)&=-\sum_{r\ge0}\frac{(-1)^r}{(2i)^{r+1}}
                         D_u^r[u^2A_{j-1}(u)],\\
 A_j(u)&=\int_0^uB_j(s)\,ds+i\Re B_j(0),
                 \qquad A_j=P_j-iQ_j .
 \end{split}                                                   \tag{PHT10}
\]
The derivative sum is finite. It satisfies
$A_j''+2iA_j'=-u^2A_{j-1}$ and both real initial conditions.
For $\|P\|_a=\sum_p|P_p|p!a^p$, $a\ge1$, differentiation
costs $a^{-r}$, its displayed inverse at most one, multiplication
by $u^2$ at most $9j^2a^2$, and integration with the constant
at most $2a$. Induction yields
$\|A_j\|_a\le(18a^3)^j(j!)^2$.
Combining this with (PHT8), at $a=256|w|/d$, bounds the
summed correction $\Psi_j$ by
\[
 |\Psi_j(w)|\le4|T_d(w)|
       \left[\frac{18\cdot256^3\ell^2}{|w|d^3}\right]^j(j!)^2.
                                                               \tag{PHT11}
\]
For example the first correction is exactly
$-\ell^2T_d'''/(6w)+\ell^2T_d''/(4w^2)-\ell^2T_d'/(4w^3)$.
For $j<J$, $(j!)^2\le J^{2j}$, and $|w|\ge W_C$
makes the full correction sum at most $4|T_d|/127<|T_d|/16$.
Only a finite, growing order is used.

Let $Q_\alpha(z)=A^{-1}e^{-\sigma z}E_k(z)$.
The subtraction in (PHT4) is retained.
For $w=u+iv$, $u\ge1$, $|v|\le d/4$, $z=w^2/(2\ell)$,
we have $\Re z\ge u^2/(4\ell)$ and $|\Im z|\le ud/(4\ell)$.
If $u\ge W_B$,
\[
             (B/A)e^{-\sigma\Re z+H|\Im z|}\le c_dd/256.
\]
Together with (PHT9)--(PHT11) this proves
\[
 |Q_\alpha(w^2/(2\ell))-T_d(w)|<3|T_d(w)|/16             \tag{PHT12}
\]
on the domain of (PHT8), with $\Re w\ge\max(2,W_C,W_R,W_B)$.

The entire Gaussian heat and its derivative series are admissible
by (MHG). Exact conjugation gives
\[
 \mathcal E_{\alpha,\tau}(z)
 =Ae^{\sigma z-\tau\sigma^2}
                   Q_{\alpha,\tau}(z-2\tau\sigma),
        \qquad Q_{\alpha,\tau}=e^{-\tau D^2}Q_\alpha.     \tag{PHT13}
\]
To control its remote integral, use the oscillator estimate on every
vertical height, before taking absolute values of the indices:
\[
 |Q_\alpha(x+i\eta)|\le3d^{-1/2}e^{\ell\eta^2/(4dx)}
                   +(B/A)e^{-\sigma x+H|\eta|},\quad x>0. \tag{PHT14}
\]
This follows by completing the square in the full Gaussian sum.

Take $|\Im w|\le d/8$, distance from the grid at least $d/8$,
$\Re w=u\ge W$, and put $z=w^2/(2\ell)$, $R=ud/(1024\ell)$.
The disk $|\zeta-z|\le R$ maps into the domain of (PHT12):
$|\sqrt{2\ell\zeta}-w|\le d/256$, with real coordinate at least $W/2$.
Equations (PHT7),(PHT12) give $|Q_\alpha(\zeta)|\le3|T_d(w)|$;
Cauchy on the half disk gives $|Q_\alpha''|\le24|T_d(w)|/R^2$.
The central symmetric Gaussian integral has error at most
$48\tau|T_d(w)|/R^2$.
The listed terms in $W$ guarantee
\[
             R^2\ge4096T/d,\qquad R\ge16TH,\qquad
                         x\ge4\ell T/d.                \tag{PHT15}
\]
The central ratio is therefore less than $1/64$.
For $|r|\ge R/2$, apply (PHT14) at $\eta=y+r$.
Using $(y+r)^2\le2y^2+2r^2$ and $x\ge4\ell T/d$, the full
theta contribution is at most
$12d^{-1/2}e^{-R^2/(64\tau)}$.
Relative to $c_dd/8$ this is at most
$400d^{-3/2}e^{-60/d}<1/64$.
For the subtraction, $R\ge16TH$ implies
$H|r|\le r^2/(8\tau)$ on this region; its complete tail is also
less than $|T_d(w)|/64$.
The omitted constant $Q_\alpha(z)$ times the Gaussian tail mass
is less than $|T_d(w)|/64$, since
$|Q_\alpha(z)|\le5|T_d(w)|/4$ and the tail mass is
at most $2e^{-R^2/(32\tau)}$.
Thus the entire heat error is less than $|T_d(w)|/16$,
uniformly for $0\le\tau\le T$, and
\[
 |Q_{\alpha,\tau}(w^2/(2\ell))-T_d(w)|<|T_d(w)|/4         \tag{PHT16}
\]
on this smaller domain. At $\tau=0$ the assertion follows directly.

For each $w_j\ge W+d/8$, Rouch\'e on $|w-w_j|=d/8$ gives
exactly one zero of $Q_{\alpha,\tau}(w^2/(2\ell))$.
Real type and conjugation invariance make it real and simple.
The quadratic map has nonzero derivative there.
Outside these disjoint disks, (PHT16) gives nonvanishing.
If $\Re z\ge X$, $|\Im z|\le Y$, set
$w=\sqrt{2\ell(z-2\tau\sigma)}$.
Then $\Re w\ge W+1+16\ell Y/d$ and
$|\Im w|=\ell|\Im z|/\Re w\le d/16$.
Every zero is consequently in one of the counted disks;
its center exceeds $W+d/8$, and its unique zero is real and simple.
Equation (PHT13) transfers the conclusion to $\mathcal E$.
Evenness handles the negative half-line.
The factor $G_{5/2}$ is holomorphic and nonzero for $|\Im z|<5/2$.
Finally $d\asymp\alpha^{-1}$, $\ell\to\pi$, $J\le33/d$,
and $W_C=O(d^{-5})$ dominates every other listed term for fixed
$T,Y$. Squaring in (PHT2) proves (PHT3).
\end{proof}

This theorem supplies a full tail exclusion for the actual family,
with growing-order error control. It does not exclude any zero
in $|\Re z|<X(\alpha,T,1/2)$.
Every fixed hypothetical off-line Xi zero persists in that region
along the convergent diagonals of (MHU). A global theta-specific
control of this growing middle region is still required.

\section{Conclusions and open questions}\label{sec:conclusions}

The exact index theorem detects every distinct off-line mirror pair,
including an infinite collection. The local smoothed theorem shows
that this detection survives every fixed finite moment prefix on
every bounded open interval. Its moment-price estimate explains
why this local density is not a uniform bound for increasingly
constrained detectors. The complete compact realization identifies
all spectral values, and the fractional family preserves the full
arithmetic expression. Its persistent index and vanishing trace
norm give a precise distinction between absolute smallness and
sign exclusion.

The adaptive completion estimate pays both infinite endpoints and
the actual selected multiplicity. Its source-growth lower bound
holds for the stated adaptive input class and every original heat.
The fixed-band theorem further gives a cumulative, multiplicity-sensitive
energy lower bound for one fixed compact input at every zero height.
Its numerator is quantitatively nonvanishing, and the joint integer
completion transfers this energy to the finite corrected gamma head.
The primitive-sector estimate pays the squarefree remainder through
\(H=X^{1/3-o(1)}\), with the stated input-norm dependence.
These are limsup and transfer assertions. Positivity at almost all modulation
heights, on a favorable subsequence of compressions, or on a
finite-dimensional family does not provide an upper bound for
that entire class.

Three sufficient global estimates remain open:
\begin{enumerate}
\item For the reflected packet \(Q_h\), prove
\[
 \forall\epsilon>0\ \exists C_\epsilon,Y_\epsilon\
 \forall y\ge Y_\epsilon:\quad
 |Q_h(y)|\le C_\epsilon e^{\epsilon y}.
\]
Lemma~\ref{lem:strip} then excludes every off-line zero.
\item For the full ordinary mean-zero space in (SHD1), prove
\[
 \exists t_j\to\infty,\quad \exists\epsilon_j\to0,\quad
 d_{t_j,1/2}\le\sigma_{t_j}e^{\epsilon_jt_j}.
\]
The endpoint detector theorem excludes every positive displacement
from this cofinal \emph{full-vector} bound.
\item Prove an arithmetic upper bound that contradicts
Theorem~\ref{thm:adaptive-completion-radius} on an admissible
input family with a selected value \(X^{-o(1)}\).
An eventual subpolynomial source bound would suffice. A bound
only on an arbitrary favorable subsequence would not contradict
the proved limsup lower. Alternatively, for the fixed input of
Theorem~\ref{thm:fixed-band-cumulative}, prove
\[
 \forall\epsilon>0\ \exists C_\epsilon,Y_\epsilon\
 \forall Y\ge Y_\epsilon:\quad
               \mathcal E_h(Y)\le C_\epsilon Y^\epsilon .
\]
The same theorem would then exclude every positive displacement.
The exact decomposition (GC18) preserves the signed main and
the high primitive complement; neither required upper bound
is established by paying their completion errors.
\end{enumerate}
Alternatively, nonnegativity of the complete Weil form, or realness
of the entire compact arithmetic spectrum, would suffice. Neither
condition follows from the positive ground vector, the factorial
block inequality, or selected positive compressions.

For the approximation results, an unbounded sequence of globally
real critical-strip spectra would exclude an off-line Xi zero by
Rouch\'e's theorem and the stated uniform convergence. Local
multiplicity repair and a real high tail do not establish this
sequence. The certified nonreal approximant zeros describe the
limits of particular constructions and are not zeta counterexamples.

\paragraph{Proof and reproducibility scope.}
The analytic statements are accompanied by complete proofs with
their explicit hypotheses. The numerical statements identify exact geometry,
outward-rounded intervals and analytic infinite-tail bounds in
their proofs. No finite computation is used to conclude anything
about all zeta heights. The source projections and finite gamma
heads are used only in their stated weighted domains. In particular,
the two pole moments do not confer unrestricted higher moments
on a separate gamma head. External endorsement and proof-assistant
formalization are not claimed.

\appendix
\section{Comparison with finite-compression and sextic estimates}\label{app:comparison}
The OpenAI preprints \cite{OpenAI78,OpenAI1112} assert
zero-free half-planes with respective boundaries $7/8$ and $11/12$.
The estimates from these preprints enter only the conditional
calculations in this appendix; none is a premise of the main theorems.
The second manuscript's Proposition 3.1, on PDF p.~11, asserts, for
fixed $0<\vartheta\le1/10$ and every $\epsilon>0$,
\[
 \sum_{0<N_Eu\le H}|A_u(D)|^2
 \ll_{\nu,S,I,\vartheta,\epsilon}\|W\|_{C^k(I)}^2D^{1+\epsilon}H,
 \quad H=D^{1+\vartheta},
\]
where $E=\mathbb Q(\omega)$, $\omega=e^{2\pi i/3}$,
$\nu$ is a fixed finite-order Hecke character, $S$ is fixed and contains
the primes above $2,3$ and its conductor, $I\Subset(0,\infty)$,
$W\in C_c^\infty(I)$, and $k=k(\vartheta,\epsilon)$.
The columns are
\[
 A_u(D)=\sum_{(\mathfrak n,S)=1}
   \mu_E(\mathfrak n)\nu(\mathfrak n)\chi_{\mathfrak n}(u)
   W(N_E\mathfrak n/D).
\]
Primary generators identify ideals with elements. The sextic character
has zero extension at nonunits. In particular
$\chi_{\mathfrak n}(p^6)=\mathbf1_{p\nmid\mathfrak n}$, not an
unrestricted constant one. The extraction on pp.~11--12 gives
\[
 |A_1(D)|^2\ll D^{11/6+5\vartheta/6+\epsilon}
                      +D^{5/3-\vartheta/3}.
\]
Indeed there are $\asymp H^{1/6}/\log H$ admissible prime witnesses,
each differing from $A_1$ by $O_W(D/H^{1/6})$; averaging their squared
values and using the asserted energy gives the displayed bound.
For each requested exponent loss one may then choose $\vartheta$
small, producing $11/12+\epsilon$. This extraction does not approach
$1/2$ by iteration of its parameter alone. Its constants and derivative
order depend on the fixed parameters as stated in the source.

The unpublished model-generated manuscript \cite{ClaudeFull}, Theorems A--C and
Theorem 5.8, asserts asymptotic critical-line and simplicity proportions,
with optimized value approximately $0.6725$, through finite Weil
compressions and their first two traces. The condensed text
\cite{ClaudeShort} uses the same reflection pair
$\{\rho,1-\overline\rho\}$ with a consistent $2\pi$ rescaling.
The two texts describe one mechanism. A public mathematical
communication is available in \cite{AlpogeFurman}; its authors
report the provenance of the model-generated proof. The proportions
are background and are not used as premises of our proofs.

The two arithmetic inputs are different. Expanding the OpenAI energy
gives products of $\mu_E\nu$ and sextic characters. The rational prime
measure entering the Claude form and the packet calculation is
\[
 \sum_{n\ge2}\frac{\Lambda(n)}{\sqrt n}\,\delta_{\log n},
 \qquad
 \Lambda(p^m)=\log p,\quad\Lambda(n)=0\text{ otherwise}.
\]
Their equality cannot be obtained by a fixed rescaling: a rational prime
$p\equiv2\bmod3$ is inert in $E$ and has ideal norm $p^2$, whereas a
split prime has two ideals of norm $p$. The Euler factors give
$\zeta_E(s)=\zeta(s)L(s,\chi_{-3})$ on $\Re s>1$, but this identity
does not turn a twisted ideal-Mobius square into a centered rational
prime-power correlation. Lemma~\ref{lem:prime} makes the latter target explicit.
Any transfer must retain the different Euler factors, all prime powers,
archimedean weights and uniform errors. No such transfer has been
established here.


\begin{thebibliography}{99}
\bibitem{ElkiesHadamard} N. D. Elkies, \emph{Hadamard's product formula},
Harvard Math 259 course note, course index dated 5 March 2007,
with the recorded 10 March correction to Jensen's formula (3).
Theorem (1), p.~1: finite-order entire functions and the degree
of the exponential polynomial factor.
\url{https://people.math.harvard.edu/~elkies/M259.06/prod.pdf}.
\bibitem{SectorReducing} D. A. Cardon, T. Forg\'acs, A. Piotrowski,
E. Sorensen, and J. C. White, \emph{On Zero-Sector Reducing Operators},
arXiv:1802.02641v1, 7 February 2018.
Theorem 2.5 and Corollary 2.6, printed p.~5 / PDF p.~5.
The fixed-sector polynomial hypothesis and normal polynomial
approximation for the entire version are both used explicitly above.
\url{https://arxiv.org/abs/1802.02641v1}.

\bibitem{FlintDocs} python-flint 0.9.0 documentation,
\emph{acb---complex numbers} and \emph{arb---real numbers},
entries \texttt{gamma\_upper}, \texttt{integral}, \texttt{arg},
\texttt{upper}, \texttt{abs\_upper}, and interval constructors,
accessed 7 October 2026.
\url{https://python-flint.readthedocs.io/en/latest/acb.html};
\url{https://python-flint.readthedocs.io/en/latest/arb.html}.
\bibitem{DLMF} NIST Digital Library of Mathematical Functions,
\S\S5.5, 5.12, 20.2, 20.5, 25.2, 25.4, 25.10,
in particular eqs.~(5.5.3), (5.12.1), (5.12.3), (20.2.2), (20.5.2),
and (20.5.9), accessed 7--8 October 2026.
\url{https://dlmf.nist.gov/5.5.E3};
\url{https://dlmf.nist.gov/5.12.E1};
\url{https://dlmf.nist.gov/5.12.E3};
\url{https://dlmf.nist.gov/20.2}; \url{https://dlmf.nist.gov/20.5};
\url{https://dlmf.nist.gov/20.2.E2};
\url{https://dlmf.nist.gov/20.5.E2};
\url{https://dlmf.nist.gov/20.5.E9};
\url{https://dlmf.nist.gov/25.2}; \url{https://dlmf.nist.gov/25.4};
\url{https://dlmf.nist.gov/25.10}.
\bibitem{DLMFGamma} NIST Digital Library of Mathematical Functions,
\S5.2, eqs.~(5.4.19), (5.7.6), (5.9.10), (5.9.16),
(5.11.1), (5.11.2), (5.11.9),
and \S5.11(ii),
accessed 7 October 2026.
\url{https://dlmf.nist.gov/5.2};
\url{https://dlmf.nist.gov/5.4.E19};
\url{https://dlmf.nist.gov/5.7.E6};
\url{https://dlmf.nist.gov/5.9.E10};
\url{https://dlmf.nist.gov/5.9.E16};
\url{https://dlmf.nist.gov/5.11.E1};
\url{https://dlmf.nist.gov/5.11.E2};
\url{https://dlmf.nist.gov/5.11.E9}; \url{https://dlmf.nist.gov/5.11#ii}.
\bibitem{TeschlQM2009} G. Teschl,
\emph{Mathematical Methods in Quantum Mechanics: With Applications
to Schr\"odinger Operators}, American Mathematical Society,
Graduate Studies in Mathematics 99 (2009).
Author-posted version dated 12 February 2009, accessed 7 October 2026;
Theorem 2.13, pp.~71--72; Theorem 4.12(i), p.~119;
Lemma 6.17, pp.~145--146; eq.~(6.28), p.~145;
Lemma 6.18, p.~146; Theorem 6.19, pp.~146--147.
\url{https://www.mat.univie.ac.at/~gerald/ftp/book-schroe/schroe.pdf}.
\bibitem{Romik2019} D. Romik,
\emph{Orthogonal polynomial expansions for the Riemann xi function},
arXiv:1902.06330v3, 5 May 2019.
Context only: \S3.4, Theorems 3.6--3.8 and Proposition 3.9,
printed pp.~26--29 (PDF pp.~32--35).
\url{https://arxiv.org/abs/1902.06330v3}.
\bibitem{FordZF} K. Ford, \emph{Zero-free regions for the Riemann zeta
function}, in \emph{Number Theory for the Millennium}, vol.~II
(2002), 25--56; corrected arXiv:1910.08205v5, 18 February 2025.
Theorem 1 and its growth hypothesis (1.2), PDF p.~1;
the correction to the old Cheng--Graham input is on PDF p.~3.
\url{https://arxiv.org/abs/1910.08205v5}.
\bibitem{FordGrowth} K. Ford, \emph{Vinogradov's integral and bounds for
the Riemann zeta function}, Proc. London Math. Soc. 85 (2002), 565--633,
Theorem 1; reliable primary HTML version arXiv:1910.08209v1,
18 October 2019 (manuscript dated 25 October 2001).
\url{https://arxiv.org/html/1910.08209v1}.
\bibitem{PolymathHeat} D. H. J. Polymath,
\emph{Effective approximation of heat flow evolution of the Riemann
$\xi$ function, and a new upper bound for the de Bruijn--Newman constant},
Research in the Mathematical Sciences 6, article 31 (2019),
DOI 10.1007/s40687-019-0193-1.
Used version: arXiv:1904.12438v2, 4 August 2019, printed pp.~4 and 6,
definitions (6)--(10), Theorem 1.3 and eqs.~(13)--(24).
The factor-two intermediate error on printed p.~57 is repaired in
the direct application above; Proposition 9.1 and its enlarged strip
are not premises of that proof.
\url{https://arxiv.org/pdf/1904.12438v2};
\url{https://doi.org/10.1007/s40687-019-0193-1}.
\bibitem{GaoZhao} P. Gao and L. Zhao,
\emph{Moments and One level density of sextic Hecke $L$-functions
of $\mathbb Q(\omega)$}, arXiv:2201.01885v2, 24 July 2022;
Functiones et Approximatio 70(1) (2024), 7--28,
DOI 10.7169/facm/2057.
\S2.1, physical PDF p.~3: $E$-primary generators,
multiplicativity and sextic reciprocity.
Only these algebraic facts are imported here;
Lemma 2.9 is not used as a large-sieve premise.
\url{https://arxiv.org/pdf/2201.01885v2}.
\bibitem{OpenAI78} OpenAI,
\emph{The Quasi-Riemann Hypothesis: A Zero-Free Half-Plane
\(\Re s>7/8\)}, OpenAI Math Release preprint, 30 September 2026.
\href{https://github.com/openai/math/blob/fd4aeeb2ee4fc729c18d98444fed42fd0529eeeb/preprints/The-Quasi-Riemann-Hypothesis-September-30-2026/paper.pdf}{Version-specific manuscript}.
\bibitem{OpenAI1112} OpenAI,
\emph{The Quasi-Riemann Hypothesis}, OpenAI Math Release preprint,
5 October 2026; written with human assistance.
Proposition 3.1, pp.~11--12, is used only conditionally.
\href{https://github.com/openai/math/blob/fd4aeeb2ee4fc729c18d98444fed42fd0529eeeb/preprints/The-Quasi-Riemann-Hypothesis-October-5-2026/paper2.pdf}{Version-specific manuscript}.
\bibitem{ClaudeFull} Claude,
\emph{More Than Two Thirds of the Zeros of the Riemann Zeta Function
Lie on the Critical Line}, unpublished model-generated manuscript,
10 August 2026, 35 pages. Theorems A--C and Theorem 5.8.
Its public mathematical communication is \cite{AlpogeFurman};
these are not interchangeable versions of a proof.
\bibitem{ClaudeShort} Claude,
condensed version of the preceding unpublished manuscript,
5 pages; reflection and Fourier conventions on pp.~1--2.
Cited only for the comparison in Appendix~\ref{app:comparison}.

\bibitem{AlpogeFurman} L. Alp\"oge and R. Furman,
\emph{More than two thirds of the zeta zeros are simple and on the
critical line}, arXiv:2608.13637v2, 19 August 2026, 21 pages.
Theorem A; provenance in Section 1.
\url{https://arxiv.org/abs/2608.13637v2}.
\bibitem{Milne} J. S. Milne, \emph{Class Field Theory},
version 4.03, 6 August 2020; V, Theorem 3.23, printed p.~164;
VII, Appendix A, Theorem A.3, printed p.~226;
VIII, \S\S5.2--5.3, printed pp.~244--245.
\url{https://www.jmilne.org/math/CourseNotes/CFT.pdf}.
\bibitem{ConnesConsani} A. Connes and C. Consani,
\emph{Spectral triples and zeta-cycles},
Enseign. Math. 69 (2023), 93--148;
Proposition 2.1 and Corollary 2.4.
\url{https://doi.org/10.4171/LEM/1049}.
\bibitem{BombieriWeil} E. Bombieri,
\emph{Remarks on Weil's quadratic functional in the theory of
prime numbers, I}, Rend. Lincei Mat. Appl. (9) 11 (2000), no.~3,
183--233 (2001). Primary bibliographic record and indexed abstract;
the full source proof is not used as a premise here.
\url{https://www.bdim.eu/item?id=RLIN_2000_9_11_3_183_0}.
\bibitem{OperatorSource} L. van Hemert,
\emph{Fixed-Parameter Spectral Approximation and Smoothed Relative
Traces for Semilocal Zeta Operators}, preprint, version 1.0.0,
25 July 2026.
Version-specific paper source:
\url{https://github.com/LeonardSEO/semilocal-zeta-operator-theory/blob/v1.0.0/paper/main.tex}.
Source and reproducibility release:
\url{https://github.com/LeonardSEO/semilocal-zeta-operator-theory/releases/tag/v1.0.0}.
The form and conventions used here are reproduced in the text.
\bibitem{Kedlaya} K. S. Kedlaya, \emph{Notes on analytic number theory},
Chapter 9, Lemma 9.4 and Remark 9.7, accessed 7 October 2026.
\url{https://kskedlaya.org/ant/chap-von-mangoldt.html}.
\bibitem{Elkies} N. D. Elkies, \emph{Math 259: Introduction to Analytic
Number Theory. The product formula for $\xi(s)$ and $\zeta(s)$; vertical
distribution of zeros}, Harvard course notes (2002), pp.~1--2 and 5.
\url{https://people.math.harvard.edu/~elkies/M259.02/zeta2.pdf}.
\bibitem{Petrow} I. Petrow, \emph{Introduction to Analytic Number Theory},
ETH Z\"urich course notes, 2018 course-directory copy;
based on notes by P. Elbaz-Vincent and P. Michel.
Theorem 7.7, printed p.~58: Hadamard factorization for an entire
function of order at most one nonzero at the origin.
\url{https://metaphor.ethz.ch/x/2018/hs/401-3100-68L/IANTnotes2018.pdf}.
\bibitem{Conrey} J. B. Conrey, \emph{More than two fifths of the zeros
of the Riemann zeta function are on the critical line},
J. reine angew. Math. 399 (1989), 1--26.
Eqs.~(8)--(12), printed p.~3; Theorem 1, printed p.~4.
Only the unconditional positive critical-line proportion is used.
\url{https://aimath.org/~kaur/publications/24.pdf}.
\bibitem{PacketSource} L. van Hemert,
\emph{Reflected-Packet Weil Kernels: Zero Expansions, Growth Criteria,
and Infinite Oscillation}, preprint, version 1.0.0, 25 July 2026.
Version-specific paper source:
\url{https://github.com/LeonardSEO/reflected-packet-growth-criteria/blob/v1.0.0/paper/main.tex}.
Source and reproducibility release:
\url{https://github.com/LeonardSEO/reflected-packet-growth-criteria/releases/tag/v1.0.0}.
\bibitem{OscillationSource} LeonardSEO,
\emph{Reflected-Packet Interactions in the Semilocal Weil Form:
An Exact Zero Expansion and Unconditional Infinite Oscillation},
preprint, version 1.0.0, 23 July 2026; author identifier as on
the version-specific title page:
\url{https://github.com/LeonardSEO/semilocal-reflected-packet-oscillation/blob/v1.0.0/paper/main.tex}.
Source and reproducibility release:
\url{https://github.com/LeonardSEO/semilocal-reflected-packet-oscillation/releases/tag/v1.0.0}.
\bibitem{XiSource} L. van Hemert,
\emph{Finite Positivity Certificates and a Cross-Endpoint Obstruction
for a Riemann Xi-Defect Kernel}, preprint, 20 July 2026;
archival release v0.2.1.
Version-specific paper source:
\url{https://github.com/LeonardSEO/certified-riemann-xi-positivity/blob/v0.2.1/paper/main.tex}.
Archival release:
\url{https://github.com/LeonardSEO/certified-riemann-xi-positivity/releases/tag/v0.2.1}.
Code and interval data:
\url{https://github.com/LeonardSEO/certified-riemann-xi-positivity/tree/v0.2.1}.
\bibitem{KernelSource} L. van Hemert,
\emph{Smoothed Positive Kernels for Finite Zeta Spectral Triples
and a Scalar Criterion for the Riemann Hypothesis},
preprint, version 1.0.0, 23 July 2026.
Version-specific paper source:
\url{https://github.com/LeonardSEO/smoothed-zeta-spectral-kernels/blob/v1.0.0/paper/main.tex}.
Source and reproducibility release:
\url{https://github.com/LeonardSEO/smoothed-zeta-spectral-kernels/releases/tag/v1.0.0}.
\bibitem{LeeLeong} E. S. Lee and N. Leong,
\emph{New explicit bounds for Mertens function and the reciprocal
of the Riemann zeta-function}, arXiv:2208.06141v5,
9 September 2026. Theorem 2.2, eq.~(23), with a fixed
admissible parameter \(W\), its stated threshold, and a reduced
positive exponential constant.
\url{https://arxiv.org/abs/2208.06141v5}.
\bibitem{Suzuki2026} M. Suzuki,
\emph{Weil's quadratic form via the screw function},
arXiv:2606.09096v3, 23 September 2026, 33 pages.
Operator-theoretic context; no theorem is imported as a premise.
\url{https://arxiv.org/abs/2606.09096v3}.
\end{thebibliography}
\end{document}
